\documentclass{article}
\usepackage[paperwidth=16cm,paperheight=14cm,margin=0cm]{geometry}
\usepackage[utf8]{inputenc}\usepackage[T1]{fontenc}
\usepackage{mathpazo}\usepackage{amsmath,amssymb}
\usepackage[most]{tcolorbox}
\pagestyle{empty}\setlength{\parindent}{0pt}

\definecolor{cnom}{HTML}{1F6FA8}
\definecolor{csta}{HTML}{C77F0A}
\definecolor{cfon}{HTML}{6A4C93}
\definecolor{cvec}{HTML}{C0503A}
\definecolor{cari}{HTML}{2E7D6B}
\definecolor{cpro}{HTML}{8A5A2B}
\begin{document}
\begin{tcolorbox}[enhanced,width=16cm,height=14cm,colback=white,colframe=cnom,boxrule=2pt,arc=0mm,
left=6mm,right=6mm,top=0mm,bottom=4mm,sharp corners,
overlay={\fill[cnom] (frame.north west) rectangle ([yshift=-1.25cm]frame.north east);
\node[anchor=west,text=white,font=\sffamily\bfseries\LARGE] at ([xshift=5mm,yshift=-0.62cm]frame.north west) {Question 1};
\node[anchor=east,text=white,font=\sffamily\itshape\large] at ([xshift=-5mm,yshift=-0.62cm]frame.north east) {Nombres et calculs};}]
\vspace*{1.75cm}
\begin{minipage}[t][6.9cm][c]{\linewidth}\LARGE\centering
\begin{minipage}{0.95\linewidth}\raggedright\linespread{1.25}\selectfont Écrire $\sqrt{50}$ sous la forme $a\sqrt{b}$, avec $b$ le plus petit possible.\end{minipage}
\end{minipage}

\vfill
\begin{tcolorbox}[enhanced,nobeforeafter,width=7.1cm,height=2.1cm,valign=center,colback=cnom!6,colframe=cnom,arc=2.5mm,boxrule=0.9pt,left=3mm]{\Large\bfseries\color{cnom}a)}\quad\begin{minipage}[c]{5.3cm}\raggedright\Large $5\sqrt{2}$\end{minipage}\end{tcolorbox}\hfill\begin{tcolorbox}[enhanced,nobeforeafter,width=7.1cm,height=2.1cm,valign=center,colback=cnom!6,colframe=cnom,arc=2.5mm,boxrule=0.9pt,left=3mm]{\Large\bfseries\color{cnom}b)}\quad\begin{minipage}[c]{5.3cm}\raggedright\Large $25\sqrt{2}$\end{minipage}\end{tcolorbox}\par\vspace{2.5mm}
\begin{tcolorbox}[enhanced,nobeforeafter,width=7.1cm,height=2.1cm,valign=center,colback=cnom!6,colframe=cnom,arc=2.5mm,boxrule=0.9pt,left=3mm]{\Large\bfseries\color{cnom}c)}\quad\begin{minipage}[c]{5.3cm}\raggedright\Large $5\sqrt{4}$\end{minipage}\end{tcolorbox}\hfill\begin{tcolorbox}[enhanced,nobeforeafter,width=7.1cm,height=2.1cm,valign=center,colback=cnom!6,colframe=cnom,arc=2.5mm,boxrule=0.9pt,left=3mm]{\Large\bfseries\color{cnom}d)}\quad\begin{minipage}[c]{5.3cm}\raggedright\Large $2\sqrt{5}$\end{minipage}\end{tcolorbox}
\end{tcolorbox}
\newpage
\begin{tcolorbox}[enhanced,width=16cm,height=14cm,colback=white,colframe=cnom,boxrule=2pt,arc=0mm,
left=6mm,right=6mm,top=0mm,bottom=4mm,sharp corners,
overlay={\fill[cnom] (frame.north west) rectangle ([yshift=-1.25cm]frame.north east);
\node[anchor=west,text=white,font=\sffamily\bfseries\LARGE] at ([xshift=5mm,yshift=-0.62cm]frame.north west) {Question 2};
\node[anchor=east,text=white,font=\sffamily\itshape\large] at ([xshift=-5mm,yshift=-0.62cm]frame.north east) {Nombres et calculs};}]
\vspace*{1.75cm}
\begin{minipage}[t][6.9cm][c]{\linewidth}\LARGE\centering
\begin{minipage}{0.95\linewidth}\raggedright\linespread{1.25}\selectfont Écrire $\sqrt{12}$ sous la forme $a\sqrt{b}$, avec $b$ le plus petit possible.\end{minipage}
\end{minipage}

\vfill
\begin{tcolorbox}[enhanced,nobeforeafter,width=7.1cm,height=2.1cm,valign=center,colback=cnom!6,colframe=cnom,arc=2.5mm,boxrule=0.9pt,left=3mm]{\Large\bfseries\color{cnom}a)}\quad\begin{minipage}[c]{5.3cm}\raggedright\Large $4\sqrt{3}$\end{minipage}\end{tcolorbox}\hfill\begin{tcolorbox}[enhanced,nobeforeafter,width=7.1cm,height=2.1cm,valign=center,colback=cnom!6,colframe=cnom,arc=2.5mm,boxrule=0.9pt,left=3mm]{\Large\bfseries\color{cnom}b)}\quad\begin{minipage}[c]{5.3cm}\raggedright\Large $2\sqrt{3}$\end{minipage}\end{tcolorbox}\par\vspace{2.5mm}
\begin{tcolorbox}[enhanced,nobeforeafter,width=7.1cm,height=2.1cm,valign=center,colback=cnom!6,colframe=cnom,arc=2.5mm,boxrule=0.9pt,left=3mm]{\Large\bfseries\color{cnom}c)}\quad\begin{minipage}[c]{5.3cm}\raggedright\Large $3\sqrt{2}$\end{minipage}\end{tcolorbox}\hfill\begin{tcolorbox}[enhanced,nobeforeafter,width=7.1cm,height=2.1cm,valign=center,colback=cnom!6,colframe=cnom,arc=2.5mm,boxrule=0.9pt,left=3mm]{\Large\bfseries\color{cnom}d)}\quad\begin{minipage}[c]{5.3cm}\raggedright\Large $2\sqrt{6}$\end{minipage}\end{tcolorbox}
\end{tcolorbox}
\newpage
\begin{tcolorbox}[enhanced,width=16cm,height=14cm,colback=white,colframe=cnom,boxrule=2pt,arc=0mm,
left=6mm,right=6mm,top=0mm,bottom=4mm,sharp corners,
overlay={\fill[cnom] (frame.north west) rectangle ([yshift=-1.25cm]frame.north east);
\node[anchor=west,text=white,font=\sffamily\bfseries\LARGE] at ([xshift=5mm,yshift=-0.62cm]frame.north west) {Question 3};
\node[anchor=east,text=white,font=\sffamily\itshape\large] at ([xshift=-5mm,yshift=-0.62cm]frame.north east) {Nombres et calculs};}]
\vspace*{1.75cm}
\begin{minipage}[t][6.9cm][c]{\linewidth}\LARGE\centering
\begin{minipage}{0.95\linewidth}\raggedright\linespread{1.25}\selectfont Que vaut $\sqrt{12}+\sqrt{27}$~?\end{minipage}
\end{minipage}

\vfill
\begin{tcolorbox}[enhanced,nobeforeafter,width=7.1cm,height=2.1cm,valign=center,colback=cnom!6,colframe=cnom,arc=2.5mm,boxrule=0.9pt,left=3mm]{\Large\bfseries\color{cnom}a)}\quad\begin{minipage}[c]{5.3cm}\raggedright\Large $\sqrt{39}$\end{minipage}\end{tcolorbox}\hfill\begin{tcolorbox}[enhanced,nobeforeafter,width=7.1cm,height=2.1cm,valign=center,colback=cnom!6,colframe=cnom,arc=2.5mm,boxrule=0.9pt,left=3mm]{\Large\bfseries\color{cnom}b)}\quad\begin{minipage}[c]{5.3cm}\raggedright\Large $6\sqrt{3}$\end{minipage}\end{tcolorbox}\par\vspace{2.5mm}
\begin{tcolorbox}[enhanced,nobeforeafter,width=7.1cm,height=2.1cm,valign=center,colback=cnom!6,colframe=cnom,arc=2.5mm,boxrule=0.9pt,left=3mm]{\Large\bfseries\color{cnom}c)}\quad\begin{minipage}[c]{5.3cm}\raggedright\Large $5\sqrt{6}$\end{minipage}\end{tcolorbox}\hfill\begin{tcolorbox}[enhanced,nobeforeafter,width=7.1cm,height=2.1cm,valign=center,colback=cnom!6,colframe=cnom,arc=2.5mm,boxrule=0.9pt,left=3mm]{\Large\bfseries\color{cnom}d)}\quad\begin{minipage}[c]{5.3cm}\raggedright\Large $5\sqrt{3}$\end{minipage}\end{tcolorbox}
\end{tcolorbox}
\newpage
\begin{tcolorbox}[enhanced,width=16cm,height=14cm,colback=white,colframe=cnom,boxrule=2pt,arc=0mm,
left=6mm,right=6mm,top=0mm,bottom=4mm,sharp corners,
overlay={\fill[cnom] (frame.north west) rectangle ([yshift=-1.25cm]frame.north east);
\node[anchor=west,text=white,font=\sffamily\bfseries\LARGE] at ([xshift=5mm,yshift=-0.62cm]frame.north west) {Question 4};
\node[anchor=east,text=white,font=\sffamily\itshape\large] at ([xshift=-5mm,yshift=-0.62cm]frame.north east) {Nombres et calculs};}]
\vspace*{1.75cm}
\begin{minipage}[t][6.9cm][c]{\linewidth}\LARGE\centering
\begin{minipage}{0.95\linewidth}\raggedright\linespread{1.25}\selectfont Que vaut $\sqrt{8}+\sqrt{18}$~?\end{minipage}
\end{minipage}

\vfill
\begin{tcolorbox}[enhanced,nobeforeafter,width=7.1cm,height=2.1cm,valign=center,colback=cnom!6,colframe=cnom,arc=2.5mm,boxrule=0.9pt,left=3mm]{\Large\bfseries\color{cnom}a)}\quad\begin{minipage}[c]{5.3cm}\raggedright\Large $\sqrt{26}$\end{minipage}\end{tcolorbox}\hfill\begin{tcolorbox}[enhanced,nobeforeafter,width=7.1cm,height=2.1cm,valign=center,colback=cnom!6,colframe=cnom,arc=2.5mm,boxrule=0.9pt,left=3mm]{\Large\bfseries\color{cnom}b)}\quad\begin{minipage}[c]{5.3cm}\raggedright\Large $5\sqrt{4}$\end{minipage}\end{tcolorbox}\par\vspace{2.5mm}
\begin{tcolorbox}[enhanced,nobeforeafter,width=7.1cm,height=2.1cm,valign=center,colback=cnom!6,colframe=cnom,arc=2.5mm,boxrule=0.9pt,left=3mm]{\Large\bfseries\color{cnom}c)}\quad\begin{minipage}[c]{5.3cm}\raggedright\Large $6\sqrt{2}$\end{minipage}\end{tcolorbox}\hfill\begin{tcolorbox}[enhanced,nobeforeafter,width=7.1cm,height=2.1cm,valign=center,colback=cnom!6,colframe=cnom,arc=2.5mm,boxrule=0.9pt,left=3mm]{\Large\bfseries\color{cnom}d)}\quad\begin{minipage}[c]{5.3cm}\raggedright\Large $5\sqrt{2}$\end{minipage}\end{tcolorbox}
\end{tcolorbox}
\newpage
\begin{tcolorbox}[enhanced,width=16cm,height=14cm,colback=white,colframe=cnom,boxrule=2pt,arc=0mm,
left=6mm,right=6mm,top=0mm,bottom=4mm,sharp corners,
overlay={\fill[cnom] (frame.north west) rectangle ([yshift=-1.25cm]frame.north east);
\node[anchor=west,text=white,font=\sffamily\bfseries\LARGE] at ([xshift=5mm,yshift=-0.62cm]frame.north west) {Question 5};
\node[anchor=east,text=white,font=\sffamily\itshape\large] at ([xshift=-5mm,yshift=-0.62cm]frame.north east) {Nombres et calculs};}]
\vspace*{1.75cm}
\begin{minipage}[t][6.9cm][c]{\linewidth}\LARGE\centering
\begin{minipage}{0.95\linewidth}\raggedright\linespread{1.25}\selectfont Développer $(x+3)^2$.\end{minipage}
\end{minipage}

\vfill
\begin{tcolorbox}[enhanced,nobeforeafter,width=7.1cm,height=2.1cm,valign=center,colback=cnom!6,colframe=cnom,arc=2.5mm,boxrule=0.9pt,left=3mm]{\Large\bfseries\color{cnom}a)}\quad\begin{minipage}[c]{5.3cm}\raggedright\Large $x^2+6x-9$\end{minipage}\end{tcolorbox}\hfill\begin{tcolorbox}[enhanced,nobeforeafter,width=7.1cm,height=2.1cm,valign=center,colback=cnom!6,colframe=cnom,arc=2.5mm,boxrule=0.9pt,left=3mm]{\Large\bfseries\color{cnom}b)}\quad\begin{minipage}[c]{5.3cm}\raggedright\Large $x^2+9$\end{minipage}\end{tcolorbox}\par\vspace{2.5mm}
\begin{tcolorbox}[enhanced,nobeforeafter,width=7.1cm,height=2.1cm,valign=center,colback=cnom!6,colframe=cnom,arc=2.5mm,boxrule=0.9pt,left=3mm]{\Large\bfseries\color{cnom}c)}\quad\begin{minipage}[c]{5.3cm}\raggedright\Large $x^2+3x+9$\end{minipage}\end{tcolorbox}\hfill\begin{tcolorbox}[enhanced,nobeforeafter,width=7.1cm,height=2.1cm,valign=center,colback=cnom!6,colframe=cnom,arc=2.5mm,boxrule=0.9pt,left=3mm]{\Large\bfseries\color{cnom}d)}\quad\begin{minipage}[c]{5.3cm}\raggedright\Large $x^2+6x+9$\end{minipage}\end{tcolorbox}
\end{tcolorbox}
\newpage
\begin{tcolorbox}[enhanced,width=16cm,height=14cm,colback=white,colframe=cnom,boxrule=2pt,arc=0mm,
left=6mm,right=6mm,top=0mm,bottom=4mm,sharp corners,
overlay={\fill[cnom] (frame.north west) rectangle ([yshift=-1.25cm]frame.north east);
\node[anchor=west,text=white,font=\sffamily\bfseries\LARGE] at ([xshift=5mm,yshift=-0.62cm]frame.north west) {Question 6};
\node[anchor=east,text=white,font=\sffamily\itshape\large] at ([xshift=-5mm,yshift=-0.62cm]frame.north east) {Nombres et calculs};}]
\vspace*{1.75cm}
\begin{minipage}[t][6.9cm][c]{\linewidth}\LARGE\centering
\begin{minipage}{0.95\linewidth}\raggedright\linespread{1.25}\selectfont Développer $(x-5)^2$.\end{minipage}
\end{minipage}

\vfill
\begin{tcolorbox}[enhanced,nobeforeafter,width=7.1cm,height=2.1cm,valign=center,colback=cnom!6,colframe=cnom,arc=2.5mm,boxrule=0.9pt,left=3mm]{\Large\bfseries\color{cnom}a)}\quad\begin{minipage}[c]{5.3cm}\raggedright\Large $x^2-5x+25$\end{minipage}\end{tcolorbox}\hfill\begin{tcolorbox}[enhanced,nobeforeafter,width=7.1cm,height=2.1cm,valign=center,colback=cnom!6,colframe=cnom,arc=2.5mm,boxrule=0.9pt,left=3mm]{\Large\bfseries\color{cnom}b)}\quad\begin{minipage}[c]{5.3cm}\raggedright\Large $x^2+25$\end{minipage}\end{tcolorbox}\par\vspace{2.5mm}
\begin{tcolorbox}[enhanced,nobeforeafter,width=7.1cm,height=2.1cm,valign=center,colback=cnom!6,colframe=cnom,arc=2.5mm,boxrule=0.9pt,left=3mm]{\Large\bfseries\color{cnom}c)}\quad\begin{minipage}[c]{5.3cm}\raggedright\Large $x^2-10x+25$\end{minipage}\end{tcolorbox}\hfill\begin{tcolorbox}[enhanced,nobeforeafter,width=7.1cm,height=2.1cm,valign=center,colback=cnom!6,colframe=cnom,arc=2.5mm,boxrule=0.9pt,left=3mm]{\Large\bfseries\color{cnom}d)}\quad\begin{minipage}[c]{5.3cm}\raggedright\Large $x^2-10x-25$\end{minipage}\end{tcolorbox}
\end{tcolorbox}
\newpage
\begin{tcolorbox}[enhanced,width=16cm,height=14cm,colback=white,colframe=cnom,boxrule=2pt,arc=0mm,
left=6mm,right=6mm,top=0mm,bottom=4mm,sharp corners,
overlay={\fill[cnom] (frame.north west) rectangle ([yshift=-1.25cm]frame.north east);
\node[anchor=west,text=white,font=\sffamily\bfseries\LARGE] at ([xshift=5mm,yshift=-0.62cm]frame.north west) {Question 7};
\node[anchor=east,text=white,font=\sffamily\itshape\large] at ([xshift=-5mm,yshift=-0.62cm]frame.north east) {Nombres et calculs};}]
\vspace*{1.75cm}
\begin{minipage}[t][6.9cm][c]{\linewidth}\LARGE\centering
\begin{minipage}{0.95\linewidth}\raggedright\linespread{1.25}\selectfont Factoriser $x^2-9$.\end{minipage}
\end{minipage}

\vfill
\begin{tcolorbox}[enhanced,nobeforeafter,width=7.1cm,height=2.1cm,valign=center,colback=cnom!6,colframe=cnom,arc=2.5mm,boxrule=0.9pt,left=3mm]{\Large\bfseries\color{cnom}a)}\quad\begin{minipage}[c]{5.3cm}\raggedright\Large $(x+3)^2$\end{minipage}\end{tcolorbox}\hfill\begin{tcolorbox}[enhanced,nobeforeafter,width=7.1cm,height=2.1cm,valign=center,colback=cnom!6,colframe=cnom,arc=2.5mm,boxrule=0.9pt,left=3mm]{\Large\bfseries\color{cnom}b)}\quad\begin{minipage}[c]{5.3cm}\raggedright\Large $(x-9)(x+9)$\end{minipage}\end{tcolorbox}\par\vspace{2.5mm}
\begin{tcolorbox}[enhanced,nobeforeafter,width=7.1cm,height=2.1cm,valign=center,colback=cnom!6,colframe=cnom,arc=2.5mm,boxrule=0.9pt,left=3mm]{\Large\bfseries\color{cnom}c)}\quad\begin{minipage}[c]{5.3cm}\raggedright\Large $(x-3)^2$\end{minipage}\end{tcolorbox}\hfill\begin{tcolorbox}[enhanced,nobeforeafter,width=7.1cm,height=2.1cm,valign=center,colback=cnom!6,colframe=cnom,arc=2.5mm,boxrule=0.9pt,left=3mm]{\Large\bfseries\color{cnom}d)}\quad\begin{minipage}[c]{5.3cm}\raggedright\Large $(x-3)(x+3)$\end{minipage}\end{tcolorbox}
\end{tcolorbox}
\newpage
\begin{tcolorbox}[enhanced,width=16cm,height=14cm,colback=white,colframe=cnom,boxrule=2pt,arc=0mm,
left=6mm,right=6mm,top=0mm,bottom=4mm,sharp corners,
overlay={\fill[cnom] (frame.north west) rectangle ([yshift=-1.25cm]frame.north east);
\node[anchor=west,text=white,font=\sffamily\bfseries\LARGE] at ([xshift=5mm,yshift=-0.62cm]frame.north west) {Question 8};
\node[anchor=east,text=white,font=\sffamily\itshape\large] at ([xshift=-5mm,yshift=-0.62cm]frame.north east) {Nombres et calculs};}]
\vspace*{1.75cm}
\begin{minipage}[t][6.9cm][c]{\linewidth}\LARGE\centering
\begin{minipage}{0.95\linewidth}\raggedright\linespread{1.25}\selectfont Factoriser $x^2-25$.\end{minipage}
\end{minipage}

\vfill
\begin{tcolorbox}[enhanced,nobeforeafter,width=7.1cm,height=2.1cm,valign=center,colback=cnom!6,colframe=cnom,arc=2.5mm,boxrule=0.9pt,left=3mm]{\Large\bfseries\color{cnom}a)}\quad\begin{minipage}[c]{5.3cm}\raggedright\Large $(x+5)^2$\end{minipage}\end{tcolorbox}\hfill\begin{tcolorbox}[enhanced,nobeforeafter,width=7.1cm,height=2.1cm,valign=center,colback=cnom!6,colframe=cnom,arc=2.5mm,boxrule=0.9pt,left=3mm]{\Large\bfseries\color{cnom}b)}\quad\begin{minipage}[c]{5.3cm}\raggedright\Large $(x-5)^2$\end{minipage}\end{tcolorbox}\par\vspace{2.5mm}
\begin{tcolorbox}[enhanced,nobeforeafter,width=7.1cm,height=2.1cm,valign=center,colback=cnom!6,colframe=cnom,arc=2.5mm,boxrule=0.9pt,left=3mm]{\Large\bfseries\color{cnom}c)}\quad\begin{minipage}[c]{5.3cm}\raggedright\Large $(x-25)(x+25)$\end{minipage}\end{tcolorbox}\hfill\begin{tcolorbox}[enhanced,nobeforeafter,width=7.1cm,height=2.1cm,valign=center,colback=cnom!6,colframe=cnom,arc=2.5mm,boxrule=0.9pt,left=3mm]{\Large\bfseries\color{cnom}d)}\quad\begin{minipage}[c]{5.3cm}\raggedright\Large $(x-5)(x+5)$\end{minipage}\end{tcolorbox}
\end{tcolorbox}
\newpage
\begin{tcolorbox}[enhanced,width=16cm,height=14cm,colback=white,colframe=cnom,boxrule=2pt,arc=0mm,
left=6mm,right=6mm,top=0mm,bottom=4mm,sharp corners,
overlay={\fill[cnom] (frame.north west) rectangle ([yshift=-1.25cm]frame.north east);
\node[anchor=west,text=white,font=\sffamily\bfseries\LARGE] at ([xshift=5mm,yshift=-0.62cm]frame.north west) {Question 9};
\node[anchor=east,text=white,font=\sffamily\itshape\large] at ([xshift=-5mm,yshift=-0.62cm]frame.north east) {Nombres et calculs};}]
\vspace*{1.75cm}
\begin{minipage}[t][6.9cm][c]{\linewidth}\LARGE\centering
\begin{minipage}{0.95\linewidth}\raggedright\linespread{1.25}\selectfont Factoriser $x^2+6x+9$.\end{minipage}
\end{minipage}

\vfill
\begin{tcolorbox}[enhanced,nobeforeafter,width=7.1cm,height=2.1cm,valign=center,colback=cnom!6,colframe=cnom,arc=2.5mm,boxrule=0.9pt,left=3mm]{\Large\bfseries\color{cnom}a)}\quad\begin{minipage}[c]{5.3cm}\raggedright\Large $(x-3)^2$\end{minipage}\end{tcolorbox}\hfill\begin{tcolorbox}[enhanced,nobeforeafter,width=7.1cm,height=2.1cm,valign=center,colback=cnom!6,colframe=cnom,arc=2.5mm,boxrule=0.9pt,left=3mm]{\Large\bfseries\color{cnom}b)}\quad\begin{minipage}[c]{5.3cm}\raggedright\Large $(x-3)(x+3)$\end{minipage}\end{tcolorbox}\par\vspace{2.5mm}
\begin{tcolorbox}[enhanced,nobeforeafter,width=7.1cm,height=2.1cm,valign=center,colback=cnom!6,colframe=cnom,arc=2.5mm,boxrule=0.9pt,left=3mm]{\Large\bfseries\color{cnom}c)}\quad\begin{minipage}[c]{5.3cm}\raggedright\Large $(x+3)^2$\end{minipage}\end{tcolorbox}\hfill\begin{tcolorbox}[enhanced,nobeforeafter,width=7.1cm,height=2.1cm,valign=center,colback=cnom!6,colframe=cnom,arc=2.5mm,boxrule=0.9pt,left=3mm]{\Large\bfseries\color{cnom}d)}\quad\begin{minipage}[c]{5.3cm}\raggedright\Large $(x+6)^2$\end{minipage}\end{tcolorbox}
\end{tcolorbox}
\newpage
\begin{tcolorbox}[enhanced,width=16cm,height=14cm,colback=white,colframe=cnom,boxrule=2pt,arc=0mm,
left=6mm,right=6mm,top=0mm,bottom=4mm,sharp corners,
overlay={\fill[cnom] (frame.north west) rectangle ([yshift=-1.25cm]frame.north east);
\node[anchor=west,text=white,font=\sffamily\bfseries\LARGE] at ([xshift=5mm,yshift=-0.62cm]frame.north west) {Question 10};
\node[anchor=east,text=white,font=\sffamily\itshape\large] at ([xshift=-5mm,yshift=-0.62cm]frame.north east) {Nombres et calculs};}]
\vspace*{1.75cm}
\begin{minipage}[t][6.9cm][c]{\linewidth}\LARGE\centering
\begin{minipage}{0.95\linewidth}\raggedright\linespread{1.25}\selectfont Factoriser $x^2-8x+16$.\end{minipage}
\end{minipage}

\vfill
\begin{tcolorbox}[enhanced,nobeforeafter,width=7.1cm,height=2.1cm,valign=center,colback=cnom!6,colframe=cnom,arc=2.5mm,boxrule=0.9pt,left=3mm]{\Large\bfseries\color{cnom}a)}\quad\begin{minipage}[c]{5.3cm}\raggedright\Large $(x+4)(x-4)$\end{minipage}\end{tcolorbox}\hfill\begin{tcolorbox}[enhanced,nobeforeafter,width=7.1cm,height=2.1cm,valign=center,colback=cnom!6,colframe=cnom,arc=2.5mm,boxrule=0.9pt,left=3mm]{\Large\bfseries\color{cnom}b)}\quad\begin{minipage}[c]{5.3cm}\raggedright\Large $(x-8)^2$\end{minipage}\end{tcolorbox}\par\vspace{2.5mm}
\begin{tcolorbox}[enhanced,nobeforeafter,width=7.1cm,height=2.1cm,valign=center,colback=cnom!6,colframe=cnom,arc=2.5mm,boxrule=0.9pt,left=3mm]{\Large\bfseries\color{cnom}c)}\quad\begin{minipage}[c]{5.3cm}\raggedright\Large $(x+4)^2$\end{minipage}\end{tcolorbox}\hfill\begin{tcolorbox}[enhanced,nobeforeafter,width=7.1cm,height=2.1cm,valign=center,colback=cnom!6,colframe=cnom,arc=2.5mm,boxrule=0.9pt,left=3mm]{\Large\bfseries\color{cnom}d)}\quad\begin{minipage}[c]{5.3cm}\raggedright\Large $(x-4)^2$\end{minipage}\end{tcolorbox}
\end{tcolorbox}
\newpage
\begin{tcolorbox}[enhanced,width=16cm,height=14cm,colback=white,colframe=cnom,boxrule=2pt,arc=0mm,
left=6mm,right=6mm,top=0mm,bottom=4mm,sharp corners,
overlay={\fill[cnom] (frame.north west) rectangle ([yshift=-1.25cm]frame.north east);
\node[anchor=west,text=white,font=\sffamily\bfseries\LARGE] at ([xshift=5mm,yshift=-0.62cm]frame.north west) {Question 11};
\node[anchor=east,text=white,font=\sffamily\itshape\large] at ([xshift=-5mm,yshift=-0.62cm]frame.north east) {Nombres et calculs};}]
\vspace*{1.75cm}
\begin{minipage}[t][6.9cm][c]{\linewidth}\LARGE\centering
\begin{minipage}{0.95\linewidth}\raggedright\linespread{1.25}\selectfont Que signifie $x\in[-2\,;\,5[$~?\end{minipage}
\end{minipage}

\vfill
\begin{tcolorbox}[enhanced,nobeforeafter,width=7.1cm,height=2.1cm,valign=center,colback=cnom!6,colframe=cnom,arc=2.5mm,boxrule=0.9pt,left=3mm]{\Large\bfseries\color{cnom}a)}\quad\begin{minipage}[c]{5.3cm}\raggedright\Large $-2\leqslant x<5$\end{minipage}\end{tcolorbox}\hfill\begin{tcolorbox}[enhanced,nobeforeafter,width=7.1cm,height=2.1cm,valign=center,colback=cnom!6,colframe=cnom,arc=2.5mm,boxrule=0.9pt,left=3mm]{\Large\bfseries\color{cnom}b)}\quad\begin{minipage}[c]{5.3cm}\raggedright\Large $-2\leqslant x\leqslant 5$\end{minipage}\end{tcolorbox}\par\vspace{2.5mm}
\begin{tcolorbox}[enhanced,nobeforeafter,width=7.1cm,height=2.1cm,valign=center,colback=cnom!6,colframe=cnom,arc=2.5mm,boxrule=0.9pt,left=3mm]{\Large\bfseries\color{cnom}c)}\quad\begin{minipage}[c]{5.3cm}\raggedright\Large $-2<x\leqslant 5$\end{minipage}\end{tcolorbox}\hfill\begin{tcolorbox}[enhanced,nobeforeafter,width=7.1cm,height=2.1cm,valign=center,colback=cnom!6,colframe=cnom,arc=2.5mm,boxrule=0.9pt,left=3mm]{\Large\bfseries\color{cnom}d)}\quad\begin{minipage}[c]{5.3cm}\raggedright\Large $-2<x<5$\end{minipage}\end{tcolorbox}
\end{tcolorbox}
\newpage
\begin{tcolorbox}[enhanced,width=16cm,height=14cm,colback=white,colframe=cnom,boxrule=2pt,arc=0mm,
left=6mm,right=6mm,top=0mm,bottom=4mm,sharp corners,
overlay={\fill[cnom] (frame.north west) rectangle ([yshift=-1.25cm]frame.north east);
\node[anchor=west,text=white,font=\sffamily\bfseries\LARGE] at ([xshift=5mm,yshift=-0.62cm]frame.north west) {Question 12};
\node[anchor=east,text=white,font=\sffamily\itshape\large] at ([xshift=-5mm,yshift=-0.62cm]frame.north east) {Nombres et calculs};}]
\vspace*{1.75cm}
\begin{minipage}[t][6.9cm][c]{\linewidth}\LARGE\centering
\begin{minipage}{0.95\linewidth}\raggedright\linespread{1.25}\selectfont Que signifie $x\in]3\,;\,+\infty[$~?\end{minipage}
\end{minipage}

\vfill
\begin{tcolorbox}[enhanced,nobeforeafter,width=7.1cm,height=2.1cm,valign=center,colback=cnom!6,colframe=cnom,arc=2.5mm,boxrule=0.9pt,left=3mm]{\Large\bfseries\color{cnom}a)}\quad\begin{minipage}[c]{5.3cm}\raggedright\Large $x\geqslant 3$\end{minipage}\end{tcolorbox}\hfill\begin{tcolorbox}[enhanced,nobeforeafter,width=7.1cm,height=2.1cm,valign=center,colback=cnom!6,colframe=cnom,arc=2.5mm,boxrule=0.9pt,left=3mm]{\Large\bfseries\color{cnom}b)}\quad\begin{minipage}[c]{5.3cm}\raggedright\Large $x>3$\end{minipage}\end{tcolorbox}\par\vspace{2.5mm}
\begin{tcolorbox}[enhanced,nobeforeafter,width=7.1cm,height=2.1cm,valign=center,colback=cnom!6,colframe=cnom,arc=2.5mm,boxrule=0.9pt,left=3mm]{\Large\bfseries\color{cnom}c)}\quad\begin{minipage}[c]{5.3cm}\raggedright\Large $x\leqslant 3$\end{minipage}\end{tcolorbox}\hfill\begin{tcolorbox}[enhanced,nobeforeafter,width=7.1cm,height=2.1cm,valign=center,colback=cnom!6,colframe=cnom,arc=2.5mm,boxrule=0.9pt,left=3mm]{\Large\bfseries\color{cnom}d)}\quad\begin{minipage}[c]{5.3cm}\raggedright\Large $x<3$\end{minipage}\end{tcolorbox}
\end{tcolorbox}
\newpage
\begin{tcolorbox}[enhanced,width=16cm,height=14cm,colback=white,colframe=cnom,boxrule=2pt,arc=0mm,
left=6mm,right=6mm,top=0mm,bottom=4mm,sharp corners,
overlay={\fill[cnom] (frame.north west) rectangle ([yshift=-1.25cm]frame.north east);
\node[anchor=west,text=white,font=\sffamily\bfseries\LARGE] at ([xshift=5mm,yshift=-0.62cm]frame.north west) {Question 13};
\node[anchor=east,text=white,font=\sffamily\itshape\large] at ([xshift=-5mm,yshift=-0.62cm]frame.north east) {Nombres et calculs};}]
\vspace*{1.75cm}
\begin{minipage}[t][6.9cm][c]{\linewidth}\LARGE\centering
\begin{minipage}{0.95\linewidth}\raggedright\linespread{1.25}\selectfont Pour quels réels $x$ a-t-on $|x-3|\leqslant 2$~?\end{minipage}
\end{minipage}

\vfill
\begin{tcolorbox}[enhanced,nobeforeafter,width=7.1cm,height=2.1cm,valign=center,colback=cnom!6,colframe=cnom,arc=2.5mm,boxrule=0.9pt,left=3mm]{\Large\bfseries\color{cnom}a)}\quad\begin{minipage}[c]{5.3cm}\raggedright\Large $x\in[-1\,;\,7]$\end{minipage}\end{tcolorbox}\hfill\begin{tcolorbox}[enhanced,nobeforeafter,width=7.1cm,height=2.1cm,valign=center,colback=cnom!6,colframe=cnom,arc=2.5mm,boxrule=0.9pt,left=3mm]{\Large\bfseries\color{cnom}b)}\quad\begin{minipage}[c]{5.3cm}\raggedright\Large $x\in[3\,;\,5]$\end{minipage}\end{tcolorbox}\par\vspace{2.5mm}
\begin{tcolorbox}[enhanced,nobeforeafter,width=7.1cm,height=2.1cm,valign=center,colback=cnom!6,colframe=cnom,arc=2.5mm,boxrule=0.9pt,left=3mm]{\Large\bfseries\color{cnom}c)}\quad\begin{minipage}[c]{5.3cm}\raggedright\Large $x\in[1\,;\,5]$\end{minipage}\end{tcolorbox}\hfill\begin{tcolorbox}[enhanced,nobeforeafter,width=7.1cm,height=2.1cm,valign=center,colback=cnom!6,colframe=cnom,arc=2.5mm,boxrule=0.9pt,left=3mm]{\Large\bfseries\color{cnom}d)}\quad\begin{minipage}[c]{5.3cm}\raggedright\Large $x\in[-5\,;\,-1]$\end{minipage}\end{tcolorbox}
\end{tcolorbox}
\newpage
\begin{tcolorbox}[enhanced,width=16cm,height=14cm,colback=white,colframe=cnom,boxrule=2pt,arc=0mm,
left=6mm,right=6mm,top=0mm,bottom=4mm,sharp corners,
overlay={\fill[cnom] (frame.north west) rectangle ([yshift=-1.25cm]frame.north east);
\node[anchor=west,text=white,font=\sffamily\bfseries\LARGE] at ([xshift=5mm,yshift=-0.62cm]frame.north west) {Question 14};
\node[anchor=east,text=white,font=\sffamily\itshape\large] at ([xshift=-5mm,yshift=-0.62cm]frame.north east) {Nombres et calculs};}]
\vspace*{1.75cm}
\begin{minipage}[t][6.9cm][c]{\linewidth}\LARGE\centering
\begin{minipage}{0.95\linewidth}\raggedright\linespread{1.25}\selectfont Pour quels réels $x$ a-t-on $|x+1|\leqslant 4$~?\end{minipage}
\end{minipage}

\vfill
\begin{tcolorbox}[enhanced,nobeforeafter,width=7.1cm,height=2.1cm,valign=center,colback=cnom!6,colframe=cnom,arc=2.5mm,boxrule=0.9pt,left=3mm]{\Large\bfseries\color{cnom}a)}\quad\begin{minipage}[c]{5.3cm}\raggedright\Large $x\in[-9\,;\,7]$\end{minipage}\end{tcolorbox}\hfill\begin{tcolorbox}[enhanced,nobeforeafter,width=7.1cm,height=2.1cm,valign=center,colback=cnom!6,colframe=cnom,arc=2.5mm,boxrule=0.9pt,left=3mm]{\Large\bfseries\color{cnom}b)}\quad\begin{minipage}[c]{5.3cm}\raggedright\Large $x\in[-3\,;\,5]$\end{minipage}\end{tcolorbox}\par\vspace{2.5mm}
\begin{tcolorbox}[enhanced,nobeforeafter,width=7.1cm,height=2.1cm,valign=center,colback=cnom!6,colframe=cnom,arc=2.5mm,boxrule=0.9pt,left=3mm]{\Large\bfseries\color{cnom}c)}\quad\begin{minipage}[c]{5.3cm}\raggedright\Large $x\in[-5\,;\,3]$\end{minipage}\end{tcolorbox}\hfill\begin{tcolorbox}[enhanced,nobeforeafter,width=7.1cm,height=2.1cm,valign=center,colback=cnom!6,colframe=cnom,arc=2.5mm,boxrule=0.9pt,left=3mm]{\Large\bfseries\color{cnom}d)}\quad\begin{minipage}[c]{5.3cm}\raggedright\Large $x\in[-1\,;\,3]$\end{minipage}\end{tcolorbox}
\end{tcolorbox}
\newpage
\begin{tcolorbox}[enhanced,width=16cm,height=14cm,colback=white,colframe=cnom,boxrule=2pt,arc=0mm,
left=6mm,right=6mm,top=0mm,bottom=4mm,sharp corners,
overlay={\fill[cnom] (frame.north west) rectangle ([yshift=-1.25cm]frame.north east);
\node[anchor=west,text=white,font=\sffamily\bfseries\LARGE] at ([xshift=5mm,yshift=-0.62cm]frame.north west) {Question 15};
\node[anchor=east,text=white,font=\sffamily\itshape\large] at ([xshift=-5mm,yshift=-0.62cm]frame.north east) {Nombres et calculs};}]
\vspace*{1.75cm}
\begin{minipage}[t][6.9cm][c]{\linewidth}\LARGE\centering
\begin{minipage}{0.95\linewidth}\raggedright\linespread{1.25}\selectfont Écrire $\dfrac{2^{5}\times2^{-3}}{2^{4}}$ sous la forme d'une puissance de $2$.\end{minipage}
\end{minipage}

\vfill
\begin{tcolorbox}[enhanced,nobeforeafter,width=7.1cm,height=2.1cm,valign=center,colback=cnom!6,colframe=cnom,arc=2.5mm,boxrule=0.9pt,left=3mm]{\Large\bfseries\color{cnom}a)}\quad\begin{minipage}[c]{5.3cm}\raggedright\Large $2^{-19}$\end{minipage}\end{tcolorbox}\hfill\begin{tcolorbox}[enhanced,nobeforeafter,width=7.1cm,height=2.1cm,valign=center,colback=cnom!6,colframe=cnom,arc=2.5mm,boxrule=0.9pt,left=3mm]{\Large\bfseries\color{cnom}b)}\quad\begin{minipage}[c]{5.3cm}\raggedright\Large $2^{4}$\end{minipage}\end{tcolorbox}\par\vspace{2.5mm}
\begin{tcolorbox}[enhanced,nobeforeafter,width=7.1cm,height=2.1cm,valign=center,colback=cnom!6,colframe=cnom,arc=2.5mm,boxrule=0.9pt,left=3mm]{\Large\bfseries\color{cnom}c)}\quad\begin{minipage}[c]{5.3cm}\raggedright\Large $2^{6}$\end{minipage}\end{tcolorbox}\hfill\begin{tcolorbox}[enhanced,nobeforeafter,width=7.1cm,height=2.1cm,valign=center,colback=cnom!6,colframe=cnom,arc=2.5mm,boxrule=0.9pt,left=3mm]{\Large\bfseries\color{cnom}d)}\quad\begin{minipage}[c]{5.3cm}\raggedright\Large $2^{-2}$\end{minipage}\end{tcolorbox}
\end{tcolorbox}
\newpage
\begin{tcolorbox}[enhanced,width=16cm,height=14cm,colback=white,colframe=cnom,boxrule=2pt,arc=0mm,
left=6mm,right=6mm,top=0mm,bottom=4mm,sharp corners,
overlay={\fill[cnom] (frame.north west) rectangle ([yshift=-1.25cm]frame.north east);
\node[anchor=west,text=white,font=\sffamily\bfseries\LARGE] at ([xshift=5mm,yshift=-0.62cm]frame.north west) {Question 16};
\node[anchor=east,text=white,font=\sffamily\itshape\large] at ([xshift=-5mm,yshift=-0.62cm]frame.north east) {Nombres et calculs};}]
\vspace*{1.75cm}
\begin{minipage}[t][6.9cm][c]{\linewidth}\LARGE\centering
\begin{minipage}{0.95\linewidth}\raggedright\linespread{1.25}\selectfont Écrire $\dfrac{3^{4}\times3^{2}}{3^{7}}$ sous la forme d'une puissance de $3$.\end{minipage}
\end{minipage}

\vfill
\begin{tcolorbox}[enhanced,nobeforeafter,width=7.1cm,height=2.1cm,valign=center,colback=cnom!6,colframe=cnom,arc=2.5mm,boxrule=0.9pt,left=3mm]{\Large\bfseries\color{cnom}a)}\quad\begin{minipage}[c]{5.3cm}\raggedright\Large $3^{1}$\end{minipage}\end{tcolorbox}\hfill\begin{tcolorbox}[enhanced,nobeforeafter,width=7.1cm,height=2.1cm,valign=center,colback=cnom!6,colframe=cnom,arc=2.5mm,boxrule=0.9pt,left=3mm]{\Large\bfseries\color{cnom}b)}\quad\begin{minipage}[c]{5.3cm}\raggedright\Large $3^{-5}$\end{minipage}\end{tcolorbox}\par\vspace{2.5mm}
\begin{tcolorbox}[enhanced,nobeforeafter,width=7.1cm,height=2.1cm,valign=center,colback=cnom!6,colframe=cnom,arc=2.5mm,boxrule=0.9pt,left=3mm]{\Large\bfseries\color{cnom}c)}\quad\begin{minipage}[c]{5.3cm}\raggedright\Large $3^{-1}$\end{minipage}\end{tcolorbox}\hfill\begin{tcolorbox}[enhanced,nobeforeafter,width=7.1cm,height=2.1cm,valign=center,colback=cnom!6,colframe=cnom,arc=2.5mm,boxrule=0.9pt,left=3mm]{\Large\bfseries\color{cnom}d)}\quad\begin{minipage}[c]{5.3cm}\raggedright\Large $3^{13}$\end{minipage}\end{tcolorbox}
\end{tcolorbox}
\newpage
\begin{tcolorbox}[enhanced,width=16cm,height=14cm,colback=white,colframe=csta,boxrule=2pt,arc=0mm,
left=6mm,right=6mm,top=0mm,bottom=4mm,sharp corners,
overlay={\fill[csta] (frame.north west) rectangle ([yshift=-1.25cm]frame.north east);
\node[anchor=west,text=white,font=\sffamily\bfseries\LARGE] at ([xshift=5mm,yshift=-0.62cm]frame.north west) {Question 17};
\node[anchor=east,text=white,font=\sffamily\itshape\large] at ([xshift=-5mm,yshift=-0.62cm]frame.north east) {Évolutions et statistiques};}]
\vspace*{1.75cm}
\begin{minipage}[t][6.9cm][c]{\linewidth}\LARGE\centering
\begin{minipage}{0.95\linewidth}\raggedright\linespread{1.25}\selectfont Quel est le coefficient multiplicateur associé à une hausse de $20\,\%$~?\end{minipage}
\end{minipage}

\vfill
\begin{tcolorbox}[enhanced,nobeforeafter,width=7.1cm,height=2.1cm,valign=center,colback=csta!6,colframe=csta,arc=2.5mm,boxrule=0.9pt,left=3mm]{\Large\bfseries\color{csta}a)}\quad\begin{minipage}[c]{5.3cm}\raggedright\Large $0{,}8$\end{minipage}\end{tcolorbox}\hfill\begin{tcolorbox}[enhanced,nobeforeafter,width=7.1cm,height=2.1cm,valign=center,colback=csta!6,colframe=csta,arc=2.5mm,boxrule=0.9pt,left=3mm]{\Large\bfseries\color{csta}b)}\quad\begin{minipage}[c]{5.3cm}\raggedright\Large $1{,}2$\end{minipage}\end{tcolorbox}\par\vspace{2.5mm}
\begin{tcolorbox}[enhanced,nobeforeafter,width=7.1cm,height=2.1cm,valign=center,colback=csta!6,colframe=csta,arc=2.5mm,boxrule=0.9pt,left=3mm]{\Large\bfseries\color{csta}c)}\quad\begin{minipage}[c]{5.3cm}\raggedright\Large $0{,}2$\end{minipage}\end{tcolorbox}\hfill\begin{tcolorbox}[enhanced,nobeforeafter,width=7.1cm,height=2.1cm,valign=center,colback=csta!6,colframe=csta,arc=2.5mm,boxrule=0.9pt,left=3mm]{\Large\bfseries\color{csta}d)}\quad\begin{minipage}[c]{5.3cm}\raggedright\Large $3$\end{minipage}\end{tcolorbox}
\end{tcolorbox}
\newpage
\begin{tcolorbox}[enhanced,width=16cm,height=14cm,colback=white,colframe=csta,boxrule=2pt,arc=0mm,
left=6mm,right=6mm,top=0mm,bottom=4mm,sharp corners,
overlay={\fill[csta] (frame.north west) rectangle ([yshift=-1.25cm]frame.north east);
\node[anchor=west,text=white,font=\sffamily\bfseries\LARGE] at ([xshift=5mm,yshift=-0.62cm]frame.north west) {Question 18};
\node[anchor=east,text=white,font=\sffamily\itshape\large] at ([xshift=-5mm,yshift=-0.62cm]frame.north east) {Évolutions et statistiques};}]
\vspace*{1.75cm}
\begin{minipage}[t][6.9cm][c]{\linewidth}\LARGE\centering
\begin{minipage}{0.95\linewidth}\raggedright\linespread{1.25}\selectfont Quel est le coefficient multiplicateur associé à une baisse de $15\,\%$~?\end{minipage}
\end{minipage}

\vfill
\begin{tcolorbox}[enhanced,nobeforeafter,width=7.1cm,height=2.1cm,valign=center,colback=csta!6,colframe=csta,arc=2.5mm,boxrule=0.9pt,left=3mm]{\Large\bfseries\color{csta}a)}\quad\begin{minipage}[c]{5.3cm}\raggedright\Large $1{,}15$\end{minipage}\end{tcolorbox}\hfill\begin{tcolorbox}[enhanced,nobeforeafter,width=7.1cm,height=2.1cm,valign=center,colback=csta!6,colframe=csta,arc=2.5mm,boxrule=0.9pt,left=3mm]{\Large\bfseries\color{csta}b)}\quad\begin{minipage}[c]{5.3cm}\raggedright\Large $-0{,}15$\end{minipage}\end{tcolorbox}\par\vspace{2.5mm}
\begin{tcolorbox}[enhanced,nobeforeafter,width=7.1cm,height=2.1cm,valign=center,colback=csta!6,colframe=csta,arc=2.5mm,boxrule=0.9pt,left=3mm]{\Large\bfseries\color{csta}c)}\quad\begin{minipage}[c]{5.3cm}\raggedright\Large $-0{,}5$\end{minipage}\end{tcolorbox}\hfill\begin{tcolorbox}[enhanced,nobeforeafter,width=7.1cm,height=2.1cm,valign=center,colback=csta!6,colframe=csta,arc=2.5mm,boxrule=0.9pt,left=3mm]{\Large\bfseries\color{csta}d)}\quad\begin{minipage}[c]{5.3cm}\raggedright\Large $0{,}85$\end{minipage}\end{tcolorbox}
\end{tcolorbox}
\newpage
\begin{tcolorbox}[enhanced,width=16cm,height=14cm,colback=white,colframe=csta,boxrule=2pt,arc=0mm,
left=6mm,right=6mm,top=0mm,bottom=4mm,sharp corners,
overlay={\fill[csta] (frame.north west) rectangle ([yshift=-1.25cm]frame.north east);
\node[anchor=west,text=white,font=\sffamily\bfseries\LARGE] at ([xshift=5mm,yshift=-0.62cm]frame.north west) {Question 19};
\node[anchor=east,text=white,font=\sffamily\itshape\large] at ([xshift=-5mm,yshift=-0.62cm]frame.north east) {Évolutions et statistiques};}]
\vspace*{1.75cm}
\begin{minipage}[t][6.9cm][c]{\linewidth}\LARGE\centering
\begin{minipage}{0.95\linewidth}\raggedright\linespread{1.25}\selectfont Un prix de $80$ € augmente de $25\,\%$. Quel est le nouveau prix~?\end{minipage}
\end{minipage}

\vfill
\begin{tcolorbox}[enhanced,nobeforeafter,width=7.1cm,height=2.1cm,valign=center,colback=csta!6,colframe=csta,arc=2.5mm,boxrule=0.9pt,left=3mm]{\Large\bfseries\color{csta}a)}\quad\begin{minipage}[c]{5.3cm}\raggedright\Large $20~\text{€}$\end{minipage}\end{tcolorbox}\hfill\begin{tcolorbox}[enhanced,nobeforeafter,width=7.1cm,height=2.1cm,valign=center,colback=csta!6,colframe=csta,arc=2.5mm,boxrule=0.9pt,left=3mm]{\Large\bfseries\color{csta}b)}\quad\begin{minipage}[c]{5.3cm}\raggedright\Large $105~\text{€}$\end{minipage}\end{tcolorbox}\par\vspace{2.5mm}
\begin{tcolorbox}[enhanced,nobeforeafter,width=7.1cm,height=2.1cm,valign=center,colback=csta!6,colframe=csta,arc=2.5mm,boxrule=0.9pt,left=3mm]{\Large\bfseries\color{csta}c)}\quad\begin{minipage}[c]{5.3cm}\raggedright\Large $60~\text{€}$\end{minipage}\end{tcolorbox}\hfill\begin{tcolorbox}[enhanced,nobeforeafter,width=7.1cm,height=2.1cm,valign=center,colback=csta!6,colframe=csta,arc=2.5mm,boxrule=0.9pt,left=3mm]{\Large\bfseries\color{csta}d)}\quad\begin{minipage}[c]{5.3cm}\raggedright\Large $100~\text{€}$\end{minipage}\end{tcolorbox}
\end{tcolorbox}
\newpage
\begin{tcolorbox}[enhanced,width=16cm,height=14cm,colback=white,colframe=csta,boxrule=2pt,arc=0mm,
left=6mm,right=6mm,top=0mm,bottom=4mm,sharp corners,
overlay={\fill[csta] (frame.north west) rectangle ([yshift=-1.25cm]frame.north east);
\node[anchor=west,text=white,font=\sffamily\bfseries\LARGE] at ([xshift=5mm,yshift=-0.62cm]frame.north west) {Question 20};
\node[anchor=east,text=white,font=\sffamily\itshape\large] at ([xshift=-5mm,yshift=-0.62cm]frame.north east) {Évolutions et statistiques};}]
\vspace*{1.75cm}
\begin{minipage}[t][6.9cm][c]{\linewidth}\LARGE\centering
\begin{minipage}{0.95\linewidth}\raggedright\linespread{1.25}\selectfont Un prix de $250$ € baisse de $20\,\%$. Quel est le nouveau prix~?\end{minipage}
\end{minipage}

\vfill
\begin{tcolorbox}[enhanced,nobeforeafter,width=7.1cm,height=2.1cm,valign=center,colback=csta!6,colframe=csta,arc=2.5mm,boxrule=0.9pt,left=3mm]{\Large\bfseries\color{csta}a)}\quad\begin{minipage}[c]{5.3cm}\raggedright\Large $300~\text{€}$\end{minipage}\end{tcolorbox}\hfill\begin{tcolorbox}[enhanced,nobeforeafter,width=7.1cm,height=2.1cm,valign=center,colback=csta!6,colframe=csta,arc=2.5mm,boxrule=0.9pt,left=3mm]{\Large\bfseries\color{csta}b)}\quad\begin{minipage}[c]{5.3cm}\raggedright\Large $200~\text{€}$\end{minipage}\end{tcolorbox}\par\vspace{2.5mm}
\begin{tcolorbox}[enhanced,nobeforeafter,width=7.1cm,height=2.1cm,valign=center,colback=csta!6,colframe=csta,arc=2.5mm,boxrule=0.9pt,left=3mm]{\Large\bfseries\color{csta}c)}\quad\begin{minipage}[c]{5.3cm}\raggedright\Large $230~\text{€}$\end{minipage}\end{tcolorbox}\hfill\begin{tcolorbox}[enhanced,nobeforeafter,width=7.1cm,height=2.1cm,valign=center,colback=csta!6,colframe=csta,arc=2.5mm,boxrule=0.9pt,left=3mm]{\Large\bfseries\color{csta}d)}\quad\begin{minipage}[c]{5.3cm}\raggedright\Large $50~\text{€}$\end{minipage}\end{tcolorbox}
\end{tcolorbox}
\newpage
\begin{tcolorbox}[enhanced,width=16cm,height=14cm,colback=white,colframe=csta,boxrule=2pt,arc=0mm,
left=6mm,right=6mm,top=0mm,bottom=4mm,sharp corners,
overlay={\fill[csta] (frame.north west) rectangle ([yshift=-1.25cm]frame.north east);
\node[anchor=west,text=white,font=\sffamily\bfseries\LARGE] at ([xshift=5mm,yshift=-0.62cm]frame.north west) {Question 21};
\node[anchor=east,text=white,font=\sffamily\itshape\large] at ([xshift=-5mm,yshift=-0.62cm]frame.north east) {Évolutions et statistiques};}]
\vspace*{1.75cm}
\begin{minipage}[t][6.9cm][c]{\linewidth}\LARGE\centering
\begin{minipage}{0.95\linewidth}\raggedright\linespread{1.25}\selectfont Une quantité passe de $50$ à $65$. Quel est son taux d'évolution~?\end{minipage}
\end{minipage}

\vfill
\begin{tcolorbox}[enhanced,nobeforeafter,width=7.1cm,height=2.1cm,valign=center,colback=csta!6,colframe=csta,arc=2.5mm,boxrule=0.9pt,left=3mm]{\Large\bfseries\color{csta}a)}\quad\begin{minipage}[c]{5.3cm}\raggedright\Large $+30\,\%$\end{minipage}\end{tcolorbox}\hfill\begin{tcolorbox}[enhanced,nobeforeafter,width=7.1cm,height=2.1cm,valign=center,colback=csta!6,colframe=csta,arc=2.5mm,boxrule=0.9pt,left=3mm]{\Large\bfseries\color{csta}b)}\quad\begin{minipage}[c]{5.3cm}\raggedright\Large $+35\,\%$\end{minipage}\end{tcolorbox}\par\vspace{2.5mm}
\begin{tcolorbox}[enhanced,nobeforeafter,width=7.1cm,height=2.1cm,valign=center,colback=csta!6,colframe=csta,arc=2.5mm,boxrule=0.9pt,left=3mm]{\Large\bfseries\color{csta}c)}\quad\begin{minipage}[c]{5.3cm}\raggedright\Large $-30\,\%$\end{minipage}\end{tcolorbox}\hfill\begin{tcolorbox}[enhanced,nobeforeafter,width=7.1cm,height=2.1cm,valign=center,colback=csta!6,colframe=csta,arc=2.5mm,boxrule=0.9pt,left=3mm]{\Large\bfseries\color{csta}d)}\quad\begin{minipage}[c]{5.3cm}\raggedright\Large $+15\,\%$\end{minipage}\end{tcolorbox}
\end{tcolorbox}
\newpage
\begin{tcolorbox}[enhanced,width=16cm,height=14cm,colback=white,colframe=csta,boxrule=2pt,arc=0mm,
left=6mm,right=6mm,top=0mm,bottom=4mm,sharp corners,
overlay={\fill[csta] (frame.north west) rectangle ([yshift=-1.25cm]frame.north east);
\node[anchor=west,text=white,font=\sffamily\bfseries\LARGE] at ([xshift=5mm,yshift=-0.62cm]frame.north west) {Question 22};
\node[anchor=east,text=white,font=\sffamily\itshape\large] at ([xshift=-5mm,yshift=-0.62cm]frame.north east) {Évolutions et statistiques};}]
\vspace*{1.75cm}
\begin{minipage}[t][6.9cm][c]{\linewidth}\LARGE\centering
\begin{minipage}{0.95\linewidth}\raggedright\linespread{1.25}\selectfont Une quantité passe de $80$ à $60$. Quel est son taux d'évolution~?\end{minipage}
\end{minipage}

\vfill
\begin{tcolorbox}[enhanced,nobeforeafter,width=7.1cm,height=2.1cm,valign=center,colback=csta!6,colframe=csta,arc=2.5mm,boxrule=0.9pt,left=3mm]{\Large\bfseries\color{csta}a)}\quad\begin{minipage}[c]{5.3cm}\raggedright\Large $-12{,}5\,\%$\end{minipage}\end{tcolorbox}\hfill\begin{tcolorbox}[enhanced,nobeforeafter,width=7.1cm,height=2.1cm,valign=center,colback=csta!6,colframe=csta,arc=2.5mm,boxrule=0.9pt,left=3mm]{\Large\bfseries\color{csta}b)}\quad\begin{minipage}[c]{5.3cm}\raggedright\Large $+25\,\%$\end{minipage}\end{tcolorbox}\par\vspace{2.5mm}
\begin{tcolorbox}[enhanced,nobeforeafter,width=7.1cm,height=2.1cm,valign=center,colback=csta!6,colframe=csta,arc=2.5mm,boxrule=0.9pt,left=3mm]{\Large\bfseries\color{csta}c)}\quad\begin{minipage}[c]{5.3cm}\raggedright\Large $-20\,\%$\end{minipage}\end{tcolorbox}\hfill\begin{tcolorbox}[enhanced,nobeforeafter,width=7.1cm,height=2.1cm,valign=center,colback=csta!6,colframe=csta,arc=2.5mm,boxrule=0.9pt,left=3mm]{\Large\bfseries\color{csta}d)}\quad\begin{minipage}[c]{5.3cm}\raggedright\Large $-25\,\%$\end{minipage}\end{tcolorbox}
\end{tcolorbox}
\newpage
\begin{tcolorbox}[enhanced,width=16cm,height=14cm,colback=white,colframe=csta,boxrule=2pt,arc=0mm,
left=6mm,right=6mm,top=0mm,bottom=4mm,sharp corners,
overlay={\fill[csta] (frame.north west) rectangle ([yshift=-1.25cm]frame.north east);
\node[anchor=west,text=white,font=\sffamily\bfseries\LARGE] at ([xshift=5mm,yshift=-0.62cm]frame.north west) {Question 23};
\node[anchor=east,text=white,font=\sffamily\itshape\large] at ([xshift=-5mm,yshift=-0.62cm]frame.north east) {Évolutions et statistiques};}]
\vspace*{1.75cm}
\begin{minipage}[t][6.9cm][c]{\linewidth}\LARGE\centering
\begin{minipage}{0.95\linewidth}\raggedright\linespread{1.25}\selectfont Un prix subit une hausse de $20\,\%$, puis une baisse de $20\,\%$. Quelle est son évolution globale~?\end{minipage}
\end{minipage}

\vfill
\begin{tcolorbox}[enhanced,nobeforeafter,width=7.1cm,height=2.1cm,valign=center,colback=csta!6,colframe=csta,arc=2.5mm,boxrule=0.9pt,left=3mm]{\Large\bfseries\color{csta}a)}\quad\begin{minipage}[c]{5.3cm}\raggedright\Large $-6\,\%$\end{minipage}\end{tcolorbox}\hfill\begin{tcolorbox}[enhanced,nobeforeafter,width=7.1cm,height=2.1cm,valign=center,colback=csta!6,colframe=csta,arc=2.5mm,boxrule=0.9pt,left=3mm]{\Large\bfseries\color{csta}b)}\quad\begin{minipage}[c]{5.3cm}\raggedright\Large $-4\,\%$\end{minipage}\end{tcolorbox}\par\vspace{2.5mm}
\begin{tcolorbox}[enhanced,nobeforeafter,width=7.1cm,height=2.1cm,valign=center,colback=csta!6,colframe=csta,arc=2.5mm,boxrule=0.9pt,left=3mm]{\Large\bfseries\color{csta}c)}\quad\begin{minipage}[c]{5.3cm}\raggedright\Large $0\,\%$\end{minipage}\end{tcolorbox}\hfill\begin{tcolorbox}[enhanced,nobeforeafter,width=7.1cm,height=2.1cm,valign=center,colback=csta!6,colframe=csta,arc=2.5mm,boxrule=0.9pt,left=3mm]{\Large\bfseries\color{csta}d)}\quad\begin{minipage}[c]{5.3cm}\raggedright\Large $-3\,\%$\end{minipage}\end{tcolorbox}
\end{tcolorbox}
\newpage
\begin{tcolorbox}[enhanced,width=16cm,height=14cm,colback=white,colframe=csta,boxrule=2pt,arc=0mm,
left=6mm,right=6mm,top=0mm,bottom=4mm,sharp corners,
overlay={\fill[csta] (frame.north west) rectangle ([yshift=-1.25cm]frame.north east);
\node[anchor=west,text=white,font=\sffamily\bfseries\LARGE] at ([xshift=5mm,yshift=-0.62cm]frame.north west) {Question 24};
\node[anchor=east,text=white,font=\sffamily\itshape\large] at ([xshift=-5mm,yshift=-0.62cm]frame.north east) {Évolutions et statistiques};}]
\vspace*{1.75cm}
\begin{minipage}[t][6.9cm][c]{\linewidth}\LARGE\centering
\begin{minipage}{0.95\linewidth}\raggedright\linespread{1.25}\selectfont Un prix subit une hausse de $10\,\%$, puis une hausse de $10\,\%$. Quelle est son évolution globale~?\end{minipage}
\end{minipage}

\vfill
\begin{tcolorbox}[enhanced,nobeforeafter,width=7.1cm,height=2.1cm,valign=center,colback=csta!6,colframe=csta,arc=2.5mm,boxrule=0.9pt,left=3mm]{\Large\bfseries\color{csta}a)}\quad\begin{minipage}[c]{5.3cm}\raggedright\Large $+20\,\%$\end{minipage}\end{tcolorbox}\hfill\begin{tcolorbox}[enhanced,nobeforeafter,width=7.1cm,height=2.1cm,valign=center,colback=csta!6,colframe=csta,arc=2.5mm,boxrule=0.9pt,left=3mm]{\Large\bfseries\color{csta}b)}\quad\begin{minipage}[c]{5.3cm}\raggedright\Large $+21\,\%$\end{minipage}\end{tcolorbox}\par\vspace{2.5mm}
\begin{tcolorbox}[enhanced,nobeforeafter,width=7.1cm,height=2.1cm,valign=center,colback=csta!6,colframe=csta,arc=2.5mm,boxrule=0.9pt,left=3mm]{\Large\bfseries\color{csta}c)}\quad\begin{minipage}[c]{5.3cm}\raggedright\Large $+22\,\%$\end{minipage}\end{tcolorbox}\hfill\begin{tcolorbox}[enhanced,nobeforeafter,width=7.1cm,height=2.1cm,valign=center,colback=csta!6,colframe=csta,arc=2.5mm,boxrule=0.9pt,left=3mm]{\Large\bfseries\color{csta}d)}\quad\begin{minipage}[c]{5.3cm}\raggedright\Large $+1\,\%$\end{minipage}\end{tcolorbox}
\end{tcolorbox}
\newpage
\begin{tcolorbox}[enhanced,width=16cm,height=14cm,colback=white,colframe=csta,boxrule=2pt,arc=0mm,
left=6mm,right=6mm,top=0mm,bottom=4mm,sharp corners,
overlay={\fill[csta] (frame.north west) rectangle ([yshift=-1.25cm]frame.north east);
\node[anchor=west,text=white,font=\sffamily\bfseries\LARGE] at ([xshift=5mm,yshift=-0.62cm]frame.north west) {Question 25};
\node[anchor=east,text=white,font=\sffamily\itshape\large] at ([xshift=-5mm,yshift=-0.62cm]frame.north east) {Évolutions et statistiques};}]
\vspace*{1.75cm}
\begin{minipage}[t][6.9cm][c]{\linewidth}\LARGE\centering
\begin{minipage}{0.95\linewidth}\raggedright\linespread{1.25}\selectfont Après une baisse de $20\,\%$, quelle évolution faut-il appliquer pour revenir au prix de départ~?\end{minipage}
\end{minipage}

\vfill
\begin{tcolorbox}[enhanced,nobeforeafter,width=7.1cm,height=2.1cm,valign=center,colback=csta!6,colframe=csta,arc=2.5mm,boxrule=0.9pt,left=3mm]{\Large\bfseries\color{csta}a)}\quad\begin{minipage}[c]{5.3cm}\raggedright\Large une baisse de $25\,\%$\end{minipage}\end{tcolorbox}\hfill\begin{tcolorbox}[enhanced,nobeforeafter,width=7.1cm,height=2.1cm,valign=center,colback=csta!6,colframe=csta,arc=2.5mm,boxrule=0.9pt,left=3mm]{\Large\bfseries\color{csta}b)}\quad\begin{minipage}[c]{5.3cm}\raggedright\Large une hausse de $20\,\%$\end{minipage}\end{tcolorbox}\par\vspace{2.5mm}
\begin{tcolorbox}[enhanced,nobeforeafter,width=7.1cm,height=2.1cm,valign=center,colback=csta!6,colframe=csta,arc=2.5mm,boxrule=0.9pt,left=3mm]{\Large\bfseries\color{csta}c)}\quad\begin{minipage}[c]{5.3cm}\raggedright\Large une hausse de $35\,\%$\end{minipage}\end{tcolorbox}\hfill\begin{tcolorbox}[enhanced,nobeforeafter,width=7.1cm,height=2.1cm,valign=center,colback=csta!6,colframe=csta,arc=2.5mm,boxrule=0.9pt,left=3mm]{\Large\bfseries\color{csta}d)}\quad\begin{minipage}[c]{5.3cm}\raggedright\Large une hausse de $25\,\%$\end{minipage}\end{tcolorbox}
\end{tcolorbox}
\newpage
\begin{tcolorbox}[enhanced,width=16cm,height=14cm,colback=white,colframe=csta,boxrule=2pt,arc=0mm,
left=6mm,right=6mm,top=0mm,bottom=4mm,sharp corners,
overlay={\fill[csta] (frame.north west) rectangle ([yshift=-1.25cm]frame.north east);
\node[anchor=west,text=white,font=\sffamily\bfseries\LARGE] at ([xshift=5mm,yshift=-0.62cm]frame.north west) {Question 26};
\node[anchor=east,text=white,font=\sffamily\itshape\large] at ([xshift=-5mm,yshift=-0.62cm]frame.north east) {Évolutions et statistiques};}]
\vspace*{1.75cm}
\begin{minipage}[t][6.9cm][c]{\linewidth}\LARGE\centering
\begin{minipage}{0.95\linewidth}\raggedright\linespread{1.25}\selectfont Après une baisse de $50\,\%$, quelle évolution faut-il appliquer pour revenir au prix de départ~?\end{minipage}
\end{minipage}

\vfill
\begin{tcolorbox}[enhanced,nobeforeafter,width=7.1cm,height=2.1cm,valign=center,colback=csta!6,colframe=csta,arc=2.5mm,boxrule=0.9pt,left=3mm]{\Large\bfseries\color{csta}a)}\quad\begin{minipage}[c]{5.3cm}\raggedright\Large une hausse de $50\,\%$\end{minipage}\end{tcolorbox}\hfill\begin{tcolorbox}[enhanced,nobeforeafter,width=7.1cm,height=2.1cm,valign=center,colback=csta!6,colframe=csta,arc=2.5mm,boxrule=0.9pt,left=3mm]{\Large\bfseries\color{csta}b)}\quad\begin{minipage}[c]{5.3cm}\raggedright\Large une hausse de $100\,\%$\end{minipage}\end{tcolorbox}\par\vspace{2.5mm}
\begin{tcolorbox}[enhanced,nobeforeafter,width=7.1cm,height=2.1cm,valign=center,colback=csta!6,colframe=csta,arc=2.5mm,boxrule=0.9pt,left=3mm]{\Large\bfseries\color{csta}c)}\quad\begin{minipage}[c]{5.3cm}\raggedright\Large une baisse de $100\,\%$\end{minipage}\end{tcolorbox}\hfill\begin{tcolorbox}[enhanced,nobeforeafter,width=7.1cm,height=2.1cm,valign=center,colback=csta!6,colframe=csta,arc=2.5mm,boxrule=0.9pt,left=3mm]{\Large\bfseries\color{csta}d)}\quad\begin{minipage}[c]{5.3cm}\raggedright\Large une hausse de $110\,\%$\end{minipage}\end{tcolorbox}
\end{tcolorbox}
\newpage
\begin{tcolorbox}[enhanced,width=16cm,height=14cm,colback=white,colframe=csta,boxrule=2pt,arc=0mm,
left=6mm,right=6mm,top=0mm,bottom=4mm,sharp corners,
overlay={\fill[csta] (frame.north west) rectangle ([yshift=-1.25cm]frame.north east);
\node[anchor=west,text=white,font=\sffamily\bfseries\LARGE] at ([xshift=5mm,yshift=-0.62cm]frame.north west) {Question 27};
\node[anchor=east,text=white,font=\sffamily\itshape\large] at ([xshift=-5mm,yshift=-0.62cm]frame.north east) {Évolutions et statistiques};}]
\vspace*{1.75cm}
\begin{minipage}[t][6.9cm][c]{\linewidth}\Large\centering
\begin{minipage}{0.95\linewidth}\raggedright\linespread{1.25}\selectfont Dans un lycée, $40\,\%$ des élèves sont internes et $25\,\%$ des internes viennent d'une autre commune. Quelle proportion des élèves du lycée sont des internes venant d'une autre commune~?\end{minipage}
\end{minipage}

\vfill
\begin{tcolorbox}[enhanced,nobeforeafter,width=7.1cm,height=2.1cm,valign=center,colback=csta!6,colframe=csta,arc=2.5mm,boxrule=0.9pt,left=3mm]{\Large\bfseries\color{csta}a)}\quad\begin{minipage}[c]{5.3cm}\raggedright\Large $15\,\%$\end{minipage}\end{tcolorbox}\hfill\begin{tcolorbox}[enhanced,nobeforeafter,width=7.1cm,height=2.1cm,valign=center,colback=csta!6,colframe=csta,arc=2.5mm,boxrule=0.9pt,left=3mm]{\Large\bfseries\color{csta}b)}\quad\begin{minipage}[c]{5.3cm}\raggedright\Large $10\,\%$\end{minipage}\end{tcolorbox}\par\vspace{2.5mm}
\begin{tcolorbox}[enhanced,nobeforeafter,width=7.1cm,height=2.1cm,valign=center,colback=csta!6,colframe=csta,arc=2.5mm,boxrule=0.9pt,left=3mm]{\Large\bfseries\color{csta}c)}\quad\begin{minipage}[c]{5.3cm}\raggedright\Large $32{,}5\,\%$\end{minipage}\end{tcolorbox}\hfill\begin{tcolorbox}[enhanced,nobeforeafter,width=7.1cm,height=2.1cm,valign=center,colback=csta!6,colframe=csta,arc=2.5mm,boxrule=0.9pt,left=3mm]{\Large\bfseries\color{csta}d)}\quad\begin{minipage}[c]{5.3cm}\raggedright\Large $65\,\%$\end{minipage}\end{tcolorbox}
\end{tcolorbox}
\newpage
\begin{tcolorbox}[enhanced,width=16cm,height=14cm,colback=white,colframe=csta,boxrule=2pt,arc=0mm,
left=6mm,right=6mm,top=0mm,bottom=4mm,sharp corners,
overlay={\fill[csta] (frame.north west) rectangle ([yshift=-1.25cm]frame.north east);
\node[anchor=west,text=white,font=\sffamily\bfseries\LARGE] at ([xshift=5mm,yshift=-0.62cm]frame.north west) {Question 28};
\node[anchor=east,text=white,font=\sffamily\itshape\large] at ([xshift=-5mm,yshift=-0.62cm]frame.north east) {Évolutions et statistiques};}]
\vspace*{1.75cm}
\begin{minipage}[t][6.9cm][c]{\linewidth}\Large\centering
\begin{minipage}{0.95\linewidth}\raggedright\linespread{1.25}\selectfont Dans un lycée, $60\,\%$ des élèves sont internes et $50\,\%$ des internes viennent d'une autre commune. Quelle proportion des élèves du lycée sont des internes venant d'une autre commune~?\end{minipage}
\end{minipage}

\vfill
\begin{tcolorbox}[enhanced,nobeforeafter,width=7.1cm,height=2.1cm,valign=center,colback=csta!6,colframe=csta,arc=2.5mm,boxrule=0.9pt,left=3mm]{\Large\bfseries\color{csta}a)}\quad\begin{minipage}[c]{5.3cm}\raggedright\Large $30\,\%$\end{minipage}\end{tcolorbox}\hfill\begin{tcolorbox}[enhanced,nobeforeafter,width=7.1cm,height=2.1cm,valign=center,colback=csta!6,colframe=csta,arc=2.5mm,boxrule=0.9pt,left=3mm]{\Large\bfseries\color{csta}b)}\quad\begin{minipage}[c]{5.3cm}\raggedright\Large $11\,\%$\end{minipage}\end{tcolorbox}\par\vspace{2.5mm}
\begin{tcolorbox}[enhanced,nobeforeafter,width=7.1cm,height=2.1cm,valign=center,colback=csta!6,colframe=csta,arc=2.5mm,boxrule=0.9pt,left=3mm]{\Large\bfseries\color{csta}c)}\quad\begin{minipage}[c]{5.3cm}\raggedright\Large $10\,\%$\end{minipage}\end{tcolorbox}\hfill\begin{tcolorbox}[enhanced,nobeforeafter,width=7.1cm,height=2.1cm,valign=center,colback=csta!6,colframe=csta,arc=2.5mm,boxrule=0.9pt,left=3mm]{\Large\bfseries\color{csta}d)}\quad\begin{minipage}[c]{5.3cm}\raggedright\Large $55\,\%$\end{minipage}\end{tcolorbox}
\end{tcolorbox}
\newpage
\begin{tcolorbox}[enhanced,width=16cm,height=14cm,colback=white,colframe=csta,boxrule=2pt,arc=0mm,
left=6mm,right=6mm,top=0mm,bottom=4mm,sharp corners,
overlay={\fill[csta] (frame.north west) rectangle ([yshift=-1.25cm]frame.north east);
\node[anchor=west,text=white,font=\sffamily\bfseries\LARGE] at ([xshift=5mm,yshift=-0.62cm]frame.north west) {Question 29};
\node[anchor=east,text=white,font=\sffamily\itshape\large] at ([xshift=-5mm,yshift=-0.62cm]frame.north east) {Évolutions et statistiques};}]
\vspace*{1.75cm}
\begin{minipage}[t][6.9cm][c]{\linewidth}\Large\centering
\begin{minipage}{0.95\linewidth}\raggedright\linespread{1.25}\selectfont Quelle est la moyenne de cette série statistique~? $\begin{array}{|l|c|c|c|}\hline \text{Valeur} & 8 & 12 & 15\\\hline \text{Effectif} & 2 & 5 & 3\\\hline\end{array}$\end{minipage}
\end{minipage}

\vfill
\begin{tcolorbox}[enhanced,nobeforeafter,width=7.1cm,height=2.1cm,valign=center,colback=csta!6,colframe=csta,arc=2.5mm,boxrule=0.9pt,left=3mm]{\Large\bfseries\color{csta}a)}\quad\begin{minipage}[c]{5.3cm}\raggedright\Large $\dfrac{35}{3}$\end{minipage}\end{tcolorbox}\hfill\begin{tcolorbox}[enhanced,nobeforeafter,width=7.1cm,height=2.1cm,valign=center,colback=csta!6,colframe=csta,arc=2.5mm,boxrule=0.9pt,left=3mm]{\Large\bfseries\color{csta}b)}\quad\begin{minipage}[c]{5.3cm}\raggedright\Large $12{,}1$\end{minipage}\end{tcolorbox}\par\vspace{2.5mm}
\begin{tcolorbox}[enhanced,nobeforeafter,width=7.1cm,height=2.1cm,valign=center,colback=csta!6,colframe=csta,arc=2.5mm,boxrule=0.9pt,left=3mm]{\Large\bfseries\color{csta}c)}\quad\begin{minipage}[c]{5.3cm}\raggedright\Large $13{,}1$\end{minipage}\end{tcolorbox}\hfill\begin{tcolorbox}[enhanced,nobeforeafter,width=7.1cm,height=2.1cm,valign=center,colback=csta!6,colframe=csta,arc=2.5mm,boxrule=0.9pt,left=3mm]{\Large\bfseries\color{csta}d)}\quad\begin{minipage}[c]{5.3cm}\raggedright\Large $\dfrac{121}{3}$\end{minipage}\end{tcolorbox}
\end{tcolorbox}
\newpage
\begin{tcolorbox}[enhanced,width=16cm,height=14cm,colback=white,colframe=csta,boxrule=2pt,arc=0mm,
left=6mm,right=6mm,top=0mm,bottom=4mm,sharp corners,
overlay={\fill[csta] (frame.north west) rectangle ([yshift=-1.25cm]frame.north east);
\node[anchor=west,text=white,font=\sffamily\bfseries\LARGE] at ([xshift=5mm,yshift=-0.62cm]frame.north west) {Question 30};
\node[anchor=east,text=white,font=\sffamily\itshape\large] at ([xshift=-5mm,yshift=-0.62cm]frame.north east) {Évolutions et statistiques};}]
\vspace*{1.75cm}
\begin{minipage}[t][6.9cm][c]{\linewidth}\Large\centering
\begin{minipage}{0.95\linewidth}\raggedright\linespread{1.25}\selectfont Quelle est la moyenne de cette série statistique~? $\begin{array}{|l|c|c|c|}\hline \text{Valeur} & 10 & 14 & 18\\\hline \text{Effectif} & 4 & 4 & 2\\\hline\end{array}$\end{minipage}
\end{minipage}

\vfill
\begin{tcolorbox}[enhanced,nobeforeafter,width=7.1cm,height=2.1cm,valign=center,colback=csta!6,colframe=csta,arc=2.5mm,boxrule=0.9pt,left=3mm]{\Large\bfseries\color{csta}a)}\quad\begin{minipage}[c]{5.3cm}\raggedright\Large $14$\end{minipage}\end{tcolorbox}\hfill\begin{tcolorbox}[enhanced,nobeforeafter,width=7.1cm,height=2.1cm,valign=center,colback=csta!6,colframe=csta,arc=2.5mm,boxrule=0.9pt,left=3mm]{\Large\bfseries\color{csta}b)}\quad\begin{minipage}[c]{5.3cm}\raggedright\Large $14{,}2$\end{minipage}\end{tcolorbox}\par\vspace{2.5mm}
\begin{tcolorbox}[enhanced,nobeforeafter,width=7.1cm,height=2.1cm,valign=center,colback=csta!6,colframe=csta,arc=2.5mm,boxrule=0.9pt,left=3mm]{\Large\bfseries\color{csta}c)}\quad\begin{minipage}[c]{5.3cm}\raggedright\Large $13{,}2$\end{minipage}\end{tcolorbox}\hfill\begin{tcolorbox}[enhanced,nobeforeafter,width=7.1cm,height=2.1cm,valign=center,colback=csta!6,colframe=csta,arc=2.5mm,boxrule=0.9pt,left=3mm]{\Large\bfseries\color{csta}d)}\quad\begin{minipage}[c]{5.3cm}\raggedright\Large $44$\end{minipage}\end{tcolorbox}
\end{tcolorbox}
\newpage
\begin{tcolorbox}[enhanced,width=16cm,height=14cm,colback=white,colframe=csta,boxrule=2pt,arc=0mm,
left=6mm,right=6mm,top=0mm,bottom=4mm,sharp corners,
overlay={\fill[csta] (frame.north west) rectangle ([yshift=-1.25cm]frame.north east);
\node[anchor=west,text=white,font=\sffamily\bfseries\LARGE] at ([xshift=5mm,yshift=-0.62cm]frame.north west) {Question 31};
\node[anchor=east,text=white,font=\sffamily\itshape\large] at ([xshift=-5mm,yshift=-0.62cm]frame.north east) {Évolutions et statistiques};}]
\vspace*{1.75cm}
\begin{minipage}[t][6.9cm][c]{\linewidth}\LARGE\centering
\begin{minipage}{0.95\linewidth}\raggedright\linespread{1.25}\selectfont Voici une série rangée dans l'ordre croissant~: $3\ ;\ 5\ ;\ 7\ ;\ 8\ ;\ 10\ ;\ 12\ ;\ 15\ ;\ 18$. Quel est son premier quartile $Q_1$~?\end{minipage}
\end{minipage}

\vfill
\begin{tcolorbox}[enhanced,nobeforeafter,width=7.1cm,height=2.1cm,valign=center,colback=csta!6,colframe=csta,arc=2.5mm,boxrule=0.9pt,left=3mm]{\Large\bfseries\color{csta}a)}\quad\begin{minipage}[c]{5.3cm}\raggedright\Large $5$\end{minipage}\end{tcolorbox}\hfill\begin{tcolorbox}[enhanced,nobeforeafter,width=7.1cm,height=2.1cm,valign=center,colback=csta!6,colframe=csta,arc=2.5mm,boxrule=0.9pt,left=3mm]{\Large\bfseries\color{csta}b)}\quad\begin{minipage}[c]{5.3cm}\raggedright\Large $7$\end{minipage}\end{tcolorbox}\par\vspace{2.5mm}
\begin{tcolorbox}[enhanced,nobeforeafter,width=7.1cm,height=2.1cm,valign=center,colback=csta!6,colframe=csta,arc=2.5mm,boxrule=0.9pt,left=3mm]{\Large\bfseries\color{csta}c)}\quad\begin{minipage}[c]{5.3cm}\raggedright\Large $12$\end{minipage}\end{tcolorbox}\hfill\begin{tcolorbox}[enhanced,nobeforeafter,width=7.1cm,height=2.1cm,valign=center,colback=csta!6,colframe=csta,arc=2.5mm,boxrule=0.9pt,left=3mm]{\Large\bfseries\color{csta}d)}\quad\begin{minipage}[c]{5.3cm}\raggedright\Large $9$\end{minipage}\end{tcolorbox}
\end{tcolorbox}
\newpage
\begin{tcolorbox}[enhanced,width=16cm,height=14cm,colback=white,colframe=csta,boxrule=2pt,arc=0mm,
left=6mm,right=6mm,top=0mm,bottom=4mm,sharp corners,
overlay={\fill[csta] (frame.north west) rectangle ([yshift=-1.25cm]frame.north east);
\node[anchor=west,text=white,font=\sffamily\bfseries\LARGE] at ([xshift=5mm,yshift=-0.62cm]frame.north west) {Question 32};
\node[anchor=east,text=white,font=\sffamily\itshape\large] at ([xshift=-5mm,yshift=-0.62cm]frame.north east) {Évolutions et statistiques};}]
\vspace*{1.75cm}
\begin{minipage}[t][6.9cm][c]{\linewidth}\LARGE\centering
\begin{minipage}{0.95\linewidth}\raggedright\linespread{1.25}\selectfont Voici une série rangée dans l'ordre croissant~: $2\ ;\ 4\ ;\ 4\ ;\ 6\ ;\ 9\ ;\ 11\ ;\ 13\ ;\ 14\ ;\ 20$. Quel est son premier quartile $Q_1$~?\end{minipage}
\end{minipage}

\vfill
\begin{tcolorbox}[enhanced,nobeforeafter,width=7.1cm,height=2.1cm,valign=center,colback=csta!6,colframe=csta,arc=2.5mm,boxrule=0.9pt,left=3mm]{\Large\bfseries\color{csta}a)}\quad\begin{minipage}[c]{5.3cm}\raggedright\Large $9$\end{minipage}\end{tcolorbox}\hfill\begin{tcolorbox}[enhanced,nobeforeafter,width=7.1cm,height=2.1cm,valign=center,colback=csta!6,colframe=csta,arc=2.5mm,boxrule=0.9pt,left=3mm]{\Large\bfseries\color{csta}b)}\quad\begin{minipage}[c]{5.3cm}\raggedright\Large $13$\end{minipage}\end{tcolorbox}\par\vspace{2.5mm}
\begin{tcolorbox}[enhanced,nobeforeafter,width=7.1cm,height=2.1cm,valign=center,colback=csta!6,colframe=csta,arc=2.5mm,boxrule=0.9pt,left=3mm]{\Large\bfseries\color{csta}c)}\quad\begin{minipage}[c]{5.3cm}\raggedright\Large $6$\end{minipage}\end{tcolorbox}\hfill\begin{tcolorbox}[enhanced,nobeforeafter,width=7.1cm,height=2.1cm,valign=center,colback=csta!6,colframe=csta,arc=2.5mm,boxrule=0.9pt,left=3mm]{\Large\bfseries\color{csta}d)}\quad\begin{minipage}[c]{5.3cm}\raggedright\Large $4$\end{minipage}\end{tcolorbox}
\end{tcolorbox}
\newpage
\begin{tcolorbox}[enhanced,width=16cm,height=14cm,colback=white,colframe=cfon,boxrule=2pt,arc=0mm,
left=6mm,right=6mm,top=0mm,bottom=4mm,sharp corners,
overlay={\fill[cfon] (frame.north west) rectangle ([yshift=-1.25cm]frame.north east);
\node[anchor=west,text=white,font=\sffamily\bfseries\LARGE] at ([xshift=5mm,yshift=-0.62cm]frame.north west) {Question 33};
\node[anchor=east,text=white,font=\sffamily\itshape\large] at ([xshift=-5mm,yshift=-0.62cm]frame.north east) {Fonctions};}]
\vspace*{1.75cm}
\begin{minipage}[t][6.9cm][c]{\linewidth}\LARGE\centering
\begin{minipage}{0.95\linewidth}\raggedright\linespread{1.25}\selectfont Soit $f(x)=2x^2-3x+1$. Que vaut $f(-2)$~?\end{minipage}
\end{minipage}

\vfill
\begin{tcolorbox}[enhanced,nobeforeafter,width=7.1cm,height=2.1cm,valign=center,colback=cfon!6,colframe=cfon,arc=2.5mm,boxrule=0.9pt,left=3mm]{\Large\bfseries\color{cfon}a)}\quad\begin{minipage}[c]{5.3cm}\raggedright\Large $15$\end{minipage}\end{tcolorbox}\hfill\begin{tcolorbox}[enhanced,nobeforeafter,width=7.1cm,height=2.1cm,valign=center,colback=cfon!6,colframe=cfon,arc=2.5mm,boxrule=0.9pt,left=3mm]{\Large\bfseries\color{cfon}b)}\quad\begin{minipage}[c]{5.3cm}\raggedright\Large $-15$\end{minipage}\end{tcolorbox}\par\vspace{2.5mm}
\begin{tcolorbox}[enhanced,nobeforeafter,width=7.1cm,height=2.1cm,valign=center,colback=cfon!6,colframe=cfon,arc=2.5mm,boxrule=0.9pt,left=3mm]{\Large\bfseries\color{cfon}c)}\quad\begin{minipage}[c]{5.3cm}\raggedright\Large $-1$\end{minipage}\end{tcolorbox}\hfill\begin{tcolorbox}[enhanced,nobeforeafter,width=7.1cm,height=2.1cm,valign=center,colback=cfon!6,colframe=cfon,arc=2.5mm,boxrule=0.9pt,left=3mm]{\Large\bfseries\color{cfon}d)}\quad\begin{minipage}[c]{5.3cm}\raggedright\Large $17$\end{minipage}\end{tcolorbox}
\end{tcolorbox}
\newpage
\begin{tcolorbox}[enhanced,width=16cm,height=14cm,colback=white,colframe=cfon,boxrule=2pt,arc=0mm,
left=6mm,right=6mm,top=0mm,bottom=4mm,sharp corners,
overlay={\fill[cfon] (frame.north west) rectangle ([yshift=-1.25cm]frame.north east);
\node[anchor=west,text=white,font=\sffamily\bfseries\LARGE] at ([xshift=5mm,yshift=-0.62cm]frame.north west) {Question 34};
\node[anchor=east,text=white,font=\sffamily\itshape\large] at ([xshift=-5mm,yshift=-0.62cm]frame.north east) {Fonctions};}]
\vspace*{1.75cm}
\begin{minipage}[t][6.9cm][c]{\linewidth}\LARGE\centering
\begin{minipage}{0.95\linewidth}\raggedright\linespread{1.25}\selectfont Soit $f(x)=x^2+4x-5$. Que vaut $f(3)$~?\end{minipage}
\end{minipage}

\vfill
\begin{tcolorbox}[enhanced,nobeforeafter,width=7.1cm,height=2.1cm,valign=center,colback=cfon!6,colframe=cfon,arc=2.5mm,boxrule=0.9pt,left=3mm]{\Large\bfseries\color{cfon}a)}\quad\begin{minipage}[c]{5.3cm}\raggedright\Large $-16$\end{minipage}\end{tcolorbox}\hfill\begin{tcolorbox}[enhanced,nobeforeafter,width=7.1cm,height=2.1cm,valign=center,colback=cfon!6,colframe=cfon,arc=2.5mm,boxrule=0.9pt,left=3mm]{\Large\bfseries\color{cfon}b)}\quad\begin{minipage}[c]{5.3cm}\raggedright\Large $16$\end{minipage}\end{tcolorbox}\par\vspace{2.5mm}
\begin{tcolorbox}[enhanced,nobeforeafter,width=7.1cm,height=2.1cm,valign=center,colback=cfon!6,colframe=cfon,arc=2.5mm,boxrule=0.9pt,left=3mm]{\Large\bfseries\color{cfon}c)}\quad\begin{minipage}[c]{5.3cm}\raggedright\Large $13$\end{minipage}\end{tcolorbox}\hfill\begin{tcolorbox}[enhanced,nobeforeafter,width=7.1cm,height=2.1cm,valign=center,colback=cfon!6,colframe=cfon,arc=2.5mm,boxrule=0.9pt,left=3mm]{\Large\bfseries\color{cfon}d)}\quad\begin{minipage}[c]{5.3cm}\raggedright\Large $-2$\end{minipage}\end{tcolorbox}
\end{tcolorbox}
\newpage
\begin{tcolorbox}[enhanced,width=16cm,height=14cm,colback=white,colframe=cfon,boxrule=2pt,arc=0mm,
left=6mm,right=6mm,top=0mm,bottom=4mm,sharp corners,
overlay={\fill[cfon] (frame.north west) rectangle ([yshift=-1.25cm]frame.north east);
\node[anchor=west,text=white,font=\sffamily\bfseries\LARGE] at ([xshift=5mm,yshift=-0.62cm]frame.north west) {Question 35};
\node[anchor=east,text=white,font=\sffamily\itshape\large] at ([xshift=-5mm,yshift=-0.62cm]frame.north east) {Fonctions};}]
\vspace*{1.75cm}
\begin{minipage}[t][6.9cm][c]{\linewidth}\LARGE\centering
\begin{minipage}{0.95\linewidth}\raggedright\linespread{1.25}\selectfont Soit $f(x)=2x+3$. Quel est l'antécédent de $11$ par $f$~?\end{minipage}
\end{minipage}

\vfill
\begin{tcolorbox}[enhanced,nobeforeafter,width=7.1cm,height=2.1cm,valign=center,colback=cfon!6,colframe=cfon,arc=2.5mm,boxrule=0.9pt,left=3mm]{\Large\bfseries\color{cfon}a)}\quad\begin{minipage}[c]{5.3cm}\raggedright\Large $7$\end{minipage}\end{tcolorbox}\hfill\begin{tcolorbox}[enhanced,nobeforeafter,width=7.1cm,height=2.1cm,valign=center,colback=cfon!6,colframe=cfon,arc=2.5mm,boxrule=0.9pt,left=3mm]{\Large\bfseries\color{cfon}b)}\quad\begin{minipage}[c]{5.3cm}\raggedright\Large $16$\end{minipage}\end{tcolorbox}\par\vspace{2.5mm}
\begin{tcolorbox}[enhanced,nobeforeafter,width=7.1cm,height=2.1cm,valign=center,colback=cfon!6,colframe=cfon,arc=2.5mm,boxrule=0.9pt,left=3mm]{\Large\bfseries\color{cfon}c)}\quad\begin{minipage}[c]{5.3cm}\raggedright\Large $25$\end{minipage}\end{tcolorbox}\hfill\begin{tcolorbox}[enhanced,nobeforeafter,width=7.1cm,height=2.1cm,valign=center,colback=cfon!6,colframe=cfon,arc=2.5mm,boxrule=0.9pt,left=3mm]{\Large\bfseries\color{cfon}d)}\quad\begin{minipage}[c]{5.3cm}\raggedright\Large $4$\end{minipage}\end{tcolorbox}
\end{tcolorbox}
\newpage
\begin{tcolorbox}[enhanced,width=16cm,height=14cm,colback=white,colframe=cfon,boxrule=2pt,arc=0mm,
left=6mm,right=6mm,top=0mm,bottom=4mm,sharp corners,
overlay={\fill[cfon] (frame.north west) rectangle ([yshift=-1.25cm]frame.north east);
\node[anchor=west,text=white,font=\sffamily\bfseries\LARGE] at ([xshift=5mm,yshift=-0.62cm]frame.north west) {Question 36};
\node[anchor=east,text=white,font=\sffamily\itshape\large] at ([xshift=-5mm,yshift=-0.62cm]frame.north east) {Fonctions};}]
\vspace*{1.75cm}
\begin{minipage}[t][6.9cm][c]{\linewidth}\LARGE\centering
\begin{minipage}{0.95\linewidth}\raggedright\linespread{1.25}\selectfont Soit $f(x)=-3x+5$. Quel est l'antécédent de $-1$ par $f$~?\end{minipage}
\end{minipage}

\vfill
\begin{tcolorbox}[enhanced,nobeforeafter,width=7.1cm,height=2.1cm,valign=center,colback=cfon!6,colframe=cfon,arc=2.5mm,boxrule=0.9pt,left=3mm]{\Large\bfseries\color{cfon}a)}\quad\begin{minipage}[c]{5.3cm}\raggedright\Large $18$\end{minipage}\end{tcolorbox}\hfill\begin{tcolorbox}[enhanced,nobeforeafter,width=7.1cm,height=2.1cm,valign=center,colback=cfon!6,colframe=cfon,arc=2.5mm,boxrule=0.9pt,left=3mm]{\Large\bfseries\color{cfon}b)}\quad\begin{minipage}[c]{5.3cm}\raggedright\Large $8$\end{minipage}\end{tcolorbox}\par\vspace{2.5mm}
\begin{tcolorbox}[enhanced,nobeforeafter,width=7.1cm,height=2.1cm,valign=center,colback=cfon!6,colframe=cfon,arc=2.5mm,boxrule=0.9pt,left=3mm]{\Large\bfseries\color{cfon}c)}\quad\begin{minipage}[c]{5.3cm}\raggedright\Large $-\dfrac{4}{3}$\end{minipage}\end{tcolorbox}\hfill\begin{tcolorbox}[enhanced,nobeforeafter,width=7.1cm,height=2.1cm,valign=center,colback=cfon!6,colframe=cfon,arc=2.5mm,boxrule=0.9pt,left=3mm]{\Large\bfseries\color{cfon}d)}\quad\begin{minipage}[c]{5.3cm}\raggedright\Large $2$\end{minipage}\end{tcolorbox}
\end{tcolorbox}
\newpage
\begin{tcolorbox}[enhanced,width=16cm,height=14cm,colback=white,colframe=cfon,boxrule=2pt,arc=0mm,
left=6mm,right=6mm,top=0mm,bottom=4mm,sharp corners,
overlay={\fill[cfon] (frame.north west) rectangle ([yshift=-1.25cm]frame.north east);
\node[anchor=west,text=white,font=\sffamily\bfseries\LARGE] at ([xshift=5mm,yshift=-0.62cm]frame.north west) {Question 37};
\node[anchor=east,text=white,font=\sffamily\itshape\large] at ([xshift=-5mm,yshift=-0.62cm]frame.north east) {Fonctions};}]
\vspace*{1.75cm}
\begin{minipage}[t][6.9cm][c]{\linewidth}\LARGE\centering
\begin{minipage}{0.95\linewidth}\raggedright\linespread{1.25}\selectfont Quel est le coefficient directeur de la droite $(AB)$, avec $A(1\,;\,2)$ et $B(3\,;\,8)$~?\end{minipage}
\end{minipage}

\vfill
\begin{tcolorbox}[enhanced,nobeforeafter,width=7.1cm,height=2.1cm,valign=center,colback=cfon!6,colframe=cfon,arc=2.5mm,boxrule=0.9pt,left=3mm]{\Large\bfseries\color{cfon}a)}\quad\begin{minipage}[c]{5.3cm}\raggedright\Large $-3$\end{minipage}\end{tcolorbox}\hfill\begin{tcolorbox}[enhanced,nobeforeafter,width=7.1cm,height=2.1cm,valign=center,colback=cfon!6,colframe=cfon,arc=2.5mm,boxrule=0.9pt,left=3mm]{\Large\bfseries\color{cfon}b)}\quad\begin{minipage}[c]{5.3cm}\raggedright\Large $\dfrac{5}{2}$\end{minipage}\end{tcolorbox}\par\vspace{2.5mm}
\begin{tcolorbox}[enhanced,nobeforeafter,width=7.1cm,height=2.1cm,valign=center,colback=cfon!6,colframe=cfon,arc=2.5mm,boxrule=0.9pt,left=3mm]{\Large\bfseries\color{cfon}c)}\quad\begin{minipage}[c]{5.3cm}\raggedright\Large $\dfrac{1}{3}$\end{minipage}\end{tcolorbox}\hfill\begin{tcolorbox}[enhanced,nobeforeafter,width=7.1cm,height=2.1cm,valign=center,colback=cfon!6,colframe=cfon,arc=2.5mm,boxrule=0.9pt,left=3mm]{\Large\bfseries\color{cfon}d)}\quad\begin{minipage}[c]{5.3cm}\raggedright\Large $3$\end{minipage}\end{tcolorbox}
\end{tcolorbox}
\newpage
\begin{tcolorbox}[enhanced,width=16cm,height=14cm,colback=white,colframe=cfon,boxrule=2pt,arc=0mm,
left=6mm,right=6mm,top=0mm,bottom=4mm,sharp corners,
overlay={\fill[cfon] (frame.north west) rectangle ([yshift=-1.25cm]frame.north east);
\node[anchor=west,text=white,font=\sffamily\bfseries\LARGE] at ([xshift=5mm,yshift=-0.62cm]frame.north west) {Question 38};
\node[anchor=east,text=white,font=\sffamily\itshape\large] at ([xshift=-5mm,yshift=-0.62cm]frame.north east) {Fonctions};}]
\vspace*{1.75cm}
\begin{minipage}[t][6.9cm][c]{\linewidth}\LARGE\centering
\begin{minipage}{0.95\linewidth}\raggedright\linespread{1.25}\selectfont Quel est le coefficient directeur de la droite $(AB)$, avec $A(-2\,;\,5)$ et $B(2\,;\,-3)$~?\end{minipage}
\end{minipage}

\vfill
\begin{tcolorbox}[enhanced,nobeforeafter,width=7.1cm,height=2.1cm,valign=center,colback=cfon!6,colframe=cfon,arc=2.5mm,boxrule=0.9pt,left=3mm]{\Large\bfseries\color{cfon}a)}\quad\begin{minipage}[c]{5.3cm}\raggedright\Large $-\dfrac{1}{2}$\end{minipage}\end{tcolorbox}\hfill\begin{tcolorbox}[enhanced,nobeforeafter,width=7.1cm,height=2.1cm,valign=center,colback=cfon!6,colframe=cfon,arc=2.5mm,boxrule=0.9pt,left=3mm]{\Large\bfseries\color{cfon}b)}\quad\begin{minipage}[c]{5.3cm}\raggedright\Large $-2$\end{minipage}\end{tcolorbox}\par\vspace{2.5mm}
\begin{tcolorbox}[enhanced,nobeforeafter,width=7.1cm,height=2.1cm,valign=center,colback=cfon!6,colframe=cfon,arc=2.5mm,boxrule=0.9pt,left=3mm]{\Large\bfseries\color{cfon}c)}\quad\begin{minipage}[c]{5.3cm}\raggedright\Large $-1$\end{minipage}\end{tcolorbox}\hfill\begin{tcolorbox}[enhanced,nobeforeafter,width=7.1cm,height=2.1cm,valign=center,colback=cfon!6,colframe=cfon,arc=2.5mm,boxrule=0.9pt,left=3mm]{\Large\bfseries\color{cfon}d)}\quad\begin{minipage}[c]{5.3cm}\raggedright\Large $2$\end{minipage}\end{tcolorbox}
\end{tcolorbox}
\newpage
\begin{tcolorbox}[enhanced,width=16cm,height=14cm,colback=white,colframe=cfon,boxrule=2pt,arc=0mm,
left=6mm,right=6mm,top=0mm,bottom=4mm,sharp corners,
overlay={\fill[cfon] (frame.north west) rectangle ([yshift=-1.25cm]frame.north east);
\node[anchor=west,text=white,font=\sffamily\bfseries\LARGE] at ([xshift=5mm,yshift=-0.62cm]frame.north west) {Question 39};
\node[anchor=east,text=white,font=\sffamily\itshape\large] at ([xshift=-5mm,yshift=-0.62cm]frame.north east) {Fonctions};}]
\vspace*{1.75cm}
\begin{minipage}[t][6.9cm][c]{\linewidth}\LARGE\centering
\begin{minipage}{0.95\linewidth}\raggedright\linespread{1.25}\selectfont Quelle est l'équation réduite de la droite de coefficient directeur $2$ passant par $A(1\,;\,5)$~?\end{minipage}
\end{minipage}

\vfill
\begin{tcolorbox}[enhanced,nobeforeafter,width=7.1cm,height=2.1cm,valign=center,colback=cfon!6,colframe=cfon,arc=2.5mm,boxrule=0.9pt,left=3mm]{\Large\bfseries\color{cfon}a)}\quad\begin{minipage}[c]{5.3cm}\raggedright\Large $y=x+3$\end{minipage}\end{tcolorbox}\hfill\begin{tcolorbox}[enhanced,nobeforeafter,width=7.1cm,height=2.1cm,valign=center,colback=cfon!6,colframe=cfon,arc=2.5mm,boxrule=0.9pt,left=3mm]{\Large\bfseries\color{cfon}b)}\quad\begin{minipage}[c]{5.3cm}\raggedright\Large $y=2x+7$\end{minipage}\end{tcolorbox}\par\vspace{2.5mm}
\begin{tcolorbox}[enhanced,nobeforeafter,width=7.1cm,height=2.1cm,valign=center,colback=cfon!6,colframe=cfon,arc=2.5mm,boxrule=0.9pt,left=3mm]{\Large\bfseries\color{cfon}c)}\quad\begin{minipage}[c]{5.3cm}\raggedright\Large $y=2x+5$\end{minipage}\end{tcolorbox}\hfill\begin{tcolorbox}[enhanced,nobeforeafter,width=7.1cm,height=2.1cm,valign=center,colback=cfon!6,colframe=cfon,arc=2.5mm,boxrule=0.9pt,left=3mm]{\Large\bfseries\color{cfon}d)}\quad\begin{minipage}[c]{5.3cm}\raggedright\Large $y=2x+3$\end{minipage}\end{tcolorbox}
\end{tcolorbox}
\newpage
\begin{tcolorbox}[enhanced,width=16cm,height=14cm,colback=white,colframe=cfon,boxrule=2pt,arc=0mm,
left=6mm,right=6mm,top=0mm,bottom=4mm,sharp corners,
overlay={\fill[cfon] (frame.north west) rectangle ([yshift=-1.25cm]frame.north east);
\node[anchor=west,text=white,font=\sffamily\bfseries\LARGE] at ([xshift=5mm,yshift=-0.62cm]frame.north west) {Question 40};
\node[anchor=east,text=white,font=\sffamily\itshape\large] at ([xshift=-5mm,yshift=-0.62cm]frame.north east) {Fonctions};}]
\vspace*{1.75cm}
\begin{minipage}[t][6.9cm][c]{\linewidth}\LARGE\centering
\begin{minipage}{0.95\linewidth}\raggedright\linespread{1.25}\selectfont Quelle est l'équation réduite de la droite de coefficient directeur $-3$ passant par $A(2\,;\,1)$~?\end{minipage}
\end{minipage}

\vfill
\begin{tcolorbox}[enhanced,nobeforeafter,width=7.1cm,height=2.1cm,valign=center,colback=cfon!6,colframe=cfon,arc=2.5mm,boxrule=0.9pt,left=3mm]{\Large\bfseries\color{cfon}a)}\quad\begin{minipage}[c]{5.3cm}\raggedright\Large $y=-3x+7$\end{minipage}\end{tcolorbox}\hfill\begin{tcolorbox}[enhanced,nobeforeafter,width=7.1cm,height=2.1cm,valign=center,colback=cfon!6,colframe=cfon,arc=2.5mm,boxrule=0.9pt,left=3mm]{\Large\bfseries\color{cfon}b)}\quad\begin{minipage}[c]{5.3cm}\raggedright\Large $y=-3x-5$\end{minipage}\end{tcolorbox}\par\vspace{2.5mm}
\begin{tcolorbox}[enhanced,nobeforeafter,width=7.1cm,height=2.1cm,valign=center,colback=cfon!6,colframe=cfon,arc=2.5mm,boxrule=0.9pt,left=3mm]{\Large\bfseries\color{cfon}c)}\quad\begin{minipage}[c]{5.3cm}\raggedright\Large $y=-3x+1$\end{minipage}\end{tcolorbox}\hfill\begin{tcolorbox}[enhanced,nobeforeafter,width=7.1cm,height=2.1cm,valign=center,colback=cfon!6,colframe=cfon,arc=2.5mm,boxrule=0.9pt,left=3mm]{\Large\bfseries\color{cfon}d)}\quad\begin{minipage}[c]{5.3cm}\raggedright\Large $y=2x+7$\end{minipage}\end{tcolorbox}
\end{tcolorbox}
\newpage
\begin{tcolorbox}[enhanced,width=16cm,height=14cm,colback=white,colframe=cfon,boxrule=2pt,arc=0mm,
left=6mm,right=6mm,top=0mm,bottom=4mm,sharp corners,
overlay={\fill[cfon] (frame.north west) rectangle ([yshift=-1.25cm]frame.north east);
\node[anchor=west,text=white,font=\sffamily\bfseries\LARGE] at ([xshift=5mm,yshift=-0.62cm]frame.north west) {Question 41};
\node[anchor=east,text=white,font=\sffamily\itshape\large] at ([xshift=-5mm,yshift=-0.62cm]frame.north east) {Fonctions};}]
\vspace*{1.75cm}
\begin{minipage}[t][6.9cm][c]{\linewidth}\LARGE\centering
\begin{minipage}{0.95\linewidth}\raggedright\linespread{1.25}\selectfont Pour quels réels $x$ a-t-on $2x-6>0$~?\end{minipage}
\end{minipage}

\vfill
\begin{tcolorbox}[enhanced,nobeforeafter,width=7.1cm,height=2.1cm,valign=center,colback=cfon!6,colframe=cfon,arc=2.5mm,boxrule=0.9pt,left=3mm]{\Large\bfseries\color{cfon}a)}\quad\begin{minipage}[c]{5.3cm}\raggedright\Large $x<-3$\end{minipage}\end{tcolorbox}\hfill\begin{tcolorbox}[enhanced,nobeforeafter,width=7.1cm,height=2.1cm,valign=center,colback=cfon!6,colframe=cfon,arc=2.5mm,boxrule=0.9pt,left=3mm]{\Large\bfseries\color{cfon}b)}\quad\begin{minipage}[c]{5.3cm}\raggedright\Large $x>3$\end{minipage}\end{tcolorbox}\par\vspace{2.5mm}
\begin{tcolorbox}[enhanced,nobeforeafter,width=7.1cm,height=2.1cm,valign=center,colback=cfon!6,colframe=cfon,arc=2.5mm,boxrule=0.9pt,left=3mm]{\Large\bfseries\color{cfon}c)}\quad\begin{minipage}[c]{5.3cm}\raggedright\Large $x<3$\end{minipage}\end{tcolorbox}\hfill\begin{tcolorbox}[enhanced,nobeforeafter,width=7.1cm,height=2.1cm,valign=center,colback=cfon!6,colframe=cfon,arc=2.5mm,boxrule=0.9pt,left=3mm]{\Large\bfseries\color{cfon}d)}\quad\begin{minipage}[c]{5.3cm}\raggedright\Large $x>-3$\end{minipage}\end{tcolorbox}
\end{tcolorbox}
\newpage
\begin{tcolorbox}[enhanced,width=16cm,height=14cm,colback=white,colframe=cfon,boxrule=2pt,arc=0mm,
left=6mm,right=6mm,top=0mm,bottom=4mm,sharp corners,
overlay={\fill[cfon] (frame.north west) rectangle ([yshift=-1.25cm]frame.north east);
\node[anchor=west,text=white,font=\sffamily\bfseries\LARGE] at ([xshift=5mm,yshift=-0.62cm]frame.north west) {Question 42};
\node[anchor=east,text=white,font=\sffamily\itshape\large] at ([xshift=-5mm,yshift=-0.62cm]frame.north east) {Fonctions};}]
\vspace*{1.75cm}
\begin{minipage}[t][6.9cm][c]{\linewidth}\LARGE\centering
\begin{minipage}{0.95\linewidth}\raggedright\linespread{1.25}\selectfont Pour quels réels $x$ a-t-on $-3x+12>0$~?\end{minipage}
\end{minipage}

\vfill
\begin{tcolorbox}[enhanced,nobeforeafter,width=7.1cm,height=2.1cm,valign=center,colback=cfon!6,colframe=cfon,arc=2.5mm,boxrule=0.9pt,left=3mm]{\Large\bfseries\color{cfon}a)}\quad\begin{minipage}[c]{5.3cm}\raggedright\Large $x>4$\end{minipage}\end{tcolorbox}\hfill\begin{tcolorbox}[enhanced,nobeforeafter,width=7.1cm,height=2.1cm,valign=center,colback=cfon!6,colframe=cfon,arc=2.5mm,boxrule=0.9pt,left=3mm]{\Large\bfseries\color{cfon}b)}\quad\begin{minipage}[c]{5.3cm}\raggedright\Large $x<4$\end{minipage}\end{tcolorbox}\par\vspace{2.5mm}
\begin{tcolorbox}[enhanced,nobeforeafter,width=7.1cm,height=2.1cm,valign=center,colback=cfon!6,colframe=cfon,arc=2.5mm,boxrule=0.9pt,left=3mm]{\Large\bfseries\color{cfon}c)}\quad\begin{minipage}[c]{5.3cm}\raggedright\Large $x<-4$\end{minipage}\end{tcolorbox}\hfill\begin{tcolorbox}[enhanced,nobeforeafter,width=7.1cm,height=2.1cm,valign=center,colback=cfon!6,colframe=cfon,arc=2.5mm,boxrule=0.9pt,left=3mm]{\Large\bfseries\color{cfon}d)}\quad\begin{minipage}[c]{5.3cm}\raggedright\Large $x>-4$\end{minipage}\end{tcolorbox}
\end{tcolorbox}
\newpage
\begin{tcolorbox}[enhanced,width=16cm,height=14cm,colback=white,colframe=cfon,boxrule=2pt,arc=0mm,
left=6mm,right=6mm,top=0mm,bottom=4mm,sharp corners,
overlay={\fill[cfon] (frame.north west) rectangle ([yshift=-1.25cm]frame.north east);
\node[anchor=west,text=white,font=\sffamily\bfseries\LARGE] at ([xshift=5mm,yshift=-0.62cm]frame.north west) {Question 43};
\node[anchor=east,text=white,font=\sffamily\itshape\large] at ([xshift=-5mm,yshift=-0.62cm]frame.north east) {Fonctions};}]
\vspace*{1.75cm}
\begin{minipage}[t][6.9cm][c]{\linewidth}\LARGE\centering
\begin{minipage}{0.95\linewidth}\raggedright\linespread{1.25}\selectfont Quelles sont les solutions de l'équation $x^2=9$~?\end{minipage}
\end{minipage}

\vfill
\begin{tcolorbox}[enhanced,nobeforeafter,width=7.1cm,height=2.1cm,valign=center,colback=cfon!6,colframe=cfon,arc=2.5mm,boxrule=0.9pt,left=3mm]{\Large\bfseries\color{cfon}a)}\quad\begin{minipage}[c]{5.3cm}\raggedright\Large $\{3\}$\end{minipage}\end{tcolorbox}\hfill\begin{tcolorbox}[enhanced,nobeforeafter,width=7.1cm,height=2.1cm,valign=center,colback=cfon!6,colframe=cfon,arc=2.5mm,boxrule=0.9pt,left=3mm]{\Large\bfseries\color{cfon}b)}\quad\begin{minipage}[c]{5.3cm}\raggedright\Large $\{-9\,;\,9\}$\end{minipage}\end{tcolorbox}\par\vspace{2.5mm}
\begin{tcolorbox}[enhanced,nobeforeafter,width=7.1cm,height=2.1cm,valign=center,colback=cfon!6,colframe=cfon,arc=2.5mm,boxrule=0.9pt,left=3mm]{\Large\bfseries\color{cfon}c)}\quad\begin{minipage}[c]{5.3cm}\raggedright\Large $\{-3\,;\,3\}$\end{minipage}\end{tcolorbox}\hfill\begin{tcolorbox}[enhanced,nobeforeafter,width=7.1cm,height=2.1cm,valign=center,colback=cfon!6,colframe=cfon,arc=2.5mm,boxrule=0.9pt,left=3mm]{\Large\bfseries\color{cfon}d)}\quad\begin{minipage}[c]{5.3cm}\raggedright\Large aucune solution\end{minipage}\end{tcolorbox}
\end{tcolorbox}
\newpage
\begin{tcolorbox}[enhanced,width=16cm,height=14cm,colback=white,colframe=cfon,boxrule=2pt,arc=0mm,
left=6mm,right=6mm,top=0mm,bottom=4mm,sharp corners,
overlay={\fill[cfon] (frame.north west) rectangle ([yshift=-1.25cm]frame.north east);
\node[anchor=west,text=white,font=\sffamily\bfseries\LARGE] at ([xshift=5mm,yshift=-0.62cm]frame.north west) {Question 44};
\node[anchor=east,text=white,font=\sffamily\itshape\large] at ([xshift=-5mm,yshift=-0.62cm]frame.north east) {Fonctions};}]
\vspace*{1.75cm}
\begin{minipage}[t][6.9cm][c]{\linewidth}\LARGE\centering
\begin{minipage}{0.95\linewidth}\raggedright\linespread{1.25}\selectfont Quelles sont les solutions de l'équation $x^2=25$~?\end{minipage}
\end{minipage}

\vfill
\begin{tcolorbox}[enhanced,nobeforeafter,width=7.1cm,height=2.1cm,valign=center,colback=cfon!6,colframe=cfon,arc=2.5mm,boxrule=0.9pt,left=3mm]{\Large\bfseries\color{cfon}a)}\quad\begin{minipage}[c]{5.3cm}\raggedright\Large aucune solution\end{minipage}\end{tcolorbox}\hfill\begin{tcolorbox}[enhanced,nobeforeafter,width=7.1cm,height=2.1cm,valign=center,colback=cfon!6,colframe=cfon,arc=2.5mm,boxrule=0.9pt,left=3mm]{\Large\bfseries\color{cfon}b)}\quad\begin{minipage}[c]{5.3cm}\raggedright\Large $\{-25\,;\,25\}$\end{minipage}\end{tcolorbox}\par\vspace{2.5mm}
\begin{tcolorbox}[enhanced,nobeforeafter,width=7.1cm,height=2.1cm,valign=center,colback=cfon!6,colframe=cfon,arc=2.5mm,boxrule=0.9pt,left=3mm]{\Large\bfseries\color{cfon}c)}\quad\begin{minipage}[c]{5.3cm}\raggedright\Large $\{5\}$\end{minipage}\end{tcolorbox}\hfill\begin{tcolorbox}[enhanced,nobeforeafter,width=7.1cm,height=2.1cm,valign=center,colback=cfon!6,colframe=cfon,arc=2.5mm,boxrule=0.9pt,left=3mm]{\Large\bfseries\color{cfon}d)}\quad\begin{minipage}[c]{5.3cm}\raggedright\Large $\{-5\,;\,5\}$\end{minipage}\end{tcolorbox}
\end{tcolorbox}
\newpage
\begin{tcolorbox}[enhanced,width=16cm,height=14cm,colback=white,colframe=cfon,boxrule=2pt,arc=0mm,
left=6mm,right=6mm,top=0mm,bottom=4mm,sharp corners,
overlay={\fill[cfon] (frame.north west) rectangle ([yshift=-1.25cm]frame.north east);
\node[anchor=west,text=white,font=\sffamily\bfseries\LARGE] at ([xshift=5mm,yshift=-0.62cm]frame.north west) {Question 45};
\node[anchor=east,text=white,font=\sffamily\itshape\large] at ([xshift=-5mm,yshift=-0.62cm]frame.north east) {Fonctions};}]
\vspace*{1.75cm}
\begin{minipage}[t][6.9cm][c]{\linewidth}\LARGE\centering
\begin{minipage}{0.95\linewidth}\raggedright\linespread{1.25}\selectfont Quel est le sens de variation de la fonction $f(x)=-3x+5$ sur $\mathbb{R}$~?\end{minipage}
\end{minipage}

\vfill
\begin{tcolorbox}[enhanced,nobeforeafter,width=7.1cm,height=2.1cm,valign=center,colback=cfon!6,colframe=cfon,arc=2.5mm,boxrule=0.9pt,left=3mm]{\Large\bfseries\color{cfon}a)}\quad\begin{minipage}[c]{5.3cm}\raggedright\large croissante sur $\mathbb{R}$\end{minipage}\end{tcolorbox}\hfill\begin{tcolorbox}[enhanced,nobeforeafter,width=7.1cm,height=2.1cm,valign=center,colback=cfon!6,colframe=cfon,arc=2.5mm,boxrule=0.9pt,left=3mm]{\Large\bfseries\color{cfon}b)}\quad\begin{minipage}[c]{5.3cm}\raggedright\large constante sur $\mathbb{R}$\end{minipage}\end{tcolorbox}\par\vspace{2.5mm}
\begin{tcolorbox}[enhanced,nobeforeafter,width=7.1cm,height=2.1cm,valign=center,colback=cfon!6,colframe=cfon,arc=2.5mm,boxrule=0.9pt,left=3mm]{\Large\bfseries\color{cfon}c)}\quad\begin{minipage}[c]{5.3cm}\raggedright\large croissante puis décroissante\end{minipage}\end{tcolorbox}\hfill\begin{tcolorbox}[enhanced,nobeforeafter,width=7.1cm,height=2.1cm,valign=center,colback=cfon!6,colframe=cfon,arc=2.5mm,boxrule=0.9pt,left=3mm]{\Large\bfseries\color{cfon}d)}\quad\begin{minipage}[c]{5.3cm}\raggedright\large décroissante sur $\mathbb{R}$\end{minipage}\end{tcolorbox}
\end{tcolorbox}
\newpage
\begin{tcolorbox}[enhanced,width=16cm,height=14cm,colback=white,colframe=cfon,boxrule=2pt,arc=0mm,
left=6mm,right=6mm,top=0mm,bottom=4mm,sharp corners,
overlay={\fill[cfon] (frame.north west) rectangle ([yshift=-1.25cm]frame.north east);
\node[anchor=west,text=white,font=\sffamily\bfseries\LARGE] at ([xshift=5mm,yshift=-0.62cm]frame.north west) {Question 46};
\node[anchor=east,text=white,font=\sffamily\itshape\large] at ([xshift=-5mm,yshift=-0.62cm]frame.north east) {Fonctions};}]
\vspace*{1.75cm}
\begin{minipage}[t][6.9cm][c]{\linewidth}\LARGE\centering
\begin{minipage}{0.95\linewidth}\raggedright\linespread{1.25}\selectfont Quel est le sens de variation de la fonction $f(x)=2x-1$ sur $\mathbb{R}$~?\end{minipage}
\end{minipage}

\vfill
\begin{tcolorbox}[enhanced,nobeforeafter,width=7.1cm,height=2.1cm,valign=center,colback=cfon!6,colframe=cfon,arc=2.5mm,boxrule=0.9pt,left=3mm]{\Large\bfseries\color{cfon}a)}\quad\begin{minipage}[c]{5.3cm}\raggedright\large croissante puis décroissante\end{minipage}\end{tcolorbox}\hfill\begin{tcolorbox}[enhanced,nobeforeafter,width=7.1cm,height=2.1cm,valign=center,colback=cfon!6,colframe=cfon,arc=2.5mm,boxrule=0.9pt,left=3mm]{\Large\bfseries\color{cfon}b)}\quad\begin{minipage}[c]{5.3cm}\raggedright\large décroissante sur $\mathbb{R}$\end{minipage}\end{tcolorbox}\par\vspace{2.5mm}
\begin{tcolorbox}[enhanced,nobeforeafter,width=7.1cm,height=2.1cm,valign=center,colback=cfon!6,colframe=cfon,arc=2.5mm,boxrule=0.9pt,left=3mm]{\Large\bfseries\color{cfon}c)}\quad\begin{minipage}[c]{5.3cm}\raggedright\large croissante sur $\mathbb{R}$\end{minipage}\end{tcolorbox}\hfill\begin{tcolorbox}[enhanced,nobeforeafter,width=7.1cm,height=2.1cm,valign=center,colback=cfon!6,colframe=cfon,arc=2.5mm,boxrule=0.9pt,left=3mm]{\Large\bfseries\color{cfon}d)}\quad\begin{minipage}[c]{5.3cm}\raggedright\large constante sur $\mathbb{R}$\end{minipage}\end{tcolorbox}
\end{tcolorbox}
\newpage
\begin{tcolorbox}[enhanced,width=16cm,height=14cm,colback=white,colframe=cfon,boxrule=2pt,arc=0mm,
left=6mm,right=6mm,top=0mm,bottom=4mm,sharp corners,
overlay={\fill[cfon] (frame.north west) rectangle ([yshift=-1.25cm]frame.north east);
\node[anchor=west,text=white,font=\sffamily\bfseries\LARGE] at ([xshift=5mm,yshift=-0.62cm]frame.north west) {Question 47};
\node[anchor=east,text=white,font=\sffamily\itshape\large] at ([xshift=-5mm,yshift=-0.62cm]frame.north east) {Fonctions};}]
\vspace*{1.75cm}
\begin{minipage}[t][6.9cm][c]{\linewidth}\LARGE\centering
\begin{minipage}{0.95\linewidth}\raggedright\linespread{1.25}\selectfont Quelles sont les solutions de l'inéquation $3x+2\leqslant x+10$~?\end{minipage}
\end{minipage}

\vfill
\begin{tcolorbox}[enhanced,nobeforeafter,width=7.1cm,height=2.1cm,valign=center,colback=cfon!6,colframe=cfon,arc=2.5mm,boxrule=0.9pt,left=3mm]{\Large\bfseries\color{cfon}a)}\quad\begin{minipage}[c]{5.3cm}\raggedright\Large $x\leqslant 6$\end{minipage}\end{tcolorbox}\hfill\begin{tcolorbox}[enhanced,nobeforeafter,width=7.1cm,height=2.1cm,valign=center,colback=cfon!6,colframe=cfon,arc=2.5mm,boxrule=0.9pt,left=3mm]{\Large\bfseries\color{cfon}b)}\quad\begin{minipage}[c]{5.3cm}\raggedright\Large $x\geqslant 4$\end{minipage}\end{tcolorbox}\par\vspace{2.5mm}
\begin{tcolorbox}[enhanced,nobeforeafter,width=7.1cm,height=2.1cm,valign=center,colback=cfon!6,colframe=cfon,arc=2.5mm,boxrule=0.9pt,left=3mm]{\Large\bfseries\color{cfon}c)}\quad\begin{minipage}[c]{5.3cm}\raggedright\Large $x\geqslant -4$\end{minipage}\end{tcolorbox}\hfill\begin{tcolorbox}[enhanced,nobeforeafter,width=7.1cm,height=2.1cm,valign=center,colback=cfon!6,colframe=cfon,arc=2.5mm,boxrule=0.9pt,left=3mm]{\Large\bfseries\color{cfon}d)}\quad\begin{minipage}[c]{5.3cm}\raggedright\Large $x\leqslant 4$\end{minipage}\end{tcolorbox}
\end{tcolorbox}
\newpage
\begin{tcolorbox}[enhanced,width=16cm,height=14cm,colback=white,colframe=cfon,boxrule=2pt,arc=0mm,
left=6mm,right=6mm,top=0mm,bottom=4mm,sharp corners,
overlay={\fill[cfon] (frame.north west) rectangle ([yshift=-1.25cm]frame.north east);
\node[anchor=west,text=white,font=\sffamily\bfseries\LARGE] at ([xshift=5mm,yshift=-0.62cm]frame.north west) {Question 48};
\node[anchor=east,text=white,font=\sffamily\itshape\large] at ([xshift=-5mm,yshift=-0.62cm]frame.north east) {Fonctions};}]
\vspace*{1.75cm}
\begin{minipage}[t][6.9cm][c]{\linewidth}\LARGE\centering
\begin{minipage}{0.95\linewidth}\raggedright\linespread{1.25}\selectfont Quelles sont les solutions de l'inéquation $2x+5\leqslant 5x-4$~?\end{minipage}
\end{minipage}

\vfill
\begin{tcolorbox}[enhanced,nobeforeafter,width=7.1cm,height=2.1cm,valign=center,colback=cfon!6,colframe=cfon,arc=2.5mm,boxrule=0.9pt,left=3mm]{\Large\bfseries\color{cfon}a)}\quad\begin{minipage}[c]{5.3cm}\raggedright\Large $x\geqslant 3$\end{minipage}\end{tcolorbox}\hfill\begin{tcolorbox}[enhanced,nobeforeafter,width=7.1cm,height=2.1cm,valign=center,colback=cfon!6,colframe=cfon,arc=2.5mm,boxrule=0.9pt,left=3mm]{\Large\bfseries\color{cfon}b)}\quad\begin{minipage}[c]{5.3cm}\raggedright\Large $x\leqslant -\dfrac{1}{3}$\end{minipage}\end{tcolorbox}\par\vspace{2.5mm}
\begin{tcolorbox}[enhanced,nobeforeafter,width=7.1cm,height=2.1cm,valign=center,colback=cfon!6,colframe=cfon,arc=2.5mm,boxrule=0.9pt,left=3mm]{\Large\bfseries\color{cfon}c)}\quad\begin{minipage}[c]{5.3cm}\raggedright\Large $x\geqslant -3$\end{minipage}\end{tcolorbox}\hfill\begin{tcolorbox}[enhanced,nobeforeafter,width=7.1cm,height=2.1cm,valign=center,colback=cfon!6,colframe=cfon,arc=2.5mm,boxrule=0.9pt,left=3mm]{\Large\bfseries\color{cfon}d)}\quad\begin{minipage}[c]{5.3cm}\raggedright\Large $x\leqslant 3$\end{minipage}\end{tcolorbox}
\end{tcolorbox}
\newpage
\begin{tcolorbox}[enhanced,width=16cm,height=14cm,colback=white,colframe=cvec,boxrule=2pt,arc=0mm,
left=6mm,right=6mm,top=0mm,bottom=4mm,sharp corners,
overlay={\fill[cvec] (frame.north west) rectangle ([yshift=-1.25cm]frame.north east);
\node[anchor=west,text=white,font=\sffamily\bfseries\LARGE] at ([xshift=5mm,yshift=-0.62cm]frame.north west) {Question 49};
\node[anchor=east,text=white,font=\sffamily\itshape\large] at ([xshift=-5mm,yshift=-0.62cm]frame.north east) {Vecteurs et repérage};}]
\vspace*{1.75cm}
\begin{minipage}[t][6.9cm][c]{\linewidth}\LARGE\centering
\begin{minipage}{0.95\linewidth}\raggedright\linespread{1.25}\selectfont Dans un repère, $A(1\,;\,2)$ et $B(4\,;\,-3)$. Quelles sont les coordonnées du vecteur $\overrightarrow{AB}$~?\end{minipage}
\end{minipage}

\vfill
\begin{tcolorbox}[enhanced,nobeforeafter,width=7.1cm,height=2.1cm,valign=center,colback=cvec!6,colframe=cvec,arc=2.5mm,boxrule=0.9pt,left=3mm]{\Large\bfseries\color{cvec}a)}\quad\begin{minipage}[c]{5.3cm}\raggedright\Large $(-5\,;\,3)$\end{minipage}\end{tcolorbox}\hfill\begin{tcolorbox}[enhanced,nobeforeafter,width=7.1cm,height=2.1cm,valign=center,colback=cvec!6,colframe=cvec,arc=2.5mm,boxrule=0.9pt,left=3mm]{\Large\bfseries\color{cvec}b)}\quad\begin{minipage}[c]{5.3cm}\raggedright\Large $(3\,;\,-5)$\end{minipage}\end{tcolorbox}\par\vspace{2.5mm}
\begin{tcolorbox}[enhanced,nobeforeafter,width=7.1cm,height=2.1cm,valign=center,colback=cvec!6,colframe=cvec,arc=2.5mm,boxrule=0.9pt,left=3mm]{\Large\bfseries\color{cvec}c)}\quad\begin{minipage}[c]{5.3cm}\raggedright\Large $(-3\,;\,5)$\end{minipage}\end{tcolorbox}\hfill\begin{tcolorbox}[enhanced,nobeforeafter,width=7.1cm,height=2.1cm,valign=center,colback=cvec!6,colframe=cvec,arc=2.5mm,boxrule=0.9pt,left=3mm]{\Large\bfseries\color{cvec}d)}\quad\begin{minipage}[c]{5.3cm}\raggedright\Large $(5\,;\,-1)$\end{minipage}\end{tcolorbox}
\end{tcolorbox}
\newpage
\begin{tcolorbox}[enhanced,width=16cm,height=14cm,colback=white,colframe=cvec,boxrule=2pt,arc=0mm,
left=6mm,right=6mm,top=0mm,bottom=4mm,sharp corners,
overlay={\fill[cvec] (frame.north west) rectangle ([yshift=-1.25cm]frame.north east);
\node[anchor=west,text=white,font=\sffamily\bfseries\LARGE] at ([xshift=5mm,yshift=-0.62cm]frame.north west) {Question 50};
\node[anchor=east,text=white,font=\sffamily\itshape\large] at ([xshift=-5mm,yshift=-0.62cm]frame.north east) {Vecteurs et repérage};}]
\vspace*{1.75cm}
\begin{minipage}[t][6.9cm][c]{\linewidth}\LARGE\centering
\begin{minipage}{0.95\linewidth}\raggedright\linespread{1.25}\selectfont Dans un repère, $A(-2\,;\,5)$ et $B(3\,;\,1)$. Quelles sont les coordonnées du vecteur $\overrightarrow{AB}$~?\end{minipage}
\end{minipage}

\vfill
\begin{tcolorbox}[enhanced,nobeforeafter,width=7.1cm,height=2.1cm,valign=center,colback=cvec!6,colframe=cvec,arc=2.5mm,boxrule=0.9pt,left=3mm]{\Large\bfseries\color{cvec}a)}\quad\begin{minipage}[c]{5.3cm}\raggedright\Large $(5\,;\,-4)$\end{minipage}\end{tcolorbox}\hfill\begin{tcolorbox}[enhanced,nobeforeafter,width=7.1cm,height=2.1cm,valign=center,colback=cvec!6,colframe=cvec,arc=2.5mm,boxrule=0.9pt,left=3mm]{\Large\bfseries\color{cvec}b)}\quad\begin{minipage}[c]{5.3cm}\raggedright\Large $(-4\,;\,5)$\end{minipage}\end{tcolorbox}\par\vspace{2.5mm}
\begin{tcolorbox}[enhanced,nobeforeafter,width=7.1cm,height=2.1cm,valign=center,colback=cvec!6,colframe=cvec,arc=2.5mm,boxrule=0.9pt,left=3mm]{\Large\bfseries\color{cvec}c)}\quad\begin{minipage}[c]{5.3cm}\raggedright\Large $(-5\,;\,4)$\end{minipage}\end{tcolorbox}\hfill\begin{tcolorbox}[enhanced,nobeforeafter,width=7.1cm,height=2.1cm,valign=center,colback=cvec!6,colframe=cvec,arc=2.5mm,boxrule=0.9pt,left=3mm]{\Large\bfseries\color{cvec}d)}\quad\begin{minipage}[c]{5.3cm}\raggedright\Large $(1\,;\,6)$\end{minipage}\end{tcolorbox}
\end{tcolorbox}
\newpage
\begin{tcolorbox}[enhanced,width=16cm,height=14cm,colback=white,colframe=cvec,boxrule=2pt,arc=0mm,
left=6mm,right=6mm,top=0mm,bottom=4mm,sharp corners,
overlay={\fill[cvec] (frame.north west) rectangle ([yshift=-1.25cm]frame.north east);
\node[anchor=west,text=white,font=\sffamily\bfseries\LARGE] at ([xshift=5mm,yshift=-0.62cm]frame.north west) {Question 51};
\node[anchor=east,text=white,font=\sffamily\itshape\large] at ([xshift=-5mm,yshift=-0.62cm]frame.north east) {Vecteurs et repérage};}]
\vspace*{1.75cm}
\begin{minipage}[t][6.9cm][c]{\linewidth}\LARGE\centering
\begin{minipage}{0.95\linewidth}\raggedright\linespread{1.25}\selectfont Dans un repère, $A(1\,;\,2)$ et $B(5\,;\,-4)$. Quelles sont les coordonnées du milieu de $[AB]$~?\end{minipage}
\end{minipage}

\vfill
\begin{tcolorbox}[enhanced,nobeforeafter,width=7.1cm,height=2.1cm,valign=center,colback=cvec!6,colframe=cvec,arc=2.5mm,boxrule=0.9pt,left=3mm]{\Large\bfseries\color{cvec}a)}\quad\begin{minipage}[c]{5.3cm}\raggedright\Large $(2\,;\,-3)$\end{minipage}\end{tcolorbox}\hfill\begin{tcolorbox}[enhanced,nobeforeafter,width=7.1cm,height=2.1cm,valign=center,colback=cvec!6,colframe=cvec,arc=2.5mm,boxrule=0.9pt,left=3mm]{\Large\bfseries\color{cvec}b)}\quad\begin{minipage}[c]{5.3cm}\raggedright\Large $(3\,;\,-1)$\end{minipage}\end{tcolorbox}\par\vspace{2.5mm}
\begin{tcolorbox}[enhanced,nobeforeafter,width=7.1cm,height=2.1cm,valign=center,colback=cvec!6,colframe=cvec,arc=2.5mm,boxrule=0.9pt,left=3mm]{\Large\bfseries\color{cvec}c)}\quad\begin{minipage}[c]{5.3cm}\raggedright\Large $(3\,;\,3)$\end{minipage}\end{tcolorbox}\hfill\begin{tcolorbox}[enhanced,nobeforeafter,width=7.1cm,height=2.1cm,valign=center,colback=cvec!6,colframe=cvec,arc=2.5mm,boxrule=0.9pt,left=3mm]{\Large\bfseries\color{cvec}d)}\quad\begin{minipage}[c]{5.3cm}\raggedright\Large $(6\,;\,-2)$\end{minipage}\end{tcolorbox}
\end{tcolorbox}
\newpage
\begin{tcolorbox}[enhanced,width=16cm,height=14cm,colback=white,colframe=cvec,boxrule=2pt,arc=0mm,
left=6mm,right=6mm,top=0mm,bottom=4mm,sharp corners,
overlay={\fill[cvec] (frame.north west) rectangle ([yshift=-1.25cm]frame.north east);
\node[anchor=west,text=white,font=\sffamily\bfseries\LARGE] at ([xshift=5mm,yshift=-0.62cm]frame.north west) {Question 52};
\node[anchor=east,text=white,font=\sffamily\itshape\large] at ([xshift=-5mm,yshift=-0.62cm]frame.north east) {Vecteurs et repérage};}]
\vspace*{1.75cm}
\begin{minipage}[t][6.9cm][c]{\linewidth}\LARGE\centering
\begin{minipage}{0.95\linewidth}\raggedright\linespread{1.25}\selectfont Dans un repère, $A(-2\,;\,5)$ et $B(4\,;\,1)$. Quelles sont les coordonnées du milieu de $[AB]$~?\end{minipage}
\end{minipage}

\vfill
\begin{tcolorbox}[enhanced,nobeforeafter,width=7.1cm,height=2.1cm,valign=center,colback=cvec!6,colframe=cvec,arc=2.5mm,boxrule=0.9pt,left=3mm]{\Large\bfseries\color{cvec}a)}\quad\begin{minipage}[c]{5.3cm}\raggedright\Large $(1\,;\,2)$\end{minipage}\end{tcolorbox}\hfill\begin{tcolorbox}[enhanced,nobeforeafter,width=7.1cm,height=2.1cm,valign=center,colback=cvec!6,colframe=cvec,arc=2.5mm,boxrule=0.9pt,left=3mm]{\Large\bfseries\color{cvec}b)}\quad\begin{minipage}[c]{5.3cm}\raggedright\Large $(2\,;\,6)$\end{minipage}\end{tcolorbox}\par\vspace{2.5mm}
\begin{tcolorbox}[enhanced,nobeforeafter,width=7.1cm,height=2.1cm,valign=center,colback=cvec!6,colframe=cvec,arc=2.5mm,boxrule=0.9pt,left=3mm]{\Large\bfseries\color{cvec}c)}\quad\begin{minipage}[c]{5.3cm}\raggedright\Large $(3\,;\,-2)$\end{minipage}\end{tcolorbox}\hfill\begin{tcolorbox}[enhanced,nobeforeafter,width=7.1cm,height=2.1cm,valign=center,colback=cvec!6,colframe=cvec,arc=2.5mm,boxrule=0.9pt,left=3mm]{\Large\bfseries\color{cvec}d)}\quad\begin{minipage}[c]{5.3cm}\raggedright\Large $(1\,;\,3)$\end{minipage}\end{tcolorbox}
\end{tcolorbox}
\newpage
\begin{tcolorbox}[enhanced,width=16cm,height=14cm,colback=white,colframe=cvec,boxrule=2pt,arc=0mm,
left=6mm,right=6mm,top=0mm,bottom=4mm,sharp corners,
overlay={\fill[cvec] (frame.north west) rectangle ([yshift=-1.25cm]frame.north east);
\node[anchor=west,text=white,font=\sffamily\bfseries\LARGE] at ([xshift=5mm,yshift=-0.62cm]frame.north west) {Question 53};
\node[anchor=east,text=white,font=\sffamily\itshape\large] at ([xshift=-5mm,yshift=-0.62cm]frame.north east) {Vecteurs et repérage};}]
\vspace*{1.75cm}
\begin{minipage}[t][6.9cm][c]{\linewidth}\LARGE\centering
\begin{minipage}{0.95\linewidth}\raggedright\linespread{1.25}\selectfont Dans un repère orthonormé, $A(1\,;\,2)$ et $B(4\,;\,6)$. Quelle est la distance $AB$~?\end{minipage}
\end{minipage}

\vfill
\begin{tcolorbox}[enhanced,nobeforeafter,width=7.1cm,height=2.1cm,valign=center,colback=cvec!6,colframe=cvec,arc=2.5mm,boxrule=0.9pt,left=3mm]{\Large\bfseries\color{cvec}a)}\quad\begin{minipage}[c]{5.3cm}\raggedright\Large $5$\end{minipage}\end{tcolorbox}\hfill\begin{tcolorbox}[enhanced,nobeforeafter,width=7.1cm,height=2.1cm,valign=center,colback=cvec!6,colframe=cvec,arc=2.5mm,boxrule=0.9pt,left=3mm]{\Large\bfseries\color{cvec}b)}\quad\begin{minipage}[c]{5.3cm}\raggedright\Large $\sqrt{7}$\end{minipage}\end{tcolorbox}\par\vspace{2.5mm}
\begin{tcolorbox}[enhanced,nobeforeafter,width=7.1cm,height=2.1cm,valign=center,colback=cvec!6,colframe=cvec,arc=2.5mm,boxrule=0.9pt,left=3mm]{\Large\bfseries\color{cvec}c)}\quad\begin{minipage}[c]{5.3cm}\raggedright\Large $25$\end{minipage}\end{tcolorbox}\hfill\begin{tcolorbox}[enhanced,nobeforeafter,width=7.1cm,height=2.1cm,valign=center,colback=cvec!6,colframe=cvec,arc=2.5mm,boxrule=0.9pt,left=3mm]{\Large\bfseries\color{cvec}d)}\quad\begin{minipage}[c]{5.3cm}\raggedright\Large $7$\end{minipage}\end{tcolorbox}
\end{tcolorbox}
\newpage
\begin{tcolorbox}[enhanced,width=16cm,height=14cm,colback=white,colframe=cvec,boxrule=2pt,arc=0mm,
left=6mm,right=6mm,top=0mm,bottom=4mm,sharp corners,
overlay={\fill[cvec] (frame.north west) rectangle ([yshift=-1.25cm]frame.north east);
\node[anchor=west,text=white,font=\sffamily\bfseries\LARGE] at ([xshift=5mm,yshift=-0.62cm]frame.north west) {Question 54};
\node[anchor=east,text=white,font=\sffamily\itshape\large] at ([xshift=-5mm,yshift=-0.62cm]frame.north east) {Vecteurs et repérage};}]
\vspace*{1.75cm}
\begin{minipage}[t][6.9cm][c]{\linewidth}\LARGE\centering
\begin{minipage}{0.95\linewidth}\raggedright\linespread{1.25}\selectfont Dans un repère orthonormé, $A(-2\,;\,1)$ et $B(3\,;\,13)$. Quelle est la distance $AB$~?\end{minipage}
\end{minipage}

\vfill
\begin{tcolorbox}[enhanced,nobeforeafter,width=7.1cm,height=2.1cm,valign=center,colback=cvec!6,colframe=cvec,arc=2.5mm,boxrule=0.9pt,left=3mm]{\Large\bfseries\color{cvec}a)}\quad\begin{minipage}[c]{5.3cm}\raggedright\Large $\sqrt{119}$\end{minipage}\end{tcolorbox}\hfill\begin{tcolorbox}[enhanced,nobeforeafter,width=7.1cm,height=2.1cm,valign=center,colback=cvec!6,colframe=cvec,arc=2.5mm,boxrule=0.9pt,left=3mm]{\Large\bfseries\color{cvec}b)}\quad\begin{minipage}[c]{5.3cm}\raggedright\Large $169$\end{minipage}\end{tcolorbox}\par\vspace{2.5mm}
\begin{tcolorbox}[enhanced,nobeforeafter,width=7.1cm,height=2.1cm,valign=center,colback=cvec!6,colframe=cvec,arc=2.5mm,boxrule=0.9pt,left=3mm]{\Large\bfseries\color{cvec}c)}\quad\begin{minipage}[c]{5.3cm}\raggedright\Large $17$\end{minipage}\end{tcolorbox}\hfill\begin{tcolorbox}[enhanced,nobeforeafter,width=7.1cm,height=2.1cm,valign=center,colback=cvec!6,colframe=cvec,arc=2.5mm,boxrule=0.9pt,left=3mm]{\Large\bfseries\color{cvec}d)}\quad\begin{minipage}[c]{5.3cm}\raggedright\Large $13$\end{minipage}\end{tcolorbox}
\end{tcolorbox}
\newpage
\begin{tcolorbox}[enhanced,width=16cm,height=14cm,colback=white,colframe=cvec,boxrule=2pt,arc=0mm,
left=6mm,right=6mm,top=0mm,bottom=4mm,sharp corners,
overlay={\fill[cvec] (frame.north west) rectangle ([yshift=-1.25cm]frame.north east);
\node[anchor=west,text=white,font=\sffamily\bfseries\LARGE] at ([xshift=5mm,yshift=-0.62cm]frame.north west) {Question 55};
\node[anchor=east,text=white,font=\sffamily\itshape\large] at ([xshift=-5mm,yshift=-0.62cm]frame.north east) {Vecteurs et repérage};}]
\vspace*{1.75cm}
\begin{minipage}[t][6.9cm][c]{\linewidth}\LARGE\centering
\begin{minipage}{0.95\linewidth}\raggedright\linespread{1.25}\selectfont Pour quelle valeur de $k$ les vecteurs $\vec u\begin{pmatrix}2\\3\end{pmatrix}$ et $\vec v\begin{pmatrix}k\\9\end{pmatrix}$ sont-ils colinéaires~?\end{minipage}
\end{minipage}

\vfill
\begin{tcolorbox}[enhanced,nobeforeafter,width=7.1cm,height=2.1cm,valign=center,colback=cvec!6,colframe=cvec,arc=2.5mm,boxrule=0.9pt,left=3mm]{\Large\bfseries\color{cvec}a)}\quad\begin{minipage}[c]{5.3cm}\raggedright\Large $k=\dfrac{27}{2}$\end{minipage}\end{tcolorbox}\hfill\begin{tcolorbox}[enhanced,nobeforeafter,width=7.1cm,height=2.1cm,valign=center,colback=cvec!6,colframe=cvec,arc=2.5mm,boxrule=0.9pt,left=3mm]{\Large\bfseries\color{cvec}b)}\quad\begin{minipage}[c]{5.3cm}\raggedright\Large $k=-6$\end{minipage}\end{tcolorbox}\par\vspace{2.5mm}
\begin{tcolorbox}[enhanced,nobeforeafter,width=7.1cm,height=2.1cm,valign=center,colback=cvec!6,colframe=cvec,arc=2.5mm,boxrule=0.9pt,left=3mm]{\Large\bfseries\color{cvec}c)}\quad\begin{minipage}[c]{5.3cm}\raggedright\Large $k=8$\end{minipage}\end{tcolorbox}\hfill\begin{tcolorbox}[enhanced,nobeforeafter,width=7.1cm,height=2.1cm,valign=center,colback=cvec!6,colframe=cvec,arc=2.5mm,boxrule=0.9pt,left=3mm]{\Large\bfseries\color{cvec}d)}\quad\begin{minipage}[c]{5.3cm}\raggedright\Large $k=6$\end{minipage}\end{tcolorbox}
\end{tcolorbox}
\newpage
\begin{tcolorbox}[enhanced,width=16cm,height=14cm,colback=white,colframe=cvec,boxrule=2pt,arc=0mm,
left=6mm,right=6mm,top=0mm,bottom=4mm,sharp corners,
overlay={\fill[cvec] (frame.north west) rectangle ([yshift=-1.25cm]frame.north east);
\node[anchor=west,text=white,font=\sffamily\bfseries\LARGE] at ([xshift=5mm,yshift=-0.62cm]frame.north west) {Question 56};
\node[anchor=east,text=white,font=\sffamily\itshape\large] at ([xshift=-5mm,yshift=-0.62cm]frame.north east) {Vecteurs et repérage};}]
\vspace*{1.75cm}
\begin{minipage}[t][6.9cm][c]{\linewidth}\LARGE\centering
\begin{minipage}{0.95\linewidth}\raggedright\linespread{1.25}\selectfont Pour quelle valeur de $k$ les vecteurs $\vec u\begin{pmatrix}4\\-2\end{pmatrix}$ et $\vec v\begin{pmatrix}k\\5\end{pmatrix}$ sont-ils colinéaires~?\end{minipage}
\end{minipage}

\vfill
\begin{tcolorbox}[enhanced,nobeforeafter,width=7.1cm,height=2.1cm,valign=center,colback=cvec!6,colframe=cvec,arc=2.5mm,boxrule=0.9pt,left=3mm]{\Large\bfseries\color{cvec}a)}\quad\begin{minipage}[c]{5.3cm}\raggedright\Large $k=10$\end{minipage}\end{tcolorbox}\hfill\begin{tcolorbox}[enhanced,nobeforeafter,width=7.1cm,height=2.1cm,valign=center,colback=cvec!6,colframe=cvec,arc=2.5mm,boxrule=0.9pt,left=3mm]{\Large\bfseries\color{cvec}b)}\quad\begin{minipage}[c]{5.3cm}\raggedright\Large $k=-\dfrac{5}{2}$\end{minipage}\end{tcolorbox}\par\vspace{2.5mm}
\begin{tcolorbox}[enhanced,nobeforeafter,width=7.1cm,height=2.1cm,valign=center,colback=cvec!6,colframe=cvec,arc=2.5mm,boxrule=0.9pt,left=3mm]{\Large\bfseries\color{cvec}c)}\quad\begin{minipage}[c]{5.3cm}\raggedright\Large $k=-10$\end{minipage}\end{tcolorbox}\hfill\begin{tcolorbox}[enhanced,nobeforeafter,width=7.1cm,height=2.1cm,valign=center,colback=cvec!6,colframe=cvec,arc=2.5mm,boxrule=0.9pt,left=3mm]{\Large\bfseries\color{cvec}d)}\quad\begin{minipage}[c]{5.3cm}\raggedright\Large $k=11$\end{minipage}\end{tcolorbox}
\end{tcolorbox}
\newpage
\begin{tcolorbox}[enhanced,width=16cm,height=14cm,colback=white,colframe=cvec,boxrule=2pt,arc=0mm,
left=6mm,right=6mm,top=0mm,bottom=4mm,sharp corners,
overlay={\fill[cvec] (frame.north west) rectangle ([yshift=-1.25cm]frame.north east);
\node[anchor=west,text=white,font=\sffamily\bfseries\LARGE] at ([xshift=5mm,yshift=-0.62cm]frame.north west) {Question 57};
\node[anchor=east,text=white,font=\sffamily\itshape\large] at ([xshift=-5mm,yshift=-0.62cm]frame.north east) {Vecteurs et repérage};}]
\vspace*{1.75cm}
\begin{minipage}[t][6.9cm][c]{\linewidth}\Large\centering
\begin{minipage}{0.95\linewidth}\raggedright\linespread{1.25}\selectfont On donne $\vec u\begin{pmatrix}1\\2\end{pmatrix}$ et $\vec v\begin{pmatrix}3\\-1\end{pmatrix}$. Quelles sont les coordonnées du vecteur $2\vec u+\vec v$~?\end{minipage}
\end{minipage}

\vfill
\begin{tcolorbox}[enhanced,nobeforeafter,width=7.1cm,height=2.1cm,valign=center,colback=cvec!6,colframe=cvec,arc=2.5mm,boxrule=0.9pt,left=3mm]{\Large\bfseries\color{cvec}a)}\quad\begin{minipage}[c]{5.3cm}\raggedright\Large $(3\,;\,5)$\end{minipage}\end{tcolorbox}\hfill\begin{tcolorbox}[enhanced,nobeforeafter,width=7.1cm,height=2.1cm,valign=center,colback=cvec!6,colframe=cvec,arc=2.5mm,boxrule=0.9pt,left=3mm]{\Large\bfseries\color{cvec}b)}\quad\begin{minipage}[c]{5.3cm}\raggedright\Large $(5\,;\,3)$\end{minipage}\end{tcolorbox}\par\vspace{2.5mm}
\begin{tcolorbox}[enhanced,nobeforeafter,width=7.1cm,height=2.1cm,valign=center,colback=cvec!6,colframe=cvec,arc=2.5mm,boxrule=0.9pt,left=3mm]{\Large\bfseries\color{cvec}c)}\quad\begin{minipage}[c]{5.3cm}\raggedright\Large $(-1\,;\,5)$\end{minipage}\end{tcolorbox}\hfill\begin{tcolorbox}[enhanced,nobeforeafter,width=7.1cm,height=2.1cm,valign=center,colback=cvec!6,colframe=cvec,arc=2.5mm,boxrule=0.9pt,left=3mm]{\Large\bfseries\color{cvec}d)}\quad\begin{minipage}[c]{5.3cm}\raggedright\Large $(4\,;\,1)$\end{minipage}\end{tcolorbox}
\end{tcolorbox}
\newpage
\begin{tcolorbox}[enhanced,width=16cm,height=14cm,colback=white,colframe=cvec,boxrule=2pt,arc=0mm,
left=6mm,right=6mm,top=0mm,bottom=4mm,sharp corners,
overlay={\fill[cvec] (frame.north west) rectangle ([yshift=-1.25cm]frame.north east);
\node[anchor=west,text=white,font=\sffamily\bfseries\LARGE] at ([xshift=5mm,yshift=-0.62cm]frame.north west) {Question 58};
\node[anchor=east,text=white,font=\sffamily\itshape\large] at ([xshift=-5mm,yshift=-0.62cm]frame.north east) {Vecteurs et repérage};}]
\vspace*{1.75cm}
\begin{minipage}[t][6.9cm][c]{\linewidth}\Large\centering
\begin{minipage}{0.95\linewidth}\raggedright\linespread{1.25}\selectfont On donne $\vec u\begin{pmatrix}4\\-3\end{pmatrix}$ et $\vec v\begin{pmatrix}1\\2\end{pmatrix}$. Quelles sont les coordonnées du vecteur $\vec u-2\vec v$~?\end{minipage}
\end{minipage}

\vfill
\begin{tcolorbox}[enhanced,nobeforeafter,width=7.1cm,height=2.1cm,valign=center,colback=cvec!6,colframe=cvec,arc=2.5mm,boxrule=0.9pt,left=3mm]{\Large\bfseries\color{cvec}a)}\quad\begin{minipage}[c]{5.3cm}\raggedright\Large $(-7\,;\,2)$\end{minipage}\end{tcolorbox}\hfill\begin{tcolorbox}[enhanced,nobeforeafter,width=7.1cm,height=2.1cm,valign=center,colback=cvec!6,colframe=cvec,arc=2.5mm,boxrule=0.9pt,left=3mm]{\Large\bfseries\color{cvec}b)}\quad\begin{minipage}[c]{5.3cm}\raggedright\Large $(5\,;\,-1)$\end{minipage}\end{tcolorbox}\par\vspace{2.5mm}
\begin{tcolorbox}[enhanced,nobeforeafter,width=7.1cm,height=2.1cm,valign=center,colback=cvec!6,colframe=cvec,arc=2.5mm,boxrule=0.9pt,left=3mm]{\Large\bfseries\color{cvec}c)}\quad\begin{minipage}[c]{5.3cm}\raggedright\Large $(6\,;\,1)$\end{minipage}\end{tcolorbox}\hfill\begin{tcolorbox}[enhanced,nobeforeafter,width=7.1cm,height=2.1cm,valign=center,colback=cvec!6,colframe=cvec,arc=2.5mm,boxrule=0.9pt,left=3mm]{\Large\bfseries\color{cvec}d)}\quad\begin{minipage}[c]{5.3cm}\raggedright\Large $(2\,;\,-7)$\end{minipage}\end{tcolorbox}
\end{tcolorbox}
\newpage
\begin{tcolorbox}[enhanced,width=16cm,height=14cm,colback=white,colframe=cvec,boxrule=2pt,arc=0mm,
left=6mm,right=6mm,top=0mm,bottom=4mm,sharp corners,
overlay={\fill[cvec] (frame.north west) rectangle ([yshift=-1.25cm]frame.north east);
\node[anchor=west,text=white,font=\sffamily\bfseries\LARGE] at ([xshift=5mm,yshift=-0.62cm]frame.north west) {Question 59};
\node[anchor=east,text=white,font=\sffamily\itshape\large] at ([xshift=-5mm,yshift=-0.62cm]frame.north east) {Vecteurs et repérage};}]
\vspace*{1.75cm}
\begin{minipage}[t][6.9cm][c]{\linewidth}\LARGE\centering
\begin{minipage}{0.95\linewidth}\raggedright\linespread{1.25}\selectfont Que vaut le déterminant $\det(\vec u,\vec v)$ pour $\vec u\begin{pmatrix}2\\3\end{pmatrix}$ et $\vec v\begin{pmatrix}4\\1\end{pmatrix}$~?\end{minipage}
\end{minipage}

\vfill
\begin{tcolorbox}[enhanced,nobeforeafter,width=7.1cm,height=2.1cm,valign=center,colback=cvec!6,colframe=cvec,arc=2.5mm,boxrule=0.9pt,left=3mm]{\Large\bfseries\color{cvec}a)}\quad\begin{minipage}[c]{5.3cm}\raggedright\Large $14$\end{minipage}\end{tcolorbox}\hfill\begin{tcolorbox}[enhanced,nobeforeafter,width=7.1cm,height=2.1cm,valign=center,colback=cvec!6,colframe=cvec,arc=2.5mm,boxrule=0.9pt,left=3mm]{\Large\bfseries\color{cvec}b)}\quad\begin{minipage}[c]{5.3cm}\raggedright\Large $10$\end{minipage}\end{tcolorbox}\par\vspace{2.5mm}
\begin{tcolorbox}[enhanced,nobeforeafter,width=7.1cm,height=2.1cm,valign=center,colback=cvec!6,colframe=cvec,arc=2.5mm,boxrule=0.9pt,left=3mm]{\Large\bfseries\color{cvec}c)}\quad\begin{minipage}[c]{5.3cm}\raggedright\Large $11$\end{minipage}\end{tcolorbox}\hfill\begin{tcolorbox}[enhanced,nobeforeafter,width=7.1cm,height=2.1cm,valign=center,colback=cvec!6,colframe=cvec,arc=2.5mm,boxrule=0.9pt,left=3mm]{\Large\bfseries\color{cvec}d)}\quad\begin{minipage}[c]{5.3cm}\raggedright\Large $-10$\end{minipage}\end{tcolorbox}
\end{tcolorbox}
\newpage
\begin{tcolorbox}[enhanced,width=16cm,height=14cm,colback=white,colframe=cvec,boxrule=2pt,arc=0mm,
left=6mm,right=6mm,top=0mm,bottom=4mm,sharp corners,
overlay={\fill[cvec] (frame.north west) rectangle ([yshift=-1.25cm]frame.north east);
\node[anchor=west,text=white,font=\sffamily\bfseries\LARGE] at ([xshift=5mm,yshift=-0.62cm]frame.north west) {Question 60};
\node[anchor=east,text=white,font=\sffamily\itshape\large] at ([xshift=-5mm,yshift=-0.62cm]frame.north east) {Vecteurs et repérage};}]
\vspace*{1.75cm}
\begin{minipage}[t][6.9cm][c]{\linewidth}\LARGE\centering
\begin{minipage}{0.95\linewidth}\raggedright\linespread{1.25}\selectfont Que vaut le déterminant $\det(\vec u,\vec v)$ pour $\vec u\begin{pmatrix}-1\\5\end{pmatrix}$ et $\vec v\begin{pmatrix}2\\-3\end{pmatrix}$~?\end{minipage}
\end{minipage}

\vfill
\begin{tcolorbox}[enhanced,nobeforeafter,width=7.1cm,height=2.1cm,valign=center,colback=cvec!6,colframe=cvec,arc=2.5mm,boxrule=0.9pt,left=3mm]{\Large\bfseries\color{cvec}a)}\quad\begin{minipage}[c]{5.3cm}\raggedright\Large $7$\end{minipage}\end{tcolorbox}\hfill\begin{tcolorbox}[enhanced,nobeforeafter,width=7.1cm,height=2.1cm,valign=center,colback=cvec!6,colframe=cvec,arc=2.5mm,boxrule=0.9pt,left=3mm]{\Large\bfseries\color{cvec}b)}\quad\begin{minipage}[c]{5.3cm}\raggedright\Large $-17$\end{minipage}\end{tcolorbox}\par\vspace{2.5mm}
\begin{tcolorbox}[enhanced,nobeforeafter,width=7.1cm,height=2.1cm,valign=center,colback=cvec!6,colframe=cvec,arc=2.5mm,boxrule=0.9pt,left=3mm]{\Large\bfseries\color{cvec}c)}\quad\begin{minipage}[c]{5.3cm}\raggedright\Large $13$\end{minipage}\end{tcolorbox}\hfill\begin{tcolorbox}[enhanced,nobeforeafter,width=7.1cm,height=2.1cm,valign=center,colback=cvec!6,colframe=cvec,arc=2.5mm,boxrule=0.9pt,left=3mm]{\Large\bfseries\color{cvec}d)}\quad\begin{minipage}[c]{5.3cm}\raggedright\Large $-7$\end{minipage}\end{tcolorbox}
\end{tcolorbox}
\newpage
\begin{tcolorbox}[enhanced,width=16cm,height=14cm,colback=white,colframe=cvec,boxrule=2pt,arc=0mm,
left=6mm,right=6mm,top=0mm,bottom=4mm,sharp corners,
overlay={\fill[cvec] (frame.north west) rectangle ([yshift=-1.25cm]frame.north east);
\node[anchor=west,text=white,font=\sffamily\bfseries\LARGE] at ([xshift=5mm,yshift=-0.62cm]frame.north west) {Question 61};
\node[anchor=east,text=white,font=\sffamily\itshape\large] at ([xshift=-5mm,yshift=-0.62cm]frame.north east) {Vecteurs et repérage};}]
\vspace*{1.75cm}
\begin{minipage}[t][6.9cm][c]{\linewidth}\LARGE\centering
\begin{minipage}{0.95\linewidth}\raggedright\linespread{1.25}\selectfont Dans un repère, $A(1\,;\,1)$, $B(4\,;\,2)$ et $C(5\,;\,5)$. Quelles sont les coordonnées du point $D$ tel que $ABCD$ soit un parallélogramme~?\end{minipage}
\end{minipage}

\vfill
\begin{tcolorbox}[enhanced,nobeforeafter,width=7.1cm,height=2.1cm,valign=center,colback=cvec!6,colframe=cvec,arc=2.5mm,boxrule=0.9pt,left=3mm]{\Large\bfseries\color{cvec}a)}\quad\begin{minipage}[c]{5.3cm}\raggedright\Large $(0\,;\,-2)$\end{minipage}\end{tcolorbox}\hfill\begin{tcolorbox}[enhanced,nobeforeafter,width=7.1cm,height=2.1cm,valign=center,colback=cvec!6,colframe=cvec,arc=2.5mm,boxrule=0.9pt,left=3mm]{\Large\bfseries\color{cvec}b)}\quad\begin{minipage}[c]{5.3cm}\raggedright\Large $(8\,;\,6)$\end{minipage}\end{tcolorbox}\par\vspace{2.5mm}
\begin{tcolorbox}[enhanced,nobeforeafter,width=7.1cm,height=2.1cm,valign=center,colback=cvec!6,colframe=cvec,arc=2.5mm,boxrule=0.9pt,left=3mm]{\Large\bfseries\color{cvec}c)}\quad\begin{minipage}[c]{5.3cm}\raggedright\Large $(2\,;\,4)$\end{minipage}\end{tcolorbox}\hfill\begin{tcolorbox}[enhanced,nobeforeafter,width=7.1cm,height=2.1cm,valign=center,colback=cvec!6,colframe=cvec,arc=2.5mm,boxrule=0.9pt,left=3mm]{\Large\bfseries\color{cvec}d)}\quad\begin{minipage}[c]{5.3cm}\raggedright\Large $(3\,;\,3)$\end{minipage}\end{tcolorbox}
\end{tcolorbox}
\newpage
\begin{tcolorbox}[enhanced,width=16cm,height=14cm,colback=white,colframe=cvec,boxrule=2pt,arc=0mm,
left=6mm,right=6mm,top=0mm,bottom=4mm,sharp corners,
overlay={\fill[cvec] (frame.north west) rectangle ([yshift=-1.25cm]frame.north east);
\node[anchor=west,text=white,font=\sffamily\bfseries\LARGE] at ([xshift=5mm,yshift=-0.62cm]frame.north west) {Question 62};
\node[anchor=east,text=white,font=\sffamily\itshape\large] at ([xshift=-5mm,yshift=-0.62cm]frame.north east) {Vecteurs et repérage};}]
\vspace*{1.75cm}
\begin{minipage}[t][6.9cm][c]{\linewidth}\LARGE\centering
\begin{minipage}{0.95\linewidth}\raggedright\linespread{1.25}\selectfont Dans un repère, $A(-2\,;\,0)$, $B(1\,;\,3)$ et $C(0\,;\,6)$. Quelles sont les coordonnées du point $D$ tel que $ABCD$ soit un parallélogramme~?\end{minipage}
\end{minipage}

\vfill
\begin{tcolorbox}[enhanced,nobeforeafter,width=7.1cm,height=2.1cm,valign=center,colback=cvec!6,colframe=cvec,arc=2.5mm,boxrule=0.9pt,left=3mm]{\Large\bfseries\color{cvec}a)}\quad\begin{minipage}[c]{5.3cm}\raggedright\Large $(3\,;\,9)$\end{minipage}\end{tcolorbox}\hfill\begin{tcolorbox}[enhanced,nobeforeafter,width=7.1cm,height=2.1cm,valign=center,colback=cvec!6,colframe=cvec,arc=2.5mm,boxrule=0.9pt,left=3mm]{\Large\bfseries\color{cvec}b)}\quad\begin{minipage}[c]{5.3cm}\raggedright\Large $(-3\,;\,3)$\end{minipage}\end{tcolorbox}\par\vspace{2.5mm}
\begin{tcolorbox}[enhanced,nobeforeafter,width=7.1cm,height=2.1cm,valign=center,colback=cvec!6,colframe=cvec,arc=2.5mm,boxrule=0.9pt,left=3mm]{\Large\bfseries\color{cvec}c)}\quad\begin{minipage}[c]{5.3cm}\raggedright\Large $(-1\,;\,3)$\end{minipage}\end{tcolorbox}\hfill\begin{tcolorbox}[enhanced,nobeforeafter,width=7.1cm,height=2.1cm,valign=center,colback=cvec!6,colframe=cvec,arc=2.5mm,boxrule=0.9pt,left=3mm]{\Large\bfseries\color{cvec}d)}\quad\begin{minipage}[c]{5.3cm}\raggedright\Large $(-1\,;\,-3)$\end{minipage}\end{tcolorbox}
\end{tcolorbox}
\newpage
\begin{tcolorbox}[enhanced,width=16cm,height=14cm,colback=white,colframe=cvec,boxrule=2pt,arc=0mm,
left=6mm,right=6mm,top=0mm,bottom=4mm,sharp corners,
overlay={\fill[cvec] (frame.north west) rectangle ([yshift=-1.25cm]frame.north east);
\node[anchor=west,text=white,font=\sffamily\bfseries\LARGE] at ([xshift=5mm,yshift=-0.62cm]frame.north west) {Question 63};
\node[anchor=east,text=white,font=\sffamily\itshape\large] at ([xshift=-5mm,yshift=-0.62cm]frame.north east) {Vecteurs et repérage};}]
\vspace*{1.75cm}
\begin{minipage}[t][6.9cm][c]{\linewidth}\LARGE\centering
\begin{minipage}{0.95\linewidth}\raggedright\linespread{1.25}\selectfont Quel est le couple solution du système $\left\{\begin{array}{l}x+y=4\\x-y=2\end{array}\right.$~?\end{minipage}
\end{minipage}

\vfill
\begin{tcolorbox}[enhanced,nobeforeafter,width=7.1cm,height=2.1cm,valign=center,colback=cvec!6,colframe=cvec,arc=2.5mm,boxrule=0.9pt,left=3mm]{\Large\bfseries\color{cvec}a)}\quad\begin{minipage}[c]{5.3cm}\raggedright\Large $(1\,;\,3)$\end{minipage}\end{tcolorbox}\hfill\begin{tcolorbox}[enhanced,nobeforeafter,width=7.1cm,height=2.1cm,valign=center,colback=cvec!6,colframe=cvec,arc=2.5mm,boxrule=0.9pt,left=3mm]{\Large\bfseries\color{cvec}b)}\quad\begin{minipage}[c]{5.3cm}\raggedright\Large $(3\,;\,-1)$\end{minipage}\end{tcolorbox}\par\vspace{2.5mm}
\begin{tcolorbox}[enhanced,nobeforeafter,width=7.1cm,height=2.1cm,valign=center,colback=cvec!6,colframe=cvec,arc=2.5mm,boxrule=0.9pt,left=3mm]{\Large\bfseries\color{cvec}c)}\quad\begin{minipage}[c]{5.3cm}\raggedright\Large $(3\,;\,1)$\end{minipage}\end{tcolorbox}\hfill\begin{tcolorbox}[enhanced,nobeforeafter,width=7.1cm,height=2.1cm,valign=center,colback=cvec!6,colframe=cvec,arc=2.5mm,boxrule=0.9pt,left=3mm]{\Large\bfseries\color{cvec}d)}\quad\begin{minipage}[c]{5.3cm}\raggedright\Large $(-3\,;\,1)$\end{minipage}\end{tcolorbox}
\end{tcolorbox}
\newpage
\begin{tcolorbox}[enhanced,width=16cm,height=14cm,colback=white,colframe=cvec,boxrule=2pt,arc=0mm,
left=6mm,right=6mm,top=0mm,bottom=4mm,sharp corners,
overlay={\fill[cvec] (frame.north west) rectangle ([yshift=-1.25cm]frame.north east);
\node[anchor=west,text=white,font=\sffamily\bfseries\LARGE] at ([xshift=5mm,yshift=-0.62cm]frame.north west) {Question 64};
\node[anchor=east,text=white,font=\sffamily\itshape\large] at ([xshift=-5mm,yshift=-0.62cm]frame.north east) {Vecteurs et repérage};}]
\vspace*{1.75cm}
\begin{minipage}[t][6.9cm][c]{\linewidth}\LARGE\centering
\begin{minipage}{0.95\linewidth}\raggedright\linespread{1.25}\selectfont Quel est le couple solution du système $\left\{\begin{array}{l}2x+y=3\\x-y=3\end{array}\right.$~?\end{minipage}
\end{minipage}

\vfill
\begin{tcolorbox}[enhanced,nobeforeafter,width=7.1cm,height=2.1cm,valign=center,colback=cvec!6,colframe=cvec,arc=2.5mm,boxrule=0.9pt,left=3mm]{\Large\bfseries\color{cvec}a)}\quad\begin{minipage}[c]{5.3cm}\raggedright\Large $(2\,;\,1)$\end{minipage}\end{tcolorbox}\hfill\begin{tcolorbox}[enhanced,nobeforeafter,width=7.1cm,height=2.1cm,valign=center,colback=cvec!6,colframe=cvec,arc=2.5mm,boxrule=0.9pt,left=3mm]{\Large\bfseries\color{cvec}b)}\quad\begin{minipage}[c]{5.3cm}\raggedright\Large $(2\,;\,-1)$\end{minipage}\end{tcolorbox}\par\vspace{2.5mm}
\begin{tcolorbox}[enhanced,nobeforeafter,width=7.1cm,height=2.1cm,valign=center,colback=cvec!6,colframe=cvec,arc=2.5mm,boxrule=0.9pt,left=3mm]{\Large\bfseries\color{cvec}c)}\quad\begin{minipage}[c]{5.3cm}\raggedright\Large $(-1\,;\,2)$\end{minipage}\end{tcolorbox}\hfill\begin{tcolorbox}[enhanced,nobeforeafter,width=7.1cm,height=2.1cm,valign=center,colback=cvec!6,colframe=cvec,arc=2.5mm,boxrule=0.9pt,left=3mm]{\Large\bfseries\color{cvec}d)}\quad\begin{minipage}[c]{5.3cm}\raggedright\Large $(-2\,;\,-1)$\end{minipage}\end{tcolorbox}
\end{tcolorbox}
\newpage
\begin{tcolorbox}[enhanced,width=16cm,height=14cm,colback=white,colframe=cari,boxrule=2pt,arc=0mm,
left=6mm,right=6mm,top=0mm,bottom=4mm,sharp corners,
overlay={\fill[cari] (frame.north west) rectangle ([yshift=-1.25cm]frame.north east);
\node[anchor=west,text=white,font=\sffamily\bfseries\LARGE] at ([xshift=5mm,yshift=-0.62cm]frame.north west) {Question 65};
\node[anchor=east,text=white,font=\sffamily\itshape\large] at ([xshift=-5mm,yshift=-0.62cm]frame.north east) {Arithmétique et ensembles de nombres};}]
\vspace*{1.75cm}
\begin{minipage}[t][6.9cm][c]{\linewidth}\LARGE\centering
\begin{minipage}{0.95\linewidth}\raggedright\linespread{1.25}\selectfont Lequel de ces nombres est premier~?\end{minipage}
\end{minipage}

\vfill
\begin{tcolorbox}[enhanced,nobeforeafter,width=7.1cm,height=2.1cm,valign=center,colback=cari!6,colframe=cari,arc=2.5mm,boxrule=0.9pt,left=3mm]{\Large\bfseries\color{cari}a)}\quad\begin{minipage}[c]{5.3cm}\raggedright\Large $87$\end{minipage}\end{tcolorbox}\hfill\begin{tcolorbox}[enhanced,nobeforeafter,width=7.1cm,height=2.1cm,valign=center,colback=cari!6,colframe=cari,arc=2.5mm,boxrule=0.9pt,left=3mm]{\Large\bfseries\color{cari}b)}\quad\begin{minipage}[c]{5.3cm}\raggedright\Large $91$\end{minipage}\end{tcolorbox}\par\vspace{2.5mm}
\begin{tcolorbox}[enhanced,nobeforeafter,width=7.1cm,height=2.1cm,valign=center,colback=cari!6,colframe=cari,arc=2.5mm,boxrule=0.9pt,left=3mm]{\Large\bfseries\color{cari}c)}\quad\begin{minipage}[c]{5.3cm}\raggedright\Large $97$\end{minipage}\end{tcolorbox}\hfill\begin{tcolorbox}[enhanced,nobeforeafter,width=7.1cm,height=2.1cm,valign=center,colback=cari!6,colframe=cari,arc=2.5mm,boxrule=0.9pt,left=3mm]{\Large\bfseries\color{cari}d)}\quad\begin{minipage}[c]{5.3cm}\raggedright\Large $93$\end{minipage}\end{tcolorbox}
\end{tcolorbox}
\newpage
\begin{tcolorbox}[enhanced,width=16cm,height=14cm,colback=white,colframe=cari,boxrule=2pt,arc=0mm,
left=6mm,right=6mm,top=0mm,bottom=4mm,sharp corners,
overlay={\fill[cari] (frame.north west) rectangle ([yshift=-1.25cm]frame.north east);
\node[anchor=west,text=white,font=\sffamily\bfseries\LARGE] at ([xshift=5mm,yshift=-0.62cm]frame.north west) {Question 66};
\node[anchor=east,text=white,font=\sffamily\itshape\large] at ([xshift=-5mm,yshift=-0.62cm]frame.north east) {Arithmétique et ensembles de nombres};}]
\vspace*{1.75cm}
\begin{minipage}[t][6.9cm][c]{\linewidth}\LARGE\centering
\begin{minipage}{0.95\linewidth}\raggedright\linespread{1.25}\selectfont Lequel de ces nombres est premier~?\end{minipage}
\end{minipage}

\vfill
\begin{tcolorbox}[enhanced,nobeforeafter,width=7.1cm,height=2.1cm,valign=center,colback=cari!6,colframe=cari,arc=2.5mm,boxrule=0.9pt,left=3mm]{\Large\bfseries\color{cari}a)}\quad\begin{minipage}[c]{5.3cm}\raggedright\Large $51$\end{minipage}\end{tcolorbox}\hfill\begin{tcolorbox}[enhanced,nobeforeafter,width=7.1cm,height=2.1cm,valign=center,colback=cari!6,colframe=cari,arc=2.5mm,boxrule=0.9pt,left=3mm]{\Large\bfseries\color{cari}b)}\quad\begin{minipage}[c]{5.3cm}\raggedright\Large $49$\end{minipage}\end{tcolorbox}\par\vspace{2.5mm}
\begin{tcolorbox}[enhanced,nobeforeafter,width=7.1cm,height=2.1cm,valign=center,colback=cari!6,colframe=cari,arc=2.5mm,boxrule=0.9pt,left=3mm]{\Large\bfseries\color{cari}c)}\quad\begin{minipage}[c]{5.3cm}\raggedright\Large $53$\end{minipage}\end{tcolorbox}\hfill\begin{tcolorbox}[enhanced,nobeforeafter,width=7.1cm,height=2.1cm,valign=center,colback=cari!6,colframe=cari,arc=2.5mm,boxrule=0.9pt,left=3mm]{\Large\bfseries\color{cari}d)}\quad\begin{minipage}[c]{5.3cm}\raggedright\Large $57$\end{minipage}\end{tcolorbox}
\end{tcolorbox}
\newpage
\begin{tcolorbox}[enhanced,width=16cm,height=14cm,colback=white,colframe=cari,boxrule=2pt,arc=0mm,
left=6mm,right=6mm,top=0mm,bottom=4mm,sharp corners,
overlay={\fill[cari] (frame.north west) rectangle ([yshift=-1.25cm]frame.north east);
\node[anchor=west,text=white,font=\sffamily\bfseries\LARGE] at ([xshift=5mm,yshift=-0.62cm]frame.north west) {Question 67};
\node[anchor=east,text=white,font=\sffamily\itshape\large] at ([xshift=-5mm,yshift=-0.62cm]frame.north east) {Arithmétique et ensembles de nombres};}]
\vspace*{1.75cm}
\begin{minipage}[t][6.9cm][c]{\linewidth}\LARGE\centering
\begin{minipage}{0.95\linewidth}\raggedright\linespread{1.25}\selectfont Quelle est la décomposition de $360$ en produit de facteurs premiers~?\end{minipage}
\end{minipage}

\vfill
\begin{tcolorbox}[enhanced,nobeforeafter,width=7.1cm,height=2.1cm,valign=center,colback=cari!6,colframe=cari,arc=2.5mm,boxrule=0.9pt,left=3mm]{\Large\bfseries\color{cari}a)}\quad\begin{minipage}[c]{5.3cm}\raggedright\Large $2^{3}\times3^{2}\times5$\end{minipage}\end{tcolorbox}\hfill\begin{tcolorbox}[enhanced,nobeforeafter,width=7.1cm,height=2.1cm,valign=center,colback=cari!6,colframe=cari,arc=2.5mm,boxrule=0.9pt,left=3mm]{\Large\bfseries\color{cari}b)}\quad\begin{minipage}[c]{5.3cm}\raggedright\large $2^{3}\times3^{2}\times5^{2}$\end{minipage}\end{tcolorbox}\par\vspace{2.5mm}
\begin{tcolorbox}[enhanced,nobeforeafter,width=7.1cm,height=2.1cm,valign=center,colback=cari!6,colframe=cari,arc=2.5mm,boxrule=0.9pt,left=3mm]{\Large\bfseries\color{cari}c)}\quad\begin{minipage}[c]{5.3cm}\raggedright\Large $24\times15$\end{minipage}\end{tcolorbox}\hfill\begin{tcolorbox}[enhanced,nobeforeafter,width=7.1cm,height=2.1cm,valign=center,colback=cari!6,colframe=cari,arc=2.5mm,boxrule=0.9pt,left=3mm]{\Large\bfseries\color{cari}d)}\quad\begin{minipage}[c]{5.3cm}\raggedright\Large $2^{4}\times3^{2}\times5$\end{minipage}\end{tcolorbox}
\end{tcolorbox}
\newpage
\begin{tcolorbox}[enhanced,width=16cm,height=14cm,colback=white,colframe=cari,boxrule=2pt,arc=0mm,
left=6mm,right=6mm,top=0mm,bottom=4mm,sharp corners,
overlay={\fill[cari] (frame.north west) rectangle ([yshift=-1.25cm]frame.north east);
\node[anchor=west,text=white,font=\sffamily\bfseries\LARGE] at ([xshift=5mm,yshift=-0.62cm]frame.north west) {Question 68};
\node[anchor=east,text=white,font=\sffamily\itshape\large] at ([xshift=-5mm,yshift=-0.62cm]frame.north east) {Arithmétique et ensembles de nombres};}]
\vspace*{1.75cm}
\begin{minipage}[t][6.9cm][c]{\linewidth}\LARGE\centering
\begin{minipage}{0.95\linewidth}\raggedright\linespread{1.25}\selectfont Quelle est la décomposition de $84$ en produit de facteurs premiers~?\end{minipage}
\end{minipage}

\vfill
\begin{tcolorbox}[enhanced,nobeforeafter,width=7.1cm,height=2.1cm,valign=center,colback=cari!6,colframe=cari,arc=2.5mm,boxrule=0.9pt,left=3mm]{\Large\bfseries\color{cari}a)}\quad\begin{minipage}[c]{5.3cm}\raggedright\Large $2^{2}\times3\times7^{2}$\end{minipage}\end{tcolorbox}\hfill\begin{tcolorbox}[enhanced,nobeforeafter,width=7.1cm,height=2.1cm,valign=center,colback=cari!6,colframe=cari,arc=2.5mm,boxrule=0.9pt,left=3mm]{\Large\bfseries\color{cari}b)}\quad\begin{minipage}[c]{5.3cm}\raggedright\Large $2^{3}\times3\times7$\end{minipage}\end{tcolorbox}\par\vspace{2.5mm}
\begin{tcolorbox}[enhanced,nobeforeafter,width=7.1cm,height=2.1cm,valign=center,colback=cari!6,colframe=cari,arc=2.5mm,boxrule=0.9pt,left=3mm]{\Large\bfseries\color{cari}c)}\quad\begin{minipage}[c]{5.3cm}\raggedright\Large $2^{2}\times3\times7$\end{minipage}\end{tcolorbox}\hfill\begin{tcolorbox}[enhanced,nobeforeafter,width=7.1cm,height=2.1cm,valign=center,colback=cari!6,colframe=cari,arc=2.5mm,boxrule=0.9pt,left=3mm]{\Large\bfseries\color{cari}d)}\quad\begin{minipage}[c]{5.3cm}\raggedright\Large $12\times7$\end{minipage}\end{tcolorbox}
\end{tcolorbox}
\newpage
\begin{tcolorbox}[enhanced,width=16cm,height=14cm,colback=white,colframe=cari,boxrule=2pt,arc=0mm,
left=6mm,right=6mm,top=0mm,bottom=4mm,sharp corners,
overlay={\fill[cari] (frame.north west) rectangle ([yshift=-1.25cm]frame.north east);
\node[anchor=west,text=white,font=\sffamily\bfseries\LARGE] at ([xshift=5mm,yshift=-0.62cm]frame.north west) {Question 69};
\node[anchor=east,text=white,font=\sffamily\itshape\large] at ([xshift=-5mm,yshift=-0.62cm]frame.north east) {Arithmétique et ensembles de nombres};}]
\vspace*{1.75cm}
\begin{minipage}[t][6.9cm][c]{\linewidth}\LARGE\centering
\begin{minipage}{0.95\linewidth}\raggedright\linespread{1.25}\selectfont Quel est le plus petit ensemble de nombres auquel appartient $-7$~?\end{minipage}
\end{minipage}

\vfill
\begin{tcolorbox}[enhanced,nobeforeafter,width=7.1cm,height=2.1cm,valign=center,colback=cari!6,colframe=cari,arc=2.5mm,boxrule=0.9pt,left=3mm]{\Large\bfseries\color{cari}a)}\quad\begin{minipage}[c]{5.3cm}\raggedright\Large $\mathbb{Z}$\end{minipage}\end{tcolorbox}\hfill\begin{tcolorbox}[enhanced,nobeforeafter,width=7.1cm,height=2.1cm,valign=center,colback=cari!6,colframe=cari,arc=2.5mm,boxrule=0.9pt,left=3mm]{\Large\bfseries\color{cari}b)}\quad\begin{minipage}[c]{5.3cm}\raggedright\Large $\mathbb{D}$\end{minipage}\end{tcolorbox}\par\vspace{2.5mm}
\begin{tcolorbox}[enhanced,nobeforeafter,width=7.1cm,height=2.1cm,valign=center,colback=cari!6,colframe=cari,arc=2.5mm,boxrule=0.9pt,left=3mm]{\Large\bfseries\color{cari}c)}\quad\begin{minipage}[c]{5.3cm}\raggedright\Large $\mathbb{N}$\end{minipage}\end{tcolorbox}\hfill\begin{tcolorbox}[enhanced,nobeforeafter,width=7.1cm,height=2.1cm,valign=center,colback=cari!6,colframe=cari,arc=2.5mm,boxrule=0.9pt,left=3mm]{\Large\bfseries\color{cari}d)}\quad\begin{minipage}[c]{5.3cm}\raggedright\Large $\mathbb{Q}$\end{minipage}\end{tcolorbox}
\end{tcolorbox}
\newpage
\begin{tcolorbox}[enhanced,width=16cm,height=14cm,colback=white,colframe=cari,boxrule=2pt,arc=0mm,
left=6mm,right=6mm,top=0mm,bottom=4mm,sharp corners,
overlay={\fill[cari] (frame.north west) rectangle ([yshift=-1.25cm]frame.north east);
\node[anchor=west,text=white,font=\sffamily\bfseries\LARGE] at ([xshift=5mm,yshift=-0.62cm]frame.north west) {Question 70};
\node[anchor=east,text=white,font=\sffamily\itshape\large] at ([xshift=-5mm,yshift=-0.62cm]frame.north east) {Arithmétique et ensembles de nombres};}]
\vspace*{1.75cm}
\begin{minipage}[t][6.9cm][c]{\linewidth}\LARGE\centering
\begin{minipage}{0.95\linewidth}\raggedright\linespread{1.25}\selectfont Quel est le plus petit ensemble de nombres auquel appartient $\dfrac{3}{4}$~?\end{minipage}
\end{minipage}

\vfill
\begin{tcolorbox}[enhanced,nobeforeafter,width=7.1cm,height=2.1cm,valign=center,colback=cari!6,colframe=cari,arc=2.5mm,boxrule=0.9pt,left=3mm]{\Large\bfseries\color{cari}a)}\quad\begin{minipage}[c]{5.3cm}\raggedright\Large $\mathbb{N}$\end{minipage}\end{tcolorbox}\hfill\begin{tcolorbox}[enhanced,nobeforeafter,width=7.1cm,height=2.1cm,valign=center,colback=cari!6,colframe=cari,arc=2.5mm,boxrule=0.9pt,left=3mm]{\Large\bfseries\color{cari}b)}\quad\begin{minipage}[c]{5.3cm}\raggedright\Large $\mathbb{D}$\end{minipage}\end{tcolorbox}\par\vspace{2.5mm}
\begin{tcolorbox}[enhanced,nobeforeafter,width=7.1cm,height=2.1cm,valign=center,colback=cari!6,colframe=cari,arc=2.5mm,boxrule=0.9pt,left=3mm]{\Large\bfseries\color{cari}c)}\quad\begin{minipage}[c]{5.3cm}\raggedright\Large $\mathbb{Q}$\end{minipage}\end{tcolorbox}\hfill\begin{tcolorbox}[enhanced,nobeforeafter,width=7.1cm,height=2.1cm,valign=center,colback=cari!6,colframe=cari,arc=2.5mm,boxrule=0.9pt,left=3mm]{\Large\bfseries\color{cari}d)}\quad\begin{minipage}[c]{5.3cm}\raggedright\Large $\mathbb{Z}$\end{minipage}\end{tcolorbox}
\end{tcolorbox}
\newpage
\begin{tcolorbox}[enhanced,width=16cm,height=14cm,colback=white,colframe=cari,boxrule=2pt,arc=0mm,
left=6mm,right=6mm,top=0mm,bottom=4mm,sharp corners,
overlay={\fill[cari] (frame.north west) rectangle ([yshift=-1.25cm]frame.north east);
\node[anchor=west,text=white,font=\sffamily\bfseries\LARGE] at ([xshift=5mm,yshift=-0.62cm]frame.north west) {Question 71};
\node[anchor=east,text=white,font=\sffamily\itshape\large] at ([xshift=-5mm,yshift=-0.62cm]frame.north east) {Arithmétique et ensembles de nombres};}]
\vspace*{1.75cm}
\begin{minipage}[t][6.9cm][c]{\linewidth}\LARGE\centering
\begin{minipage}{0.95\linewidth}\raggedright\linespread{1.25}\selectfont Quelle est la forme irréductible de $\dfrac{84}{126}$~?\end{minipage}
\end{minipage}

\vfill
\begin{tcolorbox}[enhanced,nobeforeafter,width=7.1cm,height=2.1cm,valign=center,colback=cari!6,colframe=cari,arc=2.5mm,boxrule=0.9pt,left=3mm]{\Large\bfseries\color{cari}a)}\quad\begin{minipage}[c]{5.3cm}\raggedright\Large $\dfrac{2}{3}$\end{minipage}\end{tcolorbox}\hfill\begin{tcolorbox}[enhanced,nobeforeafter,width=7.1cm,height=2.1cm,valign=center,colback=cari!6,colframe=cari,arc=2.5mm,boxrule=0.9pt,left=3mm]{\Large\bfseries\color{cari}b)}\quad\begin{minipage}[c]{5.3cm}\raggedright\Large $\dfrac{42}{63}$\end{minipage}\end{tcolorbox}\par\vspace{2.5mm}
\begin{tcolorbox}[enhanced,nobeforeafter,width=7.1cm,height=2.1cm,valign=center,colback=cari!6,colframe=cari,arc=2.5mm,boxrule=0.9pt,left=3mm]{\Large\bfseries\color{cari}c)}\quad\begin{minipage}[c]{5.3cm}\raggedright\Large $\dfrac{12}{18}$\end{minipage}\end{tcolorbox}\hfill\begin{tcolorbox}[enhanced,nobeforeafter,width=7.1cm,height=2.1cm,valign=center,colback=cari!6,colframe=cari,arc=2.5mm,boxrule=0.9pt,left=3mm]{\Large\bfseries\color{cari}d)}\quad\begin{minipage}[c]{5.3cm}\raggedright\Large $\dfrac{28}{42}$\end{minipage}\end{tcolorbox}
\end{tcolorbox}
\newpage
\begin{tcolorbox}[enhanced,width=16cm,height=14cm,colback=white,colframe=cari,boxrule=2pt,arc=0mm,
left=6mm,right=6mm,top=0mm,bottom=4mm,sharp corners,
overlay={\fill[cari] (frame.north west) rectangle ([yshift=-1.25cm]frame.north east);
\node[anchor=west,text=white,font=\sffamily\bfseries\LARGE] at ([xshift=5mm,yshift=-0.62cm]frame.north west) {Question 72};
\node[anchor=east,text=white,font=\sffamily\itshape\large] at ([xshift=-5mm,yshift=-0.62cm]frame.north east) {Arithmétique et ensembles de nombres};}]
\vspace*{1.75cm}
\begin{minipage}[t][6.9cm][c]{\linewidth}\LARGE\centering
\begin{minipage}{0.95\linewidth}\raggedright\linespread{1.25}\selectfont Quelle est la forme irréductible de $\dfrac{45}{75}$~?\end{minipage}
\end{minipage}

\vfill
\begin{tcolorbox}[enhanced,nobeforeafter,width=7.1cm,height=2.1cm,valign=center,colback=cari!6,colframe=cari,arc=2.5mm,boxrule=0.9pt,left=3mm]{\Large\bfseries\color{cari}a)}\quad\begin{minipage}[c]{5.3cm}\raggedright\Large $\dfrac{9}{15}$\end{minipage}\end{tcolorbox}\hfill\begin{tcolorbox}[enhanced,nobeforeafter,width=7.1cm,height=2.1cm,valign=center,colback=cari!6,colframe=cari,arc=2.5mm,boxrule=0.9pt,left=3mm]{\Large\bfseries\color{cari}b)}\quad\begin{minipage}[c]{5.3cm}\raggedright\Large $\dfrac{5}{3}$\end{minipage}\end{tcolorbox}\par\vspace{2.5mm}
\begin{tcolorbox}[enhanced,nobeforeafter,width=7.1cm,height=2.1cm,valign=center,colback=cari!6,colframe=cari,arc=2.5mm,boxrule=0.9pt,left=3mm]{\Large\bfseries\color{cari}c)}\quad\begin{minipage}[c]{5.3cm}\raggedright\Large $\dfrac{15}{25}$\end{minipage}\end{tcolorbox}\hfill\begin{tcolorbox}[enhanced,nobeforeafter,width=7.1cm,height=2.1cm,valign=center,colback=cari!6,colframe=cari,arc=2.5mm,boxrule=0.9pt,left=3mm]{\Large\bfseries\color{cari}d)}\quad\begin{minipage}[c]{5.3cm}\raggedright\Large $\dfrac{3}{5}$\end{minipage}\end{tcolorbox}
\end{tcolorbox}
\newpage
\begin{tcolorbox}[enhanced,width=16cm,height=14cm,colback=white,colframe=cari,boxrule=2pt,arc=0mm,
left=6mm,right=6mm,top=0mm,bottom=4mm,sharp corners,
overlay={\fill[cari] (frame.north west) rectangle ([yshift=-1.25cm]frame.north east);
\node[anchor=west,text=white,font=\sffamily\bfseries\LARGE] at ([xshift=5mm,yshift=-0.62cm]frame.north west) {Question 73};
\node[anchor=east,text=white,font=\sffamily\itshape\large] at ([xshift=-5mm,yshift=-0.62cm]frame.north east) {Arithmétique et ensembles de nombres};}]
\vspace*{1.75cm}
\begin{minipage}[t][6.9cm][c]{\linewidth}\LARGE\centering
\begin{minipage}{0.95\linewidth}\raggedright\linespread{1.25}\selectfont Pour tout entier naturel $n$, que peut-on dire de l'entier $2n+4$~?\end{minipage}
\end{minipage}

\vfill
\begin{tcolorbox}[enhanced,nobeforeafter,width=7.1cm,height=2.1cm,valign=center,colback=cari!6,colframe=cari,arc=2.5mm,boxrule=0.9pt,left=3mm]{\Large\bfseries\color{cari}a)}\quad\begin{minipage}[c]{5.3cm}\raggedright\Large il est toujours premier\end{minipage}\end{tcolorbox}\hfill\begin{tcolorbox}[enhanced,nobeforeafter,width=7.1cm,height=2.1cm,valign=center,colback=cari!6,colframe=cari,arc=2.5mm,boxrule=0.9pt,left=3mm]{\Large\bfseries\color{cari}b)}\quad\begin{minipage}[c]{5.3cm}\raggedright\large il est pair ou impair selon $n$\end{minipage}\end{tcolorbox}\par\vspace{2.5mm}
\begin{tcolorbox}[enhanced,nobeforeafter,width=7.1cm,height=2.1cm,valign=center,colback=cari!6,colframe=cari,arc=2.5mm,boxrule=0.9pt,left=3mm]{\Large\bfseries\color{cari}c)}\quad\begin{minipage}[c]{5.3cm}\raggedright\Large il est toujours pair\end{minipage}\end{tcolorbox}\hfill\begin{tcolorbox}[enhanced,nobeforeafter,width=7.1cm,height=2.1cm,valign=center,colback=cari!6,colframe=cari,arc=2.5mm,boxrule=0.9pt,left=3mm]{\Large\bfseries\color{cari}d)}\quad\begin{minipage}[c]{5.3cm}\raggedright\Large il est toujours impair\end{minipage}\end{tcolorbox}
\end{tcolorbox}
\newpage
\begin{tcolorbox}[enhanced,width=16cm,height=14cm,colback=white,colframe=cari,boxrule=2pt,arc=0mm,
left=6mm,right=6mm,top=0mm,bottom=4mm,sharp corners,
overlay={\fill[cari] (frame.north west) rectangle ([yshift=-1.25cm]frame.north east);
\node[anchor=west,text=white,font=\sffamily\bfseries\LARGE] at ([xshift=5mm,yshift=-0.62cm]frame.north west) {Question 74};
\node[anchor=east,text=white,font=\sffamily\itshape\large] at ([xshift=-5mm,yshift=-0.62cm]frame.north east) {Arithmétique et ensembles de nombres};}]
\vspace*{1.75cm}
\begin{minipage}[t][6.9cm][c]{\linewidth}\LARGE\centering
\begin{minipage}{0.95\linewidth}\raggedright\linespread{1.25}\selectfont Pour tout entier naturel $n$, que peut-on dire de l'entier $2n+1$~?\end{minipage}
\end{minipage}

\vfill
\begin{tcolorbox}[enhanced,nobeforeafter,width=7.1cm,height=2.1cm,valign=center,colback=cari!6,colframe=cari,arc=2.5mm,boxrule=0.9pt,left=3mm]{\Large\bfseries\color{cari}a)}\quad\begin{minipage}[c]{5.3cm}\raggedright\large il est pair ou impair selon $n$\end{minipage}\end{tcolorbox}\hfill\begin{tcolorbox}[enhanced,nobeforeafter,width=7.1cm,height=2.1cm,valign=center,colback=cari!6,colframe=cari,arc=2.5mm,boxrule=0.9pt,left=3mm]{\Large\bfseries\color{cari}b)}\quad\begin{minipage}[c]{5.3cm}\raggedright\Large il est toujours impair\end{minipage}\end{tcolorbox}\par\vspace{2.5mm}
\begin{tcolorbox}[enhanced,nobeforeafter,width=7.1cm,height=2.1cm,valign=center,colback=cari!6,colframe=cari,arc=2.5mm,boxrule=0.9pt,left=3mm]{\Large\bfseries\color{cari}c)}\quad\begin{minipage}[c]{5.3cm}\raggedright\Large il est toujours premier\end{minipage}\end{tcolorbox}\hfill\begin{tcolorbox}[enhanced,nobeforeafter,width=7.1cm,height=2.1cm,valign=center,colback=cari!6,colframe=cari,arc=2.5mm,boxrule=0.9pt,left=3mm]{\Large\bfseries\color{cari}d)}\quad\begin{minipage}[c]{5.3cm}\raggedright\Large il est toujours pair\end{minipage}\end{tcolorbox}
\end{tcolorbox}
\newpage
\begin{tcolorbox}[enhanced,width=16cm,height=14cm,colback=white,colframe=cari,boxrule=2pt,arc=0mm,
left=6mm,right=6mm,top=0mm,bottom=4mm,sharp corners,
overlay={\fill[cari] (frame.north west) rectangle ([yshift=-1.25cm]frame.north east);
\node[anchor=west,text=white,font=\sffamily\bfseries\LARGE] at ([xshift=5mm,yshift=-0.62cm]frame.north west) {Question 75};
\node[anchor=east,text=white,font=\sffamily\itshape\large] at ([xshift=-5mm,yshift=-0.62cm]frame.north east) {Arithmétique et ensembles de nombres};}]
\vspace*{1.75cm}
\begin{minipage}[t][6.9cm][c]{\linewidth}\LARGE\centering
\begin{minipage}{0.95\linewidth}\raggedright\linespread{1.25}\selectfont Combien y a-t-il de multiples de $7$ compris entre $1$ et $100$~?\end{minipage}
\end{minipage}

\vfill
\begin{tcolorbox}[enhanced,nobeforeafter,width=7.1cm,height=2.1cm,valign=center,colback=cari!6,colframe=cari,arc=2.5mm,boxrule=0.9pt,left=3mm]{\Large\bfseries\color{cari}a)}\quad\begin{minipage}[c]{5.3cm}\raggedright\Large $14$\end{minipage}\end{tcolorbox}\hfill\begin{tcolorbox}[enhanced,nobeforeafter,width=7.1cm,height=2.1cm,valign=center,colback=cari!6,colframe=cari,arc=2.5mm,boxrule=0.9pt,left=3mm]{\Large\bfseries\color{cari}b)}\quad\begin{minipage}[c]{5.3cm}\raggedright\Large $15$\end{minipage}\end{tcolorbox}\par\vspace{2.5mm}
\begin{tcolorbox}[enhanced,nobeforeafter,width=7.1cm,height=2.1cm,valign=center,colback=cari!6,colframe=cari,arc=2.5mm,boxrule=0.9pt,left=3mm]{\Large\bfseries\color{cari}c)}\quad\begin{minipage}[c]{5.3cm}\raggedright\Large $13$\end{minipage}\end{tcolorbox}\hfill\begin{tcolorbox}[enhanced,nobeforeafter,width=7.1cm,height=2.1cm,valign=center,colback=cari!6,colframe=cari,arc=2.5mm,boxrule=0.9pt,left=3mm]{\Large\bfseries\color{cari}d)}\quad\begin{minipage}[c]{5.3cm}\raggedright\Large $86$\end{minipage}\end{tcolorbox}
\end{tcolorbox}
\newpage
\begin{tcolorbox}[enhanced,width=16cm,height=14cm,colback=white,colframe=cari,boxrule=2pt,arc=0mm,
left=6mm,right=6mm,top=0mm,bottom=4mm,sharp corners,
overlay={\fill[cari] (frame.north west) rectangle ([yshift=-1.25cm]frame.north east);
\node[anchor=west,text=white,font=\sffamily\bfseries\LARGE] at ([xshift=5mm,yshift=-0.62cm]frame.north west) {Question 76};
\node[anchor=east,text=white,font=\sffamily\itshape\large] at ([xshift=-5mm,yshift=-0.62cm]frame.north east) {Arithmétique et ensembles de nombres};}]
\vspace*{1.75cm}
\begin{minipage}[t][6.9cm][c]{\linewidth}\LARGE\centering
\begin{minipage}{0.95\linewidth}\raggedright\linespread{1.25}\selectfont Combien y a-t-il de multiples de $9$ compris entre $1$ et $200$~?\end{minipage}
\end{minipage}

\vfill
\begin{tcolorbox}[enhanced,nobeforeafter,width=7.1cm,height=2.1cm,valign=center,colback=cari!6,colframe=cari,arc=2.5mm,boxrule=0.9pt,left=3mm]{\Large\bfseries\color{cari}a)}\quad\begin{minipage}[c]{5.3cm}\raggedright\Large $21$\end{minipage}\end{tcolorbox}\hfill\begin{tcolorbox}[enhanced,nobeforeafter,width=7.1cm,height=2.1cm,valign=center,colback=cari!6,colframe=cari,arc=2.5mm,boxrule=0.9pt,left=3mm]{\Large\bfseries\color{cari}b)}\quad\begin{minipage}[c]{5.3cm}\raggedright\Large $22$\end{minipage}\end{tcolorbox}\par\vspace{2.5mm}
\begin{tcolorbox}[enhanced,nobeforeafter,width=7.1cm,height=2.1cm,valign=center,colback=cari!6,colframe=cari,arc=2.5mm,boxrule=0.9pt,left=3mm]{\Large\bfseries\color{cari}c)}\quad\begin{minipage}[c]{5.3cm}\raggedright\Large $23$\end{minipage}\end{tcolorbox}\hfill\begin{tcolorbox}[enhanced,nobeforeafter,width=7.1cm,height=2.1cm,valign=center,colback=cari!6,colframe=cari,arc=2.5mm,boxrule=0.9pt,left=3mm]{\Large\bfseries\color{cari}d)}\quad\begin{minipage}[c]{5.3cm}\raggedright\Large $178$\end{minipage}\end{tcolorbox}
\end{tcolorbox}
\newpage
\begin{tcolorbox}[enhanced,width=16cm,height=14cm,colback=white,colframe=cari,boxrule=2pt,arc=0mm,
left=6mm,right=6mm,top=0mm,bottom=4mm,sharp corners,
overlay={\fill[cari] (frame.north west) rectangle ([yshift=-1.25cm]frame.north east);
\node[anchor=west,text=white,font=\sffamily\bfseries\LARGE] at ([xshift=5mm,yshift=-0.62cm]frame.north west) {Question 77};
\node[anchor=east,text=white,font=\sffamily\itshape\large] at ([xshift=-5mm,yshift=-0.62cm]frame.north east) {Arithmétique et ensembles de nombres};}]
\vspace*{1.75cm}
\begin{minipage}[t][6.9cm][c]{\linewidth}\LARGE\centering
\begin{minipage}{0.95\linewidth}\raggedright\linespread{1.25}\selectfont Combien l'entier $36$ a-t-il de diviseurs positifs~?\end{minipage}
\end{minipage}

\vfill
\begin{tcolorbox}[enhanced,nobeforeafter,width=7.1cm,height=2.1cm,valign=center,colback=cari!6,colframe=cari,arc=2.5mm,boxrule=0.9pt,left=3mm]{\Large\bfseries\color{cari}a)}\quad\begin{minipage}[c]{5.3cm}\raggedright\Large $10$\end{minipage}\end{tcolorbox}\hfill\begin{tcolorbox}[enhanced,nobeforeafter,width=7.1cm,height=2.1cm,valign=center,colback=cari!6,colframe=cari,arc=2.5mm,boxrule=0.9pt,left=3mm]{\Large\bfseries\color{cari}b)}\quad\begin{minipage}[c]{5.3cm}\raggedright\Large $8$\end{minipage}\end{tcolorbox}\par\vspace{2.5mm}
\begin{tcolorbox}[enhanced,nobeforeafter,width=7.1cm,height=2.1cm,valign=center,colback=cari!6,colframe=cari,arc=2.5mm,boxrule=0.9pt,left=3mm]{\Large\bfseries\color{cari}c)}\quad\begin{minipage}[c]{5.3cm}\raggedright\Large $9$\end{minipage}\end{tcolorbox}\hfill\begin{tcolorbox}[enhanced,nobeforeafter,width=7.1cm,height=2.1cm,valign=center,colback=cari!6,colframe=cari,arc=2.5mm,boxrule=0.9pt,left=3mm]{\Large\bfseries\color{cari}d)}\quad\begin{minipage}[c]{5.3cm}\raggedright\Large $11$\end{minipage}\end{tcolorbox}
\end{tcolorbox}
\newpage
\begin{tcolorbox}[enhanced,width=16cm,height=14cm,colback=white,colframe=cari,boxrule=2pt,arc=0mm,
left=6mm,right=6mm,top=0mm,bottom=4mm,sharp corners,
overlay={\fill[cari] (frame.north west) rectangle ([yshift=-1.25cm]frame.north east);
\node[anchor=west,text=white,font=\sffamily\bfseries\LARGE] at ([xshift=5mm,yshift=-0.62cm]frame.north west) {Question 78};
\node[anchor=east,text=white,font=\sffamily\itshape\large] at ([xshift=-5mm,yshift=-0.62cm]frame.north east) {Arithmétique et ensembles de nombres};}]
\vspace*{1.75cm}
\begin{minipage}[t][6.9cm][c]{\linewidth}\LARGE\centering
\begin{minipage}{0.95\linewidth}\raggedright\linespread{1.25}\selectfont Combien l'entier $24$ a-t-il de diviseurs positifs~?\end{minipage}
\end{minipage}

\vfill
\begin{tcolorbox}[enhanced,nobeforeafter,width=7.1cm,height=2.1cm,valign=center,colback=cari!6,colframe=cari,arc=2.5mm,boxrule=0.9pt,left=3mm]{\Large\bfseries\color{cari}a)}\quad\begin{minipage}[c]{5.3cm}\raggedright\Large $7$\end{minipage}\end{tcolorbox}\hfill\begin{tcolorbox}[enhanced,nobeforeafter,width=7.1cm,height=2.1cm,valign=center,colback=cari!6,colframe=cari,arc=2.5mm,boxrule=0.9pt,left=3mm]{\Large\bfseries\color{cari}b)}\quad\begin{minipage}[c]{5.3cm}\raggedright\Large $10$\end{minipage}\end{tcolorbox}\par\vspace{2.5mm}
\begin{tcolorbox}[enhanced,nobeforeafter,width=7.1cm,height=2.1cm,valign=center,colback=cari!6,colframe=cari,arc=2.5mm,boxrule=0.9pt,left=3mm]{\Large\bfseries\color{cari}c)}\quad\begin{minipage}[c]{5.3cm}\raggedright\Large $8$\end{minipage}\end{tcolorbox}\hfill\begin{tcolorbox}[enhanced,nobeforeafter,width=7.1cm,height=2.1cm,valign=center,colback=cari!6,colframe=cari,arc=2.5mm,boxrule=0.9pt,left=3mm]{\Large\bfseries\color{cari}d)}\quad\begin{minipage}[c]{5.3cm}\raggedright\Large $9$\end{minipage}\end{tcolorbox}
\end{tcolorbox}
\newpage
\begin{tcolorbox}[enhanced,width=16cm,height=14cm,colback=white,colframe=cari,boxrule=2pt,arc=0mm,
left=6mm,right=6mm,top=0mm,bottom=4mm,sharp corners,
overlay={\fill[cari] (frame.north west) rectangle ([yshift=-1.25cm]frame.north east);
\node[anchor=west,text=white,font=\sffamily\bfseries\LARGE] at ([xshift=5mm,yshift=-0.62cm]frame.north west) {Question 79};
\node[anchor=east,text=white,font=\sffamily\itshape\large] at ([xshift=-5mm,yshift=-0.62cm]frame.north east) {Arithmétique et ensembles de nombres};}]
\vspace*{1.75cm}
\begin{minipage}[t][6.9cm][c]{\linewidth}\LARGE\centering
\begin{minipage}{0.95\linewidth}\raggedright\linespread{1.25}\selectfont Que vaut l'intersection $[-3\,;\,5]\cap[1\,;\,8]$~?\end{minipage}
\end{minipage}

\vfill
\begin{tcolorbox}[enhanced,nobeforeafter,width=7.1cm,height=2.1cm,valign=center,colback=cari!6,colframe=cari,arc=2.5mm,boxrule=0.9pt,left=3mm]{\Large\bfseries\color{cari}a)}\quad\begin{minipage}[c]{5.3cm}\raggedright\Large $[-3\,;\,8]$\end{minipage}\end{tcolorbox}\hfill\begin{tcolorbox}[enhanced,nobeforeafter,width=7.1cm,height=2.1cm,valign=center,colback=cari!6,colframe=cari,arc=2.5mm,boxrule=0.9pt,left=3mm]{\Large\bfseries\color{cari}b)}\quad\begin{minipage}[c]{5.3cm}\raggedright\Large $[1\,;\,5]$\end{minipage}\end{tcolorbox}\par\vspace{2.5mm}
\begin{tcolorbox}[enhanced,nobeforeafter,width=7.1cm,height=2.1cm,valign=center,colback=cari!6,colframe=cari,arc=2.5mm,boxrule=0.9pt,left=3mm]{\Large\bfseries\color{cari}c)}\quad\begin{minipage}[c]{5.3cm}\raggedright\Large $[1\,;\,8]$\end{minipage}\end{tcolorbox}\hfill\begin{tcolorbox}[enhanced,nobeforeafter,width=7.1cm,height=2.1cm,valign=center,colback=cari!6,colframe=cari,arc=2.5mm,boxrule=0.9pt,left=3mm]{\Large\bfseries\color{cari}d)}\quad\begin{minipage}[c]{5.3cm}\raggedright\Large $[-3\,;\,5]$\end{minipage}\end{tcolorbox}
\end{tcolorbox}
\newpage
\begin{tcolorbox}[enhanced,width=16cm,height=14cm,colback=white,colframe=cari,boxrule=2pt,arc=0mm,
left=6mm,right=6mm,top=0mm,bottom=4mm,sharp corners,
overlay={\fill[cari] (frame.north west) rectangle ([yshift=-1.25cm]frame.north east);
\node[anchor=west,text=white,font=\sffamily\bfseries\LARGE] at ([xshift=5mm,yshift=-0.62cm]frame.north west) {Question 80};
\node[anchor=east,text=white,font=\sffamily\itshape\large] at ([xshift=-5mm,yshift=-0.62cm]frame.north east) {Arithmétique et ensembles de nombres};}]
\vspace*{1.75cm}
\begin{minipage}[t][6.9cm][c]{\linewidth}\LARGE\centering
\begin{minipage}{0.95\linewidth}\raggedright\linespread{1.25}\selectfont Que vaut l'intersection $[0\,;\,4[\cap[2\,;\,9]$~?\end{minipage}
\end{minipage}

\vfill
\begin{tcolorbox}[enhanced,nobeforeafter,width=7.1cm,height=2.1cm,valign=center,colback=cari!6,colframe=cari,arc=2.5mm,boxrule=0.9pt,left=3mm]{\Large\bfseries\color{cari}a)}\quad\begin{minipage}[c]{5.3cm}\raggedright\Large $[2\,;\,4]$\end{minipage}\end{tcolorbox}\hfill\begin{tcolorbox}[enhanced,nobeforeafter,width=7.1cm,height=2.1cm,valign=center,colback=cari!6,colframe=cari,arc=2.5mm,boxrule=0.9pt,left=3mm]{\Large\bfseries\color{cari}b)}\quad\begin{minipage}[c]{5.3cm}\raggedright\Large $]2\,;\,4[$\end{minipage}\end{tcolorbox}\par\vspace{2.5mm}
\begin{tcolorbox}[enhanced,nobeforeafter,width=7.1cm,height=2.1cm,valign=center,colback=cari!6,colframe=cari,arc=2.5mm,boxrule=0.9pt,left=3mm]{\Large\bfseries\color{cari}c)}\quad\begin{minipage}[c]{5.3cm}\raggedright\Large $[2\,;\,4[$\end{minipage}\end{tcolorbox}\hfill\begin{tcolorbox}[enhanced,nobeforeafter,width=7.1cm,height=2.1cm,valign=center,colback=cari!6,colframe=cari,arc=2.5mm,boxrule=0.9pt,left=3mm]{\Large\bfseries\color{cari}d)}\quad\begin{minipage}[c]{5.3cm}\raggedright\Large $[0\,;\,9]$\end{minipage}\end{tcolorbox}
\end{tcolorbox}
\newpage
\begin{tcolorbox}[enhanced,width=16cm,height=14cm,colback=white,colframe=cpro,boxrule=2pt,arc=0mm,
left=6mm,right=6mm,top=0mm,bottom=4mm,sharp corners,
overlay={\fill[cpro] (frame.north west) rectangle ([yshift=-1.25cm]frame.north east);
\node[anchor=west,text=white,font=\sffamily\bfseries\LARGE] at ([xshift=5mm,yshift=-0.62cm]frame.north west) {Question 81};
\node[anchor=east,text=white,font=\sffamily\itshape\large] at ([xshift=-5mm,yshift=-0.62cm]frame.north east) {Probabilités};}]
\vspace*{1.75cm}
\begin{minipage}[t][6.9cm][c]{\linewidth}\LARGE\centering
\begin{minipage}{0.95\linewidth}\raggedright\linespread{1.25}\selectfont On a $P(A)=0{,}3$, $P(B)=0{,}5$ et $P(A\cap B)=0{,}1$. Que vaut $P(A\cup B)$~?\end{minipage}
\end{minipage}

\vfill
\begin{tcolorbox}[enhanced,nobeforeafter,width=7.1cm,height=2.1cm,valign=center,colback=cpro!6,colframe=cpro,arc=2.5mm,boxrule=0.9pt,left=3mm]{\Large\bfseries\color{cpro}a)}\quad\begin{minipage}[c]{5.3cm}\raggedright\Large $0{,}1$\end{minipage}\end{tcolorbox}\hfill\begin{tcolorbox}[enhanced,nobeforeafter,width=7.1cm,height=2.1cm,valign=center,colback=cpro!6,colframe=cpro,arc=2.5mm,boxrule=0.9pt,left=3mm]{\Large\bfseries\color{cpro}b)}\quad\begin{minipage}[c]{5.3cm}\raggedright\Large $0{,}8$\end{minipage}\end{tcolorbox}\par\vspace{2.5mm}
\begin{tcolorbox}[enhanced,nobeforeafter,width=7.1cm,height=2.1cm,valign=center,colback=cpro!6,colframe=cpro,arc=2.5mm,boxrule=0.9pt,left=3mm]{\Large\bfseries\color{cpro}c)}\quad\begin{minipage}[c]{5.3cm}\raggedright\Large $0{,}9$\end{minipage}\end{tcolorbox}\hfill\begin{tcolorbox}[enhanced,nobeforeafter,width=7.1cm,height=2.1cm,valign=center,colback=cpro!6,colframe=cpro,arc=2.5mm,boxrule=0.9pt,left=3mm]{\Large\bfseries\color{cpro}d)}\quad\begin{minipage}[c]{5.3cm}\raggedright\Large $0{,}7$\end{minipage}\end{tcolorbox}
\end{tcolorbox}
\newpage
\begin{tcolorbox}[enhanced,width=16cm,height=14cm,colback=white,colframe=cpro,boxrule=2pt,arc=0mm,
left=6mm,right=6mm,top=0mm,bottom=4mm,sharp corners,
overlay={\fill[cpro] (frame.north west) rectangle ([yshift=-1.25cm]frame.north east);
\node[anchor=west,text=white,font=\sffamily\bfseries\LARGE] at ([xshift=5mm,yshift=-0.62cm]frame.north west) {Question 82};
\node[anchor=east,text=white,font=\sffamily\itshape\large] at ([xshift=-5mm,yshift=-0.62cm]frame.north east) {Probabilités};}]
\vspace*{1.75cm}
\begin{minipage}[t][6.9cm][c]{\linewidth}\LARGE\centering
\begin{minipage}{0.95\linewidth}\raggedright\linespread{1.25}\selectfont On a $P(A)=0{,}4$, $P(B)=0{,}35$ et $P(A\cap B)=0{,}15$. Que vaut $P(A\cup B)$~?\end{minipage}
\end{minipage}

\vfill
\begin{tcolorbox}[enhanced,nobeforeafter,width=7.1cm,height=2.1cm,valign=center,colback=cpro!6,colframe=cpro,arc=2.5mm,boxrule=0.9pt,left=3mm]{\Large\bfseries\color{cpro}a)}\quad\begin{minipage}[c]{5.3cm}\raggedright\Large $0{,}6$\end{minipage}\end{tcolorbox}\hfill\begin{tcolorbox}[enhanced,nobeforeafter,width=7.1cm,height=2.1cm,valign=center,colback=cpro!6,colframe=cpro,arc=2.5mm,boxrule=0.9pt,left=3mm]{\Large\bfseries\color{cpro}b)}\quad\begin{minipage}[c]{5.3cm}\raggedright\Large $0{,}15$\end{minipage}\end{tcolorbox}\par\vspace{2.5mm}
\begin{tcolorbox}[enhanced,nobeforeafter,width=7.1cm,height=2.1cm,valign=center,colback=cpro!6,colframe=cpro,arc=2.5mm,boxrule=0.9pt,left=3mm]{\Large\bfseries\color{cpro}c)}\quad\begin{minipage}[c]{5.3cm}\raggedright\Large $0{,}75$\end{minipage}\end{tcolorbox}\hfill\begin{tcolorbox}[enhanced,nobeforeafter,width=7.1cm,height=2.1cm,valign=center,colback=cpro!6,colframe=cpro,arc=2.5mm,boxrule=0.9pt,left=3mm]{\Large\bfseries\color{cpro}d)}\quad\begin{minipage}[c]{5.3cm}\raggedright\Large $0{,}9$\end{minipage}\end{tcolorbox}
\end{tcolorbox}
\newpage
\begin{tcolorbox}[enhanced,width=16cm,height=14cm,colback=white,colframe=cpro,boxrule=2pt,arc=0mm,
left=6mm,right=6mm,top=0mm,bottom=4mm,sharp corners,
overlay={\fill[cpro] (frame.north west) rectangle ([yshift=-1.25cm]frame.north east);
\node[anchor=west,text=white,font=\sffamily\bfseries\LARGE] at ([xshift=5mm,yshift=-0.62cm]frame.north west) {Question 83};
\node[anchor=east,text=white,font=\sffamily\itshape\large] at ([xshift=-5mm,yshift=-0.62cm]frame.north east) {Probabilités};}]
\vspace*{1.75cm}
\begin{minipage}[t][6.9cm][c]{\linewidth}\LARGE\centering
\begin{minipage}{0.95\linewidth}\raggedright\linespread{1.25}\selectfont On a $P(A)=0{,}5$, $P(B)=0{,}4$ et $P(A\cup B)=0{,}7$. Que vaut $P(A\cap B)$~?\end{minipage}
\end{minipage}

\vfill
\begin{tcolorbox}[enhanced,nobeforeafter,width=7.1cm,height=2.1cm,valign=center,colback=cpro!6,colframe=cpro,arc=2.5mm,boxrule=0.9pt,left=3mm]{\Large\bfseries\color{cpro}a)}\quad\begin{minipage}[c]{5.3cm}\raggedright\Large $0{,}2$\end{minipage}\end{tcolorbox}\hfill\begin{tcolorbox}[enhanced,nobeforeafter,width=7.1cm,height=2.1cm,valign=center,colback=cpro!6,colframe=cpro,arc=2.5mm,boxrule=0.9pt,left=3mm]{\Large\bfseries\color{cpro}b)}\quad\begin{minipage}[c]{5.3cm}\raggedright\Large $0{,}3$\end{minipage}\end{tcolorbox}\par\vspace{2.5mm}
\begin{tcolorbox}[enhanced,nobeforeafter,width=7.1cm,height=2.1cm,valign=center,colback=cpro!6,colframe=cpro,arc=2.5mm,boxrule=0.9pt,left=3mm]{\Large\bfseries\color{cpro}c)}\quad\begin{minipage}[c]{5.3cm}\raggedright\Large $0{,}7$\end{minipage}\end{tcolorbox}\hfill\begin{tcolorbox}[enhanced,nobeforeafter,width=7.1cm,height=2.1cm,valign=center,colback=cpro!6,colframe=cpro,arc=2.5mm,boxrule=0.9pt,left=3mm]{\Large\bfseries\color{cpro}d)}\quad\begin{minipage}[c]{5.3cm}\raggedright\Large $0{,}4$\end{minipage}\end{tcolorbox}
\end{tcolorbox}
\newpage
\begin{tcolorbox}[enhanced,width=16cm,height=14cm,colback=white,colframe=cpro,boxrule=2pt,arc=0mm,
left=6mm,right=6mm,top=0mm,bottom=4mm,sharp corners,
overlay={\fill[cpro] (frame.north west) rectangle ([yshift=-1.25cm]frame.north east);
\node[anchor=west,text=white,font=\sffamily\bfseries\LARGE] at ([xshift=5mm,yshift=-0.62cm]frame.north west) {Question 84};
\node[anchor=east,text=white,font=\sffamily\itshape\large] at ([xshift=-5mm,yshift=-0.62cm]frame.north east) {Probabilités};}]
\vspace*{1.75cm}
\begin{minipage}[t][6.9cm][c]{\linewidth}\LARGE\centering
\begin{minipage}{0.95\linewidth}\raggedright\linespread{1.25}\selectfont On a $P(A)=0{,}3$, $P(B)=0{,}6$ et $P(A\cup B)=0{,}8$. Que vaut $P(A\cap B)$~?\end{minipage}
\end{minipage}

\vfill
\begin{tcolorbox}[enhanced,nobeforeafter,width=7.1cm,height=2.1cm,valign=center,colback=cpro!6,colframe=cpro,arc=2.5mm,boxrule=0.9pt,left=3mm]{\Large\bfseries\color{cpro}a)}\quad\begin{minipage}[c]{5.3cm}\raggedright\Large $0{,}5$\end{minipage}\end{tcolorbox}\hfill\begin{tcolorbox}[enhanced,nobeforeafter,width=7.1cm,height=2.1cm,valign=center,colback=cpro!6,colframe=cpro,arc=2.5mm,boxrule=0.9pt,left=3mm]{\Large\bfseries\color{cpro}b)}\quad\begin{minipage}[c]{5.3cm}\raggedright\Large $0{,}18$\end{minipage}\end{tcolorbox}\par\vspace{2.5mm}
\begin{tcolorbox}[enhanced,nobeforeafter,width=7.1cm,height=2.1cm,valign=center,colback=cpro!6,colframe=cpro,arc=2.5mm,boxrule=0.9pt,left=3mm]{\Large\bfseries\color{cpro}c)}\quad\begin{minipage}[c]{5.3cm}\raggedright\Large $0{,}1$\end{minipage}\end{tcolorbox}\hfill\begin{tcolorbox}[enhanced,nobeforeafter,width=7.1cm,height=2.1cm,valign=center,colback=cpro!6,colframe=cpro,arc=2.5mm,boxrule=0.9pt,left=3mm]{\Large\bfseries\color{cpro}d)}\quad\begin{minipage}[c]{5.3cm}\raggedright\Large $0{,}2$\end{minipage}\end{tcolorbox}
\end{tcolorbox}
\newpage
\begin{tcolorbox}[enhanced,width=16cm,height=14cm,colback=white,colframe=cpro,boxrule=2pt,arc=0mm,
left=6mm,right=6mm,top=0mm,bottom=4mm,sharp corners,
overlay={\fill[cpro] (frame.north west) rectangle ([yshift=-1.25cm]frame.north east);
\node[anchor=west,text=white,font=\sffamily\bfseries\LARGE] at ([xshift=5mm,yshift=-0.62cm]frame.north west) {Question 85};
\node[anchor=east,text=white,font=\sffamily\itshape\large] at ([xshift=-5mm,yshift=-0.62cm]frame.north east) {Probabilités};}]
\vspace*{1.75cm}
\begin{minipage}[t][6.9cm][c]{\linewidth}\LARGE\centering
\begin{minipage}{0.95\linewidth}\raggedright\linespread{1.25}\selectfont On a $P(A)=\dfrac{3}{8}$. Que vaut $P(\overline{A})$, la probabilité de l'événement contraire~?\end{minipage}
\end{minipage}

\vfill
\begin{tcolorbox}[enhanced,nobeforeafter,width=7.1cm,height=2.1cm,valign=center,colback=cpro!6,colframe=cpro,arc=2.5mm,boxrule=0.9pt,left=3mm]{\Large\bfseries\color{cpro}a)}\quad\begin{minipage}[c]{5.3cm}\raggedright\Large $\dfrac{3}{8}$\end{minipage}\end{tcolorbox}\hfill\begin{tcolorbox}[enhanced,nobeforeafter,width=7.1cm,height=2.1cm,valign=center,colback=cpro!6,colframe=cpro,arc=2.5mm,boxrule=0.9pt,left=3mm]{\Large\bfseries\color{cpro}b)}\quad\begin{minipage}[c]{5.3cm}\raggedright\Large $\dfrac{5}{16}$\end{minipage}\end{tcolorbox}\par\vspace{2.5mm}
\begin{tcolorbox}[enhanced,nobeforeafter,width=7.1cm,height=2.1cm,valign=center,colback=cpro!6,colframe=cpro,arc=2.5mm,boxrule=0.9pt,left=3mm]{\Large\bfseries\color{cpro}c)}\quad\begin{minipage}[c]{5.3cm}\raggedright\Large $\dfrac{8}{3}$\end{minipage}\end{tcolorbox}\hfill\begin{tcolorbox}[enhanced,nobeforeafter,width=7.1cm,height=2.1cm,valign=center,colback=cpro!6,colframe=cpro,arc=2.5mm,boxrule=0.9pt,left=3mm]{\Large\bfseries\color{cpro}d)}\quad\begin{minipage}[c]{5.3cm}\raggedright\Large $\dfrac{5}{8}$\end{minipage}\end{tcolorbox}
\end{tcolorbox}
\newpage
\begin{tcolorbox}[enhanced,width=16cm,height=14cm,colback=white,colframe=cpro,boxrule=2pt,arc=0mm,
left=6mm,right=6mm,top=0mm,bottom=4mm,sharp corners,
overlay={\fill[cpro] (frame.north west) rectangle ([yshift=-1.25cm]frame.north east);
\node[anchor=west,text=white,font=\sffamily\bfseries\LARGE] at ([xshift=5mm,yshift=-0.62cm]frame.north west) {Question 86};
\node[anchor=east,text=white,font=\sffamily\itshape\large] at ([xshift=-5mm,yshift=-0.62cm]frame.north east) {Probabilités};}]
\vspace*{1.75cm}
\begin{minipage}[t][6.9cm][c]{\linewidth}\LARGE\centering
\begin{minipage}{0.95\linewidth}\raggedright\linespread{1.25}\selectfont On a $P(A)=\dfrac{2}{9}$. Que vaut $P(\overline{A})$, la probabilité de l'événement contraire~?\end{minipage}
\end{minipage}

\vfill
\begin{tcolorbox}[enhanced,nobeforeafter,width=7.1cm,height=2.1cm,valign=center,colback=cpro!6,colframe=cpro,arc=2.5mm,boxrule=0.9pt,left=3mm]{\Large\bfseries\color{cpro}a)}\quad\begin{minipage}[c]{5.3cm}\raggedright\Large $\dfrac{7}{18}$\end{minipage}\end{tcolorbox}\hfill\begin{tcolorbox}[enhanced,nobeforeafter,width=7.1cm,height=2.1cm,valign=center,colback=cpro!6,colframe=cpro,arc=2.5mm,boxrule=0.9pt,left=3mm]{\Large\bfseries\color{cpro}b)}\quad\begin{minipage}[c]{5.3cm}\raggedright\Large $\dfrac{7}{9}$\end{minipage}\end{tcolorbox}\par\vspace{2.5mm}
\begin{tcolorbox}[enhanced,nobeforeafter,width=7.1cm,height=2.1cm,valign=center,colback=cpro!6,colframe=cpro,arc=2.5mm,boxrule=0.9pt,left=3mm]{\Large\bfseries\color{cpro}c)}\quad\begin{minipage}[c]{5.3cm}\raggedright\Large $\dfrac{2}{9}$\end{minipage}\end{tcolorbox}\hfill\begin{tcolorbox}[enhanced,nobeforeafter,width=7.1cm,height=2.1cm,valign=center,colback=cpro!6,colframe=cpro,arc=2.5mm,boxrule=0.9pt,left=3mm]{\Large\bfseries\color{cpro}d)}\quad\begin{minipage}[c]{5.3cm}\raggedright\Large $\dfrac{9}{2}$\end{minipage}\end{tcolorbox}
\end{tcolorbox}
\newpage
\begin{tcolorbox}[enhanced,width=16cm,height=14cm,colback=white,colframe=cpro,boxrule=2pt,arc=0mm,
left=6mm,right=6mm,top=0mm,bottom=4mm,sharp corners,
overlay={\fill[cpro] (frame.north west) rectangle ([yshift=-1.25cm]frame.north east);
\node[anchor=west,text=white,font=\sffamily\bfseries\LARGE] at ([xshift=5mm,yshift=-0.62cm]frame.north west) {Question 87};
\node[anchor=east,text=white,font=\sffamily\itshape\large] at ([xshift=-5mm,yshift=-0.62cm]frame.north east) {Probabilités};}]
\vspace*{1.75cm}
\begin{minipage}[t][6.9cm][c]{\linewidth}\LARGE\centering
\begin{minipage}{0.95\linewidth}\raggedright\linespread{1.25}\selectfont On lance un dé équilibré à six faces. Quelle est la probabilité d'obtenir un nombre pair ou multiple de 3~?\end{minipage}
\end{minipage}

\vfill
\begin{tcolorbox}[enhanced,nobeforeafter,width=7.1cm,height=2.1cm,valign=center,colback=cpro!6,colframe=cpro,arc=2.5mm,boxrule=0.9pt,left=3mm]{\Large\bfseries\color{cpro}a)}\quad\begin{minipage}[c]{5.3cm}\raggedright\Large $\dfrac{2}{3}$\end{minipage}\end{tcolorbox}\hfill\begin{tcolorbox}[enhanced,nobeforeafter,width=7.1cm,height=2.1cm,valign=center,colback=cpro!6,colframe=cpro,arc=2.5mm,boxrule=0.9pt,left=3mm]{\Large\bfseries\color{cpro}b)}\quad\begin{minipage}[c]{5.3cm}\raggedright\Large $\dfrac{1}{3}$\end{minipage}\end{tcolorbox}\par\vspace{2.5mm}
\begin{tcolorbox}[enhanced,nobeforeafter,width=7.1cm,height=2.1cm,valign=center,colback=cpro!6,colframe=cpro,arc=2.5mm,boxrule=0.9pt,left=3mm]{\Large\bfseries\color{cpro}c)}\quad\begin{minipage}[c]{5.3cm}\raggedright\Large $\dfrac{1}{2}$\end{minipage}\end{tcolorbox}\hfill\begin{tcolorbox}[enhanced,nobeforeafter,width=7.1cm,height=2.1cm,valign=center,colback=cpro!6,colframe=cpro,arc=2.5mm,boxrule=0.9pt,left=3mm]{\Large\bfseries\color{cpro}d)}\quad\begin{minipage}[c]{5.3cm}\raggedright\Large $\dfrac{5}{6}$\end{minipage}\end{tcolorbox}
\end{tcolorbox}
\newpage
\begin{tcolorbox}[enhanced,width=16cm,height=14cm,colback=white,colframe=cpro,boxrule=2pt,arc=0mm,
left=6mm,right=6mm,top=0mm,bottom=4mm,sharp corners,
overlay={\fill[cpro] (frame.north west) rectangle ([yshift=-1.25cm]frame.north east);
\node[anchor=west,text=white,font=\sffamily\bfseries\LARGE] at ([xshift=5mm,yshift=-0.62cm]frame.north west) {Question 88};
\node[anchor=east,text=white,font=\sffamily\itshape\large] at ([xshift=-5mm,yshift=-0.62cm]frame.north east) {Probabilités};}]
\vspace*{1.75cm}
\begin{minipage}[t][6.9cm][c]{\linewidth}\LARGE\centering
\begin{minipage}{0.95\linewidth}\raggedright\linespread{1.25}\selectfont On lance un dé équilibré à six faces. Quelle est la probabilité d'obtenir un nombre supérieur à 4 et pair~?\end{minipage}
\end{minipage}

\vfill
\begin{tcolorbox}[enhanced,nobeforeafter,width=7.1cm,height=2.1cm,valign=center,colback=cpro!6,colframe=cpro,arc=2.5mm,boxrule=0.9pt,left=3mm]{\Large\bfseries\color{cpro}a)}\quad\begin{minipage}[c]{5.3cm}\raggedright\Large $\dfrac{1}{5}$\end{minipage}\end{tcolorbox}\hfill\begin{tcolorbox}[enhanced,nobeforeafter,width=7.1cm,height=2.1cm,valign=center,colback=cpro!6,colframe=cpro,arc=2.5mm,boxrule=0.9pt,left=3mm]{\Large\bfseries\color{cpro}b)}\quad\begin{minipage}[c]{5.3cm}\raggedright\Large $\dfrac{1}{3}$\end{minipage}\end{tcolorbox}\par\vspace{2.5mm}
\begin{tcolorbox}[enhanced,nobeforeafter,width=7.1cm,height=2.1cm,valign=center,colback=cpro!6,colframe=cpro,arc=2.5mm,boxrule=0.9pt,left=3mm]{\Large\bfseries\color{cpro}c)}\quad\begin{minipage}[c]{5.3cm}\raggedright\Large $\dfrac{1}{6}$\end{minipage}\end{tcolorbox}\hfill\begin{tcolorbox}[enhanced,nobeforeafter,width=7.1cm,height=2.1cm,valign=center,colback=cpro!6,colframe=cpro,arc=2.5mm,boxrule=0.9pt,left=3mm]{\Large\bfseries\color{cpro}d)}\quad\begin{minipage}[c]{5.3cm}\raggedright\Large $\dfrac{5}{6}$\end{minipage}\end{tcolorbox}
\end{tcolorbox}
\newpage
\begin{tcolorbox}[enhanced,width=16cm,height=14cm,colback=white,colframe=cpro,boxrule=2pt,arc=0mm,
left=6mm,right=6mm,top=0mm,bottom=4mm,sharp corners,
overlay={\fill[cpro] (frame.north west) rectangle ([yshift=-1.25cm]frame.north east);
\node[anchor=west,text=white,font=\sffamily\bfseries\LARGE] at ([xshift=5mm,yshift=-0.62cm]frame.north west) {Question 89};
\node[anchor=east,text=white,font=\sffamily\itshape\large] at ([xshift=-5mm,yshift=-0.62cm]frame.north east) {Probabilités};}]
\vspace*{1.75cm}
\begin{minipage}[t][6.9cm][c]{\linewidth}\LARGE\centering
\begin{minipage}{0.95\linewidth}\raggedright\linespread{1.25}\selectfont On lance $2$ pièces équilibrées. Quelle est la probabilité de n'obtenir que des PILE~?\end{minipage}
\end{minipage}

\vfill
\begin{tcolorbox}[enhanced,nobeforeafter,width=7.1cm,height=2.1cm,valign=center,colback=cpro!6,colframe=cpro,arc=2.5mm,boxrule=0.9pt,left=3mm]{\Large\bfseries\color{cpro}a)}\quad\begin{minipage}[c]{5.3cm}\raggedright\Large $\dfrac{1}{2}$\end{minipage}\end{tcolorbox}\hfill\begin{tcolorbox}[enhanced,nobeforeafter,width=7.1cm,height=2.1cm,valign=center,colback=cpro!6,colframe=cpro,arc=2.5mm,boxrule=0.9pt,left=3mm]{\Large\bfseries\color{cpro}b)}\quad\begin{minipage}[c]{5.3cm}\raggedright\Large $\dfrac{1}{4}$\end{minipage}\end{tcolorbox}\par\vspace{2.5mm}
\begin{tcolorbox}[enhanced,nobeforeafter,width=7.1cm,height=2.1cm,valign=center,colback=cpro!6,colframe=cpro,arc=2.5mm,boxrule=0.9pt,left=3mm]{\Large\bfseries\color{cpro}c)}\quad\begin{minipage}[c]{5.3cm}\raggedright\Large $\dfrac{3}{4}$\end{minipage}\end{tcolorbox}\hfill\begin{tcolorbox}[enhanced,nobeforeafter,width=7.1cm,height=2.1cm,valign=center,colback=cpro!6,colframe=cpro,arc=2.5mm,boxrule=0.9pt,left=3mm]{\Large\bfseries\color{cpro}d)}\quad\begin{minipage}[c]{5.3cm}\raggedright\Large $\dfrac{1}{8}$\end{minipage}\end{tcolorbox}
\end{tcolorbox}
\newpage
\begin{tcolorbox}[enhanced,width=16cm,height=14cm,colback=white,colframe=cpro,boxrule=2pt,arc=0mm,
left=6mm,right=6mm,top=0mm,bottom=4mm,sharp corners,
overlay={\fill[cpro] (frame.north west) rectangle ([yshift=-1.25cm]frame.north east);
\node[anchor=west,text=white,font=\sffamily\bfseries\LARGE] at ([xshift=5mm,yshift=-0.62cm]frame.north west) {Question 90};
\node[anchor=east,text=white,font=\sffamily\itshape\large] at ([xshift=-5mm,yshift=-0.62cm]frame.north east) {Probabilités};}]
\vspace*{1.75cm}
\begin{minipage}[t][6.9cm][c]{\linewidth}\LARGE\centering
\begin{minipage}{0.95\linewidth}\raggedright\linespread{1.25}\selectfont On lance $3$ pièces équilibrées. Quelle est la probabilité de n'obtenir que des PILE~?\end{minipage}
\end{minipage}

\vfill
\begin{tcolorbox}[enhanced,nobeforeafter,width=7.1cm,height=2.1cm,valign=center,colback=cpro!6,colframe=cpro,arc=2.5mm,boxrule=0.9pt,left=3mm]{\Large\bfseries\color{cpro}a)}\quad\begin{minipage}[c]{5.3cm}\raggedright\Large $\dfrac{7}{8}$\end{minipage}\end{tcolorbox}\hfill\begin{tcolorbox}[enhanced,nobeforeafter,width=7.1cm,height=2.1cm,valign=center,colback=cpro!6,colframe=cpro,arc=2.5mm,boxrule=0.9pt,left=3mm]{\Large\bfseries\color{cpro}b)}\quad\begin{minipage}[c]{5.3cm}\raggedright\Large $\dfrac{1}{8}$\end{minipage}\end{tcolorbox}\par\vspace{2.5mm}
\begin{tcolorbox}[enhanced,nobeforeafter,width=7.1cm,height=2.1cm,valign=center,colback=cpro!6,colframe=cpro,arc=2.5mm,boxrule=0.9pt,left=3mm]{\Large\bfseries\color{cpro}c)}\quad\begin{minipage}[c]{5.3cm}\raggedright\Large $\dfrac{1}{2}$\end{minipage}\end{tcolorbox}\hfill\begin{tcolorbox}[enhanced,nobeforeafter,width=7.1cm,height=2.1cm,valign=center,colback=cpro!6,colframe=cpro,arc=2.5mm,boxrule=0.9pt,left=3mm]{\Large\bfseries\color{cpro}d)}\quad\begin{minipage}[c]{5.3cm}\raggedright\Large $\dfrac{1}{6}$\end{minipage}\end{tcolorbox}
\end{tcolorbox}
\newpage
\begin{tcolorbox}[enhanced,width=16cm,height=14cm,colback=white,colframe=cpro,boxrule=2pt,arc=0mm,
left=6mm,right=6mm,top=0mm,bottom=4mm,sharp corners,
overlay={\fill[cpro] (frame.north west) rectangle ([yshift=-1.25cm]frame.north east);
\node[anchor=west,text=white,font=\sffamily\bfseries\LARGE] at ([xshift=5mm,yshift=-0.62cm]frame.north west) {Question 91};
\node[anchor=east,text=white,font=\sffamily\itshape\large] at ([xshift=-5mm,yshift=-0.62cm]frame.north east) {Probabilités};}]
\vspace*{1.75cm}
\begin{minipage}[t][6.9cm][c]{\linewidth}\Large\centering
\begin{minipage}{0.95\linewidth}\raggedright\linespread{1.25}\selectfont Voici la loi de probabilité d'une expérience aléatoire. Quelle est la probabilité manquante~? $\begin{array}{|l|c|c|c|c|}\hline \text{Issue} & 1 & 2 & 3 & 4\\\hline \text{Probabilité} & 0{,}2 & 0{,}35 &~? & 0{,}1\\\hline\end{array}$\end{minipage}
\end{minipage}

\vfill
\begin{tcolorbox}[enhanced,nobeforeafter,width=7.1cm,height=2.1cm,valign=center,colback=cpro!6,colframe=cpro,arc=2.5mm,boxrule=0.9pt,left=3mm]{\Large\bfseries\color{cpro}a)}\quad\begin{minipage}[c]{5.3cm}\raggedright\Large $0{,}35$\end{minipage}\end{tcolorbox}\hfill\begin{tcolorbox}[enhanced,nobeforeafter,width=7.1cm,height=2.1cm,valign=center,colback=cpro!6,colframe=cpro,arc=2.5mm,boxrule=0.9pt,left=3mm]{\Large\bfseries\color{cpro}b)}\quad\begin{minipage}[c]{5.3cm}\raggedright\Large $0{,}25$\end{minipage}\end{tcolorbox}\par\vspace{2.5mm}
\begin{tcolorbox}[enhanced,nobeforeafter,width=7.1cm,height=2.1cm,valign=center,colback=cpro!6,colframe=cpro,arc=2.5mm,boxrule=0.9pt,left=3mm]{\Large\bfseries\color{cpro}c)}\quad\begin{minipage}[c]{5.3cm}\raggedright\Large $0{,}45$\end{minipage}\end{tcolorbox}\hfill\begin{tcolorbox}[enhanced,nobeforeafter,width=7.1cm,height=2.1cm,valign=center,colback=cpro!6,colframe=cpro,arc=2.5mm,boxrule=0.9pt,left=3mm]{\Large\bfseries\color{cpro}d)}\quad\begin{minipage}[c]{5.3cm}\raggedright\Large $0{,}65$\end{minipage}\end{tcolorbox}
\end{tcolorbox}
\newpage
\begin{tcolorbox}[enhanced,width=16cm,height=14cm,colback=white,colframe=cpro,boxrule=2pt,arc=0mm,
left=6mm,right=6mm,top=0mm,bottom=4mm,sharp corners,
overlay={\fill[cpro] (frame.north west) rectangle ([yshift=-1.25cm]frame.north east);
\node[anchor=west,text=white,font=\sffamily\bfseries\LARGE] at ([xshift=5mm,yshift=-0.62cm]frame.north west) {Question 92};
\node[anchor=east,text=white,font=\sffamily\itshape\large] at ([xshift=-5mm,yshift=-0.62cm]frame.north east) {Probabilités};}]
\vspace*{1.75cm}
\begin{minipage}[t][6.9cm][c]{\linewidth}\Large\centering
\begin{minipage}{0.95\linewidth}\raggedright\linespread{1.25}\selectfont Voici la loi de probabilité d'une expérience aléatoire. Quelle est la probabilité manquante~? $\begin{array}{|l|c|c|c|c|}\hline \text{Issue} & 1 & 2 & 3 & 4\\\hline \text{Probabilité} & 0{,}15 & 0{,}4 &~? & 0{,}25\\\hline\end{array}$\end{minipage}
\end{minipage}

\vfill
\begin{tcolorbox}[enhanced,nobeforeafter,width=7.1cm,height=2.1cm,valign=center,colback=cpro!6,colframe=cpro,arc=2.5mm,boxrule=0.9pt,left=3mm]{\Large\bfseries\color{cpro}a)}\quad\begin{minipage}[c]{5.3cm}\raggedright\Large $0{,}45$\end{minipage}\end{tcolorbox}\hfill\begin{tcolorbox}[enhanced,nobeforeafter,width=7.1cm,height=2.1cm,valign=center,colback=cpro!6,colframe=cpro,arc=2.5mm,boxrule=0.9pt,left=3mm]{\Large\bfseries\color{cpro}b)}\quad\begin{minipage}[c]{5.3cm}\raggedright\Large $0{,}2$\end{minipage}\end{tcolorbox}\par\vspace{2.5mm}
\begin{tcolorbox}[enhanced,nobeforeafter,width=7.1cm,height=2.1cm,valign=center,colback=cpro!6,colframe=cpro,arc=2.5mm,boxrule=0.9pt,left=3mm]{\Large\bfseries\color{cpro}c)}\quad\begin{minipage}[c]{5.3cm}\raggedright\Large $0{,}8$\end{minipage}\end{tcolorbox}\hfill\begin{tcolorbox}[enhanced,nobeforeafter,width=7.1cm,height=2.1cm,valign=center,colback=cpro!6,colframe=cpro,arc=2.5mm,boxrule=0.9pt,left=3mm]{\Large\bfseries\color{cpro}d)}\quad\begin{minipage}[c]{5.3cm}\raggedright\Large $0{,}25$\end{minipage}\end{tcolorbox}
\end{tcolorbox}
\newpage
\begin{tcolorbox}[enhanced,width=16cm,height=14cm,colback=white,colframe=cpro,boxrule=2pt,arc=0mm,
left=6mm,right=6mm,top=0mm,bottom=4mm,sharp corners,
overlay={\fill[cpro] (frame.north west) rectangle ([yshift=-1.25cm]frame.north east);
\node[anchor=west,text=white,font=\sffamily\bfseries\LARGE] at ([xshift=5mm,yshift=-0.62cm]frame.north west) {Question 93};
\node[anchor=east,text=white,font=\sffamily\itshape\large] at ([xshift=-5mm,yshift=-0.62cm]frame.north east) {Probabilités};}]
\vspace*{1.75cm}
\begin{minipage}[t][6.9cm][c]{\linewidth}\Large\centering
\begin{minipage}{0.95\linewidth}\raggedright\linespread{1.25}\selectfont Une urne contient $5$ boules dont $3$ rouges. On tire une boule, on la remet, puis on en tire une deuxième. Quelle est la probabilité de tirer deux boules rouges~?\end{minipage}
\end{minipage}

\vfill
\begin{tcolorbox}[enhanced,nobeforeafter,width=7.1cm,height=2.1cm,valign=center,colback=cpro!6,colframe=cpro,arc=2.5mm,boxrule=0.9pt,left=3mm]{\Large\bfseries\color{cpro}a)}\quad\begin{minipage}[c]{5.3cm}\raggedright\Large $1$\end{minipage}\end{tcolorbox}\hfill\begin{tcolorbox}[enhanced,nobeforeafter,width=7.1cm,height=2.1cm,valign=center,colback=cpro!6,colframe=cpro,arc=2.5mm,boxrule=0.9pt,left=3mm]{\Large\bfseries\color{cpro}b)}\quad\begin{minipage}[c]{5.3cm}\raggedright\Large $\dfrac{9}{25}$\end{minipage}\end{tcolorbox}\par\vspace{2.5mm}
\begin{tcolorbox}[enhanced,nobeforeafter,width=7.1cm,height=2.1cm,valign=center,colback=cpro!6,colframe=cpro,arc=2.5mm,boxrule=0.9pt,left=3mm]{\Large\bfseries\color{cpro}c)}\quad\begin{minipage}[c]{5.3cm}\raggedright\Large $\dfrac{3}{10}$\end{minipage}\end{tcolorbox}\hfill\begin{tcolorbox}[enhanced,nobeforeafter,width=7.1cm,height=2.1cm,valign=center,colback=cpro!6,colframe=cpro,arc=2.5mm,boxrule=0.9pt,left=3mm]{\Large\bfseries\color{cpro}d)}\quad\begin{minipage}[c]{5.3cm}\raggedright\Large $\dfrac{3}{5}$\end{minipage}\end{tcolorbox}
\end{tcolorbox}
\newpage
\begin{tcolorbox}[enhanced,width=16cm,height=14cm,colback=white,colframe=cpro,boxrule=2pt,arc=0mm,
left=6mm,right=6mm,top=0mm,bottom=4mm,sharp corners,
overlay={\fill[cpro] (frame.north west) rectangle ([yshift=-1.25cm]frame.north east);
\node[anchor=west,text=white,font=\sffamily\bfseries\LARGE] at ([xshift=5mm,yshift=-0.62cm]frame.north west) {Question 94};
\node[anchor=east,text=white,font=\sffamily\itshape\large] at ([xshift=-5mm,yshift=-0.62cm]frame.north east) {Probabilités};}]
\vspace*{1.75cm}
\begin{minipage}[t][6.9cm][c]{\linewidth}\Large\centering
\begin{minipage}{0.95\linewidth}\raggedright\linespread{1.25}\selectfont Une urne contient $6$ boules dont $2$ rouges. On tire une boule, on la remet, puis on en tire une deuxième. Quelle est la probabilité de tirer deux boules rouges~?\end{minipage}
\end{minipage}

\vfill
\begin{tcolorbox}[enhanced,nobeforeafter,width=7.1cm,height=2.1cm,valign=center,colback=cpro!6,colframe=cpro,arc=2.5mm,boxrule=0.9pt,left=3mm]{\Large\bfseries\color{cpro}a)}\quad\begin{minipage}[c]{5.3cm}\raggedright\Large $\dfrac{1}{9}$\end{minipage}\end{tcolorbox}\hfill\begin{tcolorbox}[enhanced,nobeforeafter,width=7.1cm,height=2.1cm,valign=center,colback=cpro!6,colframe=cpro,arc=2.5mm,boxrule=0.9pt,left=3mm]{\Large\bfseries\color{cpro}b)}\quad\begin{minipage}[c]{5.3cm}\raggedright\Large $\dfrac{1}{15}$\end{minipage}\end{tcolorbox}\par\vspace{2.5mm}
\begin{tcolorbox}[enhanced,nobeforeafter,width=7.1cm,height=2.1cm,valign=center,colback=cpro!6,colframe=cpro,arc=2.5mm,boxrule=0.9pt,left=3mm]{\Large\bfseries\color{cpro}c)}\quad\begin{minipage}[c]{5.3cm}\raggedright\Large $\dfrac{2}{3}$\end{minipage}\end{tcolorbox}\hfill\begin{tcolorbox}[enhanced,nobeforeafter,width=7.1cm,height=2.1cm,valign=center,colback=cpro!6,colframe=cpro,arc=2.5mm,boxrule=0.9pt,left=3mm]{\Large\bfseries\color{cpro}d)}\quad\begin{minipage}[c]{5.3cm}\raggedright\Large $\dfrac{1}{3}$\end{minipage}\end{tcolorbox}
\end{tcolorbox}
\newpage
\begin{tcolorbox}[enhanced,width=16cm,height=14cm,colback=white,colframe=cpro,boxrule=2pt,arc=0mm,
left=6mm,right=6mm,top=0mm,bottom=4mm,sharp corners,
overlay={\fill[cpro] (frame.north west) rectangle ([yshift=-1.25cm]frame.north east);
\node[anchor=west,text=white,font=\sffamily\bfseries\LARGE] at ([xshift=5mm,yshift=-0.62cm]frame.north west) {Question 95};
\node[anchor=east,text=white,font=\sffamily\itshape\large] at ([xshift=-5mm,yshift=-0.62cm]frame.north east) {Probabilités};}]
\vspace*{1.75cm}
\begin{minipage}[t][6.9cm][c]{\linewidth}\LARGE\centering
\begin{minipage}{0.95\linewidth}\raggedright\linespread{1.25}\selectfont $A$ et $B$ sont deux événements incompatibles, avec $P(A)=0{,}3$ et $P(B)=0{,}45$. Que vaut $P(A\cup B)$~?\end{minipage}
\end{minipage}

\vfill
\begin{tcolorbox}[enhanced,nobeforeafter,width=7.1cm,height=2.1cm,valign=center,colback=cpro!6,colframe=cpro,arc=2.5mm,boxrule=0.9pt,left=3mm]{\Large\bfseries\color{cpro}a)}\quad\begin{minipage}[c]{5.3cm}\raggedright\Large $0{,}615$\end{minipage}\end{tcolorbox}\hfill\begin{tcolorbox}[enhanced,nobeforeafter,width=7.1cm,height=2.1cm,valign=center,colback=cpro!6,colframe=cpro,arc=2.5mm,boxrule=0.9pt,left=3mm]{\Large\bfseries\color{cpro}b)}\quad\begin{minipage}[c]{5.3cm}\raggedright\Large $0{,}135$\end{minipage}\end{tcolorbox}\par\vspace{2.5mm}
\begin{tcolorbox}[enhanced,nobeforeafter,width=7.1cm,height=2.1cm,valign=center,colback=cpro!6,colframe=cpro,arc=2.5mm,boxrule=0.9pt,left=3mm]{\Large\bfseries\color{cpro}c)}\quad\begin{minipage}[c]{5.3cm}\raggedright\Large $0{,}15$\end{minipage}\end{tcolorbox}\hfill\begin{tcolorbox}[enhanced,nobeforeafter,width=7.1cm,height=2.1cm,valign=center,colback=cpro!6,colframe=cpro,arc=2.5mm,boxrule=0.9pt,left=3mm]{\Large\bfseries\color{cpro}d)}\quad\begin{minipage}[c]{5.3cm}\raggedright\Large $0{,}75$\end{minipage}\end{tcolorbox}
\end{tcolorbox}
\newpage
\begin{tcolorbox}[enhanced,width=16cm,height=14cm,colback=white,colframe=cpro,boxrule=2pt,arc=0mm,
left=6mm,right=6mm,top=0mm,bottom=4mm,sharp corners,
overlay={\fill[cpro] (frame.north west) rectangle ([yshift=-1.25cm]frame.north east);
\node[anchor=west,text=white,font=\sffamily\bfseries\LARGE] at ([xshift=5mm,yshift=-0.62cm]frame.north west) {Question 96};
\node[anchor=east,text=white,font=\sffamily\itshape\large] at ([xshift=-5mm,yshift=-0.62cm]frame.north east) {Probabilités};}]
\vspace*{1.75cm}
\begin{minipage}[t][6.9cm][c]{\linewidth}\LARGE\centering
\begin{minipage}{0.95\linewidth}\raggedright\linespread{1.25}\selectfont $A$ et $B$ sont deux événements incompatibles, avec $P(A)=\dfrac{1}{4}$ et $P(B)=\dfrac{1}{3}$. Que vaut $P(A\cup B)$~?\end{minipage}
\end{minipage}

\vfill
\begin{tcolorbox}[enhanced,nobeforeafter,width=7.1cm,height=2.1cm,valign=center,colback=cpro!6,colframe=cpro,arc=2.5mm,boxrule=0.9pt,left=3mm]{\Large\bfseries\color{cpro}a)}\quad\begin{minipage}[c]{5.3cm}\raggedright\Large $\dfrac{1}{2}$\end{minipage}\end{tcolorbox}\hfill\begin{tcolorbox}[enhanced,nobeforeafter,width=7.1cm,height=2.1cm,valign=center,colback=cpro!6,colframe=cpro,arc=2.5mm,boxrule=0.9pt,left=3mm]{\Large\bfseries\color{cpro}b)}\quad\begin{minipage}[c]{5.3cm}\raggedright\Large $\dfrac{7}{12}$\end{minipage}\end{tcolorbox}\par\vspace{2.5mm}
\begin{tcolorbox}[enhanced,nobeforeafter,width=7.1cm,height=2.1cm,valign=center,colback=cpro!6,colframe=cpro,arc=2.5mm,boxrule=0.9pt,left=3mm]{\Large\bfseries\color{cpro}c)}\quad\begin{minipage}[c]{5.3cm}\raggedright\Large $\dfrac{5}{12}$\end{minipage}\end{tcolorbox}\hfill\begin{tcolorbox}[enhanced,nobeforeafter,width=7.1cm,height=2.1cm,valign=center,colback=cpro!6,colframe=cpro,arc=2.5mm,boxrule=0.9pt,left=3mm]{\Large\bfseries\color{cpro}d)}\quad\begin{minipage}[c]{5.3cm}\raggedright\Large $\dfrac{1}{12}$\end{minipage}\end{tcolorbox}
\end{tcolorbox}
\newpage
\begin{tcolorbox}[enhanced,width=16cm,height=14cm,colback=white,colframe=cnom,boxrule=2pt,arc=0mm,
left=6mm,right=6mm,top=0mm,bottom=4mm,sharp corners,
overlay={\fill[cnom] (frame.north west) rectangle ([yshift=-1.25cm]frame.north east);
\node[anchor=west,text=white,font=\sffamily\bfseries\LARGE] at ([xshift=5mm,yshift=-0.62cm]frame.north west) {Question 1 B};
\node[anchor=east,text=white,font=\sffamily\itshape\large] at ([xshift=-5mm,yshift=-0.62cm]frame.north east) {Nombres et calculs};}]
\vspace*{1.75cm}
\begin{minipage}[t][6.9cm][c]{\linewidth}\LARGE\centering
\begin{minipage}{0.95\linewidth}\raggedright\linespread{1.25}\selectfont Écrire $\sqrt{72}$ sous la forme $a\sqrt{b}$, avec $b$ le plus petit possible.\end{minipage}
\end{minipage}

\vfill
\begin{tcolorbox}[enhanced,nobeforeafter,width=7.1cm,height=2.1cm,valign=center,colback=cnom!6,colframe=cnom,arc=2.5mm,boxrule=0.9pt,left=3mm]{\Large\bfseries\color{cnom}a)}\quad\begin{minipage}[c]{5.3cm}\raggedright\Large $2\sqrt{6}$\end{minipage}\end{tcolorbox}\hfill\begin{tcolorbox}[enhanced,nobeforeafter,width=7.1cm,height=2.1cm,valign=center,colback=cnom!6,colframe=cnom,arc=2.5mm,boxrule=0.9pt,left=3mm]{\Large\bfseries\color{cnom}b)}\quad\begin{minipage}[c]{5.3cm}\raggedright\Large $36\sqrt{2}$\end{minipage}\end{tcolorbox}\par\vspace{2.5mm}
\begin{tcolorbox}[enhanced,nobeforeafter,width=7.1cm,height=2.1cm,valign=center,colback=cnom!6,colframe=cnom,arc=2.5mm,boxrule=0.9pt,left=3mm]{\Large\bfseries\color{cnom}c)}\quad\begin{minipage}[c]{5.3cm}\raggedright\Large $6\sqrt{4}$\end{minipage}\end{tcolorbox}\hfill\begin{tcolorbox}[enhanced,nobeforeafter,width=7.1cm,height=2.1cm,valign=center,colback=cnom!6,colframe=cnom,arc=2.5mm,boxrule=0.9pt,left=3mm]{\Large\bfseries\color{cnom}d)}\quad\begin{minipage}[c]{5.3cm}\raggedright\Large $6\sqrt{2}$\end{minipage}\end{tcolorbox}
\end{tcolorbox}
\newpage
\begin{tcolorbox}[enhanced,width=16cm,height=14cm,colback=white,colframe=cnom,boxrule=2pt,arc=0mm,
left=6mm,right=6mm,top=0mm,bottom=4mm,sharp corners,
overlay={\fill[cnom] (frame.north west) rectangle ([yshift=-1.25cm]frame.north east);
\node[anchor=west,text=white,font=\sffamily\bfseries\LARGE] at ([xshift=5mm,yshift=-0.62cm]frame.north west) {Question 2 B};
\node[anchor=east,text=white,font=\sffamily\itshape\large] at ([xshift=-5mm,yshift=-0.62cm]frame.north east) {Nombres et calculs};}]
\vspace*{1.75cm}
\begin{minipage}[t][6.9cm][c]{\linewidth}\LARGE\centering
\begin{minipage}{0.95\linewidth}\raggedright\linespread{1.25}\selectfont Écrire $\sqrt{18}$ sous la forme $a\sqrt{b}$, avec $b$ le plus petit possible.\end{minipage}
\end{minipage}

\vfill
\begin{tcolorbox}[enhanced,nobeforeafter,width=7.1cm,height=2.1cm,valign=center,colback=cnom!6,colframe=cnom,arc=2.5mm,boxrule=0.9pt,left=3mm]{\Large\bfseries\color{cnom}a)}\quad\begin{minipage}[c]{5.3cm}\raggedright\Large $3\sqrt{2}$\end{minipage}\end{tcolorbox}\hfill\begin{tcolorbox}[enhanced,nobeforeafter,width=7.1cm,height=2.1cm,valign=center,colback=cnom!6,colframe=cnom,arc=2.5mm,boxrule=0.9pt,left=3mm]{\Large\bfseries\color{cnom}b)}\quad\begin{minipage}[c]{5.3cm}\raggedright\Large $9\sqrt{2}$\end{minipage}\end{tcolorbox}\par\vspace{2.5mm}
\begin{tcolorbox}[enhanced,nobeforeafter,width=7.1cm,height=2.1cm,valign=center,colback=cnom!6,colframe=cnom,arc=2.5mm,boxrule=0.9pt,left=3mm]{\Large\bfseries\color{cnom}c)}\quad\begin{minipage}[c]{5.3cm}\raggedright\Large $2\sqrt{3}$\end{minipage}\end{tcolorbox}\hfill\begin{tcolorbox}[enhanced,nobeforeafter,width=7.1cm,height=2.1cm,valign=center,colback=cnom!6,colframe=cnom,arc=2.5mm,boxrule=0.9pt,left=3mm]{\Large\bfseries\color{cnom}d)}\quad\begin{minipage}[c]{5.3cm}\raggedright\Large $3\sqrt{4}$\end{minipage}\end{tcolorbox}
\end{tcolorbox}
\newpage
\begin{tcolorbox}[enhanced,width=16cm,height=14cm,colback=white,colframe=cnom,boxrule=2pt,arc=0mm,
left=6mm,right=6mm,top=0mm,bottom=4mm,sharp corners,
overlay={\fill[cnom] (frame.north west) rectangle ([yshift=-1.25cm]frame.north east);
\node[anchor=west,text=white,font=\sffamily\bfseries\LARGE] at ([xshift=5mm,yshift=-0.62cm]frame.north west) {Question 3 B};
\node[anchor=east,text=white,font=\sffamily\itshape\large] at ([xshift=-5mm,yshift=-0.62cm]frame.north east) {Nombres et calculs};}]
\vspace*{1.75cm}
\begin{minipage}[t][6.9cm][c]{\linewidth}\LARGE\centering
\begin{minipage}{0.95\linewidth}\raggedright\linespread{1.25}\selectfont Que vaut $\sqrt{20}+\sqrt{45}$~?\end{minipage}
\end{minipage}

\vfill
\begin{tcolorbox}[enhanced,nobeforeafter,width=7.1cm,height=2.1cm,valign=center,colback=cnom!6,colframe=cnom,arc=2.5mm,boxrule=0.9pt,left=3mm]{\Large\bfseries\color{cnom}a)}\quad\begin{minipage}[c]{5.3cm}\raggedright\Large $\sqrt{65}$\end{minipage}\end{tcolorbox}\hfill\begin{tcolorbox}[enhanced,nobeforeafter,width=7.1cm,height=2.1cm,valign=center,colback=cnom!6,colframe=cnom,arc=2.5mm,boxrule=0.9pt,left=3mm]{\Large\bfseries\color{cnom}b)}\quad\begin{minipage}[c]{5.3cm}\raggedright\Large $6\sqrt{5}$\end{minipage}\end{tcolorbox}\par\vspace{2.5mm}
\begin{tcolorbox}[enhanced,nobeforeafter,width=7.1cm,height=2.1cm,valign=center,colback=cnom!6,colframe=cnom,arc=2.5mm,boxrule=0.9pt,left=3mm]{\Large\bfseries\color{cnom}c)}\quad\begin{minipage}[c]{5.3cm}\raggedright\Large $5\sqrt{5}$\end{minipage}\end{tcolorbox}\hfill\begin{tcolorbox}[enhanced,nobeforeafter,width=7.1cm,height=2.1cm,valign=center,colback=cnom!6,colframe=cnom,arc=2.5mm,boxrule=0.9pt,left=3mm]{\Large\bfseries\color{cnom}d)}\quad\begin{minipage}[c]{5.3cm}\raggedright\Large $5\sqrt{10}$\end{minipage}\end{tcolorbox}
\end{tcolorbox}
\newpage
\begin{tcolorbox}[enhanced,width=16cm,height=14cm,colback=white,colframe=cnom,boxrule=2pt,arc=0mm,
left=6mm,right=6mm,top=0mm,bottom=4mm,sharp corners,
overlay={\fill[cnom] (frame.north west) rectangle ([yshift=-1.25cm]frame.north east);
\node[anchor=west,text=white,font=\sffamily\bfseries\LARGE] at ([xshift=5mm,yshift=-0.62cm]frame.north west) {Question 4 B};
\node[anchor=east,text=white,font=\sffamily\itshape\large] at ([xshift=-5mm,yshift=-0.62cm]frame.north east) {Nombres et calculs};}]
\vspace*{1.75cm}
\begin{minipage}[t][6.9cm][c]{\linewidth}\LARGE\centering
\begin{minipage}{0.95\linewidth}\raggedright\linespread{1.25}\selectfont Que vaut $\sqrt{50}+\sqrt{32}$~?\end{minipage}
\end{minipage}

\vfill
\begin{tcolorbox}[enhanced,nobeforeafter,width=7.1cm,height=2.1cm,valign=center,colback=cnom!6,colframe=cnom,arc=2.5mm,boxrule=0.9pt,left=3mm]{\Large\bfseries\color{cnom}a)}\quad\begin{minipage}[c]{5.3cm}\raggedright\Large $9\sqrt{4}$\end{minipage}\end{tcolorbox}\hfill\begin{tcolorbox}[enhanced,nobeforeafter,width=7.1cm,height=2.1cm,valign=center,colback=cnom!6,colframe=cnom,arc=2.5mm,boxrule=0.9pt,left=3mm]{\Large\bfseries\color{cnom}b)}\quad\begin{minipage}[c]{5.3cm}\raggedright\Large $9\sqrt{2}$\end{minipage}\end{tcolorbox}\par\vspace{2.5mm}
\begin{tcolorbox}[enhanced,nobeforeafter,width=7.1cm,height=2.1cm,valign=center,colback=cnom!6,colframe=cnom,arc=2.5mm,boxrule=0.9pt,left=3mm]{\Large\bfseries\color{cnom}c)}\quad\begin{minipage}[c]{5.3cm}\raggedright\Large $20\sqrt{2}$\end{minipage}\end{tcolorbox}\hfill\begin{tcolorbox}[enhanced,nobeforeafter,width=7.1cm,height=2.1cm,valign=center,colback=cnom!6,colframe=cnom,arc=2.5mm,boxrule=0.9pt,left=3mm]{\Large\bfseries\color{cnom}d)}\quad\begin{minipage}[c]{5.3cm}\raggedright\Large $\sqrt{82}$\end{minipage}\end{tcolorbox}
\end{tcolorbox}
\newpage
\begin{tcolorbox}[enhanced,width=16cm,height=14cm,colback=white,colframe=cnom,boxrule=2pt,arc=0mm,
left=6mm,right=6mm,top=0mm,bottom=4mm,sharp corners,
overlay={\fill[cnom] (frame.north west) rectangle ([yshift=-1.25cm]frame.north east);
\node[anchor=west,text=white,font=\sffamily\bfseries\LARGE] at ([xshift=5mm,yshift=-0.62cm]frame.north west) {Question 5 B};
\node[anchor=east,text=white,font=\sffamily\itshape\large] at ([xshift=-5mm,yshift=-0.62cm]frame.north east) {Nombres et calculs};}]
\vspace*{1.75cm}
\begin{minipage}[t][6.9cm][c]{\linewidth}\LARGE\centering
\begin{minipage}{0.95\linewidth}\raggedright\linespread{1.25}\selectfont Développer $(2x+1)^2$.\end{minipage}
\end{minipage}

\vfill
\begin{tcolorbox}[enhanced,nobeforeafter,width=7.1cm,height=2.1cm,valign=center,colback=cnom!6,colframe=cnom,arc=2.5mm,boxrule=0.9pt,left=3mm]{\Large\bfseries\color{cnom}a)}\quad\begin{minipage}[c]{5.3cm}\raggedright\Large $4x^2+4x+1$\end{minipage}\end{tcolorbox}\hfill\begin{tcolorbox}[enhanced,nobeforeafter,width=7.1cm,height=2.1cm,valign=center,colback=cnom!6,colframe=cnom,arc=2.5mm,boxrule=0.9pt,left=3mm]{\Large\bfseries\color{cnom}b)}\quad\begin{minipage}[c]{5.3cm}\raggedright\Large $4x^2+2x+1$\end{minipage}\end{tcolorbox}\par\vspace{2.5mm}
\begin{tcolorbox}[enhanced,nobeforeafter,width=7.1cm,height=2.1cm,valign=center,colback=cnom!6,colframe=cnom,arc=2.5mm,boxrule=0.9pt,left=3mm]{\Large\bfseries\color{cnom}c)}\quad\begin{minipage}[c]{5.3cm}\raggedright\Large $2x^2+4x+1$\end{minipage}\end{tcolorbox}\hfill\begin{tcolorbox}[enhanced,nobeforeafter,width=7.1cm,height=2.1cm,valign=center,colback=cnom!6,colframe=cnom,arc=2.5mm,boxrule=0.9pt,left=3mm]{\Large\bfseries\color{cnom}d)}\quad\begin{minipage}[c]{5.3cm}\raggedright\Large $4x^2+1$\end{minipage}\end{tcolorbox}
\end{tcolorbox}
\newpage
\begin{tcolorbox}[enhanced,width=16cm,height=14cm,colback=white,colframe=cnom,boxrule=2pt,arc=0mm,
left=6mm,right=6mm,top=0mm,bottom=4mm,sharp corners,
overlay={\fill[cnom] (frame.north west) rectangle ([yshift=-1.25cm]frame.north east);
\node[anchor=west,text=white,font=\sffamily\bfseries\LARGE] at ([xshift=5mm,yshift=-0.62cm]frame.north west) {Question 6 B};
\node[anchor=east,text=white,font=\sffamily\itshape\large] at ([xshift=-5mm,yshift=-0.62cm]frame.north east) {Nombres et calculs};}]
\vspace*{1.75cm}
\begin{minipage}[t][6.9cm][c]{\linewidth}\LARGE\centering
\begin{minipage}{0.95\linewidth}\raggedright\linespread{1.25}\selectfont Développer $(3x-2)^2$.\end{minipage}
\end{minipage}

\vfill
\begin{tcolorbox}[enhanced,nobeforeafter,width=7.1cm,height=2.1cm,valign=center,colback=cnom!6,colframe=cnom,arc=2.5mm,boxrule=0.9pt,left=3mm]{\Large\bfseries\color{cnom}a)}\quad\begin{minipage}[c]{5.3cm}\raggedright\Large $9x^2-6x+4$\end{minipage}\end{tcolorbox}\hfill\begin{tcolorbox}[enhanced,nobeforeafter,width=7.1cm,height=2.1cm,valign=center,colback=cnom!6,colframe=cnom,arc=2.5mm,boxrule=0.9pt,left=3mm]{\Large\bfseries\color{cnom}b)}\quad\begin{minipage}[c]{5.3cm}\raggedright\Large $3x^2-12x+4$\end{minipage}\end{tcolorbox}\par\vspace{2.5mm}
\begin{tcolorbox}[enhanced,nobeforeafter,width=7.1cm,height=2.1cm,valign=center,colback=cnom!6,colframe=cnom,arc=2.5mm,boxrule=0.9pt,left=3mm]{\Large\bfseries\color{cnom}c)}\quad\begin{minipage}[c]{5.3cm}\raggedright\Large $9x^2+4$\end{minipage}\end{tcolorbox}\hfill\begin{tcolorbox}[enhanced,nobeforeafter,width=7.1cm,height=2.1cm,valign=center,colback=cnom!6,colframe=cnom,arc=2.5mm,boxrule=0.9pt,left=3mm]{\Large\bfseries\color{cnom}d)}\quad\begin{minipage}[c]{5.3cm}\raggedright\Large $9x^2-12x+4$\end{minipage}\end{tcolorbox}
\end{tcolorbox}
\newpage
\begin{tcolorbox}[enhanced,width=16cm,height=14cm,colback=white,colframe=cnom,boxrule=2pt,arc=0mm,
left=6mm,right=6mm,top=0mm,bottom=4mm,sharp corners,
overlay={\fill[cnom] (frame.north west) rectangle ([yshift=-1.25cm]frame.north east);
\node[anchor=west,text=white,font=\sffamily\bfseries\LARGE] at ([xshift=5mm,yshift=-0.62cm]frame.north west) {Question 7 B};
\node[anchor=east,text=white,font=\sffamily\itshape\large] at ([xshift=-5mm,yshift=-0.62cm]frame.north east) {Nombres et calculs};}]
\vspace*{1.75cm}
\begin{minipage}[t][6.9cm][c]{\linewidth}\LARGE\centering
\begin{minipage}{0.95\linewidth}\raggedright\linespread{1.25}\selectfont Factoriser $4x^2-1$.\end{minipage}
\end{minipage}

\vfill
\begin{tcolorbox}[enhanced,nobeforeafter,width=7.1cm,height=2.1cm,valign=center,colback=cnom!6,colframe=cnom,arc=2.5mm,boxrule=0.9pt,left=3mm]{\Large\bfseries\color{cnom}a)}\quad\begin{minipage}[c]{5.3cm}\raggedright\Large $(2x-1)(2x+1)$\end{minipage}\end{tcolorbox}\hfill\begin{tcolorbox}[enhanced,nobeforeafter,width=7.1cm,height=2.1cm,valign=center,colback=cnom!6,colframe=cnom,arc=2.5mm,boxrule=0.9pt,left=3mm]{\Large\bfseries\color{cnom}b)}\quad\begin{minipage}[c]{5.3cm}\raggedright\Large $(4x-1)(4x+1)$\end{minipage}\end{tcolorbox}\par\vspace{2.5mm}
\begin{tcolorbox}[enhanced,nobeforeafter,width=7.1cm,height=2.1cm,valign=center,colback=cnom!6,colframe=cnom,arc=2.5mm,boxrule=0.9pt,left=3mm]{\Large\bfseries\color{cnom}c)}\quad\begin{minipage}[c]{5.3cm}\raggedright\Large $(2x-1)^2$\end{minipage}\end{tcolorbox}\hfill\begin{tcolorbox}[enhanced,nobeforeafter,width=7.1cm,height=2.1cm,valign=center,colback=cnom!6,colframe=cnom,arc=2.5mm,boxrule=0.9pt,left=3mm]{\Large\bfseries\color{cnom}d)}\quad\begin{minipage}[c]{5.3cm}\raggedright\Large $(2x+1)^2$\end{minipage}\end{tcolorbox}
\end{tcolorbox}
\newpage
\begin{tcolorbox}[enhanced,width=16cm,height=14cm,colback=white,colframe=cnom,boxrule=2pt,arc=0mm,
left=6mm,right=6mm,top=0mm,bottom=4mm,sharp corners,
overlay={\fill[cnom] (frame.north west) rectangle ([yshift=-1.25cm]frame.north east);
\node[anchor=west,text=white,font=\sffamily\bfseries\LARGE] at ([xshift=5mm,yshift=-0.62cm]frame.north west) {Question 8 B};
\node[anchor=east,text=white,font=\sffamily\itshape\large] at ([xshift=-5mm,yshift=-0.62cm]frame.north east) {Nombres et calculs};}]
\vspace*{1.75cm}
\begin{minipage}[t][6.9cm][c]{\linewidth}\LARGE\centering
\begin{minipage}{0.95\linewidth}\raggedright\linespread{1.25}\selectfont Factoriser $9x^2-16$.\end{minipage}
\end{minipage}

\vfill
\begin{tcolorbox}[enhanced,nobeforeafter,width=7.1cm,height=2.1cm,valign=center,colback=cnom!6,colframe=cnom,arc=2.5mm,boxrule=0.9pt,left=3mm]{\Large\bfseries\color{cnom}a)}\quad\begin{minipage}[c]{5.3cm}\raggedright\Large $(3x+4)^2$\end{minipage}\end{tcolorbox}\hfill\begin{tcolorbox}[enhanced,nobeforeafter,width=7.1cm,height=2.1cm,valign=center,colback=cnom!6,colframe=cnom,arc=2.5mm,boxrule=0.9pt,left=3mm]{\Large\bfseries\color{cnom}b)}\quad\begin{minipage}[c]{5.3cm}\raggedright\Large $(3x-4)(3x+4)$\end{minipage}\end{tcolorbox}\par\vspace{2.5mm}
\begin{tcolorbox}[enhanced,nobeforeafter,width=7.1cm,height=2.1cm,valign=center,colback=cnom!6,colframe=cnom,arc=2.5mm,boxrule=0.9pt,left=3mm]{\Large\bfseries\color{cnom}c)}\quad\begin{minipage}[c]{5.3cm}\raggedright\Large $(9x-4)(9x+4)$\end{minipage}\end{tcolorbox}\hfill\begin{tcolorbox}[enhanced,nobeforeafter,width=7.1cm,height=2.1cm,valign=center,colback=cnom!6,colframe=cnom,arc=2.5mm,boxrule=0.9pt,left=3mm]{\Large\bfseries\color{cnom}d)}\quad\begin{minipage}[c]{5.3cm}\raggedright\Large $(3x-4)^2$\end{minipage}\end{tcolorbox}
\end{tcolorbox}
\newpage
\begin{tcolorbox}[enhanced,width=16cm,height=14cm,colback=white,colframe=cnom,boxrule=2pt,arc=0mm,
left=6mm,right=6mm,top=0mm,bottom=4mm,sharp corners,
overlay={\fill[cnom] (frame.north west) rectangle ([yshift=-1.25cm]frame.north east);
\node[anchor=west,text=white,font=\sffamily\bfseries\LARGE] at ([xshift=5mm,yshift=-0.62cm]frame.north west) {Question 9 B};
\node[anchor=east,text=white,font=\sffamily\itshape\large] at ([xshift=-5mm,yshift=-0.62cm]frame.north east) {Nombres et calculs};}]
\vspace*{1.75cm}
\begin{minipage}[t][6.9cm][c]{\linewidth}\LARGE\centering
\begin{minipage}{0.95\linewidth}\raggedright\linespread{1.25}\selectfont Factoriser $9x^2-12x+4$.\end{minipage}
\end{minipage}

\vfill
\begin{tcolorbox}[enhanced,nobeforeafter,width=7.1cm,height=2.1cm,valign=center,colback=cnom!6,colframe=cnom,arc=2.5mm,boxrule=0.9pt,left=3mm]{\Large\bfseries\color{cnom}a)}\quad\begin{minipage}[c]{5.3cm}\raggedright\Large $(9x-2)^2$\end{minipage}\end{tcolorbox}\hfill\begin{tcolorbox}[enhanced,nobeforeafter,width=7.1cm,height=2.1cm,valign=center,colback=cnom!6,colframe=cnom,arc=2.5mm,boxrule=0.9pt,left=3mm]{\Large\bfseries\color{cnom}b)}\quad\begin{minipage}[c]{5.3cm}\raggedright\Large $(3x+2)^2$\end{minipage}\end{tcolorbox}\par\vspace{2.5mm}
\begin{tcolorbox}[enhanced,nobeforeafter,width=7.1cm,height=2.1cm,valign=center,colback=cnom!6,colframe=cnom,arc=2.5mm,boxrule=0.9pt,left=3mm]{\Large\bfseries\color{cnom}c)}\quad\begin{minipage}[c]{5.3cm}\raggedright\Large $(3x+2)(3x-2)$\end{minipage}\end{tcolorbox}\hfill\begin{tcolorbox}[enhanced,nobeforeafter,width=7.1cm,height=2.1cm,valign=center,colback=cnom!6,colframe=cnom,arc=2.5mm,boxrule=0.9pt,left=3mm]{\Large\bfseries\color{cnom}d)}\quad\begin{minipage}[c]{5.3cm}\raggedright\Large $(3x-2)^2$\end{minipage}\end{tcolorbox}
\end{tcolorbox}
\newpage
\begin{tcolorbox}[enhanced,width=16cm,height=14cm,colback=white,colframe=cnom,boxrule=2pt,arc=0mm,
left=6mm,right=6mm,top=0mm,bottom=4mm,sharp corners,
overlay={\fill[cnom] (frame.north west) rectangle ([yshift=-1.25cm]frame.north east);
\node[anchor=west,text=white,font=\sffamily\bfseries\LARGE] at ([xshift=5mm,yshift=-0.62cm]frame.north west) {Question 10 B};
\node[anchor=east,text=white,font=\sffamily\itshape\large] at ([xshift=-5mm,yshift=-0.62cm]frame.north east) {Nombres et calculs};}]
\vspace*{1.75cm}
\begin{minipage}[t][6.9cm][c]{\linewidth}\LARGE\centering
\begin{minipage}{0.95\linewidth}\raggedright\linespread{1.25}\selectfont Factoriser $4x^2+20x+25$.\end{minipage}
\end{minipage}

\vfill
\begin{tcolorbox}[enhanced,nobeforeafter,width=7.1cm,height=2.1cm,valign=center,colback=cnom!6,colframe=cnom,arc=2.5mm,boxrule=0.9pt,left=3mm]{\Large\bfseries\color{cnom}a)}\quad\begin{minipage}[c]{5.3cm}\raggedright\Large $(4x+5)^2$\end{minipage}\end{tcolorbox}\hfill\begin{tcolorbox}[enhanced,nobeforeafter,width=7.1cm,height=2.1cm,valign=center,colback=cnom!6,colframe=cnom,arc=2.5mm,boxrule=0.9pt,left=3mm]{\Large\bfseries\color{cnom}b)}\quad\begin{minipage}[c]{5.3cm}\raggedright\Large $(2x-5)^2$\end{minipage}\end{tcolorbox}\par\vspace{2.5mm}
\begin{tcolorbox}[enhanced,nobeforeafter,width=7.1cm,height=2.1cm,valign=center,colback=cnom!6,colframe=cnom,arc=2.5mm,boxrule=0.9pt,left=3mm]{\Large\bfseries\color{cnom}c)}\quad\begin{minipage}[c]{5.3cm}\raggedright\Large $(2x-5)(2x+5)$\end{minipage}\end{tcolorbox}\hfill\begin{tcolorbox}[enhanced,nobeforeafter,width=7.1cm,height=2.1cm,valign=center,colback=cnom!6,colframe=cnom,arc=2.5mm,boxrule=0.9pt,left=3mm]{\Large\bfseries\color{cnom}d)}\quad\begin{minipage}[c]{5.3cm}\raggedright\Large $(2x+5)^2$\end{minipage}\end{tcolorbox}
\end{tcolorbox}
\newpage
\begin{tcolorbox}[enhanced,width=16cm,height=14cm,colback=white,colframe=cnom,boxrule=2pt,arc=0mm,
left=6mm,right=6mm,top=0mm,bottom=4mm,sharp corners,
overlay={\fill[cnom] (frame.north west) rectangle ([yshift=-1.25cm]frame.north east);
\node[anchor=west,text=white,font=\sffamily\bfseries\LARGE] at ([xshift=5mm,yshift=-0.62cm]frame.north west) {Question 11 B};
\node[anchor=east,text=white,font=\sffamily\itshape\large] at ([xshift=-5mm,yshift=-0.62cm]frame.north east) {Nombres et calculs};}]
\vspace*{1.75cm}
\begin{minipage}[t][6.9cm][c]{\linewidth}\LARGE\centering
\begin{minipage}{0.95\linewidth}\raggedright\linespread{1.25}\selectfont Que signifie $x\in]-\infty\,;\,4]$~?\end{minipage}
\end{minipage}

\vfill
\begin{tcolorbox}[enhanced,nobeforeafter,width=7.1cm,height=2.1cm,valign=center,colback=cnom!6,colframe=cnom,arc=2.5mm,boxrule=0.9pt,left=3mm]{\Large\bfseries\color{cnom}a)}\quad\begin{minipage}[c]{5.3cm}\raggedright\Large $x>4$\end{minipage}\end{tcolorbox}\hfill\begin{tcolorbox}[enhanced,nobeforeafter,width=7.1cm,height=2.1cm,valign=center,colback=cnom!6,colframe=cnom,arc=2.5mm,boxrule=0.9pt,left=3mm]{\Large\bfseries\color{cnom}b)}\quad\begin{minipage}[c]{5.3cm}\raggedright\Large $x\leqslant 4$\end{minipage}\end{tcolorbox}\par\vspace{2.5mm}
\begin{tcolorbox}[enhanced,nobeforeafter,width=7.1cm,height=2.1cm,valign=center,colback=cnom!6,colframe=cnom,arc=2.5mm,boxrule=0.9pt,left=3mm]{\Large\bfseries\color{cnom}c)}\quad\begin{minipage}[c]{5.3cm}\raggedright\Large $x<4$\end{minipage}\end{tcolorbox}\hfill\begin{tcolorbox}[enhanced,nobeforeafter,width=7.1cm,height=2.1cm,valign=center,colback=cnom!6,colframe=cnom,arc=2.5mm,boxrule=0.9pt,left=3mm]{\Large\bfseries\color{cnom}d)}\quad\begin{minipage}[c]{5.3cm}\raggedright\Large $x\geqslant 4$\end{minipage}\end{tcolorbox}
\end{tcolorbox}
\newpage
\begin{tcolorbox}[enhanced,width=16cm,height=14cm,colback=white,colframe=cnom,boxrule=2pt,arc=0mm,
left=6mm,right=6mm,top=0mm,bottom=4mm,sharp corners,
overlay={\fill[cnom] (frame.north west) rectangle ([yshift=-1.25cm]frame.north east);
\node[anchor=west,text=white,font=\sffamily\bfseries\LARGE] at ([xshift=5mm,yshift=-0.62cm]frame.north west) {Question 12 B};
\node[anchor=east,text=white,font=\sffamily\itshape\large] at ([xshift=-5mm,yshift=-0.62cm]frame.north east) {Nombres et calculs};}]
\vspace*{1.75cm}
\begin{minipage}[t][6.9cm][c]{\linewidth}\LARGE\centering
\begin{minipage}{0.95\linewidth}\raggedright\linespread{1.25}\selectfont Que signifie $x\in[0\,;\,1]$~?\end{minipage}
\end{minipage}

\vfill
\begin{tcolorbox}[enhanced,nobeforeafter,width=7.1cm,height=2.1cm,valign=center,colback=cnom!6,colframe=cnom,arc=2.5mm,boxrule=0.9pt,left=3mm]{\Large\bfseries\color{cnom}a)}\quad\begin{minipage}[c]{5.3cm}\raggedright\Large $0\leqslant x<1$\end{minipage}\end{tcolorbox}\hfill\begin{tcolorbox}[enhanced,nobeforeafter,width=7.1cm,height=2.1cm,valign=center,colback=cnom!6,colframe=cnom,arc=2.5mm,boxrule=0.9pt,left=3mm]{\Large\bfseries\color{cnom}b)}\quad\begin{minipage}[c]{5.3cm}\raggedright\Large $0<x\leqslant 1$\end{minipage}\end{tcolorbox}\par\vspace{2.5mm}
\begin{tcolorbox}[enhanced,nobeforeafter,width=7.1cm,height=2.1cm,valign=center,colback=cnom!6,colframe=cnom,arc=2.5mm,boxrule=0.9pt,left=3mm]{\Large\bfseries\color{cnom}c)}\quad\begin{minipage}[c]{5.3cm}\raggedright\Large $0\leqslant x\leqslant 1$\end{minipage}\end{tcolorbox}\hfill\begin{tcolorbox}[enhanced,nobeforeafter,width=7.1cm,height=2.1cm,valign=center,colback=cnom!6,colframe=cnom,arc=2.5mm,boxrule=0.9pt,left=3mm]{\Large\bfseries\color{cnom}d)}\quad\begin{minipage}[c]{5.3cm}\raggedright\Large $0<x<1$\end{minipage}\end{tcolorbox}
\end{tcolorbox}
\newpage
\begin{tcolorbox}[enhanced,width=16cm,height=14cm,colback=white,colframe=cnom,boxrule=2pt,arc=0mm,
left=6mm,right=6mm,top=0mm,bottom=4mm,sharp corners,
overlay={\fill[cnom] (frame.north west) rectangle ([yshift=-1.25cm]frame.north east);
\node[anchor=west,text=white,font=\sffamily\bfseries\LARGE] at ([xshift=5mm,yshift=-0.62cm]frame.north west) {Question 13 B};
\node[anchor=east,text=white,font=\sffamily\itshape\large] at ([xshift=-5mm,yshift=-0.62cm]frame.north east) {Nombres et calculs};}]
\vspace*{1.75cm}
\begin{minipage}[t][6.9cm][c]{\linewidth}\LARGE\centering
\begin{minipage}{0.95\linewidth}\raggedright\linespread{1.25}\selectfont Pour quels réels $x$ a-t-on $|x|\leqslant 5$~?\end{minipage}
\end{minipage}

\vfill
\begin{tcolorbox}[enhanced,nobeforeafter,width=7.1cm,height=2.1cm,valign=center,colback=cnom!6,colframe=cnom,arc=2.5mm,boxrule=0.9pt,left=3mm]{\Large\bfseries\color{cnom}a)}\quad\begin{minipage}[c]{5.3cm}\raggedright\Large $x\in[-5\,;\,5]$\end{minipage}\end{tcolorbox}\hfill\begin{tcolorbox}[enhanced,nobeforeafter,width=7.1cm,height=2.1cm,valign=center,colback=cnom!6,colframe=cnom,arc=2.5mm,boxrule=0.9pt,left=3mm]{\Large\bfseries\color{cnom}b)}\quad\begin{minipage}[c]{5.3cm}\raggedright\Large $x\in[-10\,;\,10]$\end{minipage}\end{tcolorbox}\par\vspace{2.5mm}
\begin{tcolorbox}[enhanced,nobeforeafter,width=7.1cm,height=2.1cm,valign=center,colback=cnom!6,colframe=cnom,arc=2.5mm,boxrule=0.9pt,left=3mm]{\Large\bfseries\color{cnom}c)}\quad\begin{minipage}[c]{5.3cm}\raggedright\Large $x\in[0\,;\,5]$\end{minipage}\end{tcolorbox}\hfill\begin{tcolorbox}[enhanced,nobeforeafter,width=7.1cm,height=2.1cm,valign=center,colback=cnom!6,colframe=cnom,arc=2.5mm,boxrule=0.9pt,left=3mm]{\Large\bfseries\color{cnom}d)}\quad\begin{minipage}[c]{5.3cm}\raggedright\Large $x\in]-5\,;\,5[$\end{minipage}\end{tcolorbox}
\end{tcolorbox}
\newpage
\begin{tcolorbox}[enhanced,width=16cm,height=14cm,colback=white,colframe=cnom,boxrule=2pt,arc=0mm,
left=6mm,right=6mm,top=0mm,bottom=4mm,sharp corners,
overlay={\fill[cnom] (frame.north west) rectangle ([yshift=-1.25cm]frame.north east);
\node[anchor=west,text=white,font=\sffamily\bfseries\LARGE] at ([xshift=5mm,yshift=-0.62cm]frame.north west) {Question 14 B};
\node[anchor=east,text=white,font=\sffamily\itshape\large] at ([xshift=-5mm,yshift=-0.62cm]frame.north east) {Nombres et calculs};}]
\vspace*{1.75cm}
\begin{minipage}[t][6.9cm][c]{\linewidth}\LARGE\centering
\begin{minipage}{0.95\linewidth}\raggedright\linespread{1.25}\selectfont Pour quels réels $x$ a-t-on $|x-2|\leqslant 7$~?\end{minipage}
\end{minipage}

\vfill
\begin{tcolorbox}[enhanced,nobeforeafter,width=7.1cm,height=2.1cm,valign=center,colback=cnom!6,colframe=cnom,arc=2.5mm,boxrule=0.9pt,left=3mm]{\Large\bfseries\color{cnom}a)}\quad\begin{minipage}[c]{5.3cm}\raggedright\Large $x\in[-12\,;\,16]$\end{minipage}\end{tcolorbox}\hfill\begin{tcolorbox}[enhanced,nobeforeafter,width=7.1cm,height=2.1cm,valign=center,colback=cnom!6,colframe=cnom,arc=2.5mm,boxrule=0.9pt,left=3mm]{\Large\bfseries\color{cnom}b)}\quad\begin{minipage}[c]{5.3cm}\raggedright\Large $x\in[-9\,;\,5]$\end{minipage}\end{tcolorbox}\par\vspace{2.5mm}
\begin{tcolorbox}[enhanced,nobeforeafter,width=7.1cm,height=2.1cm,valign=center,colback=cnom!6,colframe=cnom,arc=2.5mm,boxrule=0.9pt,left=3mm]{\Large\bfseries\color{cnom}c)}\quad\begin{minipage}[c]{5.3cm}\raggedright\Large $x\in[-5\,;\,9]$\end{minipage}\end{tcolorbox}\hfill\begin{tcolorbox}[enhanced,nobeforeafter,width=7.1cm,height=2.1cm,valign=center,colback=cnom!6,colframe=cnom,arc=2.5mm,boxrule=0.9pt,left=3mm]{\Large\bfseries\color{cnom}d)}\quad\begin{minipage}[c]{5.3cm}\raggedright\Large $x\in[2\,;\,9]$\end{minipage}\end{tcolorbox}
\end{tcolorbox}
\newpage
\begin{tcolorbox}[enhanced,width=16cm,height=14cm,colback=white,colframe=cnom,boxrule=2pt,arc=0mm,
left=6mm,right=6mm,top=0mm,bottom=4mm,sharp corners,
overlay={\fill[cnom] (frame.north west) rectangle ([yshift=-1.25cm]frame.north east);
\node[anchor=west,text=white,font=\sffamily\bfseries\LARGE] at ([xshift=5mm,yshift=-0.62cm]frame.north west) {Question 15 B};
\node[anchor=east,text=white,font=\sffamily\itshape\large] at ([xshift=-5mm,yshift=-0.62cm]frame.north east) {Nombres et calculs};}]
\vspace*{1.75cm}
\begin{minipage}[t][6.9cm][c]{\linewidth}\LARGE\centering
\begin{minipage}{0.95\linewidth}\raggedright\linespread{1.25}\selectfont Écrire $\dfrac{5^{-2}\times5^{6}}{5^{1}}$ sous la forme d'une puissance de $5$.\end{minipage}
\end{minipage}

\vfill
\begin{tcolorbox}[enhanced,nobeforeafter,width=7.1cm,height=2.1cm,valign=center,colback=cnom!6,colframe=cnom,arc=2.5mm,boxrule=0.9pt,left=3mm]{\Large\bfseries\color{cnom}a)}\quad\begin{minipage}[c]{5.3cm}\raggedright\Large $5^{-13}$\end{minipage}\end{tcolorbox}\hfill\begin{tcolorbox}[enhanced,nobeforeafter,width=7.1cm,height=2.1cm,valign=center,colback=cnom!6,colframe=cnom,arc=2.5mm,boxrule=0.9pt,left=3mm]{\Large\bfseries\color{cnom}b)}\quad\begin{minipage}[c]{5.3cm}\raggedright\Large $5^{5}$\end{minipage}\end{tcolorbox}\par\vspace{2.5mm}
\begin{tcolorbox}[enhanced,nobeforeafter,width=7.1cm,height=2.1cm,valign=center,colback=cnom!6,colframe=cnom,arc=2.5mm,boxrule=0.9pt,left=3mm]{\Large\bfseries\color{cnom}c)}\quad\begin{minipage}[c]{5.3cm}\raggedright\Large $5^{-9}$\end{minipage}\end{tcolorbox}\hfill\begin{tcolorbox}[enhanced,nobeforeafter,width=7.1cm,height=2.1cm,valign=center,colback=cnom!6,colframe=cnom,arc=2.5mm,boxrule=0.9pt,left=3mm]{\Large\bfseries\color{cnom}d)}\quad\begin{minipage}[c]{5.3cm}\raggedright\Large $5^{3}$\end{minipage}\end{tcolorbox}
\end{tcolorbox}
\newpage
\begin{tcolorbox}[enhanced,width=16cm,height=14cm,colback=white,colframe=cnom,boxrule=2pt,arc=0mm,
left=6mm,right=6mm,top=0mm,bottom=4mm,sharp corners,
overlay={\fill[cnom] (frame.north west) rectangle ([yshift=-1.25cm]frame.north east);
\node[anchor=west,text=white,font=\sffamily\bfseries\LARGE] at ([xshift=5mm,yshift=-0.62cm]frame.north west) {Question 16 B};
\node[anchor=east,text=white,font=\sffamily\itshape\large] at ([xshift=-5mm,yshift=-0.62cm]frame.north east) {Nombres et calculs};}]
\vspace*{1.75cm}
\begin{minipage}[t][6.9cm][c]{\linewidth}\LARGE\centering
\begin{minipage}{0.95\linewidth}\raggedright\linespread{1.25}\selectfont Écrire $\dfrac{10^{3}\times10^{-5}}{10^{-4}}$ sous la forme d'une puissance de $10$.\end{minipage}
\end{minipage}

\vfill
\begin{tcolorbox}[enhanced,nobeforeafter,width=7.1cm,height=2.1cm,valign=center,colback=cnom!6,colframe=cnom,arc=2.5mm,boxrule=0.9pt,left=3mm]{\Large\bfseries\color{cnom}a)}\quad\begin{minipage}[c]{5.3cm}\raggedright\Large $10^{2}$\end{minipage}\end{tcolorbox}\hfill\begin{tcolorbox}[enhanced,nobeforeafter,width=7.1cm,height=2.1cm,valign=center,colback=cnom!6,colframe=cnom,arc=2.5mm,boxrule=0.9pt,left=3mm]{\Large\bfseries\color{cnom}b)}\quad\begin{minipage}[c]{5.3cm}\raggedright\Large $10^{12}$\end{minipage}\end{tcolorbox}\par\vspace{2.5mm}
\begin{tcolorbox}[enhanced,nobeforeafter,width=7.1cm,height=2.1cm,valign=center,colback=cnom!6,colframe=cnom,arc=2.5mm,boxrule=0.9pt,left=3mm]{\Large\bfseries\color{cnom}c)}\quad\begin{minipage}[c]{5.3cm}\raggedright\Large $10^{-6}$\end{minipage}\end{tcolorbox}\hfill\begin{tcolorbox}[enhanced,nobeforeafter,width=7.1cm,height=2.1cm,valign=center,colback=cnom!6,colframe=cnom,arc=2.5mm,boxrule=0.9pt,left=3mm]{\Large\bfseries\color{cnom}d)}\quad\begin{minipage}[c]{5.3cm}\raggedright\Large $10^{-11}$\end{minipage}\end{tcolorbox}
\end{tcolorbox}
\newpage
\begin{tcolorbox}[enhanced,width=16cm,height=14cm,colback=white,colframe=csta,boxrule=2pt,arc=0mm,
left=6mm,right=6mm,top=0mm,bottom=4mm,sharp corners,
overlay={\fill[csta] (frame.north west) rectangle ([yshift=-1.25cm]frame.north east);
\node[anchor=west,text=white,font=\sffamily\bfseries\LARGE] at ([xshift=5mm,yshift=-0.62cm]frame.north west) {Question 17 B};
\node[anchor=east,text=white,font=\sffamily\itshape\large] at ([xshift=-5mm,yshift=-0.62cm]frame.north east) {Évolutions et statistiques};}]
\vspace*{1.75cm}
\begin{minipage}[t][6.9cm][c]{\linewidth}\LARGE\centering
\begin{minipage}{0.95\linewidth}\raggedright\linespread{1.25}\selectfont Quel est le coefficient multiplicateur associé à une hausse de $5\,\%$~?\end{minipage}
\end{minipage}

\vfill
\begin{tcolorbox}[enhanced,nobeforeafter,width=7.1cm,height=2.1cm,valign=center,colback=csta!6,colframe=csta,arc=2.5mm,boxrule=0.9pt,left=3mm]{\Large\bfseries\color{csta}a)}\quad\begin{minipage}[c]{5.3cm}\raggedright\Large $1{,}5$\end{minipage}\end{tcolorbox}\hfill\begin{tcolorbox}[enhanced,nobeforeafter,width=7.1cm,height=2.1cm,valign=center,colback=csta!6,colframe=csta,arc=2.5mm,boxrule=0.9pt,left=3mm]{\Large\bfseries\color{csta}b)}\quad\begin{minipage}[c]{5.3cm}\raggedright\Large $0{,}95$\end{minipage}\end{tcolorbox}\par\vspace{2.5mm}
\begin{tcolorbox}[enhanced,nobeforeafter,width=7.1cm,height=2.1cm,valign=center,colback=csta!6,colframe=csta,arc=2.5mm,boxrule=0.9pt,left=3mm]{\Large\bfseries\color{csta}c)}\quad\begin{minipage}[c]{5.3cm}\raggedright\Large $1{,}05$\end{minipage}\end{tcolorbox}\hfill\begin{tcolorbox}[enhanced,nobeforeafter,width=7.1cm,height=2.1cm,valign=center,colback=csta!6,colframe=csta,arc=2.5mm,boxrule=0.9pt,left=3mm]{\Large\bfseries\color{csta}d)}\quad\begin{minipage}[c]{5.3cm}\raggedright\Large $0{,}05$\end{minipage}\end{tcolorbox}
\end{tcolorbox}
\newpage
\begin{tcolorbox}[enhanced,width=16cm,height=14cm,colback=white,colframe=csta,boxrule=2pt,arc=0mm,
left=6mm,right=6mm,top=0mm,bottom=4mm,sharp corners,
overlay={\fill[csta] (frame.north west) rectangle ([yshift=-1.25cm]frame.north east);
\node[anchor=west,text=white,font=\sffamily\bfseries\LARGE] at ([xshift=5mm,yshift=-0.62cm]frame.north west) {Question 18 B};
\node[anchor=east,text=white,font=\sffamily\itshape\large] at ([xshift=-5mm,yshift=-0.62cm]frame.north east) {Évolutions et statistiques};}]
\vspace*{1.75cm}
\begin{minipage}[t][6.9cm][c]{\linewidth}\LARGE\centering
\begin{minipage}{0.95\linewidth}\raggedright\linespread{1.25}\selectfont Quel est le coefficient multiplicateur associé à une baisse de $30\,\%$~?\end{minipage}
\end{minipage}

\vfill
\begin{tcolorbox}[enhanced,nobeforeafter,width=7.1cm,height=2.1cm,valign=center,colback=csta!6,colframe=csta,arc=2.5mm,boxrule=0.9pt,left=3mm]{\Large\bfseries\color{csta}a)}\quad\begin{minipage}[c]{5.3cm}\raggedright\Large $-2$\end{minipage}\end{tcolorbox}\hfill\begin{tcolorbox}[enhanced,nobeforeafter,width=7.1cm,height=2.1cm,valign=center,colback=csta!6,colframe=csta,arc=2.5mm,boxrule=0.9pt,left=3mm]{\Large\bfseries\color{csta}b)}\quad\begin{minipage}[c]{5.3cm}\raggedright\Large $-0{,}3$\end{minipage}\end{tcolorbox}\par\vspace{2.5mm}
\begin{tcolorbox}[enhanced,nobeforeafter,width=7.1cm,height=2.1cm,valign=center,colback=csta!6,colframe=csta,arc=2.5mm,boxrule=0.9pt,left=3mm]{\Large\bfseries\color{csta}c)}\quad\begin{minipage}[c]{5.3cm}\raggedright\Large $0{,}7$\end{minipage}\end{tcolorbox}\hfill\begin{tcolorbox}[enhanced,nobeforeafter,width=7.1cm,height=2.1cm,valign=center,colback=csta!6,colframe=csta,arc=2.5mm,boxrule=0.9pt,left=3mm]{\Large\bfseries\color{csta}d)}\quad\begin{minipage}[c]{5.3cm}\raggedright\Large $1{,}3$\end{minipage}\end{tcolorbox}
\end{tcolorbox}
\newpage
\begin{tcolorbox}[enhanced,width=16cm,height=14cm,colback=white,colframe=csta,boxrule=2pt,arc=0mm,
left=6mm,right=6mm,top=0mm,bottom=4mm,sharp corners,
overlay={\fill[csta] (frame.north west) rectangle ([yshift=-1.25cm]frame.north east);
\node[anchor=west,text=white,font=\sffamily\bfseries\LARGE] at ([xshift=5mm,yshift=-0.62cm]frame.north west) {Question 19 B};
\node[anchor=east,text=white,font=\sffamily\itshape\large] at ([xshift=-5mm,yshift=-0.62cm]frame.north east) {Évolutions et statistiques};}]
\vspace*{1.75cm}
\begin{minipage}[t][6.9cm][c]{\linewidth}\LARGE\centering
\begin{minipage}{0.95\linewidth}\raggedright\linespread{1.25}\selectfont Un prix de $1\,200$ € augmente de $15\,\%$. Quel est le nouveau prix~?\end{minipage}
\end{minipage}

\vfill
\begin{tcolorbox}[enhanced,nobeforeafter,width=7.1cm,height=2.1cm,valign=center,colback=csta!6,colframe=csta,arc=2.5mm,boxrule=0.9pt,left=3mm]{\Large\bfseries\color{csta}a)}\quad\begin{minipage}[c]{5.3cm}\raggedright\Large $1\,215~\text{€}$\end{minipage}\end{tcolorbox}\hfill\begin{tcolorbox}[enhanced,nobeforeafter,width=7.1cm,height=2.1cm,valign=center,colback=csta!6,colframe=csta,arc=2.5mm,boxrule=0.9pt,left=3mm]{\Large\bfseries\color{csta}b)}\quad\begin{minipage}[c]{5.3cm}\raggedright\Large $180~\text{€}$\end{minipage}\end{tcolorbox}\par\vspace{2.5mm}
\begin{tcolorbox}[enhanced,nobeforeafter,width=7.1cm,height=2.1cm,valign=center,colback=csta!6,colframe=csta,arc=2.5mm,boxrule=0.9pt,left=3mm]{\Large\bfseries\color{csta}c)}\quad\begin{minipage}[c]{5.3cm}\raggedright\Large $1\,020~\text{€}$\end{minipage}\end{tcolorbox}\hfill\begin{tcolorbox}[enhanced,nobeforeafter,width=7.1cm,height=2.1cm,valign=center,colback=csta!6,colframe=csta,arc=2.5mm,boxrule=0.9pt,left=3mm]{\Large\bfseries\color{csta}d)}\quad\begin{minipage}[c]{5.3cm}\raggedright\Large $1\,380~\text{€}$\end{minipage}\end{tcolorbox}
\end{tcolorbox}
\newpage
\begin{tcolorbox}[enhanced,width=16cm,height=14cm,colback=white,colframe=csta,boxrule=2pt,arc=0mm,
left=6mm,right=6mm,top=0mm,bottom=4mm,sharp corners,
overlay={\fill[csta] (frame.north west) rectangle ([yshift=-1.25cm]frame.north east);
\node[anchor=west,text=white,font=\sffamily\bfseries\LARGE] at ([xshift=5mm,yshift=-0.62cm]frame.north west) {Question 20 B};
\node[anchor=east,text=white,font=\sffamily\itshape\large] at ([xshift=-5mm,yshift=-0.62cm]frame.north east) {Évolutions et statistiques};}]
\vspace*{1.75cm}
\begin{minipage}[t][6.9cm][c]{\linewidth}\LARGE\centering
\begin{minipage}{0.95\linewidth}\raggedright\linespread{1.25}\selectfont Un prix de $45$ € baisse de $10\,\%$. Quel est le nouveau prix~?\end{minipage}
\end{minipage}

\vfill
\begin{tcolorbox}[enhanced,nobeforeafter,width=7.1cm,height=2.1cm,valign=center,colback=csta!6,colframe=csta,arc=2.5mm,boxrule=0.9pt,left=3mm]{\Large\bfseries\color{csta}a)}\quad\begin{minipage}[c]{5.3cm}\raggedright\Large $35~\text{€}$\end{minipage}\end{tcolorbox}\hfill\begin{tcolorbox}[enhanced,nobeforeafter,width=7.1cm,height=2.1cm,valign=center,colback=csta!6,colframe=csta,arc=2.5mm,boxrule=0.9pt,left=3mm]{\Large\bfseries\color{csta}b)}\quad\begin{minipage}[c]{5.3cm}\raggedright\Large $4{,}5~\text{€}$\end{minipage}\end{tcolorbox}\par\vspace{2.5mm}
\begin{tcolorbox}[enhanced,nobeforeafter,width=7.1cm,height=2.1cm,valign=center,colback=csta!6,colframe=csta,arc=2.5mm,boxrule=0.9pt,left=3mm]{\Large\bfseries\color{csta}c)}\quad\begin{minipage}[c]{5.3cm}\raggedright\Large $49{,}5~\text{€}$\end{minipage}\end{tcolorbox}\hfill\begin{tcolorbox}[enhanced,nobeforeafter,width=7.1cm,height=2.1cm,valign=center,colback=csta!6,colframe=csta,arc=2.5mm,boxrule=0.9pt,left=3mm]{\Large\bfseries\color{csta}d)}\quad\begin{minipage}[c]{5.3cm}\raggedright\Large $40{,}5~\text{€}$\end{minipage}\end{tcolorbox}
\end{tcolorbox}
\newpage
\begin{tcolorbox}[enhanced,width=16cm,height=14cm,colback=white,colframe=csta,boxrule=2pt,arc=0mm,
left=6mm,right=6mm,top=0mm,bottom=4mm,sharp corners,
overlay={\fill[csta] (frame.north west) rectangle ([yshift=-1.25cm]frame.north east);
\node[anchor=west,text=white,font=\sffamily\bfseries\LARGE] at ([xshift=5mm,yshift=-0.62cm]frame.north west) {Question 21 B};
\node[anchor=east,text=white,font=\sffamily\itshape\large] at ([xshift=-5mm,yshift=-0.62cm]frame.north east) {Évolutions et statistiques};}]
\vspace*{1.75cm}
\begin{minipage}[t][6.9cm][c]{\linewidth}\LARGE\centering
\begin{minipage}{0.95\linewidth}\raggedright\linespread{1.25}\selectfont Une quantité passe de $200$ à $230$. Quel est son taux d'évolution~?\end{minipage}
\end{minipage}

\vfill
\begin{tcolorbox}[enhanced,nobeforeafter,width=7.1cm,height=2.1cm,valign=center,colback=csta!6,colframe=csta,arc=2.5mm,boxrule=0.9pt,left=3mm]{\Large\bfseries\color{csta}a)}\quad\begin{minipage}[c]{5.3cm}\raggedright\Large $+15\,\%$\end{minipage}\end{tcolorbox}\hfill\begin{tcolorbox}[enhanced,nobeforeafter,width=7.1cm,height=2.1cm,valign=center,colback=csta!6,colframe=csta,arc=2.5mm,boxrule=0.9pt,left=3mm]{\Large\bfseries\color{csta}b)}\quad\begin{minipage}[c]{5.3cm}\raggedright\Large $+30\,\%$\end{minipage}\end{tcolorbox}\par\vspace{2.5mm}
\begin{tcolorbox}[enhanced,nobeforeafter,width=7.1cm,height=2.1cm,valign=center,colback=csta!6,colframe=csta,arc=2.5mm,boxrule=0.9pt,left=3mm]{\Large\bfseries\color{csta}c)}\quad\begin{minipage}[c]{5.3cm}\raggedright\Large $-15\,\%$\end{minipage}\end{tcolorbox}\hfill\begin{tcolorbox}[enhanced,nobeforeafter,width=7.1cm,height=2.1cm,valign=center,colback=csta!6,colframe=csta,arc=2.5mm,boxrule=0.9pt,left=3mm]{\Large\bfseries\color{csta}d)}\quad\begin{minipage}[c]{5.3cm}\raggedright\Large $+7{,}5\,\%$\end{minipage}\end{tcolorbox}
\end{tcolorbox}
\newpage
\begin{tcolorbox}[enhanced,width=16cm,height=14cm,colback=white,colframe=csta,boxrule=2pt,arc=0mm,
left=6mm,right=6mm,top=0mm,bottom=4mm,sharp corners,
overlay={\fill[csta] (frame.north west) rectangle ([yshift=-1.25cm]frame.north east);
\node[anchor=west,text=white,font=\sffamily\bfseries\LARGE] at ([xshift=5mm,yshift=-0.62cm]frame.north west) {Question 22 B};
\node[anchor=east,text=white,font=\sffamily\itshape\large] at ([xshift=-5mm,yshift=-0.62cm]frame.north east) {Évolutions et statistiques};}]
\vspace*{1.75cm}
\begin{minipage}[t][6.9cm][c]{\linewidth}\LARGE\centering
\begin{minipage}{0.95\linewidth}\raggedright\linespread{1.25}\selectfont Une quantité passe de $40$ à $50$. Quel est son taux d'évolution~?\end{minipage}
\end{minipage}

\vfill
\begin{tcolorbox}[enhanced,nobeforeafter,width=7.1cm,height=2.1cm,valign=center,colback=csta!6,colframe=csta,arc=2.5mm,boxrule=0.9pt,left=3mm]{\Large\bfseries\color{csta}a)}\quad\begin{minipage}[c]{5.3cm}\raggedright\Large $+25\,\%$\end{minipage}\end{tcolorbox}\hfill\begin{tcolorbox}[enhanced,nobeforeafter,width=7.1cm,height=2.1cm,valign=center,colback=csta!6,colframe=csta,arc=2.5mm,boxrule=0.9pt,left=3mm]{\Large\bfseries\color{csta}b)}\quad\begin{minipage}[c]{5.3cm}\raggedright\Large $-25\,\%$\end{minipage}\end{tcolorbox}\par\vspace{2.5mm}
\begin{tcolorbox}[enhanced,nobeforeafter,width=7.1cm,height=2.1cm,valign=center,colback=csta!6,colframe=csta,arc=2.5mm,boxrule=0.9pt,left=3mm]{\Large\bfseries\color{csta}c)}\quad\begin{minipage}[c]{5.3cm}\raggedright\Large $+12{,}5\,\%$\end{minipage}\end{tcolorbox}\hfill\begin{tcolorbox}[enhanced,nobeforeafter,width=7.1cm,height=2.1cm,valign=center,colback=csta!6,colframe=csta,arc=2.5mm,boxrule=0.9pt,left=3mm]{\Large\bfseries\color{csta}d)}\quad\begin{minipage}[c]{5.3cm}\raggedright\Large $+10\,\%$\end{minipage}\end{tcolorbox}
\end{tcolorbox}
\newpage
\begin{tcolorbox}[enhanced,width=16cm,height=14cm,colback=white,colframe=csta,boxrule=2pt,arc=0mm,
left=6mm,right=6mm,top=0mm,bottom=4mm,sharp corners,
overlay={\fill[csta] (frame.north west) rectangle ([yshift=-1.25cm]frame.north east);
\node[anchor=west,text=white,font=\sffamily\bfseries\LARGE] at ([xshift=5mm,yshift=-0.62cm]frame.north west) {Question 23 B};
\node[anchor=east,text=white,font=\sffamily\itshape\large] at ([xshift=-5mm,yshift=-0.62cm]frame.north east) {Évolutions et statistiques};}]
\vspace*{1.75cm}
\begin{minipage}[t][6.9cm][c]{\linewidth}\LARGE\centering
\begin{minipage}{0.95\linewidth}\raggedright\linespread{1.25}\selectfont Un prix subit une baisse de $50\,\%$, puis une hausse de $50\,\%$. Quelle est son évolution globale~?\end{minipage}
\end{minipage}

\vfill
\begin{tcolorbox}[enhanced,nobeforeafter,width=7.1cm,height=2.1cm,valign=center,colback=csta!6,colframe=csta,arc=2.5mm,boxrule=0.9pt,left=3mm]{\Large\bfseries\color{csta}a)}\quad\begin{minipage}[c]{5.3cm}\raggedright\Large $-24\,\%$\end{minipage}\end{tcolorbox}\hfill\begin{tcolorbox}[enhanced,nobeforeafter,width=7.1cm,height=2.1cm,valign=center,colback=csta!6,colframe=csta,arc=2.5mm,boxrule=0.9pt,left=3mm]{\Large\bfseries\color{csta}b)}\quad\begin{minipage}[c]{5.3cm}\raggedright\Large $-27\,\%$\end{minipage}\end{tcolorbox}\par\vspace{2.5mm}
\begin{tcolorbox}[enhanced,nobeforeafter,width=7.1cm,height=2.1cm,valign=center,colback=csta!6,colframe=csta,arc=2.5mm,boxrule=0.9pt,left=3mm]{\Large\bfseries\color{csta}c)}\quad\begin{minipage}[c]{5.3cm}\raggedright\Large $0\,\%$\end{minipage}\end{tcolorbox}\hfill\begin{tcolorbox}[enhanced,nobeforeafter,width=7.1cm,height=2.1cm,valign=center,colback=csta!6,colframe=csta,arc=2.5mm,boxrule=0.9pt,left=3mm]{\Large\bfseries\color{csta}d)}\quad\begin{minipage}[c]{5.3cm}\raggedright\Large $-25\,\%$\end{minipage}\end{tcolorbox}
\end{tcolorbox}
\newpage
\begin{tcolorbox}[enhanced,width=16cm,height=14cm,colback=white,colframe=csta,boxrule=2pt,arc=0mm,
left=6mm,right=6mm,top=0mm,bottom=4mm,sharp corners,
overlay={\fill[csta] (frame.north west) rectangle ([yshift=-1.25cm]frame.north east);
\node[anchor=west,text=white,font=\sffamily\bfseries\LARGE] at ([xshift=5mm,yshift=-0.62cm]frame.north west) {Question 24 B};
\node[anchor=east,text=white,font=\sffamily\itshape\large] at ([xshift=-5mm,yshift=-0.62cm]frame.north east) {Évolutions et statistiques};}]
\vspace*{1.75cm}
\begin{minipage}[t][6.9cm][c]{\linewidth}\LARGE\centering
\begin{minipage}{0.95\linewidth}\raggedright\linespread{1.25}\selectfont Un prix subit une hausse de $30\,\%$, puis une baisse de $10\,\%$. Quelle est son évolution globale~?\end{minipage}
\end{minipage}

\vfill
\begin{tcolorbox}[enhanced,nobeforeafter,width=7.1cm,height=2.1cm,valign=center,colback=csta!6,colframe=csta,arc=2.5mm,boxrule=0.9pt,left=3mm]{\Large\bfseries\color{csta}a)}\quad\begin{minipage}[c]{5.3cm}\raggedright\Large $+18\,\%$\end{minipage}\end{tcolorbox}\hfill\begin{tcolorbox}[enhanced,nobeforeafter,width=7.1cm,height=2.1cm,valign=center,colback=csta!6,colframe=csta,arc=2.5mm,boxrule=0.9pt,left=3mm]{\Large\bfseries\color{csta}b)}\quad\begin{minipage}[c]{5.3cm}\raggedright\Large $+20\,\%$\end{minipage}\end{tcolorbox}\par\vspace{2.5mm}
\begin{tcolorbox}[enhanced,nobeforeafter,width=7.1cm,height=2.1cm,valign=center,colback=csta!6,colframe=csta,arc=2.5mm,boxrule=0.9pt,left=3mm]{\Large\bfseries\color{csta}c)}\quad\begin{minipage}[c]{5.3cm}\raggedright\Large $-3\,\%$\end{minipage}\end{tcolorbox}\hfill\begin{tcolorbox}[enhanced,nobeforeafter,width=7.1cm,height=2.1cm,valign=center,colback=csta!6,colframe=csta,arc=2.5mm,boxrule=0.9pt,left=3mm]{\Large\bfseries\color{csta}d)}\quad\begin{minipage}[c]{5.3cm}\raggedright\Large $+17\,\%$\end{minipage}\end{tcolorbox}
\end{tcolorbox}
\newpage
\begin{tcolorbox}[enhanced,width=16cm,height=14cm,colback=white,colframe=csta,boxrule=2pt,arc=0mm,
left=6mm,right=6mm,top=0mm,bottom=4mm,sharp corners,
overlay={\fill[csta] (frame.north west) rectangle ([yshift=-1.25cm]frame.north east);
\node[anchor=west,text=white,font=\sffamily\bfseries\LARGE] at ([xshift=5mm,yshift=-0.62cm]frame.north west) {Question 25 B};
\node[anchor=east,text=white,font=\sffamily\itshape\large] at ([xshift=-5mm,yshift=-0.62cm]frame.north east) {Évolutions et statistiques};}]
\vspace*{1.75cm}
\begin{minipage}[t][6.9cm][c]{\linewidth}\LARGE\centering
\begin{minipage}{0.95\linewidth}\raggedright\linespread{1.25}\selectfont Après une hausse de $25\,\%$, quelle évolution faut-il appliquer pour revenir au prix de départ~?\end{minipage}
\end{minipage}

\vfill
\begin{tcolorbox}[enhanced,nobeforeafter,width=7.1cm,height=2.1cm,valign=center,colback=csta!6,colframe=csta,arc=2.5mm,boxrule=0.9pt,left=3mm]{\Large\bfseries\color{csta}a)}\quad\begin{minipage}[c]{5.3cm}\raggedright\Large une baisse de $20\,\%$\end{minipage}\end{tcolorbox}\hfill\begin{tcolorbox}[enhanced,nobeforeafter,width=7.1cm,height=2.1cm,valign=center,colback=csta!6,colframe=csta,arc=2.5mm,boxrule=0.9pt,left=3mm]{\Large\bfseries\color{csta}b)}\quad\begin{minipage}[c]{5.3cm}\raggedright\Large une baisse de $25\,\%$\end{minipage}\end{tcolorbox}\par\vspace{2.5mm}
\begin{tcolorbox}[enhanced,nobeforeafter,width=7.1cm,height=2.1cm,valign=center,colback=csta!6,colframe=csta,arc=2.5mm,boxrule=0.9pt,left=3mm]{\Large\bfseries\color{csta}c)}\quad\begin{minipage}[c]{5.3cm}\raggedright\Large une baisse de $30\,\%$\end{minipage}\end{tcolorbox}\hfill\begin{tcolorbox}[enhanced,nobeforeafter,width=7.1cm,height=2.1cm,valign=center,colback=csta!6,colframe=csta,arc=2.5mm,boxrule=0.9pt,left=3mm]{\Large\bfseries\color{csta}d)}\quad\begin{minipage}[c]{5.3cm}\raggedright\Large une hausse de $20\,\%$\end{minipage}\end{tcolorbox}
\end{tcolorbox}
\newpage
\begin{tcolorbox}[enhanced,width=16cm,height=14cm,colback=white,colframe=csta,boxrule=2pt,arc=0mm,
left=6mm,right=6mm,top=0mm,bottom=4mm,sharp corners,
overlay={\fill[csta] (frame.north west) rectangle ([yshift=-1.25cm]frame.north east);
\node[anchor=west,text=white,font=\sffamily\bfseries\LARGE] at ([xshift=5mm,yshift=-0.62cm]frame.north west) {Question 26 B};
\node[anchor=east,text=white,font=\sffamily\itshape\large] at ([xshift=-5mm,yshift=-0.62cm]frame.north east) {Évolutions et statistiques};}]
\vspace*{1.75cm}
\begin{minipage}[t][6.9cm][c]{\linewidth}\LARGE\centering
\begin{minipage}{0.95\linewidth}\raggedright\linespread{1.25}\selectfont Après une baisse de $60\,\%$, quelle évolution faut-il appliquer pour revenir au prix de départ~?\end{minipage}
\end{minipage}

\vfill
\begin{tcolorbox}[enhanced,nobeforeafter,width=7.1cm,height=2.1cm,valign=center,colback=csta!6,colframe=csta,arc=2.5mm,boxrule=0.9pt,left=3mm]{\Large\bfseries\color{csta}a)}\quad\begin{minipage}[c]{5.3cm}\raggedright\Large une hausse de $60\,\%$\end{minipage}\end{tcolorbox}\hfill\begin{tcolorbox}[enhanced,nobeforeafter,width=7.1cm,height=2.1cm,valign=center,colback=csta!6,colframe=csta,arc=2.5mm,boxrule=0.9pt,left=3mm]{\Large\bfseries\color{csta}b)}\quad\begin{minipage}[c]{5.3cm}\raggedright\Large une baisse de $150\,\%$\end{minipage}\end{tcolorbox}\par\vspace{2.5mm}
\begin{tcolorbox}[enhanced,nobeforeafter,width=7.1cm,height=2.1cm,valign=center,colback=csta!6,colframe=csta,arc=2.5mm,boxrule=0.9pt,left=3mm]{\Large\bfseries\color{csta}c)}\quad\begin{minipage}[c]{5.3cm}\raggedright\Large une hausse de $160\,\%$\end{minipage}\end{tcolorbox}\hfill\begin{tcolorbox}[enhanced,nobeforeafter,width=7.1cm,height=2.1cm,valign=center,colback=csta!6,colframe=csta,arc=2.5mm,boxrule=0.9pt,left=3mm]{\Large\bfseries\color{csta}d)}\quad\begin{minipage}[c]{5.3cm}\raggedright\Large une hausse de $150\,\%$\end{minipage}\end{tcolorbox}
\end{tcolorbox}
\newpage
\begin{tcolorbox}[enhanced,width=16cm,height=14cm,colback=white,colframe=csta,boxrule=2pt,arc=0mm,
left=6mm,right=6mm,top=0mm,bottom=4mm,sharp corners,
overlay={\fill[csta] (frame.north west) rectangle ([yshift=-1.25cm]frame.north east);
\node[anchor=west,text=white,font=\sffamily\bfseries\LARGE] at ([xshift=5mm,yshift=-0.62cm]frame.north west) {Question 27 B};
\node[anchor=east,text=white,font=\sffamily\itshape\large] at ([xshift=-5mm,yshift=-0.62cm]frame.north east) {Évolutions et statistiques};}]
\vspace*{1.75cm}
\begin{minipage}[t][6.9cm][c]{\linewidth}\Large\centering
\begin{minipage}{0.95\linewidth}\raggedright\linespread{1.25}\selectfont Dans un lycée, $30\,\%$ des élèves sont internes et $20\,\%$ des internes viennent d'une autre commune. Quelle proportion des élèves du lycée sont des internes venant d'une autre commune~?\end{minipage}
\end{minipage}

\vfill
\begin{tcolorbox}[enhanced,nobeforeafter,width=7.1cm,height=2.1cm,valign=center,colback=csta!6,colframe=csta,arc=2.5mm,boxrule=0.9pt,left=3mm]{\Large\bfseries\color{csta}a)}\quad\begin{minipage}[c]{5.3cm}\raggedright\Large $25\,\%$\end{minipage}\end{tcolorbox}\hfill\begin{tcolorbox}[enhanced,nobeforeafter,width=7.1cm,height=2.1cm,valign=center,colback=csta!6,colframe=csta,arc=2.5mm,boxrule=0.9pt,left=3mm]{\Large\bfseries\color{csta}b)}\quad\begin{minipage}[c]{5.3cm}\raggedright\Large $6\,\%$\end{minipage}\end{tcolorbox}\par\vspace{2.5mm}
\begin{tcolorbox}[enhanced,nobeforeafter,width=7.1cm,height=2.1cm,valign=center,colback=csta!6,colframe=csta,arc=2.5mm,boxrule=0.9pt,left=3mm]{\Large\bfseries\color{csta}c)}\quad\begin{minipage}[c]{5.3cm}\raggedright\Large $50\,\%$\end{minipage}\end{tcolorbox}\hfill\begin{tcolorbox}[enhanced,nobeforeafter,width=7.1cm,height=2.1cm,valign=center,colback=csta!6,colframe=csta,arc=2.5mm,boxrule=0.9pt,left=3mm]{\Large\bfseries\color{csta}d)}\quad\begin{minipage}[c]{5.3cm}\raggedright\Large $10\,\%$\end{minipage}\end{tcolorbox}
\end{tcolorbox}
\newpage
\begin{tcolorbox}[enhanced,width=16cm,height=14cm,colback=white,colframe=csta,boxrule=2pt,arc=0mm,
left=6mm,right=6mm,top=0mm,bottom=4mm,sharp corners,
overlay={\fill[csta] (frame.north west) rectangle ([yshift=-1.25cm]frame.north east);
\node[anchor=west,text=white,font=\sffamily\bfseries\LARGE] at ([xshift=5mm,yshift=-0.62cm]frame.north west) {Question 28 B};
\node[anchor=east,text=white,font=\sffamily\itshape\large] at ([xshift=-5mm,yshift=-0.62cm]frame.north east) {Évolutions et statistiques};}]
\vspace*{1.75cm}
\begin{minipage}[t][6.9cm][c]{\linewidth}\Large\centering
\begin{minipage}{0.95\linewidth}\raggedright\linespread{1.25}\selectfont Dans un lycée, $80\,\%$ des élèves sont internes et $15\,\%$ des internes viennent d'une autre commune. Quelle proportion des élèves du lycée sont des internes venant d'une autre commune~?\end{minipage}
\end{minipage}

\vfill
\begin{tcolorbox}[enhanced,nobeforeafter,width=7.1cm,height=2.1cm,valign=center,colback=csta!6,colframe=csta,arc=2.5mm,boxrule=0.9pt,left=3mm]{\Large\bfseries\color{csta}a)}\quad\begin{minipage}[c]{5.3cm}\raggedright\Large $47{,}5\,\%$\end{minipage}\end{tcolorbox}\hfill\begin{tcolorbox}[enhanced,nobeforeafter,width=7.1cm,height=2.1cm,valign=center,colback=csta!6,colframe=csta,arc=2.5mm,boxrule=0.9pt,left=3mm]{\Large\bfseries\color{csta}b)}\quad\begin{minipage}[c]{5.3cm}\raggedright\Large $65\,\%$\end{minipage}\end{tcolorbox}\par\vspace{2.5mm}
\begin{tcolorbox}[enhanced,nobeforeafter,width=7.1cm,height=2.1cm,valign=center,colback=csta!6,colframe=csta,arc=2.5mm,boxrule=0.9pt,left=3mm]{\Large\bfseries\color{csta}c)}\quad\begin{minipage}[c]{5.3cm}\raggedright\Large $95\,\%$\end{minipage}\end{tcolorbox}\hfill\begin{tcolorbox}[enhanced,nobeforeafter,width=7.1cm,height=2.1cm,valign=center,colback=csta!6,colframe=csta,arc=2.5mm,boxrule=0.9pt,left=3mm]{\Large\bfseries\color{csta}d)}\quad\begin{minipage}[c]{5.3cm}\raggedright\Large $12\,\%$\end{minipage}\end{tcolorbox}
\end{tcolorbox}
\newpage
\begin{tcolorbox}[enhanced,width=16cm,height=14cm,colback=white,colframe=csta,boxrule=2pt,arc=0mm,
left=6mm,right=6mm,top=0mm,bottom=4mm,sharp corners,
overlay={\fill[csta] (frame.north west) rectangle ([yshift=-1.25cm]frame.north east);
\node[anchor=west,text=white,font=\sffamily\bfseries\LARGE] at ([xshift=5mm,yshift=-0.62cm]frame.north west) {Question 29 B};
\node[anchor=east,text=white,font=\sffamily\itshape\large] at ([xshift=-5mm,yshift=-0.62cm]frame.north east) {Évolutions et statistiques};}]
\vspace*{1.75cm}
\begin{minipage}[t][6.9cm][c]{\linewidth}\Large\centering
\begin{minipage}{0.95\linewidth}\raggedright\linespread{1.25}\selectfont Quelle est la moyenne de cette série statistique~? $\begin{array}{|l|c|c|c|}\hline \text{Valeur} & 5 & 10 & 20\\\hline \text{Effectif} & 1 & 3 & 1\\\hline\end{array}$\end{minipage}
\end{minipage}

\vfill
\begin{tcolorbox}[enhanced,nobeforeafter,width=7.1cm,height=2.1cm,valign=center,colback=csta!6,colframe=csta,arc=2.5mm,boxrule=0.9pt,left=3mm]{\Large\bfseries\color{csta}a)}\quad\begin{minipage}[c]{5.3cm}\raggedright\Large $11$\end{minipage}\end{tcolorbox}\hfill\begin{tcolorbox}[enhanced,nobeforeafter,width=7.1cm,height=2.1cm,valign=center,colback=csta!6,colframe=csta,arc=2.5mm,boxrule=0.9pt,left=3mm]{\Large\bfseries\color{csta}b)}\quad\begin{minipage}[c]{5.3cm}\raggedright\Large $\dfrac{55}{3}$\end{minipage}\end{tcolorbox}\par\vspace{2.5mm}
\begin{tcolorbox}[enhanced,nobeforeafter,width=7.1cm,height=2.1cm,valign=center,colback=csta!6,colframe=csta,arc=2.5mm,boxrule=0.9pt,left=3mm]{\Large\bfseries\color{csta}c)}\quad\begin{minipage}[c]{5.3cm}\raggedright\Large $12$\end{minipage}\end{tcolorbox}\hfill\begin{tcolorbox}[enhanced,nobeforeafter,width=7.1cm,height=2.1cm,valign=center,colback=csta!6,colframe=csta,arc=2.5mm,boxrule=0.9pt,left=3mm]{\Large\bfseries\color{csta}d)}\quad\begin{minipage}[c]{5.3cm}\raggedright\Large $\dfrac{35}{3}$\end{minipage}\end{tcolorbox}
\end{tcolorbox}
\newpage
\begin{tcolorbox}[enhanced,width=16cm,height=14cm,colback=white,colframe=csta,boxrule=2pt,arc=0mm,
left=6mm,right=6mm,top=0mm,bottom=4mm,sharp corners,
overlay={\fill[csta] (frame.north west) rectangle ([yshift=-1.25cm]frame.north east);
\node[anchor=west,text=white,font=\sffamily\bfseries\LARGE] at ([xshift=5mm,yshift=-0.62cm]frame.north west) {Question 30 B};
\node[anchor=east,text=white,font=\sffamily\itshape\large] at ([xshift=-5mm,yshift=-0.62cm]frame.north east) {Évolutions et statistiques};}]
\vspace*{1.75cm}
\begin{minipage}[t][6.9cm][c]{\linewidth}\Large\centering
\begin{minipage}{0.95\linewidth}\raggedright\linespread{1.25}\selectfont Quelle est la moyenne de cette série statistique~? $\begin{array}{|l|c|c|c|}\hline \text{Valeur} & 0 & 1 & 2\\\hline \text{Effectif} & 2 & 5 & 3\\\hline\end{array}$\end{minipage}
\end{minipage}

\vfill
\begin{tcolorbox}[enhanced,nobeforeafter,width=7.1cm,height=2.1cm,valign=center,colback=csta!6,colframe=csta,arc=2.5mm,boxrule=0.9pt,left=3mm]{\Large\bfseries\color{csta}a)}\quad\begin{minipage}[c]{5.3cm}\raggedright\Large $1{,}1$\end{minipage}\end{tcolorbox}\hfill\begin{tcolorbox}[enhanced,nobeforeafter,width=7.1cm,height=2.1cm,valign=center,colback=csta!6,colframe=csta,arc=2.5mm,boxrule=0.9pt,left=3mm]{\Large\bfseries\color{csta}b)}\quad\begin{minipage}[c]{5.3cm}\raggedright\Large $\dfrac{11}{3}$\end{minipage}\end{tcolorbox}\par\vspace{2.5mm}
\begin{tcolorbox}[enhanced,nobeforeafter,width=7.1cm,height=2.1cm,valign=center,colback=csta!6,colframe=csta,arc=2.5mm,boxrule=0.9pt,left=3mm]{\Large\bfseries\color{csta}c)}\quad\begin{minipage}[c]{5.3cm}\raggedright\Large $2{,}1$\end{minipage}\end{tcolorbox}\hfill\begin{tcolorbox}[enhanced,nobeforeafter,width=7.1cm,height=2.1cm,valign=center,colback=csta!6,colframe=csta,arc=2.5mm,boxrule=0.9pt,left=3mm]{\Large\bfseries\color{csta}d)}\quad\begin{minipage}[c]{5.3cm}\raggedright\Large $1$\end{minipage}\end{tcolorbox}
\end{tcolorbox}
\newpage
\begin{tcolorbox}[enhanced,width=16cm,height=14cm,colback=white,colframe=csta,boxrule=2pt,arc=0mm,
left=6mm,right=6mm,top=0mm,bottom=4mm,sharp corners,
overlay={\fill[csta] (frame.north west) rectangle ([yshift=-1.25cm]frame.north east);
\node[anchor=west,text=white,font=\sffamily\bfseries\LARGE] at ([xshift=5mm,yshift=-0.62cm]frame.north west) {Question 31 B};
\node[anchor=east,text=white,font=\sffamily\itshape\large] at ([xshift=-5mm,yshift=-0.62cm]frame.north east) {Évolutions et statistiques};}]
\vspace*{1.75cm}
\begin{minipage}[t][6.9cm][c]{\linewidth}\Large\centering
\begin{minipage}{0.95\linewidth}\raggedright\linespread{1.25}\selectfont Voici une série rangée dans l'ordre croissant~: $10\ ;\ 12\ ;\ 13\ ;\ 15\ ;\ 16\ ;\ 18\ ;\ 21\ ;\ 22\ ;\ 25\ ;\ 30$. Quel est son premier quartile $Q_1$~?\end{minipage}
\end{minipage}

\vfill
\begin{tcolorbox}[enhanced,nobeforeafter,width=7.1cm,height=2.1cm,valign=center,colback=csta!6,colframe=csta,arc=2.5mm,boxrule=0.9pt,left=3mm]{\Large\bfseries\color{csta}a)}\quad\begin{minipage}[c]{5.3cm}\raggedright\Large $13$\end{minipage}\end{tcolorbox}\hfill\begin{tcolorbox}[enhanced,nobeforeafter,width=7.1cm,height=2.1cm,valign=center,colback=csta!6,colframe=csta,arc=2.5mm,boxrule=0.9pt,left=3mm]{\Large\bfseries\color{csta}b)}\quad\begin{minipage}[c]{5.3cm}\raggedright\Large $22$\end{minipage}\end{tcolorbox}\par\vspace{2.5mm}
\begin{tcolorbox}[enhanced,nobeforeafter,width=7.1cm,height=2.1cm,valign=center,colback=csta!6,colframe=csta,arc=2.5mm,boxrule=0.9pt,left=3mm]{\Large\bfseries\color{csta}c)}\quad\begin{minipage}[c]{5.3cm}\raggedright\Large $17$\end{minipage}\end{tcolorbox}\hfill\begin{tcolorbox}[enhanced,nobeforeafter,width=7.1cm,height=2.1cm,valign=center,colback=csta!6,colframe=csta,arc=2.5mm,boxrule=0.9pt,left=3mm]{\Large\bfseries\color{csta}d)}\quad\begin{minipage}[c]{5.3cm}\raggedright\Large $15$\end{minipage}\end{tcolorbox}
\end{tcolorbox}
\newpage
\begin{tcolorbox}[enhanced,width=16cm,height=14cm,colback=white,colframe=csta,boxrule=2pt,arc=0mm,
left=6mm,right=6mm,top=0mm,bottom=4mm,sharp corners,
overlay={\fill[csta] (frame.north west) rectangle ([yshift=-1.25cm]frame.north east);
\node[anchor=west,text=white,font=\sffamily\bfseries\LARGE] at ([xshift=5mm,yshift=-0.62cm]frame.north west) {Question 32 B};
\node[anchor=east,text=white,font=\sffamily\itshape\large] at ([xshift=-5mm,yshift=-0.62cm]frame.north east) {Évolutions et statistiques};}]
\vspace*{1.75cm}
\begin{minipage}[t][6.9cm][c]{\linewidth}\LARGE\centering
\begin{minipage}{0.95\linewidth}\raggedright\linespread{1.25}\selectfont Voici une série rangée dans l'ordre croissant~: $1\ ;\ 1\ ;\ 2\ ;\ 3\ ;\ 5\ ;\ 8\ ;\ 13\ ;\ 21$. Quel est son premier quartile $Q_1$~?\end{minipage}
\end{minipage}

\vfill
\begin{tcolorbox}[enhanced,nobeforeafter,width=7.1cm,height=2.1cm,valign=center,colback=csta!6,colframe=csta,arc=2.5mm,boxrule=0.9pt,left=3mm]{\Large\bfseries\color{csta}a)}\quad\begin{minipage}[c]{5.3cm}\raggedright\Large $2$\end{minipage}\end{tcolorbox}\hfill\begin{tcolorbox}[enhanced,nobeforeafter,width=7.1cm,height=2.1cm,valign=center,colback=csta!6,colframe=csta,arc=2.5mm,boxrule=0.9pt,left=3mm]{\Large\bfseries\color{csta}b)}\quad\begin{minipage}[c]{5.3cm}\raggedright\Large $1$\end{minipage}\end{tcolorbox}\par\vspace{2.5mm}
\begin{tcolorbox}[enhanced,nobeforeafter,width=7.1cm,height=2.1cm,valign=center,colback=csta!6,colframe=csta,arc=2.5mm,boxrule=0.9pt,left=3mm]{\Large\bfseries\color{csta}c)}\quad\begin{minipage}[c]{5.3cm}\raggedright\Large $8$\end{minipage}\end{tcolorbox}\hfill\begin{tcolorbox}[enhanced,nobeforeafter,width=7.1cm,height=2.1cm,valign=center,colback=csta!6,colframe=csta,arc=2.5mm,boxrule=0.9pt,left=3mm]{\Large\bfseries\color{csta}d)}\quad\begin{minipage}[c]{5.3cm}\raggedright\Large $4$\end{minipage}\end{tcolorbox}
\end{tcolorbox}
\newpage
\begin{tcolorbox}[enhanced,width=16cm,height=14cm,colback=white,colframe=cfon,boxrule=2pt,arc=0mm,
left=6mm,right=6mm,top=0mm,bottom=4mm,sharp corners,
overlay={\fill[cfon] (frame.north west) rectangle ([yshift=-1.25cm]frame.north east);
\node[anchor=west,text=white,font=\sffamily\bfseries\LARGE] at ([xshift=5mm,yshift=-0.62cm]frame.north west) {Question 33 B};
\node[anchor=east,text=white,font=\sffamily\itshape\large] at ([xshift=-5mm,yshift=-0.62cm]frame.north east) {Fonctions};}]
\vspace*{1.75cm}
\begin{minipage}[t][6.9cm][c]{\linewidth}\LARGE\centering
\begin{minipage}{0.95\linewidth}\raggedright\linespread{1.25}\selectfont Soit $f(x)=-x^2+2x+6$. Que vaut $f(-3)$~?\end{minipage}
\end{minipage}

\vfill
\begin{tcolorbox}[enhanced,nobeforeafter,width=7.1cm,height=2.1cm,valign=center,colback=cfon!6,colframe=cfon,arc=2.5mm,boxrule=0.9pt,left=3mm]{\Large\bfseries\color{cfon}a)}\quad\begin{minipage}[c]{5.3cm}\raggedright\Large $6$\end{minipage}\end{tcolorbox}\hfill\begin{tcolorbox}[enhanced,nobeforeafter,width=7.1cm,height=2.1cm,valign=center,colback=cfon!6,colframe=cfon,arc=2.5mm,boxrule=0.9pt,left=3mm]{\Large\bfseries\color{cfon}b)}\quad\begin{minipage}[c]{5.3cm}\raggedright\Large $9$\end{minipage}\end{tcolorbox}\par\vspace{2.5mm}
\begin{tcolorbox}[enhanced,nobeforeafter,width=7.1cm,height=2.1cm,valign=center,colback=cfon!6,colframe=cfon,arc=2.5mm,boxrule=0.9pt,left=3mm]{\Large\bfseries\color{cfon}c)}\quad\begin{minipage}[c]{5.3cm}\raggedright\Large $-7$\end{minipage}\end{tcolorbox}\hfill\begin{tcolorbox}[enhanced,nobeforeafter,width=7.1cm,height=2.1cm,valign=center,colback=cfon!6,colframe=cfon,arc=2.5mm,boxrule=0.9pt,left=3mm]{\Large\bfseries\color{cfon}d)}\quad\begin{minipage}[c]{5.3cm}\raggedright\Large $-9$\end{minipage}\end{tcolorbox}
\end{tcolorbox}
\newpage
\begin{tcolorbox}[enhanced,width=16cm,height=14cm,colback=white,colframe=cfon,boxrule=2pt,arc=0mm,
left=6mm,right=6mm,top=0mm,bottom=4mm,sharp corners,
overlay={\fill[cfon] (frame.north west) rectangle ([yshift=-1.25cm]frame.north east);
\node[anchor=west,text=white,font=\sffamily\bfseries\LARGE] at ([xshift=5mm,yshift=-0.62cm]frame.north west) {Question 34 B};
\node[anchor=east,text=white,font=\sffamily\itshape\large] at ([xshift=-5mm,yshift=-0.62cm]frame.north east) {Fonctions};}]
\vspace*{1.75cm}
\begin{minipage}[t][6.9cm][c]{\linewidth}\LARGE\centering
\begin{minipage}{0.95\linewidth}\raggedright\linespread{1.25}\selectfont Soit $f(x)=3x^2-7$. Que vaut $f(2)$~?\end{minipage}
\end{minipage}

\vfill
\begin{tcolorbox}[enhanced,nobeforeafter,width=7.1cm,height=2.1cm,valign=center,colback=cfon!6,colframe=cfon,arc=2.5mm,boxrule=0.9pt,left=3mm]{\Large\bfseries\color{cfon}a)}\quad\begin{minipage}[c]{5.3cm}\raggedright\Large $5$\end{minipage}\end{tcolorbox}\hfill\begin{tcolorbox}[enhanced,nobeforeafter,width=7.1cm,height=2.1cm,valign=center,colback=cfon!6,colframe=cfon,arc=2.5mm,boxrule=0.9pt,left=3mm]{\Large\bfseries\color{cfon}b)}\quad\begin{minipage}[c]{5.3cm}\raggedright\Large $-19$\end{minipage}\end{tcolorbox}\par\vspace{2.5mm}
\begin{tcolorbox}[enhanced,nobeforeafter,width=7.1cm,height=2.1cm,valign=center,colback=cfon!6,colframe=cfon,arc=2.5mm,boxrule=0.9pt,left=3mm]{\Large\bfseries\color{cfon}c)}\quad\begin{minipage}[c]{5.3cm}\raggedright\Large $7$\end{minipage}\end{tcolorbox}\hfill\begin{tcolorbox}[enhanced,nobeforeafter,width=7.1cm,height=2.1cm,valign=center,colback=cfon!6,colframe=cfon,arc=2.5mm,boxrule=0.9pt,left=3mm]{\Large\bfseries\color{cfon}d)}\quad\begin{minipage}[c]{5.3cm}\raggedright\Large $-5$\end{minipage}\end{tcolorbox}
\end{tcolorbox}
\newpage
\begin{tcolorbox}[enhanced,width=16cm,height=14cm,colback=white,colframe=cfon,boxrule=2pt,arc=0mm,
left=6mm,right=6mm,top=0mm,bottom=4mm,sharp corners,
overlay={\fill[cfon] (frame.north west) rectangle ([yshift=-1.25cm]frame.north east);
\node[anchor=west,text=white,font=\sffamily\bfseries\LARGE] at ([xshift=5mm,yshift=-0.62cm]frame.north west) {Question 35 B};
\node[anchor=east,text=white,font=\sffamily\itshape\large] at ([xshift=-5mm,yshift=-0.62cm]frame.north east) {Fonctions};}]
\vspace*{1.75cm}
\begin{minipage}[t][6.9cm][c]{\linewidth}\LARGE\centering
\begin{minipage}{0.95\linewidth}\raggedright\linespread{1.25}\selectfont Soit $f(x)=4x-7$. Quel est l'antécédent de $5$ par $f$~?\end{minipage}
\end{minipage}

\vfill
\begin{tcolorbox}[enhanced,nobeforeafter,width=7.1cm,height=2.1cm,valign=center,colback=cfon!6,colframe=cfon,arc=2.5mm,boxrule=0.9pt,left=3mm]{\Large\bfseries\color{cfon}a)}\quad\begin{minipage}[c]{5.3cm}\raggedright\Large $13$\end{minipage}\end{tcolorbox}\hfill\begin{tcolorbox}[enhanced,nobeforeafter,width=7.1cm,height=2.1cm,valign=center,colback=cfon!6,colframe=cfon,arc=2.5mm,boxrule=0.9pt,left=3mm]{\Large\bfseries\color{cfon}b)}\quad\begin{minipage}[c]{5.3cm}\raggedright\Large $-\dfrac{1}{2}$\end{minipage}\end{tcolorbox}\par\vspace{2.5mm}
\begin{tcolorbox}[enhanced,nobeforeafter,width=7.1cm,height=2.1cm,valign=center,colback=cfon!6,colframe=cfon,arc=2.5mm,boxrule=0.9pt,left=3mm]{\Large\bfseries\color{cfon}c)}\quad\begin{minipage}[c]{5.3cm}\raggedright\Large $48$\end{minipage}\end{tcolorbox}\hfill\begin{tcolorbox}[enhanced,nobeforeafter,width=7.1cm,height=2.1cm,valign=center,colback=cfon!6,colframe=cfon,arc=2.5mm,boxrule=0.9pt,left=3mm]{\Large\bfseries\color{cfon}d)}\quad\begin{minipage}[c]{5.3cm}\raggedright\Large $3$\end{minipage}\end{tcolorbox}
\end{tcolorbox}
\newpage
\begin{tcolorbox}[enhanced,width=16cm,height=14cm,colback=white,colframe=cfon,boxrule=2pt,arc=0mm,
left=6mm,right=6mm,top=0mm,bottom=4mm,sharp corners,
overlay={\fill[cfon] (frame.north west) rectangle ([yshift=-1.25cm]frame.north east);
\node[anchor=west,text=white,font=\sffamily\bfseries\LARGE] at ([xshift=5mm,yshift=-0.62cm]frame.north west) {Question 36 B};
\node[anchor=east,text=white,font=\sffamily\itshape\large] at ([xshift=-5mm,yshift=-0.62cm]frame.north east) {Fonctions};}]
\vspace*{1.75cm}
\begin{minipage}[t][6.9cm][c]{\linewidth}\LARGE\centering
\begin{minipage}{0.95\linewidth}\raggedright\linespread{1.25}\selectfont Soit $f(x)=5x+2$. Quel est l'antécédent de $-8$ par $f$~?\end{minipage}
\end{minipage}

\vfill
\begin{tcolorbox}[enhanced,nobeforeafter,width=7.1cm,height=2.1cm,valign=center,colback=cfon!6,colframe=cfon,arc=2.5mm,boxrule=0.9pt,left=3mm]{\Large\bfseries\color{cfon}a)}\quad\begin{minipage}[c]{5.3cm}\raggedright\Large $-38$\end{minipage}\end{tcolorbox}\hfill\begin{tcolorbox}[enhanced,nobeforeafter,width=7.1cm,height=2.1cm,valign=center,colback=cfon!6,colframe=cfon,arc=2.5mm,boxrule=0.9pt,left=3mm]{\Large\bfseries\color{cfon}b)}\quad\begin{minipage}[c]{5.3cm}\raggedright\Large $-50$\end{minipage}\end{tcolorbox}\par\vspace{2.5mm}
\begin{tcolorbox}[enhanced,nobeforeafter,width=7.1cm,height=2.1cm,valign=center,colback=cfon!6,colframe=cfon,arc=2.5mm,boxrule=0.9pt,left=3mm]{\Large\bfseries\color{cfon}c)}\quad\begin{minipage}[c]{5.3cm}\raggedright\Large $-2$\end{minipage}\end{tcolorbox}\hfill\begin{tcolorbox}[enhanced,nobeforeafter,width=7.1cm,height=2.1cm,valign=center,colback=cfon!6,colframe=cfon,arc=2.5mm,boxrule=0.9pt,left=3mm]{\Large\bfseries\color{cfon}d)}\quad\begin{minipage}[c]{5.3cm}\raggedright\Large $-\dfrac{6}{5}$\end{minipage}\end{tcolorbox}
\end{tcolorbox}
\newpage
\begin{tcolorbox}[enhanced,width=16cm,height=14cm,colback=white,colframe=cfon,boxrule=2pt,arc=0mm,
left=6mm,right=6mm,top=0mm,bottom=4mm,sharp corners,
overlay={\fill[cfon] (frame.north west) rectangle ([yshift=-1.25cm]frame.north east);
\node[anchor=west,text=white,font=\sffamily\bfseries\LARGE] at ([xshift=5mm,yshift=-0.62cm]frame.north west) {Question 37 B};
\node[anchor=east,text=white,font=\sffamily\itshape\large] at ([xshift=-5mm,yshift=-0.62cm]frame.north east) {Fonctions};}]
\vspace*{1.75cm}
\begin{minipage}[t][6.9cm][c]{\linewidth}\LARGE\centering
\begin{minipage}{0.95\linewidth}\raggedright\linespread{1.25}\selectfont Quel est le coefficient directeur de la droite $(AB)$, avec $A(0\,;\,1)$ et $B(4\,;\,3)$~?\end{minipage}
\end{minipage}

\vfill
\begin{tcolorbox}[enhanced,nobeforeafter,width=7.1cm,height=2.1cm,valign=center,colback=cfon!6,colframe=cfon,arc=2.5mm,boxrule=0.9pt,left=3mm]{\Large\bfseries\color{cfon}a)}\quad\begin{minipage}[c]{5.3cm}\raggedright\Large $1$\end{minipage}\end{tcolorbox}\hfill\begin{tcolorbox}[enhanced,nobeforeafter,width=7.1cm,height=2.1cm,valign=center,colback=cfon!6,colframe=cfon,arc=2.5mm,boxrule=0.9pt,left=3mm]{\Large\bfseries\color{cfon}b)}\quad\begin{minipage}[c]{5.3cm}\raggedright\Large $-\dfrac{1}{2}$\end{minipage}\end{tcolorbox}\par\vspace{2.5mm}
\begin{tcolorbox}[enhanced,nobeforeafter,width=7.1cm,height=2.1cm,valign=center,colback=cfon!6,colframe=cfon,arc=2.5mm,boxrule=0.9pt,left=3mm]{\Large\bfseries\color{cfon}c)}\quad\begin{minipage}[c]{5.3cm}\raggedright\Large $2$\end{minipage}\end{tcolorbox}\hfill\begin{tcolorbox}[enhanced,nobeforeafter,width=7.1cm,height=2.1cm,valign=center,colback=cfon!6,colframe=cfon,arc=2.5mm,boxrule=0.9pt,left=3mm]{\Large\bfseries\color{cfon}d)}\quad\begin{minipage}[c]{5.3cm}\raggedright\Large $\dfrac{1}{2}$\end{minipage}\end{tcolorbox}
\end{tcolorbox}
\newpage
\begin{tcolorbox}[enhanced,width=16cm,height=14cm,colback=white,colframe=cfon,boxrule=2pt,arc=0mm,
left=6mm,right=6mm,top=0mm,bottom=4mm,sharp corners,
overlay={\fill[cfon] (frame.north west) rectangle ([yshift=-1.25cm]frame.north east);
\node[anchor=west,text=white,font=\sffamily\bfseries\LARGE] at ([xshift=5mm,yshift=-0.62cm]frame.north west) {Question 38 B};
\node[anchor=east,text=white,font=\sffamily\itshape\large] at ([xshift=-5mm,yshift=-0.62cm]frame.north east) {Fonctions};}]
\vspace*{1.75cm}
\begin{minipage}[t][6.9cm][c]{\linewidth}\LARGE\centering
\begin{minipage}{0.95\linewidth}\raggedright\linespread{1.25}\selectfont Quel est le coefficient directeur de la droite $(AB)$, avec $A(1\,;\,-4)$ et $B(5\,;\,6)$~?\end{minipage}
\end{minipage}

\vfill
\begin{tcolorbox}[enhanced,nobeforeafter,width=7.1cm,height=2.1cm,valign=center,colback=cfon!6,colframe=cfon,arc=2.5mm,boxrule=0.9pt,left=3mm]{\Large\bfseries\color{cfon}a)}\quad\begin{minipage}[c]{5.3cm}\raggedright\Large $-\dfrac{5}{2}$\end{minipage}\end{tcolorbox}\hfill\begin{tcolorbox}[enhanced,nobeforeafter,width=7.1cm,height=2.1cm,valign=center,colback=cfon!6,colframe=cfon,arc=2.5mm,boxrule=0.9pt,left=3mm]{\Large\bfseries\color{cfon}b)}\quad\begin{minipage}[c]{5.3cm}\raggedright\Large $\dfrac{1}{3}$\end{minipage}\end{tcolorbox}\par\vspace{2.5mm}
\begin{tcolorbox}[enhanced,nobeforeafter,width=7.1cm,height=2.1cm,valign=center,colback=cfon!6,colframe=cfon,arc=2.5mm,boxrule=0.9pt,left=3mm]{\Large\bfseries\color{cfon}c)}\quad\begin{minipage}[c]{5.3cm}\raggedright\Large $\dfrac{5}{2}$\end{minipage}\end{tcolorbox}\hfill\begin{tcolorbox}[enhanced,nobeforeafter,width=7.1cm,height=2.1cm,valign=center,colback=cfon!6,colframe=cfon,arc=2.5mm,boxrule=0.9pt,left=3mm]{\Large\bfseries\color{cfon}d)}\quad\begin{minipage}[c]{5.3cm}\raggedright\Large $\dfrac{2}{5}$\end{minipage}\end{tcolorbox}
\end{tcolorbox}
\newpage
\begin{tcolorbox}[enhanced,width=16cm,height=14cm,colback=white,colframe=cfon,boxrule=2pt,arc=0mm,
left=6mm,right=6mm,top=0mm,bottom=4mm,sharp corners,
overlay={\fill[cfon] (frame.north west) rectangle ([yshift=-1.25cm]frame.north east);
\node[anchor=west,text=white,font=\sffamily\bfseries\LARGE] at ([xshift=5mm,yshift=-0.62cm]frame.north west) {Question 39 B};
\node[anchor=east,text=white,font=\sffamily\itshape\large] at ([xshift=-5mm,yshift=-0.62cm]frame.north east) {Fonctions};}]
\vspace*{1.75cm}
\begin{minipage}[t][6.9cm][c]{\linewidth}\LARGE\centering
\begin{minipage}{0.95\linewidth}\raggedright\linespread{1.25}\selectfont Quelle est l'équation réduite de la droite de coefficient directeur $\dfrac{1}{2}$ passant par $A(4\,;\,0)$~?\end{minipage}
\end{minipage}

\vfill
\begin{tcolorbox}[enhanced,nobeforeafter,width=7.1cm,height=2.1cm,valign=center,colback=cfon!6,colframe=cfon,arc=2.5mm,boxrule=0.9pt,left=3mm]{\Large\bfseries\color{cfon}a)}\quad\begin{minipage}[c]{5.3cm}\raggedright\Large $y=\dfrac{1}{2}x$\end{minipage}\end{tcolorbox}\hfill\begin{tcolorbox}[enhanced,nobeforeafter,width=7.1cm,height=2.1cm,valign=center,colback=cfon!6,colframe=cfon,arc=2.5mm,boxrule=0.9pt,left=3mm]{\Large\bfseries\color{cfon}b)}\quad\begin{minipage}[c]{5.3cm}\raggedright\Large $y=\dfrac{1}{2}x-2$\end{minipage}\end{tcolorbox}\par\vspace{2.5mm}
\begin{tcolorbox}[enhanced,nobeforeafter,width=7.1cm,height=2.1cm,valign=center,colback=cfon!6,colframe=cfon,arc=2.5mm,boxrule=0.9pt,left=3mm]{\Large\bfseries\color{cfon}c)}\quad\begin{minipage}[c]{5.3cm}\raggedright\Large $y=4x-2$\end{minipage}\end{tcolorbox}\hfill\begin{tcolorbox}[enhanced,nobeforeafter,width=7.1cm,height=2.1cm,valign=center,colback=cfon!6,colframe=cfon,arc=2.5mm,boxrule=0.9pt,left=3mm]{\Large\bfseries\color{cfon}d)}\quad\begin{minipage}[c]{5.3cm}\raggedright\Large $y=\dfrac{1}{2}x+2$\end{minipage}\end{tcolorbox}
\end{tcolorbox}
\newpage
\begin{tcolorbox}[enhanced,width=16cm,height=14cm,colback=white,colframe=cfon,boxrule=2pt,arc=0mm,
left=6mm,right=6mm,top=0mm,bottom=4mm,sharp corners,
overlay={\fill[cfon] (frame.north west) rectangle ([yshift=-1.25cm]frame.north east);
\node[anchor=west,text=white,font=\sffamily\bfseries\LARGE] at ([xshift=5mm,yshift=-0.62cm]frame.north west) {Question 40 B};
\node[anchor=east,text=white,font=\sffamily\itshape\large] at ([xshift=-5mm,yshift=-0.62cm]frame.north east) {Fonctions};}]
\vspace*{1.75cm}
\begin{minipage}[t][6.9cm][c]{\linewidth}\LARGE\centering
\begin{minipage}{0.95\linewidth}\raggedright\linespread{1.25}\selectfont Quelle est l'équation réduite de la droite de coefficient directeur $4$ passant par $A(-1\,;\,-2)$~?\end{minipage}
\end{minipage}

\vfill
\begin{tcolorbox}[enhanced,nobeforeafter,width=7.1cm,height=2.1cm,valign=center,colback=cfon!6,colframe=cfon,arc=2.5mm,boxrule=0.9pt,left=3mm]{\Large\bfseries\color{cfon}a)}\quad\begin{minipage}[c]{5.3cm}\raggedright\Large $y=4x+2$\end{minipage}\end{tcolorbox}\hfill\begin{tcolorbox}[enhanced,nobeforeafter,width=7.1cm,height=2.1cm,valign=center,colback=cfon!6,colframe=cfon,arc=2.5mm,boxrule=0.9pt,left=3mm]{\Large\bfseries\color{cfon}b)}\quad\begin{minipage}[c]{5.3cm}\raggedright\Large $y=4x-6$\end{minipage}\end{tcolorbox}\par\vspace{2.5mm}
\begin{tcolorbox}[enhanced,nobeforeafter,width=7.1cm,height=2.1cm,valign=center,colback=cfon!6,colframe=cfon,arc=2.5mm,boxrule=0.9pt,left=3mm]{\Large\bfseries\color{cfon}c)}\quad\begin{minipage}[c]{5.3cm}\raggedright\Large $y=4x-2$\end{minipage}\end{tcolorbox}\hfill\begin{tcolorbox}[enhanced,nobeforeafter,width=7.1cm,height=2.1cm,valign=center,colback=cfon!6,colframe=cfon,arc=2.5mm,boxrule=0.9pt,left=3mm]{\Large\bfseries\color{cfon}d)}\quad\begin{minipage}[c]{5.3cm}\raggedright\Large $y=-x+2$\end{minipage}\end{tcolorbox}
\end{tcolorbox}
\newpage
\begin{tcolorbox}[enhanced,width=16cm,height=14cm,colback=white,colframe=cfon,boxrule=2pt,arc=0mm,
left=6mm,right=6mm,top=0mm,bottom=4mm,sharp corners,
overlay={\fill[cfon] (frame.north west) rectangle ([yshift=-1.25cm]frame.north east);
\node[anchor=west,text=white,font=\sffamily\bfseries\LARGE] at ([xshift=5mm,yshift=-0.62cm]frame.north west) {Question 41 B};
\node[anchor=east,text=white,font=\sffamily\itshape\large] at ([xshift=-5mm,yshift=-0.62cm]frame.north east) {Fonctions};}]
\vspace*{1.75cm}
\begin{minipage}[t][6.9cm][c]{\linewidth}\LARGE\centering
\begin{minipage}{0.95\linewidth}\raggedright\linespread{1.25}\selectfont Pour quels réels $x$ a-t-on $5x+10>0$~?\end{minipage}
\end{minipage}

\vfill
\begin{tcolorbox}[enhanced,nobeforeafter,width=7.1cm,height=2.1cm,valign=center,colback=cfon!6,colframe=cfon,arc=2.5mm,boxrule=0.9pt,left=3mm]{\Large\bfseries\color{cfon}a)}\quad\begin{minipage}[c]{5.3cm}\raggedright\Large $x>2$\end{minipage}\end{tcolorbox}\hfill\begin{tcolorbox}[enhanced,nobeforeafter,width=7.1cm,height=2.1cm,valign=center,colback=cfon!6,colframe=cfon,arc=2.5mm,boxrule=0.9pt,left=3mm]{\Large\bfseries\color{cfon}b)}\quad\begin{minipage}[c]{5.3cm}\raggedright\Large $x<2$\end{minipage}\end{tcolorbox}\par\vspace{2.5mm}
\begin{tcolorbox}[enhanced,nobeforeafter,width=7.1cm,height=2.1cm,valign=center,colback=cfon!6,colframe=cfon,arc=2.5mm,boxrule=0.9pt,left=3mm]{\Large\bfseries\color{cfon}c)}\quad\begin{minipage}[c]{5.3cm}\raggedright\Large $x<-2$\end{minipage}\end{tcolorbox}\hfill\begin{tcolorbox}[enhanced,nobeforeafter,width=7.1cm,height=2.1cm,valign=center,colback=cfon!6,colframe=cfon,arc=2.5mm,boxrule=0.9pt,left=3mm]{\Large\bfseries\color{cfon}d)}\quad\begin{minipage}[c]{5.3cm}\raggedright\Large $x>-2$\end{minipage}\end{tcolorbox}
\end{tcolorbox}
\newpage
\begin{tcolorbox}[enhanced,width=16cm,height=14cm,colback=white,colframe=cfon,boxrule=2pt,arc=0mm,
left=6mm,right=6mm,top=0mm,bottom=4mm,sharp corners,
overlay={\fill[cfon] (frame.north west) rectangle ([yshift=-1.25cm]frame.north east);
\node[anchor=west,text=white,font=\sffamily\bfseries\LARGE] at ([xshift=5mm,yshift=-0.62cm]frame.north west) {Question 42 B};
\node[anchor=east,text=white,font=\sffamily\itshape\large] at ([xshift=-5mm,yshift=-0.62cm]frame.north east) {Fonctions};}]
\vspace*{1.75cm}
\begin{minipage}[t][6.9cm][c]{\linewidth}\LARGE\centering
\begin{minipage}{0.95\linewidth}\raggedright\linespread{1.25}\selectfont Pour quels réels $x$ a-t-on $-4x-8>0$~?\end{minipage}
\end{minipage}

\vfill
\begin{tcolorbox}[enhanced,nobeforeafter,width=7.1cm,height=2.1cm,valign=center,colback=cfon!6,colframe=cfon,arc=2.5mm,boxrule=0.9pt,left=3mm]{\Large\bfseries\color{cfon}a)}\quad\begin{minipage}[c]{5.3cm}\raggedright\Large $x>2$\end{minipage}\end{tcolorbox}\hfill\begin{tcolorbox}[enhanced,nobeforeafter,width=7.1cm,height=2.1cm,valign=center,colback=cfon!6,colframe=cfon,arc=2.5mm,boxrule=0.9pt,left=3mm]{\Large\bfseries\color{cfon}b)}\quad\begin{minipage}[c]{5.3cm}\raggedright\Large $x<-2$\end{minipage}\end{tcolorbox}\par\vspace{2.5mm}
\begin{tcolorbox}[enhanced,nobeforeafter,width=7.1cm,height=2.1cm,valign=center,colback=cfon!6,colframe=cfon,arc=2.5mm,boxrule=0.9pt,left=3mm]{\Large\bfseries\color{cfon}c)}\quad\begin{minipage}[c]{5.3cm}\raggedright\Large $x<2$\end{minipage}\end{tcolorbox}\hfill\begin{tcolorbox}[enhanced,nobeforeafter,width=7.1cm,height=2.1cm,valign=center,colback=cfon!6,colframe=cfon,arc=2.5mm,boxrule=0.9pt,left=3mm]{\Large\bfseries\color{cfon}d)}\quad\begin{minipage}[c]{5.3cm}\raggedright\Large $x>-2$\end{minipage}\end{tcolorbox}
\end{tcolorbox}
\newpage
\begin{tcolorbox}[enhanced,width=16cm,height=14cm,colback=white,colframe=cfon,boxrule=2pt,arc=0mm,
left=6mm,right=6mm,top=0mm,bottom=4mm,sharp corners,
overlay={\fill[cfon] (frame.north west) rectangle ([yshift=-1.25cm]frame.north east);
\node[anchor=west,text=white,font=\sffamily\bfseries\LARGE] at ([xshift=5mm,yshift=-0.62cm]frame.north west) {Question 43 B};
\node[anchor=east,text=white,font=\sffamily\itshape\large] at ([xshift=-5mm,yshift=-0.62cm]frame.north east) {Fonctions};}]
\vspace*{1.75cm}
\begin{minipage}[t][6.9cm][c]{\linewidth}\LARGE\centering
\begin{minipage}{0.95\linewidth}\raggedright\linespread{1.25}\selectfont Quelles sont les solutions de l'équation $x^2=-4$~?\end{minipage}
\end{minipage}

\vfill
\begin{tcolorbox}[enhanced,nobeforeafter,width=7.1cm,height=2.1cm,valign=center,colback=cfon!6,colframe=cfon,arc=2.5mm,boxrule=0.9pt,left=3mm]{\Large\bfseries\color{cfon}a)}\quad\begin{minipage}[c]{5.3cm}\raggedright\Large aucune solution\end{minipage}\end{tcolorbox}\hfill\begin{tcolorbox}[enhanced,nobeforeafter,width=7.1cm,height=2.1cm,valign=center,colback=cfon!6,colframe=cfon,arc=2.5mm,boxrule=0.9pt,left=3mm]{\Large\bfseries\color{cfon}b)}\quad\begin{minipage}[c]{5.3cm}\raggedright\Large $\{-2\,;\,2\}$\end{minipage}\end{tcolorbox}\par\vspace{2.5mm}
\begin{tcolorbox}[enhanced,nobeforeafter,width=7.1cm,height=2.1cm,valign=center,colback=cfon!6,colframe=cfon,arc=2.5mm,boxrule=0.9pt,left=3mm]{\Large\bfseries\color{cfon}c)}\quad\begin{minipage}[c]{5.3cm}\raggedright\Large $\{-2\}$\end{minipage}\end{tcolorbox}\hfill\begin{tcolorbox}[enhanced,nobeforeafter,width=7.1cm,height=2.1cm,valign=center,colback=cfon!6,colframe=cfon,arc=2.5mm,boxrule=0.9pt,left=3mm]{\Large\bfseries\color{cfon}d)}\quad\begin{minipage}[c]{5.3cm}\raggedright\Large $\{-4\}$\end{minipage}\end{tcolorbox}
\end{tcolorbox}
\newpage
\begin{tcolorbox}[enhanced,width=16cm,height=14cm,colback=white,colframe=cfon,boxrule=2pt,arc=0mm,
left=6mm,right=6mm,top=0mm,bottom=4mm,sharp corners,
overlay={\fill[cfon] (frame.north west) rectangle ([yshift=-1.25cm]frame.north east);
\node[anchor=west,text=white,font=\sffamily\bfseries\LARGE] at ([xshift=5mm,yshift=-0.62cm]frame.north west) {Question 44 B};
\node[anchor=east,text=white,font=\sffamily\itshape\large] at ([xshift=-5mm,yshift=-0.62cm]frame.north east) {Fonctions};}]
\vspace*{1.75cm}
\begin{minipage}[t][6.9cm][c]{\linewidth}\LARGE\centering
\begin{minipage}{0.95\linewidth}\raggedright\linespread{1.25}\selectfont Quelles sont les solutions de l'équation $x^2=0$~?\end{minipage}
\end{minipage}

\vfill
\begin{tcolorbox}[enhanced,nobeforeafter,width=7.1cm,height=2.1cm,valign=center,colback=cfon!6,colframe=cfon,arc=2.5mm,boxrule=0.9pt,left=3mm]{\Large\bfseries\color{cfon}a)}\quad\begin{minipage}[c]{5.3cm}\raggedright\Large aucune solution\end{minipage}\end{tcolorbox}\hfill\begin{tcolorbox}[enhanced,nobeforeafter,width=7.1cm,height=2.1cm,valign=center,colback=cfon!6,colframe=cfon,arc=2.5mm,boxrule=0.9pt,left=3mm]{\Large\bfseries\color{cfon}b)}\quad\begin{minipage}[c]{5.3cm}\raggedright\Large $\{-1\,;\,1\}$\end{minipage}\end{tcolorbox}\par\vspace{2.5mm}
\begin{tcolorbox}[enhanced,nobeforeafter,width=7.1cm,height=2.1cm,valign=center,colback=cfon!6,colframe=cfon,arc=2.5mm,boxrule=0.9pt,left=3mm]{\Large\bfseries\color{cfon}c)}\quad\begin{minipage}[c]{5.3cm}\raggedright\Large $\{0\}$\end{minipage}\end{tcolorbox}\hfill\begin{tcolorbox}[enhanced,nobeforeafter,width=7.1cm,height=2.1cm,valign=center,colback=cfon!6,colframe=cfon,arc=2.5mm,boxrule=0.9pt,left=3mm]{\Large\bfseries\color{cfon}d)}\quad\begin{minipage}[c]{5.3cm}\raggedright\Large $\{0\,;\,1\}$\end{minipage}\end{tcolorbox}
\end{tcolorbox}
\newpage
\begin{tcolorbox}[enhanced,width=16cm,height=14cm,colback=white,colframe=cfon,boxrule=2pt,arc=0mm,
left=6mm,right=6mm,top=0mm,bottom=4mm,sharp corners,
overlay={\fill[cfon] (frame.north west) rectangle ([yshift=-1.25cm]frame.north east);
\node[anchor=west,text=white,font=\sffamily\bfseries\LARGE] at ([xshift=5mm,yshift=-0.62cm]frame.north west) {Question 45 B};
\node[anchor=east,text=white,font=\sffamily\itshape\large] at ([xshift=-5mm,yshift=-0.62cm]frame.north east) {Fonctions};}]
\vspace*{1.75cm}
\begin{minipage}[t][6.9cm][c]{\linewidth}\LARGE\centering
\begin{minipage}{0.95\linewidth}\raggedright\linespread{1.25}\selectfont Quel est le sens de variation de la fonction $f(x)=4$ sur $\mathbb{R}$~?\end{minipage}
\end{minipage}

\vfill
\begin{tcolorbox}[enhanced,nobeforeafter,width=7.1cm,height=2.1cm,valign=center,colback=cfon!6,colframe=cfon,arc=2.5mm,boxrule=0.9pt,left=3mm]{\Large\bfseries\color{cfon}a)}\quad\begin{minipage}[c]{5.3cm}\raggedright\large croissante sur $\mathbb{R}$\end{minipage}\end{tcolorbox}\hfill\begin{tcolorbox}[enhanced,nobeforeafter,width=7.1cm,height=2.1cm,valign=center,colback=cfon!6,colframe=cfon,arc=2.5mm,boxrule=0.9pt,left=3mm]{\Large\bfseries\color{cfon}b)}\quad\begin{minipage}[c]{5.3cm}\raggedright\large croissante puis décroissante\end{minipage}\end{tcolorbox}\par\vspace{2.5mm}
\begin{tcolorbox}[enhanced,nobeforeafter,width=7.1cm,height=2.1cm,valign=center,colback=cfon!6,colframe=cfon,arc=2.5mm,boxrule=0.9pt,left=3mm]{\Large\bfseries\color{cfon}c)}\quad\begin{minipage}[c]{5.3cm}\raggedright\large décroissante sur $\mathbb{R}$\end{minipage}\end{tcolorbox}\hfill\begin{tcolorbox}[enhanced,nobeforeafter,width=7.1cm,height=2.1cm,valign=center,colback=cfon!6,colframe=cfon,arc=2.5mm,boxrule=0.9pt,left=3mm]{\Large\bfseries\color{cfon}d)}\quad\begin{minipage}[c]{5.3cm}\raggedright\large constante sur $\mathbb{R}$\end{minipage}\end{tcolorbox}
\end{tcolorbox}
\newpage
\begin{tcolorbox}[enhanced,width=16cm,height=14cm,colback=white,colframe=cfon,boxrule=2pt,arc=0mm,
left=6mm,right=6mm,top=0mm,bottom=4mm,sharp corners,
overlay={\fill[cfon] (frame.north west) rectangle ([yshift=-1.25cm]frame.north east);
\node[anchor=west,text=white,font=\sffamily\bfseries\LARGE] at ([xshift=5mm,yshift=-0.62cm]frame.north west) {Question 46 B};
\node[anchor=east,text=white,font=\sffamily\itshape\large] at ([xshift=-5mm,yshift=-0.62cm]frame.north east) {Fonctions};}]
\vspace*{1.75cm}
\begin{minipage}[t][6.9cm][c]{\linewidth}\LARGE\centering
\begin{minipage}{0.95\linewidth}\raggedright\linespread{1.25}\selectfont Quel est le sens de variation de la fonction $f(x)=-\dfrac{1}{2}x+2$ sur $\mathbb{R}$~?\end{minipage}
\end{minipage}

\vfill
\begin{tcolorbox}[enhanced,nobeforeafter,width=7.1cm,height=2.1cm,valign=center,colback=cfon!6,colframe=cfon,arc=2.5mm,boxrule=0.9pt,left=3mm]{\Large\bfseries\color{cfon}a)}\quad\begin{minipage}[c]{5.3cm}\raggedright\large décroissante sur $\mathbb{R}$\end{minipage}\end{tcolorbox}\hfill\begin{tcolorbox}[enhanced,nobeforeafter,width=7.1cm,height=2.1cm,valign=center,colback=cfon!6,colframe=cfon,arc=2.5mm,boxrule=0.9pt,left=3mm]{\Large\bfseries\color{cfon}b)}\quad\begin{minipage}[c]{5.3cm}\raggedright\large constante sur $\mathbb{R}$\end{minipage}\end{tcolorbox}\par\vspace{2.5mm}
\begin{tcolorbox}[enhanced,nobeforeafter,width=7.1cm,height=2.1cm,valign=center,colback=cfon!6,colframe=cfon,arc=2.5mm,boxrule=0.9pt,left=3mm]{\Large\bfseries\color{cfon}c)}\quad\begin{minipage}[c]{5.3cm}\raggedright\large croissante puis décroissante\end{minipage}\end{tcolorbox}\hfill\begin{tcolorbox}[enhanced,nobeforeafter,width=7.1cm,height=2.1cm,valign=center,colback=cfon!6,colframe=cfon,arc=2.5mm,boxrule=0.9pt,left=3mm]{\Large\bfseries\color{cfon}d)}\quad\begin{minipage}[c]{5.3cm}\raggedright\large croissante sur $\mathbb{R}$\end{minipage}\end{tcolorbox}
\end{tcolorbox}
\newpage
\begin{tcolorbox}[enhanced,width=16cm,height=14cm,colback=white,colframe=cfon,boxrule=2pt,arc=0mm,
left=6mm,right=6mm,top=0mm,bottom=4mm,sharp corners,
overlay={\fill[cfon] (frame.north west) rectangle ([yshift=-1.25cm]frame.north east);
\node[anchor=west,text=white,font=\sffamily\bfseries\LARGE] at ([xshift=5mm,yshift=-0.62cm]frame.north west) {Question 47 B};
\node[anchor=east,text=white,font=\sffamily\itshape\large] at ([xshift=-5mm,yshift=-0.62cm]frame.north east) {Fonctions};}]
\vspace*{1.75cm}
\begin{minipage}[t][6.9cm][c]{\linewidth}\LARGE\centering
\begin{minipage}{0.95\linewidth}\raggedright\linespread{1.25}\selectfont Quelles sont les solutions de l'inéquation $-x+4\leqslant 2x-8$~?\end{minipage}
\end{minipage}

\vfill
\begin{tcolorbox}[enhanced,nobeforeafter,width=7.1cm,height=2.1cm,valign=center,colback=cfon!6,colframe=cfon,arc=2.5mm,boxrule=0.9pt,left=3mm]{\Large\bfseries\color{cfon}a)}\quad\begin{minipage}[c]{5.3cm}\raggedright\Large $x\geqslant 4$\end{minipage}\end{tcolorbox}\hfill\begin{tcolorbox}[enhanced,nobeforeafter,width=7.1cm,height=2.1cm,valign=center,colback=cfon!6,colframe=cfon,arc=2.5mm,boxrule=0.9pt,left=3mm]{\Large\bfseries\color{cfon}b)}\quad\begin{minipage}[c]{5.3cm}\raggedright\Large $x\geqslant -4$\end{minipage}\end{tcolorbox}\par\vspace{2.5mm}
\begin{tcolorbox}[enhanced,nobeforeafter,width=7.1cm,height=2.1cm,valign=center,colback=cfon!6,colframe=cfon,arc=2.5mm,boxrule=0.9pt,left=3mm]{\Large\bfseries\color{cfon}c)}\quad\begin{minipage}[c]{5.3cm}\raggedright\Large $x\leqslant \dfrac{4}{3}$\end{minipage}\end{tcolorbox}\hfill\begin{tcolorbox}[enhanced,nobeforeafter,width=7.1cm,height=2.1cm,valign=center,colback=cfon!6,colframe=cfon,arc=2.5mm,boxrule=0.9pt,left=3mm]{\Large\bfseries\color{cfon}d)}\quad\begin{minipage}[c]{5.3cm}\raggedright\Large $x\leqslant 4$\end{minipage}\end{tcolorbox}
\end{tcolorbox}
\newpage
\begin{tcolorbox}[enhanced,width=16cm,height=14cm,colback=white,colframe=cfon,boxrule=2pt,arc=0mm,
left=6mm,right=6mm,top=0mm,bottom=4mm,sharp corners,
overlay={\fill[cfon] (frame.north west) rectangle ([yshift=-1.25cm]frame.north east);
\node[anchor=west,text=white,font=\sffamily\bfseries\LARGE] at ([xshift=5mm,yshift=-0.62cm]frame.north west) {Question 48 B};
\node[anchor=east,text=white,font=\sffamily\itshape\large] at ([xshift=-5mm,yshift=-0.62cm]frame.north east) {Fonctions};}]
\vspace*{1.75cm}
\begin{minipage}[t][6.9cm][c]{\linewidth}\LARGE\centering
\begin{minipage}{0.95\linewidth}\raggedright\linespread{1.25}\selectfont Quelles sont les solutions de l'inéquation $4x-3\leqslant x+6$~?\end{minipage}
\end{minipage}

\vfill
\begin{tcolorbox}[enhanced,nobeforeafter,width=7.1cm,height=2.1cm,valign=center,colback=cfon!6,colframe=cfon,arc=2.5mm,boxrule=0.9pt,left=3mm]{\Large\bfseries\color{cfon}a)}\quad\begin{minipage}[c]{5.3cm}\raggedright\Large $x\leqslant 1$\end{minipage}\end{tcolorbox}\hfill\begin{tcolorbox}[enhanced,nobeforeafter,width=7.1cm,height=2.1cm,valign=center,colback=cfon!6,colframe=cfon,arc=2.5mm,boxrule=0.9pt,left=3mm]{\Large\bfseries\color{cfon}b)}\quad\begin{minipage}[c]{5.3cm}\raggedright\Large $x\geqslant -3$\end{minipage}\end{tcolorbox}\par\vspace{2.5mm}
\begin{tcolorbox}[enhanced,nobeforeafter,width=7.1cm,height=2.1cm,valign=center,colback=cfon!6,colframe=cfon,arc=2.5mm,boxrule=0.9pt,left=3mm]{\Large\bfseries\color{cfon}c)}\quad\begin{minipage}[c]{5.3cm}\raggedright\Large $x\leqslant 3$\end{minipage}\end{tcolorbox}\hfill\begin{tcolorbox}[enhanced,nobeforeafter,width=7.1cm,height=2.1cm,valign=center,colback=cfon!6,colframe=cfon,arc=2.5mm,boxrule=0.9pt,left=3mm]{\Large\bfseries\color{cfon}d)}\quad\begin{minipage}[c]{5.3cm}\raggedright\Large $x\geqslant 3$\end{minipage}\end{tcolorbox}
\end{tcolorbox}
\newpage
\begin{tcolorbox}[enhanced,width=16cm,height=14cm,colback=white,colframe=cvec,boxrule=2pt,arc=0mm,
left=6mm,right=6mm,top=0mm,bottom=4mm,sharp corners,
overlay={\fill[cvec] (frame.north west) rectangle ([yshift=-1.25cm]frame.north east);
\node[anchor=west,text=white,font=\sffamily\bfseries\LARGE] at ([xshift=5mm,yshift=-0.62cm]frame.north west) {Question 49 B};
\node[anchor=east,text=white,font=\sffamily\itshape\large] at ([xshift=-5mm,yshift=-0.62cm]frame.north east) {Vecteurs et repérage};}]
\vspace*{1.75cm}
\begin{minipage}[t][6.9cm][c]{\linewidth}\LARGE\centering
\begin{minipage}{0.95\linewidth}\raggedright\linespread{1.25}\selectfont Dans un repère, $A(0\,;\,-4)$ et $B(-3\,;\,2)$. Quelles sont les coordonnées du vecteur $\overrightarrow{AB}$~?\end{minipage}
\end{minipage}

\vfill
\begin{tcolorbox}[enhanced,nobeforeafter,width=7.1cm,height=2.1cm,valign=center,colback=cvec!6,colframe=cvec,arc=2.5mm,boxrule=0.9pt,left=3mm]{\Large\bfseries\color{cvec}a)}\quad\begin{minipage}[c]{5.3cm}\raggedright\Large $(3\,;\,-6)$\end{minipage}\end{tcolorbox}\hfill\begin{tcolorbox}[enhanced,nobeforeafter,width=7.1cm,height=2.1cm,valign=center,colback=cvec!6,colframe=cvec,arc=2.5mm,boxrule=0.9pt,left=3mm]{\Large\bfseries\color{cvec}b)}\quad\begin{minipage}[c]{5.3cm}\raggedright\Large $(6\,;\,-3)$\end{minipage}\end{tcolorbox}\par\vspace{2.5mm}
\begin{tcolorbox}[enhanced,nobeforeafter,width=7.1cm,height=2.1cm,valign=center,colback=cvec!6,colframe=cvec,arc=2.5mm,boxrule=0.9pt,left=3mm]{\Large\bfseries\color{cvec}c)}\quad\begin{minipage}[c]{5.3cm}\raggedright\Large $(-3\,;\,6)$\end{minipage}\end{tcolorbox}\hfill\begin{tcolorbox}[enhanced,nobeforeafter,width=7.1cm,height=2.1cm,valign=center,colback=cvec!6,colframe=cvec,arc=2.5mm,boxrule=0.9pt,left=3mm]{\Large\bfseries\color{cvec}d)}\quad\begin{minipage}[c]{5.3cm}\raggedright\Large $(-3\,;\,-2)$\end{minipage}\end{tcolorbox}
\end{tcolorbox}
\newpage
\begin{tcolorbox}[enhanced,width=16cm,height=14cm,colback=white,colframe=cvec,boxrule=2pt,arc=0mm,
left=6mm,right=6mm,top=0mm,bottom=4mm,sharp corners,
overlay={\fill[cvec] (frame.north west) rectangle ([yshift=-1.25cm]frame.north east);
\node[anchor=west,text=white,font=\sffamily\bfseries\LARGE] at ([xshift=5mm,yshift=-0.62cm]frame.north west) {Question 50 B};
\node[anchor=east,text=white,font=\sffamily\itshape\large] at ([xshift=-5mm,yshift=-0.62cm]frame.north east) {Vecteurs et repérage};}]
\vspace*{1.75cm}
\begin{minipage}[t][6.9cm][c]{\linewidth}\LARGE\centering
\begin{minipage}{0.95\linewidth}\raggedright\linespread{1.25}\selectfont Dans un repère, $A(6\,;\,1)$ et $B(2\,;\,7)$. Quelles sont les coordonnées du vecteur $\overrightarrow{AB}$~?\end{minipage}
\end{minipage}

\vfill
\begin{tcolorbox}[enhanced,nobeforeafter,width=7.1cm,height=2.1cm,valign=center,colback=cvec!6,colframe=cvec,arc=2.5mm,boxrule=0.9pt,left=3mm]{\Large\bfseries\color{cvec}a)}\quad\begin{minipage}[c]{5.3cm}\raggedright\Large $(8\,;\,8)$\end{minipage}\end{tcolorbox}\hfill\begin{tcolorbox}[enhanced,nobeforeafter,width=7.1cm,height=2.1cm,valign=center,colback=cvec!6,colframe=cvec,arc=2.5mm,boxrule=0.9pt,left=3mm]{\Large\bfseries\color{cvec}b)}\quad\begin{minipage}[c]{5.3cm}\raggedright\Large $(-4\,;\,6)$\end{minipage}\end{tcolorbox}\par\vspace{2.5mm}
\begin{tcolorbox}[enhanced,nobeforeafter,width=7.1cm,height=2.1cm,valign=center,colback=cvec!6,colframe=cvec,arc=2.5mm,boxrule=0.9pt,left=3mm]{\Large\bfseries\color{cvec}c)}\quad\begin{minipage}[c]{5.3cm}\raggedright\Large $(6\,;\,-4)$\end{minipage}\end{tcolorbox}\hfill\begin{tcolorbox}[enhanced,nobeforeafter,width=7.1cm,height=2.1cm,valign=center,colback=cvec!6,colframe=cvec,arc=2.5mm,boxrule=0.9pt,left=3mm]{\Large\bfseries\color{cvec}d)}\quad\begin{minipage}[c]{5.3cm}\raggedright\Large $(4\,;\,-6)$\end{minipage}\end{tcolorbox}
\end{tcolorbox}
\newpage
\begin{tcolorbox}[enhanced,width=16cm,height=14cm,colback=white,colframe=cvec,boxrule=2pt,arc=0mm,
left=6mm,right=6mm,top=0mm,bottom=4mm,sharp corners,
overlay={\fill[cvec] (frame.north west) rectangle ([yshift=-1.25cm]frame.north east);
\node[anchor=west,text=white,font=\sffamily\bfseries\LARGE] at ([xshift=5mm,yshift=-0.62cm]frame.north west) {Question 51 B};
\node[anchor=east,text=white,font=\sffamily\itshape\large] at ([xshift=-5mm,yshift=-0.62cm]frame.north east) {Vecteurs et repérage};}]
\vspace*{1.75cm}
\begin{minipage}[t][6.9cm][c]{\linewidth}\LARGE\centering
\begin{minipage}{0.95\linewidth}\raggedright\linespread{1.25}\selectfont Dans un repère, $A(0\,;\,-3)$ et $B(-6\,;\,7)$. Quelles sont les coordonnées du milieu de $[AB]$~?\end{minipage}
\end{minipage}

\vfill
\begin{tcolorbox}[enhanced,nobeforeafter,width=7.1cm,height=2.1cm,valign=center,colback=cvec!6,colframe=cvec,arc=2.5mm,boxrule=0.9pt,left=3mm]{\Large\bfseries\color{cvec}a)}\quad\begin{minipage}[c]{5.3cm}\raggedright\Large $(-3\,;\,2)$\end{minipage}\end{tcolorbox}\hfill\begin{tcolorbox}[enhanced,nobeforeafter,width=7.1cm,height=2.1cm,valign=center,colback=cvec!6,colframe=cvec,arc=2.5mm,boxrule=0.9pt,left=3mm]{\Large\bfseries\color{cvec}b)}\quad\begin{minipage}[c]{5.3cm}\raggedright\Large $(-3\,;\,5)$\end{minipage}\end{tcolorbox}\par\vspace{2.5mm}
\begin{tcolorbox}[enhanced,nobeforeafter,width=7.1cm,height=2.1cm,valign=center,colback=cvec!6,colframe=cvec,arc=2.5mm,boxrule=0.9pt,left=3mm]{\Large\bfseries\color{cvec}c)}\quad\begin{minipage}[c]{5.3cm}\raggedright\Large $(-3\,;\,-5)$\end{minipage}\end{tcolorbox}\hfill\begin{tcolorbox}[enhanced,nobeforeafter,width=7.1cm,height=2.1cm,valign=center,colback=cvec!6,colframe=cvec,arc=2.5mm,boxrule=0.9pt,left=3mm]{\Large\bfseries\color{cvec}d)}\quad\begin{minipage}[c]{5.3cm}\raggedright\Large $(-6\,;\,4)$\end{minipage}\end{tcolorbox}
\end{tcolorbox}
\newpage
\begin{tcolorbox}[enhanced,width=16cm,height=14cm,colback=white,colframe=cvec,boxrule=2pt,arc=0mm,
left=6mm,right=6mm,top=0mm,bottom=4mm,sharp corners,
overlay={\fill[cvec] (frame.north west) rectangle ([yshift=-1.25cm]frame.north east);
\node[anchor=west,text=white,font=\sffamily\bfseries\LARGE] at ([xshift=5mm,yshift=-0.62cm]frame.north west) {Question 52 B};
\node[anchor=east,text=white,font=\sffamily\itshape\large] at ([xshift=-5mm,yshift=-0.62cm]frame.north east) {Vecteurs et repérage};}]
\vspace*{1.75cm}
\begin{minipage}[t][6.9cm][c]{\linewidth}\LARGE\centering
\begin{minipage}{0.95\linewidth}\raggedright\linespread{1.25}\selectfont Dans un repère, $A(3\,;\,3)$ et $B(8\,;\,-1)$. Quelles sont les coordonnées du milieu de $[AB]$~?\end{minipage}
\end{minipage}

\vfill
\begin{tcolorbox}[enhanced,nobeforeafter,width=7.1cm,height=2.1cm,valign=center,colback=cvec!6,colframe=cvec,arc=2.5mm,boxrule=0.9pt,left=3mm]{\Large\bfseries\color{cvec}a)}\quad\begin{minipage}[c]{5.3cm}\raggedright\Large $(2{,}5\,;\,-2)$\end{minipage}\end{tcolorbox}\hfill\begin{tcolorbox}[enhanced,nobeforeafter,width=7.1cm,height=2.1cm,valign=center,colback=cvec!6,colframe=cvec,arc=2.5mm,boxrule=0.9pt,left=3mm]{\Large\bfseries\color{cvec}b)}\quad\begin{minipage}[c]{5.3cm}\raggedright\Large $(11\,;\,2)$\end{minipage}\end{tcolorbox}\par\vspace{2.5mm}
\begin{tcolorbox}[enhanced,nobeforeafter,width=7.1cm,height=2.1cm,valign=center,colback=cvec!6,colframe=cvec,arc=2.5mm,boxrule=0.9pt,left=3mm]{\Large\bfseries\color{cvec}c)}\quad\begin{minipage}[c]{5.3cm}\raggedright\Large $(5{,}5\,;\,1)$\end{minipage}\end{tcolorbox}\hfill\begin{tcolorbox}[enhanced,nobeforeafter,width=7.1cm,height=2.1cm,valign=center,colback=cvec!6,colframe=cvec,arc=2.5mm,boxrule=0.9pt,left=3mm]{\Large\bfseries\color{cvec}d)}\quad\begin{minipage}[c]{5.3cm}\raggedright\Large $(5{,}5\,;\,2)$\end{minipage}\end{tcolorbox}
\end{tcolorbox}
\newpage
\begin{tcolorbox}[enhanced,width=16cm,height=14cm,colback=white,colframe=cvec,boxrule=2pt,arc=0mm,
left=6mm,right=6mm,top=0mm,bottom=4mm,sharp corners,
overlay={\fill[cvec] (frame.north west) rectangle ([yshift=-1.25cm]frame.north east);
\node[anchor=west,text=white,font=\sffamily\bfseries\LARGE] at ([xshift=5mm,yshift=-0.62cm]frame.north west) {Question 53 B};
\node[anchor=east,text=white,font=\sffamily\itshape\large] at ([xshift=-5mm,yshift=-0.62cm]frame.north east) {Vecteurs et repérage};}]
\vspace*{1.75cm}
\begin{minipage}[t][6.9cm][c]{\linewidth}\LARGE\centering
\begin{minipage}{0.95\linewidth}\raggedright\linespread{1.25}\selectfont Dans un repère orthonormé, $A(0\,;\,0)$ et $B(2\,;\,3)$. Quelle est la distance $AB$~?\end{minipage}
\end{minipage}

\vfill
\begin{tcolorbox}[enhanced,nobeforeafter,width=7.1cm,height=2.1cm,valign=center,colback=cvec!6,colframe=cvec,arc=2.5mm,boxrule=0.9pt,left=3mm]{\Large\bfseries\color{cvec}a)}\quad\begin{minipage}[c]{5.3cm}\raggedright\Large $13$\end{minipage}\end{tcolorbox}\hfill\begin{tcolorbox}[enhanced,nobeforeafter,width=7.1cm,height=2.1cm,valign=center,colback=cvec!6,colframe=cvec,arc=2.5mm,boxrule=0.9pt,left=3mm]{\Large\bfseries\color{cvec}b)}\quad\begin{minipage}[c]{5.3cm}\raggedright\Large $\sqrt{5}$\end{minipage}\end{tcolorbox}\par\vspace{2.5mm}
\begin{tcolorbox}[enhanced,nobeforeafter,width=7.1cm,height=2.1cm,valign=center,colback=cvec!6,colframe=cvec,arc=2.5mm,boxrule=0.9pt,left=3mm]{\Large\bfseries\color{cvec}c)}\quad\begin{minipage}[c]{5.3cm}\raggedright\Large $5$\end{minipage}\end{tcolorbox}\hfill\begin{tcolorbox}[enhanced,nobeforeafter,width=7.1cm,height=2.1cm,valign=center,colback=cvec!6,colframe=cvec,arc=2.5mm,boxrule=0.9pt,left=3mm]{\Large\bfseries\color{cvec}d)}\quad\begin{minipage}[c]{5.3cm}\raggedright\Large $\sqrt{13}$\end{minipage}\end{tcolorbox}
\end{tcolorbox}
\newpage
\begin{tcolorbox}[enhanced,width=16cm,height=14cm,colback=white,colframe=cvec,boxrule=2pt,arc=0mm,
left=6mm,right=6mm,top=0mm,bottom=4mm,sharp corners,
overlay={\fill[cvec] (frame.north west) rectangle ([yshift=-1.25cm]frame.north east);
\node[anchor=west,text=white,font=\sffamily\bfseries\LARGE] at ([xshift=5mm,yshift=-0.62cm]frame.north west) {Question 54 B};
\node[anchor=east,text=white,font=\sffamily\itshape\large] at ([xshift=-5mm,yshift=-0.62cm]frame.north east) {Vecteurs et repérage};}]
\vspace*{1.75cm}
\begin{minipage}[t][6.9cm][c]{\linewidth}\LARGE\centering
\begin{minipage}{0.95\linewidth}\raggedright\linespread{1.25}\selectfont Dans un repère orthonormé, $A(-3\,;\,4)$ et $B(3\,;\,-4)$. Quelle est la distance $AB$~?\end{minipage}
\end{minipage}

\vfill
\begin{tcolorbox}[enhanced,nobeforeafter,width=7.1cm,height=2.1cm,valign=center,colback=cvec!6,colframe=cvec,arc=2.5mm,boxrule=0.9pt,left=3mm]{\Large\bfseries\color{cvec}a)}\quad\begin{minipage}[c]{5.3cm}\raggedright\Large $2\sqrt{7}$\end{minipage}\end{tcolorbox}\hfill\begin{tcolorbox}[enhanced,nobeforeafter,width=7.1cm,height=2.1cm,valign=center,colback=cvec!6,colframe=cvec,arc=2.5mm,boxrule=0.9pt,left=3mm]{\Large\bfseries\color{cvec}b)}\quad\begin{minipage}[c]{5.3cm}\raggedright\Large $100$\end{minipage}\end{tcolorbox}\par\vspace{2.5mm}
\begin{tcolorbox}[enhanced,nobeforeafter,width=7.1cm,height=2.1cm,valign=center,colback=cvec!6,colframe=cvec,arc=2.5mm,boxrule=0.9pt,left=3mm]{\Large\bfseries\color{cvec}c)}\quad\begin{minipage}[c]{5.3cm}\raggedright\Large $10$\end{minipage}\end{tcolorbox}\hfill\begin{tcolorbox}[enhanced,nobeforeafter,width=7.1cm,height=2.1cm,valign=center,colback=cvec!6,colframe=cvec,arc=2.5mm,boxrule=0.9pt,left=3mm]{\Large\bfseries\color{cvec}d)}\quad\begin{minipage}[c]{5.3cm}\raggedright\Large $14$\end{minipage}\end{tcolorbox}
\end{tcolorbox}
\newpage
\begin{tcolorbox}[enhanced,width=16cm,height=14cm,colback=white,colframe=cvec,boxrule=2pt,arc=0mm,
left=6mm,right=6mm,top=0mm,bottom=4mm,sharp corners,
overlay={\fill[cvec] (frame.north west) rectangle ([yshift=-1.25cm]frame.north east);
\node[anchor=west,text=white,font=\sffamily\bfseries\LARGE] at ([xshift=5mm,yshift=-0.62cm]frame.north west) {Question 55 B};
\node[anchor=east,text=white,font=\sffamily\itshape\large] at ([xshift=-5mm,yshift=-0.62cm]frame.north east) {Vecteurs et repérage};}]
\vspace*{1.75cm}
\begin{minipage}[t][6.9cm][c]{\linewidth}\LARGE\centering
\begin{minipage}{0.95\linewidth}\raggedright\linespread{1.25}\selectfont Pour quelle valeur de $k$ les vecteurs $\vec u\begin{pmatrix}1\\5\end{pmatrix}$ et $\vec v\begin{pmatrix}k\\-10\end{pmatrix}$ sont-ils colinéaires~?\end{minipage}
\end{minipage}

\vfill
\begin{tcolorbox}[enhanced,nobeforeafter,width=7.1cm,height=2.1cm,valign=center,colback=cvec!6,colframe=cvec,arc=2.5mm,boxrule=0.9pt,left=3mm]{\Large\bfseries\color{cvec}a)}\quad\begin{minipage}[c]{5.3cm}\raggedright\Large $k=2$\end{minipage}\end{tcolorbox}\hfill\begin{tcolorbox}[enhanced,nobeforeafter,width=7.1cm,height=2.1cm,valign=center,colback=cvec!6,colframe=cvec,arc=2.5mm,boxrule=0.9pt,left=3mm]{\Large\bfseries\color{cvec}b)}\quad\begin{minipage}[c]{5.3cm}\raggedright\Large $k=-2$\end{minipage}\end{tcolorbox}\par\vspace{2.5mm}
\begin{tcolorbox}[enhanced,nobeforeafter,width=7.1cm,height=2.1cm,valign=center,colback=cvec!6,colframe=cvec,arc=2.5mm,boxrule=0.9pt,left=3mm]{\Large\bfseries\color{cvec}c)}\quad\begin{minipage}[c]{5.3cm}\raggedright\Large $k=-14$\end{minipage}\end{tcolorbox}\hfill\begin{tcolorbox}[enhanced,nobeforeafter,width=7.1cm,height=2.1cm,valign=center,colback=cvec!6,colframe=cvec,arc=2.5mm,boxrule=0.9pt,left=3mm]{\Large\bfseries\color{cvec}d)}\quad\begin{minipage}[c]{5.3cm}\raggedright\Large $k=-50$\end{minipage}\end{tcolorbox}
\end{tcolorbox}
\newpage
\begin{tcolorbox}[enhanced,width=16cm,height=14cm,colback=white,colframe=cvec,boxrule=2pt,arc=0mm,
left=6mm,right=6mm,top=0mm,bottom=4mm,sharp corners,
overlay={\fill[cvec] (frame.north west) rectangle ([yshift=-1.25cm]frame.north east);
\node[anchor=west,text=white,font=\sffamily\bfseries\LARGE] at ([xshift=5mm,yshift=-0.62cm]frame.north west) {Question 56 B};
\node[anchor=east,text=white,font=\sffamily\itshape\large] at ([xshift=-5mm,yshift=-0.62cm]frame.north east) {Vecteurs et repérage};}]
\vspace*{1.75cm}
\begin{minipage}[t][6.9cm][c]{\linewidth}\LARGE\centering
\begin{minipage}{0.95\linewidth}\raggedright\linespread{1.25}\selectfont Pour quelle valeur de $k$ les vecteurs $\vec u\begin{pmatrix}-3\\6\end{pmatrix}$ et $\vec v\begin{pmatrix}k\\4\end{pmatrix}$ sont-ils colinéaires~?\end{minipage}
\end{minipage}

\vfill
\begin{tcolorbox}[enhanced,nobeforeafter,width=7.1cm,height=2.1cm,valign=center,colback=cvec!6,colframe=cvec,arc=2.5mm,boxrule=0.9pt,left=3mm]{\Large\bfseries\color{cvec}a)}\quad\begin{minipage}[c]{5.3cm}\raggedright\Large $k=-5$\end{minipage}\end{tcolorbox}\hfill\begin{tcolorbox}[enhanced,nobeforeafter,width=7.1cm,height=2.1cm,valign=center,colback=cvec!6,colframe=cvec,arc=2.5mm,boxrule=0.9pt,left=3mm]{\Large\bfseries\color{cvec}b)}\quad\begin{minipage}[c]{5.3cm}\raggedright\Large $k=2$\end{minipage}\end{tcolorbox}\par\vspace{2.5mm}
\begin{tcolorbox}[enhanced,nobeforeafter,width=7.1cm,height=2.1cm,valign=center,colback=cvec!6,colframe=cvec,arc=2.5mm,boxrule=0.9pt,left=3mm]{\Large\bfseries\color{cvec}c)}\quad\begin{minipage}[c]{5.3cm}\raggedright\Large $k=-2$\end{minipage}\end{tcolorbox}\hfill\begin{tcolorbox}[enhanced,nobeforeafter,width=7.1cm,height=2.1cm,valign=center,colback=cvec!6,colframe=cvec,arc=2.5mm,boxrule=0.9pt,left=3mm]{\Large\bfseries\color{cvec}d)}\quad\begin{minipage}[c]{5.3cm}\raggedright\Large $k=-8$\end{minipage}\end{tcolorbox}
\end{tcolorbox}
\newpage
\begin{tcolorbox}[enhanced,width=16cm,height=14cm,colback=white,colframe=cvec,boxrule=2pt,arc=0mm,
left=6mm,right=6mm,top=0mm,bottom=4mm,sharp corners,
overlay={\fill[cvec] (frame.north west) rectangle ([yshift=-1.25cm]frame.north east);
\node[anchor=west,text=white,font=\sffamily\bfseries\LARGE] at ([xshift=5mm,yshift=-0.62cm]frame.north west) {Question 57 B};
\node[anchor=east,text=white,font=\sffamily\itshape\large] at ([xshift=-5mm,yshift=-0.62cm]frame.north east) {Vecteurs et repérage};}]
\vspace*{1.75cm}
\begin{minipage}[t][6.9cm][c]{\linewidth}\Large\centering
\begin{minipage}{0.95\linewidth}\raggedright\linespread{1.25}\selectfont On donne $\vec u\begin{pmatrix}-2\\5\end{pmatrix}$ et $\vec v\begin{pmatrix}0\\3\end{pmatrix}$. Quelles sont les coordonnées du vecteur $3\vec u+\vec v$~?\end{minipage}
\end{minipage}

\vfill
\begin{tcolorbox}[enhanced,nobeforeafter,width=7.1cm,height=2.1cm,valign=center,colback=cvec!6,colframe=cvec,arc=2.5mm,boxrule=0.9pt,left=3mm]{\Large\bfseries\color{cvec}a)}\quad\begin{minipage}[c]{5.3cm}\raggedright\Large $(18\,;\,-6)$\end{minipage}\end{tcolorbox}\hfill\begin{tcolorbox}[enhanced,nobeforeafter,width=7.1cm,height=2.1cm,valign=center,colback=cvec!6,colframe=cvec,arc=2.5mm,boxrule=0.9pt,left=3mm]{\Large\bfseries\color{cvec}b)}\quad\begin{minipage}[c]{5.3cm}\raggedright\Large $(-6\,;\,12)$\end{minipage}\end{tcolorbox}\par\vspace{2.5mm}
\begin{tcolorbox}[enhanced,nobeforeafter,width=7.1cm,height=2.1cm,valign=center,colback=cvec!6,colframe=cvec,arc=2.5mm,boxrule=0.9pt,left=3mm]{\Large\bfseries\color{cvec}c)}\quad\begin{minipage}[c]{5.3cm}\raggedright\Large $(-2\,;\,8)$\end{minipage}\end{tcolorbox}\hfill\begin{tcolorbox}[enhanced,nobeforeafter,width=7.1cm,height=2.1cm,valign=center,colback=cvec!6,colframe=cvec,arc=2.5mm,boxrule=0.9pt,left=3mm]{\Large\bfseries\color{cvec}d)}\quad\begin{minipage}[c]{5.3cm}\raggedright\Large $(-6\,;\,18)$\end{minipage}\end{tcolorbox}
\end{tcolorbox}
\newpage
\begin{tcolorbox}[enhanced,width=16cm,height=14cm,colback=white,colframe=cvec,boxrule=2pt,arc=0mm,
left=6mm,right=6mm,top=0mm,bottom=4mm,sharp corners,
overlay={\fill[cvec] (frame.north west) rectangle ([yshift=-1.25cm]frame.north east);
\node[anchor=west,text=white,font=\sffamily\bfseries\LARGE] at ([xshift=5mm,yshift=-0.62cm]frame.north west) {Question 58 B};
\node[anchor=east,text=white,font=\sffamily\itshape\large] at ([xshift=-5mm,yshift=-0.62cm]frame.north east) {Vecteurs et repérage};}]
\vspace*{1.75cm}
\begin{minipage}[t][6.9cm][c]{\linewidth}\Large\centering
\begin{minipage}{0.95\linewidth}\raggedright\linespread{1.25}\selectfont On donne $\vec u\begin{pmatrix}2\\2\end{pmatrix}$ et $\vec v\begin{pmatrix}-1\\4\end{pmatrix}$. Quelles sont les coordonnées du vecteur $-\vec u+2\vec v$~?\end{minipage}
\end{minipage}

\vfill
\begin{tcolorbox}[enhanced,nobeforeafter,width=7.1cm,height=2.1cm,valign=center,colback=cvec!6,colframe=cvec,arc=2.5mm,boxrule=0.9pt,left=3mm]{\Large\bfseries\color{cvec}a)}\quad\begin{minipage}[c]{5.3cm}\raggedright\Large $(0\,;\,-10)$\end{minipage}\end{tcolorbox}\hfill\begin{tcolorbox}[enhanced,nobeforeafter,width=7.1cm,height=2.1cm,valign=center,colback=cvec!6,colframe=cvec,arc=2.5mm,boxrule=0.9pt,left=3mm]{\Large\bfseries\color{cvec}b)}\quad\begin{minipage}[c]{5.3cm}\raggedright\Large $(-4\,;\,6)$\end{minipage}\end{tcolorbox}\par\vspace{2.5mm}
\begin{tcolorbox}[enhanced,nobeforeafter,width=7.1cm,height=2.1cm,valign=center,colback=cvec!6,colframe=cvec,arc=2.5mm,boxrule=0.9pt,left=3mm]{\Large\bfseries\color{cvec}c)}\quad\begin{minipage}[c]{5.3cm}\raggedright\Large $(1\,;\,6)$\end{minipage}\end{tcolorbox}\hfill\begin{tcolorbox}[enhanced,nobeforeafter,width=7.1cm,height=2.1cm,valign=center,colback=cvec!6,colframe=cvec,arc=2.5mm,boxrule=0.9pt,left=3mm]{\Large\bfseries\color{cvec}d)}\quad\begin{minipage}[c]{5.3cm}\raggedright\Large $(-3\,;\,2)$\end{minipage}\end{tcolorbox}
\end{tcolorbox}
\newpage
\begin{tcolorbox}[enhanced,width=16cm,height=14cm,colback=white,colframe=cvec,boxrule=2pt,arc=0mm,
left=6mm,right=6mm,top=0mm,bottom=4mm,sharp corners,
overlay={\fill[cvec] (frame.north west) rectangle ([yshift=-1.25cm]frame.north east);
\node[anchor=west,text=white,font=\sffamily\bfseries\LARGE] at ([xshift=5mm,yshift=-0.62cm]frame.north west) {Question 59 B};
\node[anchor=east,text=white,font=\sffamily\itshape\large] at ([xshift=-5mm,yshift=-0.62cm]frame.north east) {Vecteurs et repérage};}]
\vspace*{1.75cm}
\begin{minipage}[t][6.9cm][c]{\linewidth}\LARGE\centering
\begin{minipage}{0.95\linewidth}\raggedright\linespread{1.25}\selectfont Que vaut le déterminant $\det(\vec u,\vec v)$ pour $\vec u\begin{pmatrix}3\\-2\end{pmatrix}$ et $\vec v\begin{pmatrix}6\\-4\end{pmatrix}$~?\end{minipage}
\end{minipage}

\vfill
\begin{tcolorbox}[enhanced,nobeforeafter,width=7.1cm,height=2.1cm,valign=center,colback=cvec!6,colframe=cvec,arc=2.5mm,boxrule=0.9pt,left=3mm]{\Large\bfseries\color{cvec}a)}\quad\begin{minipage}[c]{5.3cm}\raggedright\Large $1$\end{minipage}\end{tcolorbox}\hfill\begin{tcolorbox}[enhanced,nobeforeafter,width=7.1cm,height=2.1cm,valign=center,colback=cvec!6,colframe=cvec,arc=2.5mm,boxrule=0.9pt,left=3mm]{\Large\bfseries\color{cvec}b)}\quad\begin{minipage}[c]{5.3cm}\raggedright\Large $0$\end{minipage}\end{tcolorbox}\par\vspace{2.5mm}
\begin{tcolorbox}[enhanced,nobeforeafter,width=7.1cm,height=2.1cm,valign=center,colback=cvec!6,colframe=cvec,arc=2.5mm,boxrule=0.9pt,left=3mm]{\Large\bfseries\color{cvec}c)}\quad\begin{minipage}[c]{5.3cm}\raggedright\Large $-24$\end{minipage}\end{tcolorbox}\hfill\begin{tcolorbox}[enhanced,nobeforeafter,width=7.1cm,height=2.1cm,valign=center,colback=cvec!6,colframe=cvec,arc=2.5mm,boxrule=0.9pt,left=3mm]{\Large\bfseries\color{cvec}d)}\quad\begin{minipage}[c]{5.3cm}\raggedright\Large $26$\end{minipage}\end{tcolorbox}
\end{tcolorbox}
\newpage
\begin{tcolorbox}[enhanced,width=16cm,height=14cm,colback=white,colframe=cvec,boxrule=2pt,arc=0mm,
left=6mm,right=6mm,top=0mm,bottom=4mm,sharp corners,
overlay={\fill[cvec] (frame.north west) rectangle ([yshift=-1.25cm]frame.north east);
\node[anchor=west,text=white,font=\sffamily\bfseries\LARGE] at ([xshift=5mm,yshift=-0.62cm]frame.north west) {Question 60 B};
\node[anchor=east,text=white,font=\sffamily\itshape\large] at ([xshift=-5mm,yshift=-0.62cm]frame.north east) {Vecteurs et repérage};}]
\vspace*{1.75cm}
\begin{minipage}[t][6.9cm][c]{\linewidth}\LARGE\centering
\begin{minipage}{0.95\linewidth}\raggedright\linespread{1.25}\selectfont Que vaut le déterminant $\det(\vec u,\vec v)$ pour $\vec u\begin{pmatrix}4\\1\end{pmatrix}$ et $\vec v\begin{pmatrix}-2\\5\end{pmatrix}$~?\end{minipage}
\end{minipage}

\vfill
\begin{tcolorbox}[enhanced,nobeforeafter,width=7.1cm,height=2.1cm,valign=center,colback=cvec!6,colframe=cvec,arc=2.5mm,boxrule=0.9pt,left=3mm]{\Large\bfseries\color{cvec}a)}\quad\begin{minipage}[c]{5.3cm}\raggedright\Large $-3$\end{minipage}\end{tcolorbox}\hfill\begin{tcolorbox}[enhanced,nobeforeafter,width=7.1cm,height=2.1cm,valign=center,colback=cvec!6,colframe=cvec,arc=2.5mm,boxrule=0.9pt,left=3mm]{\Large\bfseries\color{cvec}b)}\quad\begin{minipage}[c]{5.3cm}\raggedright\Large $22$\end{minipage}\end{tcolorbox}\par\vspace{2.5mm}
\begin{tcolorbox}[enhanced,nobeforeafter,width=7.1cm,height=2.1cm,valign=center,colback=cvec!6,colframe=cvec,arc=2.5mm,boxrule=0.9pt,left=3mm]{\Large\bfseries\color{cvec}c)}\quad\begin{minipage}[c]{5.3cm}\raggedright\Large $-22$\end{minipage}\end{tcolorbox}\hfill\begin{tcolorbox}[enhanced,nobeforeafter,width=7.1cm,height=2.1cm,valign=center,colback=cvec!6,colframe=cvec,arc=2.5mm,boxrule=0.9pt,left=3mm]{\Large\bfseries\color{cvec}d)}\quad\begin{minipage}[c]{5.3cm}\raggedright\Large $18$\end{minipage}\end{tcolorbox}
\end{tcolorbox}
\newpage
\begin{tcolorbox}[enhanced,width=16cm,height=14cm,colback=white,colframe=cvec,boxrule=2pt,arc=0mm,
left=6mm,right=6mm,top=0mm,bottom=4mm,sharp corners,
overlay={\fill[cvec] (frame.north west) rectangle ([yshift=-1.25cm]frame.north east);
\node[anchor=west,text=white,font=\sffamily\bfseries\LARGE] at ([xshift=5mm,yshift=-0.62cm]frame.north west) {Question 61 B};
\node[anchor=east,text=white,font=\sffamily\itshape\large] at ([xshift=-5mm,yshift=-0.62cm]frame.north east) {Vecteurs et repérage};}]
\vspace*{1.75cm}
\begin{minipage}[t][6.9cm][c]{\linewidth}\LARGE\centering
\begin{minipage}{0.95\linewidth}\raggedright\linespread{1.25}\selectfont Dans un repère, $A(0\,;\,-1)$, $B(3\,;\,-2)$ et $C(4\,;\,2)$. Quelles sont les coordonnées du point $D$ tel que $ABCD$ soit un parallélogramme~?\end{minipage}
\end{minipage}

\vfill
\begin{tcolorbox}[enhanced,nobeforeafter,width=7.1cm,height=2.1cm,valign=center,colback=cvec!6,colframe=cvec,arc=2.5mm,boxrule=0.9pt,left=3mm]{\Large\bfseries\color{cvec}a)}\quad\begin{minipage}[c]{5.3cm}\raggedright\Large $(1\,;\,3)$\end{minipage}\end{tcolorbox}\hfill\begin{tcolorbox}[enhanced,nobeforeafter,width=7.1cm,height=2.1cm,valign=center,colback=cvec!6,colframe=cvec,arc=2.5mm,boxrule=0.9pt,left=3mm]{\Large\bfseries\color{cvec}b)}\quad\begin{minipage}[c]{5.3cm}\raggedright\Large $(2\,;\,0{,}5)$\end{minipage}\end{tcolorbox}\par\vspace{2.5mm}
\begin{tcolorbox}[enhanced,nobeforeafter,width=7.1cm,height=2.1cm,valign=center,colback=cvec!6,colframe=cvec,arc=2.5mm,boxrule=0.9pt,left=3mm]{\Large\bfseries\color{cvec}c)}\quad\begin{minipage}[c]{5.3cm}\raggedright\Large $(-1\,;\,-5)$\end{minipage}\end{tcolorbox}\hfill\begin{tcolorbox}[enhanced,nobeforeafter,width=7.1cm,height=2.1cm,valign=center,colback=cvec!6,colframe=cvec,arc=2.5mm,boxrule=0.9pt,left=3mm]{\Large\bfseries\color{cvec}d)}\quad\begin{minipage}[c]{5.3cm}\raggedright\Large $(7\,;\,1)$\end{minipage}\end{tcolorbox}
\end{tcolorbox}
\newpage
\begin{tcolorbox}[enhanced,width=16cm,height=14cm,colback=white,colframe=cvec,boxrule=2pt,arc=0mm,
left=6mm,right=6mm,top=0mm,bottom=4mm,sharp corners,
overlay={\fill[cvec] (frame.north west) rectangle ([yshift=-1.25cm]frame.north east);
\node[anchor=west,text=white,font=\sffamily\bfseries\LARGE] at ([xshift=5mm,yshift=-0.62cm]frame.north west) {Question 62 B};
\node[anchor=east,text=white,font=\sffamily\itshape\large] at ([xshift=-5mm,yshift=-0.62cm]frame.north east) {Vecteurs et repérage};}]
\vspace*{1.75cm}
\begin{minipage}[t][6.9cm][c]{\linewidth}\LARGE\centering
\begin{minipage}{0.95\linewidth}\raggedright\linespread{1.25}\selectfont Dans un repère, $A(2\,;\,3)$, $B(6\,;\,1)$ et $C(5\,;\,-2)$. Quelles sont les coordonnées du point $D$ tel que $ABCD$ soit un parallélogramme~?\end{minipage}
\end{minipage}

\vfill
\begin{tcolorbox}[enhanced,nobeforeafter,width=7.1cm,height=2.1cm,valign=center,colback=cvec!6,colframe=cvec,arc=2.5mm,boxrule=0.9pt,left=3mm]{\Large\bfseries\color{cvec}a)}\quad\begin{minipage}[c]{5.3cm}\raggedright\Large $(9\,;\,-4)$\end{minipage}\end{tcolorbox}\hfill\begin{tcolorbox}[enhanced,nobeforeafter,width=7.1cm,height=2.1cm,valign=center,colback=cvec!6,colframe=cvec,arc=2.5mm,boxrule=0.9pt,left=3mm]{\Large\bfseries\color{cvec}b)}\quad\begin{minipage}[c]{5.3cm}\raggedright\Large $(1\,;\,0)$\end{minipage}\end{tcolorbox}\par\vspace{2.5mm}
\begin{tcolorbox}[enhanced,nobeforeafter,width=7.1cm,height=2.1cm,valign=center,colback=cvec!6,colframe=cvec,arc=2.5mm,boxrule=0.9pt,left=3mm]{\Large\bfseries\color{cvec}c)}\quad\begin{minipage}[c]{5.3cm}\raggedright\Large $(3{,}5\,;\,0{,}5)$\end{minipage}\end{tcolorbox}\hfill\begin{tcolorbox}[enhanced,nobeforeafter,width=7.1cm,height=2.1cm,valign=center,colback=cvec!6,colframe=cvec,arc=2.5mm,boxrule=0.9pt,left=3mm]{\Large\bfseries\color{cvec}d)}\quad\begin{minipage}[c]{5.3cm}\raggedright\Large $(3\,;\,6)$\end{minipage}\end{tcolorbox}
\end{tcolorbox}
\newpage
\begin{tcolorbox}[enhanced,width=16cm,height=14cm,colback=white,colframe=cvec,boxrule=2pt,arc=0mm,
left=6mm,right=6mm,top=0mm,bottom=4mm,sharp corners,
overlay={\fill[cvec] (frame.north west) rectangle ([yshift=-1.25cm]frame.north east);
\node[anchor=west,text=white,font=\sffamily\bfseries\LARGE] at ([xshift=5mm,yshift=-0.62cm]frame.north west) {Question 63 B};
\node[anchor=east,text=white,font=\sffamily\itshape\large] at ([xshift=-5mm,yshift=-0.62cm]frame.north east) {Vecteurs et repérage};}]
\vspace*{1.75cm}
\begin{minipage}[t][6.9cm][c]{\linewidth}\LARGE\centering
\begin{minipage}{0.95\linewidth}\raggedright\linespread{1.25}\selectfont Quel est le couple solution du système $\left\{\begin{array}{l}x+2y=3\\3x-y=-5\end{array}\right.$~?\end{minipage}
\end{minipage}

\vfill
\begin{tcolorbox}[enhanced,nobeforeafter,width=7.1cm,height=2.1cm,valign=center,colback=cvec!6,colframe=cvec,arc=2.5mm,boxrule=0.9pt,left=3mm]{\Large\bfseries\color{cvec}a)}\quad\begin{minipage}[c]{5.3cm}\raggedright\Large $(2\,;\,-1)$\end{minipage}\end{tcolorbox}\hfill\begin{tcolorbox}[enhanced,nobeforeafter,width=7.1cm,height=2.1cm,valign=center,colback=cvec!6,colframe=cvec,arc=2.5mm,boxrule=0.9pt,left=3mm]{\Large\bfseries\color{cvec}b)}\quad\begin{minipage}[c]{5.3cm}\raggedright\Large $(1\,;\,2)$\end{minipage}\end{tcolorbox}\par\vspace{2.5mm}
\begin{tcolorbox}[enhanced,nobeforeafter,width=7.1cm,height=2.1cm,valign=center,colback=cvec!6,colframe=cvec,arc=2.5mm,boxrule=0.9pt,left=3mm]{\Large\bfseries\color{cvec}c)}\quad\begin{minipage}[c]{5.3cm}\raggedright\Large $(-1\,;\,2)$\end{minipage}\end{tcolorbox}\hfill\begin{tcolorbox}[enhanced,nobeforeafter,width=7.1cm,height=2.1cm,valign=center,colback=cvec!6,colframe=cvec,arc=2.5mm,boxrule=0.9pt,left=3mm]{\Large\bfseries\color{cvec}d)}\quad\begin{minipage}[c]{5.3cm}\raggedright\Large $(-1\,;\,-2)$\end{minipage}\end{tcolorbox}
\end{tcolorbox}
\newpage
\begin{tcolorbox}[enhanced,width=16cm,height=14cm,colback=white,colframe=cvec,boxrule=2pt,arc=0mm,
left=6mm,right=6mm,top=0mm,bottom=4mm,sharp corners,
overlay={\fill[cvec] (frame.north west) rectangle ([yshift=-1.25cm]frame.north east);
\node[anchor=west,text=white,font=\sffamily\bfseries\LARGE] at ([xshift=5mm,yshift=-0.62cm]frame.north west) {Question 64 B};
\node[anchor=east,text=white,font=\sffamily\itshape\large] at ([xshift=-5mm,yshift=-0.62cm]frame.north east) {Vecteurs et repérage};}]
\vspace*{1.75cm}
\begin{minipage}[t][6.9cm][c]{\linewidth}\LARGE\centering
\begin{minipage}{0.95\linewidth}\raggedright\linespread{1.25}\selectfont Quel est le couple solution du système $\left\{\begin{array}{l}3x+y=9\\x+y=1\end{array}\right.$~?\end{minipage}
\end{minipage}

\vfill
\begin{tcolorbox}[enhanced,nobeforeafter,width=7.1cm,height=2.1cm,valign=center,colback=cvec!6,colframe=cvec,arc=2.5mm,boxrule=0.9pt,left=3mm]{\Large\bfseries\color{cvec}a)}\quad\begin{minipage}[c]{5.3cm}\raggedright\Large $(4\,;\,-3)$\end{minipage}\end{tcolorbox}\hfill\begin{tcolorbox}[enhanced,nobeforeafter,width=7.1cm,height=2.1cm,valign=center,colback=cvec!6,colframe=cvec,arc=2.5mm,boxrule=0.9pt,left=3mm]{\Large\bfseries\color{cvec}b)}\quad\begin{minipage}[c]{5.3cm}\raggedright\Large $(-3\,;\,4)$\end{minipage}\end{tcolorbox}\par\vspace{2.5mm}
\begin{tcolorbox}[enhanced,nobeforeafter,width=7.1cm,height=2.1cm,valign=center,colback=cvec!6,colframe=cvec,arc=2.5mm,boxrule=0.9pt,left=3mm]{\Large\bfseries\color{cvec}c)}\quad\begin{minipage}[c]{5.3cm}\raggedright\Large $(4\,;\,3)$\end{minipage}\end{tcolorbox}\hfill\begin{tcolorbox}[enhanced,nobeforeafter,width=7.1cm,height=2.1cm,valign=center,colback=cvec!6,colframe=cvec,arc=2.5mm,boxrule=0.9pt,left=3mm]{\Large\bfseries\color{cvec}d)}\quad\begin{minipage}[c]{5.3cm}\raggedright\Large $(-4\,;\,-3)$\end{minipage}\end{tcolorbox}
\end{tcolorbox}
\newpage
\begin{tcolorbox}[enhanced,width=16cm,height=14cm,colback=white,colframe=cari,boxrule=2pt,arc=0mm,
left=6mm,right=6mm,top=0mm,bottom=4mm,sharp corners,
overlay={\fill[cari] (frame.north west) rectangle ([yshift=-1.25cm]frame.north east);
\node[anchor=west,text=white,font=\sffamily\bfseries\LARGE] at ([xshift=5mm,yshift=-0.62cm]frame.north west) {Question 65 B};
\node[anchor=east,text=white,font=\sffamily\itshape\large] at ([xshift=-5mm,yshift=-0.62cm]frame.north east) {Arithmétique et ensembles de nombres};}]
\vspace*{1.75cm}
\begin{minipage}[t][6.9cm][c]{\linewidth}\LARGE\centering
\begin{minipage}{0.95\linewidth}\raggedright\linespread{1.25}\selectfont Lequel de ces nombres est premier~?\end{minipage}
\end{minipage}

\vfill
\begin{tcolorbox}[enhanced,nobeforeafter,width=7.1cm,height=2.1cm,valign=center,colback=cari!6,colframe=cari,arc=2.5mm,boxrule=0.9pt,left=3mm]{\Large\bfseries\color{cari}a)}\quad\begin{minipage}[c]{5.3cm}\raggedright\Large $81$\end{minipage}\end{tcolorbox}\hfill\begin{tcolorbox}[enhanced,nobeforeafter,width=7.1cm,height=2.1cm,valign=center,colback=cari!6,colframe=cari,arc=2.5mm,boxrule=0.9pt,left=3mm]{\Large\bfseries\color{cari}b)}\quad\begin{minipage}[c]{5.3cm}\raggedright\Large $87$\end{minipage}\end{tcolorbox}\par\vspace{2.5mm}
\begin{tcolorbox}[enhanced,nobeforeafter,width=7.1cm,height=2.1cm,valign=center,colback=cari!6,colframe=cari,arc=2.5mm,boxrule=0.9pt,left=3mm]{\Large\bfseries\color{cari}c)}\quad\begin{minipage}[c]{5.3cm}\raggedright\Large $89$\end{minipage}\end{tcolorbox}\hfill\begin{tcolorbox}[enhanced,nobeforeafter,width=7.1cm,height=2.1cm,valign=center,colback=cari!6,colframe=cari,arc=2.5mm,boxrule=0.9pt,left=3mm]{\Large\bfseries\color{cari}d)}\quad\begin{minipage}[c]{5.3cm}\raggedright\Large $91$\end{minipage}\end{tcolorbox}
\end{tcolorbox}
\newpage
\begin{tcolorbox}[enhanced,width=16cm,height=14cm,colback=white,colframe=cari,boxrule=2pt,arc=0mm,
left=6mm,right=6mm,top=0mm,bottom=4mm,sharp corners,
overlay={\fill[cari] (frame.north west) rectangle ([yshift=-1.25cm]frame.north east);
\node[anchor=west,text=white,font=\sffamily\bfseries\LARGE] at ([xshift=5mm,yshift=-0.62cm]frame.north west) {Question 66 B};
\node[anchor=east,text=white,font=\sffamily\itshape\large] at ([xshift=-5mm,yshift=-0.62cm]frame.north east) {Arithmétique et ensembles de nombres};}]
\vspace*{1.75cm}
\begin{minipage}[t][6.9cm][c]{\linewidth}\LARGE\centering
\begin{minipage}{0.95\linewidth}\raggedright\linespread{1.25}\selectfont Lequel de ces nombres est premier~?\end{minipage}
\end{minipage}

\vfill
\begin{tcolorbox}[enhanced,nobeforeafter,width=7.1cm,height=2.1cm,valign=center,colback=cari!6,colframe=cari,arc=2.5mm,boxrule=0.9pt,left=3mm]{\Large\bfseries\color{cari}a)}\quad\begin{minipage}[c]{5.3cm}\raggedright\Large $61$\end{minipage}\end{tcolorbox}\hfill\begin{tcolorbox}[enhanced,nobeforeafter,width=7.1cm,height=2.1cm,valign=center,colback=cari!6,colframe=cari,arc=2.5mm,boxrule=0.9pt,left=3mm]{\Large\bfseries\color{cari}b)}\quad\begin{minipage}[c]{5.3cm}\raggedright\Large $63$\end{minipage}\end{tcolorbox}\par\vspace{2.5mm}
\begin{tcolorbox}[enhanced,nobeforeafter,width=7.1cm,height=2.1cm,valign=center,colback=cari!6,colframe=cari,arc=2.5mm,boxrule=0.9pt,left=3mm]{\Large\bfseries\color{cari}c)}\quad\begin{minipage}[c]{5.3cm}\raggedright\Large $57$\end{minipage}\end{tcolorbox}\hfill\begin{tcolorbox}[enhanced,nobeforeafter,width=7.1cm,height=2.1cm,valign=center,colback=cari!6,colframe=cari,arc=2.5mm,boxrule=0.9pt,left=3mm]{\Large\bfseries\color{cari}d)}\quad\begin{minipage}[c]{5.3cm}\raggedright\Large $51$\end{minipage}\end{tcolorbox}
\end{tcolorbox}
\newpage
\begin{tcolorbox}[enhanced,width=16cm,height=14cm,colback=white,colframe=cari,boxrule=2pt,arc=0mm,
left=6mm,right=6mm,top=0mm,bottom=4mm,sharp corners,
overlay={\fill[cari] (frame.north west) rectangle ([yshift=-1.25cm]frame.north east);
\node[anchor=west,text=white,font=\sffamily\bfseries\LARGE] at ([xshift=5mm,yshift=-0.62cm]frame.north west) {Question 67 B};
\node[anchor=east,text=white,font=\sffamily\itshape\large] at ([xshift=-5mm,yshift=-0.62cm]frame.north east) {Arithmétique et ensembles de nombres};}]
\vspace*{1.75cm}
\begin{minipage}[t][6.9cm][c]{\linewidth}\LARGE\centering
\begin{minipage}{0.95\linewidth}\raggedright\linespread{1.25}\selectfont Quelle est la décomposition de $150$ en produit de facteurs premiers~?\end{minipage}
\end{minipage}

\vfill
\begin{tcolorbox}[enhanced,nobeforeafter,width=7.1cm,height=2.1cm,valign=center,colback=cari!6,colframe=cari,arc=2.5mm,boxrule=0.9pt,left=3mm]{\Large\bfseries\color{cari}a)}\quad\begin{minipage}[c]{5.3cm}\raggedright\Large $2^{2}\times3\times5^{2}$\end{minipage}\end{tcolorbox}\hfill\begin{tcolorbox}[enhanced,nobeforeafter,width=7.1cm,height=2.1cm,valign=center,colback=cari!6,colframe=cari,arc=2.5mm,boxrule=0.9pt,left=3mm]{\Large\bfseries\color{cari}b)}\quad\begin{minipage}[c]{5.3cm}\raggedright\Large $2\times3\times5^{2}$\end{minipage}\end{tcolorbox}\par\vspace{2.5mm}
\begin{tcolorbox}[enhanced,nobeforeafter,width=7.1cm,height=2.1cm,valign=center,colback=cari!6,colframe=cari,arc=2.5mm,boxrule=0.9pt,left=3mm]{\Large\bfseries\color{cari}c)}\quad\begin{minipage}[c]{5.3cm}\raggedright\Large $2\times3\times5^{3}$\end{minipage}\end{tcolorbox}\hfill\begin{tcolorbox}[enhanced,nobeforeafter,width=7.1cm,height=2.1cm,valign=center,colback=cari!6,colframe=cari,arc=2.5mm,boxrule=0.9pt,left=3mm]{\Large\bfseries\color{cari}d)}\quad\begin{minipage}[c]{5.3cm}\raggedright\Large $6\times25$\end{minipage}\end{tcolorbox}
\end{tcolorbox}
\newpage
\begin{tcolorbox}[enhanced,width=16cm,height=14cm,colback=white,colframe=cari,boxrule=2pt,arc=0mm,
left=6mm,right=6mm,top=0mm,bottom=4mm,sharp corners,
overlay={\fill[cari] (frame.north west) rectangle ([yshift=-1.25cm]frame.north east);
\node[anchor=west,text=white,font=\sffamily\bfseries\LARGE] at ([xshift=5mm,yshift=-0.62cm]frame.north west) {Question 68 B};
\node[anchor=east,text=white,font=\sffamily\itshape\large] at ([xshift=-5mm,yshift=-0.62cm]frame.north east) {Arithmétique et ensembles de nombres};}]
\vspace*{1.75cm}
\begin{minipage}[t][6.9cm][c]{\linewidth}\LARGE\centering
\begin{minipage}{0.95\linewidth}\raggedright\linespread{1.25}\selectfont Quelle est la décomposition de $1\,000$ en produit de facteurs premiers~?\end{minipage}
\end{minipage}

\vfill
\begin{tcolorbox}[enhanced,nobeforeafter,width=7.1cm,height=2.1cm,valign=center,colback=cari!6,colframe=cari,arc=2.5mm,boxrule=0.9pt,left=3mm]{\Large\bfseries\color{cari}a)}\quad\begin{minipage}[c]{5.3cm}\raggedright\Large $2^{3}\times5^{3}$\end{minipage}\end{tcolorbox}\hfill\begin{tcolorbox}[enhanced,nobeforeafter,width=7.1cm,height=2.1cm,valign=center,colback=cari!6,colframe=cari,arc=2.5mm,boxrule=0.9pt,left=3mm]{\Large\bfseries\color{cari}b)}\quad\begin{minipage}[c]{5.3cm}\raggedright\Large $2^{3}\times5^{4}$\end{minipage}\end{tcolorbox}\par\vspace{2.5mm}
\begin{tcolorbox}[enhanced,nobeforeafter,width=7.1cm,height=2.1cm,valign=center,colback=cari!6,colframe=cari,arc=2.5mm,boxrule=0.9pt,left=3mm]{\Large\bfseries\color{cari}c)}\quad\begin{minipage}[c]{5.3cm}\raggedright\Large $40\times25$\end{minipage}\end{tcolorbox}\hfill\begin{tcolorbox}[enhanced,nobeforeafter,width=7.1cm,height=2.1cm,valign=center,colback=cari!6,colframe=cari,arc=2.5mm,boxrule=0.9pt,left=3mm]{\Large\bfseries\color{cari}d)}\quad\begin{minipage}[c]{5.3cm}\raggedright\Large $2^{4}\times5^{3}$\end{minipage}\end{tcolorbox}
\end{tcolorbox}
\newpage
\begin{tcolorbox}[enhanced,width=16cm,height=14cm,colback=white,colframe=cari,boxrule=2pt,arc=0mm,
left=6mm,right=6mm,top=0mm,bottom=4mm,sharp corners,
overlay={\fill[cari] (frame.north west) rectangle ([yshift=-1.25cm]frame.north east);
\node[anchor=west,text=white,font=\sffamily\bfseries\LARGE] at ([xshift=5mm,yshift=-0.62cm]frame.north west) {Question 69 B};
\node[anchor=east,text=white,font=\sffamily\itshape\large] at ([xshift=-5mm,yshift=-0.62cm]frame.north east) {Arithmétique et ensembles de nombres};}]
\vspace*{1.75cm}
\begin{minipage}[t][6.9cm][c]{\linewidth}\LARGE\centering
\begin{minipage}{0.95\linewidth}\raggedright\linespread{1.25}\selectfont Quel est le plus petit ensemble de nombres auquel appartient $\dfrac{1}{3}$~?\end{minipage}
\end{minipage}

\vfill
\begin{tcolorbox}[enhanced,nobeforeafter,width=7.1cm,height=2.1cm,valign=center,colback=cari!6,colframe=cari,arc=2.5mm,boxrule=0.9pt,left=3mm]{\Large\bfseries\color{cari}a)}\quad\begin{minipage}[c]{5.3cm}\raggedright\Large $\mathbb{D}$\end{minipage}\end{tcolorbox}\hfill\begin{tcolorbox}[enhanced,nobeforeafter,width=7.1cm,height=2.1cm,valign=center,colback=cari!6,colframe=cari,arc=2.5mm,boxrule=0.9pt,left=3mm]{\Large\bfseries\color{cari}b)}\quad\begin{minipage}[c]{5.3cm}\raggedright\Large $\mathbb{N}$\end{minipage}\end{tcolorbox}\par\vspace{2.5mm}
\begin{tcolorbox}[enhanced,nobeforeafter,width=7.1cm,height=2.1cm,valign=center,colback=cari!6,colframe=cari,arc=2.5mm,boxrule=0.9pt,left=3mm]{\Large\bfseries\color{cari}c)}\quad\begin{minipage}[c]{5.3cm}\raggedright\Large $\mathbb{Z}$\end{minipage}\end{tcolorbox}\hfill\begin{tcolorbox}[enhanced,nobeforeafter,width=7.1cm,height=2.1cm,valign=center,colback=cari!6,colframe=cari,arc=2.5mm,boxrule=0.9pt,left=3mm]{\Large\bfseries\color{cari}d)}\quad\begin{minipage}[c]{5.3cm}\raggedright\Large $\mathbb{Q}$\end{minipage}\end{tcolorbox}
\end{tcolorbox}
\newpage
\begin{tcolorbox}[enhanced,width=16cm,height=14cm,colback=white,colframe=cari,boxrule=2pt,arc=0mm,
left=6mm,right=6mm,top=0mm,bottom=4mm,sharp corners,
overlay={\fill[cari] (frame.north west) rectangle ([yshift=-1.25cm]frame.north east);
\node[anchor=west,text=white,font=\sffamily\bfseries\LARGE] at ([xshift=5mm,yshift=-0.62cm]frame.north west) {Question 70 B};
\node[anchor=east,text=white,font=\sffamily\itshape\large] at ([xshift=-5mm,yshift=-0.62cm]frame.north east) {Arithmétique et ensembles de nombres};}]
\vspace*{1.75cm}
\begin{minipage}[t][6.9cm][c]{\linewidth}\LARGE\centering
\begin{minipage}{0.95\linewidth}\raggedright\linespread{1.25}\selectfont Quel est le plus petit ensemble de nombres auquel appartient $\sqrt{2}$~?\end{minipage}
\end{minipage}

\vfill
\begin{tcolorbox}[enhanced,nobeforeafter,width=7.1cm,height=2.1cm,valign=center,colback=cari!6,colframe=cari,arc=2.5mm,boxrule=0.9pt,left=3mm]{\Large\bfseries\color{cari}a)}\quad\begin{minipage}[c]{5.3cm}\raggedright\Large $\mathbb{R}$\end{minipage}\end{tcolorbox}\hfill\begin{tcolorbox}[enhanced,nobeforeafter,width=7.1cm,height=2.1cm,valign=center,colback=cari!6,colframe=cari,arc=2.5mm,boxrule=0.9pt,left=3mm]{\Large\bfseries\color{cari}b)}\quad\begin{minipage}[c]{5.3cm}\raggedright\Large $\mathbb{D}$\end{minipage}\end{tcolorbox}\par\vspace{2.5mm}
\begin{tcolorbox}[enhanced,nobeforeafter,width=7.1cm,height=2.1cm,valign=center,colback=cari!6,colframe=cari,arc=2.5mm,boxrule=0.9pt,left=3mm]{\Large\bfseries\color{cari}c)}\quad\begin{minipage}[c]{5.3cm}\raggedright\Large $\mathbb{Z}$\end{minipage}\end{tcolorbox}\hfill\begin{tcolorbox}[enhanced,nobeforeafter,width=7.1cm,height=2.1cm,valign=center,colback=cari!6,colframe=cari,arc=2.5mm,boxrule=0.9pt,left=3mm]{\Large\bfseries\color{cari}d)}\quad\begin{minipage}[c]{5.3cm}\raggedright\Large $\mathbb{N}$\end{minipage}\end{tcolorbox}
\end{tcolorbox}
\newpage
\begin{tcolorbox}[enhanced,width=16cm,height=14cm,colback=white,colframe=cari,boxrule=2pt,arc=0mm,
left=6mm,right=6mm,top=0mm,bottom=4mm,sharp corners,
overlay={\fill[cari] (frame.north west) rectangle ([yshift=-1.25cm]frame.north east);
\node[anchor=west,text=white,font=\sffamily\bfseries\LARGE] at ([xshift=5mm,yshift=-0.62cm]frame.north west) {Question 71 B};
\node[anchor=east,text=white,font=\sffamily\itshape\large] at ([xshift=-5mm,yshift=-0.62cm]frame.north east) {Arithmétique et ensembles de nombres};}]
\vspace*{1.75cm}
\begin{minipage}[t][6.9cm][c]{\linewidth}\LARGE\centering
\begin{minipage}{0.95\linewidth}\raggedright\linespread{1.25}\selectfont Quelle est la forme irréductible de $\dfrac{60}{144}$~?\end{minipage}
\end{minipage}

\vfill
\begin{tcolorbox}[enhanced,nobeforeafter,width=7.1cm,height=2.1cm,valign=center,colback=cari!6,colframe=cari,arc=2.5mm,boxrule=0.9pt,left=3mm]{\Large\bfseries\color{cari}a)}\quad\begin{minipage}[c]{5.3cm}\raggedright\Large $\dfrac{30}{72}$\end{minipage}\end{tcolorbox}\hfill\begin{tcolorbox}[enhanced,nobeforeafter,width=7.1cm,height=2.1cm,valign=center,colback=cari!6,colframe=cari,arc=2.5mm,boxrule=0.9pt,left=3mm]{\Large\bfseries\color{cari}b)}\quad\begin{minipage}[c]{5.3cm}\raggedright\Large $\dfrac{20}{48}$\end{minipage}\end{tcolorbox}\par\vspace{2.5mm}
\begin{tcolorbox}[enhanced,nobeforeafter,width=7.1cm,height=2.1cm,valign=center,colback=cari!6,colframe=cari,arc=2.5mm,boxrule=0.9pt,left=3mm]{\Large\bfseries\color{cari}c)}\quad\begin{minipage}[c]{5.3cm}\raggedright\Large $\dfrac{12}{5}$\end{minipage}\end{tcolorbox}\hfill\begin{tcolorbox}[enhanced,nobeforeafter,width=7.1cm,height=2.1cm,valign=center,colback=cari!6,colframe=cari,arc=2.5mm,boxrule=0.9pt,left=3mm]{\Large\bfseries\color{cari}d)}\quad\begin{minipage}[c]{5.3cm}\raggedright\Large $\dfrac{5}{12}$\end{minipage}\end{tcolorbox}
\end{tcolorbox}
\newpage
\begin{tcolorbox}[enhanced,width=16cm,height=14cm,colback=white,colframe=cari,boxrule=2pt,arc=0mm,
left=6mm,right=6mm,top=0mm,bottom=4mm,sharp corners,
overlay={\fill[cari] (frame.north west) rectangle ([yshift=-1.25cm]frame.north east);
\node[anchor=west,text=white,font=\sffamily\bfseries\LARGE] at ([xshift=5mm,yshift=-0.62cm]frame.north west) {Question 72 B};
\node[anchor=east,text=white,font=\sffamily\itshape\large] at ([xshift=-5mm,yshift=-0.62cm]frame.north east) {Arithmétique et ensembles de nombres};}]
\vspace*{1.75cm}
\begin{minipage}[t][6.9cm][c]{\linewidth}\LARGE\centering
\begin{minipage}{0.95\linewidth}\raggedright\linespread{1.25}\selectfont Quelle est la forme irréductible de $\dfrac{98}{70}$~?\end{minipage}
\end{minipage}

\vfill
\begin{tcolorbox}[enhanced,nobeforeafter,width=7.1cm,height=2.1cm,valign=center,colback=cari!6,colframe=cari,arc=2.5mm,boxrule=0.9pt,left=3mm]{\Large\bfseries\color{cari}a)}\quad\begin{minipage}[c]{5.3cm}\raggedright\Large $\dfrac{49}{35}$\end{minipage}\end{tcolorbox}\hfill\begin{tcolorbox}[enhanced,nobeforeafter,width=7.1cm,height=2.1cm,valign=center,colback=cari!6,colframe=cari,arc=2.5mm,boxrule=0.9pt,left=3mm]{\Large\bfseries\color{cari}b)}\quad\begin{minipage}[c]{5.3cm}\raggedright\Large $\dfrac{14}{10}$\end{minipage}\end{tcolorbox}\par\vspace{2.5mm}
\begin{tcolorbox}[enhanced,nobeforeafter,width=7.1cm,height=2.1cm,valign=center,colback=cari!6,colframe=cari,arc=2.5mm,boxrule=0.9pt,left=3mm]{\Large\bfseries\color{cari}c)}\quad\begin{minipage}[c]{5.3cm}\raggedright\Large $\dfrac{5}{7}$\end{minipage}\end{tcolorbox}\hfill\begin{tcolorbox}[enhanced,nobeforeafter,width=7.1cm,height=2.1cm,valign=center,colback=cari!6,colframe=cari,arc=2.5mm,boxrule=0.9pt,left=3mm]{\Large\bfseries\color{cari}d)}\quad\begin{minipage}[c]{5.3cm}\raggedright\Large $\dfrac{7}{5}$\end{minipage}\end{tcolorbox}
\end{tcolorbox}
\newpage
\begin{tcolorbox}[enhanced,width=16cm,height=14cm,colback=white,colframe=cari,boxrule=2pt,arc=0mm,
left=6mm,right=6mm,top=0mm,bottom=4mm,sharp corners,
overlay={\fill[cari] (frame.north west) rectangle ([yshift=-1.25cm]frame.north east);
\node[anchor=west,text=white,font=\sffamily\bfseries\LARGE] at ([xshift=5mm,yshift=-0.62cm]frame.north west) {Question 73 B};
\node[anchor=east,text=white,font=\sffamily\itshape\large] at ([xshift=-5mm,yshift=-0.62cm]frame.north east) {Arithmétique et ensembles de nombres};}]
\vspace*{1.75cm}
\begin{minipage}[t][6.9cm][c]{\linewidth}\LARGE\centering
\begin{minipage}{0.95\linewidth}\raggedright\linespread{1.25}\selectfont Pour tout entier naturel $n$, que peut-on dire de l'entier $n(n+1)$~?\end{minipage}
\end{minipage}

\vfill
\begin{tcolorbox}[enhanced,nobeforeafter,width=7.1cm,height=2.1cm,valign=center,colback=cari!6,colframe=cari,arc=2.5mm,boxrule=0.9pt,left=3mm]{\Large\bfseries\color{cari}a)}\quad\begin{minipage}[c]{5.3cm}\raggedright\Large il est toujours impair\end{minipage}\end{tcolorbox}\hfill\begin{tcolorbox}[enhanced,nobeforeafter,width=7.1cm,height=2.1cm,valign=center,colback=cari!6,colframe=cari,arc=2.5mm,boxrule=0.9pt,left=3mm]{\Large\bfseries\color{cari}b)}\quad\begin{minipage}[c]{5.3cm}\raggedright\Large il est toujours pair\end{minipage}\end{tcolorbox}\par\vspace{2.5mm}
\begin{tcolorbox}[enhanced,nobeforeafter,width=7.1cm,height=2.1cm,valign=center,colback=cari!6,colframe=cari,arc=2.5mm,boxrule=0.9pt,left=3mm]{\Large\bfseries\color{cari}c)}\quad\begin{minipage}[c]{5.3cm}\raggedright\large il est pair ou impair selon $n$\end{minipage}\end{tcolorbox}\hfill\begin{tcolorbox}[enhanced,nobeforeafter,width=7.1cm,height=2.1cm,valign=center,colback=cari!6,colframe=cari,arc=2.5mm,boxrule=0.9pt,left=3mm]{\Large\bfseries\color{cari}d)}\quad\begin{minipage}[c]{5.3cm}\raggedright\Large il est toujours premier\end{minipage}\end{tcolorbox}
\end{tcolorbox}
\newpage
\begin{tcolorbox}[enhanced,width=16cm,height=14cm,colback=white,colframe=cari,boxrule=2pt,arc=0mm,
left=6mm,right=6mm,top=0mm,bottom=4mm,sharp corners,
overlay={\fill[cari] (frame.north west) rectangle ([yshift=-1.25cm]frame.north east);
\node[anchor=west,text=white,font=\sffamily\bfseries\LARGE] at ([xshift=5mm,yshift=-0.62cm]frame.north west) {Question 74 B};
\node[anchor=east,text=white,font=\sffamily\itshape\large] at ([xshift=-5mm,yshift=-0.62cm]frame.north east) {Arithmétique et ensembles de nombres};}]
\vspace*{1.75cm}
\begin{minipage}[t][6.9cm][c]{\linewidth}\LARGE\centering
\begin{minipage}{0.95\linewidth}\raggedright\linespread{1.25}\selectfont Pour tout entier naturel $n$, que peut-on dire de l'entier $2n+3$~?\end{minipage}
\end{minipage}

\vfill
\begin{tcolorbox}[enhanced,nobeforeafter,width=7.1cm,height=2.1cm,valign=center,colback=cari!6,colframe=cari,arc=2.5mm,boxrule=0.9pt,left=3mm]{\Large\bfseries\color{cari}a)}\quad\begin{minipage}[c]{5.3cm}\raggedright\Large il est toujours impair\end{minipage}\end{tcolorbox}\hfill\begin{tcolorbox}[enhanced,nobeforeafter,width=7.1cm,height=2.1cm,valign=center,colback=cari!6,colframe=cari,arc=2.5mm,boxrule=0.9pt,left=3mm]{\Large\bfseries\color{cari}b)}\quad\begin{minipage}[c]{5.3cm}\raggedright\Large il est toujours premier\end{minipage}\end{tcolorbox}\par\vspace{2.5mm}
\begin{tcolorbox}[enhanced,nobeforeafter,width=7.1cm,height=2.1cm,valign=center,colback=cari!6,colframe=cari,arc=2.5mm,boxrule=0.9pt,left=3mm]{\Large\bfseries\color{cari}c)}\quad\begin{minipage}[c]{5.3cm}\raggedright\large il est pair ou impair selon $n$\end{minipage}\end{tcolorbox}\hfill\begin{tcolorbox}[enhanced,nobeforeafter,width=7.1cm,height=2.1cm,valign=center,colback=cari!6,colframe=cari,arc=2.5mm,boxrule=0.9pt,left=3mm]{\Large\bfseries\color{cari}d)}\quad\begin{minipage}[c]{5.3cm}\raggedright\Large il est toujours pair\end{minipage}\end{tcolorbox}
\end{tcolorbox}
\newpage
\begin{tcolorbox}[enhanced,width=16cm,height=14cm,colback=white,colframe=cari,boxrule=2pt,arc=0mm,
left=6mm,right=6mm,top=0mm,bottom=4mm,sharp corners,
overlay={\fill[cari] (frame.north west) rectangle ([yshift=-1.25cm]frame.north east);
\node[anchor=west,text=white,font=\sffamily\bfseries\LARGE] at ([xshift=5mm,yshift=-0.62cm]frame.north west) {Question 75 B};
\node[anchor=east,text=white,font=\sffamily\itshape\large] at ([xshift=-5mm,yshift=-0.62cm]frame.north east) {Arithmétique et ensembles de nombres};}]
\vspace*{1.75cm}
\begin{minipage}[t][6.9cm][c]{\linewidth}\LARGE\centering
\begin{minipage}{0.95\linewidth}\raggedright\linespread{1.25}\selectfont Combien y a-t-il de multiples de $12$ compris entre $1$ et $150$~?\end{minipage}
\end{minipage}

\vfill
\begin{tcolorbox}[enhanced,nobeforeafter,width=7.1cm,height=2.1cm,valign=center,colback=cari!6,colframe=cari,arc=2.5mm,boxrule=0.9pt,left=3mm]{\Large\bfseries\color{cari}a)}\quad\begin{minipage}[c]{5.3cm}\raggedright\Large $12$\end{minipage}\end{tcolorbox}\hfill\begin{tcolorbox}[enhanced,nobeforeafter,width=7.1cm,height=2.1cm,valign=center,colback=cari!6,colframe=cari,arc=2.5mm,boxrule=0.9pt,left=3mm]{\Large\bfseries\color{cari}b)}\quad\begin{minipage}[c]{5.3cm}\raggedright\Large $13$\end{minipage}\end{tcolorbox}\par\vspace{2.5mm}
\begin{tcolorbox}[enhanced,nobeforeafter,width=7.1cm,height=2.1cm,valign=center,colback=cari!6,colframe=cari,arc=2.5mm,boxrule=0.9pt,left=3mm]{\Large\bfseries\color{cari}c)}\quad\begin{minipage}[c]{5.3cm}\raggedright\Large $11$\end{minipage}\end{tcolorbox}\hfill\begin{tcolorbox}[enhanced,nobeforeafter,width=7.1cm,height=2.1cm,valign=center,colback=cari!6,colframe=cari,arc=2.5mm,boxrule=0.9pt,left=3mm]{\Large\bfseries\color{cari}d)}\quad\begin{minipage}[c]{5.3cm}\raggedright\Large $138$\end{minipage}\end{tcolorbox}
\end{tcolorbox}
\newpage
\begin{tcolorbox}[enhanced,width=16cm,height=14cm,colback=white,colframe=cari,boxrule=2pt,arc=0mm,
left=6mm,right=6mm,top=0mm,bottom=4mm,sharp corners,
overlay={\fill[cari] (frame.north west) rectangle ([yshift=-1.25cm]frame.north east);
\node[anchor=west,text=white,font=\sffamily\bfseries\LARGE] at ([xshift=5mm,yshift=-0.62cm]frame.north west) {Question 76 B};
\node[anchor=east,text=white,font=\sffamily\itshape\large] at ([xshift=-5mm,yshift=-0.62cm]frame.north east) {Arithmétique et ensembles de nombres};}]
\vspace*{1.75cm}
\begin{minipage}[t][6.9cm][c]{\linewidth}\LARGE\centering
\begin{minipage}{0.95\linewidth}\raggedright\linespread{1.25}\selectfont Combien y a-t-il de multiples de $5$ compris entre $1$ et $73$~?\end{minipage}
\end{minipage}

\vfill
\begin{tcolorbox}[enhanced,nobeforeafter,width=7.1cm,height=2.1cm,valign=center,colback=cari!6,colframe=cari,arc=2.5mm,boxrule=0.9pt,left=3mm]{\Large\bfseries\color{cari}a)}\quad\begin{minipage}[c]{5.3cm}\raggedright\Large $15$\end{minipage}\end{tcolorbox}\hfill\begin{tcolorbox}[enhanced,nobeforeafter,width=7.1cm,height=2.1cm,valign=center,colback=cari!6,colframe=cari,arc=2.5mm,boxrule=0.9pt,left=3mm]{\Large\bfseries\color{cari}b)}\quad\begin{minipage}[c]{5.3cm}\raggedright\Large $14$\end{minipage}\end{tcolorbox}\par\vspace{2.5mm}
\begin{tcolorbox}[enhanced,nobeforeafter,width=7.1cm,height=2.1cm,valign=center,colback=cari!6,colframe=cari,arc=2.5mm,boxrule=0.9pt,left=3mm]{\Large\bfseries\color{cari}c)}\quad\begin{minipage}[c]{5.3cm}\raggedright\Large $59$\end{minipage}\end{tcolorbox}\hfill\begin{tcolorbox}[enhanced,nobeforeafter,width=7.1cm,height=2.1cm,valign=center,colback=cari!6,colframe=cari,arc=2.5mm,boxrule=0.9pt,left=3mm]{\Large\bfseries\color{cari}d)}\quad\begin{minipage}[c]{5.3cm}\raggedright\Large $13$\end{minipage}\end{tcolorbox}
\end{tcolorbox}
\newpage
\begin{tcolorbox}[enhanced,width=16cm,height=14cm,colback=white,colframe=cari,boxrule=2pt,arc=0mm,
left=6mm,right=6mm,top=0mm,bottom=4mm,sharp corners,
overlay={\fill[cari] (frame.north west) rectangle ([yshift=-1.25cm]frame.north east);
\node[anchor=west,text=white,font=\sffamily\bfseries\LARGE] at ([xshift=5mm,yshift=-0.62cm]frame.north west) {Question 77 B};
\node[anchor=east,text=white,font=\sffamily\itshape\large] at ([xshift=-5mm,yshift=-0.62cm]frame.north east) {Arithmétique et ensembles de nombres};}]
\vspace*{1.75cm}
\begin{minipage}[t][6.9cm][c]{\linewidth}\LARGE\centering
\begin{minipage}{0.95\linewidth}\raggedright\linespread{1.25}\selectfont Combien l'entier $49$ a-t-il de diviseurs positifs~?\end{minipage}
\end{minipage}

\vfill
\begin{tcolorbox}[enhanced,nobeforeafter,width=7.1cm,height=2.1cm,valign=center,colback=cari!6,colframe=cari,arc=2.5mm,boxrule=0.9pt,left=3mm]{\Large\bfseries\color{cari}a)}\quad\begin{minipage}[c]{5.3cm}\raggedright\Large $2$\end{minipage}\end{tcolorbox}\hfill\begin{tcolorbox}[enhanced,nobeforeafter,width=7.1cm,height=2.1cm,valign=center,colback=cari!6,colframe=cari,arc=2.5mm,boxrule=0.9pt,left=3mm]{\Large\bfseries\color{cari}b)}\quad\begin{minipage}[c]{5.3cm}\raggedright\Large $5$\end{minipage}\end{tcolorbox}\par\vspace{2.5mm}
\begin{tcolorbox}[enhanced,nobeforeafter,width=7.1cm,height=2.1cm,valign=center,colback=cari!6,colframe=cari,arc=2.5mm,boxrule=0.9pt,left=3mm]{\Large\bfseries\color{cari}c)}\quad\begin{minipage}[c]{5.3cm}\raggedright\Large $3$\end{minipage}\end{tcolorbox}\hfill\begin{tcolorbox}[enhanced,nobeforeafter,width=7.1cm,height=2.1cm,valign=center,colback=cari!6,colframe=cari,arc=2.5mm,boxrule=0.9pt,left=3mm]{\Large\bfseries\color{cari}d)}\quad\begin{minipage}[c]{5.3cm}\raggedright\Large $4$\end{minipage}\end{tcolorbox}
\end{tcolorbox}
\newpage
\begin{tcolorbox}[enhanced,width=16cm,height=14cm,colback=white,colframe=cari,boxrule=2pt,arc=0mm,
left=6mm,right=6mm,top=0mm,bottom=4mm,sharp corners,
overlay={\fill[cari] (frame.north west) rectangle ([yshift=-1.25cm]frame.north east);
\node[anchor=west,text=white,font=\sffamily\bfseries\LARGE] at ([xshift=5mm,yshift=-0.62cm]frame.north west) {Question 78 B};
\node[anchor=east,text=white,font=\sffamily\itshape\large] at ([xshift=-5mm,yshift=-0.62cm]frame.north east) {Arithmétique et ensembles de nombres};}]
\vspace*{1.75cm}
\begin{minipage}[t][6.9cm][c]{\linewidth}\LARGE\centering
\begin{minipage}{0.95\linewidth}\raggedright\linespread{1.25}\selectfont Combien l'entier $72$ a-t-il de diviseurs positifs~?\end{minipage}
\end{minipage}

\vfill
\begin{tcolorbox}[enhanced,nobeforeafter,width=7.1cm,height=2.1cm,valign=center,colback=cari!6,colframe=cari,arc=2.5mm,boxrule=0.9pt,left=3mm]{\Large\bfseries\color{cari}a)}\quad\begin{minipage}[c]{5.3cm}\raggedright\Large $14$\end{minipage}\end{tcolorbox}\hfill\begin{tcolorbox}[enhanced,nobeforeafter,width=7.1cm,height=2.1cm,valign=center,colback=cari!6,colframe=cari,arc=2.5mm,boxrule=0.9pt,left=3mm]{\Large\bfseries\color{cari}b)}\quad\begin{minipage}[c]{5.3cm}\raggedright\Large $12$\end{minipage}\end{tcolorbox}\par\vspace{2.5mm}
\begin{tcolorbox}[enhanced,nobeforeafter,width=7.1cm,height=2.1cm,valign=center,colback=cari!6,colframe=cari,arc=2.5mm,boxrule=0.9pt,left=3mm]{\Large\bfseries\color{cari}c)}\quad\begin{minipage}[c]{5.3cm}\raggedright\Large $13$\end{minipage}\end{tcolorbox}\hfill\begin{tcolorbox}[enhanced,nobeforeafter,width=7.1cm,height=2.1cm,valign=center,colback=cari!6,colframe=cari,arc=2.5mm,boxrule=0.9pt,left=3mm]{\Large\bfseries\color{cari}d)}\quad\begin{minipage}[c]{5.3cm}\raggedright\Large $11$\end{minipage}\end{tcolorbox}
\end{tcolorbox}
\newpage
\begin{tcolorbox}[enhanced,width=16cm,height=14cm,colback=white,colframe=cari,boxrule=2pt,arc=0mm,
left=6mm,right=6mm,top=0mm,bottom=4mm,sharp corners,
overlay={\fill[cari] (frame.north west) rectangle ([yshift=-1.25cm]frame.north east);
\node[anchor=west,text=white,font=\sffamily\bfseries\LARGE] at ([xshift=5mm,yshift=-0.62cm]frame.north west) {Question 79 B};
\node[anchor=east,text=white,font=\sffamily\itshape\large] at ([xshift=-5mm,yshift=-0.62cm]frame.north east) {Arithmétique et ensembles de nombres};}]
\vspace*{1.75cm}
\begin{minipage}[t][6.9cm][c]{\linewidth}\LARGE\centering
\begin{minipage}{0.95\linewidth}\raggedright\linespread{1.25}\selectfont Que vaut l'intersection $]-\infty\,;\,3]\cap[-1\,;\,+\infty[$~?\end{minipage}
\end{minipage}

\vfill
\begin{tcolorbox}[enhanced,nobeforeafter,width=7.1cm,height=2.1cm,valign=center,colback=cari!6,colframe=cari,arc=2.5mm,boxrule=0.9pt,left=3mm]{\Large\bfseries\color{cari}a)}\quad\begin{minipage}[c]{5.3cm}\raggedright\Large $[-1\,;\,3]$\end{minipage}\end{tcolorbox}\hfill\begin{tcolorbox}[enhanced,nobeforeafter,width=7.1cm,height=2.1cm,valign=center,colback=cari!6,colframe=cari,arc=2.5mm,boxrule=0.9pt,left=3mm]{\Large\bfseries\color{cari}b)}\quad\begin{minipage}[c]{5.3cm}\raggedright\Large $[3\,;\,+\infty[$\end{minipage}\end{tcolorbox}\par\vspace{2.5mm}
\begin{tcolorbox}[enhanced,nobeforeafter,width=7.1cm,height=2.1cm,valign=center,colback=cari!6,colframe=cari,arc=2.5mm,boxrule=0.9pt,left=3mm]{\Large\bfseries\color{cari}c)}\quad\begin{minipage}[c]{5.3cm}\raggedright\Large $]-1\,;\,3[$\end{minipage}\end{tcolorbox}\hfill\begin{tcolorbox}[enhanced,nobeforeafter,width=7.1cm,height=2.1cm,valign=center,colback=cari!6,colframe=cari,arc=2.5mm,boxrule=0.9pt,left=3mm]{\Large\bfseries\color{cari}d)}\quad\begin{minipage}[c]{5.3cm}\raggedright\Large $]-\infty\,;\,+\infty[$\end{minipage}\end{tcolorbox}
\end{tcolorbox}
\newpage
\begin{tcolorbox}[enhanced,width=16cm,height=14cm,colback=white,colframe=cari,boxrule=2pt,arc=0mm,
left=6mm,right=6mm,top=0mm,bottom=4mm,sharp corners,
overlay={\fill[cari] (frame.north west) rectangle ([yshift=-1.25cm]frame.north east);
\node[anchor=west,text=white,font=\sffamily\bfseries\LARGE] at ([xshift=5mm,yshift=-0.62cm]frame.north west) {Question 80 B};
\node[anchor=east,text=white,font=\sffamily\itshape\large] at ([xshift=-5mm,yshift=-0.62cm]frame.north east) {Arithmétique et ensembles de nombres};}]
\vspace*{1.75cm}
\begin{minipage}[t][6.9cm][c]{\linewidth}\LARGE\centering
\begin{minipage}{0.95\linewidth}\raggedright\linespread{1.25}\selectfont Que vaut l'intersection $[-5\,;\,2]\cap]0\,;\,6]$~?\end{minipage}
\end{minipage}

\vfill
\begin{tcolorbox}[enhanced,nobeforeafter,width=7.1cm,height=2.1cm,valign=center,colback=cari!6,colframe=cari,arc=2.5mm,boxrule=0.9pt,left=3mm]{\Large\bfseries\color{cari}a)}\quad\begin{minipage}[c]{5.3cm}\raggedright\Large $[0\,;\,2]$\end{minipage}\end{tcolorbox}\hfill\begin{tcolorbox}[enhanced,nobeforeafter,width=7.1cm,height=2.1cm,valign=center,colback=cari!6,colframe=cari,arc=2.5mm,boxrule=0.9pt,left=3mm]{\Large\bfseries\color{cari}b)}\quad\begin{minipage}[c]{5.3cm}\raggedright\Large $]0\,;\,2[$\end{minipage}\end{tcolorbox}\par\vspace{2.5mm}
\begin{tcolorbox}[enhanced,nobeforeafter,width=7.1cm,height=2.1cm,valign=center,colback=cari!6,colframe=cari,arc=2.5mm,boxrule=0.9pt,left=3mm]{\Large\bfseries\color{cari}c)}\quad\begin{minipage}[c]{5.3cm}\raggedright\Large $]0\,;\,2]$\end{minipage}\end{tcolorbox}\hfill\begin{tcolorbox}[enhanced,nobeforeafter,width=7.1cm,height=2.1cm,valign=center,colback=cari!6,colframe=cari,arc=2.5mm,boxrule=0.9pt,left=3mm]{\Large\bfseries\color{cari}d)}\quad\begin{minipage}[c]{5.3cm}\raggedright\Large $[-5\,;\,6]$\end{minipage}\end{tcolorbox}
\end{tcolorbox}
\newpage
\begin{tcolorbox}[enhanced,width=16cm,height=14cm,colback=white,colframe=cpro,boxrule=2pt,arc=0mm,
left=6mm,right=6mm,top=0mm,bottom=4mm,sharp corners,
overlay={\fill[cpro] (frame.north west) rectangle ([yshift=-1.25cm]frame.north east);
\node[anchor=west,text=white,font=\sffamily\bfseries\LARGE] at ([xshift=5mm,yshift=-0.62cm]frame.north west) {Question 81 B};
\node[anchor=east,text=white,font=\sffamily\itshape\large] at ([xshift=-5mm,yshift=-0.62cm]frame.north east) {Probabilités};}]
\vspace*{1.75cm}
\begin{minipage}[t][6.9cm][c]{\linewidth}\LARGE\centering
\begin{minipage}{0.95\linewidth}\raggedright\linespread{1.25}\selectfont On a $P(A)=0{,}25$, $P(B)=0{,}6$ et $P(A\cap B)=0{,}2$. Que vaut $P(A\cup B)$~?\end{minipage}
\end{minipage}

\vfill
\begin{tcolorbox}[enhanced,nobeforeafter,width=7.1cm,height=2.1cm,valign=center,colback=cpro!6,colframe=cpro,arc=2.5mm,boxrule=0.9pt,left=3mm]{\Large\bfseries\color{cpro}a)}\quad\begin{minipage}[c]{5.3cm}\raggedright\Large $0{,}85$\end{minipage}\end{tcolorbox}\hfill\begin{tcolorbox}[enhanced,nobeforeafter,width=7.1cm,height=2.1cm,valign=center,colback=cpro!6,colframe=cpro,arc=2.5mm,boxrule=0.9pt,left=3mm]{\Large\bfseries\color{cpro}b)}\quad\begin{minipage}[c]{5.3cm}\raggedright\Large $0{,}65$\end{minipage}\end{tcolorbox}\par\vspace{2.5mm}
\begin{tcolorbox}[enhanced,nobeforeafter,width=7.1cm,height=2.1cm,valign=center,colback=cpro!6,colframe=cpro,arc=2.5mm,boxrule=0.9pt,left=3mm]{\Large\bfseries\color{cpro}c)}\quad\begin{minipage}[c]{5.3cm}\raggedright\Large $0{,}15$\end{minipage}\end{tcolorbox}\hfill\begin{tcolorbox}[enhanced,nobeforeafter,width=7.1cm,height=2.1cm,valign=center,colback=cpro!6,colframe=cpro,arc=2.5mm,boxrule=0.9pt,left=3mm]{\Large\bfseries\color{cpro}d)}\quad\begin{minipage}[c]{5.3cm}\raggedright\Large $0{,}2$\end{minipage}\end{tcolorbox}
\end{tcolorbox}
\newpage
\begin{tcolorbox}[enhanced,width=16cm,height=14cm,colback=white,colframe=cpro,boxrule=2pt,arc=0mm,
left=6mm,right=6mm,top=0mm,bottom=4mm,sharp corners,
overlay={\fill[cpro] (frame.north west) rectangle ([yshift=-1.25cm]frame.north east);
\node[anchor=west,text=white,font=\sffamily\bfseries\LARGE] at ([xshift=5mm,yshift=-0.62cm]frame.north west) {Question 82 B};
\node[anchor=east,text=white,font=\sffamily\itshape\large] at ([xshift=-5mm,yshift=-0.62cm]frame.north east) {Probabilités};}]
\vspace*{1.75cm}
\begin{minipage}[t][6.9cm][c]{\linewidth}\LARGE\centering
\begin{minipage}{0.95\linewidth}\raggedright\linespread{1.25}\selectfont On a $P(A)=0{,}12$, $P(B)=0{,}4$ et $P(A\cap B)=0{,}08$. Que vaut $P(A\cup B)$~?\end{minipage}
\end{minipage}

\vfill
\begin{tcolorbox}[enhanced,nobeforeafter,width=7.1cm,height=2.1cm,valign=center,colback=cpro!6,colframe=cpro,arc=2.5mm,boxrule=0.9pt,left=3mm]{\Large\bfseries\color{cpro}a)}\quad\begin{minipage}[c]{5.3cm}\raggedright\Large $0{,}6$\end{minipage}\end{tcolorbox}\hfill\begin{tcolorbox}[enhanced,nobeforeafter,width=7.1cm,height=2.1cm,valign=center,colback=cpro!6,colframe=cpro,arc=2.5mm,boxrule=0.9pt,left=3mm]{\Large\bfseries\color{cpro}b)}\quad\begin{minipage}[c]{5.3cm}\raggedright\Large $0{,}44$\end{minipage}\end{tcolorbox}\par\vspace{2.5mm}
\begin{tcolorbox}[enhanced,nobeforeafter,width=7.1cm,height=2.1cm,valign=center,colback=cpro!6,colframe=cpro,arc=2.5mm,boxrule=0.9pt,left=3mm]{\Large\bfseries\color{cpro}c)}\quad\begin{minipage}[c]{5.3cm}\raggedright\Large $0{,}08$\end{minipage}\end{tcolorbox}\hfill\begin{tcolorbox}[enhanced,nobeforeafter,width=7.1cm,height=2.1cm,valign=center,colback=cpro!6,colframe=cpro,arc=2.5mm,boxrule=0.9pt,left=3mm]{\Large\bfseries\color{cpro}d)}\quad\begin{minipage}[c]{5.3cm}\raggedright\Large $0{,}52$\end{minipage}\end{tcolorbox}
\end{tcolorbox}
\newpage
\begin{tcolorbox}[enhanced,width=16cm,height=14cm,colback=white,colframe=cpro,boxrule=2pt,arc=0mm,
left=6mm,right=6mm,top=0mm,bottom=4mm,sharp corners,
overlay={\fill[cpro] (frame.north west) rectangle ([yshift=-1.25cm]frame.north east);
\node[anchor=west,text=white,font=\sffamily\bfseries\LARGE] at ([xshift=5mm,yshift=-0.62cm]frame.north west) {Question 83 B};
\node[anchor=east,text=white,font=\sffamily\itshape\large] at ([xshift=-5mm,yshift=-0.62cm]frame.north east) {Probabilités};}]
\vspace*{1.75cm}
\begin{minipage}[t][6.9cm][c]{\linewidth}\LARGE\centering
\begin{minipage}{0.95\linewidth}\raggedright\linespread{1.25}\selectfont On a $P(A)=0{,}25$, $P(B)=0{,}5$ et $P(A\cup B)=0{,}6$. Que vaut $P(A\cap B)$~?\end{minipage}
\end{minipage}

\vfill
\begin{tcolorbox}[enhanced,nobeforeafter,width=7.1cm,height=2.1cm,valign=center,colback=cpro!6,colframe=cpro,arc=2.5mm,boxrule=0.9pt,left=3mm]{\Large\bfseries\color{cpro}a)}\quad\begin{minipage}[c]{5.3cm}\raggedright\Large $0{,}15$\end{minipage}\end{tcolorbox}\hfill\begin{tcolorbox}[enhanced,nobeforeafter,width=7.1cm,height=2.1cm,valign=center,colback=cpro!6,colframe=cpro,arc=2.5mm,boxrule=0.9pt,left=3mm]{\Large\bfseries\color{cpro}b)}\quad\begin{minipage}[c]{5.3cm}\raggedright\Large $0{,}125$\end{minipage}\end{tcolorbox}\par\vspace{2.5mm}
\begin{tcolorbox}[enhanced,nobeforeafter,width=7.1cm,height=2.1cm,valign=center,colback=cpro!6,colframe=cpro,arc=2.5mm,boxrule=0.9pt,left=3mm]{\Large\bfseries\color{cpro}c)}\quad\begin{minipage}[c]{5.3cm}\raggedright\Large $0{,}1$\end{minipage}\end{tcolorbox}\hfill\begin{tcolorbox}[enhanced,nobeforeafter,width=7.1cm,height=2.1cm,valign=center,colback=cpro!6,colframe=cpro,arc=2.5mm,boxrule=0.9pt,left=3mm]{\Large\bfseries\color{cpro}d)}\quad\begin{minipage}[c]{5.3cm}\raggedright\Large $0{,}35$\end{minipage}\end{tcolorbox}
\end{tcolorbox}
\newpage
\begin{tcolorbox}[enhanced,width=16cm,height=14cm,colback=white,colframe=cpro,boxrule=2pt,arc=0mm,
left=6mm,right=6mm,top=0mm,bottom=4mm,sharp corners,
overlay={\fill[cpro] (frame.north west) rectangle ([yshift=-1.25cm]frame.north east);
\node[anchor=west,text=white,font=\sffamily\bfseries\LARGE] at ([xshift=5mm,yshift=-0.62cm]frame.north west) {Question 84 B};
\node[anchor=east,text=white,font=\sffamily\itshape\large] at ([xshift=-5mm,yshift=-0.62cm]frame.north east) {Probabilités};}]
\vspace*{1.75cm}
\begin{minipage}[t][6.9cm][c]{\linewidth}\LARGE\centering
\begin{minipage}{0.95\linewidth}\raggedright\linespread{1.25}\selectfont On a $P(A)=0{,}6$, $P(B)=0{,}55$ et $P(A\cup B)=0{,}85$. Que vaut $P(A\cap B)$~?\end{minipage}
\end{minipage}

\vfill
\begin{tcolorbox}[enhanced,nobeforeafter,width=7.1cm,height=2.1cm,valign=center,colback=cpro!6,colframe=cpro,arc=2.5mm,boxrule=0.9pt,left=3mm]{\Large\bfseries\color{cpro}a)}\quad\begin{minipage}[c]{5.3cm}\raggedright\Large $0{,}4$\end{minipage}\end{tcolorbox}\hfill\begin{tcolorbox}[enhanced,nobeforeafter,width=7.1cm,height=2.1cm,valign=center,colback=cpro!6,colframe=cpro,arc=2.5mm,boxrule=0.9pt,left=3mm]{\Large\bfseries\color{cpro}b)}\quad\begin{minipage}[c]{5.3cm}\raggedright\Large $0{,}25$\end{minipage}\end{tcolorbox}\par\vspace{2.5mm}
\begin{tcolorbox}[enhanced,nobeforeafter,width=7.1cm,height=2.1cm,valign=center,colback=cpro!6,colframe=cpro,arc=2.5mm,boxrule=0.9pt,left=3mm]{\Large\bfseries\color{cpro}c)}\quad\begin{minipage}[c]{5.3cm}\raggedright\Large $0{,}3$\end{minipage}\end{tcolorbox}\hfill\begin{tcolorbox}[enhanced,nobeforeafter,width=7.1cm,height=2.1cm,valign=center,colback=cpro!6,colframe=cpro,arc=2.5mm,boxrule=0.9pt,left=3mm]{\Large\bfseries\color{cpro}d)}\quad\begin{minipage}[c]{5.3cm}\raggedright\Large $0{,}33$\end{minipage}\end{tcolorbox}
\end{tcolorbox}
\newpage
\begin{tcolorbox}[enhanced,width=16cm,height=14cm,colback=white,colframe=cpro,boxrule=2pt,arc=0mm,
left=6mm,right=6mm,top=0mm,bottom=4mm,sharp corners,
overlay={\fill[cpro] (frame.north west) rectangle ([yshift=-1.25cm]frame.north east);
\node[anchor=west,text=white,font=\sffamily\bfseries\LARGE] at ([xshift=5mm,yshift=-0.62cm]frame.north west) {Question 85 B};
\node[anchor=east,text=white,font=\sffamily\itshape\large] at ([xshift=-5mm,yshift=-0.62cm]frame.north east) {Probabilités};}]
\vspace*{1.75cm}
\begin{minipage}[t][6.9cm][c]{\linewidth}\LARGE\centering
\begin{minipage}{0.95\linewidth}\raggedright\linespread{1.25}\selectfont On a $P(A)=\dfrac{7}{12}$. Que vaut $P(\overline{A})$, la probabilité de l'événement contraire~?\end{minipage}
\end{minipage}

\vfill
\begin{tcolorbox}[enhanced,nobeforeafter,width=7.1cm,height=2.1cm,valign=center,colback=cpro!6,colframe=cpro,arc=2.5mm,boxrule=0.9pt,left=3mm]{\Large\bfseries\color{cpro}a)}\quad\begin{minipage}[c]{5.3cm}\raggedright\Large $\dfrac{5}{24}$\end{minipage}\end{tcolorbox}\hfill\begin{tcolorbox}[enhanced,nobeforeafter,width=7.1cm,height=2.1cm,valign=center,colback=cpro!6,colframe=cpro,arc=2.5mm,boxrule=0.9pt,left=3mm]{\Large\bfseries\color{cpro}b)}\quad\begin{minipage}[c]{5.3cm}\raggedright\Large $\dfrac{7}{12}$\end{minipage}\end{tcolorbox}\par\vspace{2.5mm}
\begin{tcolorbox}[enhanced,nobeforeafter,width=7.1cm,height=2.1cm,valign=center,colback=cpro!6,colframe=cpro,arc=2.5mm,boxrule=0.9pt,left=3mm]{\Large\bfseries\color{cpro}c)}\quad\begin{minipage}[c]{5.3cm}\raggedright\Large $\dfrac{12}{7}$\end{minipage}\end{tcolorbox}\hfill\begin{tcolorbox}[enhanced,nobeforeafter,width=7.1cm,height=2.1cm,valign=center,colback=cpro!6,colframe=cpro,arc=2.5mm,boxrule=0.9pt,left=3mm]{\Large\bfseries\color{cpro}d)}\quad\begin{minipage}[c]{5.3cm}\raggedright\Large $\dfrac{5}{12}$\end{minipage}\end{tcolorbox}
\end{tcolorbox}
\newpage
\begin{tcolorbox}[enhanced,width=16cm,height=14cm,colback=white,colframe=cpro,boxrule=2pt,arc=0mm,
left=6mm,right=6mm,top=0mm,bottom=4mm,sharp corners,
overlay={\fill[cpro] (frame.north west) rectangle ([yshift=-1.25cm]frame.north east);
\node[anchor=west,text=white,font=\sffamily\bfseries\LARGE] at ([xshift=5mm,yshift=-0.62cm]frame.north west) {Question 86 B};
\node[anchor=east,text=white,font=\sffamily\itshape\large] at ([xshift=-5mm,yshift=-0.62cm]frame.north east) {Probabilités};}]
\vspace*{1.75cm}
\begin{minipage}[t][6.9cm][c]{\linewidth}\LARGE\centering
\begin{minipage}{0.95\linewidth}\raggedright\linespread{1.25}\selectfont On a $P(A)=\dfrac{1}{6}$. Que vaut $P(\overline{A})$, la probabilité de l'événement contraire~?\end{minipage}
\end{minipage}

\vfill
\begin{tcolorbox}[enhanced,nobeforeafter,width=7.1cm,height=2.1cm,valign=center,colback=cpro!6,colframe=cpro,arc=2.5mm,boxrule=0.9pt,left=3mm]{\Large\bfseries\color{cpro}a)}\quad\begin{minipage}[c]{5.3cm}\raggedright\Large $\dfrac{1}{6}$\end{minipage}\end{tcolorbox}\hfill\begin{tcolorbox}[enhanced,nobeforeafter,width=7.1cm,height=2.1cm,valign=center,colback=cpro!6,colframe=cpro,arc=2.5mm,boxrule=0.9pt,left=3mm]{\Large\bfseries\color{cpro}b)}\quad\begin{minipage}[c]{5.3cm}\raggedright\Large $6$\end{minipage}\end{tcolorbox}\par\vspace{2.5mm}
\begin{tcolorbox}[enhanced,nobeforeafter,width=7.1cm,height=2.1cm,valign=center,colback=cpro!6,colframe=cpro,arc=2.5mm,boxrule=0.9pt,left=3mm]{\Large\bfseries\color{cpro}c)}\quad\begin{minipage}[c]{5.3cm}\raggedright\Large $\dfrac{5}{6}$\end{minipage}\end{tcolorbox}\hfill\begin{tcolorbox}[enhanced,nobeforeafter,width=7.1cm,height=2.1cm,valign=center,colback=cpro!6,colframe=cpro,arc=2.5mm,boxrule=0.9pt,left=3mm]{\Large\bfseries\color{cpro}d)}\quad\begin{minipage}[c]{5.3cm}\raggedright\Large $\dfrac{5}{12}$\end{minipage}\end{tcolorbox}
\end{tcolorbox}
\newpage
\begin{tcolorbox}[enhanced,width=16cm,height=14cm,colback=white,colframe=cpro,boxrule=2pt,arc=0mm,
left=6mm,right=6mm,top=0mm,bottom=4mm,sharp corners,
overlay={\fill[cpro] (frame.north west) rectangle ([yshift=-1.25cm]frame.north east);
\node[anchor=west,text=white,font=\sffamily\bfseries\LARGE] at ([xshift=5mm,yshift=-0.62cm]frame.north west) {Question 87 B};
\node[anchor=east,text=white,font=\sffamily\itshape\large] at ([xshift=-5mm,yshift=-0.62cm]frame.north east) {Probabilités};}]
\vspace*{1.75cm}
\begin{minipage}[t][6.9cm][c]{\linewidth}\LARGE\centering
\begin{minipage}{0.95\linewidth}\raggedright\linespread{1.25}\selectfont On lance un dé équilibré à six faces. Quelle est la probabilité d'obtenir un nombre premier~?\end{minipage}
\end{minipage}

\vfill
\begin{tcolorbox}[enhanced,nobeforeafter,width=7.1cm,height=2.1cm,valign=center,colback=cpro!6,colframe=cpro,arc=2.5mm,boxrule=0.9pt,left=3mm]{\Large\bfseries\color{cpro}a)}\quad\begin{minipage}[c]{5.3cm}\raggedright\Large $\dfrac{1}{3}$\end{minipage}\end{tcolorbox}\hfill\begin{tcolorbox}[enhanced,nobeforeafter,width=7.1cm,height=2.1cm,valign=center,colback=cpro!6,colframe=cpro,arc=2.5mm,boxrule=0.9pt,left=3mm]{\Large\bfseries\color{cpro}b)}\quad\begin{minipage}[c]{5.3cm}\raggedright\Large $\dfrac{2}{3}$\end{minipage}\end{tcolorbox}\par\vspace{2.5mm}
\begin{tcolorbox}[enhanced,nobeforeafter,width=7.1cm,height=2.1cm,valign=center,colback=cpro!6,colframe=cpro,arc=2.5mm,boxrule=0.9pt,left=3mm]{\Large\bfseries\color{cpro}c)}\quad\begin{minipage}[c]{5.3cm}\raggedright\Large $\dfrac{1}{2}$\end{minipage}\end{tcolorbox}\hfill\begin{tcolorbox}[enhanced,nobeforeafter,width=7.1cm,height=2.1cm,valign=center,colback=cpro!6,colframe=cpro,arc=2.5mm,boxrule=0.9pt,left=3mm]{\Large\bfseries\color{cpro}d)}\quad\begin{minipage}[c]{5.3cm}\raggedright\Large $\dfrac{3}{5}$\end{minipage}\end{tcolorbox}
\end{tcolorbox}
\newpage
\begin{tcolorbox}[enhanced,width=16cm,height=14cm,colback=white,colframe=cpro,boxrule=2pt,arc=0mm,
left=6mm,right=6mm,top=0mm,bottom=4mm,sharp corners,
overlay={\fill[cpro] (frame.north west) rectangle ([yshift=-1.25cm]frame.north east);
\node[anchor=west,text=white,font=\sffamily\bfseries\LARGE] at ([xshift=5mm,yshift=-0.62cm]frame.north west) {Question 88 B};
\node[anchor=east,text=white,font=\sffamily\itshape\large] at ([xshift=-5mm,yshift=-0.62cm]frame.north east) {Probabilités};}]
\vspace*{1.75cm}
\begin{minipage}[t][6.9cm][c]{\linewidth}\LARGE\centering
\begin{minipage}{0.95\linewidth}\raggedright\linespread{1.25}\selectfont On lance un dé équilibré à six faces. Quelle est la probabilité d'obtenir un nombre impair ou supérieur à 4~?\end{minipage}
\end{minipage}

\vfill
\begin{tcolorbox}[enhanced,nobeforeafter,width=7.1cm,height=2.1cm,valign=center,colback=cpro!6,colframe=cpro,arc=2.5mm,boxrule=0.9pt,left=3mm]{\Large\bfseries\color{cpro}a)}\quad\begin{minipage}[c]{5.3cm}\raggedright\Large $\dfrac{5}{6}$\end{minipage}\end{tcolorbox}\hfill\begin{tcolorbox}[enhanced,nobeforeafter,width=7.1cm,height=2.1cm,valign=center,colback=cpro!6,colframe=cpro,arc=2.5mm,boxrule=0.9pt,left=3mm]{\Large\bfseries\color{cpro}b)}\quad\begin{minipage}[c]{5.3cm}\raggedright\Large $\dfrac{1}{2}$\end{minipage}\end{tcolorbox}\par\vspace{2.5mm}
\begin{tcolorbox}[enhanced,nobeforeafter,width=7.1cm,height=2.1cm,valign=center,colback=cpro!6,colframe=cpro,arc=2.5mm,boxrule=0.9pt,left=3mm]{\Large\bfseries\color{cpro}c)}\quad\begin{minipage}[c]{5.3cm}\raggedright\Large $\dfrac{2}{3}$\end{minipage}\end{tcolorbox}\hfill\begin{tcolorbox}[enhanced,nobeforeafter,width=7.1cm,height=2.1cm,valign=center,colback=cpro!6,colframe=cpro,arc=2.5mm,boxrule=0.9pt,left=3mm]{\Large\bfseries\color{cpro}d)}\quad\begin{minipage}[c]{5.3cm}\raggedright\Large $\dfrac{1}{3}$\end{minipage}\end{tcolorbox}
\end{tcolorbox}
\newpage
\begin{tcolorbox}[enhanced,width=16cm,height=14cm,colback=white,colframe=cpro,boxrule=2pt,arc=0mm,
left=6mm,right=6mm,top=0mm,bottom=4mm,sharp corners,
overlay={\fill[cpro] (frame.north west) rectangle ([yshift=-1.25cm]frame.north east);
\node[anchor=west,text=white,font=\sffamily\bfseries\LARGE] at ([xshift=5mm,yshift=-0.62cm]frame.north west) {Question 89 B};
\node[anchor=east,text=white,font=\sffamily\itshape\large] at ([xshift=-5mm,yshift=-0.62cm]frame.north east) {Probabilités};}]
\vspace*{1.75cm}
\begin{minipage}[t][6.9cm][c]{\linewidth}\LARGE\centering
\begin{minipage}{0.95\linewidth}\raggedright\linespread{1.25}\selectfont On lance $2$ pièces équilibrées. Quelle est la probabilité d'obtenir au moins une fois FACE~?\end{minipage}
\end{minipage}

\vfill
\begin{tcolorbox}[enhanced,nobeforeafter,width=7.1cm,height=2.1cm,valign=center,colback=cpro!6,colframe=cpro,arc=2.5mm,boxrule=0.9pt,left=3mm]{\Large\bfseries\color{cpro}a)}\quad\begin{minipage}[c]{5.3cm}\raggedright\Large $\dfrac{3}{4}$\end{minipage}\end{tcolorbox}\hfill\begin{tcolorbox}[enhanced,nobeforeafter,width=7.1cm,height=2.1cm,valign=center,colback=cpro!6,colframe=cpro,arc=2.5mm,boxrule=0.9pt,left=3mm]{\Large\bfseries\color{cpro}b)}\quad\begin{minipage}[c]{5.3cm}\raggedright\Large $\dfrac{1}{4}$\end{minipage}\end{tcolorbox}\par\vspace{2.5mm}
\begin{tcolorbox}[enhanced,nobeforeafter,width=7.1cm,height=2.1cm,valign=center,colback=cpro!6,colframe=cpro,arc=2.5mm,boxrule=0.9pt,left=3mm]{\Large\bfseries\color{cpro}c)}\quad\begin{minipage}[c]{5.3cm}\raggedright\Large $\dfrac{1}{8}$\end{minipage}\end{tcolorbox}\hfill\begin{tcolorbox}[enhanced,nobeforeafter,width=7.1cm,height=2.1cm,valign=center,colback=cpro!6,colframe=cpro,arc=2.5mm,boxrule=0.9pt,left=3mm]{\Large\bfseries\color{cpro}d)}\quad\begin{minipage}[c]{5.3cm}\raggedright\Large $\dfrac{1}{2}$\end{minipage}\end{tcolorbox}
\end{tcolorbox}
\newpage
\begin{tcolorbox}[enhanced,width=16cm,height=14cm,colback=white,colframe=cpro,boxrule=2pt,arc=0mm,
left=6mm,right=6mm,top=0mm,bottom=4mm,sharp corners,
overlay={\fill[cpro] (frame.north west) rectangle ([yshift=-1.25cm]frame.north east);
\node[anchor=west,text=white,font=\sffamily\bfseries\LARGE] at ([xshift=5mm,yshift=-0.62cm]frame.north west) {Question 90 B};
\node[anchor=east,text=white,font=\sffamily\itshape\large] at ([xshift=-5mm,yshift=-0.62cm]frame.north east) {Probabilités};}]
\vspace*{1.75cm}
\begin{minipage}[t][6.9cm][c]{\linewidth}\LARGE\centering
\begin{minipage}{0.95\linewidth}\raggedright\linespread{1.25}\selectfont On lance $3$ pièces équilibrées. Quelle est la probabilité d'obtenir au moins une fois FACE~?\end{minipage}
\end{minipage}

\vfill
\begin{tcolorbox}[enhanced,nobeforeafter,width=7.1cm,height=2.1cm,valign=center,colback=cpro!6,colframe=cpro,arc=2.5mm,boxrule=0.9pt,left=3mm]{\Large\bfseries\color{cpro}a)}\quad\begin{minipage}[c]{5.3cm}\raggedright\Large $\dfrac{1}{6}$\end{minipage}\end{tcolorbox}\hfill\begin{tcolorbox}[enhanced,nobeforeafter,width=7.1cm,height=2.1cm,valign=center,colback=cpro!6,colframe=cpro,arc=2.5mm,boxrule=0.9pt,left=3mm]{\Large\bfseries\color{cpro}b)}\quad\begin{minipage}[c]{5.3cm}\raggedright\Large $\dfrac{1}{8}$\end{minipage}\end{tcolorbox}\par\vspace{2.5mm}
\begin{tcolorbox}[enhanced,nobeforeafter,width=7.1cm,height=2.1cm,valign=center,colback=cpro!6,colframe=cpro,arc=2.5mm,boxrule=0.9pt,left=3mm]{\Large\bfseries\color{cpro}c)}\quad\begin{minipage}[c]{5.3cm}\raggedright\Large $\dfrac{1}{2}$\end{minipage}\end{tcolorbox}\hfill\begin{tcolorbox}[enhanced,nobeforeafter,width=7.1cm,height=2.1cm,valign=center,colback=cpro!6,colframe=cpro,arc=2.5mm,boxrule=0.9pt,left=3mm]{\Large\bfseries\color{cpro}d)}\quad\begin{minipage}[c]{5.3cm}\raggedright\Large $\dfrac{7}{8}$\end{minipage}\end{tcolorbox}
\end{tcolorbox}
\newpage
\begin{tcolorbox}[enhanced,width=16cm,height=14cm,colback=white,colframe=cpro,boxrule=2pt,arc=0mm,
left=6mm,right=6mm,top=0mm,bottom=4mm,sharp corners,
overlay={\fill[cpro] (frame.north west) rectangle ([yshift=-1.25cm]frame.north east);
\node[anchor=west,text=white,font=\sffamily\bfseries\LARGE] at ([xshift=5mm,yshift=-0.62cm]frame.north west) {Question 91 B};
\node[anchor=east,text=white,font=\sffamily\itshape\large] at ([xshift=-5mm,yshift=-0.62cm]frame.north east) {Probabilités};}]
\vspace*{1.75cm}
\begin{minipage}[t][6.9cm][c]{\linewidth}\Large\centering
\begin{minipage}{0.95\linewidth}\raggedright\linespread{1.25}\selectfont Voici la loi de probabilité d'une expérience aléatoire. Quelle est la probabilité manquante~? $\begin{array}{|l|c|c|c|c|}\hline \text{Issue} & 1 & 2 & 3 & 4\\\hline \text{Probabilité} & 0{,}3 & 0{,}3 &~? & 0{,}12\\\hline\end{array}$\end{minipage}
\end{minipage}

\vfill
\begin{tcolorbox}[enhanced,nobeforeafter,width=7.1cm,height=2.1cm,valign=center,colback=cpro!6,colframe=cpro,arc=2.5mm,boxrule=0.9pt,left=3mm]{\Large\bfseries\color{cpro}a)}\quad\begin{minipage}[c]{5.3cm}\raggedright\Large $0{,}25$\end{minipage}\end{tcolorbox}\hfill\begin{tcolorbox}[enhanced,nobeforeafter,width=7.1cm,height=2.1cm,valign=center,colback=cpro!6,colframe=cpro,arc=2.5mm,boxrule=0.9pt,left=3mm]{\Large\bfseries\color{cpro}b)}\quad\begin{minipage}[c]{5.3cm}\raggedright\Large $0{,}4$\end{minipage}\end{tcolorbox}\par\vspace{2.5mm}
\begin{tcolorbox}[enhanced,nobeforeafter,width=7.1cm,height=2.1cm,valign=center,colback=cpro!6,colframe=cpro,arc=2.5mm,boxrule=0.9pt,left=3mm]{\Large\bfseries\color{cpro}c)}\quad\begin{minipage}[c]{5.3cm}\raggedright\Large $0{,}28$\end{minipage}\end{tcolorbox}\hfill\begin{tcolorbox}[enhanced,nobeforeafter,width=7.1cm,height=2.1cm,valign=center,colback=cpro!6,colframe=cpro,arc=2.5mm,boxrule=0.9pt,left=3mm]{\Large\bfseries\color{cpro}d)}\quad\begin{minipage}[c]{5.3cm}\raggedright\Large $0{,}72$\end{minipage}\end{tcolorbox}
\end{tcolorbox}
\newpage
\begin{tcolorbox}[enhanced,width=16cm,height=14cm,colback=white,colframe=cpro,boxrule=2pt,arc=0mm,
left=6mm,right=6mm,top=0mm,bottom=4mm,sharp corners,
overlay={\fill[cpro] (frame.north west) rectangle ([yshift=-1.25cm]frame.north east);
\node[anchor=west,text=white,font=\sffamily\bfseries\LARGE] at ([xshift=5mm,yshift=-0.62cm]frame.north west) {Question 92 B};
\node[anchor=east,text=white,font=\sffamily\itshape\large] at ([xshift=-5mm,yshift=-0.62cm]frame.north east) {Probabilités};}]
\vspace*{1.75cm}
\begin{minipage}[t][6.9cm][c]{\linewidth}\Large\centering
\begin{minipage}{0.95\linewidth}\raggedright\linespread{1.25}\selectfont Voici la loi de probabilité d'une expérience aléatoire. Quelle est la probabilité manquante~? $\begin{array}{|l|c|c|c|c|}\hline \text{Issue} & 1 & 2 & 3 & 4\\\hline \text{Probabilité} & 0{,}05 & 0{,}6 &~? & 0{,}15\\\hline\end{array}$\end{minipage}
\end{minipage}

\vfill
\begin{tcolorbox}[enhanced,nobeforeafter,width=7.1cm,height=2.1cm,valign=center,colback=cpro!6,colframe=cpro,arc=2.5mm,boxrule=0.9pt,left=3mm]{\Large\bfseries\color{cpro}a)}\quad\begin{minipage}[c]{5.3cm}\raggedright\Large $0{,}2$\end{minipage}\end{tcolorbox}\hfill\begin{tcolorbox}[enhanced,nobeforeafter,width=7.1cm,height=2.1cm,valign=center,colback=cpro!6,colframe=cpro,arc=2.5mm,boxrule=0.9pt,left=3mm]{\Large\bfseries\color{cpro}b)}\quad\begin{minipage}[c]{5.3cm}\raggedright\Large $0{,}35$\end{minipage}\end{tcolorbox}\par\vspace{2.5mm}
\begin{tcolorbox}[enhanced,nobeforeafter,width=7.1cm,height=2.1cm,valign=center,colback=cpro!6,colframe=cpro,arc=2.5mm,boxrule=0.9pt,left=3mm]{\Large\bfseries\color{cpro}c)}\quad\begin{minipage}[c]{5.3cm}\raggedright\Large $0{,}25$\end{minipage}\end{tcolorbox}\hfill\begin{tcolorbox}[enhanced,nobeforeafter,width=7.1cm,height=2.1cm,valign=center,colback=cpro!6,colframe=cpro,arc=2.5mm,boxrule=0.9pt,left=3mm]{\Large\bfseries\color{cpro}d)}\quad\begin{minipage}[c]{5.3cm}\raggedright\Large $0{,}8$\end{minipage}\end{tcolorbox}
\end{tcolorbox}
\newpage
\begin{tcolorbox}[enhanced,width=16cm,height=14cm,colback=white,colframe=cpro,boxrule=2pt,arc=0mm,
left=6mm,right=6mm,top=0mm,bottom=4mm,sharp corners,
overlay={\fill[cpro] (frame.north west) rectangle ([yshift=-1.25cm]frame.north east);
\node[anchor=west,text=white,font=\sffamily\bfseries\LARGE] at ([xshift=5mm,yshift=-0.62cm]frame.north west) {Question 93 B};
\node[anchor=east,text=white,font=\sffamily\itshape\large] at ([xshift=-5mm,yshift=-0.62cm]frame.north east) {Probabilités};}]
\vspace*{1.75cm}
\begin{minipage}[t][6.9cm][c]{\linewidth}\Large\centering
\begin{minipage}{0.95\linewidth}\raggedright\linespread{1.25}\selectfont Une urne contient $10$ boules dont $4$ rouges. On tire une boule, on la remet, puis on en tire une deuxième. Quelle est la probabilité de tirer deux boules rouges~?\end{minipage}
\end{minipage}

\vfill
\begin{tcolorbox}[enhanced,nobeforeafter,width=7.1cm,height=2.1cm,valign=center,colback=cpro!6,colframe=cpro,arc=2.5mm,boxrule=0.9pt,left=3mm]{\Large\bfseries\color{cpro}a)}\quad\begin{minipage}[c]{5.3cm}\raggedright\Large $\dfrac{2}{15}$\end{minipage}\end{tcolorbox}\hfill\begin{tcolorbox}[enhanced,nobeforeafter,width=7.1cm,height=2.1cm,valign=center,colback=cpro!6,colframe=cpro,arc=2.5mm,boxrule=0.9pt,left=3mm]{\Large\bfseries\color{cpro}b)}\quad\begin{minipage}[c]{5.3cm}\raggedright\Large $\dfrac{2}{5}$\end{minipage}\end{tcolorbox}\par\vspace{2.5mm}
\begin{tcolorbox}[enhanced,nobeforeafter,width=7.1cm,height=2.1cm,valign=center,colback=cpro!6,colframe=cpro,arc=2.5mm,boxrule=0.9pt,left=3mm]{\Large\bfseries\color{cpro}c)}\quad\begin{minipage}[c]{5.3cm}\raggedright\Large $\dfrac{4}{25}$\end{minipage}\end{tcolorbox}\hfill\begin{tcolorbox}[enhanced,nobeforeafter,width=7.1cm,height=2.1cm,valign=center,colback=cpro!6,colframe=cpro,arc=2.5mm,boxrule=0.9pt,left=3mm]{\Large\bfseries\color{cpro}d)}\quad\begin{minipage}[c]{5.3cm}\raggedright\Large $\dfrac{4}{5}$\end{minipage}\end{tcolorbox}
\end{tcolorbox}
\newpage
\begin{tcolorbox}[enhanced,width=16cm,height=14cm,colback=white,colframe=cpro,boxrule=2pt,arc=0mm,
left=6mm,right=6mm,top=0mm,bottom=4mm,sharp corners,
overlay={\fill[cpro] (frame.north west) rectangle ([yshift=-1.25cm]frame.north east);
\node[anchor=west,text=white,font=\sffamily\bfseries\LARGE] at ([xshift=5mm,yshift=-0.62cm]frame.north west) {Question 94 B};
\node[anchor=east,text=white,font=\sffamily\itshape\large] at ([xshift=-5mm,yshift=-0.62cm]frame.north east) {Probabilités};}]
\vspace*{1.75cm}
\begin{minipage}[t][6.9cm][c]{\linewidth}\Large\centering
\begin{minipage}{0.95\linewidth}\raggedright\linespread{1.25}\selectfont Une urne contient $4$ boules dont $1$ rouges. On tire une boule, on la remet, puis on en tire une deuxième. Quelle est la probabilité de tirer deux boules rouges~?\end{minipage}
\end{minipage}

\vfill
\begin{tcolorbox}[enhanced,nobeforeafter,width=7.1cm,height=2.1cm,valign=center,colback=cpro!6,colframe=cpro,arc=2.5mm,boxrule=0.9pt,left=3mm]{\Large\bfseries\color{cpro}a)}\quad\begin{minipage}[c]{5.3cm}\raggedright\Large $\dfrac{1}{16}$\end{minipage}\end{tcolorbox}\hfill\begin{tcolorbox}[enhanced,nobeforeafter,width=7.1cm,height=2.1cm,valign=center,colback=cpro!6,colframe=cpro,arc=2.5mm,boxrule=0.9pt,left=3mm]{\Large\bfseries\color{cpro}b)}\quad\begin{minipage}[c]{5.3cm}\raggedright\Large $\dfrac{1}{4}$\end{minipage}\end{tcolorbox}\par\vspace{2.5mm}
\begin{tcolorbox}[enhanced,nobeforeafter,width=7.1cm,height=2.1cm,valign=center,colback=cpro!6,colframe=cpro,arc=2.5mm,boxrule=0.9pt,left=3mm]{\Large\bfseries\color{cpro}c)}\quad\begin{minipage}[c]{5.3cm}\raggedright\Large $\dfrac{1}{8}$\end{minipage}\end{tcolorbox}\hfill\begin{tcolorbox}[enhanced,nobeforeafter,width=7.1cm,height=2.1cm,valign=center,colback=cpro!6,colframe=cpro,arc=2.5mm,boxrule=0.9pt,left=3mm]{\Large\bfseries\color{cpro}d)}\quad\begin{minipage}[c]{5.3cm}\raggedright\Large $\dfrac{1}{2}$\end{minipage}\end{tcolorbox}
\end{tcolorbox}
\newpage
\begin{tcolorbox}[enhanced,width=16cm,height=14cm,colback=white,colframe=cpro,boxrule=2pt,arc=0mm,
left=6mm,right=6mm,top=0mm,bottom=4mm,sharp corners,
overlay={\fill[cpro] (frame.north west) rectangle ([yshift=-1.25cm]frame.north east);
\node[anchor=west,text=white,font=\sffamily\bfseries\LARGE] at ([xshift=5mm,yshift=-0.62cm]frame.north west) {Question 95 B};
\node[anchor=east,text=white,font=\sffamily\itshape\large] at ([xshift=-5mm,yshift=-0.62cm]frame.north east) {Probabilités};}]
\vspace*{1.75cm}
\begin{minipage}[t][6.9cm][c]{\linewidth}\LARGE\centering
\begin{minipage}{0.95\linewidth}\raggedright\linespread{1.25}\selectfont $A$ et $B$ sont deux événements incompatibles, avec $P(A)=0{,}15$ et $P(B)=0{,}6$. Que vaut $P(A\cup B)$~?\end{minipage}
\end{minipage}

\vfill
\begin{tcolorbox}[enhanced,nobeforeafter,width=7.1cm,height=2.1cm,valign=center,colback=cpro!6,colframe=cpro,arc=2.5mm,boxrule=0.9pt,left=3mm]{\Large\bfseries\color{cpro}a)}\quad\begin{minipage}[c]{5.3cm}\raggedright\Large $0{,}75$\end{minipage}\end{tcolorbox}\hfill\begin{tcolorbox}[enhanced,nobeforeafter,width=7.1cm,height=2.1cm,valign=center,colback=cpro!6,colframe=cpro,arc=2.5mm,boxrule=0.9pt,left=3mm]{\Large\bfseries\color{cpro}b)}\quad\begin{minipage}[c]{5.3cm}\raggedright\Large $0{,}66$\end{minipage}\end{tcolorbox}\par\vspace{2.5mm}
\begin{tcolorbox}[enhanced,nobeforeafter,width=7.1cm,height=2.1cm,valign=center,colback=cpro!6,colframe=cpro,arc=2.5mm,boxrule=0.9pt,left=3mm]{\Large\bfseries\color{cpro}c)}\quad\begin{minipage}[c]{5.3cm}\raggedright\Large $0{,}09$\end{minipage}\end{tcolorbox}\hfill\begin{tcolorbox}[enhanced,nobeforeafter,width=7.1cm,height=2.1cm,valign=center,colback=cpro!6,colframe=cpro,arc=2.5mm,boxrule=0.9pt,left=3mm]{\Large\bfseries\color{cpro}d)}\quad\begin{minipage}[c]{5.3cm}\raggedright\Large $0{,}45$\end{minipage}\end{tcolorbox}
\end{tcolorbox}
\newpage
\begin{tcolorbox}[enhanced,width=16cm,height=14cm,colback=white,colframe=cpro,boxrule=2pt,arc=0mm,
left=6mm,right=6mm,top=0mm,bottom=4mm,sharp corners,
overlay={\fill[cpro] (frame.north west) rectangle ([yshift=-1.25cm]frame.north east);
\node[anchor=west,text=white,font=\sffamily\bfseries\LARGE] at ([xshift=5mm,yshift=-0.62cm]frame.north west) {Question 96 B};
\node[anchor=east,text=white,font=\sffamily\itshape\large] at ([xshift=-5mm,yshift=-0.62cm]frame.north east) {Probabilités};}]
\vspace*{1.75cm}
\begin{minipage}[t][6.9cm][c]{\linewidth}\LARGE\centering
\begin{minipage}{0.95\linewidth}\raggedright\linespread{1.25}\selectfont $A$ et $B$ sont deux événements incompatibles, avec $P(A)=\dfrac{2}{5}$ et $P(B)=\dfrac{1}{6}$. Que vaut $P(A\cup B)$~?\end{minipage}
\end{minipage}

\vfill
\begin{tcolorbox}[enhanced,nobeforeafter,width=7.1cm,height=2.1cm,valign=center,colback=cpro!6,colframe=cpro,arc=2.5mm,boxrule=0.9pt,left=3mm]{\Large\bfseries\color{cpro}a)}\quad\begin{minipage}[c]{5.3cm}\raggedright\Large $\dfrac{7}{30}$\end{minipage}\end{tcolorbox}\hfill\begin{tcolorbox}[enhanced,nobeforeafter,width=7.1cm,height=2.1cm,valign=center,colback=cpro!6,colframe=cpro,arc=2.5mm,boxrule=0.9pt,left=3mm]{\Large\bfseries\color{cpro}b)}\quad\begin{minipage}[c]{5.3cm}\raggedright\Large $\dfrac{1}{15}$\end{minipage}\end{tcolorbox}\par\vspace{2.5mm}
\begin{tcolorbox}[enhanced,nobeforeafter,width=7.1cm,height=2.1cm,valign=center,colback=cpro!6,colframe=cpro,arc=2.5mm,boxrule=0.9pt,left=3mm]{\Large\bfseries\color{cpro}c)}\quad\begin{minipage}[c]{5.3cm}\raggedright\Large $\dfrac{17}{30}$\end{minipage}\end{tcolorbox}\hfill\begin{tcolorbox}[enhanced,nobeforeafter,width=7.1cm,height=2.1cm,valign=center,colback=cpro!6,colframe=cpro,arc=2.5mm,boxrule=0.9pt,left=3mm]{\Large\bfseries\color{cpro}d)}\quad\begin{minipage}[c]{5.3cm}\raggedright\Large $\dfrac{1}{2}$\end{minipage}\end{tcolorbox}
\end{tcolorbox}
\newpage
\begin{tcolorbox}[enhanced,width=16cm,height=14cm,colback=white,colframe=cnom,boxrule=2pt,arc=0mm,
left=6mm,right=6mm,top=0mm,bottom=4mm,sharp corners,
overlay={\fill[cnom] (frame.north west) rectangle ([yshift=-1.25cm]frame.north east);
\node[anchor=west,text=white,font=\sffamily\bfseries\LARGE] at ([xshift=5mm,yshift=-0.62cm]frame.north west) {Question 1 C};
\node[anchor=east,text=white,font=\sffamily\itshape\large] at ([xshift=-5mm,yshift=-0.62cm]frame.north east) {Nombres et calculs};}]
\vspace*{1.75cm}
\begin{minipage}[t][6.9cm][c]{\linewidth}\LARGE\centering
\begin{minipage}{0.95\linewidth}\raggedright\linespread{1.25}\selectfont Écrire $\sqrt{45}$ sous la forme $a\sqrt{b}$, avec $b$ le plus petit possible.\end{minipage}
\end{minipage}

\vfill
\begin{tcolorbox}[enhanced,nobeforeafter,width=7.1cm,height=2.1cm,valign=center,colback=cnom!6,colframe=cnom,arc=2.5mm,boxrule=0.9pt,left=3mm]{\Large\bfseries\color{cnom}a)}\quad\begin{minipage}[c]{5.3cm}\raggedright\Large $5\sqrt{3}$\end{minipage}\end{tcolorbox}\hfill\begin{tcolorbox}[enhanced,nobeforeafter,width=7.1cm,height=2.1cm,valign=center,colback=cnom!6,colframe=cnom,arc=2.5mm,boxrule=0.9pt,left=3mm]{\Large\bfseries\color{cnom}b)}\quad\begin{minipage}[c]{5.3cm}\raggedright\Large $3\sqrt{10}$\end{minipage}\end{tcolorbox}\par\vspace{2.5mm}
\begin{tcolorbox}[enhanced,nobeforeafter,width=7.1cm,height=2.1cm,valign=center,colback=cnom!6,colframe=cnom,arc=2.5mm,boxrule=0.9pt,left=3mm]{\Large\bfseries\color{cnom}c)}\quad\begin{minipage}[c]{5.3cm}\raggedright\Large $9\sqrt{5}$\end{minipage}\end{tcolorbox}\hfill\begin{tcolorbox}[enhanced,nobeforeafter,width=7.1cm,height=2.1cm,valign=center,colback=cnom!6,colframe=cnom,arc=2.5mm,boxrule=0.9pt,left=3mm]{\Large\bfseries\color{cnom}d)}\quad\begin{minipage}[c]{5.3cm}\raggedright\Large $3\sqrt{5}$\end{minipage}\end{tcolorbox}
\end{tcolorbox}
\newpage
\begin{tcolorbox}[enhanced,width=16cm,height=14cm,colback=white,colframe=cnom,boxrule=2pt,arc=0mm,
left=6mm,right=6mm,top=0mm,bottom=4mm,sharp corners,
overlay={\fill[cnom] (frame.north west) rectangle ([yshift=-1.25cm]frame.north east);
\node[anchor=west,text=white,font=\sffamily\bfseries\LARGE] at ([xshift=5mm,yshift=-0.62cm]frame.north west) {Question 2 C};
\node[anchor=east,text=white,font=\sffamily\itshape\large] at ([xshift=-5mm,yshift=-0.62cm]frame.north east) {Nombres et calculs};}]
\vspace*{1.75cm}
\begin{minipage}[t][6.9cm][c]{\linewidth}\LARGE\centering
\begin{minipage}{0.95\linewidth}\raggedright\linespread{1.25}\selectfont Écrire $\sqrt{98}$ sous la forme $a\sqrt{b}$, avec $b$ le plus petit possible.\end{minipage}
\end{minipage}

\vfill
\begin{tcolorbox}[enhanced,nobeforeafter,width=7.1cm,height=2.1cm,valign=center,colback=cnom!6,colframe=cnom,arc=2.5mm,boxrule=0.9pt,left=3mm]{\Large\bfseries\color{cnom}a)}\quad\begin{minipage}[c]{5.3cm}\raggedright\Large $2\sqrt{7}$\end{minipage}\end{tcolorbox}\hfill\begin{tcolorbox}[enhanced,nobeforeafter,width=7.1cm,height=2.1cm,valign=center,colback=cnom!6,colframe=cnom,arc=2.5mm,boxrule=0.9pt,left=3mm]{\Large\bfseries\color{cnom}b)}\quad\begin{minipage}[c]{5.3cm}\raggedright\Large $7\sqrt{4}$\end{minipage}\end{tcolorbox}\par\vspace{2.5mm}
\begin{tcolorbox}[enhanced,nobeforeafter,width=7.1cm,height=2.1cm,valign=center,colback=cnom!6,colframe=cnom,arc=2.5mm,boxrule=0.9pt,left=3mm]{\Large\bfseries\color{cnom}c)}\quad\begin{minipage}[c]{5.3cm}\raggedright\Large $49\sqrt{2}$\end{minipage}\end{tcolorbox}\hfill\begin{tcolorbox}[enhanced,nobeforeafter,width=7.1cm,height=2.1cm,valign=center,colback=cnom!6,colframe=cnom,arc=2.5mm,boxrule=0.9pt,left=3mm]{\Large\bfseries\color{cnom}d)}\quad\begin{minipage}[c]{5.3cm}\raggedright\Large $7\sqrt{2}$\end{minipage}\end{tcolorbox}
\end{tcolorbox}
\newpage
\begin{tcolorbox}[enhanced,width=16cm,height=14cm,colback=white,colframe=cnom,boxrule=2pt,arc=0mm,
left=6mm,right=6mm,top=0mm,bottom=4mm,sharp corners,
overlay={\fill[cnom] (frame.north west) rectangle ([yshift=-1.25cm]frame.north east);
\node[anchor=west,text=white,font=\sffamily\bfseries\LARGE] at ([xshift=5mm,yshift=-0.62cm]frame.north west) {Question 3 C};
\node[anchor=east,text=white,font=\sffamily\itshape\large] at ([xshift=-5mm,yshift=-0.62cm]frame.north east) {Nombres et calculs};}]
\vspace*{1.75cm}
\begin{minipage}[t][6.9cm][c]{\linewidth}\LARGE\centering
\begin{minipage}{0.95\linewidth}\raggedright\linespread{1.25}\selectfont Que vaut $\sqrt{27}+\sqrt{75}$~?\end{minipage}
\end{minipage}

\vfill
\begin{tcolorbox}[enhanced,nobeforeafter,width=7.1cm,height=2.1cm,valign=center,colback=cnom!6,colframe=cnom,arc=2.5mm,boxrule=0.9pt,left=3mm]{\Large\bfseries\color{cnom}a)}\quad\begin{minipage}[c]{5.3cm}\raggedright\Large $8\sqrt{3}$\end{minipage}\end{tcolorbox}\hfill\begin{tcolorbox}[enhanced,nobeforeafter,width=7.1cm,height=2.1cm,valign=center,colback=cnom!6,colframe=cnom,arc=2.5mm,boxrule=0.9pt,left=3mm]{\Large\bfseries\color{cnom}b)}\quad\begin{minipage}[c]{5.3cm}\raggedright\Large $8\sqrt{6}$\end{minipage}\end{tcolorbox}\par\vspace{2.5mm}
\begin{tcolorbox}[enhanced,nobeforeafter,width=7.1cm,height=2.1cm,valign=center,colback=cnom!6,colframe=cnom,arc=2.5mm,boxrule=0.9pt,left=3mm]{\Large\bfseries\color{cnom}c)}\quad\begin{minipage}[c]{5.3cm}\raggedright\Large $15\sqrt{3}$\end{minipage}\end{tcolorbox}\hfill\begin{tcolorbox}[enhanced,nobeforeafter,width=7.1cm,height=2.1cm,valign=center,colback=cnom!6,colframe=cnom,arc=2.5mm,boxrule=0.9pt,left=3mm]{\Large\bfseries\color{cnom}d)}\quad\begin{minipage}[c]{5.3cm}\raggedright\Large $\sqrt{102}$\end{minipage}\end{tcolorbox}
\end{tcolorbox}
\newpage
\begin{tcolorbox}[enhanced,width=16cm,height=14cm,colback=white,colframe=cnom,boxrule=2pt,arc=0mm,
left=6mm,right=6mm,top=0mm,bottom=4mm,sharp corners,
overlay={\fill[cnom] (frame.north west) rectangle ([yshift=-1.25cm]frame.north east);
\node[anchor=west,text=white,font=\sffamily\bfseries\LARGE] at ([xshift=5mm,yshift=-0.62cm]frame.north west) {Question 4 C};
\node[anchor=east,text=white,font=\sffamily\itshape\large] at ([xshift=-5mm,yshift=-0.62cm]frame.north east) {Nombres et calculs};}]
\vspace*{1.75cm}
\begin{minipage}[t][6.9cm][c]{\linewidth}\LARGE\centering
\begin{minipage}{0.95\linewidth}\raggedright\linespread{1.25}\selectfont Que vaut $\sqrt{18}+\sqrt{72}$~?\end{minipage}
\end{minipage}

\vfill
\begin{tcolorbox}[enhanced,nobeforeafter,width=7.1cm,height=2.1cm,valign=center,colback=cnom!6,colframe=cnom,arc=2.5mm,boxrule=0.9pt,left=3mm]{\Large\bfseries\color{cnom}a)}\quad\begin{minipage}[c]{5.3cm}\raggedright\Large $9\sqrt{2}$\end{minipage}\end{tcolorbox}\hfill\begin{tcolorbox}[enhanced,nobeforeafter,width=7.1cm,height=2.1cm,valign=center,colback=cnom!6,colframe=cnom,arc=2.5mm,boxrule=0.9pt,left=3mm]{\Large\bfseries\color{cnom}b)}\quad\begin{minipage}[c]{5.3cm}\raggedright\Large $18\sqrt{2}$\end{minipage}\end{tcolorbox}\par\vspace{2.5mm}
\begin{tcolorbox}[enhanced,nobeforeafter,width=7.1cm,height=2.1cm,valign=center,colback=cnom!6,colframe=cnom,arc=2.5mm,boxrule=0.9pt,left=3mm]{\Large\bfseries\color{cnom}c)}\quad\begin{minipage}[c]{5.3cm}\raggedright\Large $9\sqrt{4}$\end{minipage}\end{tcolorbox}\hfill\begin{tcolorbox}[enhanced,nobeforeafter,width=7.1cm,height=2.1cm,valign=center,colback=cnom!6,colframe=cnom,arc=2.5mm,boxrule=0.9pt,left=3mm]{\Large\bfseries\color{cnom}d)}\quad\begin{minipage}[c]{5.3cm}\raggedright\Large $\sqrt{90}$\end{minipage}\end{tcolorbox}
\end{tcolorbox}
\newpage
\begin{tcolorbox}[enhanced,width=16cm,height=14cm,colback=white,colframe=cnom,boxrule=2pt,arc=0mm,
left=6mm,right=6mm,top=0mm,bottom=4mm,sharp corners,
overlay={\fill[cnom] (frame.north west) rectangle ([yshift=-1.25cm]frame.north east);
\node[anchor=west,text=white,font=\sffamily\bfseries\LARGE] at ([xshift=5mm,yshift=-0.62cm]frame.north west) {Question 5 C};
\node[anchor=east,text=white,font=\sffamily\itshape\large] at ([xshift=-5mm,yshift=-0.62cm]frame.north east) {Nombres et calculs};}]
\vspace*{1.75cm}
\begin{minipage}[t][6.9cm][c]{\linewidth}\LARGE\centering
\begin{minipage}{0.95\linewidth}\raggedright\linespread{1.25}\selectfont Développer $(x+7)^2$.\end{minipage}
\end{minipage}

\vfill
\begin{tcolorbox}[enhanced,nobeforeafter,width=7.1cm,height=2.1cm,valign=center,colback=cnom!6,colframe=cnom,arc=2.5mm,boxrule=0.9pt,left=3mm]{\Large\bfseries\color{cnom}a)}\quad\begin{minipage}[c]{5.3cm}\raggedright\Large $x^2+14x-49$\end{minipage}\end{tcolorbox}\hfill\begin{tcolorbox}[enhanced,nobeforeafter,width=7.1cm,height=2.1cm,valign=center,colback=cnom!6,colframe=cnom,arc=2.5mm,boxrule=0.9pt,left=3mm]{\Large\bfseries\color{cnom}b)}\quad\begin{minipage}[c]{5.3cm}\raggedright\Large $x^2+14x+49$\end{minipage}\end{tcolorbox}\par\vspace{2.5mm}
\begin{tcolorbox}[enhanced,nobeforeafter,width=7.1cm,height=2.1cm,valign=center,colback=cnom!6,colframe=cnom,arc=2.5mm,boxrule=0.9pt,left=3mm]{\Large\bfseries\color{cnom}c)}\quad\begin{minipage}[c]{5.3cm}\raggedright\Large $x^2+7x+49$\end{minipage}\end{tcolorbox}\hfill\begin{tcolorbox}[enhanced,nobeforeafter,width=7.1cm,height=2.1cm,valign=center,colback=cnom!6,colframe=cnom,arc=2.5mm,boxrule=0.9pt,left=3mm]{\Large\bfseries\color{cnom}d)}\quad\begin{minipage}[c]{5.3cm}\raggedright\Large $x^2+49$\end{minipage}\end{tcolorbox}
\end{tcolorbox}
\newpage
\begin{tcolorbox}[enhanced,width=16cm,height=14cm,colback=white,colframe=cnom,boxrule=2pt,arc=0mm,
left=6mm,right=6mm,top=0mm,bottom=4mm,sharp corners,
overlay={\fill[cnom] (frame.north west) rectangle ([yshift=-1.25cm]frame.north east);
\node[anchor=west,text=white,font=\sffamily\bfseries\LARGE] at ([xshift=5mm,yshift=-0.62cm]frame.north west) {Question 6 C};
\node[anchor=east,text=white,font=\sffamily\itshape\large] at ([xshift=-5mm,yshift=-0.62cm]frame.north east) {Nombres et calculs};}]
\vspace*{1.75cm}
\begin{minipage}[t][6.9cm][c]{\linewidth}\LARGE\centering
\begin{minipage}{0.95\linewidth}\raggedright\linespread{1.25}\selectfont Développer $(2x-5)^2$.\end{minipage}
\end{minipage}

\vfill
\begin{tcolorbox}[enhanced,nobeforeafter,width=7.1cm,height=2.1cm,valign=center,colback=cnom!6,colframe=cnom,arc=2.5mm,boxrule=0.9pt,left=3mm]{\Large\bfseries\color{cnom}a)}\quad\begin{minipage}[c]{5.3cm}\raggedright\Large $4x^2+25$\end{minipage}\end{tcolorbox}\hfill\begin{tcolorbox}[enhanced,nobeforeafter,width=7.1cm,height=2.1cm,valign=center,colback=cnom!6,colframe=cnom,arc=2.5mm,boxrule=0.9pt,left=3mm]{\Large\bfseries\color{cnom}b)}\quad\begin{minipage}[c]{5.3cm}\raggedright\Large $2x^2-20x+25$\end{minipage}\end{tcolorbox}\par\vspace{2.5mm}
\begin{tcolorbox}[enhanced,nobeforeafter,width=7.1cm,height=2.1cm,valign=center,colback=cnom!6,colframe=cnom,arc=2.5mm,boxrule=0.9pt,left=3mm]{\Large\bfseries\color{cnom}c)}\quad\begin{minipage}[c]{5.3cm}\raggedright\Large $4x^2-20x+25$\end{minipage}\end{tcolorbox}\hfill\begin{tcolorbox}[enhanced,nobeforeafter,width=7.1cm,height=2.1cm,valign=center,colback=cnom!6,colframe=cnom,arc=2.5mm,boxrule=0.9pt,left=3mm]{\Large\bfseries\color{cnom}d)}\quad\begin{minipage}[c]{5.3cm}\raggedright\Large $4x^2-10x+25$\end{minipage}\end{tcolorbox}
\end{tcolorbox}
\newpage
\begin{tcolorbox}[enhanced,width=16cm,height=14cm,colback=white,colframe=cnom,boxrule=2pt,arc=0mm,
left=6mm,right=6mm,top=0mm,bottom=4mm,sharp corners,
overlay={\fill[cnom] (frame.north west) rectangle ([yshift=-1.25cm]frame.north east);
\node[anchor=west,text=white,font=\sffamily\bfseries\LARGE] at ([xshift=5mm,yshift=-0.62cm]frame.north west) {Question 7 C};
\node[anchor=east,text=white,font=\sffamily\itshape\large] at ([xshift=-5mm,yshift=-0.62cm]frame.north east) {Nombres et calculs};}]
\vspace*{1.75cm}
\begin{minipage}[t][6.9cm][c]{\linewidth}\LARGE\centering
\begin{minipage}{0.95\linewidth}\raggedright\linespread{1.25}\selectfont Factoriser $x^2-49$.\end{minipage}
\end{minipage}

\vfill
\begin{tcolorbox}[enhanced,nobeforeafter,width=7.1cm,height=2.1cm,valign=center,colback=cnom!6,colframe=cnom,arc=2.5mm,boxrule=0.9pt,left=3mm]{\Large\bfseries\color{cnom}a)}\quad\begin{minipage}[c]{5.3cm}\raggedright\Large $(x-7)(x+7)$\end{minipage}\end{tcolorbox}\hfill\begin{tcolorbox}[enhanced,nobeforeafter,width=7.1cm,height=2.1cm,valign=center,colback=cnom!6,colframe=cnom,arc=2.5mm,boxrule=0.9pt,left=3mm]{\Large\bfseries\color{cnom}b)}\quad\begin{minipage}[c]{5.3cm}\raggedright\Large $(x-49)(x+49)$\end{minipage}\end{tcolorbox}\par\vspace{2.5mm}
\begin{tcolorbox}[enhanced,nobeforeafter,width=7.1cm,height=2.1cm,valign=center,colback=cnom!6,colframe=cnom,arc=2.5mm,boxrule=0.9pt,left=3mm]{\Large\bfseries\color{cnom}c)}\quad\begin{minipage}[c]{5.3cm}\raggedright\Large $(x+7)^2$\end{minipage}\end{tcolorbox}\hfill\begin{tcolorbox}[enhanced,nobeforeafter,width=7.1cm,height=2.1cm,valign=center,colback=cnom!6,colframe=cnom,arc=2.5mm,boxrule=0.9pt,left=3mm]{\Large\bfseries\color{cnom}d)}\quad\begin{minipage}[c]{5.3cm}\raggedright\Large $(x-7)^2$\end{minipage}\end{tcolorbox}
\end{tcolorbox}
\newpage
\begin{tcolorbox}[enhanced,width=16cm,height=14cm,colback=white,colframe=cnom,boxrule=2pt,arc=0mm,
left=6mm,right=6mm,top=0mm,bottom=4mm,sharp corners,
overlay={\fill[cnom] (frame.north west) rectangle ([yshift=-1.25cm]frame.north east);
\node[anchor=west,text=white,font=\sffamily\bfseries\LARGE] at ([xshift=5mm,yshift=-0.62cm]frame.north west) {Question 8 C};
\node[anchor=east,text=white,font=\sffamily\itshape\large] at ([xshift=-5mm,yshift=-0.62cm]frame.north east) {Nombres et calculs};}]
\vspace*{1.75cm}
\begin{minipage}[t][6.9cm][c]{\linewidth}\LARGE\centering
\begin{minipage}{0.95\linewidth}\raggedright\linespread{1.25}\selectfont Factoriser $25x^2-4$.\end{minipage}
\end{minipage}

\vfill
\begin{tcolorbox}[enhanced,nobeforeafter,width=7.1cm,height=2.1cm,valign=center,colback=cnom!6,colframe=cnom,arc=2.5mm,boxrule=0.9pt,left=3mm]{\Large\bfseries\color{cnom}a)}\quad\begin{minipage}[c]{5.3cm}\raggedright\Large $(5x+2)^2$\end{minipage}\end{tcolorbox}\hfill\begin{tcolorbox}[enhanced,nobeforeafter,width=7.1cm,height=2.1cm,valign=center,colback=cnom!6,colframe=cnom,arc=2.5mm,boxrule=0.9pt,left=3mm]{\Large\bfseries\color{cnom}b)}\quad\begin{minipage}[c]{5.3cm}\raggedright\Large $(5x-2)^2$\end{minipage}\end{tcolorbox}\par\vspace{2.5mm}
\begin{tcolorbox}[enhanced,nobeforeafter,width=7.1cm,height=2.1cm,valign=center,colback=cnom!6,colframe=cnom,arc=2.5mm,boxrule=0.9pt,left=3mm]{\Large\bfseries\color{cnom}c)}\quad\begin{minipage}[c]{5.3cm}\raggedright\Large $(5x-2)(5x+2)$\end{minipage}\end{tcolorbox}\hfill\begin{tcolorbox}[enhanced,nobeforeafter,width=7.1cm,height=2.1cm,valign=center,colback=cnom!6,colframe=cnom,arc=2.5mm,boxrule=0.9pt,left=3mm]{\Large\bfseries\color{cnom}d)}\quad\begin{minipage}[c]{5.3cm}\raggedright\Large $(25x-2)(25x+2)$\end{minipage}\end{tcolorbox}
\end{tcolorbox}
\newpage
\begin{tcolorbox}[enhanced,width=16cm,height=14cm,colback=white,colframe=cnom,boxrule=2pt,arc=0mm,
left=6mm,right=6mm,top=0mm,bottom=4mm,sharp corners,
overlay={\fill[cnom] (frame.north west) rectangle ([yshift=-1.25cm]frame.north east);
\node[anchor=west,text=white,font=\sffamily\bfseries\LARGE] at ([xshift=5mm,yshift=-0.62cm]frame.north west) {Question 9 C};
\node[anchor=east,text=white,font=\sffamily\itshape\large] at ([xshift=-5mm,yshift=-0.62cm]frame.north east) {Nombres et calculs};}]
\vspace*{1.75cm}
\begin{minipage}[t][6.9cm][c]{\linewidth}\LARGE\centering
\begin{minipage}{0.95\linewidth}\raggedright\linespread{1.25}\selectfont Factoriser $25x^2-10x+1$.\end{minipage}
\end{minipage}

\vfill
\begin{tcolorbox}[enhanced,nobeforeafter,width=7.1cm,height=2.1cm,valign=center,colback=cnom!6,colframe=cnom,arc=2.5mm,boxrule=0.9pt,left=3mm]{\Large\bfseries\color{cnom}a)}\quad\begin{minipage}[c]{5.3cm}\raggedright\Large $(5x-1)^2$\end{minipage}\end{tcolorbox}\hfill\begin{tcolorbox}[enhanced,nobeforeafter,width=7.1cm,height=2.1cm,valign=center,colback=cnom!6,colframe=cnom,arc=2.5mm,boxrule=0.9pt,left=3mm]{\Large\bfseries\color{cnom}b)}\quad\begin{minipage}[c]{5.3cm}\raggedright\Large $(5x+1)(5x-1)$\end{minipage}\end{tcolorbox}\par\vspace{2.5mm}
\begin{tcolorbox}[enhanced,nobeforeafter,width=7.1cm,height=2.1cm,valign=center,colback=cnom!6,colframe=cnom,arc=2.5mm,boxrule=0.9pt,left=3mm]{\Large\bfseries\color{cnom}c)}\quad\begin{minipage}[c]{5.3cm}\raggedright\Large $(5x+1)^2$\end{minipage}\end{tcolorbox}\hfill\begin{tcolorbox}[enhanced,nobeforeafter,width=7.1cm,height=2.1cm,valign=center,colback=cnom!6,colframe=cnom,arc=2.5mm,boxrule=0.9pt,left=3mm]{\Large\bfseries\color{cnom}d)}\quad\begin{minipage}[c]{5.3cm}\raggedright\Large $(25x-1)^2$\end{minipage}\end{tcolorbox}
\end{tcolorbox}
\newpage
\begin{tcolorbox}[enhanced,width=16cm,height=14cm,colback=white,colframe=cnom,boxrule=2pt,arc=0mm,
left=6mm,right=6mm,top=0mm,bottom=4mm,sharp corners,
overlay={\fill[cnom] (frame.north west) rectangle ([yshift=-1.25cm]frame.north east);
\node[anchor=west,text=white,font=\sffamily\bfseries\LARGE] at ([xshift=5mm,yshift=-0.62cm]frame.north west) {Question 10 C};
\node[anchor=east,text=white,font=\sffamily\itshape\large] at ([xshift=-5mm,yshift=-0.62cm]frame.north east) {Nombres et calculs};}]
\vspace*{1.75cm}
\begin{minipage}[t][6.9cm][c]{\linewidth}\LARGE\centering
\begin{minipage}{0.95\linewidth}\raggedright\linespread{1.25}\selectfont Factoriser $16x^2+24x+9$.\end{minipage}
\end{minipage}

\vfill
\begin{tcolorbox}[enhanced,nobeforeafter,width=7.1cm,height=2.1cm,valign=center,colback=cnom!6,colframe=cnom,arc=2.5mm,boxrule=0.9pt,left=3mm]{\Large\bfseries\color{cnom}a)}\quad\begin{minipage}[c]{5.3cm}\raggedright\Large $(16x+3)^2$\end{minipage}\end{tcolorbox}\hfill\begin{tcolorbox}[enhanced,nobeforeafter,width=7.1cm,height=2.1cm,valign=center,colback=cnom!6,colframe=cnom,arc=2.5mm,boxrule=0.9pt,left=3mm]{\Large\bfseries\color{cnom}b)}\quad\begin{minipage}[c]{5.3cm}\raggedright\Large $(4x+3)^2$\end{minipage}\end{tcolorbox}\par\vspace{2.5mm}
\begin{tcolorbox}[enhanced,nobeforeafter,width=7.1cm,height=2.1cm,valign=center,colback=cnom!6,colframe=cnom,arc=2.5mm,boxrule=0.9pt,left=3mm]{\Large\bfseries\color{cnom}c)}\quad\begin{minipage}[c]{5.3cm}\raggedright\Large $(4x-3)(4x+3)$\end{minipage}\end{tcolorbox}\hfill\begin{tcolorbox}[enhanced,nobeforeafter,width=7.1cm,height=2.1cm,valign=center,colback=cnom!6,colframe=cnom,arc=2.5mm,boxrule=0.9pt,left=3mm]{\Large\bfseries\color{cnom}d)}\quad\begin{minipage}[c]{5.3cm}\raggedright\Large $(4x-3)^2$\end{minipage}\end{tcolorbox}
\end{tcolorbox}
\newpage
\begin{tcolorbox}[enhanced,width=16cm,height=14cm,colback=white,colframe=cnom,boxrule=2pt,arc=0mm,
left=6mm,right=6mm,top=0mm,bottom=4mm,sharp corners,
overlay={\fill[cnom] (frame.north west) rectangle ([yshift=-1.25cm]frame.north east);
\node[anchor=west,text=white,font=\sffamily\bfseries\LARGE] at ([xshift=5mm,yshift=-0.62cm]frame.north west) {Question 11 C};
\node[anchor=east,text=white,font=\sffamily\itshape\large] at ([xshift=-5mm,yshift=-0.62cm]frame.north east) {Nombres et calculs};}]
\vspace*{1.75cm}
\begin{minipage}[t][6.9cm][c]{\linewidth}\LARGE\centering
\begin{minipage}{0.95\linewidth}\raggedright\linespread{1.25}\selectfont Que signifie $x\in]-1\,;\,7]$~?\end{minipage}
\end{minipage}

\vfill
\begin{tcolorbox}[enhanced,nobeforeafter,width=7.1cm,height=2.1cm,valign=center,colback=cnom!6,colframe=cnom,arc=2.5mm,boxrule=0.9pt,left=3mm]{\Large\bfseries\color{cnom}a)}\quad\begin{minipage}[c]{5.3cm}\raggedright\Large $-1\leqslant x\leqslant 7$\end{minipage}\end{tcolorbox}\hfill\begin{tcolorbox}[enhanced,nobeforeafter,width=7.1cm,height=2.1cm,valign=center,colback=cnom!6,colframe=cnom,arc=2.5mm,boxrule=0.9pt,left=3mm]{\Large\bfseries\color{cnom}b)}\quad\begin{minipage}[c]{5.3cm}\raggedright\Large $-1\leqslant x<7$\end{minipage}\end{tcolorbox}\par\vspace{2.5mm}
\begin{tcolorbox}[enhanced,nobeforeafter,width=7.1cm,height=2.1cm,valign=center,colback=cnom!6,colframe=cnom,arc=2.5mm,boxrule=0.9pt,left=3mm]{\Large\bfseries\color{cnom}c)}\quad\begin{minipage}[c]{5.3cm}\raggedright\Large $-1<x<7$\end{minipage}\end{tcolorbox}\hfill\begin{tcolorbox}[enhanced,nobeforeafter,width=7.1cm,height=2.1cm,valign=center,colback=cnom!6,colframe=cnom,arc=2.5mm,boxrule=0.9pt,left=3mm]{\Large\bfseries\color{cnom}d)}\quad\begin{minipage}[c]{5.3cm}\raggedright\Large $-1<x\leqslant 7$\end{minipage}\end{tcolorbox}
\end{tcolorbox}
\newpage
\begin{tcolorbox}[enhanced,width=16cm,height=14cm,colback=white,colframe=cnom,boxrule=2pt,arc=0mm,
left=6mm,right=6mm,top=0mm,bottom=4mm,sharp corners,
overlay={\fill[cnom] (frame.north west) rectangle ([yshift=-1.25cm]frame.north east);
\node[anchor=west,text=white,font=\sffamily\bfseries\LARGE] at ([xshift=5mm,yshift=-0.62cm]frame.north west) {Question 12 C};
\node[anchor=east,text=white,font=\sffamily\itshape\large] at ([xshift=-5mm,yshift=-0.62cm]frame.north east) {Nombres et calculs};}]
\vspace*{1.75cm}
\begin{minipage}[t][6.9cm][c]{\linewidth}\LARGE\centering
\begin{minipage}{0.95\linewidth}\raggedright\linespread{1.25}\selectfont Que signifie $x\in[-6\,;\,+\infty[$~?\end{minipage}
\end{minipage}

\vfill
\begin{tcolorbox}[enhanced,nobeforeafter,width=7.1cm,height=2.1cm,valign=center,colback=cnom!6,colframe=cnom,arc=2.5mm,boxrule=0.9pt,left=3mm]{\Large\bfseries\color{cnom}a)}\quad\begin{minipage}[c]{5.3cm}\raggedright\Large $x\geqslant -6$\end{minipage}\end{tcolorbox}\hfill\begin{tcolorbox}[enhanced,nobeforeafter,width=7.1cm,height=2.1cm,valign=center,colback=cnom!6,colframe=cnom,arc=2.5mm,boxrule=0.9pt,left=3mm]{\Large\bfseries\color{cnom}b)}\quad\begin{minipage}[c]{5.3cm}\raggedright\Large $x<-6$\end{minipage}\end{tcolorbox}\par\vspace{2.5mm}
\begin{tcolorbox}[enhanced,nobeforeafter,width=7.1cm,height=2.1cm,valign=center,colback=cnom!6,colframe=cnom,arc=2.5mm,boxrule=0.9pt,left=3mm]{\Large\bfseries\color{cnom}c)}\quad\begin{minipage}[c]{5.3cm}\raggedright\Large $x\leqslant -6$\end{minipage}\end{tcolorbox}\hfill\begin{tcolorbox}[enhanced,nobeforeafter,width=7.1cm,height=2.1cm,valign=center,colback=cnom!6,colframe=cnom,arc=2.5mm,boxrule=0.9pt,left=3mm]{\Large\bfseries\color{cnom}d)}\quad\begin{minipage}[c]{5.3cm}\raggedright\Large $x>-6$\end{minipage}\end{tcolorbox}
\end{tcolorbox}
\newpage
\begin{tcolorbox}[enhanced,width=16cm,height=14cm,colback=white,colframe=cnom,boxrule=2pt,arc=0mm,
left=6mm,right=6mm,top=0mm,bottom=4mm,sharp corners,
overlay={\fill[cnom] (frame.north west) rectangle ([yshift=-1.25cm]frame.north east);
\node[anchor=west,text=white,font=\sffamily\bfseries\LARGE] at ([xshift=5mm,yshift=-0.62cm]frame.north west) {Question 13 C};
\node[anchor=east,text=white,font=\sffamily\itshape\large] at ([xshift=-5mm,yshift=-0.62cm]frame.north east) {Nombres et calculs};}]
\vspace*{1.75cm}
\begin{minipage}[t][6.9cm][c]{\linewidth}\LARGE\centering
\begin{minipage}{0.95\linewidth}\raggedright\linespread{1.25}\selectfont Pour quels réels $x$ a-t-on $|x+4|\leqslant 1$~?\end{minipage}
\end{minipage}

\vfill
\begin{tcolorbox}[enhanced,nobeforeafter,width=7.1cm,height=2.1cm,valign=center,colback=cnom!6,colframe=cnom,arc=2.5mm,boxrule=0.9pt,left=3mm]{\Large\bfseries\color{cnom}a)}\quad\begin{minipage}[c]{5.3cm}\raggedright\Large $x\in[-6\,;\,-2]$\end{minipage}\end{tcolorbox}\hfill\begin{tcolorbox}[enhanced,nobeforeafter,width=7.1cm,height=2.1cm,valign=center,colback=cnom!6,colframe=cnom,arc=2.5mm,boxrule=0.9pt,left=3mm]{\Large\bfseries\color{cnom}b)}\quad\begin{minipage}[c]{5.3cm}\raggedright\Large $x\in[3\,;\,5]$\end{minipage}\end{tcolorbox}\par\vspace{2.5mm}
\begin{tcolorbox}[enhanced,nobeforeafter,width=7.1cm,height=2.1cm,valign=center,colback=cnom!6,colframe=cnom,arc=2.5mm,boxrule=0.9pt,left=3mm]{\Large\bfseries\color{cnom}c)}\quad\begin{minipage}[c]{5.3cm}\raggedright\Large $x\in[-4\,;\,-3]$\end{minipage}\end{tcolorbox}\hfill\begin{tcolorbox}[enhanced,nobeforeafter,width=7.1cm,height=2.1cm,valign=center,colback=cnom!6,colframe=cnom,arc=2.5mm,boxrule=0.9pt,left=3mm]{\Large\bfseries\color{cnom}d)}\quad\begin{minipage}[c]{5.3cm}\raggedright\Large $x\in[-5\,;\,-3]$\end{minipage}\end{tcolorbox}
\end{tcolorbox}
\newpage
\begin{tcolorbox}[enhanced,width=16cm,height=14cm,colback=white,colframe=cnom,boxrule=2pt,arc=0mm,
left=6mm,right=6mm,top=0mm,bottom=4mm,sharp corners,
overlay={\fill[cnom] (frame.north west) rectangle ([yshift=-1.25cm]frame.north east);
\node[anchor=west,text=white,font=\sffamily\bfseries\LARGE] at ([xshift=5mm,yshift=-0.62cm]frame.north west) {Question 14 C};
\node[anchor=east,text=white,font=\sffamily\itshape\large] at ([xshift=-5mm,yshift=-0.62cm]frame.north east) {Nombres et calculs};}]
\vspace*{1.75cm}
\begin{minipage}[t][6.9cm][c]{\linewidth}\LARGE\centering
\begin{minipage}{0.95\linewidth}\raggedright\linespread{1.25}\selectfont Pour quels réels $x$ a-t-on $|x-6|\leqslant 3$~?\end{minipage}
\end{minipage}

\vfill
\begin{tcolorbox}[enhanced,nobeforeafter,width=7.1cm,height=2.1cm,valign=center,colback=cnom!6,colframe=cnom,arc=2.5mm,boxrule=0.9pt,left=3mm]{\Large\bfseries\color{cnom}a)}\quad\begin{minipage}[c]{5.3cm}\raggedright\Large $x\in[-9\,;\,-3]$\end{minipage}\end{tcolorbox}\hfill\begin{tcolorbox}[enhanced,nobeforeafter,width=7.1cm,height=2.1cm,valign=center,colback=cnom!6,colframe=cnom,arc=2.5mm,boxrule=0.9pt,left=3mm]{\Large\bfseries\color{cnom}b)}\quad\begin{minipage}[c]{5.3cm}\raggedright\Large $x\in[6\,;\,9]$\end{minipage}\end{tcolorbox}\par\vspace{2.5mm}
\begin{tcolorbox}[enhanced,nobeforeafter,width=7.1cm,height=2.1cm,valign=center,colback=cnom!6,colframe=cnom,arc=2.5mm,boxrule=0.9pt,left=3mm]{\Large\bfseries\color{cnom}c)}\quad\begin{minipage}[c]{5.3cm}\raggedright\Large $x\in[3\,;\,9]$\end{minipage}\end{tcolorbox}\hfill\begin{tcolorbox}[enhanced,nobeforeafter,width=7.1cm,height=2.1cm,valign=center,colback=cnom!6,colframe=cnom,arc=2.5mm,boxrule=0.9pt,left=3mm]{\Large\bfseries\color{cnom}d)}\quad\begin{minipage}[c]{5.3cm}\raggedright\Large $x\in[0\,;\,12]$\end{minipage}\end{tcolorbox}
\end{tcolorbox}
\newpage
\begin{tcolorbox}[enhanced,width=16cm,height=14cm,colback=white,colframe=cnom,boxrule=2pt,arc=0mm,
left=6mm,right=6mm,top=0mm,bottom=4mm,sharp corners,
overlay={\fill[cnom] (frame.north west) rectangle ([yshift=-1.25cm]frame.north east);
\node[anchor=west,text=white,font=\sffamily\bfseries\LARGE] at ([xshift=5mm,yshift=-0.62cm]frame.north west) {Question 15 C};
\node[anchor=east,text=white,font=\sffamily\itshape\large] at ([xshift=-5mm,yshift=-0.62cm]frame.north east) {Nombres et calculs};}]
\vspace*{1.75cm}
\begin{minipage}[t][6.9cm][c]{\linewidth}\LARGE\centering
\begin{minipage}{0.95\linewidth}\raggedright\linespread{1.25}\selectfont Écrire $\dfrac{7^{2}\times7^{2}}{7^{6}}$ sous la forme d'une puissance de $7$.\end{minipage}
\end{minipage}

\vfill
\begin{tcolorbox}[enhanced,nobeforeafter,width=7.1cm,height=2.1cm,valign=center,colback=cnom!6,colframe=cnom,arc=2.5mm,boxrule=0.9pt,left=3mm]{\Large\bfseries\color{cnom}a)}\quad\begin{minipage}[c]{5.3cm}\raggedright\Large $7^{-1}$\end{minipage}\end{tcolorbox}\hfill\begin{tcolorbox}[enhanced,nobeforeafter,width=7.1cm,height=2.1cm,valign=center,colback=cnom!6,colframe=cnom,arc=2.5mm,boxrule=0.9pt,left=3mm]{\Large\bfseries\color{cnom}b)}\quad\begin{minipage}[c]{5.3cm}\raggedright\Large $7^{-2}$\end{minipage}\end{tcolorbox}\par\vspace{2.5mm}
\begin{tcolorbox}[enhanced,nobeforeafter,width=7.1cm,height=2.1cm,valign=center,colback=cnom!6,colframe=cnom,arc=2.5mm,boxrule=0.9pt,left=3mm]{\Large\bfseries\color{cnom}c)}\quad\begin{minipage}[c]{5.3cm}\raggedright\Large $7^{-6}$\end{minipage}\end{tcolorbox}\hfill\begin{tcolorbox}[enhanced,nobeforeafter,width=7.1cm,height=2.1cm,valign=center,colback=cnom!6,colframe=cnom,arc=2.5mm,boxrule=0.9pt,left=3mm]{\Large\bfseries\color{cnom}d)}\quad\begin{minipage}[c]{5.3cm}\raggedright\Large $7^{10}$\end{minipage}\end{tcolorbox}
\end{tcolorbox}
\newpage
\begin{tcolorbox}[enhanced,width=16cm,height=14cm,colback=white,colframe=cnom,boxrule=2pt,arc=0mm,
left=6mm,right=6mm,top=0mm,bottom=4mm,sharp corners,
overlay={\fill[cnom] (frame.north west) rectangle ([yshift=-1.25cm]frame.north east);
\node[anchor=west,text=white,font=\sffamily\bfseries\LARGE] at ([xshift=5mm,yshift=-0.62cm]frame.north west) {Question 16 C};
\node[anchor=east,text=white,font=\sffamily\itshape\large] at ([xshift=-5mm,yshift=-0.62cm]frame.north east) {Nombres et calculs};}]
\vspace*{1.75cm}
\begin{minipage}[t][6.9cm][c]{\linewidth}\LARGE\centering
\begin{minipage}{0.95\linewidth}\raggedright\linespread{1.25}\selectfont Écrire $\dfrac{2^{-1}\times2^{-4}}{2^{-8}}$ sous la forme d'une puissance de $2$.\end{minipage}
\end{minipage}

\vfill
\begin{tcolorbox}[enhanced,nobeforeafter,width=7.1cm,height=2.1cm,valign=center,colback=cnom!6,colframe=cnom,arc=2.5mm,boxrule=0.9pt,left=3mm]{\Large\bfseries\color{cnom}a)}\quad\begin{minipage}[c]{5.3cm}\raggedright\Large $2^{11}$\end{minipage}\end{tcolorbox}\hfill\begin{tcolorbox}[enhanced,nobeforeafter,width=7.1cm,height=2.1cm,valign=center,colback=cnom!6,colframe=cnom,arc=2.5mm,boxrule=0.9pt,left=3mm]{\Large\bfseries\color{cnom}b)}\quad\begin{minipage}[c]{5.3cm}\raggedright\Large $2^{3}$\end{minipage}\end{tcolorbox}\par\vspace{2.5mm}
\begin{tcolorbox}[enhanced,nobeforeafter,width=7.1cm,height=2.1cm,valign=center,colback=cnom!6,colframe=cnom,arc=2.5mm,boxrule=0.9pt,left=3mm]{\Large\bfseries\color{cnom}c)}\quad\begin{minipage}[c]{5.3cm}\raggedright\Large $2^{12}$\end{minipage}\end{tcolorbox}\hfill\begin{tcolorbox}[enhanced,nobeforeafter,width=7.1cm,height=2.1cm,valign=center,colback=cnom!6,colframe=cnom,arc=2.5mm,boxrule=0.9pt,left=3mm]{\Large\bfseries\color{cnom}d)}\quad\begin{minipage}[c]{5.3cm}\raggedright\Large $2^{-13}$\end{minipage}\end{tcolorbox}
\end{tcolorbox}
\newpage
\begin{tcolorbox}[enhanced,width=16cm,height=14cm,colback=white,colframe=csta,boxrule=2pt,arc=0mm,
left=6mm,right=6mm,top=0mm,bottom=4mm,sharp corners,
overlay={\fill[csta] (frame.north west) rectangle ([yshift=-1.25cm]frame.north east);
\node[anchor=west,text=white,font=\sffamily\bfseries\LARGE] at ([xshift=5mm,yshift=-0.62cm]frame.north west) {Question 17 C};
\node[anchor=east,text=white,font=\sffamily\itshape\large] at ([xshift=-5mm,yshift=-0.62cm]frame.north east) {Évolutions et statistiques};}]
\vspace*{1.75cm}
\begin{minipage}[t][6.9cm][c]{\linewidth}\LARGE\centering
\begin{minipage}{0.95\linewidth}\raggedright\linespread{1.25}\selectfont Quel est le coefficient multiplicateur associé à une hausse de $150\,\%$~?\end{minipage}
\end{minipage}

\vfill
\begin{tcolorbox}[enhanced,nobeforeafter,width=7.1cm,height=2.1cm,valign=center,colback=csta!6,colframe=csta,arc=2.5mm,boxrule=0.9pt,left=3mm]{\Large\bfseries\color{csta}a)}\quad\begin{minipage}[c]{5.3cm}\raggedright\Large $-0{,}5$\end{minipage}\end{tcolorbox}\hfill\begin{tcolorbox}[enhanced,nobeforeafter,width=7.1cm,height=2.1cm,valign=center,colback=csta!6,colframe=csta,arc=2.5mm,boxrule=0.9pt,left=3mm]{\Large\bfseries\color{csta}b)}\quad\begin{minipage}[c]{5.3cm}\raggedright\Large $2{,}5$\end{minipage}\end{tcolorbox}\par\vspace{2.5mm}
\begin{tcolorbox}[enhanced,nobeforeafter,width=7.1cm,height=2.1cm,valign=center,colback=csta!6,colframe=csta,arc=2.5mm,boxrule=0.9pt,left=3mm]{\Large\bfseries\color{csta}c)}\quad\begin{minipage}[c]{5.3cm}\raggedright\Large $16$\end{minipage}\end{tcolorbox}\hfill\begin{tcolorbox}[enhanced,nobeforeafter,width=7.1cm,height=2.1cm,valign=center,colback=csta!6,colframe=csta,arc=2.5mm,boxrule=0.9pt,left=3mm]{\Large\bfseries\color{csta}d)}\quad\begin{minipage}[c]{5.3cm}\raggedright\Large $1{,}5$\end{minipage}\end{tcolorbox}
\end{tcolorbox}
\newpage
\begin{tcolorbox}[enhanced,width=16cm,height=14cm,colback=white,colframe=csta,boxrule=2pt,arc=0mm,
left=6mm,right=6mm,top=0mm,bottom=4mm,sharp corners,
overlay={\fill[csta] (frame.north west) rectangle ([yshift=-1.25cm]frame.north east);
\node[anchor=west,text=white,font=\sffamily\bfseries\LARGE] at ([xshift=5mm,yshift=-0.62cm]frame.north west) {Question 18 C};
\node[anchor=east,text=white,font=\sffamily\itshape\large] at ([xshift=-5mm,yshift=-0.62cm]frame.north east) {Évolutions et statistiques};}]
\vspace*{1.75cm}
\begin{minipage}[t][6.9cm][c]{\linewidth}\LARGE\centering
\begin{minipage}{0.95\linewidth}\raggedright\linespread{1.25}\selectfont Quel est le coefficient multiplicateur associé à une baisse de $2\,\%$~?\end{minipage}
\end{minipage}

\vfill
\begin{tcolorbox}[enhanced,nobeforeafter,width=7.1cm,height=2.1cm,valign=center,colback=csta!6,colframe=csta,arc=2.5mm,boxrule=0.9pt,left=3mm]{\Large\bfseries\color{csta}a)}\quad\begin{minipage}[c]{5.3cm}\raggedright\Large $0{,}98$\end{minipage}\end{tcolorbox}\hfill\begin{tcolorbox}[enhanced,nobeforeafter,width=7.1cm,height=2.1cm,valign=center,colback=csta!6,colframe=csta,arc=2.5mm,boxrule=0.9pt,left=3mm]{\Large\bfseries\color{csta}b)}\quad\begin{minipage}[c]{5.3cm}\raggedright\Large $-0{,}02$\end{minipage}\end{tcolorbox}\par\vspace{2.5mm}
\begin{tcolorbox}[enhanced,nobeforeafter,width=7.1cm,height=2.1cm,valign=center,colback=csta!6,colframe=csta,arc=2.5mm,boxrule=0.9pt,left=3mm]{\Large\bfseries\color{csta}c)}\quad\begin{minipage}[c]{5.3cm}\raggedright\Large $1{,}02$\end{minipage}\end{tcolorbox}\hfill\begin{tcolorbox}[enhanced,nobeforeafter,width=7.1cm,height=2.1cm,valign=center,colback=csta!6,colframe=csta,arc=2.5mm,boxrule=0.9pt,left=3mm]{\Large\bfseries\color{csta}d)}\quad\begin{minipage}[c]{5.3cm}\raggedright\Large $0{,}8$\end{minipage}\end{tcolorbox}
\end{tcolorbox}
\newpage
\begin{tcolorbox}[enhanced,width=16cm,height=14cm,colback=white,colframe=csta,boxrule=2pt,arc=0mm,
left=6mm,right=6mm,top=0mm,bottom=4mm,sharp corners,
overlay={\fill[csta] (frame.north west) rectangle ([yshift=-1.25cm]frame.north east);
\node[anchor=west,text=white,font=\sffamily\bfseries\LARGE] at ([xshift=5mm,yshift=-0.62cm]frame.north west) {Question 19 C};
\node[anchor=east,text=white,font=\sffamily\itshape\large] at ([xshift=-5mm,yshift=-0.62cm]frame.north east) {Évolutions et statistiques};}]
\vspace*{1.75cm}
\begin{minipage}[t][6.9cm][c]{\linewidth}\LARGE\centering
\begin{minipage}{0.95\linewidth}\raggedright\linespread{1.25}\selectfont Un prix de $600$ € augmente de $8\,\%$. Quel est le nouveau prix~?\end{minipage}
\end{minipage}

\vfill
\begin{tcolorbox}[enhanced,nobeforeafter,width=7.1cm,height=2.1cm,valign=center,colback=csta!6,colframe=csta,arc=2.5mm,boxrule=0.9pt,left=3mm]{\Large\bfseries\color{csta}a)}\quad\begin{minipage}[c]{5.3cm}\raggedright\Large $608~\text{€}$\end{minipage}\end{tcolorbox}\hfill\begin{tcolorbox}[enhanced,nobeforeafter,width=7.1cm,height=2.1cm,valign=center,colback=csta!6,colframe=csta,arc=2.5mm,boxrule=0.9pt,left=3mm]{\Large\bfseries\color{csta}b)}\quad\begin{minipage}[c]{5.3cm}\raggedright\Large $648~\text{€}$\end{minipage}\end{tcolorbox}\par\vspace{2.5mm}
\begin{tcolorbox}[enhanced,nobeforeafter,width=7.1cm,height=2.1cm,valign=center,colback=csta!6,colframe=csta,arc=2.5mm,boxrule=0.9pt,left=3mm]{\Large\bfseries\color{csta}c)}\quad\begin{minipage}[c]{5.3cm}\raggedright\Large $48~\text{€}$\end{minipage}\end{tcolorbox}\hfill\begin{tcolorbox}[enhanced,nobeforeafter,width=7.1cm,height=2.1cm,valign=center,colback=csta!6,colframe=csta,arc=2.5mm,boxrule=0.9pt,left=3mm]{\Large\bfseries\color{csta}d)}\quad\begin{minipage}[c]{5.3cm}\raggedright\Large $552~\text{€}$\end{minipage}\end{tcolorbox}
\end{tcolorbox}
\newpage
\begin{tcolorbox}[enhanced,width=16cm,height=14cm,colback=white,colframe=csta,boxrule=2pt,arc=0mm,
left=6mm,right=6mm,top=0mm,bottom=4mm,sharp corners,
overlay={\fill[csta] (frame.north west) rectangle ([yshift=-1.25cm]frame.north east);
\node[anchor=west,text=white,font=\sffamily\bfseries\LARGE] at ([xshift=5mm,yshift=-0.62cm]frame.north west) {Question 20 C};
\node[anchor=east,text=white,font=\sffamily\itshape\large] at ([xshift=-5mm,yshift=-0.62cm]frame.north east) {Évolutions et statistiques};}]
\vspace*{1.75cm}
\begin{minipage}[t][6.9cm][c]{\linewidth}\LARGE\centering
\begin{minipage}{0.95\linewidth}\raggedright\linespread{1.25}\selectfont Un prix de $3\,500$ € baisse de $35\,\%$. Quel est le nouveau prix~?\end{minipage}
\end{minipage}

\vfill
\begin{tcolorbox}[enhanced,nobeforeafter,width=7.1cm,height=2.1cm,valign=center,colback=csta!6,colframe=csta,arc=2.5mm,boxrule=0.9pt,left=3mm]{\Large\bfseries\color{csta}a)}\quad\begin{minipage}[c]{5.3cm}\raggedright\Large $4\,725~\text{€}$\end{minipage}\end{tcolorbox}\hfill\begin{tcolorbox}[enhanced,nobeforeafter,width=7.1cm,height=2.1cm,valign=center,colback=csta!6,colframe=csta,arc=2.5mm,boxrule=0.9pt,left=3mm]{\Large\bfseries\color{csta}b)}\quad\begin{minipage}[c]{5.3cm}\raggedright\Large $3\,465~\text{€}$\end{minipage}\end{tcolorbox}\par\vspace{2.5mm}
\begin{tcolorbox}[enhanced,nobeforeafter,width=7.1cm,height=2.1cm,valign=center,colback=csta!6,colframe=csta,arc=2.5mm,boxrule=0.9pt,left=3mm]{\Large\bfseries\color{csta}c)}\quad\begin{minipage}[c]{5.3cm}\raggedright\Large $2\,275~\text{€}$\end{minipage}\end{tcolorbox}\hfill\begin{tcolorbox}[enhanced,nobeforeafter,width=7.1cm,height=2.1cm,valign=center,colback=csta!6,colframe=csta,arc=2.5mm,boxrule=0.9pt,left=3mm]{\Large\bfseries\color{csta}d)}\quad\begin{minipage}[c]{5.3cm}\raggedright\Large $1\,225~\text{€}$\end{minipage}\end{tcolorbox}
\end{tcolorbox}
\newpage
\begin{tcolorbox}[enhanced,width=16cm,height=14cm,colback=white,colframe=csta,boxrule=2pt,arc=0mm,
left=6mm,right=6mm,top=0mm,bottom=4mm,sharp corners,
overlay={\fill[csta] (frame.north west) rectangle ([yshift=-1.25cm]frame.north east);
\node[anchor=west,text=white,font=\sffamily\bfseries\LARGE] at ([xshift=5mm,yshift=-0.62cm]frame.north west) {Question 21 C};
\node[anchor=east,text=white,font=\sffamily\itshape\large] at ([xshift=-5mm,yshift=-0.62cm]frame.north east) {Évolutions et statistiques};}]
\vspace*{1.75cm}
\begin{minipage}[t][6.9cm][c]{\linewidth}\LARGE\centering
\begin{minipage}{0.95\linewidth}\raggedright\linespread{1.25}\selectfont Une quantité passe de $125$ à $100$. Quel est son taux d'évolution~?\end{minipage}
\end{minipage}

\vfill
\begin{tcolorbox}[enhanced,nobeforeafter,width=7.1cm,height=2.1cm,valign=center,colback=csta!6,colframe=csta,arc=2.5mm,boxrule=0.9pt,left=3mm]{\Large\bfseries\color{csta}a)}\quad\begin{minipage}[c]{5.3cm}\raggedright\Large $+20\,\%$\end{minipage}\end{tcolorbox}\hfill\begin{tcolorbox}[enhanced,nobeforeafter,width=7.1cm,height=2.1cm,valign=center,colback=csta!6,colframe=csta,arc=2.5mm,boxrule=0.9pt,left=3mm]{\Large\bfseries\color{csta}b)}\quad\begin{minipage}[c]{5.3cm}\raggedright\Large $-20\,\%$\end{minipage}\end{tcolorbox}\par\vspace{2.5mm}
\begin{tcolorbox}[enhanced,nobeforeafter,width=7.1cm,height=2.1cm,valign=center,colback=csta!6,colframe=csta,arc=2.5mm,boxrule=0.9pt,left=3mm]{\Large\bfseries\color{csta}c)}\quad\begin{minipage}[c]{5.3cm}\raggedright\Large $-10\,\%$\end{minipage}\end{tcolorbox}\hfill\begin{tcolorbox}[enhanced,nobeforeafter,width=7.1cm,height=2.1cm,valign=center,colback=csta!6,colframe=csta,arc=2.5mm,boxrule=0.9pt,left=3mm]{\Large\bfseries\color{csta}d)}\quad\begin{minipage}[c]{5.3cm}\raggedright\Large $-25\,\%$\end{minipage}\end{tcolorbox}
\end{tcolorbox}
\newpage
\begin{tcolorbox}[enhanced,width=16cm,height=14cm,colback=white,colframe=csta,boxrule=2pt,arc=0mm,
left=6mm,right=6mm,top=0mm,bottom=4mm,sharp corners,
overlay={\fill[csta] (frame.north west) rectangle ([yshift=-1.25cm]frame.north east);
\node[anchor=west,text=white,font=\sffamily\bfseries\LARGE] at ([xshift=5mm,yshift=-0.62cm]frame.north west) {Question 22 C};
\node[anchor=east,text=white,font=\sffamily\itshape\large] at ([xshift=-5mm,yshift=-0.62cm]frame.north east) {Évolutions et statistiques};}]
\vspace*{1.75cm}
\begin{minipage}[t][6.9cm][c]{\linewidth}\LARGE\centering
\begin{minipage}{0.95\linewidth}\raggedright\linespread{1.25}\selectfont Une quantité passe de $250$ à $280$. Quel est son taux d'évolution~?\end{minipage}
\end{minipage}

\vfill
\begin{tcolorbox}[enhanced,nobeforeafter,width=7.1cm,height=2.1cm,valign=center,colback=csta!6,colframe=csta,arc=2.5mm,boxrule=0.9pt,left=3mm]{\Large\bfseries\color{csta}a)}\quad\begin{minipage}[c]{5.3cm}\raggedright\Large $+12\,\%$\end{minipage}\end{tcolorbox}\hfill\begin{tcolorbox}[enhanced,nobeforeafter,width=7.1cm,height=2.1cm,valign=center,colback=csta!6,colframe=csta,arc=2.5mm,boxrule=0.9pt,left=3mm]{\Large\bfseries\color{csta}b)}\quad\begin{minipage}[c]{5.3cm}\raggedright\Large $+6\,\%$\end{minipage}\end{tcolorbox}\par\vspace{2.5mm}
\begin{tcolorbox}[enhanced,nobeforeafter,width=7.1cm,height=2.1cm,valign=center,colback=csta!6,colframe=csta,arc=2.5mm,boxrule=0.9pt,left=3mm]{\Large\bfseries\color{csta}c)}\quad\begin{minipage}[c]{5.3cm}\raggedright\Large $+30\,\%$\end{minipage}\end{tcolorbox}\hfill\begin{tcolorbox}[enhanced,nobeforeafter,width=7.1cm,height=2.1cm,valign=center,colback=csta!6,colframe=csta,arc=2.5mm,boxrule=0.9pt,left=3mm]{\Large\bfseries\color{csta}d)}\quad\begin{minipage}[c]{5.3cm}\raggedright\Large $-12\,\%$\end{minipage}\end{tcolorbox}
\end{tcolorbox}
\newpage
\begin{tcolorbox}[enhanced,width=16cm,height=14cm,colback=white,colframe=csta,boxrule=2pt,arc=0mm,
left=6mm,right=6mm,top=0mm,bottom=4mm,sharp corners,
overlay={\fill[csta] (frame.north west) rectangle ([yshift=-1.25cm]frame.north east);
\node[anchor=west,text=white,font=\sffamily\bfseries\LARGE] at ([xshift=5mm,yshift=-0.62cm]frame.north west) {Question 23 C};
\node[anchor=east,text=white,font=\sffamily\itshape\large] at ([xshift=-5mm,yshift=-0.62cm]frame.north east) {Évolutions et statistiques};}]
\vspace*{1.75cm}
\begin{minipage}[t][6.9cm][c]{\linewidth}\LARGE\centering
\begin{minipage}{0.95\linewidth}\raggedright\linespread{1.25}\selectfont Un prix subit une baisse de $20\,\%$, puis une baisse de $20\,\%$. Quelle est son évolution globale~?\end{minipage}
\end{minipage}

\vfill
\begin{tcolorbox}[enhanced,nobeforeafter,width=7.1cm,height=2.1cm,valign=center,colback=csta!6,colframe=csta,arc=2.5mm,boxrule=0.9pt,left=3mm]{\Large\bfseries\color{csta}a)}\quad\begin{minipage}[c]{5.3cm}\raggedright\Large $+4\,\%$\end{minipage}\end{tcolorbox}\hfill\begin{tcolorbox}[enhanced,nobeforeafter,width=7.1cm,height=2.1cm,valign=center,colback=csta!6,colframe=csta,arc=2.5mm,boxrule=0.9pt,left=3mm]{\Large\bfseries\color{csta}b)}\quad\begin{minipage}[c]{5.3cm}\raggedright\Large $-40\,\%$\end{minipage}\end{tcolorbox}\par\vspace{2.5mm}
\begin{tcolorbox}[enhanced,nobeforeafter,width=7.1cm,height=2.1cm,valign=center,colback=csta!6,colframe=csta,arc=2.5mm,boxrule=0.9pt,left=3mm]{\Large\bfseries\color{csta}c)}\quad\begin{minipage}[c]{5.3cm}\raggedright\Large $-36\,\%$\end{minipage}\end{tcolorbox}\hfill\begin{tcolorbox}[enhanced,nobeforeafter,width=7.1cm,height=2.1cm,valign=center,colback=csta!6,colframe=csta,arc=2.5mm,boxrule=0.9pt,left=3mm]{\Large\bfseries\color{csta}d)}\quad\begin{minipage}[c]{5.3cm}\raggedright\Large $-35\,\%$\end{minipage}\end{tcolorbox}
\end{tcolorbox}
\newpage
\begin{tcolorbox}[enhanced,width=16cm,height=14cm,colback=white,colframe=csta,boxrule=2pt,arc=0mm,
left=6mm,right=6mm,top=0mm,bottom=4mm,sharp corners,
overlay={\fill[csta] (frame.north west) rectangle ([yshift=-1.25cm]frame.north east);
\node[anchor=west,text=white,font=\sffamily\bfseries\LARGE] at ([xshift=5mm,yshift=-0.62cm]frame.north west) {Question 24 C};
\node[anchor=east,text=white,font=\sffamily\itshape\large] at ([xshift=-5mm,yshift=-0.62cm]frame.north east) {Évolutions et statistiques};}]
\vspace*{1.75cm}
\begin{minipage}[t][6.9cm][c]{\linewidth}\LARGE\centering
\begin{minipage}{0.95\linewidth}\raggedright\linespread{1.25}\selectfont Un prix subit une hausse de $25\,\%$, puis une baisse de $20\,\%$. Quelle est son évolution globale~?\end{minipage}
\end{minipage}

\vfill
\begin{tcolorbox}[enhanced,nobeforeafter,width=7.1cm,height=2.1cm,valign=center,colback=csta!6,colframe=csta,arc=2.5mm,boxrule=0.9pt,left=3mm]{\Large\bfseries\color{csta}a)}\quad\begin{minipage}[c]{5.3cm}\raggedright\Large $-5\,\%$\end{minipage}\end{tcolorbox}\hfill\begin{tcolorbox}[enhanced,nobeforeafter,width=7.1cm,height=2.1cm,valign=center,colback=csta!6,colframe=csta,arc=2.5mm,boxrule=0.9pt,left=3mm]{\Large\bfseries\color{csta}b)}\quad\begin{minipage}[c]{5.3cm}\raggedright\Large $+1\,\%$\end{minipage}\end{tcolorbox}\par\vspace{2.5mm}
\begin{tcolorbox}[enhanced,nobeforeafter,width=7.1cm,height=2.1cm,valign=center,colback=csta!6,colframe=csta,arc=2.5mm,boxrule=0.9pt,left=3mm]{\Large\bfseries\color{csta}c)}\quad\begin{minipage}[c]{5.3cm}\raggedright\Large $0\,\%$\end{minipage}\end{tcolorbox}\hfill\begin{tcolorbox}[enhanced,nobeforeafter,width=7.1cm,height=2.1cm,valign=center,colback=csta!6,colframe=csta,arc=2.5mm,boxrule=0.9pt,left=3mm]{\Large\bfseries\color{csta}d)}\quad\begin{minipage}[c]{5.3cm}\raggedright\Large $+5\,\%$\end{minipage}\end{tcolorbox}
\end{tcolorbox}
\newpage
\begin{tcolorbox}[enhanced,width=16cm,height=14cm,colback=white,colframe=csta,boxrule=2pt,arc=0mm,
left=6mm,right=6mm,top=0mm,bottom=4mm,sharp corners,
overlay={\fill[csta] (frame.north west) rectangle ([yshift=-1.25cm]frame.north east);
\node[anchor=west,text=white,font=\sffamily\bfseries\LARGE] at ([xshift=5mm,yshift=-0.62cm]frame.north west) {Question 25 C};
\node[anchor=east,text=white,font=\sffamily\itshape\large] at ([xshift=-5mm,yshift=-0.62cm]frame.north east) {Évolutions et statistiques};}]
\vspace*{1.75cm}
\begin{minipage}[t][6.9cm][c]{\linewidth}\LARGE\centering
\begin{minipage}{0.95\linewidth}\raggedright\linespread{1.25}\selectfont Après une hausse de $100\,\%$, quelle évolution faut-il appliquer pour revenir au prix de départ~?\end{minipage}
\end{minipage}

\vfill
\begin{tcolorbox}[enhanced,nobeforeafter,width=7.1cm,height=2.1cm,valign=center,colback=csta!6,colframe=csta,arc=2.5mm,boxrule=0.9pt,left=3mm]{\Large\bfseries\color{csta}a)}\quad\begin{minipage}[c]{5.3cm}\raggedright\Large une baisse de $50\,\%$\end{minipage}\end{tcolorbox}\hfill\begin{tcolorbox}[enhanced,nobeforeafter,width=7.1cm,height=2.1cm,valign=center,colback=csta!6,colframe=csta,arc=2.5mm,boxrule=0.9pt,left=3mm]{\Large\bfseries\color{csta}b)}\quad\begin{minipage}[c]{5.3cm}\raggedright\Large une hausse de $50\,\%$\end{minipage}\end{tcolorbox}\par\vspace{2.5mm}
\begin{tcolorbox}[enhanced,nobeforeafter,width=7.1cm,height=2.1cm,valign=center,colback=csta!6,colframe=csta,arc=2.5mm,boxrule=0.9pt,left=3mm]{\Large\bfseries\color{csta}c)}\quad\begin{minipage}[c]{5.3cm}\raggedright\Large une baisse de $60\,\%$\end{minipage}\end{tcolorbox}\hfill\begin{tcolorbox}[enhanced,nobeforeafter,width=7.1cm,height=2.1cm,valign=center,colback=csta!6,colframe=csta,arc=2.5mm,boxrule=0.9pt,left=3mm]{\Large\bfseries\color{csta}d)}\quad\begin{minipage}[c]{5.3cm}\raggedright\Large une baisse de $100\,\%$\end{minipage}\end{tcolorbox}
\end{tcolorbox}
\newpage
\begin{tcolorbox}[enhanced,width=16cm,height=14cm,colback=white,colframe=csta,boxrule=2pt,arc=0mm,
left=6mm,right=6mm,top=0mm,bottom=4mm,sharp corners,
overlay={\fill[csta] (frame.north west) rectangle ([yshift=-1.25cm]frame.north east);
\node[anchor=west,text=white,font=\sffamily\bfseries\LARGE] at ([xshift=5mm,yshift=-0.62cm]frame.north west) {Question 26 C};
\node[anchor=east,text=white,font=\sffamily\itshape\large] at ([xshift=-5mm,yshift=-0.62cm]frame.north east) {Évolutions et statistiques};}]
\vspace*{1.75cm}
\begin{minipage}[t][6.9cm][c]{\linewidth}\LARGE\centering
\begin{minipage}{0.95\linewidth}\raggedright\linespread{1.25}\selectfont Après une baisse de $37{,}5\,\%$, quelle évolution faut-il appliquer pour revenir au prix de départ~?\end{minipage}
\end{minipage}

\vfill
\begin{tcolorbox}[enhanced,nobeforeafter,width=7.1cm,height=2.1cm,valign=center,colback=csta!6,colframe=csta,arc=2.5mm,boxrule=0.9pt,left=3mm]{\Large\bfseries\color{csta}a)}\quad\begin{minipage}[c]{5.3cm}\raggedright\Large une baisse de $60\,\%$\end{minipage}\end{tcolorbox}\hfill\begin{tcolorbox}[enhanced,nobeforeafter,width=7.1cm,height=2.1cm,valign=center,colback=csta!6,colframe=csta,arc=2.5mm,boxrule=0.9pt,left=3mm]{\Large\bfseries\color{csta}b)}\quad\begin{minipage}[c]{5.3cm}\raggedright\Large une hausse de $70\,\%$\end{minipage}\end{tcolorbox}\par\vspace{2.5mm}
\begin{tcolorbox}[enhanced,nobeforeafter,width=7.1cm,height=2.1cm,valign=center,colback=csta!6,colframe=csta,arc=2.5mm,boxrule=0.9pt,left=3mm]{\Large\bfseries\color{csta}c)}\quad\begin{minipage}[c]{5.3cm}\raggedright\large une hausse de $37{,}5\,\%$\end{minipage}\end{tcolorbox}\hfill\begin{tcolorbox}[enhanced,nobeforeafter,width=7.1cm,height=2.1cm,valign=center,colback=csta!6,colframe=csta,arc=2.5mm,boxrule=0.9pt,left=3mm]{\Large\bfseries\color{csta}d)}\quad\begin{minipage}[c]{5.3cm}\raggedright\Large une hausse de $60\,\%$\end{minipage}\end{tcolorbox}
\end{tcolorbox}
\newpage
\begin{tcolorbox}[enhanced,width=16cm,height=14cm,colback=white,colframe=csta,boxrule=2pt,arc=0mm,
left=6mm,right=6mm,top=0mm,bottom=4mm,sharp corners,
overlay={\fill[csta] (frame.north west) rectangle ([yshift=-1.25cm]frame.north east);
\node[anchor=west,text=white,font=\sffamily\bfseries\LARGE] at ([xshift=5mm,yshift=-0.62cm]frame.north west) {Question 27 C};
\node[anchor=east,text=white,font=\sffamily\itshape\large] at ([xshift=-5mm,yshift=-0.62cm]frame.north east) {Évolutions et statistiques};}]
\vspace*{1.75cm}
\begin{minipage}[t][6.9cm][c]{\linewidth}\Large\centering
\begin{minipage}{0.95\linewidth}\raggedright\linespread{1.25}\selectfont Dans un lycée, $25\,\%$ des élèves sont internes et $60\,\%$ des internes viennent d'une autre commune. Quelle proportion des élèves du lycée sont des internes venant d'une autre commune~?\end{minipage}
\end{minipage}

\vfill
\begin{tcolorbox}[enhanced,nobeforeafter,width=7.1cm,height=2.1cm,valign=center,colback=csta!6,colframe=csta,arc=2.5mm,boxrule=0.9pt,left=3mm]{\Large\bfseries\color{csta}a)}\quad\begin{minipage}[c]{5.3cm}\raggedright\Large $85\,\%$\end{minipage}\end{tcolorbox}\hfill\begin{tcolorbox}[enhanced,nobeforeafter,width=7.1cm,height=2.1cm,valign=center,colback=csta!6,colframe=csta,arc=2.5mm,boxrule=0.9pt,left=3mm]{\Large\bfseries\color{csta}b)}\quad\begin{minipage}[c]{5.3cm}\raggedright\Large $42{,}5\,\%$\end{minipage}\end{tcolorbox}\par\vspace{2.5mm}
\begin{tcolorbox}[enhanced,nobeforeafter,width=7.1cm,height=2.1cm,valign=center,colback=csta!6,colframe=csta,arc=2.5mm,boxrule=0.9pt,left=3mm]{\Large\bfseries\color{csta}c)}\quad\begin{minipage}[c]{5.3cm}\raggedright\Large $15\,\%$\end{minipage}\end{tcolorbox}\hfill\begin{tcolorbox}[enhanced,nobeforeafter,width=7.1cm,height=2.1cm,valign=center,colback=csta!6,colframe=csta,arc=2.5mm,boxrule=0.9pt,left=3mm]{\Large\bfseries\color{csta}d)}\quad\begin{minipage}[c]{5.3cm}\raggedright\Large $35\,\%$\end{minipage}\end{tcolorbox}
\end{tcolorbox}
\newpage
\begin{tcolorbox}[enhanced,width=16cm,height=14cm,colback=white,colframe=csta,boxrule=2pt,arc=0mm,
left=6mm,right=6mm,top=0mm,bottom=4mm,sharp corners,
overlay={\fill[csta] (frame.north west) rectangle ([yshift=-1.25cm]frame.north east);
\node[anchor=west,text=white,font=\sffamily\bfseries\LARGE] at ([xshift=5mm,yshift=-0.62cm]frame.north west) {Question 28 C};
\node[anchor=east,text=white,font=\sffamily\itshape\large] at ([xshift=-5mm,yshift=-0.62cm]frame.north east) {Évolutions et statistiques};}]
\vspace*{1.75cm}
\begin{minipage}[t][6.9cm][c]{\linewidth}\Large\centering
\begin{minipage}{0.95\linewidth}\raggedright\linespread{1.25}\selectfont Dans un lycée, $90\,\%$ des élèves sont internes et $10\,\%$ des internes viennent d'une autre commune. Quelle proportion des élèves du lycée sont des internes venant d'une autre commune~?\end{minipage}
\end{minipage}

\vfill
\begin{tcolorbox}[enhanced,nobeforeafter,width=7.1cm,height=2.1cm,valign=center,colback=csta!6,colframe=csta,arc=2.5mm,boxrule=0.9pt,left=3mm]{\Large\bfseries\color{csta}a)}\quad\begin{minipage}[c]{5.3cm}\raggedright\Large $50\,\%$\end{minipage}\end{tcolorbox}\hfill\begin{tcolorbox}[enhanced,nobeforeafter,width=7.1cm,height=2.1cm,valign=center,colback=csta!6,colframe=csta,arc=2.5mm,boxrule=0.9pt,left=3mm]{\Large\bfseries\color{csta}b)}\quad\begin{minipage}[c]{5.3cm}\raggedright\Large $100\,\%$\end{minipage}\end{tcolorbox}\par\vspace{2.5mm}
\begin{tcolorbox}[enhanced,nobeforeafter,width=7.1cm,height=2.1cm,valign=center,colback=csta!6,colframe=csta,arc=2.5mm,boxrule=0.9pt,left=3mm]{\Large\bfseries\color{csta}c)}\quad\begin{minipage}[c]{5.3cm}\raggedright\Large $9\,\%$\end{minipage}\end{tcolorbox}\hfill\begin{tcolorbox}[enhanced,nobeforeafter,width=7.1cm,height=2.1cm,valign=center,colback=csta!6,colframe=csta,arc=2.5mm,boxrule=0.9pt,left=3mm]{\Large\bfseries\color{csta}d)}\quad\begin{minipage}[c]{5.3cm}\raggedright\Large $80\,\%$\end{minipage}\end{tcolorbox}
\end{tcolorbox}
\newpage
\begin{tcolorbox}[enhanced,width=16cm,height=14cm,colback=white,colframe=csta,boxrule=2pt,arc=0mm,
left=6mm,right=6mm,top=0mm,bottom=4mm,sharp corners,
overlay={\fill[csta] (frame.north west) rectangle ([yshift=-1.25cm]frame.north east);
\node[anchor=west,text=white,font=\sffamily\bfseries\LARGE] at ([xshift=5mm,yshift=-0.62cm]frame.north west) {Question 29 C};
\node[anchor=east,text=white,font=\sffamily\itshape\large] at ([xshift=-5mm,yshift=-0.62cm]frame.north east) {Évolutions et statistiques};}]
\vspace*{1.75cm}
\begin{minipage}[t][6.9cm][c]{\linewidth}\Large\centering
\begin{minipage}{0.95\linewidth}\raggedright\linespread{1.25}\selectfont Quelle est la moyenne de cette série statistique~? $\begin{array}{|l|c|c|}\hline \text{Valeur} & 12 & 16\\\hline \text{Effectif} & 6 & 2\\\hline\end{array}$\end{minipage}
\end{minipage}

\vfill
\begin{tcolorbox}[enhanced,nobeforeafter,width=7.1cm,height=2.1cm,valign=center,colback=csta!6,colframe=csta,arc=2.5mm,boxrule=0.9pt,left=3mm]{\Large\bfseries\color{csta}a)}\quad\begin{minipage}[c]{5.3cm}\raggedright\Large $52$\end{minipage}\end{tcolorbox}\hfill\begin{tcolorbox}[enhanced,nobeforeafter,width=7.1cm,height=2.1cm,valign=center,colback=csta!6,colframe=csta,arc=2.5mm,boxrule=0.9pt,left=3mm]{\Large\bfseries\color{csta}b)}\quad\begin{minipage}[c]{5.3cm}\raggedright\Large $13$\end{minipage}\end{tcolorbox}\par\vspace{2.5mm}
\begin{tcolorbox}[enhanced,nobeforeafter,width=7.1cm,height=2.1cm,valign=center,colback=csta!6,colframe=csta,arc=2.5mm,boxrule=0.9pt,left=3mm]{\Large\bfseries\color{csta}c)}\quad\begin{minipage}[c]{5.3cm}\raggedright\Large $14$\end{minipage}\end{tcolorbox}\hfill\begin{tcolorbox}[enhanced,nobeforeafter,width=7.1cm,height=2.1cm,valign=center,colback=csta!6,colframe=csta,arc=2.5mm,boxrule=0.9pt,left=3mm]{\Large\bfseries\color{csta}d)}\quad\begin{minipage}[c]{5.3cm}\raggedright\Large $12{,}5$\end{minipage}\end{tcolorbox}
\end{tcolorbox}
\newpage
\begin{tcolorbox}[enhanced,width=16cm,height=14cm,colback=white,colframe=csta,boxrule=2pt,arc=0mm,
left=6mm,right=6mm,top=0mm,bottom=4mm,sharp corners,
overlay={\fill[csta] (frame.north west) rectangle ([yshift=-1.25cm]frame.north east);
\node[anchor=west,text=white,font=\sffamily\bfseries\LARGE] at ([xshift=5mm,yshift=-0.62cm]frame.north west) {Question 30 C};
\node[anchor=east,text=white,font=\sffamily\itshape\large] at ([xshift=-5mm,yshift=-0.62cm]frame.north east) {Évolutions et statistiques};}]
\vspace*{1.75cm}
\begin{minipage}[t][6.9cm][c]{\linewidth}\Large\centering
\begin{minipage}{0.95\linewidth}\raggedright\linespread{1.25}\selectfont Quelle est la moyenne de cette série statistique~? $\begin{array}{|l|c|c|c|c|}\hline \text{Valeur} & 1 & 2 & 3 & 4\\\hline \text{Effectif} & 3 & 4 & 2 & 1\\\hline\end{array}$\end{minipage}
\end{minipage}

\vfill
\begin{tcolorbox}[enhanced,nobeforeafter,width=7.1cm,height=2.1cm,valign=center,colback=csta!6,colframe=csta,arc=2.5mm,boxrule=0.9pt,left=3mm]{\Large\bfseries\color{csta}a)}\quad\begin{minipage}[c]{5.3cm}\raggedright\Large $2{,}5$\end{minipage}\end{tcolorbox}\hfill\begin{tcolorbox}[enhanced,nobeforeafter,width=7.1cm,height=2.1cm,valign=center,colback=csta!6,colframe=csta,arc=2.5mm,boxrule=0.9pt,left=3mm]{\Large\bfseries\color{csta}b)}\quad\begin{minipage}[c]{5.3cm}\raggedright\Large $5{,}25$\end{minipage}\end{tcolorbox}\par\vspace{2.5mm}
\begin{tcolorbox}[enhanced,nobeforeafter,width=7.1cm,height=2.1cm,valign=center,colback=csta!6,colframe=csta,arc=2.5mm,boxrule=0.9pt,left=3mm]{\Large\bfseries\color{csta}c)}\quad\begin{minipage}[c]{5.3cm}\raggedright\Large $2{,}1$\end{minipage}\end{tcolorbox}\hfill\begin{tcolorbox}[enhanced,nobeforeafter,width=7.1cm,height=2.1cm,valign=center,colback=csta!6,colframe=csta,arc=2.5mm,boxrule=0.9pt,left=3mm]{\Large\bfseries\color{csta}d)}\quad\begin{minipage}[c]{5.3cm}\raggedright\Large $3{,}1$\end{minipage}\end{tcolorbox}
\end{tcolorbox}
\newpage
\begin{tcolorbox}[enhanced,width=16cm,height=14cm,colback=white,colframe=csta,boxrule=2pt,arc=0mm,
left=6mm,right=6mm,top=0mm,bottom=4mm,sharp corners,
overlay={\fill[csta] (frame.north west) rectangle ([yshift=-1.25cm]frame.north east);
\node[anchor=west,text=white,font=\sffamily\bfseries\LARGE] at ([xshift=5mm,yshift=-0.62cm]frame.north west) {Question 31 C};
\node[anchor=east,text=white,font=\sffamily\itshape\large] at ([xshift=-5mm,yshift=-0.62cm]frame.north east) {Évolutions et statistiques};}]
\vspace*{1.75cm}
\begin{minipage}[t][6.9cm][c]{\linewidth}\Large\centering
\begin{minipage}{0.95\linewidth}\raggedright\linespread{1.25}\selectfont Voici une série rangée dans l'ordre croissant~: $7\ ;\ 9\ ;\ 10\ ;\ 10\ ;\ 11\ ;\ 14\ ;\ 15\ ;\ 16\ ;\ 17\ ;\ 19\ ;\ 20\ ;\ 24$. Quel est son premier quartile $Q_1$~?\end{minipage}
\end{minipage}

\vfill
\begin{tcolorbox}[enhanced,nobeforeafter,width=7.1cm,height=2.1cm,valign=center,colback=csta!6,colframe=csta,arc=2.5mm,boxrule=0.9pt,left=3mm]{\Large\bfseries\color{csta}a)}\quad\begin{minipage}[c]{5.3cm}\raggedright\Large $14{,}5$\end{minipage}\end{tcolorbox}\hfill\begin{tcolorbox}[enhanced,nobeforeafter,width=7.1cm,height=2.1cm,valign=center,colback=csta!6,colframe=csta,arc=2.5mm,boxrule=0.9pt,left=3mm]{\Large\bfseries\color{csta}b)}\quad\begin{minipage}[c]{5.3cm}\raggedright\Large $7$\end{minipage}\end{tcolorbox}\par\vspace{2.5mm}
\begin{tcolorbox}[enhanced,nobeforeafter,width=7.1cm,height=2.1cm,valign=center,colback=csta!6,colframe=csta,arc=2.5mm,boxrule=0.9pt,left=3mm]{\Large\bfseries\color{csta}c)}\quad\begin{minipage}[c]{5.3cm}\raggedright\Large $10$\end{minipage}\end{tcolorbox}\hfill\begin{tcolorbox}[enhanced,nobeforeafter,width=7.1cm,height=2.1cm,valign=center,colback=csta!6,colframe=csta,arc=2.5mm,boxrule=0.9pt,left=3mm]{\Large\bfseries\color{csta}d)}\quad\begin{minipage}[c]{5.3cm}\raggedright\Large $17$\end{minipage}\end{tcolorbox}
\end{tcolorbox}
\newpage
\begin{tcolorbox}[enhanced,width=16cm,height=14cm,colback=white,colframe=csta,boxrule=2pt,arc=0mm,
left=6mm,right=6mm,top=0mm,bottom=4mm,sharp corners,
overlay={\fill[csta] (frame.north west) rectangle ([yshift=-1.25cm]frame.north east);
\node[anchor=west,text=white,font=\sffamily\bfseries\LARGE] at ([xshift=5mm,yshift=-0.62cm]frame.north west) {Question 32 C};
\node[anchor=east,text=white,font=\sffamily\itshape\large] at ([xshift=-5mm,yshift=-0.62cm]frame.north east) {Évolutions et statistiques};}]
\vspace*{1.75cm}
\begin{minipage}[t][6.9cm][c]{\linewidth}\LARGE\centering
\begin{minipage}{0.95\linewidth}\raggedright\linespread{1.25}\selectfont Voici une série rangée dans l'ordre croissant~: $5\ ;\ 6\ ;\ 8\ ;\ 9\ ;\ 12\ ;\ 14\ ;\ 15$. Quel est son premier quartile $Q_1$~?\end{minipage}
\end{minipage}

\vfill
\begin{tcolorbox}[enhanced,nobeforeafter,width=7.1cm,height=2.1cm,valign=center,colback=csta!6,colframe=csta,arc=2.5mm,boxrule=0.9pt,left=3mm]{\Large\bfseries\color{csta}a)}\quad\begin{minipage}[c]{5.3cm}\raggedright\Large $14$\end{minipage}\end{tcolorbox}\hfill\begin{tcolorbox}[enhanced,nobeforeafter,width=7.1cm,height=2.1cm,valign=center,colback=csta!6,colframe=csta,arc=2.5mm,boxrule=0.9pt,left=3mm]{\Large\bfseries\color{csta}b)}\quad\begin{minipage}[c]{5.3cm}\raggedright\Large $9$\end{minipage}\end{tcolorbox}\par\vspace{2.5mm}
\begin{tcolorbox}[enhanced,nobeforeafter,width=7.1cm,height=2.1cm,valign=center,colback=csta!6,colframe=csta,arc=2.5mm,boxrule=0.9pt,left=3mm]{\Large\bfseries\color{csta}c)}\quad\begin{minipage}[c]{5.3cm}\raggedright\Large $8$\end{minipage}\end{tcolorbox}\hfill\begin{tcolorbox}[enhanced,nobeforeafter,width=7.1cm,height=2.1cm,valign=center,colback=csta!6,colframe=csta,arc=2.5mm,boxrule=0.9pt,left=3mm]{\Large\bfseries\color{csta}d)}\quad\begin{minipage}[c]{5.3cm}\raggedright\Large $6$\end{minipage}\end{tcolorbox}
\end{tcolorbox}
\newpage
\begin{tcolorbox}[enhanced,width=16cm,height=14cm,colback=white,colframe=cfon,boxrule=2pt,arc=0mm,
left=6mm,right=6mm,top=0mm,bottom=4mm,sharp corners,
overlay={\fill[cfon] (frame.north west) rectangle ([yshift=-1.25cm]frame.north east);
\node[anchor=west,text=white,font=\sffamily\bfseries\LARGE] at ([xshift=5mm,yshift=-0.62cm]frame.north west) {Question 33 C};
\node[anchor=east,text=white,font=\sffamily\itshape\large] at ([xshift=-5mm,yshift=-0.62cm]frame.north east) {Fonctions};}]
\vspace*{1.75cm}
\begin{minipage}[t][6.9cm][c]{\linewidth}\LARGE\centering
\begin{minipage}{0.95\linewidth}\raggedright\linespread{1.25}\selectfont Soit $f(x)=-2x^2-x+4$. Que vaut $f(1)$~?\end{minipage}
\end{minipage}

\vfill
\begin{tcolorbox}[enhanced,nobeforeafter,width=7.1cm,height=2.1cm,valign=center,colback=cfon!6,colframe=cfon,arc=2.5mm,boxrule=0.9pt,left=3mm]{\Large\bfseries\color{cfon}a)}\quad\begin{minipage}[c]{5.3cm}\raggedright\Large $-1$\end{minipage}\end{tcolorbox}\hfill\begin{tcolorbox}[enhanced,nobeforeafter,width=7.1cm,height=2.1cm,valign=center,colback=cfon!6,colframe=cfon,arc=2.5mm,boxrule=0.9pt,left=3mm]{\Large\bfseries\color{cfon}b)}\quad\begin{minipage}[c]{5.3cm}\raggedright\Large $3$\end{minipage}\end{tcolorbox}\par\vspace{2.5mm}
\begin{tcolorbox}[enhanced,nobeforeafter,width=7.1cm,height=2.1cm,valign=center,colback=cfon!6,colframe=cfon,arc=2.5mm,boxrule=0.9pt,left=3mm]{\Large\bfseries\color{cfon}c)}\quad\begin{minipage}[c]{5.3cm}\raggedright\Large $1$\end{minipage}\end{tcolorbox}\hfill\begin{tcolorbox}[enhanced,nobeforeafter,width=7.1cm,height=2.1cm,valign=center,colback=cfon!6,colframe=cfon,arc=2.5mm,boxrule=0.9pt,left=3mm]{\Large\bfseries\color{cfon}d)}\quad\begin{minipage}[c]{5.3cm}\raggedright\Large $5$\end{minipage}\end{tcolorbox}
\end{tcolorbox}
\newpage
\begin{tcolorbox}[enhanced,width=16cm,height=14cm,colback=white,colframe=cfon,boxrule=2pt,arc=0mm,
left=6mm,right=6mm,top=0mm,bottom=4mm,sharp corners,
overlay={\fill[cfon] (frame.north west) rectangle ([yshift=-1.25cm]frame.north east);
\node[anchor=west,text=white,font=\sffamily\bfseries\LARGE] at ([xshift=5mm,yshift=-0.62cm]frame.north west) {Question 34 C};
\node[anchor=east,text=white,font=\sffamily\itshape\large] at ([xshift=-5mm,yshift=-0.62cm]frame.north east) {Fonctions};}]
\vspace*{1.75cm}
\begin{minipage}[t][6.9cm][c]{\linewidth}\LARGE\centering
\begin{minipage}{0.95\linewidth}\raggedright\linespread{1.25}\selectfont Soit $f(x)=x^2-6x+9$. Que vaut $f(5)$~?\end{minipage}
\end{minipage}

\vfill
\begin{tcolorbox}[enhanced,nobeforeafter,width=7.1cm,height=2.1cm,valign=center,colback=cfon!6,colframe=cfon,arc=2.5mm,boxrule=0.9pt,left=3mm]{\Large\bfseries\color{cfon}a)}\quad\begin{minipage}[c]{5.3cm}\raggedright\Large $4$\end{minipage}\end{tcolorbox}\hfill\begin{tcolorbox}[enhanced,nobeforeafter,width=7.1cm,height=2.1cm,valign=center,colback=cfon!6,colframe=cfon,arc=2.5mm,boxrule=0.9pt,left=3mm]{\Large\bfseries\color{cfon}b)}\quad\begin{minipage}[c]{5.3cm}\raggedright\Large $-46$\end{minipage}\end{tcolorbox}\par\vspace{2.5mm}
\begin{tcolorbox}[enhanced,nobeforeafter,width=7.1cm,height=2.1cm,valign=center,colback=cfon!6,colframe=cfon,arc=2.5mm,boxrule=0.9pt,left=3mm]{\Large\bfseries\color{cfon}c)}\quad\begin{minipage}[c]{5.3cm}\raggedright\Large $-4$\end{minipage}\end{tcolorbox}\hfill\begin{tcolorbox}[enhanced,nobeforeafter,width=7.1cm,height=2.1cm,valign=center,colback=cfon!6,colframe=cfon,arc=2.5mm,boxrule=0.9pt,left=3mm]{\Large\bfseries\color{cfon}d)}\quad\begin{minipage}[c]{5.3cm}\raggedright\Large $-11$\end{minipage}\end{tcolorbox}
\end{tcolorbox}
\newpage
\begin{tcolorbox}[enhanced,width=16cm,height=14cm,colback=white,colframe=cfon,boxrule=2pt,arc=0mm,
left=6mm,right=6mm,top=0mm,bottom=4mm,sharp corners,
overlay={\fill[cfon] (frame.north west) rectangle ([yshift=-1.25cm]frame.north east);
\node[anchor=west,text=white,font=\sffamily\bfseries\LARGE] at ([xshift=5mm,yshift=-0.62cm]frame.north west) {Question 35 C};
\node[anchor=east,text=white,font=\sffamily\itshape\large] at ([xshift=-5mm,yshift=-0.62cm]frame.north east) {Fonctions};}]
\vspace*{1.75cm}
\begin{minipage}[t][6.9cm][c]{\linewidth}\LARGE\centering
\begin{minipage}{0.95\linewidth}\raggedright\linespread{1.25}\selectfont Soit $f(x)=-2x+9$. Quel est l'antécédent de $15$ par $f$~?\end{minipage}
\end{minipage}

\vfill
\begin{tcolorbox}[enhanced,nobeforeafter,width=7.1cm,height=2.1cm,valign=center,colback=cfon!6,colframe=cfon,arc=2.5mm,boxrule=0.9pt,left=3mm]{\Large\bfseries\color{cfon}a)}\quad\begin{minipage}[c]{5.3cm}\raggedright\Large $-2$\end{minipage}\end{tcolorbox}\hfill\begin{tcolorbox}[enhanced,nobeforeafter,width=7.1cm,height=2.1cm,valign=center,colback=cfon!6,colframe=cfon,arc=2.5mm,boxrule=0.9pt,left=3mm]{\Large\bfseries\color{cfon}b)}\quad\begin{minipage}[c]{5.3cm}\raggedright\Large $-21$\end{minipage}\end{tcolorbox}\par\vspace{2.5mm}
\begin{tcolorbox}[enhanced,nobeforeafter,width=7.1cm,height=2.1cm,valign=center,colback=cfon!6,colframe=cfon,arc=2.5mm,boxrule=0.9pt,left=3mm]{\Large\bfseries\color{cfon}c)}\quad\begin{minipage}[c]{5.3cm}\raggedright\Large $-12$\end{minipage}\end{tcolorbox}\hfill\begin{tcolorbox}[enhanced,nobeforeafter,width=7.1cm,height=2.1cm,valign=center,colback=cfon!6,colframe=cfon,arc=2.5mm,boxrule=0.9pt,left=3mm]{\Large\bfseries\color{cfon}d)}\quad\begin{minipage}[c]{5.3cm}\raggedright\Large $-3$\end{minipage}\end{tcolorbox}
\end{tcolorbox}
\newpage
\begin{tcolorbox}[enhanced,width=16cm,height=14cm,colback=white,colframe=cfon,boxrule=2pt,arc=0mm,
left=6mm,right=6mm,top=0mm,bottom=4mm,sharp corners,
overlay={\fill[cfon] (frame.north west) rectangle ([yshift=-1.25cm]frame.north east);
\node[anchor=west,text=white,font=\sffamily\bfseries\LARGE] at ([xshift=5mm,yshift=-0.62cm]frame.north west) {Question 36 C};
\node[anchor=east,text=white,font=\sffamily\itshape\large] at ([xshift=-5mm,yshift=-0.62cm]frame.north east) {Fonctions};}]
\vspace*{1.75cm}
\begin{minipage}[t][6.9cm][c]{\linewidth}\LARGE\centering
\begin{minipage}{0.95\linewidth}\raggedright\linespread{1.25}\selectfont Soit $f(x)=3x-4$. Quel est l'antécédent de $-4$ par $f$~?\end{minipage}
\end{minipage}

\vfill
\begin{tcolorbox}[enhanced,nobeforeafter,width=7.1cm,height=2.1cm,valign=center,colback=cfon!6,colframe=cfon,arc=2.5mm,boxrule=0.9pt,left=3mm]{\Large\bfseries\color{cfon}a)}\quad\begin{minipage}[c]{5.3cm}\raggedright\Large $0$\end{minipage}\end{tcolorbox}\hfill\begin{tcolorbox}[enhanced,nobeforeafter,width=7.1cm,height=2.1cm,valign=center,colback=cfon!6,colframe=cfon,arc=2.5mm,boxrule=0.9pt,left=3mm]{\Large\bfseries\color{cfon}b)}\quad\begin{minipage}[c]{5.3cm}\raggedright\Large $1$\end{minipage}\end{tcolorbox}\par\vspace{2.5mm}
\begin{tcolorbox}[enhanced,nobeforeafter,width=7.1cm,height=2.1cm,valign=center,colback=cfon!6,colframe=cfon,arc=2.5mm,boxrule=0.9pt,left=3mm]{\Large\bfseries\color{cfon}c)}\quad\begin{minipage}[c]{5.3cm}\raggedright\Large $-\dfrac{8}{3}$\end{minipage}\end{tcolorbox}\hfill\begin{tcolorbox}[enhanced,nobeforeafter,width=7.1cm,height=2.1cm,valign=center,colback=cfon!6,colframe=cfon,arc=2.5mm,boxrule=0.9pt,left=3mm]{\Large\bfseries\color{cfon}d)}\quad\begin{minipage}[c]{5.3cm}\raggedright\Large $-16$\end{minipage}\end{tcolorbox}
\end{tcolorbox}
\newpage
\begin{tcolorbox}[enhanced,width=16cm,height=14cm,colback=white,colframe=cfon,boxrule=2pt,arc=0mm,
left=6mm,right=6mm,top=0mm,bottom=4mm,sharp corners,
overlay={\fill[cfon] (frame.north west) rectangle ([yshift=-1.25cm]frame.north east);
\node[anchor=west,text=white,font=\sffamily\bfseries\LARGE] at ([xshift=5mm,yshift=-0.62cm]frame.north west) {Question 37 C};
\node[anchor=east,text=white,font=\sffamily\itshape\large] at ([xshift=-5mm,yshift=-0.62cm]frame.north east) {Fonctions};}]
\vspace*{1.75cm}
\begin{minipage}[t][6.9cm][c]{\linewidth}\LARGE\centering
\begin{minipage}{0.95\linewidth}\raggedright\linespread{1.25}\selectfont Quel est le coefficient directeur de la droite $(AB)$, avec $A(-3\,;\,-1)$ et $B(1\,;\,7)$~?\end{minipage}
\end{minipage}

\vfill
\begin{tcolorbox}[enhanced,nobeforeafter,width=7.1cm,height=2.1cm,valign=center,colback=cfon!6,colframe=cfon,arc=2.5mm,boxrule=0.9pt,left=3mm]{\Large\bfseries\color{cfon}a)}\quad\begin{minipage}[c]{5.3cm}\raggedright\Large $2$\end{minipage}\end{tcolorbox}\hfill\begin{tcolorbox}[enhanced,nobeforeafter,width=7.1cm,height=2.1cm,valign=center,colback=cfon!6,colframe=cfon,arc=2.5mm,boxrule=0.9pt,left=3mm]{\Large\bfseries\color{cfon}b)}\quad\begin{minipage}[c]{5.3cm}\raggedright\Large $-2$\end{minipage}\end{tcolorbox}\par\vspace{2.5mm}
\begin{tcolorbox}[enhanced,nobeforeafter,width=7.1cm,height=2.1cm,valign=center,colback=cfon!6,colframe=cfon,arc=2.5mm,boxrule=0.9pt,left=3mm]{\Large\bfseries\color{cfon}c)}\quad\begin{minipage}[c]{5.3cm}\raggedright\Large $\dfrac{1}{2}$\end{minipage}\end{tcolorbox}\hfill\begin{tcolorbox}[enhanced,nobeforeafter,width=7.1cm,height=2.1cm,valign=center,colback=cfon!6,colframe=cfon,arc=2.5mm,boxrule=0.9pt,left=3mm]{\Large\bfseries\color{cfon}d)}\quad\begin{minipage}[c]{5.3cm}\raggedright\Large $-3$\end{minipage}\end{tcolorbox}
\end{tcolorbox}
\newpage
\begin{tcolorbox}[enhanced,width=16cm,height=14cm,colback=white,colframe=cfon,boxrule=2pt,arc=0mm,
left=6mm,right=6mm,top=0mm,bottom=4mm,sharp corners,
overlay={\fill[cfon] (frame.north west) rectangle ([yshift=-1.25cm]frame.north east);
\node[anchor=west,text=white,font=\sffamily\bfseries\LARGE] at ([xshift=5mm,yshift=-0.62cm]frame.north west) {Question 38 C};
\node[anchor=east,text=white,font=\sffamily\itshape\large] at ([xshift=-5mm,yshift=-0.62cm]frame.north east) {Fonctions};}]
\vspace*{1.75cm}
\begin{minipage}[t][6.9cm][c]{\linewidth}\LARGE\centering
\begin{minipage}{0.95\linewidth}\raggedright\linespread{1.25}\selectfont Quel est le coefficient directeur de la droite $(AB)$, avec $A(2\,;\,6)$ et $B(8\,;\,-3)$~?\end{minipage}
\end{minipage}

\vfill
\begin{tcolorbox}[enhanced,nobeforeafter,width=7.1cm,height=2.1cm,valign=center,colback=cfon!6,colframe=cfon,arc=2.5mm,boxrule=0.9pt,left=3mm]{\Large\bfseries\color{cfon}a)}\quad\begin{minipage}[c]{5.3cm}\raggedright\Large $\dfrac{3}{2}$\end{minipage}\end{tcolorbox}\hfill\begin{tcolorbox}[enhanced,nobeforeafter,width=7.1cm,height=2.1cm,valign=center,colback=cfon!6,colframe=cfon,arc=2.5mm,boxrule=0.9pt,left=3mm]{\Large\bfseries\color{cfon}b)}\quad\begin{minipage}[c]{5.3cm}\raggedright\Large $\dfrac{3}{10}$\end{minipage}\end{tcolorbox}\par\vspace{2.5mm}
\begin{tcolorbox}[enhanced,nobeforeafter,width=7.1cm,height=2.1cm,valign=center,colback=cfon!6,colframe=cfon,arc=2.5mm,boxrule=0.9pt,left=3mm]{\Large\bfseries\color{cfon}c)}\quad\begin{minipage}[c]{5.3cm}\raggedright\Large $-\dfrac{2}{3}$\end{minipage}\end{tcolorbox}\hfill\begin{tcolorbox}[enhanced,nobeforeafter,width=7.1cm,height=2.1cm,valign=center,colback=cfon!6,colframe=cfon,arc=2.5mm,boxrule=0.9pt,left=3mm]{\Large\bfseries\color{cfon}d)}\quad\begin{minipage}[c]{5.3cm}\raggedright\Large $-\dfrac{3}{2}$\end{minipage}\end{tcolorbox}
\end{tcolorbox}
\newpage
\begin{tcolorbox}[enhanced,width=16cm,height=14cm,colback=white,colframe=cfon,boxrule=2pt,arc=0mm,
left=6mm,right=6mm,top=0mm,bottom=4mm,sharp corners,
overlay={\fill[cfon] (frame.north west) rectangle ([yshift=-1.25cm]frame.north east);
\node[anchor=west,text=white,font=\sffamily\bfseries\LARGE] at ([xshift=5mm,yshift=-0.62cm]frame.north west) {Question 39 C};
\node[anchor=east,text=white,font=\sffamily\itshape\large] at ([xshift=-5mm,yshift=-0.62cm]frame.north east) {Fonctions};}]
\vspace*{1.75cm}
\begin{minipage}[t][6.9cm][c]{\linewidth}\LARGE\centering
\begin{minipage}{0.95\linewidth}\raggedright\linespread{1.25}\selectfont Quelle est l'équation réduite de la droite de coefficient directeur $-1$ passant par $A(3\,;\,3)$~?\end{minipage}
\end{minipage}

\vfill
\begin{tcolorbox}[enhanced,nobeforeafter,width=7.1cm,height=2.1cm,valign=center,colback=cfon!6,colframe=cfon,arc=2.5mm,boxrule=0.9pt,left=3mm]{\Large\bfseries\color{cfon}a)}\quad\begin{minipage}[c]{5.3cm}\raggedright\Large $y=-x$\end{minipage}\end{tcolorbox}\hfill\begin{tcolorbox}[enhanced,nobeforeafter,width=7.1cm,height=2.1cm,valign=center,colback=cfon!6,colframe=cfon,arc=2.5mm,boxrule=0.9pt,left=3mm]{\Large\bfseries\color{cfon}b)}\quad\begin{minipage}[c]{5.3cm}\raggedright\Large $y=3x+6$\end{minipage}\end{tcolorbox}\par\vspace{2.5mm}
\begin{tcolorbox}[enhanced,nobeforeafter,width=7.1cm,height=2.1cm,valign=center,colback=cfon!6,colframe=cfon,arc=2.5mm,boxrule=0.9pt,left=3mm]{\Large\bfseries\color{cfon}c)}\quad\begin{minipage}[c]{5.3cm}\raggedright\Large $y=-x+3$\end{minipage}\end{tcolorbox}\hfill\begin{tcolorbox}[enhanced,nobeforeafter,width=7.1cm,height=2.1cm,valign=center,colback=cfon!6,colframe=cfon,arc=2.5mm,boxrule=0.9pt,left=3mm]{\Large\bfseries\color{cfon}d)}\quad\begin{minipage}[c]{5.3cm}\raggedright\Large $y=-x+6$\end{minipage}\end{tcolorbox}
\end{tcolorbox}
\newpage
\begin{tcolorbox}[enhanced,width=16cm,height=14cm,colback=white,colframe=cfon,boxrule=2pt,arc=0mm,
left=6mm,right=6mm,top=0mm,bottom=4mm,sharp corners,
overlay={\fill[cfon] (frame.north west) rectangle ([yshift=-1.25cm]frame.north east);
\node[anchor=west,text=white,font=\sffamily\bfseries\LARGE] at ([xshift=5mm,yshift=-0.62cm]frame.north west) {Question 40 C};
\node[anchor=east,text=white,font=\sffamily\itshape\large] at ([xshift=-5mm,yshift=-0.62cm]frame.north east) {Fonctions};}]
\vspace*{1.75cm}
\begin{minipage}[t][6.9cm][c]{\linewidth}\LARGE\centering
\begin{minipage}{0.95\linewidth}\raggedright\linespread{1.25}\selectfont Quelle est l'équation réduite de la droite de coefficient directeur $5$ passant par $A(0\,;\,-4)$~?\end{minipage}
\end{minipage}

\vfill
\begin{tcolorbox}[enhanced,nobeforeafter,width=7.1cm,height=2.1cm,valign=center,colback=cfon!6,colframe=cfon,arc=2.5mm,boxrule=0.9pt,left=3mm]{\Large\bfseries\color{cfon}a)}\quad\begin{minipage}[c]{5.3cm}\raggedright\Large $y=5x-4$\end{minipage}\end{tcolorbox}\hfill\begin{tcolorbox}[enhanced,nobeforeafter,width=7.1cm,height=2.1cm,valign=center,colback=cfon!6,colframe=cfon,arc=2.5mm,boxrule=0.9pt,left=3mm]{\Large\bfseries\color{cfon}b)}\quad\begin{minipage}[c]{5.3cm}\raggedright\Large $y=-5x-4$\end{minipage}\end{tcolorbox}\par\vspace{2.5mm}
\begin{tcolorbox}[enhanced,nobeforeafter,width=7.1cm,height=2.1cm,valign=center,colback=cfon!6,colframe=cfon,arc=2.5mm,boxrule=0.9pt,left=3mm]{\Large\bfseries\color{cfon}c)}\quad\begin{minipage}[c]{5.3cm}\raggedright\Large $y=5x+4$\end{minipage}\end{tcolorbox}\hfill\begin{tcolorbox}[enhanced,nobeforeafter,width=7.1cm,height=2.1cm,valign=center,colback=cfon!6,colframe=cfon,arc=2.5mm,boxrule=0.9pt,left=3mm]{\Large\bfseries\color{cfon}d)}\quad\begin{minipage}[c]{5.3cm}\raggedright\Large $y=5x-3$\end{minipage}\end{tcolorbox}
\end{tcolorbox}
\newpage
\begin{tcolorbox}[enhanced,width=16cm,height=14cm,colback=white,colframe=cfon,boxrule=2pt,arc=0mm,
left=6mm,right=6mm,top=0mm,bottom=4mm,sharp corners,
overlay={\fill[cfon] (frame.north west) rectangle ([yshift=-1.25cm]frame.north east);
\node[anchor=west,text=white,font=\sffamily\bfseries\LARGE] at ([xshift=5mm,yshift=-0.62cm]frame.north west) {Question 41 C};
\node[anchor=east,text=white,font=\sffamily\itshape\large] at ([xshift=-5mm,yshift=-0.62cm]frame.north east) {Fonctions};}]
\vspace*{1.75cm}
\begin{minipage}[t][6.9cm][c]{\linewidth}\LARGE\centering
\begin{minipage}{0.95\linewidth}\raggedright\linespread{1.25}\selectfont Pour quels réels $x$ a-t-on $x-7>0$~?\end{minipage}
\end{minipage}

\vfill
\begin{tcolorbox}[enhanced,nobeforeafter,width=7.1cm,height=2.1cm,valign=center,colback=cfon!6,colframe=cfon,arc=2.5mm,boxrule=0.9pt,left=3mm]{\Large\bfseries\color{cfon}a)}\quad\begin{minipage}[c]{5.3cm}\raggedright\Large $x<-7$\end{minipage}\end{tcolorbox}\hfill\begin{tcolorbox}[enhanced,nobeforeafter,width=7.1cm,height=2.1cm,valign=center,colback=cfon!6,colframe=cfon,arc=2.5mm,boxrule=0.9pt,left=3mm]{\Large\bfseries\color{cfon}b)}\quad\begin{minipage}[c]{5.3cm}\raggedright\Large $x>-7$\end{minipage}\end{tcolorbox}\par\vspace{2.5mm}
\begin{tcolorbox}[enhanced,nobeforeafter,width=7.1cm,height=2.1cm,valign=center,colback=cfon!6,colframe=cfon,arc=2.5mm,boxrule=0.9pt,left=3mm]{\Large\bfseries\color{cfon}c)}\quad\begin{minipage}[c]{5.3cm}\raggedright\Large $x>7$\end{minipage}\end{tcolorbox}\hfill\begin{tcolorbox}[enhanced,nobeforeafter,width=7.1cm,height=2.1cm,valign=center,colback=cfon!6,colframe=cfon,arc=2.5mm,boxrule=0.9pt,left=3mm]{\Large\bfseries\color{cfon}d)}\quad\begin{minipage}[c]{5.3cm}\raggedright\Large $x<7$\end{minipage}\end{tcolorbox}
\end{tcolorbox}
\newpage
\begin{tcolorbox}[enhanced,width=16cm,height=14cm,colback=white,colframe=cfon,boxrule=2pt,arc=0mm,
left=6mm,right=6mm,top=0mm,bottom=4mm,sharp corners,
overlay={\fill[cfon] (frame.north west) rectangle ([yshift=-1.25cm]frame.north east);
\node[anchor=west,text=white,font=\sffamily\bfseries\LARGE] at ([xshift=5mm,yshift=-0.62cm]frame.north west) {Question 42 C};
\node[anchor=east,text=white,font=\sffamily\itshape\large] at ([xshift=-5mm,yshift=-0.62cm]frame.north east) {Fonctions};}]
\vspace*{1.75cm}
\begin{minipage}[t][6.9cm][c]{\linewidth}\LARGE\centering
\begin{minipage}{0.95\linewidth}\raggedright\linespread{1.25}\selectfont Pour quels réels $x$ a-t-on $-2x+5>0$~?\end{minipage}
\end{minipage}

\vfill
\begin{tcolorbox}[enhanced,nobeforeafter,width=7.1cm,height=2.1cm,valign=center,colback=cfon!6,colframe=cfon,arc=2.5mm,boxrule=0.9pt,left=3mm]{\Large\bfseries\color{cfon}a)}\quad\begin{minipage}[c]{5.3cm}\raggedright\Large $x<-\dfrac{5}{2}$\end{minipage}\end{tcolorbox}\hfill\begin{tcolorbox}[enhanced,nobeforeafter,width=7.1cm,height=2.1cm,valign=center,colback=cfon!6,colframe=cfon,arc=2.5mm,boxrule=0.9pt,left=3mm]{\Large\bfseries\color{cfon}b)}\quad\begin{minipage}[c]{5.3cm}\raggedright\Large $x>-\dfrac{5}{2}$\end{minipage}\end{tcolorbox}\par\vspace{2.5mm}
\begin{tcolorbox}[enhanced,nobeforeafter,width=7.1cm,height=2.1cm,valign=center,colback=cfon!6,colframe=cfon,arc=2.5mm,boxrule=0.9pt,left=3mm]{\Large\bfseries\color{cfon}c)}\quad\begin{minipage}[c]{5.3cm}\raggedright\Large $x>\dfrac{5}{2}$\end{minipage}\end{tcolorbox}\hfill\begin{tcolorbox}[enhanced,nobeforeafter,width=7.1cm,height=2.1cm,valign=center,colback=cfon!6,colframe=cfon,arc=2.5mm,boxrule=0.9pt,left=3mm]{\Large\bfseries\color{cfon}d)}\quad\begin{minipage}[c]{5.3cm}\raggedright\Large $x<\dfrac{5}{2}$\end{minipage}\end{tcolorbox}
\end{tcolorbox}
\newpage
\begin{tcolorbox}[enhanced,width=16cm,height=14cm,colback=white,colframe=cfon,boxrule=2pt,arc=0mm,
left=6mm,right=6mm,top=0mm,bottom=4mm,sharp corners,
overlay={\fill[cfon] (frame.north west) rectangle ([yshift=-1.25cm]frame.north east);
\node[anchor=west,text=white,font=\sffamily\bfseries\LARGE] at ([xshift=5mm,yshift=-0.62cm]frame.north west) {Question 43 C};
\node[anchor=east,text=white,font=\sffamily\itshape\large] at ([xshift=-5mm,yshift=-0.62cm]frame.north east) {Fonctions};}]
\vspace*{1.75cm}
\begin{minipage}[t][6.9cm][c]{\linewidth}\LARGE\centering
\begin{minipage}{0.95\linewidth}\raggedright\linespread{1.25}\selectfont Quelles sont les solutions de l'équation $x^2=2$~?\end{minipage}
\end{minipage}

\vfill
\begin{tcolorbox}[enhanced,nobeforeafter,width=7.1cm,height=2.1cm,valign=center,colback=cfon!6,colframe=cfon,arc=2.5mm,boxrule=0.9pt,left=3mm]{\Large\bfseries\color{cfon}a)}\quad\begin{minipage}[c]{5.3cm}\raggedright\Large $\{-2\,;\,2\}$\end{minipage}\end{tcolorbox}\hfill\begin{tcolorbox}[enhanced,nobeforeafter,width=7.1cm,height=2.1cm,valign=center,colback=cfon!6,colframe=cfon,arc=2.5mm,boxrule=0.9pt,left=3mm]{\Large\bfseries\color{cfon}b)}\quad\begin{minipage}[c]{5.3cm}\raggedright\Large $\{\sqrt{2}\}$\end{minipage}\end{tcolorbox}\par\vspace{2.5mm}
\begin{tcolorbox}[enhanced,nobeforeafter,width=7.1cm,height=2.1cm,valign=center,colback=cfon!6,colframe=cfon,arc=2.5mm,boxrule=0.9pt,left=3mm]{\Large\bfseries\color{cfon}c)}\quad\begin{minipage}[c]{5.3cm}\raggedright\large $\{-\sqrt{2}\,;\,\sqrt{2}\}$\end{minipage}\end{tcolorbox}\hfill\begin{tcolorbox}[enhanced,nobeforeafter,width=7.1cm,height=2.1cm,valign=center,colback=cfon!6,colframe=cfon,arc=2.5mm,boxrule=0.9pt,left=3mm]{\Large\bfseries\color{cfon}d)}\quad\begin{minipage}[c]{5.3cm}\raggedright\Large aucune solution\end{minipage}\end{tcolorbox}
\end{tcolorbox}
\newpage
\begin{tcolorbox}[enhanced,width=16cm,height=14cm,colback=white,colframe=cfon,boxrule=2pt,arc=0mm,
left=6mm,right=6mm,top=0mm,bottom=4mm,sharp corners,
overlay={\fill[cfon] (frame.north west) rectangle ([yshift=-1.25cm]frame.north east);
\node[anchor=west,text=white,font=\sffamily\bfseries\LARGE] at ([xshift=5mm,yshift=-0.62cm]frame.north west) {Question 44 C};
\node[anchor=east,text=white,font=\sffamily\itshape\large] at ([xshift=-5mm,yshift=-0.62cm]frame.north east) {Fonctions};}]
\vspace*{1.75cm}
\begin{minipage}[t][6.9cm][c]{\linewidth}\LARGE\centering
\begin{minipage}{0.95\linewidth}\raggedright\linespread{1.25}\selectfont Quelles sont les solutions de l'équation $x^2=49$~?\end{minipage}
\end{minipage}

\vfill
\begin{tcolorbox}[enhanced,nobeforeafter,width=7.1cm,height=2.1cm,valign=center,colback=cfon!6,colframe=cfon,arc=2.5mm,boxrule=0.9pt,left=3mm]{\Large\bfseries\color{cfon}a)}\quad\begin{minipage}[c]{5.3cm}\raggedright\Large $\{-49\,;\,49\}$\end{minipage}\end{tcolorbox}\hfill\begin{tcolorbox}[enhanced,nobeforeafter,width=7.1cm,height=2.1cm,valign=center,colback=cfon!6,colframe=cfon,arc=2.5mm,boxrule=0.9pt,left=3mm]{\Large\bfseries\color{cfon}b)}\quad\begin{minipage}[c]{5.3cm}\raggedright\Large $\{-7\,;\,7\}$\end{minipage}\end{tcolorbox}\par\vspace{2.5mm}
\begin{tcolorbox}[enhanced,nobeforeafter,width=7.1cm,height=2.1cm,valign=center,colback=cfon!6,colframe=cfon,arc=2.5mm,boxrule=0.9pt,left=3mm]{\Large\bfseries\color{cfon}c)}\quad\begin{minipage}[c]{5.3cm}\raggedright\Large aucune solution\end{minipage}\end{tcolorbox}\hfill\begin{tcolorbox}[enhanced,nobeforeafter,width=7.1cm,height=2.1cm,valign=center,colback=cfon!6,colframe=cfon,arc=2.5mm,boxrule=0.9pt,left=3mm]{\Large\bfseries\color{cfon}d)}\quad\begin{minipage}[c]{5.3cm}\raggedright\Large $\{7\}$\end{minipage}\end{tcolorbox}
\end{tcolorbox}
\newpage
\begin{tcolorbox}[enhanced,width=16cm,height=14cm,colback=white,colframe=cfon,boxrule=2pt,arc=0mm,
left=6mm,right=6mm,top=0mm,bottom=4mm,sharp corners,
overlay={\fill[cfon] (frame.north west) rectangle ([yshift=-1.25cm]frame.north east);
\node[anchor=west,text=white,font=\sffamily\bfseries\LARGE] at ([xshift=5mm,yshift=-0.62cm]frame.north west) {Question 45 C};
\node[anchor=east,text=white,font=\sffamily\itshape\large] at ([xshift=-5mm,yshift=-0.62cm]frame.north east) {Fonctions};}]
\vspace*{1.75cm}
\begin{minipage}[t][6.9cm][c]{\linewidth}\LARGE\centering
\begin{minipage}{0.95\linewidth}\raggedright\linespread{1.25}\selectfont Quel est le sens de variation de la fonction $f(x)=7x$ sur $\mathbb{R}$~?\end{minipage}
\end{minipage}

\vfill
\begin{tcolorbox}[enhanced,nobeforeafter,width=7.1cm,height=2.1cm,valign=center,colback=cfon!6,colframe=cfon,arc=2.5mm,boxrule=0.9pt,left=3mm]{\Large\bfseries\color{cfon}a)}\quad\begin{minipage}[c]{5.3cm}\raggedright\large décroissante sur $\mathbb{R}$\end{minipage}\end{tcolorbox}\hfill\begin{tcolorbox}[enhanced,nobeforeafter,width=7.1cm,height=2.1cm,valign=center,colback=cfon!6,colframe=cfon,arc=2.5mm,boxrule=0.9pt,left=3mm]{\Large\bfseries\color{cfon}b)}\quad\begin{minipage}[c]{5.3cm}\raggedright\large croissante sur $\mathbb{R}$\end{minipage}\end{tcolorbox}\par\vspace{2.5mm}
\begin{tcolorbox}[enhanced,nobeforeafter,width=7.1cm,height=2.1cm,valign=center,colback=cfon!6,colframe=cfon,arc=2.5mm,boxrule=0.9pt,left=3mm]{\Large\bfseries\color{cfon}c)}\quad\begin{minipage}[c]{5.3cm}\raggedright\large croissante puis décroissante\end{minipage}\end{tcolorbox}\hfill\begin{tcolorbox}[enhanced,nobeforeafter,width=7.1cm,height=2.1cm,valign=center,colback=cfon!6,colframe=cfon,arc=2.5mm,boxrule=0.9pt,left=3mm]{\Large\bfseries\color{cfon}d)}\quad\begin{minipage}[c]{5.3cm}\raggedright\large constante sur $\mathbb{R}$\end{minipage}\end{tcolorbox}
\end{tcolorbox}
\newpage
\begin{tcolorbox}[enhanced,width=16cm,height=14cm,colback=white,colframe=cfon,boxrule=2pt,arc=0mm,
left=6mm,right=6mm,top=0mm,bottom=4mm,sharp corners,
overlay={\fill[cfon] (frame.north west) rectangle ([yshift=-1.25cm]frame.north east);
\node[anchor=west,text=white,font=\sffamily\bfseries\LARGE] at ([xshift=5mm,yshift=-0.62cm]frame.north west) {Question 46 C};
\node[anchor=east,text=white,font=\sffamily\itshape\large] at ([xshift=-5mm,yshift=-0.62cm]frame.north east) {Fonctions};}]
\vspace*{1.75cm}
\begin{minipage}[t][6.9cm][c]{\linewidth}\LARGE\centering
\begin{minipage}{0.95\linewidth}\raggedright\linespread{1.25}\selectfont Quel est le sens de variation de la fonction $f(x)=x-3$ sur $\mathbb{R}$~?\end{minipage}
\end{minipage}

\vfill
\begin{tcolorbox}[enhanced,nobeforeafter,width=7.1cm,height=2.1cm,valign=center,colback=cfon!6,colframe=cfon,arc=2.5mm,boxrule=0.9pt,left=3mm]{\Large\bfseries\color{cfon}a)}\quad\begin{minipage}[c]{5.3cm}\raggedright\large croissante sur $\mathbb{R}$\end{minipage}\end{tcolorbox}\hfill\begin{tcolorbox}[enhanced,nobeforeafter,width=7.1cm,height=2.1cm,valign=center,colback=cfon!6,colframe=cfon,arc=2.5mm,boxrule=0.9pt,left=3mm]{\Large\bfseries\color{cfon}b)}\quad\begin{minipage}[c]{5.3cm}\raggedright\large constante sur $\mathbb{R}$\end{minipage}\end{tcolorbox}\par\vspace{2.5mm}
\begin{tcolorbox}[enhanced,nobeforeafter,width=7.1cm,height=2.1cm,valign=center,colback=cfon!6,colframe=cfon,arc=2.5mm,boxrule=0.9pt,left=3mm]{\Large\bfseries\color{cfon}c)}\quad\begin{minipage}[c]{5.3cm}\raggedright\large décroissante sur $\mathbb{R}$\end{minipage}\end{tcolorbox}\hfill\begin{tcolorbox}[enhanced,nobeforeafter,width=7.1cm,height=2.1cm,valign=center,colback=cfon!6,colframe=cfon,arc=2.5mm,boxrule=0.9pt,left=3mm]{\Large\bfseries\color{cfon}d)}\quad\begin{minipage}[c]{5.3cm}\raggedright\large croissante puis décroissante\end{minipage}\end{tcolorbox}
\end{tcolorbox}
\newpage
\begin{tcolorbox}[enhanced,width=16cm,height=14cm,colback=white,colframe=cfon,boxrule=2pt,arc=0mm,
left=6mm,right=6mm,top=0mm,bottom=4mm,sharp corners,
overlay={\fill[cfon] (frame.north west) rectangle ([yshift=-1.25cm]frame.north east);
\node[anchor=west,text=white,font=\sffamily\bfseries\LARGE] at ([xshift=5mm,yshift=-0.62cm]frame.north west) {Question 47 C};
\node[anchor=east,text=white,font=\sffamily\itshape\large] at ([xshift=-5mm,yshift=-0.62cm]frame.north east) {Fonctions};}]
\vspace*{1.75cm}
\begin{minipage}[t][6.9cm][c]{\linewidth}\LARGE\centering
\begin{minipage}{0.95\linewidth}\raggedright\linespread{1.25}\selectfont Quelles sont les solutions de l'inéquation $x+7\leqslant 3x+1$~?\end{minipage}
\end{minipage}

\vfill
\begin{tcolorbox}[enhanced,nobeforeafter,width=7.1cm,height=2.1cm,valign=center,colback=cfon!6,colframe=cfon,arc=2.5mm,boxrule=0.9pt,left=3mm]{\Large\bfseries\color{cfon}a)}\quad\begin{minipage}[c]{5.3cm}\raggedright\Large $x\leqslant -4$\end{minipage}\end{tcolorbox}\hfill\begin{tcolorbox}[enhanced,nobeforeafter,width=7.1cm,height=2.1cm,valign=center,colback=cfon!6,colframe=cfon,arc=2.5mm,boxrule=0.9pt,left=3mm]{\Large\bfseries\color{cfon}b)}\quad\begin{minipage}[c]{5.3cm}\raggedright\Large $x\leqslant 3$\end{minipage}\end{tcolorbox}\par\vspace{2.5mm}
\begin{tcolorbox}[enhanced,nobeforeafter,width=7.1cm,height=2.1cm,valign=center,colback=cfon!6,colframe=cfon,arc=2.5mm,boxrule=0.9pt,left=3mm]{\Large\bfseries\color{cfon}c)}\quad\begin{minipage}[c]{5.3cm}\raggedright\Large $x\geqslant -3$\end{minipage}\end{tcolorbox}\hfill\begin{tcolorbox}[enhanced,nobeforeafter,width=7.1cm,height=2.1cm,valign=center,colback=cfon!6,colframe=cfon,arc=2.5mm,boxrule=0.9pt,left=3mm]{\Large\bfseries\color{cfon}d)}\quad\begin{minipage}[c]{5.3cm}\raggedright\Large $x\geqslant 3$\end{minipage}\end{tcolorbox}
\end{tcolorbox}
\newpage
\begin{tcolorbox}[enhanced,width=16cm,height=14cm,colback=white,colframe=cfon,boxrule=2pt,arc=0mm,
left=6mm,right=6mm,top=0mm,bottom=4mm,sharp corners,
overlay={\fill[cfon] (frame.north west) rectangle ([yshift=-1.25cm]frame.north east);
\node[anchor=west,text=white,font=\sffamily\bfseries\LARGE] at ([xshift=5mm,yshift=-0.62cm]frame.north west) {Question 48 C};
\node[anchor=east,text=white,font=\sffamily\itshape\large] at ([xshift=-5mm,yshift=-0.62cm]frame.north east) {Fonctions};}]
\vspace*{1.75cm}
\begin{minipage}[t][6.9cm][c]{\linewidth}\LARGE\centering
\begin{minipage}{0.95\linewidth}\raggedright\linespread{1.25}\selectfont Quelles sont les solutions de l'inéquation $5x\leqslant 2x-9$~?\end{minipage}
\end{minipage}

\vfill
\begin{tcolorbox}[enhanced,nobeforeafter,width=7.1cm,height=2.1cm,valign=center,colback=cfon!6,colframe=cfon,arc=2.5mm,boxrule=0.9pt,left=3mm]{\Large\bfseries\color{cfon}a)}\quad\begin{minipage}[c]{5.3cm}\raggedright\Large $x\leqslant -2$\end{minipage}\end{tcolorbox}\hfill\begin{tcolorbox}[enhanced,nobeforeafter,width=7.1cm,height=2.1cm,valign=center,colback=cfon!6,colframe=cfon,arc=2.5mm,boxrule=0.9pt,left=3mm]{\Large\bfseries\color{cfon}b)}\quad\begin{minipage}[c]{5.3cm}\raggedright\Large $x\leqslant -3$\end{minipage}\end{tcolorbox}\par\vspace{2.5mm}
\begin{tcolorbox}[enhanced,nobeforeafter,width=7.1cm,height=2.1cm,valign=center,colback=cfon!6,colframe=cfon,arc=2.5mm,boxrule=0.9pt,left=3mm]{\Large\bfseries\color{cfon}c)}\quad\begin{minipage}[c]{5.3cm}\raggedright\Large $x\geqslant 3$\end{minipage}\end{tcolorbox}\hfill\begin{tcolorbox}[enhanced,nobeforeafter,width=7.1cm,height=2.1cm,valign=center,colback=cfon!6,colframe=cfon,arc=2.5mm,boxrule=0.9pt,left=3mm]{\Large\bfseries\color{cfon}d)}\quad\begin{minipage}[c]{5.3cm}\raggedright\Large $x\geqslant -3$\end{minipage}\end{tcolorbox}
\end{tcolorbox}
\newpage
\begin{tcolorbox}[enhanced,width=16cm,height=14cm,colback=white,colframe=cvec,boxrule=2pt,arc=0mm,
left=6mm,right=6mm,top=0mm,bottom=4mm,sharp corners,
overlay={\fill[cvec] (frame.north west) rectangle ([yshift=-1.25cm]frame.north east);
\node[anchor=west,text=white,font=\sffamily\bfseries\LARGE] at ([xshift=5mm,yshift=-0.62cm]frame.north west) {Question 49 C};
\node[anchor=east,text=white,font=\sffamily\itshape\large] at ([xshift=-5mm,yshift=-0.62cm]frame.north east) {Vecteurs et repérage};}]
\vspace*{1.75cm}
\begin{minipage}[t][6.9cm][c]{\linewidth}\LARGE\centering
\begin{minipage}{0.95\linewidth}\raggedright\linespread{1.25}\selectfont Dans un repère, $A(-5\,;\,-2)$ et $B(1\,;\,-6)$. Quelles sont les coordonnées du vecteur $\overrightarrow{AB}$~?\end{minipage}
\end{minipage}

\vfill
\begin{tcolorbox}[enhanced,nobeforeafter,width=7.1cm,height=2.1cm,valign=center,colback=cvec!6,colframe=cvec,arc=2.5mm,boxrule=0.9pt,left=3mm]{\Large\bfseries\color{cvec}a)}\quad\begin{minipage}[c]{5.3cm}\raggedright\Large $(-4\,;\,-8)$\end{minipage}\end{tcolorbox}\hfill\begin{tcolorbox}[enhanced,nobeforeafter,width=7.1cm,height=2.1cm,valign=center,colback=cvec!6,colframe=cvec,arc=2.5mm,boxrule=0.9pt,left=3mm]{\Large\bfseries\color{cvec}b)}\quad\begin{minipage}[c]{5.3cm}\raggedright\Large $(-4\,;\,6)$\end{minipage}\end{tcolorbox}\par\vspace{2.5mm}
\begin{tcolorbox}[enhanced,nobeforeafter,width=7.1cm,height=2.1cm,valign=center,colback=cvec!6,colframe=cvec,arc=2.5mm,boxrule=0.9pt,left=3mm]{\Large\bfseries\color{cvec}c)}\quad\begin{minipage}[c]{5.3cm}\raggedright\Large $(6\,;\,-4)$\end{minipage}\end{tcolorbox}\hfill\begin{tcolorbox}[enhanced,nobeforeafter,width=7.1cm,height=2.1cm,valign=center,colback=cvec!6,colframe=cvec,arc=2.5mm,boxrule=0.9pt,left=3mm]{\Large\bfseries\color{cvec}d)}\quad\begin{minipage}[c]{5.3cm}\raggedright\Large $(-6\,;\,4)$\end{minipage}\end{tcolorbox}
\end{tcolorbox}
\newpage
\begin{tcolorbox}[enhanced,width=16cm,height=14cm,colback=white,colframe=cvec,boxrule=2pt,arc=0mm,
left=6mm,right=6mm,top=0mm,bottom=4mm,sharp corners,
overlay={\fill[cvec] (frame.north west) rectangle ([yshift=-1.25cm]frame.north east);
\node[anchor=west,text=white,font=\sffamily\bfseries\LARGE] at ([xshift=5mm,yshift=-0.62cm]frame.north west) {Question 50 C};
\node[anchor=east,text=white,font=\sffamily\itshape\large] at ([xshift=-5mm,yshift=-0.62cm]frame.north east) {Vecteurs et repérage};}]
\vspace*{1.75cm}
\begin{minipage}[t][6.9cm][c]{\linewidth}\LARGE\centering
\begin{minipage}{0.95\linewidth}\raggedright\linespread{1.25}\selectfont Dans un repère, $A(3\,;\,3)$ et $B(-1\,;\,8)$. Quelles sont les coordonnées du vecteur $\overrightarrow{AB}$~?\end{minipage}
\end{minipage}

\vfill
\begin{tcolorbox}[enhanced,nobeforeafter,width=7.1cm,height=2.1cm,valign=center,colback=cvec!6,colframe=cvec,arc=2.5mm,boxrule=0.9pt,left=3mm]{\Large\bfseries\color{cvec}a)}\quad\begin{minipage}[c]{5.3cm}\raggedright\Large $(5\,;\,-4)$\end{minipage}\end{tcolorbox}\hfill\begin{tcolorbox}[enhanced,nobeforeafter,width=7.1cm,height=2.1cm,valign=center,colback=cvec!6,colframe=cvec,arc=2.5mm,boxrule=0.9pt,left=3mm]{\Large\bfseries\color{cvec}b)}\quad\begin{minipage}[c]{5.3cm}\raggedright\Large $(-4\,;\,5)$\end{minipage}\end{tcolorbox}\par\vspace{2.5mm}
\begin{tcolorbox}[enhanced,nobeforeafter,width=7.1cm,height=2.1cm,valign=center,colback=cvec!6,colframe=cvec,arc=2.5mm,boxrule=0.9pt,left=3mm]{\Large\bfseries\color{cvec}c)}\quad\begin{minipage}[c]{5.3cm}\raggedright\Large $(4\,;\,-5)$\end{minipage}\end{tcolorbox}\hfill\begin{tcolorbox}[enhanced,nobeforeafter,width=7.1cm,height=2.1cm,valign=center,colback=cvec!6,colframe=cvec,arc=2.5mm,boxrule=0.9pt,left=3mm]{\Large\bfseries\color{cvec}d)}\quad\begin{minipage}[c]{5.3cm}\raggedright\Large $(2\,;\,11)$\end{minipage}\end{tcolorbox}
\end{tcolorbox}
\newpage
\begin{tcolorbox}[enhanced,width=16cm,height=14cm,colback=white,colframe=cvec,boxrule=2pt,arc=0mm,
left=6mm,right=6mm,top=0mm,bottom=4mm,sharp corners,
overlay={\fill[cvec] (frame.north west) rectangle ([yshift=-1.25cm]frame.north east);
\node[anchor=west,text=white,font=\sffamily\bfseries\LARGE] at ([xshift=5mm,yshift=-0.62cm]frame.north west) {Question 51 C};
\node[anchor=east,text=white,font=\sffamily\itshape\large] at ([xshift=-5mm,yshift=-0.62cm]frame.north east) {Vecteurs et repérage};}]
\vspace*{1.75cm}
\begin{minipage}[t][6.9cm][c]{\linewidth}\LARGE\centering
\begin{minipage}{0.95\linewidth}\raggedright\linespread{1.25}\selectfont Dans un repère, $A(-5\,;\,-2)$ et $B(1\,;\,6)$. Quelles sont les coordonnées du milieu de $[AB]$~?\end{minipage}
\end{minipage}

\vfill
\begin{tcolorbox}[enhanced,nobeforeafter,width=7.1cm,height=2.1cm,valign=center,colback=cvec!6,colframe=cvec,arc=2.5mm,boxrule=0.9pt,left=3mm]{\Large\bfseries\color{cvec}a)}\quad\begin{minipage}[c]{5.3cm}\raggedright\Large $(3\,;\,4)$\end{minipage}\end{tcolorbox}\hfill\begin{tcolorbox}[enhanced,nobeforeafter,width=7.1cm,height=2.1cm,valign=center,colback=cvec!6,colframe=cvec,arc=2.5mm,boxrule=0.9pt,left=3mm]{\Large\bfseries\color{cvec}b)}\quad\begin{minipage}[c]{5.3cm}\raggedright\Large $(-2\,;\,2)$\end{minipage}\end{tcolorbox}\par\vspace{2.5mm}
\begin{tcolorbox}[enhanced,nobeforeafter,width=7.1cm,height=2.1cm,valign=center,colback=cvec!6,colframe=cvec,arc=2.5mm,boxrule=0.9pt,left=3mm]{\Large\bfseries\color{cvec}c)}\quad\begin{minipage}[c]{5.3cm}\raggedright\Large $(-2\,;\,-4)$\end{minipage}\end{tcolorbox}\hfill\begin{tcolorbox}[enhanced,nobeforeafter,width=7.1cm,height=2.1cm,valign=center,colback=cvec!6,colframe=cvec,arc=2.5mm,boxrule=0.9pt,left=3mm]{\Large\bfseries\color{cvec}d)}\quad\begin{minipage}[c]{5.3cm}\raggedright\Large $(-4\,;\,4)$\end{minipage}\end{tcolorbox}
\end{tcolorbox}
\newpage
\begin{tcolorbox}[enhanced,width=16cm,height=14cm,colback=white,colframe=cvec,boxrule=2pt,arc=0mm,
left=6mm,right=6mm,top=0mm,bottom=4mm,sharp corners,
overlay={\fill[cvec] (frame.north west) rectangle ([yshift=-1.25cm]frame.north east);
\node[anchor=west,text=white,font=\sffamily\bfseries\LARGE] at ([xshift=5mm,yshift=-0.62cm]frame.north west) {Question 52 C};
\node[anchor=east,text=white,font=\sffamily\itshape\large] at ([xshift=-5mm,yshift=-0.62cm]frame.north east) {Vecteurs et repérage};}]
\vspace*{1.75cm}
\begin{minipage}[t][6.9cm][c]{\linewidth}\LARGE\centering
\begin{minipage}{0.95\linewidth}\raggedright\linespread{1.25}\selectfont Dans un repère, $A(2\,;\,9)$ et $B(-4\,;\,1)$. Quelles sont les coordonnées du milieu de $[AB]$~?\end{minipage}
\end{minipage}

\vfill
\begin{tcolorbox}[enhanced,nobeforeafter,width=7.1cm,height=2.1cm,valign=center,colback=cvec!6,colframe=cvec,arc=2.5mm,boxrule=0.9pt,left=3mm]{\Large\bfseries\color{cvec}a)}\quad\begin{minipage}[c]{5.3cm}\raggedright\Large $(-1\,;\,5)$\end{minipage}\end{tcolorbox}\hfill\begin{tcolorbox}[enhanced,nobeforeafter,width=7.1cm,height=2.1cm,valign=center,colback=cvec!6,colframe=cvec,arc=2.5mm,boxrule=0.9pt,left=3mm]{\Large\bfseries\color{cvec}b)}\quad\begin{minipage}[c]{5.3cm}\raggedright\Large $(-1\,;\,4)$\end{minipage}\end{tcolorbox}\par\vspace{2.5mm}
\begin{tcolorbox}[enhanced,nobeforeafter,width=7.1cm,height=2.1cm,valign=center,colback=cvec!6,colframe=cvec,arc=2.5mm,boxrule=0.9pt,left=3mm]{\Large\bfseries\color{cvec}c)}\quad\begin{minipage}[c]{5.3cm}\raggedright\Large $(-2\,;\,10)$\end{minipage}\end{tcolorbox}\hfill\begin{tcolorbox}[enhanced,nobeforeafter,width=7.1cm,height=2.1cm,valign=center,colback=cvec!6,colframe=cvec,arc=2.5mm,boxrule=0.9pt,left=3mm]{\Large\bfseries\color{cvec}d)}\quad\begin{minipage}[c]{5.3cm}\raggedright\Large $(-3\,;\,-4)$\end{minipage}\end{tcolorbox}
\end{tcolorbox}
\newpage
\begin{tcolorbox}[enhanced,width=16cm,height=14cm,colback=white,colframe=cvec,boxrule=2pt,arc=0mm,
left=6mm,right=6mm,top=0mm,bottom=4mm,sharp corners,
overlay={\fill[cvec] (frame.north west) rectangle ([yshift=-1.25cm]frame.north east);
\node[anchor=west,text=white,font=\sffamily\bfseries\LARGE] at ([xshift=5mm,yshift=-0.62cm]frame.north west) {Question 53 C};
\node[anchor=east,text=white,font=\sffamily\itshape\large] at ([xshift=-5mm,yshift=-0.62cm]frame.north east) {Vecteurs et repérage};}]
\vspace*{1.75cm}
\begin{minipage}[t][6.9cm][c]{\linewidth}\LARGE\centering
\begin{minipage}{0.95\linewidth}\raggedright\linespread{1.25}\selectfont Dans un repère orthonormé, $A(2\,;\,5)$ et $B(4\,;\,1)$. Quelle est la distance $AB$~?\end{minipage}
\end{minipage}

\vfill
\begin{tcolorbox}[enhanced,nobeforeafter,width=7.1cm,height=2.1cm,valign=center,colback=cvec!6,colframe=cvec,arc=2.5mm,boxrule=0.9pt,left=3mm]{\Large\bfseries\color{cvec}a)}\quad\begin{minipage}[c]{5.3cm}\raggedright\Large $2\sqrt{3}$\end{minipage}\end{tcolorbox}\hfill\begin{tcolorbox}[enhanced,nobeforeafter,width=7.1cm,height=2.1cm,valign=center,colback=cvec!6,colframe=cvec,arc=2.5mm,boxrule=0.9pt,left=3mm]{\Large\bfseries\color{cvec}b)}\quad\begin{minipage}[c]{5.3cm}\raggedright\Large $20$\end{minipage}\end{tcolorbox}\par\vspace{2.5mm}
\begin{tcolorbox}[enhanced,nobeforeafter,width=7.1cm,height=2.1cm,valign=center,colback=cvec!6,colframe=cvec,arc=2.5mm,boxrule=0.9pt,left=3mm]{\Large\bfseries\color{cvec}c)}\quad\begin{minipage}[c]{5.3cm}\raggedright\Large $2\sqrt{5}$\end{minipage}\end{tcolorbox}\hfill\begin{tcolorbox}[enhanced,nobeforeafter,width=7.1cm,height=2.1cm,valign=center,colback=cvec!6,colframe=cvec,arc=2.5mm,boxrule=0.9pt,left=3mm]{\Large\bfseries\color{cvec}d)}\quad\begin{minipage}[c]{5.3cm}\raggedright\Large $6$\end{minipage}\end{tcolorbox}
\end{tcolorbox}
\newpage
\begin{tcolorbox}[enhanced,width=16cm,height=14cm,colback=white,colframe=cvec,boxrule=2pt,arc=0mm,
left=6mm,right=6mm,top=0mm,bottom=4mm,sharp corners,
overlay={\fill[cvec] (frame.north west) rectangle ([yshift=-1.25cm]frame.north east);
\node[anchor=west,text=white,font=\sffamily\bfseries\LARGE] at ([xshift=5mm,yshift=-0.62cm]frame.north west) {Question 54 C};
\node[anchor=east,text=white,font=\sffamily\itshape\large] at ([xshift=-5mm,yshift=-0.62cm]frame.north east) {Vecteurs et repérage};}]
\vspace*{1.75cm}
\begin{minipage}[t][6.9cm][c]{\linewidth}\LARGE\centering
\begin{minipage}{0.95\linewidth}\raggedright\linespread{1.25}\selectfont Dans un repère orthonormé, $A(-1\,;\,-2)$ et $B(7\,;\,4)$. Quelle est la distance $AB$~?\end{minipage}
\end{minipage}

\vfill
\begin{tcolorbox}[enhanced,nobeforeafter,width=7.1cm,height=2.1cm,valign=center,colback=cvec!6,colframe=cvec,arc=2.5mm,boxrule=0.9pt,left=3mm]{\Large\bfseries\color{cvec}a)}\quad\begin{minipage}[c]{5.3cm}\raggedright\Large $14$\end{minipage}\end{tcolorbox}\hfill\begin{tcolorbox}[enhanced,nobeforeafter,width=7.1cm,height=2.1cm,valign=center,colback=cvec!6,colframe=cvec,arc=2.5mm,boxrule=0.9pt,left=3mm]{\Large\bfseries\color{cvec}b)}\quad\begin{minipage}[c]{5.3cm}\raggedright\Large $2\sqrt{7}$\end{minipage}\end{tcolorbox}\par\vspace{2.5mm}
\begin{tcolorbox}[enhanced,nobeforeafter,width=7.1cm,height=2.1cm,valign=center,colback=cvec!6,colframe=cvec,arc=2.5mm,boxrule=0.9pt,left=3mm]{\Large\bfseries\color{cvec}c)}\quad\begin{minipage}[c]{5.3cm}\raggedright\Large $100$\end{minipage}\end{tcolorbox}\hfill\begin{tcolorbox}[enhanced,nobeforeafter,width=7.1cm,height=2.1cm,valign=center,colback=cvec!6,colframe=cvec,arc=2.5mm,boxrule=0.9pt,left=3mm]{\Large\bfseries\color{cvec}d)}\quad\begin{minipage}[c]{5.3cm}\raggedright\Large $10$\end{minipage}\end{tcolorbox}
\end{tcolorbox}
\newpage
\begin{tcolorbox}[enhanced,width=16cm,height=14cm,colback=white,colframe=cvec,boxrule=2pt,arc=0mm,
left=6mm,right=6mm,top=0mm,bottom=4mm,sharp corners,
overlay={\fill[cvec] (frame.north west) rectangle ([yshift=-1.25cm]frame.north east);
\node[anchor=west,text=white,font=\sffamily\bfseries\LARGE] at ([xshift=5mm,yshift=-0.62cm]frame.north west) {Question 55 C};
\node[anchor=east,text=white,font=\sffamily\itshape\large] at ([xshift=-5mm,yshift=-0.62cm]frame.north east) {Vecteurs et repérage};}]
\vspace*{1.75cm}
\begin{minipage}[t][6.9cm][c]{\linewidth}\LARGE\centering
\begin{minipage}{0.95\linewidth}\raggedright\linespread{1.25}\selectfont Pour quelle valeur de $k$ les vecteurs $\vec u\begin{pmatrix}6\\4\end{pmatrix}$ et $\vec v\begin{pmatrix}k\\10\end{pmatrix}$ sont-ils colinéaires~?\end{minipage}
\end{minipage}

\vfill
\begin{tcolorbox}[enhanced,nobeforeafter,width=7.1cm,height=2.1cm,valign=center,colback=cvec!6,colframe=cvec,arc=2.5mm,boxrule=0.9pt,left=3mm]{\Large\bfseries\color{cvec}a)}\quad\begin{minipage}[c]{5.3cm}\raggedright\Large $k=-15$\end{minipage}\end{tcolorbox}\hfill\begin{tcolorbox}[enhanced,nobeforeafter,width=7.1cm,height=2.1cm,valign=center,colback=cvec!6,colframe=cvec,arc=2.5mm,boxrule=0.9pt,left=3mm]{\Large\bfseries\color{cvec}b)}\quad\begin{minipage}[c]{5.3cm}\raggedright\Large $k=\dfrac{20}{3}$\end{minipage}\end{tcolorbox}\par\vspace{2.5mm}
\begin{tcolorbox}[enhanced,nobeforeafter,width=7.1cm,height=2.1cm,valign=center,colback=cvec!6,colframe=cvec,arc=2.5mm,boxrule=0.9pt,left=3mm]{\Large\bfseries\color{cvec}c)}\quad\begin{minipage}[c]{5.3cm}\raggedright\Large $k=15$\end{minipage}\end{tcolorbox}\hfill\begin{tcolorbox}[enhanced,nobeforeafter,width=7.1cm,height=2.1cm,valign=center,colback=cvec!6,colframe=cvec,arc=2.5mm,boxrule=0.9pt,left=3mm]{\Large\bfseries\color{cvec}d)}\quad\begin{minipage}[c]{5.3cm}\raggedright\Large $k=12$\end{minipage}\end{tcolorbox}
\end{tcolorbox}
\newpage
\begin{tcolorbox}[enhanced,width=16cm,height=14cm,colback=white,colframe=cvec,boxrule=2pt,arc=0mm,
left=6mm,right=6mm,top=0mm,bottom=4mm,sharp corners,
overlay={\fill[cvec] (frame.north west) rectangle ([yshift=-1.25cm]frame.north east);
\node[anchor=west,text=white,font=\sffamily\bfseries\LARGE] at ([xshift=5mm,yshift=-0.62cm]frame.north west) {Question 56 C};
\node[anchor=east,text=white,font=\sffamily\itshape\large] at ([xshift=-5mm,yshift=-0.62cm]frame.north east) {Vecteurs et repérage};}]
\vspace*{1.75cm}
\begin{minipage}[t][6.9cm][c]{\linewidth}\LARGE\centering
\begin{minipage}{0.95\linewidth}\raggedright\linespread{1.25}\selectfont Pour quelle valeur de $k$ les vecteurs $\vec u\begin{pmatrix}5\\2\end{pmatrix}$ et $\vec v\begin{pmatrix}k\\-3\end{pmatrix}$ sont-ils colinéaires~?\end{minipage}
\end{minipage}

\vfill
\begin{tcolorbox}[enhanced,nobeforeafter,width=7.1cm,height=2.1cm,valign=center,colback=cvec!6,colframe=cvec,arc=2.5mm,boxrule=0.9pt,left=3mm]{\Large\bfseries\color{cvec}a)}\quad\begin{minipage}[c]{5.3cm}\raggedright\Large $k=0$\end{minipage}\end{tcolorbox}\hfill\begin{tcolorbox}[enhanced,nobeforeafter,width=7.1cm,height=2.1cm,valign=center,colback=cvec!6,colframe=cvec,arc=2.5mm,boxrule=0.9pt,left=3mm]{\Large\bfseries\color{cvec}b)}\quad\begin{minipage}[c]{5.3cm}\raggedright\Large $k=\dfrac{15}{2}$\end{minipage}\end{tcolorbox}\par\vspace{2.5mm}
\begin{tcolorbox}[enhanced,nobeforeafter,width=7.1cm,height=2.1cm,valign=center,colback=cvec!6,colframe=cvec,arc=2.5mm,boxrule=0.9pt,left=3mm]{\Large\bfseries\color{cvec}c)}\quad\begin{minipage}[c]{5.3cm}\raggedright\Large $k=-\dfrac{15}{2}$\end{minipage}\end{tcolorbox}\hfill\begin{tcolorbox}[enhanced,nobeforeafter,width=7.1cm,height=2.1cm,valign=center,colback=cvec!6,colframe=cvec,arc=2.5mm,boxrule=0.9pt,left=3mm]{\Large\bfseries\color{cvec}d)}\quad\begin{minipage}[c]{5.3cm}\raggedright\Large $k=-\dfrac{6}{5}$\end{minipage}\end{tcolorbox}
\end{tcolorbox}
\newpage
\begin{tcolorbox}[enhanced,width=16cm,height=14cm,colback=white,colframe=cvec,boxrule=2pt,arc=0mm,
left=6mm,right=6mm,top=0mm,bottom=4mm,sharp corners,
overlay={\fill[cvec] (frame.north west) rectangle ([yshift=-1.25cm]frame.north east);
\node[anchor=west,text=white,font=\sffamily\bfseries\LARGE] at ([xshift=5mm,yshift=-0.62cm]frame.north west) {Question 57 C};
\node[anchor=east,text=white,font=\sffamily\itshape\large] at ([xshift=-5mm,yshift=-0.62cm]frame.north east) {Vecteurs et repérage};}]
\vspace*{1.75cm}
\begin{minipage}[t][6.9cm][c]{\linewidth}\Large\centering
\begin{minipage}{0.95\linewidth}\raggedright\linespread{1.25}\selectfont On donne $\vec u\begin{pmatrix}3\\-1\end{pmatrix}$ et $\vec v\begin{pmatrix}2\\5\end{pmatrix}$. Quelles sont les coordonnées du vecteur $2\vec u-3\vec v$~?\end{minipage}
\end{minipage}

\vfill
\begin{tcolorbox}[enhanced,nobeforeafter,width=7.1cm,height=2.1cm,valign=center,colback=cvec!6,colframe=cvec,arc=2.5mm,boxrule=0.9pt,left=3mm]{\Large\bfseries\color{cvec}a)}\quad\begin{minipage}[c]{5.3cm}\raggedright\Large $(5\,;\,4)$\end{minipage}\end{tcolorbox}\hfill\begin{tcolorbox}[enhanced,nobeforeafter,width=7.1cm,height=2.1cm,valign=center,colback=cvec!6,colframe=cvec,arc=2.5mm,boxrule=0.9pt,left=3mm]{\Large\bfseries\color{cvec}b)}\quad\begin{minipage}[c]{5.3cm}\raggedright\Large $(12\,;\,13)$\end{minipage}\end{tcolorbox}\par\vspace{2.5mm}
\begin{tcolorbox}[enhanced,nobeforeafter,width=7.1cm,height=2.1cm,valign=center,colback=cvec!6,colframe=cvec,arc=2.5mm,boxrule=0.9pt,left=3mm]{\Large\bfseries\color{cvec}c)}\quad\begin{minipage}[c]{5.3cm}\raggedright\Large $(8\,;\,3)$\end{minipage}\end{tcolorbox}\hfill\begin{tcolorbox}[enhanced,nobeforeafter,width=7.1cm,height=2.1cm,valign=center,colback=cvec!6,colframe=cvec,arc=2.5mm,boxrule=0.9pt,left=3mm]{\Large\bfseries\color{cvec}d)}\quad\begin{minipage}[c]{5.3cm}\raggedright\Large $(0\,;\,-17)$\end{minipage}\end{tcolorbox}
\end{tcolorbox}
\newpage
\begin{tcolorbox}[enhanced,width=16cm,height=14cm,colback=white,colframe=cvec,boxrule=2pt,arc=0mm,
left=6mm,right=6mm,top=0mm,bottom=4mm,sharp corners,
overlay={\fill[cvec] (frame.north west) rectangle ([yshift=-1.25cm]frame.north east);
\node[anchor=west,text=white,font=\sffamily\bfseries\LARGE] at ([xshift=5mm,yshift=-0.62cm]frame.north west) {Question 58 C};
\node[anchor=east,text=white,font=\sffamily\itshape\large] at ([xshift=-5mm,yshift=-0.62cm]frame.north east) {Vecteurs et repérage};}]
\vspace*{1.75cm}
\begin{minipage}[t][6.9cm][c]{\linewidth}\Large\centering
\begin{minipage}{0.95\linewidth}\raggedright\linespread{1.25}\selectfont On donne $\vec u\begin{pmatrix}0\\4\end{pmatrix}$ et $\vec v\begin{pmatrix}5\\-2\end{pmatrix}$. Quelles sont les coordonnées du vecteur $\vec u+3\vec v$~?\end{minipage}
\end{minipage}

\vfill
\begin{tcolorbox}[enhanced,nobeforeafter,width=7.1cm,height=2.1cm,valign=center,colback=cvec!6,colframe=cvec,arc=2.5mm,boxrule=0.9pt,left=3mm]{\Large\bfseries\color{cvec}a)}\quad\begin{minipage}[c]{5.3cm}\raggedright\Large $(5\,;\,2)$\end{minipage}\end{tcolorbox}\hfill\begin{tcolorbox}[enhanced,nobeforeafter,width=7.1cm,height=2.1cm,valign=center,colback=cvec!6,colframe=cvec,arc=2.5mm,boxrule=0.9pt,left=3mm]{\Large\bfseries\color{cvec}b)}\quad\begin{minipage}[c]{5.3cm}\raggedright\Large $(15\,;\,-2)$\end{minipage}\end{tcolorbox}\par\vspace{2.5mm}
\begin{tcolorbox}[enhanced,nobeforeafter,width=7.1cm,height=2.1cm,valign=center,colback=cvec!6,colframe=cvec,arc=2.5mm,boxrule=0.9pt,left=3mm]{\Large\bfseries\color{cvec}c)}\quad\begin{minipage}[c]{5.3cm}\raggedright\Large $(-2\,;\,15)$\end{minipage}\end{tcolorbox}\hfill\begin{tcolorbox}[enhanced,nobeforeafter,width=7.1cm,height=2.1cm,valign=center,colback=cvec!6,colframe=cvec,arc=2.5mm,boxrule=0.9pt,left=3mm]{\Large\bfseries\color{cvec}d)}\quad\begin{minipage}[c]{5.3cm}\raggedright\Large $(-15\,;\,10)$\end{minipage}\end{tcolorbox}
\end{tcolorbox}
\newpage
\begin{tcolorbox}[enhanced,width=16cm,height=14cm,colback=white,colframe=cvec,boxrule=2pt,arc=0mm,
left=6mm,right=6mm,top=0mm,bottom=4mm,sharp corners,
overlay={\fill[cvec] (frame.north west) rectangle ([yshift=-1.25cm]frame.north east);
\node[anchor=west,text=white,font=\sffamily\bfseries\LARGE] at ([xshift=5mm,yshift=-0.62cm]frame.north west) {Question 59 C};
\node[anchor=east,text=white,font=\sffamily\itshape\large] at ([xshift=-5mm,yshift=-0.62cm]frame.north east) {Vecteurs et repérage};}]
\vspace*{1.75cm}
\begin{minipage}[t][6.9cm][c]{\linewidth}\LARGE\centering
\begin{minipage}{0.95\linewidth}\raggedright\linespread{1.25}\selectfont Que vaut le déterminant $\det(\vec u,\vec v)$ pour $\vec u\begin{pmatrix}0\\3\end{pmatrix}$ et $\vec v\begin{pmatrix}2\\7\end{pmatrix}$~?\end{minipage}
\end{minipage}

\vfill
\begin{tcolorbox}[enhanced,nobeforeafter,width=7.1cm,height=2.1cm,valign=center,colback=cvec!6,colframe=cvec,arc=2.5mm,boxrule=0.9pt,left=3mm]{\Large\bfseries\color{cvec}a)}\quad\begin{minipage}[c]{5.3cm}\raggedright\Large $-5$\end{minipage}\end{tcolorbox}\hfill\begin{tcolorbox}[enhanced,nobeforeafter,width=7.1cm,height=2.1cm,valign=center,colback=cvec!6,colframe=cvec,arc=2.5mm,boxrule=0.9pt,left=3mm]{\Large\bfseries\color{cvec}b)}\quad\begin{minipage}[c]{5.3cm}\raggedright\Large $-6$\end{minipage}\end{tcolorbox}\par\vspace{2.5mm}
\begin{tcolorbox}[enhanced,nobeforeafter,width=7.1cm,height=2.1cm,valign=center,colback=cvec!6,colframe=cvec,arc=2.5mm,boxrule=0.9pt,left=3mm]{\Large\bfseries\color{cvec}c)}\quad\begin{minipage}[c]{5.3cm}\raggedright\Large $21$\end{minipage}\end{tcolorbox}\hfill\begin{tcolorbox}[enhanced,nobeforeafter,width=7.1cm,height=2.1cm,valign=center,colback=cvec!6,colframe=cvec,arc=2.5mm,boxrule=0.9pt,left=3mm]{\Large\bfseries\color{cvec}d)}\quad\begin{minipage}[c]{5.3cm}\raggedright\Large $6$\end{minipage}\end{tcolorbox}
\end{tcolorbox}
\newpage
\begin{tcolorbox}[enhanced,width=16cm,height=14cm,colback=white,colframe=cvec,boxrule=2pt,arc=0mm,
left=6mm,right=6mm,top=0mm,bottom=4mm,sharp corners,
overlay={\fill[cvec] (frame.north west) rectangle ([yshift=-1.25cm]frame.north east);
\node[anchor=west,text=white,font=\sffamily\bfseries\LARGE] at ([xshift=5mm,yshift=-0.62cm]frame.north west) {Question 60 C};
\node[anchor=east,text=white,font=\sffamily\itshape\large] at ([xshift=-5mm,yshift=-0.62cm]frame.north east) {Vecteurs et repérage};}]
\vspace*{1.75cm}
\begin{minipage}[t][6.9cm][c]{\linewidth}\LARGE\centering
\begin{minipage}{0.95\linewidth}\raggedright\linespread{1.25}\selectfont Que vaut le déterminant $\det(\vec u,\vec v)$ pour $\vec u\begin{pmatrix}5\\-1\end{pmatrix}$ et $\vec v\begin{pmatrix}3\\2\end{pmatrix}$~?\end{minipage}
\end{minipage}

\vfill
\begin{tcolorbox}[enhanced,nobeforeafter,width=7.1cm,height=2.1cm,valign=center,colback=cvec!6,colframe=cvec,arc=2.5mm,boxrule=0.9pt,left=3mm]{\Large\bfseries\color{cvec}a)}\quad\begin{minipage}[c]{5.3cm}\raggedright\Large $14$\end{minipage}\end{tcolorbox}\hfill\begin{tcolorbox}[enhanced,nobeforeafter,width=7.1cm,height=2.1cm,valign=center,colback=cvec!6,colframe=cvec,arc=2.5mm,boxrule=0.9pt,left=3mm]{\Large\bfseries\color{cvec}b)}\quad\begin{minipage}[c]{5.3cm}\raggedright\Large $-13$\end{minipage}\end{tcolorbox}\par\vspace{2.5mm}
\begin{tcolorbox}[enhanced,nobeforeafter,width=7.1cm,height=2.1cm,valign=center,colback=cvec!6,colframe=cvec,arc=2.5mm,boxrule=0.9pt,left=3mm]{\Large\bfseries\color{cvec}c)}\quad\begin{minipage}[c]{5.3cm}\raggedright\Large $13$\end{minipage}\end{tcolorbox}\hfill\begin{tcolorbox}[enhanced,nobeforeafter,width=7.1cm,height=2.1cm,valign=center,colback=cvec!6,colframe=cvec,arc=2.5mm,boxrule=0.9pt,left=3mm]{\Large\bfseries\color{cvec}d)}\quad\begin{minipage}[c]{5.3cm}\raggedright\Large $7$\end{minipage}\end{tcolorbox}
\end{tcolorbox}
\newpage
\begin{tcolorbox}[enhanced,width=16cm,height=14cm,colback=white,colframe=cvec,boxrule=2pt,arc=0mm,
left=6mm,right=6mm,top=0mm,bottom=4mm,sharp corners,
overlay={\fill[cvec] (frame.north west) rectangle ([yshift=-1.25cm]frame.north east);
\node[anchor=west,text=white,font=\sffamily\bfseries\LARGE] at ([xshift=5mm,yshift=-0.62cm]frame.north west) {Question 61 C};
\node[anchor=east,text=white,font=\sffamily\itshape\large] at ([xshift=-5mm,yshift=-0.62cm]frame.north east) {Vecteurs et repérage};}]
\vspace*{1.75cm}
\begin{minipage}[t][6.9cm][c]{\linewidth}\LARGE\centering
\begin{minipage}{0.95\linewidth}\raggedright\linespread{1.25}\selectfont Dans un repère, $A(-3\,;\,-2)$, $B(0\,;\,-4)$ et $C(2\,;\,1)$. Quelles sont les coordonnées du point $D$ tel que $ABCD$ soit un parallélogramme~?\end{minipage}
\end{minipage}

\vfill
\begin{tcolorbox}[enhanced,nobeforeafter,width=7.1cm,height=2.1cm,valign=center,colback=cvec!6,colframe=cvec,arc=2.5mm,boxrule=0.9pt,left=3mm]{\Large\bfseries\color{cvec}a)}\quad\begin{minipage}[c]{5.3cm}\raggedright\Large $(-5\,;\,-7)$\end{minipage}\end{tcolorbox}\hfill\begin{tcolorbox}[enhanced,nobeforeafter,width=7.1cm,height=2.1cm,valign=center,colback=cvec!6,colframe=cvec,arc=2.5mm,boxrule=0.9pt,left=3mm]{\Large\bfseries\color{cvec}b)}\quad\begin{minipage}[c]{5.3cm}\raggedright\Large $(5\,;\,-1)$\end{minipage}\end{tcolorbox}\par\vspace{2.5mm}
\begin{tcolorbox}[enhanced,nobeforeafter,width=7.1cm,height=2.1cm,valign=center,colback=cvec!6,colframe=cvec,arc=2.5mm,boxrule=0.9pt,left=3mm]{\Large\bfseries\color{cvec}c)}\quad\begin{minipage}[c]{5.3cm}\raggedright\Large $(-1\,;\,3)$\end{minipage}\end{tcolorbox}\hfill\begin{tcolorbox}[enhanced,nobeforeafter,width=7.1cm,height=2.1cm,valign=center,colback=cvec!6,colframe=cvec,arc=2.5mm,boxrule=0.9pt,left=3mm]{\Large\bfseries\color{cvec}d)}\quad\begin{minipage}[c]{5.3cm}\raggedright\Large $(-0{,}5\,;\,-0{,}5)$\end{minipage}\end{tcolorbox}
\end{tcolorbox}
\newpage
\begin{tcolorbox}[enhanced,width=16cm,height=14cm,colback=white,colframe=cvec,boxrule=2pt,arc=0mm,
left=6mm,right=6mm,top=0mm,bottom=4mm,sharp corners,
overlay={\fill[cvec] (frame.north west) rectangle ([yshift=-1.25cm]frame.north east);
\node[anchor=west,text=white,font=\sffamily\bfseries\LARGE] at ([xshift=5mm,yshift=-0.62cm]frame.north west) {Question 62 C};
\node[anchor=east,text=white,font=\sffamily\itshape\large] at ([xshift=-5mm,yshift=-0.62cm]frame.north east) {Vecteurs et repérage};}]
\vspace*{1.75cm}
\begin{minipage}[t][6.9cm][c]{\linewidth}\LARGE\centering
\begin{minipage}{0.95\linewidth}\raggedright\linespread{1.25}\selectfont Dans un repère, $A(1\,;\,4)$, $B(2\,;\,0)$ et $C(6\,;\,1)$. Quelles sont les coordonnées du point $D$ tel que $ABCD$ soit un parallélogramme~?\end{minipage}
\end{minipage}

\vfill
\begin{tcolorbox}[enhanced,nobeforeafter,width=7.1cm,height=2.1cm,valign=center,colback=cvec!6,colframe=cvec,arc=2.5mm,boxrule=0.9pt,left=3mm]{\Large\bfseries\color{cvec}a)}\quad\begin{minipage}[c]{5.3cm}\raggedright\Large $(7\,;\,-3)$\end{minipage}\end{tcolorbox}\hfill\begin{tcolorbox}[enhanced,nobeforeafter,width=7.1cm,height=2.1cm,valign=center,colback=cvec!6,colframe=cvec,arc=2.5mm,boxrule=0.9pt,left=3mm]{\Large\bfseries\color{cvec}b)}\quad\begin{minipage}[c]{5.3cm}\raggedright\Large $(3{,}5\,;\,2{,}5)$\end{minipage}\end{tcolorbox}\par\vspace{2.5mm}
\begin{tcolorbox}[enhanced,nobeforeafter,width=7.1cm,height=2.1cm,valign=center,colback=cvec!6,colframe=cvec,arc=2.5mm,boxrule=0.9pt,left=3mm]{\Large\bfseries\color{cvec}c)}\quad\begin{minipage}[c]{5.3cm}\raggedright\Large $(5\,;\,5)$\end{minipage}\end{tcolorbox}\hfill\begin{tcolorbox}[enhanced,nobeforeafter,width=7.1cm,height=2.1cm,valign=center,colback=cvec!6,colframe=cvec,arc=2.5mm,boxrule=0.9pt,left=3mm]{\Large\bfseries\color{cvec}d)}\quad\begin{minipage}[c]{5.3cm}\raggedright\Large $(-3\,;\,3)$\end{minipage}\end{tcolorbox}
\end{tcolorbox}
\newpage
\begin{tcolorbox}[enhanced,width=16cm,height=14cm,colback=white,colframe=cvec,boxrule=2pt,arc=0mm,
left=6mm,right=6mm,top=0mm,bottom=4mm,sharp corners,
overlay={\fill[cvec] (frame.north west) rectangle ([yshift=-1.25cm]frame.north east);
\node[anchor=west,text=white,font=\sffamily\bfseries\LARGE] at ([xshift=5mm,yshift=-0.62cm]frame.north west) {Question 63 C};
\node[anchor=east,text=white,font=\sffamily\itshape\large] at ([xshift=-5mm,yshift=-0.62cm]frame.north east) {Vecteurs et repérage};}]
\vspace*{1.75cm}
\begin{minipage}[t][6.9cm][c]{\linewidth}\LARGE\centering
\begin{minipage}{0.95\linewidth}\raggedright\linespread{1.25}\selectfont Quel est le couple solution du système $\left\{\begin{array}{l}x-y=3\\2x+3y=16\end{array}\right.$~?\end{minipage}
\end{minipage}

\vfill
\begin{tcolorbox}[enhanced,nobeforeafter,width=7.1cm,height=2.1cm,valign=center,colback=cvec!6,colframe=cvec,arc=2.5mm,boxrule=0.9pt,left=3mm]{\Large\bfseries\color{cvec}a)}\quad\begin{minipage}[c]{5.3cm}\raggedright\Large $(5\,;\,2)$\end{minipage}\end{tcolorbox}\hfill\begin{tcolorbox}[enhanced,nobeforeafter,width=7.1cm,height=2.1cm,valign=center,colback=cvec!6,colframe=cvec,arc=2.5mm,boxrule=0.9pt,left=3mm]{\Large\bfseries\color{cvec}b)}\quad\begin{minipage}[c]{5.3cm}\raggedright\Large $(2\,;\,5)$\end{minipage}\end{tcolorbox}\par\vspace{2.5mm}
\begin{tcolorbox}[enhanced,nobeforeafter,width=7.1cm,height=2.1cm,valign=center,colback=cvec!6,colframe=cvec,arc=2.5mm,boxrule=0.9pt,left=3mm]{\Large\bfseries\color{cvec}c)}\quad\begin{minipage}[c]{5.3cm}\raggedright\Large $(5\,;\,-2)$\end{minipage}\end{tcolorbox}\hfill\begin{tcolorbox}[enhanced,nobeforeafter,width=7.1cm,height=2.1cm,valign=center,colback=cvec!6,colframe=cvec,arc=2.5mm,boxrule=0.9pt,left=3mm]{\Large\bfseries\color{cvec}d)}\quad\begin{minipage}[c]{5.3cm}\raggedright\Large $(-5\,;\,2)$\end{minipage}\end{tcolorbox}
\end{tcolorbox}
\newpage
\begin{tcolorbox}[enhanced,width=16cm,height=14cm,colback=white,colframe=cvec,boxrule=2pt,arc=0mm,
left=6mm,right=6mm,top=0mm,bottom=4mm,sharp corners,
overlay={\fill[cvec] (frame.north west) rectangle ([yshift=-1.25cm]frame.north east);
\node[anchor=west,text=white,font=\sffamily\bfseries\LARGE] at ([xshift=5mm,yshift=-0.62cm]frame.north west) {Question 64 C};
\node[anchor=east,text=white,font=\sffamily\itshape\large] at ([xshift=-5mm,yshift=-0.62cm]frame.north east) {Vecteurs et repérage};}]
\vspace*{1.75cm}
\begin{minipage}[t][6.9cm][c]{\linewidth}\LARGE\centering
\begin{minipage}{0.95\linewidth}\raggedright\linespread{1.25}\selectfont Quel est le couple solution du système $\left\{\begin{array}{l}4x+y=7\\x-2y=-5\end{array}\right.$~?\end{minipage}
\end{minipage}

\vfill
\begin{tcolorbox}[enhanced,nobeforeafter,width=7.1cm,height=2.1cm,valign=center,colback=cvec!6,colframe=cvec,arc=2.5mm,boxrule=0.9pt,left=3mm]{\Large\bfseries\color{cvec}a)}\quad\begin{minipage}[c]{5.3cm}\raggedright\Large $(3\,;\,1)$\end{minipage}\end{tcolorbox}\hfill\begin{tcolorbox}[enhanced,nobeforeafter,width=7.1cm,height=2.1cm,valign=center,colback=cvec!6,colframe=cvec,arc=2.5mm,boxrule=0.9pt,left=3mm]{\Large\bfseries\color{cvec}b)}\quad\begin{minipage}[c]{5.3cm}\raggedright\Large $(-1\,;\,3)$\end{minipage}\end{tcolorbox}\par\vspace{2.5mm}
\begin{tcolorbox}[enhanced,nobeforeafter,width=7.1cm,height=2.1cm,valign=center,colback=cvec!6,colframe=cvec,arc=2.5mm,boxrule=0.9pt,left=3mm]{\Large\bfseries\color{cvec}c)}\quad\begin{minipage}[c]{5.3cm}\raggedright\Large $(1\,;\,-3)$\end{minipage}\end{tcolorbox}\hfill\begin{tcolorbox}[enhanced,nobeforeafter,width=7.1cm,height=2.1cm,valign=center,colback=cvec!6,colframe=cvec,arc=2.5mm,boxrule=0.9pt,left=3mm]{\Large\bfseries\color{cvec}d)}\quad\begin{minipage}[c]{5.3cm}\raggedright\Large $(1\,;\,3)$\end{minipage}\end{tcolorbox}
\end{tcolorbox}
\newpage
\begin{tcolorbox}[enhanced,width=16cm,height=14cm,colback=white,colframe=cari,boxrule=2pt,arc=0mm,
left=6mm,right=6mm,top=0mm,bottom=4mm,sharp corners,
overlay={\fill[cari] (frame.north west) rectangle ([yshift=-1.25cm]frame.north east);
\node[anchor=west,text=white,font=\sffamily\bfseries\LARGE] at ([xshift=5mm,yshift=-0.62cm]frame.north west) {Question 65 C};
\node[anchor=east,text=white,font=\sffamily\itshape\large] at ([xshift=-5mm,yshift=-0.62cm]frame.north east) {Arithmétique et ensembles de nombres};}]
\vspace*{1.75cm}
\begin{minipage}[t][6.9cm][c]{\linewidth}\LARGE\centering
\begin{minipage}{0.95\linewidth}\raggedright\linespread{1.25}\selectfont Lequel de ces nombres est premier~?\end{minipage}
\end{minipage}

\vfill
\begin{tcolorbox}[enhanced,nobeforeafter,width=7.1cm,height=2.1cm,valign=center,colback=cari!6,colframe=cari,arc=2.5mm,boxrule=0.9pt,left=3mm]{\Large\bfseries\color{cari}a)}\quad\begin{minipage}[c]{5.3cm}\raggedright\Large $111$\end{minipage}\end{tcolorbox}\hfill\begin{tcolorbox}[enhanced,nobeforeafter,width=7.1cm,height=2.1cm,valign=center,colback=cari!6,colframe=cari,arc=2.5mm,boxrule=0.9pt,left=3mm]{\Large\bfseries\color{cari}b)}\quad\begin{minipage}[c]{5.3cm}\raggedright\Large $101$\end{minipage}\end{tcolorbox}\par\vspace{2.5mm}
\begin{tcolorbox}[enhanced,nobeforeafter,width=7.1cm,height=2.1cm,valign=center,colback=cari!6,colframe=cari,arc=2.5mm,boxrule=0.9pt,left=3mm]{\Large\bfseries\color{cari}c)}\quad\begin{minipage}[c]{5.3cm}\raggedright\Large $91$\end{minipage}\end{tcolorbox}\hfill\begin{tcolorbox}[enhanced,nobeforeafter,width=7.1cm,height=2.1cm,valign=center,colback=cari!6,colframe=cari,arc=2.5mm,boxrule=0.9pt,left=3mm]{\Large\bfseries\color{cari}d)}\quad\begin{minipage}[c]{5.3cm}\raggedright\Large $121$\end{minipage}\end{tcolorbox}
\end{tcolorbox}
\newpage
\begin{tcolorbox}[enhanced,width=16cm,height=14cm,colback=white,colframe=cari,boxrule=2pt,arc=0mm,
left=6mm,right=6mm,top=0mm,bottom=4mm,sharp corners,
overlay={\fill[cari] (frame.north west) rectangle ([yshift=-1.25cm]frame.north east);
\node[anchor=west,text=white,font=\sffamily\bfseries\LARGE] at ([xshift=5mm,yshift=-0.62cm]frame.north west) {Question 66 C};
\node[anchor=east,text=white,font=\sffamily\itshape\large] at ([xshift=-5mm,yshift=-0.62cm]frame.north east) {Arithmétique et ensembles de nombres};}]
\vspace*{1.75cm}
\begin{minipage}[t][6.9cm][c]{\linewidth}\LARGE\centering
\begin{minipage}{0.95\linewidth}\raggedright\linespread{1.25}\selectfont Lequel de ces nombres est premier~?\end{minipage}
\end{minipage}

\vfill
\begin{tcolorbox}[enhanced,nobeforeafter,width=7.1cm,height=2.1cm,valign=center,colback=cari!6,colframe=cari,arc=2.5mm,boxrule=0.9pt,left=3mm]{\Large\bfseries\color{cari}a)}\quad\begin{minipage}[c]{5.3cm}\raggedright\Large $69$\end{minipage}\end{tcolorbox}\hfill\begin{tcolorbox}[enhanced,nobeforeafter,width=7.1cm,height=2.1cm,valign=center,colback=cari!6,colframe=cari,arc=2.5mm,boxrule=0.9pt,left=3mm]{\Large\bfseries\color{cari}b)}\quad\begin{minipage}[c]{5.3cm}\raggedright\Large $81$\end{minipage}\end{tcolorbox}\par\vspace{2.5mm}
\begin{tcolorbox}[enhanced,nobeforeafter,width=7.1cm,height=2.1cm,valign=center,colback=cari!6,colframe=cari,arc=2.5mm,boxrule=0.9pt,left=3mm]{\Large\bfseries\color{cari}c)}\quad\begin{minipage}[c]{5.3cm}\raggedright\Large $73$\end{minipage}\end{tcolorbox}\hfill\begin{tcolorbox}[enhanced,nobeforeafter,width=7.1cm,height=2.1cm,valign=center,colback=cari!6,colframe=cari,arc=2.5mm,boxrule=0.9pt,left=3mm]{\Large\bfseries\color{cari}d)}\quad\begin{minipage}[c]{5.3cm}\raggedright\Large $77$\end{minipage}\end{tcolorbox}
\end{tcolorbox}
\newpage
\begin{tcolorbox}[enhanced,width=16cm,height=14cm,colback=white,colframe=cari,boxrule=2pt,arc=0mm,
left=6mm,right=6mm,top=0mm,bottom=4mm,sharp corners,
overlay={\fill[cari] (frame.north west) rectangle ([yshift=-1.25cm]frame.north east);
\node[anchor=west,text=white,font=\sffamily\bfseries\LARGE] at ([xshift=5mm,yshift=-0.62cm]frame.north west) {Question 67 C};
\node[anchor=east,text=white,font=\sffamily\itshape\large] at ([xshift=-5mm,yshift=-0.62cm]frame.north east) {Arithmétique et ensembles de nombres};}]
\vspace*{1.75cm}
\begin{minipage}[t][6.9cm][c]{\linewidth}\LARGE\centering
\begin{minipage}{0.95\linewidth}\raggedright\linespread{1.25}\selectfont Quelle est la décomposition de $252$ en produit de facteurs premiers~?\end{minipage}
\end{minipage}

\vfill
\begin{tcolorbox}[enhanced,nobeforeafter,width=7.1cm,height=2.1cm,valign=center,colback=cari!6,colframe=cari,arc=2.5mm,boxrule=0.9pt,left=3mm]{\Large\bfseries\color{cari}a)}\quad\begin{minipage}[c]{5.3cm}\raggedright\Large $2^{2}\times3^{2}\times7$\end{minipage}\end{tcolorbox}\hfill\begin{tcolorbox}[enhanced,nobeforeafter,width=7.1cm,height=2.1cm,valign=center,colback=cari!6,colframe=cari,arc=2.5mm,boxrule=0.9pt,left=3mm]{\Large\bfseries\color{cari}b)}\quad\begin{minipage}[c]{5.3cm}\raggedright\Large $2^{3}\times3^{2}\times7$\end{minipage}\end{tcolorbox}\par\vspace{2.5mm}
\begin{tcolorbox}[enhanced,nobeforeafter,width=7.1cm,height=2.1cm,valign=center,colback=cari!6,colframe=cari,arc=2.5mm,boxrule=0.9pt,left=3mm]{\Large\bfseries\color{cari}c)}\quad\begin{minipage}[c]{5.3cm}\raggedright\Large $12\times21$\end{minipage}\end{tcolorbox}\hfill\begin{tcolorbox}[enhanced,nobeforeafter,width=7.1cm,height=2.1cm,valign=center,colback=cari!6,colframe=cari,arc=2.5mm,boxrule=0.9pt,left=3mm]{\Large\bfseries\color{cari}d)}\quad\begin{minipage}[c]{5.3cm}\raggedright\large $2^{2}\times3^{2}\times7^{2}$\end{minipage}\end{tcolorbox}
\end{tcolorbox}
\newpage
\begin{tcolorbox}[enhanced,width=16cm,height=14cm,colback=white,colframe=cari,boxrule=2pt,arc=0mm,
left=6mm,right=6mm,top=0mm,bottom=4mm,sharp corners,
overlay={\fill[cari] (frame.north west) rectangle ([yshift=-1.25cm]frame.north east);
\node[anchor=west,text=white,font=\sffamily\bfseries\LARGE] at ([xshift=5mm,yshift=-0.62cm]frame.north west) {Question 68 C};
\node[anchor=east,text=white,font=\sffamily\itshape\large] at ([xshift=-5mm,yshift=-0.62cm]frame.north east) {Arithmétique et ensembles de nombres};}]
\vspace*{1.75cm}
\begin{minipage}[t][6.9cm][c]{\linewidth}\LARGE\centering
\begin{minipage}{0.95\linewidth}\raggedright\linespread{1.25}\selectfont Quelle est la décomposition de $90$ en produit de facteurs premiers~?\end{minipage}
\end{minipage}

\vfill
\begin{tcolorbox}[enhanced,nobeforeafter,width=7.1cm,height=2.1cm,valign=center,colback=cari!6,colframe=cari,arc=2.5mm,boxrule=0.9pt,left=3mm]{\Large\bfseries\color{cari}a)}\quad\begin{minipage}[c]{5.3cm}\raggedright\Large $6\times15$\end{minipage}\end{tcolorbox}\hfill\begin{tcolorbox}[enhanced,nobeforeafter,width=7.1cm,height=2.1cm,valign=center,colback=cari!6,colframe=cari,arc=2.5mm,boxrule=0.9pt,left=3mm]{\Large\bfseries\color{cari}b)}\quad\begin{minipage}[c]{5.3cm}\raggedright\Large $2\times3^{2}\times5^{2}$\end{minipage}\end{tcolorbox}\par\vspace{2.5mm}
\begin{tcolorbox}[enhanced,nobeforeafter,width=7.1cm,height=2.1cm,valign=center,colback=cari!6,colframe=cari,arc=2.5mm,boxrule=0.9pt,left=3mm]{\Large\bfseries\color{cari}c)}\quad\begin{minipage}[c]{5.3cm}\raggedright\Large $2\times3^{2}\times5$\end{minipage}\end{tcolorbox}\hfill\begin{tcolorbox}[enhanced,nobeforeafter,width=7.1cm,height=2.1cm,valign=center,colback=cari!6,colframe=cari,arc=2.5mm,boxrule=0.9pt,left=3mm]{\Large\bfseries\color{cari}d)}\quad\begin{minipage}[c]{5.3cm}\raggedright\Large $2^{2}\times3^{2}\times5$\end{minipage}\end{tcolorbox}
\end{tcolorbox}
\newpage
\begin{tcolorbox}[enhanced,width=16cm,height=14cm,colback=white,colframe=cari,boxrule=2pt,arc=0mm,
left=6mm,right=6mm,top=0mm,bottom=4mm,sharp corners,
overlay={\fill[cari] (frame.north west) rectangle ([yshift=-1.25cm]frame.north east);
\node[anchor=west,text=white,font=\sffamily\bfseries\LARGE] at ([xshift=5mm,yshift=-0.62cm]frame.north west) {Question 69 C};
\node[anchor=east,text=white,font=\sffamily\itshape\large] at ([xshift=-5mm,yshift=-0.62cm]frame.north east) {Arithmétique et ensembles de nombres};}]
\vspace*{1.75cm}
\begin{minipage}[t][6.9cm][c]{\linewidth}\LARGE\centering
\begin{minipage}{0.95\linewidth}\raggedright\linespread{1.25}\selectfont Quel est le plus petit ensemble de nombres auquel appartient $\sqrt{49}$~?\end{minipage}
\end{minipage}

\vfill
\begin{tcolorbox}[enhanced,nobeforeafter,width=7.1cm,height=2.1cm,valign=center,colback=cari!6,colframe=cari,arc=2.5mm,boxrule=0.9pt,left=3mm]{\Large\bfseries\color{cari}a)}\quad\begin{minipage}[c]{5.3cm}\raggedright\Large $\mathbb{N}$\end{minipage}\end{tcolorbox}\hfill\begin{tcolorbox}[enhanced,nobeforeafter,width=7.1cm,height=2.1cm,valign=center,colback=cari!6,colframe=cari,arc=2.5mm,boxrule=0.9pt,left=3mm]{\Large\bfseries\color{cari}b)}\quad\begin{minipage}[c]{5.3cm}\raggedright\Large $\mathbb{Z}$\end{minipage}\end{tcolorbox}\par\vspace{2.5mm}
\begin{tcolorbox}[enhanced,nobeforeafter,width=7.1cm,height=2.1cm,valign=center,colback=cari!6,colframe=cari,arc=2.5mm,boxrule=0.9pt,left=3mm]{\Large\bfseries\color{cari}c)}\quad\begin{minipage}[c]{5.3cm}\raggedright\Large $\mathbb{D}$\end{minipage}\end{tcolorbox}\hfill\begin{tcolorbox}[enhanced,nobeforeafter,width=7.1cm,height=2.1cm,valign=center,colback=cari!6,colframe=cari,arc=2.5mm,boxrule=0.9pt,left=3mm]{\Large\bfseries\color{cari}d)}\quad\begin{minipage}[c]{5.3cm}\raggedright\Large $\mathbb{Q}$\end{minipage}\end{tcolorbox}
\end{tcolorbox}
\newpage
\begin{tcolorbox}[enhanced,width=16cm,height=14cm,colback=white,colframe=cari,boxrule=2pt,arc=0mm,
left=6mm,right=6mm,top=0mm,bottom=4mm,sharp corners,
overlay={\fill[cari] (frame.north west) rectangle ([yshift=-1.25cm]frame.north east);
\node[anchor=west,text=white,font=\sffamily\bfseries\LARGE] at ([xshift=5mm,yshift=-0.62cm]frame.north west) {Question 70 C};
\node[anchor=east,text=white,font=\sffamily\itshape\large] at ([xshift=-5mm,yshift=-0.62cm]frame.north east) {Arithmétique et ensembles de nombres};}]
\vspace*{1.75cm}
\begin{minipage}[t][6.9cm][c]{\linewidth}\LARGE\centering
\begin{minipage}{0.95\linewidth}\raggedright\linespread{1.25}\selectfont Quel est le plus petit ensemble de nombres auquel appartient $-2{,}5$~?\end{minipage}
\end{minipage}

\vfill
\begin{tcolorbox}[enhanced,nobeforeafter,width=7.1cm,height=2.1cm,valign=center,colback=cari!6,colframe=cari,arc=2.5mm,boxrule=0.9pt,left=3mm]{\Large\bfseries\color{cari}a)}\quad\begin{minipage}[c]{5.3cm}\raggedright\Large $\mathbb{Q}$\end{minipage}\end{tcolorbox}\hfill\begin{tcolorbox}[enhanced,nobeforeafter,width=7.1cm,height=2.1cm,valign=center,colback=cari!6,colframe=cari,arc=2.5mm,boxrule=0.9pt,left=3mm]{\Large\bfseries\color{cari}b)}\quad\begin{minipage}[c]{5.3cm}\raggedright\Large $\mathbb{Z}$\end{minipage}\end{tcolorbox}\par\vspace{2.5mm}
\begin{tcolorbox}[enhanced,nobeforeafter,width=7.1cm,height=2.1cm,valign=center,colback=cari!6,colframe=cari,arc=2.5mm,boxrule=0.9pt,left=3mm]{\Large\bfseries\color{cari}c)}\quad\begin{minipage}[c]{5.3cm}\raggedright\Large $\mathbb{N}$\end{minipage}\end{tcolorbox}\hfill\begin{tcolorbox}[enhanced,nobeforeafter,width=7.1cm,height=2.1cm,valign=center,colback=cari!6,colframe=cari,arc=2.5mm,boxrule=0.9pt,left=3mm]{\Large\bfseries\color{cari}d)}\quad\begin{minipage}[c]{5.3cm}\raggedright\Large $\mathbb{D}$\end{minipage}\end{tcolorbox}
\end{tcolorbox}
\newpage
\begin{tcolorbox}[enhanced,width=16cm,height=14cm,colback=white,colframe=cari,boxrule=2pt,arc=0mm,
left=6mm,right=6mm,top=0mm,bottom=4mm,sharp corners,
overlay={\fill[cari] (frame.north west) rectangle ([yshift=-1.25cm]frame.north east);
\node[anchor=west,text=white,font=\sffamily\bfseries\LARGE] at ([xshift=5mm,yshift=-0.62cm]frame.north west) {Question 71 C};
\node[anchor=east,text=white,font=\sffamily\itshape\large] at ([xshift=-5mm,yshift=-0.62cm]frame.north east) {Arithmétique et ensembles de nombres};}]
\vspace*{1.75cm}
\begin{minipage}[t][6.9cm][c]{\linewidth}\LARGE\centering
\begin{minipage}{0.95\linewidth}\raggedright\linespread{1.25}\selectfont Quelle est la forme irréductible de $\dfrac{36}{120}$~?\end{minipage}
\end{minipage}

\vfill
\begin{tcolorbox}[enhanced,nobeforeafter,width=7.1cm,height=2.1cm,valign=center,colback=cari!6,colframe=cari,arc=2.5mm,boxrule=0.9pt,left=3mm]{\Large\bfseries\color{cari}a)}\quad\begin{minipage}[c]{5.3cm}\raggedright\Large $\dfrac{18}{60}$\end{minipage}\end{tcolorbox}\hfill\begin{tcolorbox}[enhanced,nobeforeafter,width=7.1cm,height=2.1cm,valign=center,colback=cari!6,colframe=cari,arc=2.5mm,boxrule=0.9pt,left=3mm]{\Large\bfseries\color{cari}b)}\quad\begin{minipage}[c]{5.3cm}\raggedright\Large $\dfrac{10}{3}$\end{minipage}\end{tcolorbox}\par\vspace{2.5mm}
\begin{tcolorbox}[enhanced,nobeforeafter,width=7.1cm,height=2.1cm,valign=center,colback=cari!6,colframe=cari,arc=2.5mm,boxrule=0.9pt,left=3mm]{\Large\bfseries\color{cari}c)}\quad\begin{minipage}[c]{5.3cm}\raggedright\Large $\dfrac{3}{10}$\end{minipage}\end{tcolorbox}\hfill\begin{tcolorbox}[enhanced,nobeforeafter,width=7.1cm,height=2.1cm,valign=center,colback=cari!6,colframe=cari,arc=2.5mm,boxrule=0.9pt,left=3mm]{\Large\bfseries\color{cari}d)}\quad\begin{minipage}[c]{5.3cm}\raggedright\Large $\dfrac{12}{40}$\end{minipage}\end{tcolorbox}
\end{tcolorbox}
\newpage
\begin{tcolorbox}[enhanced,width=16cm,height=14cm,colback=white,colframe=cari,boxrule=2pt,arc=0mm,
left=6mm,right=6mm,top=0mm,bottom=4mm,sharp corners,
overlay={\fill[cari] (frame.north west) rectangle ([yshift=-1.25cm]frame.north east);
\node[anchor=west,text=white,font=\sffamily\bfseries\LARGE] at ([xshift=5mm,yshift=-0.62cm]frame.north west) {Question 72 C};
\node[anchor=east,text=white,font=\sffamily\itshape\large] at ([xshift=-5mm,yshift=-0.62cm]frame.north east) {Arithmétique et ensembles de nombres};}]
\vspace*{1.75cm}
\begin{minipage}[t][6.9cm][c]{\linewidth}\LARGE\centering
\begin{minipage}{0.95\linewidth}\raggedright\linespread{1.25}\selectfont Quelle est la forme irréductible de $\dfrac{150}{210}$~?\end{minipage}
\end{minipage}

\vfill
\begin{tcolorbox}[enhanced,nobeforeafter,width=7.1cm,height=2.1cm,valign=center,colback=cari!6,colframe=cari,arc=2.5mm,boxrule=0.9pt,left=3mm]{\Large\bfseries\color{cari}a)}\quad\begin{minipage}[c]{5.3cm}\raggedright\Large $\dfrac{30}{42}$\end{minipage}\end{tcolorbox}\hfill\begin{tcolorbox}[enhanced,nobeforeafter,width=7.1cm,height=2.1cm,valign=center,colback=cari!6,colframe=cari,arc=2.5mm,boxrule=0.9pt,left=3mm]{\Large\bfseries\color{cari}b)}\quad\begin{minipage}[c]{5.3cm}\raggedright\Large $\dfrac{5}{7}$\end{minipage}\end{tcolorbox}\par\vspace{2.5mm}
\begin{tcolorbox}[enhanced,nobeforeafter,width=7.1cm,height=2.1cm,valign=center,colback=cari!6,colframe=cari,arc=2.5mm,boxrule=0.9pt,left=3mm]{\Large\bfseries\color{cari}c)}\quad\begin{minipage}[c]{5.3cm}\raggedright\Large $\dfrac{75}{105}$\end{minipage}\end{tcolorbox}\hfill\begin{tcolorbox}[enhanced,nobeforeafter,width=7.1cm,height=2.1cm,valign=center,colback=cari!6,colframe=cari,arc=2.5mm,boxrule=0.9pt,left=3mm]{\Large\bfseries\color{cari}d)}\quad\begin{minipage}[c]{5.3cm}\raggedright\Large $\dfrac{50}{70}$\end{minipage}\end{tcolorbox}
\end{tcolorbox}
\newpage
\begin{tcolorbox}[enhanced,width=16cm,height=14cm,colback=white,colframe=cari,boxrule=2pt,arc=0mm,
left=6mm,right=6mm,top=0mm,bottom=4mm,sharp corners,
overlay={\fill[cari] (frame.north west) rectangle ([yshift=-1.25cm]frame.north east);
\node[anchor=west,text=white,font=\sffamily\bfseries\LARGE] at ([xshift=5mm,yshift=-0.62cm]frame.north west) {Question 73 C};
\node[anchor=east,text=white,font=\sffamily\itshape\large] at ([xshift=-5mm,yshift=-0.62cm]frame.north east) {Arithmétique et ensembles de nombres};}]
\vspace*{1.75cm}
\begin{minipage}[t][6.9cm][c]{\linewidth}\LARGE\centering
\begin{minipage}{0.95\linewidth}\raggedright\linespread{1.25}\selectfont Pour tout entier naturel $n$, que peut-on dire de l'entier $4n+2$~?\end{minipage}
\end{minipage}

\vfill
\begin{tcolorbox}[enhanced,nobeforeafter,width=7.1cm,height=2.1cm,valign=center,colback=cari!6,colframe=cari,arc=2.5mm,boxrule=0.9pt,left=3mm]{\Large\bfseries\color{cari}a)}\quad\begin{minipage}[c]{5.3cm}\raggedright\large il est pair ou impair selon $n$\end{minipage}\end{tcolorbox}\hfill\begin{tcolorbox}[enhanced,nobeforeafter,width=7.1cm,height=2.1cm,valign=center,colback=cari!6,colframe=cari,arc=2.5mm,boxrule=0.9pt,left=3mm]{\Large\bfseries\color{cari}b)}\quad\begin{minipage}[c]{5.3cm}\raggedright\Large il est toujours pair\end{minipage}\end{tcolorbox}\par\vspace{2.5mm}
\begin{tcolorbox}[enhanced,nobeforeafter,width=7.1cm,height=2.1cm,valign=center,colback=cari!6,colframe=cari,arc=2.5mm,boxrule=0.9pt,left=3mm]{\Large\bfseries\color{cari}c)}\quad\begin{minipage}[c]{5.3cm}\raggedright\Large il est toujours premier\end{minipage}\end{tcolorbox}\hfill\begin{tcolorbox}[enhanced,nobeforeafter,width=7.1cm,height=2.1cm,valign=center,colback=cari!6,colframe=cari,arc=2.5mm,boxrule=0.9pt,left=3mm]{\Large\bfseries\color{cari}d)}\quad\begin{minipage}[c]{5.3cm}\raggedright\Large il est toujours impair\end{minipage}\end{tcolorbox}
\end{tcolorbox}
\newpage
\begin{tcolorbox}[enhanced,width=16cm,height=14cm,colback=white,colframe=cari,boxrule=2pt,arc=0mm,
left=6mm,right=6mm,top=0mm,bottom=4mm,sharp corners,
overlay={\fill[cari] (frame.north west) rectangle ([yshift=-1.25cm]frame.north east);
\node[anchor=west,text=white,font=\sffamily\bfseries\LARGE] at ([xshift=5mm,yshift=-0.62cm]frame.north west) {Question 74 C};
\node[anchor=east,text=white,font=\sffamily\itshape\large] at ([xshift=-5mm,yshift=-0.62cm]frame.north east) {Arithmétique et ensembles de nombres};}]
\vspace*{1.75cm}
\begin{minipage}[t][6.9cm][c]{\linewidth}\LARGE\centering
\begin{minipage}{0.95\linewidth}\raggedright\linespread{1.25}\selectfont Pour tout entier naturel $n$, que peut-on dire de l'entier $n^2+n+1$~?\end{minipage}
\end{minipage}

\vfill
\begin{tcolorbox}[enhanced,nobeforeafter,width=7.1cm,height=2.1cm,valign=center,colback=cari!6,colframe=cari,arc=2.5mm,boxrule=0.9pt,left=3mm]{\Large\bfseries\color{cari}a)}\quad\begin{minipage}[c]{5.3cm}\raggedright\Large il est toujours impair\end{minipage}\end{tcolorbox}\hfill\begin{tcolorbox}[enhanced,nobeforeafter,width=7.1cm,height=2.1cm,valign=center,colback=cari!6,colframe=cari,arc=2.5mm,boxrule=0.9pt,left=3mm]{\Large\bfseries\color{cari}b)}\quad\begin{minipage}[c]{5.3cm}\raggedright\Large il est toujours premier\end{minipage}\end{tcolorbox}\par\vspace{2.5mm}
\begin{tcolorbox}[enhanced,nobeforeafter,width=7.1cm,height=2.1cm,valign=center,colback=cari!6,colframe=cari,arc=2.5mm,boxrule=0.9pt,left=3mm]{\Large\bfseries\color{cari}c)}\quad\begin{minipage}[c]{5.3cm}\raggedright\large il est pair ou impair selon $n$\end{minipage}\end{tcolorbox}\hfill\begin{tcolorbox}[enhanced,nobeforeafter,width=7.1cm,height=2.1cm,valign=center,colback=cari!6,colframe=cari,arc=2.5mm,boxrule=0.9pt,left=3mm]{\Large\bfseries\color{cari}d)}\quad\begin{minipage}[c]{5.3cm}\raggedright\Large il est toujours pair\end{minipage}\end{tcolorbox}
\end{tcolorbox}
\newpage
\begin{tcolorbox}[enhanced,width=16cm,height=14cm,colback=white,colframe=cari,boxrule=2pt,arc=0mm,
left=6mm,right=6mm,top=0mm,bottom=4mm,sharp corners,
overlay={\fill[cari] (frame.north west) rectangle ([yshift=-1.25cm]frame.north east);
\node[anchor=west,text=white,font=\sffamily\bfseries\LARGE] at ([xshift=5mm,yshift=-0.62cm]frame.north west) {Question 75 C};
\node[anchor=east,text=white,font=\sffamily\itshape\large] at ([xshift=-5mm,yshift=-0.62cm]frame.north east) {Arithmétique et ensembles de nombres};}]
\vspace*{1.75cm}
\begin{minipage}[t][6.9cm][c]{\linewidth}\LARGE\centering
\begin{minipage}{0.95\linewidth}\raggedright\linespread{1.25}\selectfont Combien y a-t-il de multiples de $11$ compris entre $1$ et $500$~?\end{minipage}
\end{minipage}

\vfill
\begin{tcolorbox}[enhanced,nobeforeafter,width=7.1cm,height=2.1cm,valign=center,colback=cari!6,colframe=cari,arc=2.5mm,boxrule=0.9pt,left=3mm]{\Large\bfseries\color{cari}a)}\quad\begin{minipage}[c]{5.3cm}\raggedright\Large $45$\end{minipage}\end{tcolorbox}\hfill\begin{tcolorbox}[enhanced,nobeforeafter,width=7.1cm,height=2.1cm,valign=center,colback=cari!6,colframe=cari,arc=2.5mm,boxrule=0.9pt,left=3mm]{\Large\bfseries\color{cari}b)}\quad\begin{minipage}[c]{5.3cm}\raggedright\Large $455$\end{minipage}\end{tcolorbox}\par\vspace{2.5mm}
\begin{tcolorbox}[enhanced,nobeforeafter,width=7.1cm,height=2.1cm,valign=center,colback=cari!6,colframe=cari,arc=2.5mm,boxrule=0.9pt,left=3mm]{\Large\bfseries\color{cari}c)}\quad\begin{minipage}[c]{5.3cm}\raggedright\Large $44$\end{minipage}\end{tcolorbox}\hfill\begin{tcolorbox}[enhanced,nobeforeafter,width=7.1cm,height=2.1cm,valign=center,colback=cari!6,colframe=cari,arc=2.5mm,boxrule=0.9pt,left=3mm]{\Large\bfseries\color{cari}d)}\quad\begin{minipage}[c]{5.3cm}\raggedright\Large $46$\end{minipage}\end{tcolorbox}
\end{tcolorbox}
\newpage
\begin{tcolorbox}[enhanced,width=16cm,height=14cm,colback=white,colframe=cari,boxrule=2pt,arc=0mm,
left=6mm,right=6mm,top=0mm,bottom=4mm,sharp corners,
overlay={\fill[cari] (frame.north west) rectangle ([yshift=-1.25cm]frame.north east);
\node[anchor=west,text=white,font=\sffamily\bfseries\LARGE] at ([xshift=5mm,yshift=-0.62cm]frame.north west) {Question 76 C};
\node[anchor=east,text=white,font=\sffamily\itshape\large] at ([xshift=-5mm,yshift=-0.62cm]frame.north east) {Arithmétique et ensembles de nombres};}]
\vspace*{1.75cm}
\begin{minipage}[t][6.9cm][c]{\linewidth}\LARGE\centering
\begin{minipage}{0.95\linewidth}\raggedright\linespread{1.25}\selectfont Combien y a-t-il de multiples de $4$ compris entre $1$ et $99$~?\end{minipage}
\end{minipage}

\vfill
\begin{tcolorbox}[enhanced,nobeforeafter,width=7.1cm,height=2.1cm,valign=center,colback=cari!6,colframe=cari,arc=2.5mm,boxrule=0.9pt,left=3mm]{\Large\bfseries\color{cari}a)}\quad\begin{minipage}[c]{5.3cm}\raggedright\Large $25$\end{minipage}\end{tcolorbox}\hfill\begin{tcolorbox}[enhanced,nobeforeafter,width=7.1cm,height=2.1cm,valign=center,colback=cari!6,colframe=cari,arc=2.5mm,boxrule=0.9pt,left=3mm]{\Large\bfseries\color{cari}b)}\quad\begin{minipage}[c]{5.3cm}\raggedright\Large $75$\end{minipage}\end{tcolorbox}\par\vspace{2.5mm}
\begin{tcolorbox}[enhanced,nobeforeafter,width=7.1cm,height=2.1cm,valign=center,colback=cari!6,colframe=cari,arc=2.5mm,boxrule=0.9pt,left=3mm]{\Large\bfseries\color{cari}c)}\quad\begin{minipage}[c]{5.3cm}\raggedright\Large $23$\end{minipage}\end{tcolorbox}\hfill\begin{tcolorbox}[enhanced,nobeforeafter,width=7.1cm,height=2.1cm,valign=center,colback=cari!6,colframe=cari,arc=2.5mm,boxrule=0.9pt,left=3mm]{\Large\bfseries\color{cari}d)}\quad\begin{minipage}[c]{5.3cm}\raggedright\Large $24$\end{minipage}\end{tcolorbox}
\end{tcolorbox}
\newpage
\begin{tcolorbox}[enhanced,width=16cm,height=14cm,colback=white,colframe=cari,boxrule=2pt,arc=0mm,
left=6mm,right=6mm,top=0mm,bottom=4mm,sharp corners,
overlay={\fill[cari] (frame.north west) rectangle ([yshift=-1.25cm]frame.north east);
\node[anchor=west,text=white,font=\sffamily\bfseries\LARGE] at ([xshift=5mm,yshift=-0.62cm]frame.north west) {Question 77 C};
\node[anchor=east,text=white,font=\sffamily\itshape\large] at ([xshift=-5mm,yshift=-0.62cm]frame.north east) {Arithmétique et ensembles de nombres};}]
\vspace*{1.75cm}
\begin{minipage}[t][6.9cm][c]{\linewidth}\LARGE\centering
\begin{minipage}{0.95\linewidth}\raggedright\linespread{1.25}\selectfont Combien l'entier $17$ a-t-il de diviseurs positifs~?\end{minipage}
\end{minipage}

\vfill
\begin{tcolorbox}[enhanced,nobeforeafter,width=7.1cm,height=2.1cm,valign=center,colback=cari!6,colframe=cari,arc=2.5mm,boxrule=0.9pt,left=3mm]{\Large\bfseries\color{cari}a)}\quad\begin{minipage}[c]{5.3cm}\raggedright\Large $2$\end{minipage}\end{tcolorbox}\hfill\begin{tcolorbox}[enhanced,nobeforeafter,width=7.1cm,height=2.1cm,valign=center,colback=cari!6,colframe=cari,arc=2.5mm,boxrule=0.9pt,left=3mm]{\Large\bfseries\color{cari}b)}\quad\begin{minipage}[c]{5.3cm}\raggedright\Large $4$\end{minipage}\end{tcolorbox}\par\vspace{2.5mm}
\begin{tcolorbox}[enhanced,nobeforeafter,width=7.1cm,height=2.1cm,valign=center,colback=cari!6,colframe=cari,arc=2.5mm,boxrule=0.9pt,left=3mm]{\Large\bfseries\color{cari}c)}\quad\begin{minipage}[c]{5.3cm}\raggedright\Large $1$\end{minipage}\end{tcolorbox}\hfill\begin{tcolorbox}[enhanced,nobeforeafter,width=7.1cm,height=2.1cm,valign=center,colback=cari!6,colframe=cari,arc=2.5mm,boxrule=0.9pt,left=3mm]{\Large\bfseries\color{cari}d)}\quad\begin{minipage}[c]{5.3cm}\raggedright\Large $3$\end{minipage}\end{tcolorbox}
\end{tcolorbox}
\newpage
\begin{tcolorbox}[enhanced,width=16cm,height=14cm,colback=white,colframe=cari,boxrule=2pt,arc=0mm,
left=6mm,right=6mm,top=0mm,bottom=4mm,sharp corners,
overlay={\fill[cari] (frame.north west) rectangle ([yshift=-1.25cm]frame.north east);
\node[anchor=west,text=white,font=\sffamily\bfseries\LARGE] at ([xshift=5mm,yshift=-0.62cm]frame.north west) {Question 78 C};
\node[anchor=east,text=white,font=\sffamily\itshape\large] at ([xshift=-5mm,yshift=-0.62cm]frame.north east) {Arithmétique et ensembles de nombres};}]
\vspace*{1.75cm}
\begin{minipage}[t][6.9cm][c]{\linewidth}\LARGE\centering
\begin{minipage}{0.95\linewidth}\raggedright\linespread{1.25}\selectfont Combien l'entier $64$ a-t-il de diviseurs positifs~?\end{minipage}
\end{minipage}

\vfill
\begin{tcolorbox}[enhanced,nobeforeafter,width=7.1cm,height=2.1cm,valign=center,colback=cari!6,colframe=cari,arc=2.5mm,boxrule=0.9pt,left=3mm]{\Large\bfseries\color{cari}a)}\quad\begin{minipage}[c]{5.3cm}\raggedright\Large $6$\end{minipage}\end{tcolorbox}\hfill\begin{tcolorbox}[enhanced,nobeforeafter,width=7.1cm,height=2.1cm,valign=center,colback=cari!6,colframe=cari,arc=2.5mm,boxrule=0.9pt,left=3mm]{\Large\bfseries\color{cari}b)}\quad\begin{minipage}[c]{5.3cm}\raggedright\Large $8$\end{minipage}\end{tcolorbox}\par\vspace{2.5mm}
\begin{tcolorbox}[enhanced,nobeforeafter,width=7.1cm,height=2.1cm,valign=center,colback=cari!6,colframe=cari,arc=2.5mm,boxrule=0.9pt,left=3mm]{\Large\bfseries\color{cari}c)}\quad\begin{minipage}[c]{5.3cm}\raggedright\Large $9$\end{minipage}\end{tcolorbox}\hfill\begin{tcolorbox}[enhanced,nobeforeafter,width=7.1cm,height=2.1cm,valign=center,colback=cari!6,colframe=cari,arc=2.5mm,boxrule=0.9pt,left=3mm]{\Large\bfseries\color{cari}d)}\quad\begin{minipage}[c]{5.3cm}\raggedright\Large $7$\end{minipage}\end{tcolorbox}
\end{tcolorbox}
\newpage
\begin{tcolorbox}[enhanced,width=16cm,height=14cm,colback=white,colframe=cari,boxrule=2pt,arc=0mm,
left=6mm,right=6mm,top=0mm,bottom=4mm,sharp corners,
overlay={\fill[cari] (frame.north west) rectangle ([yshift=-1.25cm]frame.north east);
\node[anchor=west,text=white,font=\sffamily\bfseries\LARGE] at ([xshift=5mm,yshift=-0.62cm]frame.north west) {Question 79 C};
\node[anchor=east,text=white,font=\sffamily\itshape\large] at ([xshift=-5mm,yshift=-0.62cm]frame.north east) {Arithmétique et ensembles de nombres};}]
\vspace*{1.75cm}
\begin{minipage}[t][6.9cm][c]{\linewidth}\LARGE\centering
\begin{minipage}{0.95\linewidth}\raggedright\linespread{1.25}\selectfont Que vaut l'intersection $[1\,;\,7]\cap[3\,;\,4]$~?\end{minipage}
\end{minipage}

\vfill
\begin{tcolorbox}[enhanced,nobeforeafter,width=7.1cm,height=2.1cm,valign=center,colback=cari!6,colframe=cari,arc=2.5mm,boxrule=0.9pt,left=3mm]{\Large\bfseries\color{cari}a)}\quad\begin{minipage}[c]{5.3cm}\raggedright\Large $[4\,;\,7]$\end{minipage}\end{tcolorbox}\hfill\begin{tcolorbox}[enhanced,nobeforeafter,width=7.1cm,height=2.1cm,valign=center,colback=cari!6,colframe=cari,arc=2.5mm,boxrule=0.9pt,left=3mm]{\Large\bfseries\color{cari}b)}\quad\begin{minipage}[c]{5.3cm}\raggedright\Large $[1\,;\,3]$\end{minipage}\end{tcolorbox}\par\vspace{2.5mm}
\begin{tcolorbox}[enhanced,nobeforeafter,width=7.1cm,height=2.1cm,valign=center,colback=cari!6,colframe=cari,arc=2.5mm,boxrule=0.9pt,left=3mm]{\Large\bfseries\color{cari}c)}\quad\begin{minipage}[c]{5.3cm}\raggedright\Large $[3\,;\,4]$\end{minipage}\end{tcolorbox}\hfill\begin{tcolorbox}[enhanced,nobeforeafter,width=7.1cm,height=2.1cm,valign=center,colback=cari!6,colframe=cari,arc=2.5mm,boxrule=0.9pt,left=3mm]{\Large\bfseries\color{cari}d)}\quad\begin{minipage}[c]{5.3cm}\raggedright\Large $[1\,;\,7]$\end{minipage}\end{tcolorbox}
\end{tcolorbox}
\newpage
\begin{tcolorbox}[enhanced,width=16cm,height=14cm,colback=white,colframe=cari,boxrule=2pt,arc=0mm,
left=6mm,right=6mm,top=0mm,bottom=4mm,sharp corners,
overlay={\fill[cari] (frame.north west) rectangle ([yshift=-1.25cm]frame.north east);
\node[anchor=west,text=white,font=\sffamily\bfseries\LARGE] at ([xshift=5mm,yshift=-0.62cm]frame.north west) {Question 80 C};
\node[anchor=east,text=white,font=\sffamily\itshape\large] at ([xshift=-5mm,yshift=-0.62cm]frame.north east) {Arithmétique et ensembles de nombres};}]
\vspace*{1.75cm}
\begin{minipage}[t][6.9cm][c]{\linewidth}\LARGE\centering
\begin{minipage}{0.95\linewidth}\raggedright\linespread{1.25}\selectfont Que vaut l'intersection $]-2\,;\,6[\cap[-4\,;\,1]$~?\end{minipage}
\end{minipage}

\vfill
\begin{tcolorbox}[enhanced,nobeforeafter,width=7.1cm,height=2.1cm,valign=center,colback=cari!6,colframe=cari,arc=2.5mm,boxrule=0.9pt,left=3mm]{\Large\bfseries\color{cari}a)}\quad\begin{minipage}[c]{5.3cm}\raggedright\Large $[-2\,;\,1]$\end{minipage}\end{tcolorbox}\hfill\begin{tcolorbox}[enhanced,nobeforeafter,width=7.1cm,height=2.1cm,valign=center,colback=cari!6,colframe=cari,arc=2.5mm,boxrule=0.9pt,left=3mm]{\Large\bfseries\color{cari}b)}\quad\begin{minipage}[c]{5.3cm}\raggedright\Large $]-2\,;\,1[$\end{minipage}\end{tcolorbox}\par\vspace{2.5mm}
\begin{tcolorbox}[enhanced,nobeforeafter,width=7.1cm,height=2.1cm,valign=center,colback=cari!6,colframe=cari,arc=2.5mm,boxrule=0.9pt,left=3mm]{\Large\bfseries\color{cari}c)}\quad\begin{minipage}[c]{5.3cm}\raggedright\Large $[-4\,;\,6[$\end{minipage}\end{tcolorbox}\hfill\begin{tcolorbox}[enhanced,nobeforeafter,width=7.1cm,height=2.1cm,valign=center,colback=cari!6,colframe=cari,arc=2.5mm,boxrule=0.9pt,left=3mm]{\Large\bfseries\color{cari}d)}\quad\begin{minipage}[c]{5.3cm}\raggedright\Large $]-2\,;\,1]$\end{minipage}\end{tcolorbox}
\end{tcolorbox}
\newpage
\begin{tcolorbox}[enhanced,width=16cm,height=14cm,colback=white,colframe=cpro,boxrule=2pt,arc=0mm,
left=6mm,right=6mm,top=0mm,bottom=4mm,sharp corners,
overlay={\fill[cpro] (frame.north west) rectangle ([yshift=-1.25cm]frame.north east);
\node[anchor=west,text=white,font=\sffamily\bfseries\LARGE] at ([xshift=5mm,yshift=-0.62cm]frame.north west) {Question 81 C};
\node[anchor=east,text=white,font=\sffamily\itshape\large] at ([xshift=-5mm,yshift=-0.62cm]frame.north east) {Probabilités};}]
\vspace*{1.75cm}
\begin{minipage}[t][6.9cm][c]{\linewidth}\LARGE\centering
\begin{minipage}{0.95\linewidth}\raggedright\linespread{1.25}\selectfont On a $P(A)=0{,}7$, $P(B)=0{,}2$ et $P(A\cap B)=0{,}1$. Que vaut $P(A\cup B)$~?\end{minipage}
\end{minipage}

\vfill
\begin{tcolorbox}[enhanced,nobeforeafter,width=7.1cm,height=2.1cm,valign=center,colback=cpro!6,colframe=cpro,arc=2.5mm,boxrule=0.9pt,left=3mm]{\Large\bfseries\color{cpro}a)}\quad\begin{minipage}[c]{5.3cm}\raggedright\Large $0{,}9$\end{minipage}\end{tcolorbox}\hfill\begin{tcolorbox}[enhanced,nobeforeafter,width=7.1cm,height=2.1cm,valign=center,colback=cpro!6,colframe=cpro,arc=2.5mm,boxrule=0.9pt,left=3mm]{\Large\bfseries\color{cpro}b)}\quad\begin{minipage}[c]{5.3cm}\raggedright\Large $1$\end{minipage}\end{tcolorbox}\par\vspace{2.5mm}
\begin{tcolorbox}[enhanced,nobeforeafter,width=7.1cm,height=2.1cm,valign=center,colback=cpro!6,colframe=cpro,arc=2.5mm,boxrule=0.9pt,left=3mm]{\Large\bfseries\color{cpro}c)}\quad\begin{minipage}[c]{5.3cm}\raggedright\Large $0{,}1$\end{minipage}\end{tcolorbox}\hfill\begin{tcolorbox}[enhanced,nobeforeafter,width=7.1cm,height=2.1cm,valign=center,colback=cpro!6,colframe=cpro,arc=2.5mm,boxrule=0.9pt,left=3mm]{\Large\bfseries\color{cpro}d)}\quad\begin{minipage}[c]{5.3cm}\raggedright\Large $0{,}8$\end{minipage}\end{tcolorbox}
\end{tcolorbox}
\newpage
\begin{tcolorbox}[enhanced,width=16cm,height=14cm,colback=white,colframe=cpro,boxrule=2pt,arc=0mm,
left=6mm,right=6mm,top=0mm,bottom=4mm,sharp corners,
overlay={\fill[cpro] (frame.north west) rectangle ([yshift=-1.25cm]frame.north east);
\node[anchor=west,text=white,font=\sffamily\bfseries\LARGE] at ([xshift=5mm,yshift=-0.62cm]frame.north west) {Question 82 C};
\node[anchor=east,text=white,font=\sffamily\itshape\large] at ([xshift=-5mm,yshift=-0.62cm]frame.north east) {Probabilités};}]
\vspace*{1.75cm}
\begin{minipage}[t][6.9cm][c]{\linewidth}\LARGE\centering
\begin{minipage}{0.95\linewidth}\raggedright\linespread{1.25}\selectfont On a $P(A)=0{,}6$, $P(B)=0{,}6$ et $P(A\cap B)=0{,}45$. Que vaut $P(A\cup B)$~?\end{minipage}
\end{minipage}

\vfill
\begin{tcolorbox}[enhanced,nobeforeafter,width=7.1cm,height=2.1cm,valign=center,colback=cpro!6,colframe=cpro,arc=2.5mm,boxrule=0.9pt,left=3mm]{\Large\bfseries\color{cpro}a)}\quad\begin{minipage}[c]{5.3cm}\raggedright\Large $0{,}3$\end{minipage}\end{tcolorbox}\hfill\begin{tcolorbox}[enhanced,nobeforeafter,width=7.1cm,height=2.1cm,valign=center,colback=cpro!6,colframe=cpro,arc=2.5mm,boxrule=0.9pt,left=3mm]{\Large\bfseries\color{cpro}b)}\quad\begin{minipage}[c]{5.3cm}\raggedright\Large $0{,}36$\end{minipage}\end{tcolorbox}\par\vspace{2.5mm}
\begin{tcolorbox}[enhanced,nobeforeafter,width=7.1cm,height=2.1cm,valign=center,colback=cpro!6,colframe=cpro,arc=2.5mm,boxrule=0.9pt,left=3mm]{\Large\bfseries\color{cpro}c)}\quad\begin{minipage}[c]{5.3cm}\raggedright\Large $0{,}45$\end{minipage}\end{tcolorbox}\hfill\begin{tcolorbox}[enhanced,nobeforeafter,width=7.1cm,height=2.1cm,valign=center,colback=cpro!6,colframe=cpro,arc=2.5mm,boxrule=0.9pt,left=3mm]{\Large\bfseries\color{cpro}d)}\quad\begin{minipage}[c]{5.3cm}\raggedright\Large $0{,}75$\end{minipage}\end{tcolorbox}
\end{tcolorbox}
\newpage
\begin{tcolorbox}[enhanced,width=16cm,height=14cm,colback=white,colframe=cpro,boxrule=2pt,arc=0mm,
left=6mm,right=6mm,top=0mm,bottom=4mm,sharp corners,
overlay={\fill[cpro] (frame.north west) rectangle ([yshift=-1.25cm]frame.north east);
\node[anchor=west,text=white,font=\sffamily\bfseries\LARGE] at ([xshift=5mm,yshift=-0.62cm]frame.north west) {Question 83 C};
\node[anchor=east,text=white,font=\sffamily\itshape\large] at ([xshift=-5mm,yshift=-0.62cm]frame.north east) {Probabilités};}]
\vspace*{1.75cm}
\begin{minipage}[t][6.9cm][c]{\linewidth}\LARGE\centering
\begin{minipage}{0.95\linewidth}\raggedright\linespread{1.25}\selectfont On a $P(A)=0{,}2$, $P(B)=0{,}3$ et $P(A\cup B)=0{,}45$. Que vaut $P(A\cap B)$~?\end{minipage}
\end{minipage}

\vfill
\begin{tcolorbox}[enhanced,nobeforeafter,width=7.1cm,height=2.1cm,valign=center,colback=cpro!6,colframe=cpro,arc=2.5mm,boxrule=0.9pt,left=3mm]{\Large\bfseries\color{cpro}a)}\quad\begin{minipage}[c]{5.3cm}\raggedright\Large $0{,}15$\end{minipage}\end{tcolorbox}\hfill\begin{tcolorbox}[enhanced,nobeforeafter,width=7.1cm,height=2.1cm,valign=center,colback=cpro!6,colframe=cpro,arc=2.5mm,boxrule=0.9pt,left=3mm]{\Large\bfseries\color{cpro}b)}\quad\begin{minipage}[c]{5.3cm}\raggedright\Large $0{,}06$\end{minipage}\end{tcolorbox}\par\vspace{2.5mm}
\begin{tcolorbox}[enhanced,nobeforeafter,width=7.1cm,height=2.1cm,valign=center,colback=cpro!6,colframe=cpro,arc=2.5mm,boxrule=0.9pt,left=3mm]{\Large\bfseries\color{cpro}c)}\quad\begin{minipage}[c]{5.3cm}\raggedright\Large $0{,}05$\end{minipage}\end{tcolorbox}\hfill\begin{tcolorbox}[enhanced,nobeforeafter,width=7.1cm,height=2.1cm,valign=center,colback=cpro!6,colframe=cpro,arc=2.5mm,boxrule=0.9pt,left=3mm]{\Large\bfseries\color{cpro}d)}\quad\begin{minipage}[c]{5.3cm}\raggedright\Large $0{,}25$\end{minipage}\end{tcolorbox}
\end{tcolorbox}
\newpage
\begin{tcolorbox}[enhanced,width=16cm,height=14cm,colback=white,colframe=cpro,boxrule=2pt,arc=0mm,
left=6mm,right=6mm,top=0mm,bottom=4mm,sharp corners,
overlay={\fill[cpro] (frame.north west) rectangle ([yshift=-1.25cm]frame.north east);
\node[anchor=west,text=white,font=\sffamily\bfseries\LARGE] at ([xshift=5mm,yshift=-0.62cm]frame.north west) {Question 84 C};
\node[anchor=east,text=white,font=\sffamily\itshape\large] at ([xshift=-5mm,yshift=-0.62cm]frame.north east) {Probabilités};}]
\vspace*{1.75cm}
\begin{minipage}[t][6.9cm][c]{\linewidth}\LARGE\centering
\begin{minipage}{0.95\linewidth}\raggedright\linespread{1.25}\selectfont On a $P(A)=0{,}8$, $P(B)=0{,}5$ et $P(A\cup B)=0{,}95$. Que vaut $P(A\cap B)$~?\end{minipage}
\end{minipage}

\vfill
\begin{tcolorbox}[enhanced,nobeforeafter,width=7.1cm,height=2.1cm,valign=center,colback=cpro!6,colframe=cpro,arc=2.5mm,boxrule=0.9pt,left=3mm]{\Large\bfseries\color{cpro}a)}\quad\begin{minipage}[c]{5.3cm}\raggedright\Large $0{,}4$\end{minipage}\end{tcolorbox}\hfill\begin{tcolorbox}[enhanced,nobeforeafter,width=7.1cm,height=2.1cm,valign=center,colback=cpro!6,colframe=cpro,arc=2.5mm,boxrule=0.9pt,left=3mm]{\Large\bfseries\color{cpro}b)}\quad\begin{minipage}[c]{5.3cm}\raggedright\Large $0{,}35$\end{minipage}\end{tcolorbox}\par\vspace{2.5mm}
\begin{tcolorbox}[enhanced,nobeforeafter,width=7.1cm,height=2.1cm,valign=center,colback=cpro!6,colframe=cpro,arc=2.5mm,boxrule=0.9pt,left=3mm]{\Large\bfseries\color{cpro}c)}\quad\begin{minipage}[c]{5.3cm}\raggedright\Large $0{,}45$\end{minipage}\end{tcolorbox}\hfill\begin{tcolorbox}[enhanced,nobeforeafter,width=7.1cm,height=2.1cm,valign=center,colback=cpro!6,colframe=cpro,arc=2.5mm,boxrule=0.9pt,left=3mm]{\Large\bfseries\color{cpro}d)}\quad\begin{minipage}[c]{5.3cm}\raggedright\Large $0{,}15$\end{minipage}\end{tcolorbox}
\end{tcolorbox}
\newpage
\begin{tcolorbox}[enhanced,width=16cm,height=14cm,colback=white,colframe=cpro,boxrule=2pt,arc=0mm,
left=6mm,right=6mm,top=0mm,bottom=4mm,sharp corners,
overlay={\fill[cpro] (frame.north west) rectangle ([yshift=-1.25cm]frame.north east);
\node[anchor=west,text=white,font=\sffamily\bfseries\LARGE] at ([xshift=5mm,yshift=-0.62cm]frame.north west) {Question 85 C};
\node[anchor=east,text=white,font=\sffamily\itshape\large] at ([xshift=-5mm,yshift=-0.62cm]frame.north east) {Probabilités};}]
\vspace*{1.75cm}
\begin{minipage}[t][6.9cm][c]{\linewidth}\LARGE\centering
\begin{minipage}{0.95\linewidth}\raggedright\linespread{1.25}\selectfont On a $P(A)=\dfrac{4}{15}$. Que vaut $P(\overline{A})$, la probabilité de l'événement contraire~?\end{minipage}
\end{minipage}

\vfill
\begin{tcolorbox}[enhanced,nobeforeafter,width=7.1cm,height=2.1cm,valign=center,colback=cpro!6,colframe=cpro,arc=2.5mm,boxrule=0.9pt,left=3mm]{\Large\bfseries\color{cpro}a)}\quad\begin{minipage}[c]{5.3cm}\raggedright\Large $\dfrac{11}{30}$\end{minipage}\end{tcolorbox}\hfill\begin{tcolorbox}[enhanced,nobeforeafter,width=7.1cm,height=2.1cm,valign=center,colback=cpro!6,colframe=cpro,arc=2.5mm,boxrule=0.9pt,left=3mm]{\Large\bfseries\color{cpro}b)}\quad\begin{minipage}[c]{5.3cm}\raggedright\Large $\dfrac{11}{15}$\end{minipage}\end{tcolorbox}\par\vspace{2.5mm}
\begin{tcolorbox}[enhanced,nobeforeafter,width=7.1cm,height=2.1cm,valign=center,colback=cpro!6,colframe=cpro,arc=2.5mm,boxrule=0.9pt,left=3mm]{\Large\bfseries\color{cpro}c)}\quad\begin{minipage}[c]{5.3cm}\raggedright\Large $\dfrac{15}{4}$\end{minipage}\end{tcolorbox}\hfill\begin{tcolorbox}[enhanced,nobeforeafter,width=7.1cm,height=2.1cm,valign=center,colback=cpro!6,colframe=cpro,arc=2.5mm,boxrule=0.9pt,left=3mm]{\Large\bfseries\color{cpro}d)}\quad\begin{minipage}[c]{5.3cm}\raggedright\Large $\dfrac{4}{15}$\end{minipage}\end{tcolorbox}
\end{tcolorbox}
\newpage
\begin{tcolorbox}[enhanced,width=16cm,height=14cm,colback=white,colframe=cpro,boxrule=2pt,arc=0mm,
left=6mm,right=6mm,top=0mm,bottom=4mm,sharp corners,
overlay={\fill[cpro] (frame.north west) rectangle ([yshift=-1.25cm]frame.north east);
\node[anchor=west,text=white,font=\sffamily\bfseries\LARGE] at ([xshift=5mm,yshift=-0.62cm]frame.north west) {Question 86 C};
\node[anchor=east,text=white,font=\sffamily\itshape\large] at ([xshift=-5mm,yshift=-0.62cm]frame.north east) {Probabilités};}]
\vspace*{1.75cm}
\begin{minipage}[t][6.9cm][c]{\linewidth}\LARGE\centering
\begin{minipage}{0.95\linewidth}\raggedright\linespread{1.25}\selectfont On a $P(A)=\dfrac{5}{11}$. Que vaut $P(\overline{A})$, la probabilité de l'événement contraire~?\end{minipage}
\end{minipage}

\vfill
\begin{tcolorbox}[enhanced,nobeforeafter,width=7.1cm,height=2.1cm,valign=center,colback=cpro!6,colframe=cpro,arc=2.5mm,boxrule=0.9pt,left=3mm]{\Large\bfseries\color{cpro}a)}\quad\begin{minipage}[c]{5.3cm}\raggedright\Large $\dfrac{3}{11}$\end{minipage}\end{tcolorbox}\hfill\begin{tcolorbox}[enhanced,nobeforeafter,width=7.1cm,height=2.1cm,valign=center,colback=cpro!6,colframe=cpro,arc=2.5mm,boxrule=0.9pt,left=3mm]{\Large\bfseries\color{cpro}b)}\quad\begin{minipage}[c]{5.3cm}\raggedright\Large $\dfrac{6}{11}$\end{minipage}\end{tcolorbox}\par\vspace{2.5mm}
\begin{tcolorbox}[enhanced,nobeforeafter,width=7.1cm,height=2.1cm,valign=center,colback=cpro!6,colframe=cpro,arc=2.5mm,boxrule=0.9pt,left=3mm]{\Large\bfseries\color{cpro}c)}\quad\begin{minipage}[c]{5.3cm}\raggedright\Large $\dfrac{5}{11}$\end{minipage}\end{tcolorbox}\hfill\begin{tcolorbox}[enhanced,nobeforeafter,width=7.1cm,height=2.1cm,valign=center,colback=cpro!6,colframe=cpro,arc=2.5mm,boxrule=0.9pt,left=3mm]{\Large\bfseries\color{cpro}d)}\quad\begin{minipage}[c]{5.3cm}\raggedright\Large $\dfrac{11}{5}$\end{minipage}\end{tcolorbox}
\end{tcolorbox}
\newpage
\begin{tcolorbox}[enhanced,width=16cm,height=14cm,colback=white,colframe=cpro,boxrule=2pt,arc=0mm,
left=6mm,right=6mm,top=0mm,bottom=4mm,sharp corners,
overlay={\fill[cpro] (frame.north west) rectangle ([yshift=-1.25cm]frame.north east);
\node[anchor=west,text=white,font=\sffamily\bfseries\LARGE] at ([xshift=5mm,yshift=-0.62cm]frame.north west) {Question 87 C};
\node[anchor=east,text=white,font=\sffamily\itshape\large] at ([xshift=-5mm,yshift=-0.62cm]frame.north east) {Probabilités};}]
\vspace*{1.75cm}
\begin{minipage}[t][6.9cm][c]{\linewidth}\LARGE\centering
\begin{minipage}{0.95\linewidth}\raggedright\linespread{1.25}\selectfont On lance un dé équilibré à six faces. Quelle est la probabilité d'obtenir un nombre diviseur de 12~?\end{minipage}
\end{minipage}

\vfill
\begin{tcolorbox}[enhanced,nobeforeafter,width=7.1cm,height=2.1cm,valign=center,colback=cpro!6,colframe=cpro,arc=2.5mm,boxrule=0.9pt,left=3mm]{\Large\bfseries\color{cpro}a)}\quad\begin{minipage}[c]{5.3cm}\raggedright\Large $\dfrac{5}{6}$\end{minipage}\end{tcolorbox}\hfill\begin{tcolorbox}[enhanced,nobeforeafter,width=7.1cm,height=2.1cm,valign=center,colback=cpro!6,colframe=cpro,arc=2.5mm,boxrule=0.9pt,left=3mm]{\Large\bfseries\color{cpro}b)}\quad\begin{minipage}[c]{5.3cm}\raggedright\Large $\dfrac{2}{3}$\end{minipage}\end{tcolorbox}\par\vspace{2.5mm}
\begin{tcolorbox}[enhanced,nobeforeafter,width=7.1cm,height=2.1cm,valign=center,colback=cpro!6,colframe=cpro,arc=2.5mm,boxrule=0.9pt,left=3mm]{\Large\bfseries\color{cpro}c)}\quad\begin{minipage}[c]{5.3cm}\raggedright\Large $\dfrac{1}{6}$\end{minipage}\end{tcolorbox}\hfill\begin{tcolorbox}[enhanced,nobeforeafter,width=7.1cm,height=2.1cm,valign=center,colback=cpro!6,colframe=cpro,arc=2.5mm,boxrule=0.9pt,left=3mm]{\Large\bfseries\color{cpro}d)}\quad\begin{minipage}[c]{5.3cm}\raggedright\Large $1$\end{minipage}\end{tcolorbox}
\end{tcolorbox}
\newpage
\begin{tcolorbox}[enhanced,width=16cm,height=14cm,colback=white,colframe=cpro,boxrule=2pt,arc=0mm,
left=6mm,right=6mm,top=0mm,bottom=4mm,sharp corners,
overlay={\fill[cpro] (frame.north west) rectangle ([yshift=-1.25cm]frame.north east);
\node[anchor=west,text=white,font=\sffamily\bfseries\LARGE] at ([xshift=5mm,yshift=-0.62cm]frame.north west) {Question 88 C};
\node[anchor=east,text=white,font=\sffamily\itshape\large] at ([xshift=-5mm,yshift=-0.62cm]frame.north east) {Probabilités};}]
\vspace*{1.75cm}
\begin{minipage}[t][6.9cm][c]{\linewidth}\LARGE\centering
\begin{minipage}{0.95\linewidth}\raggedright\linespread{1.25}\selectfont On lance un dé équilibré à six faces. Quelle est la probabilité d'obtenir un nombre pair et inférieur à 5~?\end{minipage}
\end{minipage}

\vfill
\begin{tcolorbox}[enhanced,nobeforeafter,width=7.1cm,height=2.1cm,valign=center,colback=cpro!6,colframe=cpro,arc=2.5mm,boxrule=0.9pt,left=3mm]{\Large\bfseries\color{cpro}a)}\quad\begin{minipage}[c]{5.3cm}\raggedright\Large $\dfrac{2}{3}$\end{minipage}\end{tcolorbox}\hfill\begin{tcolorbox}[enhanced,nobeforeafter,width=7.1cm,height=2.1cm,valign=center,colback=cpro!6,colframe=cpro,arc=2.5mm,boxrule=0.9pt,left=3mm]{\Large\bfseries\color{cpro}b)}\quad\begin{minipage}[c]{5.3cm}\raggedright\Large $\dfrac{1}{6}$\end{minipage}\end{tcolorbox}\par\vspace{2.5mm}
\begin{tcolorbox}[enhanced,nobeforeafter,width=7.1cm,height=2.1cm,valign=center,colback=cpro!6,colframe=cpro,arc=2.5mm,boxrule=0.9pt,left=3mm]{\Large\bfseries\color{cpro}c)}\quad\begin{minipage}[c]{5.3cm}\raggedright\Large $\dfrac{1}{2}$\end{minipage}\end{tcolorbox}\hfill\begin{tcolorbox}[enhanced,nobeforeafter,width=7.1cm,height=2.1cm,valign=center,colback=cpro!6,colframe=cpro,arc=2.5mm,boxrule=0.9pt,left=3mm]{\Large\bfseries\color{cpro}d)}\quad\begin{minipage}[c]{5.3cm}\raggedright\Large $\dfrac{1}{3}$\end{minipage}\end{tcolorbox}
\end{tcolorbox}
\newpage
\begin{tcolorbox}[enhanced,width=16cm,height=14cm,colback=white,colframe=cpro,boxrule=2pt,arc=0mm,
left=6mm,right=6mm,top=0mm,bottom=4mm,sharp corners,
overlay={\fill[cpro] (frame.north west) rectangle ([yshift=-1.25cm]frame.north east);
\node[anchor=west,text=white,font=\sffamily\bfseries\LARGE] at ([xshift=5mm,yshift=-0.62cm]frame.north west) {Question 89 C};
\node[anchor=east,text=white,font=\sffamily\itshape\large] at ([xshift=-5mm,yshift=-0.62cm]frame.north east) {Probabilités};}]
\vspace*{1.75cm}
\begin{minipage}[t][6.9cm][c]{\linewidth}\LARGE\centering
\begin{minipage}{0.95\linewidth}\raggedright\linespread{1.25}\selectfont On lance $4$ pièces équilibrées. Quelle est la probabilité de n'obtenir que des PILE~?\end{minipage}
\end{minipage}

\vfill
\begin{tcolorbox}[enhanced,nobeforeafter,width=7.1cm,height=2.1cm,valign=center,colback=cpro!6,colframe=cpro,arc=2.5mm,boxrule=0.9pt,left=3mm]{\Large\bfseries\color{cpro}a)}\quad\begin{minipage}[c]{5.3cm}\raggedright\Large $\dfrac{1}{16}$\end{minipage}\end{tcolorbox}\hfill\begin{tcolorbox}[enhanced,nobeforeafter,width=7.1cm,height=2.1cm,valign=center,colback=cpro!6,colframe=cpro,arc=2.5mm,boxrule=0.9pt,left=3mm]{\Large\bfseries\color{cpro}b)}\quad\begin{minipage}[c]{5.3cm}\raggedright\Large $\dfrac{1}{8}$\end{minipage}\end{tcolorbox}\par\vspace{2.5mm}
\begin{tcolorbox}[enhanced,nobeforeafter,width=7.1cm,height=2.1cm,valign=center,colback=cpro!6,colframe=cpro,arc=2.5mm,boxrule=0.9pt,left=3mm]{\Large\bfseries\color{cpro}c)}\quad\begin{minipage}[c]{5.3cm}\raggedright\Large $\dfrac{1}{2}$\end{minipage}\end{tcolorbox}\hfill\begin{tcolorbox}[enhanced,nobeforeafter,width=7.1cm,height=2.1cm,valign=center,colback=cpro!6,colframe=cpro,arc=2.5mm,boxrule=0.9pt,left=3mm]{\Large\bfseries\color{cpro}d)}\quad\begin{minipage}[c]{5.3cm}\raggedright\Large $\dfrac{15}{16}$\end{minipage}\end{tcolorbox}
\end{tcolorbox}
\newpage
\begin{tcolorbox}[enhanced,width=16cm,height=14cm,colback=white,colframe=cpro,boxrule=2pt,arc=0mm,
left=6mm,right=6mm,top=0mm,bottom=4mm,sharp corners,
overlay={\fill[cpro] (frame.north west) rectangle ([yshift=-1.25cm]frame.north east);
\node[anchor=west,text=white,font=\sffamily\bfseries\LARGE] at ([xshift=5mm,yshift=-0.62cm]frame.north west) {Question 90 C};
\node[anchor=east,text=white,font=\sffamily\itshape\large] at ([xshift=-5mm,yshift=-0.62cm]frame.north east) {Probabilités};}]
\vspace*{1.75cm}
\begin{minipage}[t][6.9cm][c]{\linewidth}\LARGE\centering
\begin{minipage}{0.95\linewidth}\raggedright\linespread{1.25}\selectfont On lance $4$ pièces équilibrées. Quelle est la probabilité d'obtenir au moins une fois FACE~?\end{minipage}
\end{minipage}

\vfill
\begin{tcolorbox}[enhanced,nobeforeafter,width=7.1cm,height=2.1cm,valign=center,colback=cpro!6,colframe=cpro,arc=2.5mm,boxrule=0.9pt,left=3mm]{\Large\bfseries\color{cpro}a)}\quad\begin{minipage}[c]{5.3cm}\raggedright\Large $\dfrac{1}{2}$\end{minipage}\end{tcolorbox}\hfill\begin{tcolorbox}[enhanced,nobeforeafter,width=7.1cm,height=2.1cm,valign=center,colback=cpro!6,colframe=cpro,arc=2.5mm,boxrule=0.9pt,left=3mm]{\Large\bfseries\color{cpro}b)}\quad\begin{minipage}[c]{5.3cm}\raggedright\Large $\dfrac{1}{16}$\end{minipage}\end{tcolorbox}\par\vspace{2.5mm}
\begin{tcolorbox}[enhanced,nobeforeafter,width=7.1cm,height=2.1cm,valign=center,colback=cpro!6,colframe=cpro,arc=2.5mm,boxrule=0.9pt,left=3mm]{\Large\bfseries\color{cpro}c)}\quad\begin{minipage}[c]{5.3cm}\raggedright\Large $\dfrac{15}{16}$\end{minipage}\end{tcolorbox}\hfill\begin{tcolorbox}[enhanced,nobeforeafter,width=7.1cm,height=2.1cm,valign=center,colback=cpro!6,colframe=cpro,arc=2.5mm,boxrule=0.9pt,left=3mm]{\Large\bfseries\color{cpro}d)}\quad\begin{minipage}[c]{5.3cm}\raggedright\Large $\dfrac{1}{8}$\end{minipage}\end{tcolorbox}
\end{tcolorbox}
\newpage
\begin{tcolorbox}[enhanced,width=16cm,height=14cm,colback=white,colframe=cpro,boxrule=2pt,arc=0mm,
left=6mm,right=6mm,top=0mm,bottom=4mm,sharp corners,
overlay={\fill[cpro] (frame.north west) rectangle ([yshift=-1.25cm]frame.north east);
\node[anchor=west,text=white,font=\sffamily\bfseries\LARGE] at ([xshift=5mm,yshift=-0.62cm]frame.north west) {Question 91 C};
\node[anchor=east,text=white,font=\sffamily\itshape\large] at ([xshift=-5mm,yshift=-0.62cm]frame.north east) {Probabilités};}]
\vspace*{1.75cm}
\begin{minipage}[t][6.9cm][c]{\linewidth}\Large\centering
\begin{minipage}{0.95\linewidth}\raggedright\linespread{1.25}\selectfont Voici la loi de probabilité d'une expérience aléatoire. Quelle est la probabilité manquante~? $\begin{array}{|l|c|c|c|c|}\hline \text{Issue} & 1 & 2 & 3 & 4\\\hline \text{Probabilité} & 0{,}45 & 0{,}2 &~? & 0{,}1\\\hline\end{array}$\end{minipage}
\end{minipage}

\vfill
\begin{tcolorbox}[enhanced,nobeforeafter,width=7.1cm,height=2.1cm,valign=center,colback=cpro!6,colframe=cpro,arc=2.5mm,boxrule=0.9pt,left=3mm]{\Large\bfseries\color{cpro}a)}\quad\begin{minipage}[c]{5.3cm}\raggedright\Large $0{,}25$\end{minipage}\end{tcolorbox}\hfill\begin{tcolorbox}[enhanced,nobeforeafter,width=7.1cm,height=2.1cm,valign=center,colback=cpro!6,colframe=cpro,arc=2.5mm,boxrule=0.9pt,left=3mm]{\Large\bfseries\color{cpro}b)}\quad\begin{minipage}[c]{5.3cm}\raggedright\Large $0{,}75$\end{minipage}\end{tcolorbox}\par\vspace{2.5mm}
\begin{tcolorbox}[enhanced,nobeforeafter,width=7.1cm,height=2.1cm,valign=center,colback=cpro!6,colframe=cpro,arc=2.5mm,boxrule=0.9pt,left=3mm]{\Large\bfseries\color{cpro}c)}\quad\begin{minipage}[c]{5.3cm}\raggedright\Large $0{,}35$\end{minipage}\end{tcolorbox}\hfill\begin{tcolorbox}[enhanced,nobeforeafter,width=7.1cm,height=2.1cm,valign=center,colback=cpro!6,colframe=cpro,arc=2.5mm,boxrule=0.9pt,left=3mm]{\Large\bfseries\color{cpro}d)}\quad\begin{minipage}[c]{5.3cm}\raggedright\Large $0{,}5$\end{minipage}\end{tcolorbox}
\end{tcolorbox}
\newpage
\begin{tcolorbox}[enhanced,width=16cm,height=14cm,colback=white,colframe=cpro,boxrule=2pt,arc=0mm,
left=6mm,right=6mm,top=0mm,bottom=4mm,sharp corners,
overlay={\fill[cpro] (frame.north west) rectangle ([yshift=-1.25cm]frame.north east);
\node[anchor=west,text=white,font=\sffamily\bfseries\LARGE] at ([xshift=5mm,yshift=-0.62cm]frame.north west) {Question 92 C};
\node[anchor=east,text=white,font=\sffamily\itshape\large] at ([xshift=-5mm,yshift=-0.62cm]frame.north east) {Probabilités};}]
\vspace*{1.75cm}
\begin{minipage}[t][6.9cm][c]{\linewidth}\Large\centering
\begin{minipage}{0.95\linewidth}\raggedright\linespread{1.25}\selectfont Voici la loi de probabilité d'une expérience aléatoire. Quelle est la probabilité manquante~? $\begin{array}{|l|c|c|c|c|}\hline \text{Issue} & 1 & 2 & 3 & 4\\\hline \text{Probabilité} & 0{,}1 & 0{,}1 &~? & 0{,}5\\\hline\end{array}$\end{minipage}
\end{minipage}

\vfill
\begin{tcolorbox}[enhanced,nobeforeafter,width=7.1cm,height=2.1cm,valign=center,colback=cpro!6,colframe=cpro,arc=2.5mm,boxrule=0.9pt,left=3mm]{\Large\bfseries\color{cpro}a)}\quad\begin{minipage}[c]{5.3cm}\raggedright\Large $0{,}7$\end{minipage}\end{tcolorbox}\hfill\begin{tcolorbox}[enhanced,nobeforeafter,width=7.1cm,height=2.1cm,valign=center,colback=cpro!6,colframe=cpro,arc=2.5mm,boxrule=0.9pt,left=3mm]{\Large\bfseries\color{cpro}b)}\quad\begin{minipage}[c]{5.3cm}\raggedright\Large $0{,}25$\end{minipage}\end{tcolorbox}\par\vspace{2.5mm}
\begin{tcolorbox}[enhanced,nobeforeafter,width=7.1cm,height=2.1cm,valign=center,colback=cpro!6,colframe=cpro,arc=2.5mm,boxrule=0.9pt,left=3mm]{\Large\bfseries\color{cpro}c)}\quad\begin{minipage}[c]{5.3cm}\raggedright\Large $0{,}8$\end{minipage}\end{tcolorbox}\hfill\begin{tcolorbox}[enhanced,nobeforeafter,width=7.1cm,height=2.1cm,valign=center,colback=cpro!6,colframe=cpro,arc=2.5mm,boxrule=0.9pt,left=3mm]{\Large\bfseries\color{cpro}d)}\quad\begin{minipage}[c]{5.3cm}\raggedright\Large $0{,}3$\end{minipage}\end{tcolorbox}
\end{tcolorbox}
\newpage
\begin{tcolorbox}[enhanced,width=16cm,height=14cm,colback=white,colframe=cpro,boxrule=2pt,arc=0mm,
left=6mm,right=6mm,top=0mm,bottom=4mm,sharp corners,
overlay={\fill[cpro] (frame.north west) rectangle ([yshift=-1.25cm]frame.north east);
\node[anchor=west,text=white,font=\sffamily\bfseries\LARGE] at ([xshift=5mm,yshift=-0.62cm]frame.north west) {Question 93 C};
\node[anchor=east,text=white,font=\sffamily\itshape\large] at ([xshift=-5mm,yshift=-0.62cm]frame.north east) {Probabilités};}]
\vspace*{1.75cm}
\begin{minipage}[t][6.9cm][c]{\linewidth}\Large\centering
\begin{minipage}{0.95\linewidth}\raggedright\linespread{1.25}\selectfont Une urne contient $8$ boules dont $5$ rouges. On tire une boule, on la remet, puis on en tire une deuxième. Quelle est la probabilité de tirer deux boules rouges~?\end{minipage}
\end{minipage}

\vfill
\begin{tcolorbox}[enhanced,nobeforeafter,width=7.1cm,height=2.1cm,valign=center,colback=cpro!6,colframe=cpro,arc=2.5mm,boxrule=0.9pt,left=3mm]{\Large\bfseries\color{cpro}a)}\quad\begin{minipage}[c]{5.3cm}\raggedright\Large $\dfrac{25}{64}$\end{minipage}\end{tcolorbox}\hfill\begin{tcolorbox}[enhanced,nobeforeafter,width=7.1cm,height=2.1cm,valign=center,colback=cpro!6,colframe=cpro,arc=2.5mm,boxrule=0.9pt,left=3mm]{\Large\bfseries\color{cpro}b)}\quad\begin{minipage}[c]{5.3cm}\raggedright\Large $\dfrac{5}{8}$\end{minipage}\end{tcolorbox}\par\vspace{2.5mm}
\begin{tcolorbox}[enhanced,nobeforeafter,width=7.1cm,height=2.1cm,valign=center,colback=cpro!6,colframe=cpro,arc=2.5mm,boxrule=0.9pt,left=3mm]{\Large\bfseries\color{cpro}c)}\quad\begin{minipage}[c]{5.3cm}\raggedright\Large $\dfrac{5}{14}$\end{minipage}\end{tcolorbox}\hfill\begin{tcolorbox}[enhanced,nobeforeafter,width=7.1cm,height=2.1cm,valign=center,colback=cpro!6,colframe=cpro,arc=2.5mm,boxrule=0.9pt,left=3mm]{\Large\bfseries\color{cpro}d)}\quad\begin{minipage}[c]{5.3cm}\raggedright\Large $1$\end{minipage}\end{tcolorbox}
\end{tcolorbox}
\newpage
\begin{tcolorbox}[enhanced,width=16cm,height=14cm,colback=white,colframe=cpro,boxrule=2pt,arc=0mm,
left=6mm,right=6mm,top=0mm,bottom=4mm,sharp corners,
overlay={\fill[cpro] (frame.north west) rectangle ([yshift=-1.25cm]frame.north east);
\node[anchor=west,text=white,font=\sffamily\bfseries\LARGE] at ([xshift=5mm,yshift=-0.62cm]frame.north west) {Question 94 C};
\node[anchor=east,text=white,font=\sffamily\itshape\large] at ([xshift=-5mm,yshift=-0.62cm]frame.north east) {Probabilités};}]
\vspace*{1.75cm}
\begin{minipage}[t][6.9cm][c]{\linewidth}\Large\centering
\begin{minipage}{0.95\linewidth}\raggedright\linespread{1.25}\selectfont Une urne contient $10$ boules dont $3$ rouges. On tire une boule, on la remet, puis on en tire une deuxième. Quelle est la probabilité de tirer deux boules rouges~?\end{minipage}
\end{minipage}

\vfill
\begin{tcolorbox}[enhanced,nobeforeafter,width=7.1cm,height=2.1cm,valign=center,colback=cpro!6,colframe=cpro,arc=2.5mm,boxrule=0.9pt,left=3mm]{\Large\bfseries\color{cpro}a)}\quad\begin{minipage}[c]{5.3cm}\raggedright\Large $\dfrac{9}{100}$\end{minipage}\end{tcolorbox}\hfill\begin{tcolorbox}[enhanced,nobeforeafter,width=7.1cm,height=2.1cm,valign=center,colback=cpro!6,colframe=cpro,arc=2.5mm,boxrule=0.9pt,left=3mm]{\Large\bfseries\color{cpro}b)}\quad\begin{minipage}[c]{5.3cm}\raggedright\Large $\dfrac{3}{10}$\end{minipage}\end{tcolorbox}\par\vspace{2.5mm}
\begin{tcolorbox}[enhanced,nobeforeafter,width=7.1cm,height=2.1cm,valign=center,colback=cpro!6,colframe=cpro,arc=2.5mm,boxrule=0.9pt,left=3mm]{\Large\bfseries\color{cpro}c)}\quad\begin{minipage}[c]{5.3cm}\raggedright\Large $\dfrac{1}{15}$\end{minipage}\end{tcolorbox}\hfill\begin{tcolorbox}[enhanced,nobeforeafter,width=7.1cm,height=2.1cm,valign=center,colback=cpro!6,colframe=cpro,arc=2.5mm,boxrule=0.9pt,left=3mm]{\Large\bfseries\color{cpro}d)}\quad\begin{minipage}[c]{5.3cm}\raggedright\Large $\dfrac{3}{5}$\end{minipage}\end{tcolorbox}
\end{tcolorbox}
\newpage
\begin{tcolorbox}[enhanced,width=16cm,height=14cm,colback=white,colframe=cpro,boxrule=2pt,arc=0mm,
left=6mm,right=6mm,top=0mm,bottom=4mm,sharp corners,
overlay={\fill[cpro] (frame.north west) rectangle ([yshift=-1.25cm]frame.north east);
\node[anchor=west,text=white,font=\sffamily\bfseries\LARGE] at ([xshift=5mm,yshift=-0.62cm]frame.north west) {Question 95 C};
\node[anchor=east,text=white,font=\sffamily\itshape\large] at ([xshift=-5mm,yshift=-0.62cm]frame.north east) {Probabilités};}]
\vspace*{1.75cm}
\begin{minipage}[t][6.9cm][c]{\linewidth}\LARGE\centering
\begin{minipage}{0.95\linewidth}\raggedright\linespread{1.25}\selectfont $A$ et $B$ sont deux événements incompatibles, avec $P(A)=0{,}08$ et $P(B)=0{,}5$. Que vaut $P(A\cup B)$~?\end{minipage}
\end{minipage}

\vfill
\begin{tcolorbox}[enhanced,nobeforeafter,width=7.1cm,height=2.1cm,valign=center,colback=cpro!6,colframe=cpro,arc=2.5mm,boxrule=0.9pt,left=3mm]{\Large\bfseries\color{cpro}a)}\quad\begin{minipage}[c]{5.3cm}\raggedright\Large $0{,}04$\end{minipage}\end{tcolorbox}\hfill\begin{tcolorbox}[enhanced,nobeforeafter,width=7.1cm,height=2.1cm,valign=center,colback=cpro!6,colframe=cpro,arc=2.5mm,boxrule=0.9pt,left=3mm]{\Large\bfseries\color{cpro}b)}\quad\begin{minipage}[c]{5.3cm}\raggedright\Large $0{,}58$\end{minipage}\end{tcolorbox}\par\vspace{2.5mm}
\begin{tcolorbox}[enhanced,nobeforeafter,width=7.1cm,height=2.1cm,valign=center,colback=cpro!6,colframe=cpro,arc=2.5mm,boxrule=0.9pt,left=3mm]{\Large\bfseries\color{cpro}c)}\quad\begin{minipage}[c]{5.3cm}\raggedright\Large $0{,}54$\end{minipage}\end{tcolorbox}\hfill\begin{tcolorbox}[enhanced,nobeforeafter,width=7.1cm,height=2.1cm,valign=center,colback=cpro!6,colframe=cpro,arc=2.5mm,boxrule=0.9pt,left=3mm]{\Large\bfseries\color{cpro}d)}\quad\begin{minipage}[c]{5.3cm}\raggedright\Large $0{,}42$\end{minipage}\end{tcolorbox}
\end{tcolorbox}
\newpage
\begin{tcolorbox}[enhanced,width=16cm,height=14cm,colback=white,colframe=cpro,boxrule=2pt,arc=0mm,
left=6mm,right=6mm,top=0mm,bottom=4mm,sharp corners,
overlay={\fill[cpro] (frame.north west) rectangle ([yshift=-1.25cm]frame.north east);
\node[anchor=west,text=white,font=\sffamily\bfseries\LARGE] at ([xshift=5mm,yshift=-0.62cm]frame.north west) {Question 96 C};
\node[anchor=east,text=white,font=\sffamily\itshape\large] at ([xshift=-5mm,yshift=-0.62cm]frame.north east) {Probabilités};}]
\vspace*{1.75cm}
\begin{minipage}[t][6.9cm][c]{\linewidth}\LARGE\centering
\begin{minipage}{0.95\linewidth}\raggedright\linespread{1.25}\selectfont $A$ et $B$ sont deux événements incompatibles, avec $P(A)=\dfrac{1}{8}$ et $P(B)=\dfrac{3}{4}$. Que vaut $P(A\cup B)$~?\end{minipage}
\end{minipage}

\vfill
\begin{tcolorbox}[enhanced,nobeforeafter,width=7.1cm,height=2.1cm,valign=center,colback=cpro!6,colframe=cpro,arc=2.5mm,boxrule=0.9pt,left=3mm]{\Large\bfseries\color{cpro}a)}\quad\begin{minipage}[c]{5.3cm}\raggedright\Large $\dfrac{7}{8}$\end{minipage}\end{tcolorbox}\hfill\begin{tcolorbox}[enhanced,nobeforeafter,width=7.1cm,height=2.1cm,valign=center,colback=cpro!6,colframe=cpro,arc=2.5mm,boxrule=0.9pt,left=3mm]{\Large\bfseries\color{cpro}b)}\quad\begin{minipage}[c]{5.3cm}\raggedright\Large $\dfrac{5}{8}$\end{minipage}\end{tcolorbox}\par\vspace{2.5mm}
\begin{tcolorbox}[enhanced,nobeforeafter,width=7.1cm,height=2.1cm,valign=center,colback=cpro!6,colframe=cpro,arc=2.5mm,boxrule=0.9pt,left=3mm]{\Large\bfseries\color{cpro}c)}\quad\begin{minipage}[c]{5.3cm}\raggedright\Large $\dfrac{3}{32}$\end{minipage}\end{tcolorbox}\hfill\begin{tcolorbox}[enhanced,nobeforeafter,width=7.1cm,height=2.1cm,valign=center,colback=cpro!6,colframe=cpro,arc=2.5mm,boxrule=0.9pt,left=3mm]{\Large\bfseries\color{cpro}d)}\quad\begin{minipage}[c]{5.3cm}\raggedright\Large $\dfrac{25}{32}$\end{minipage}\end{tcolorbox}
\end{tcolorbox}
\newpage
\begin{tcolorbox}[enhanced,width=16cm,height=14cm,colback=white,colframe=cnom,boxrule=2pt,arc=0mm,
left=6mm,right=6mm,top=0mm,bottom=4mm,sharp corners,
overlay={\fill[cnom] (frame.north west) rectangle ([yshift=-1.25cm]frame.north east);
\node[anchor=west,text=white,font=\sffamily\bfseries\LARGE] at ([xshift=5mm,yshift=-0.62cm]frame.north west) {Question 1 D};
\node[anchor=east,text=white,font=\sffamily\itshape\large] at ([xshift=-5mm,yshift=-0.62cm]frame.north east) {Nombres et calculs};}]
\vspace*{1.75cm}
\begin{minipage}[t][6.9cm][c]{\linewidth}\LARGE\centering
\begin{minipage}{0.95\linewidth}\raggedright\linespread{1.25}\selectfont Écrire $\sqrt{200}$ sous la forme $a\sqrt{b}$, avec $b$ le plus petit possible.\end{minipage}
\end{minipage}

\vfill
\begin{tcolorbox}[enhanced,nobeforeafter,width=7.1cm,height=2.1cm,valign=center,colback=cnom!6,colframe=cnom,arc=2.5mm,boxrule=0.9pt,left=3mm]{\Large\bfseries\color{cnom}a)}\quad\begin{minipage}[c]{5.3cm}\raggedright\Large $10\sqrt{2}$\end{minipage}\end{tcolorbox}\hfill\begin{tcolorbox}[enhanced,nobeforeafter,width=7.1cm,height=2.1cm,valign=center,colback=cnom!6,colframe=cnom,arc=2.5mm,boxrule=0.9pt,left=3mm]{\Large\bfseries\color{cnom}b)}\quad\begin{minipage}[c]{5.3cm}\raggedright\Large $100\sqrt{2}$\end{minipage}\end{tcolorbox}\par\vspace{2.5mm}
\begin{tcolorbox}[enhanced,nobeforeafter,width=7.1cm,height=2.1cm,valign=center,colback=cnom!6,colframe=cnom,arc=2.5mm,boxrule=0.9pt,left=3mm]{\Large\bfseries\color{cnom}c)}\quad\begin{minipage}[c]{5.3cm}\raggedright\Large $2\sqrt{10}$\end{minipage}\end{tcolorbox}\hfill\begin{tcolorbox}[enhanced,nobeforeafter,width=7.1cm,height=2.1cm,valign=center,colback=cnom!6,colframe=cnom,arc=2.5mm,boxrule=0.9pt,left=3mm]{\Large\bfseries\color{cnom}d)}\quad\begin{minipage}[c]{5.3cm}\raggedright\Large $10\sqrt{4}$\end{minipage}\end{tcolorbox}
\end{tcolorbox}
\newpage
\begin{tcolorbox}[enhanced,width=16cm,height=14cm,colback=white,colframe=cnom,boxrule=2pt,arc=0mm,
left=6mm,right=6mm,top=0mm,bottom=4mm,sharp corners,
overlay={\fill[cnom] (frame.north west) rectangle ([yshift=-1.25cm]frame.north east);
\node[anchor=west,text=white,font=\sffamily\bfseries\LARGE] at ([xshift=5mm,yshift=-0.62cm]frame.north west) {Question 2 D};
\node[anchor=east,text=white,font=\sffamily\itshape\large] at ([xshift=-5mm,yshift=-0.62cm]frame.north east) {Nombres et calculs};}]
\vspace*{1.75cm}
\begin{minipage}[t][6.9cm][c]{\linewidth}\LARGE\centering
\begin{minipage}{0.95\linewidth}\raggedright\linespread{1.25}\selectfont Écrire $\sqrt{75}$ sous la forme $a\sqrt{b}$, avec $b$ le plus petit possible.\end{minipage}
\end{minipage}

\vfill
\begin{tcolorbox}[enhanced,nobeforeafter,width=7.1cm,height=2.1cm,valign=center,colback=cnom!6,colframe=cnom,arc=2.5mm,boxrule=0.9pt,left=3mm]{\Large\bfseries\color{cnom}a)}\quad\begin{minipage}[c]{5.3cm}\raggedright\Large $3\sqrt{5}$\end{minipage}\end{tcolorbox}\hfill\begin{tcolorbox}[enhanced,nobeforeafter,width=7.1cm,height=2.1cm,valign=center,colback=cnom!6,colframe=cnom,arc=2.5mm,boxrule=0.9pt,left=3mm]{\Large\bfseries\color{cnom}b)}\quad\begin{minipage}[c]{5.3cm}\raggedright\Large $25\sqrt{3}$\end{minipage}\end{tcolorbox}\par\vspace{2.5mm}
\begin{tcolorbox}[enhanced,nobeforeafter,width=7.1cm,height=2.1cm,valign=center,colback=cnom!6,colframe=cnom,arc=2.5mm,boxrule=0.9pt,left=3mm]{\Large\bfseries\color{cnom}c)}\quad\begin{minipage}[c]{5.3cm}\raggedright\Large $5\sqrt{3}$\end{minipage}\end{tcolorbox}\hfill\begin{tcolorbox}[enhanced,nobeforeafter,width=7.1cm,height=2.1cm,valign=center,colback=cnom!6,colframe=cnom,arc=2.5mm,boxrule=0.9pt,left=3mm]{\Large\bfseries\color{cnom}d)}\quad\begin{minipage}[c]{5.3cm}\raggedright\Large $5\sqrt{6}$\end{minipage}\end{tcolorbox}
\end{tcolorbox}
\newpage
\begin{tcolorbox}[enhanced,width=16cm,height=14cm,colback=white,colframe=cnom,boxrule=2pt,arc=0mm,
left=6mm,right=6mm,top=0mm,bottom=4mm,sharp corners,
overlay={\fill[cnom] (frame.north west) rectangle ([yshift=-1.25cm]frame.north east);
\node[anchor=west,text=white,font=\sffamily\bfseries\LARGE] at ([xshift=5mm,yshift=-0.62cm]frame.north west) {Question 3 D};
\node[anchor=east,text=white,font=\sffamily\itshape\large] at ([xshift=-5mm,yshift=-0.62cm]frame.north east) {Nombres et calculs};}]
\vspace*{1.75cm}
\begin{minipage}[t][6.9cm][c]{\linewidth}\LARGE\centering
\begin{minipage}{0.95\linewidth}\raggedright\linespread{1.25}\selectfont Que vaut $\sqrt{28}+\sqrt{63}$~?\end{minipage}
\end{minipage}

\vfill
\begin{tcolorbox}[enhanced,nobeforeafter,width=7.1cm,height=2.1cm,valign=center,colback=cnom!6,colframe=cnom,arc=2.5mm,boxrule=0.9pt,left=3mm]{\Large\bfseries\color{cnom}a)}\quad\begin{minipage}[c]{5.3cm}\raggedright\Large $\sqrt{91}$\end{minipage}\end{tcolorbox}\hfill\begin{tcolorbox}[enhanced,nobeforeafter,width=7.1cm,height=2.1cm,valign=center,colback=cnom!6,colframe=cnom,arc=2.5mm,boxrule=0.9pt,left=3mm]{\Large\bfseries\color{cnom}b)}\quad\begin{minipage}[c]{5.3cm}\raggedright\Large $5\sqrt{7}$\end{minipage}\end{tcolorbox}\par\vspace{2.5mm}
\begin{tcolorbox}[enhanced,nobeforeafter,width=7.1cm,height=2.1cm,valign=center,colback=cnom!6,colframe=cnom,arc=2.5mm,boxrule=0.9pt,left=3mm]{\Large\bfseries\color{cnom}c)}\quad\begin{minipage}[c]{5.3cm}\raggedright\Large $6\sqrt{7}$\end{minipage}\end{tcolorbox}\hfill\begin{tcolorbox}[enhanced,nobeforeafter,width=7.1cm,height=2.1cm,valign=center,colback=cnom!6,colframe=cnom,arc=2.5mm,boxrule=0.9pt,left=3mm]{\Large\bfseries\color{cnom}d)}\quad\begin{minipage}[c]{5.3cm}\raggedright\Large $5\sqrt{14}$\end{minipage}\end{tcolorbox}
\end{tcolorbox}
\newpage
\begin{tcolorbox}[enhanced,width=16cm,height=14cm,colback=white,colframe=cnom,boxrule=2pt,arc=0mm,
left=6mm,right=6mm,top=0mm,bottom=4mm,sharp corners,
overlay={\fill[cnom] (frame.north west) rectangle ([yshift=-1.25cm]frame.north east);
\node[anchor=west,text=white,font=\sffamily\bfseries\LARGE] at ([xshift=5mm,yshift=-0.62cm]frame.north west) {Question 4 D};
\node[anchor=east,text=white,font=\sffamily\itshape\large] at ([xshift=-5mm,yshift=-0.62cm]frame.north east) {Nombres et calculs};}]
\vspace*{1.75cm}
\begin{minipage}[t][6.9cm][c]{\linewidth}\LARGE\centering
\begin{minipage}{0.95\linewidth}\raggedright\linespread{1.25}\selectfont Que vaut $\sqrt{48}+\sqrt{12}$~?\end{minipage}
\end{minipage}

\vfill
\begin{tcolorbox}[enhanced,nobeforeafter,width=7.1cm,height=2.1cm,valign=center,colback=cnom!6,colframe=cnom,arc=2.5mm,boxrule=0.9pt,left=3mm]{\Large\bfseries\color{cnom}a)}\quad\begin{minipage}[c]{5.3cm}\raggedright\Large $8\sqrt{3}$\end{minipage}\end{tcolorbox}\hfill\begin{tcolorbox}[enhanced,nobeforeafter,width=7.1cm,height=2.1cm,valign=center,colback=cnom!6,colframe=cnom,arc=2.5mm,boxrule=0.9pt,left=3mm]{\Large\bfseries\color{cnom}b)}\quad\begin{minipage}[c]{5.3cm}\raggedright\Large $\sqrt{60}$\end{minipage}\end{tcolorbox}\par\vspace{2.5mm}
\begin{tcolorbox}[enhanced,nobeforeafter,width=7.1cm,height=2.1cm,valign=center,colback=cnom!6,colframe=cnom,arc=2.5mm,boxrule=0.9pt,left=3mm]{\Large\bfseries\color{cnom}c)}\quad\begin{minipage}[c]{5.3cm}\raggedright\Large $6\sqrt{3}$\end{minipage}\end{tcolorbox}\hfill\begin{tcolorbox}[enhanced,nobeforeafter,width=7.1cm,height=2.1cm,valign=center,colback=cnom!6,colframe=cnom,arc=2.5mm,boxrule=0.9pt,left=3mm]{\Large\bfseries\color{cnom}d)}\quad\begin{minipage}[c]{5.3cm}\raggedright\Large $6\sqrt{6}$\end{minipage}\end{tcolorbox}
\end{tcolorbox}
\newpage
\begin{tcolorbox}[enhanced,width=16cm,height=14cm,colback=white,colframe=cnom,boxrule=2pt,arc=0mm,
left=6mm,right=6mm,top=0mm,bottom=4mm,sharp corners,
overlay={\fill[cnom] (frame.north west) rectangle ([yshift=-1.25cm]frame.north east);
\node[anchor=west,text=white,font=\sffamily\bfseries\LARGE] at ([xshift=5mm,yshift=-0.62cm]frame.north west) {Question 5 D};
\node[anchor=east,text=white,font=\sffamily\itshape\large] at ([xshift=-5mm,yshift=-0.62cm]frame.north east) {Nombres et calculs};}]
\vspace*{1.75cm}
\begin{minipage}[t][6.9cm][c]{\linewidth}\LARGE\centering
\begin{minipage}{0.95\linewidth}\raggedright\linespread{1.25}\selectfont Développer $(4x+1)^2$.\end{minipage}
\end{minipage}

\vfill
\begin{tcolorbox}[enhanced,nobeforeafter,width=7.1cm,height=2.1cm,valign=center,colback=cnom!6,colframe=cnom,arc=2.5mm,boxrule=0.9pt,left=3mm]{\Large\bfseries\color{cnom}a)}\quad\begin{minipage}[c]{5.3cm}\raggedright\Large $16x^2+1$\end{minipage}\end{tcolorbox}\hfill\begin{tcolorbox}[enhanced,nobeforeafter,width=7.1cm,height=2.1cm,valign=center,colback=cnom!6,colframe=cnom,arc=2.5mm,boxrule=0.9pt,left=3mm]{\Large\bfseries\color{cnom}b)}\quad\begin{minipage}[c]{5.3cm}\raggedright\Large $16x^2+4x+1$\end{minipage}\end{tcolorbox}\par\vspace{2.5mm}
\begin{tcolorbox}[enhanced,nobeforeafter,width=7.1cm,height=2.1cm,valign=center,colback=cnom!6,colframe=cnom,arc=2.5mm,boxrule=0.9pt,left=3mm]{\Large\bfseries\color{cnom}c)}\quad\begin{minipage}[c]{5.3cm}\raggedright\Large $16x^2+8x+1$\end{minipage}\end{tcolorbox}\hfill\begin{tcolorbox}[enhanced,nobeforeafter,width=7.1cm,height=2.1cm,valign=center,colback=cnom!6,colframe=cnom,arc=2.5mm,boxrule=0.9pt,left=3mm]{\Large\bfseries\color{cnom}d)}\quad\begin{minipage}[c]{5.3cm}\raggedright\Large $4x^2+8x+1$\end{minipage}\end{tcolorbox}
\end{tcolorbox}
\newpage
\begin{tcolorbox}[enhanced,width=16cm,height=14cm,colback=white,colframe=cnom,boxrule=2pt,arc=0mm,
left=6mm,right=6mm,top=0mm,bottom=4mm,sharp corners,
overlay={\fill[cnom] (frame.north west) rectangle ([yshift=-1.25cm]frame.north east);
\node[anchor=west,text=white,font=\sffamily\bfseries\LARGE] at ([xshift=5mm,yshift=-0.62cm]frame.north west) {Question 6 D};
\node[anchor=east,text=white,font=\sffamily\itshape\large] at ([xshift=-5mm,yshift=-0.62cm]frame.north east) {Nombres et calculs};}]
\vspace*{1.75cm}
\begin{minipage}[t][6.9cm][c]{\linewidth}\LARGE\centering
\begin{minipage}{0.95\linewidth}\raggedright\linespread{1.25}\selectfont Développer $(3x+4)^2$.\end{minipage}
\end{minipage}

\vfill
\begin{tcolorbox}[enhanced,nobeforeafter,width=7.1cm,height=2.1cm,valign=center,colback=cnom!6,colframe=cnom,arc=2.5mm,boxrule=0.9pt,left=3mm]{\Large\bfseries\color{cnom}a)}\quad\begin{minipage}[c]{5.3cm}\raggedright\Large $9x^2+16$\end{minipage}\end{tcolorbox}\hfill\begin{tcolorbox}[enhanced,nobeforeafter,width=7.1cm,height=2.1cm,valign=center,colback=cnom!6,colframe=cnom,arc=2.5mm,boxrule=0.9pt,left=3mm]{\Large\bfseries\color{cnom}b)}\quad\begin{minipage}[c]{5.3cm}\raggedright\Large $3x^2+24x+16$\end{minipage}\end{tcolorbox}\par\vspace{2.5mm}
\begin{tcolorbox}[enhanced,nobeforeafter,width=7.1cm,height=2.1cm,valign=center,colback=cnom!6,colframe=cnom,arc=2.5mm,boxrule=0.9pt,left=3mm]{\Large\bfseries\color{cnom}c)}\quad\begin{minipage}[c]{5.3cm}\raggedright\Large $9x^2+24x+16$\end{minipage}\end{tcolorbox}\hfill\begin{tcolorbox}[enhanced,nobeforeafter,width=7.1cm,height=2.1cm,valign=center,colback=cnom!6,colframe=cnom,arc=2.5mm,boxrule=0.9pt,left=3mm]{\Large\bfseries\color{cnom}d)}\quad\begin{minipage}[c]{5.3cm}\raggedright\Large $9x^2+12x+16$\end{minipage}\end{tcolorbox}
\end{tcolorbox}
\newpage
\begin{tcolorbox}[enhanced,width=16cm,height=14cm,colback=white,colframe=cnom,boxrule=2pt,arc=0mm,
left=6mm,right=6mm,top=0mm,bottom=4mm,sharp corners,
overlay={\fill[cnom] (frame.north west) rectangle ([yshift=-1.25cm]frame.north east);
\node[anchor=west,text=white,font=\sffamily\bfseries\LARGE] at ([xshift=5mm,yshift=-0.62cm]frame.north west) {Question 7 D};
\node[anchor=east,text=white,font=\sffamily\itshape\large] at ([xshift=-5mm,yshift=-0.62cm]frame.north east) {Nombres et calculs};}]
\vspace*{1.75cm}
\begin{minipage}[t][6.9cm][c]{\linewidth}\LARGE\centering
\begin{minipage}{0.95\linewidth}\raggedright\linespread{1.25}\selectfont Factoriser $16x^2-9$.\end{minipage}
\end{minipage}

\vfill
\begin{tcolorbox}[enhanced,nobeforeafter,width=7.1cm,height=2.1cm,valign=center,colback=cnom!6,colframe=cnom,arc=2.5mm,boxrule=0.9pt,left=3mm]{\Large\bfseries\color{cnom}a)}\quad\begin{minipage}[c]{5.3cm}\raggedright\Large $(16x-3)(16x+3)$\end{minipage}\end{tcolorbox}\hfill\begin{tcolorbox}[enhanced,nobeforeafter,width=7.1cm,height=2.1cm,valign=center,colback=cnom!6,colframe=cnom,arc=2.5mm,boxrule=0.9pt,left=3mm]{\Large\bfseries\color{cnom}b)}\quad\begin{minipage}[c]{5.3cm}\raggedright\Large $(4x-3)(4x+3)$\end{minipage}\end{tcolorbox}\par\vspace{2.5mm}
\begin{tcolorbox}[enhanced,nobeforeafter,width=7.1cm,height=2.1cm,valign=center,colback=cnom!6,colframe=cnom,arc=2.5mm,boxrule=0.9pt,left=3mm]{\Large\bfseries\color{cnom}c)}\quad\begin{minipage}[c]{5.3cm}\raggedright\Large $(4x+3)^2$\end{minipage}\end{tcolorbox}\hfill\begin{tcolorbox}[enhanced,nobeforeafter,width=7.1cm,height=2.1cm,valign=center,colback=cnom!6,colframe=cnom,arc=2.5mm,boxrule=0.9pt,left=3mm]{\Large\bfseries\color{cnom}d)}\quad\begin{minipage}[c]{5.3cm}\raggedright\Large $(4x-3)^2$\end{minipage}\end{tcolorbox}
\end{tcolorbox}
\newpage
\begin{tcolorbox}[enhanced,width=16cm,height=14cm,colback=white,colframe=cnom,boxrule=2pt,arc=0mm,
left=6mm,right=6mm,top=0mm,bottom=4mm,sharp corners,
overlay={\fill[cnom] (frame.north west) rectangle ([yshift=-1.25cm]frame.north east);
\node[anchor=west,text=white,font=\sffamily\bfseries\LARGE] at ([xshift=5mm,yshift=-0.62cm]frame.north west) {Question 8 D};
\node[anchor=east,text=white,font=\sffamily\itshape\large] at ([xshift=-5mm,yshift=-0.62cm]frame.north east) {Nombres et calculs};}]
\vspace*{1.75cm}
\begin{minipage}[t][6.9cm][c]{\linewidth}\LARGE\centering
\begin{minipage}{0.95\linewidth}\raggedright\linespread{1.25}\selectfont Factoriser $x^2-64$.\end{minipage}
\end{minipage}

\vfill
\begin{tcolorbox}[enhanced,nobeforeafter,width=7.1cm,height=2.1cm,valign=center,colback=cnom!6,colframe=cnom,arc=2.5mm,boxrule=0.9pt,left=3mm]{\Large\bfseries\color{cnom}a)}\quad\begin{minipage}[c]{5.3cm}\raggedright\Large $(x-8)^2$\end{minipage}\end{tcolorbox}\hfill\begin{tcolorbox}[enhanced,nobeforeafter,width=7.1cm,height=2.1cm,valign=center,colback=cnom!6,colframe=cnom,arc=2.5mm,boxrule=0.9pt,left=3mm]{\Large\bfseries\color{cnom}b)}\quad\begin{minipage}[c]{5.3cm}\raggedright\Large $(x+8)^2$\end{minipage}\end{tcolorbox}\par\vspace{2.5mm}
\begin{tcolorbox}[enhanced,nobeforeafter,width=7.1cm,height=2.1cm,valign=center,colback=cnom!6,colframe=cnom,arc=2.5mm,boxrule=0.9pt,left=3mm]{\Large\bfseries\color{cnom}c)}\quad\begin{minipage}[c]{5.3cm}\raggedright\Large $(x-64)(x+64)$\end{minipage}\end{tcolorbox}\hfill\begin{tcolorbox}[enhanced,nobeforeafter,width=7.1cm,height=2.1cm,valign=center,colback=cnom!6,colframe=cnom,arc=2.5mm,boxrule=0.9pt,left=3mm]{\Large\bfseries\color{cnom}d)}\quad\begin{minipage}[c]{5.3cm}\raggedright\Large $(x-8)(x+8)$\end{minipage}\end{tcolorbox}
\end{tcolorbox}
\newpage
\begin{tcolorbox}[enhanced,width=16cm,height=14cm,colback=white,colframe=cnom,boxrule=2pt,arc=0mm,
left=6mm,right=6mm,top=0mm,bottom=4mm,sharp corners,
overlay={\fill[cnom] (frame.north west) rectangle ([yshift=-1.25cm]frame.north east);
\node[anchor=west,text=white,font=\sffamily\bfseries\LARGE] at ([xshift=5mm,yshift=-0.62cm]frame.north west) {Question 9 D};
\node[anchor=east,text=white,font=\sffamily\itshape\large] at ([xshift=-5mm,yshift=-0.62cm]frame.north east) {Nombres et calculs};}]
\vspace*{1.75cm}
\begin{minipage}[t][6.9cm][c]{\linewidth}\LARGE\centering
\begin{minipage}{0.95\linewidth}\raggedright\linespread{1.25}\selectfont Factoriser $x^2-12x+36$.\end{minipage}
\end{minipage}

\vfill
\begin{tcolorbox}[enhanced,nobeforeafter,width=7.1cm,height=2.1cm,valign=center,colback=cnom!6,colframe=cnom,arc=2.5mm,boxrule=0.9pt,left=3mm]{\Large\bfseries\color{cnom}a)}\quad\begin{minipage}[c]{5.3cm}\raggedright\Large $(x-6)^2$\end{minipage}\end{tcolorbox}\hfill\begin{tcolorbox}[enhanced,nobeforeafter,width=7.1cm,height=2.1cm,valign=center,colback=cnom!6,colframe=cnom,arc=2.5mm,boxrule=0.9pt,left=3mm]{\Large\bfseries\color{cnom}b)}\quad\begin{minipage}[c]{5.3cm}\raggedright\Large $(x-12)^2$\end{minipage}\end{tcolorbox}\par\vspace{2.5mm}
\begin{tcolorbox}[enhanced,nobeforeafter,width=7.1cm,height=2.1cm,valign=center,colback=cnom!6,colframe=cnom,arc=2.5mm,boxrule=0.9pt,left=3mm]{\Large\bfseries\color{cnom}c)}\quad\begin{minipage}[c]{5.3cm}\raggedright\Large $(x+6)(x-6)$\end{minipage}\end{tcolorbox}\hfill\begin{tcolorbox}[enhanced,nobeforeafter,width=7.1cm,height=2.1cm,valign=center,colback=cnom!6,colframe=cnom,arc=2.5mm,boxrule=0.9pt,left=3mm]{\Large\bfseries\color{cnom}d)}\quad\begin{minipage}[c]{5.3cm}\raggedright\Large $(x+6)^2$\end{minipage}\end{tcolorbox}
\end{tcolorbox}
\newpage
\begin{tcolorbox}[enhanced,width=16cm,height=14cm,colback=white,colframe=cnom,boxrule=2pt,arc=0mm,
left=6mm,right=6mm,top=0mm,bottom=4mm,sharp corners,
overlay={\fill[cnom] (frame.north west) rectangle ([yshift=-1.25cm]frame.north east);
\node[anchor=west,text=white,font=\sffamily\bfseries\LARGE] at ([xshift=5mm,yshift=-0.62cm]frame.north west) {Question 10 D};
\node[anchor=east,text=white,font=\sffamily\itshape\large] at ([xshift=-5mm,yshift=-0.62cm]frame.north east) {Nombres et calculs};}]
\vspace*{1.75cm}
\begin{minipage}[t][6.9cm][c]{\linewidth}\LARGE\centering
\begin{minipage}{0.95\linewidth}\raggedright\linespread{1.25}\selectfont Factoriser $9x^2+30x+25$.\end{minipage}
\end{minipage}

\vfill
\begin{tcolorbox}[enhanced,nobeforeafter,width=7.1cm,height=2.1cm,valign=center,colback=cnom!6,colframe=cnom,arc=2.5mm,boxrule=0.9pt,left=3mm]{\Large\bfseries\color{cnom}a)}\quad\begin{minipage}[c]{5.3cm}\raggedright\Large $(9x+5)^2$\end{minipage}\end{tcolorbox}\hfill\begin{tcolorbox}[enhanced,nobeforeafter,width=7.1cm,height=2.1cm,valign=center,colback=cnom!6,colframe=cnom,arc=2.5mm,boxrule=0.9pt,left=3mm]{\Large\bfseries\color{cnom}b)}\quad\begin{minipage}[c]{5.3cm}\raggedright\Large $(3x+5)^2$\end{minipage}\end{tcolorbox}\par\vspace{2.5mm}
\begin{tcolorbox}[enhanced,nobeforeafter,width=7.1cm,height=2.1cm,valign=center,colback=cnom!6,colframe=cnom,arc=2.5mm,boxrule=0.9pt,left=3mm]{\Large\bfseries\color{cnom}c)}\quad\begin{minipage}[c]{5.3cm}\raggedright\Large $(3x-5)^2$\end{minipage}\end{tcolorbox}\hfill\begin{tcolorbox}[enhanced,nobeforeafter,width=7.1cm,height=2.1cm,valign=center,colback=cnom!6,colframe=cnom,arc=2.5mm,boxrule=0.9pt,left=3mm]{\Large\bfseries\color{cnom}d)}\quad\begin{minipage}[c]{5.3cm}\raggedright\Large $(3x-5)(3x+5)$\end{minipage}\end{tcolorbox}
\end{tcolorbox}
\newpage
\begin{tcolorbox}[enhanced,width=16cm,height=14cm,colback=white,colframe=cnom,boxrule=2pt,arc=0mm,
left=6mm,right=6mm,top=0mm,bottom=4mm,sharp corners,
overlay={\fill[cnom] (frame.north west) rectangle ([yshift=-1.25cm]frame.north east);
\node[anchor=west,text=white,font=\sffamily\bfseries\LARGE] at ([xshift=5mm,yshift=-0.62cm]frame.north west) {Question 11 D};
\node[anchor=east,text=white,font=\sffamily\itshape\large] at ([xshift=-5mm,yshift=-0.62cm]frame.north east) {Nombres et calculs};}]
\vspace*{1.75cm}
\begin{minipage}[t][6.9cm][c]{\linewidth}\LARGE\centering
\begin{minipage}{0.95\linewidth}\raggedright\linespread{1.25}\selectfont Que signifie $x\in]-\infty\,;\,-3[$~?\end{minipage}
\end{minipage}

\vfill
\begin{tcolorbox}[enhanced,nobeforeafter,width=7.1cm,height=2.1cm,valign=center,colback=cnom!6,colframe=cnom,arc=2.5mm,boxrule=0.9pt,left=3mm]{\Large\bfseries\color{cnom}a)}\quad\begin{minipage}[c]{5.3cm}\raggedright\Large $x<-3$\end{minipage}\end{tcolorbox}\hfill\begin{tcolorbox}[enhanced,nobeforeafter,width=7.1cm,height=2.1cm,valign=center,colback=cnom!6,colframe=cnom,arc=2.5mm,boxrule=0.9pt,left=3mm]{\Large\bfseries\color{cnom}b)}\quad\begin{minipage}[c]{5.3cm}\raggedright\Large $x\leqslant -3$\end{minipage}\end{tcolorbox}\par\vspace{2.5mm}
\begin{tcolorbox}[enhanced,nobeforeafter,width=7.1cm,height=2.1cm,valign=center,colback=cnom!6,colframe=cnom,arc=2.5mm,boxrule=0.9pt,left=3mm]{\Large\bfseries\color{cnom}c)}\quad\begin{minipage}[c]{5.3cm}\raggedright\Large $x>-3$\end{minipage}\end{tcolorbox}\hfill\begin{tcolorbox}[enhanced,nobeforeafter,width=7.1cm,height=2.1cm,valign=center,colback=cnom!6,colframe=cnom,arc=2.5mm,boxrule=0.9pt,left=3mm]{\Large\bfseries\color{cnom}d)}\quad\begin{minipage}[c]{5.3cm}\raggedright\Large $x\geqslant -3$\end{minipage}\end{tcolorbox}
\end{tcolorbox}
\newpage
\begin{tcolorbox}[enhanced,width=16cm,height=14cm,colback=white,colframe=cnom,boxrule=2pt,arc=0mm,
left=6mm,right=6mm,top=0mm,bottom=4mm,sharp corners,
overlay={\fill[cnom] (frame.north west) rectangle ([yshift=-1.25cm]frame.north east);
\node[anchor=west,text=white,font=\sffamily\bfseries\LARGE] at ([xshift=5mm,yshift=-0.62cm]frame.north west) {Question 12 D};
\node[anchor=east,text=white,font=\sffamily\itshape\large] at ([xshift=-5mm,yshift=-0.62cm]frame.north east) {Nombres et calculs};}]
\vspace*{1.75cm}
\begin{minipage}[t][6.9cm][c]{\linewidth}\LARGE\centering
\begin{minipage}{0.95\linewidth}\raggedright\linespread{1.25}\selectfont Que signifie $x\in]2\,;\,9[$~?\end{minipage}
\end{minipage}

\vfill
\begin{tcolorbox}[enhanced,nobeforeafter,width=7.1cm,height=2.1cm,valign=center,colback=cnom!6,colframe=cnom,arc=2.5mm,boxrule=0.9pt,left=3mm]{\Large\bfseries\color{cnom}a)}\quad\begin{minipage}[c]{5.3cm}\raggedright\Large $2<x\leqslant 9$\end{minipage}\end{tcolorbox}\hfill\begin{tcolorbox}[enhanced,nobeforeafter,width=7.1cm,height=2.1cm,valign=center,colback=cnom!6,colframe=cnom,arc=2.5mm,boxrule=0.9pt,left=3mm]{\Large\bfseries\color{cnom}b)}\quad\begin{minipage}[c]{5.3cm}\raggedright\Large $2<x<9$\end{minipage}\end{tcolorbox}\par\vspace{2.5mm}
\begin{tcolorbox}[enhanced,nobeforeafter,width=7.1cm,height=2.1cm,valign=center,colback=cnom!6,colframe=cnom,arc=2.5mm,boxrule=0.9pt,left=3mm]{\Large\bfseries\color{cnom}c)}\quad\begin{minipage}[c]{5.3cm}\raggedright\Large $2\leqslant x<9$\end{minipage}\end{tcolorbox}\hfill\begin{tcolorbox}[enhanced,nobeforeafter,width=7.1cm,height=2.1cm,valign=center,colback=cnom!6,colframe=cnom,arc=2.5mm,boxrule=0.9pt,left=3mm]{\Large\bfseries\color{cnom}d)}\quad\begin{minipage}[c]{5.3cm}\raggedright\Large $2\leqslant x\leqslant 9$\end{minipage}\end{tcolorbox}
\end{tcolorbox}
\newpage
\begin{tcolorbox}[enhanced,width=16cm,height=14cm,colback=white,colframe=cnom,boxrule=2pt,arc=0mm,
left=6mm,right=6mm,top=0mm,bottom=4mm,sharp corners,
overlay={\fill[cnom] (frame.north west) rectangle ([yshift=-1.25cm]frame.north east);
\node[anchor=west,text=white,font=\sffamily\bfseries\LARGE] at ([xshift=5mm,yshift=-0.62cm]frame.north west) {Question 13 D};
\node[anchor=east,text=white,font=\sffamily\itshape\large] at ([xshift=-5mm,yshift=-0.62cm]frame.north east) {Nombres et calculs};}]
\vspace*{1.75cm}
\begin{minipage}[t][6.9cm][c]{\linewidth}\LARGE\centering
\begin{minipage}{0.95\linewidth}\raggedright\linespread{1.25}\selectfont Pour quels réels $x$ a-t-on $|x+2|\leqslant 2$~?\end{minipage}
\end{minipage}

\vfill
\begin{tcolorbox}[enhanced,nobeforeafter,width=7.1cm,height=2.1cm,valign=center,colback=cnom!6,colframe=cnom,arc=2.5mm,boxrule=0.9pt,left=3mm]{\Large\bfseries\color{cnom}a)}\quad\begin{minipage}[c]{5.3cm}\raggedright\Large $x\in[0\,;\,4]$\end{minipage}\end{tcolorbox}\hfill\begin{tcolorbox}[enhanced,nobeforeafter,width=7.1cm,height=2.1cm,valign=center,colback=cnom!6,colframe=cnom,arc=2.5mm,boxrule=0.9pt,left=3mm]{\Large\bfseries\color{cnom}b)}\quad\begin{minipage}[c]{5.3cm}\raggedright\Large $x\in[-6\,;\,2]$\end{minipage}\end{tcolorbox}\par\vspace{2.5mm}
\begin{tcolorbox}[enhanced,nobeforeafter,width=7.1cm,height=2.1cm,valign=center,colback=cnom!6,colframe=cnom,arc=2.5mm,boxrule=0.9pt,left=3mm]{\Large\bfseries\color{cnom}c)}\quad\begin{minipage}[c]{5.3cm}\raggedright\Large $x\in[-2\,;\,0]$\end{minipage}\end{tcolorbox}\hfill\begin{tcolorbox}[enhanced,nobeforeafter,width=7.1cm,height=2.1cm,valign=center,colback=cnom!6,colframe=cnom,arc=2.5mm,boxrule=0.9pt,left=3mm]{\Large\bfseries\color{cnom}d)}\quad\begin{minipage}[c]{5.3cm}\raggedright\Large $x\in[-4\,;\,0]$\end{minipage}\end{tcolorbox}
\end{tcolorbox}
\newpage
\begin{tcolorbox}[enhanced,width=16cm,height=14cm,colback=white,colframe=cnom,boxrule=2pt,arc=0mm,
left=6mm,right=6mm,top=0mm,bottom=4mm,sharp corners,
overlay={\fill[cnom] (frame.north west) rectangle ([yshift=-1.25cm]frame.north east);
\node[anchor=west,text=white,font=\sffamily\bfseries\LARGE] at ([xshift=5mm,yshift=-0.62cm]frame.north west) {Question 14 D};
\node[anchor=east,text=white,font=\sffamily\itshape\large] at ([xshift=-5mm,yshift=-0.62cm]frame.north east) {Nombres et calculs};}]
\vspace*{1.75cm}
\begin{minipage}[t][6.9cm][c]{\linewidth}\LARGE\centering
\begin{minipage}{0.95\linewidth}\raggedright\linespread{1.25}\selectfont Pour quels réels $x$ a-t-on $|x-1|\leqslant 6$~?\end{minipage}
\end{minipage}

\vfill
\begin{tcolorbox}[enhanced,nobeforeafter,width=7.1cm,height=2.1cm,valign=center,colback=cnom!6,colframe=cnom,arc=2.5mm,boxrule=0.9pt,left=3mm]{\Large\bfseries\color{cnom}a)}\quad\begin{minipage}[c]{5.3cm}\raggedright\Large $x\in[-11\,;\,13]$\end{minipage}\end{tcolorbox}\hfill\begin{tcolorbox}[enhanced,nobeforeafter,width=7.1cm,height=2.1cm,valign=center,colback=cnom!6,colframe=cnom,arc=2.5mm,boxrule=0.9pt,left=3mm]{\Large\bfseries\color{cnom}b)}\quad\begin{minipage}[c]{5.3cm}\raggedright\Large $x\in[1\,;\,7]$\end{minipage}\end{tcolorbox}\par\vspace{2.5mm}
\begin{tcolorbox}[enhanced,nobeforeafter,width=7.1cm,height=2.1cm,valign=center,colback=cnom!6,colframe=cnom,arc=2.5mm,boxrule=0.9pt,left=3mm]{\Large\bfseries\color{cnom}c)}\quad\begin{minipage}[c]{5.3cm}\raggedright\Large $x\in[-7\,;\,5]$\end{minipage}\end{tcolorbox}\hfill\begin{tcolorbox}[enhanced,nobeforeafter,width=7.1cm,height=2.1cm,valign=center,colback=cnom!6,colframe=cnom,arc=2.5mm,boxrule=0.9pt,left=3mm]{\Large\bfseries\color{cnom}d)}\quad\begin{minipage}[c]{5.3cm}\raggedright\Large $x\in[-5\,;\,7]$\end{minipage}\end{tcolorbox}
\end{tcolorbox}
\newpage
\begin{tcolorbox}[enhanced,width=16cm,height=14cm,colback=white,colframe=cnom,boxrule=2pt,arc=0mm,
left=6mm,right=6mm,top=0mm,bottom=4mm,sharp corners,
overlay={\fill[cnom] (frame.north west) rectangle ([yshift=-1.25cm]frame.north east);
\node[anchor=west,text=white,font=\sffamily\bfseries\LARGE] at ([xshift=5mm,yshift=-0.62cm]frame.north west) {Question 15 D};
\node[anchor=east,text=white,font=\sffamily\itshape\large] at ([xshift=-5mm,yshift=-0.62cm]frame.north east) {Nombres et calculs};}]
\vspace*{1.75cm}
\begin{minipage}[t][6.9cm][c]{\linewidth}\LARGE\centering
\begin{minipage}{0.95\linewidth}\raggedright\linespread{1.25}\selectfont Écrire $\dfrac{3^{6}\times3^{-2}}{3^{3}}$ sous la forme d'une puissance de $3$.\end{minipage}
\end{minipage}

\vfill
\begin{tcolorbox}[enhanced,nobeforeafter,width=7.1cm,height=2.1cm,valign=center,colback=cnom!6,colframe=cnom,arc=2.5mm,boxrule=0.9pt,left=3mm]{\Large\bfseries\color{cnom}a)}\quad\begin{minipage}[c]{5.3cm}\raggedright\Large $3^{-15}$\end{minipage}\end{tcolorbox}\hfill\begin{tcolorbox}[enhanced,nobeforeafter,width=7.1cm,height=2.1cm,valign=center,colback=cnom!6,colframe=cnom,arc=2.5mm,boxrule=0.9pt,left=3mm]{\Large\bfseries\color{cnom}b)}\quad\begin{minipage}[c]{5.3cm}\raggedright\Large $3^{1}$\end{minipage}\end{tcolorbox}\par\vspace{2.5mm}
\begin{tcolorbox}[enhanced,nobeforeafter,width=7.1cm,height=2.1cm,valign=center,colback=cnom!6,colframe=cnom,arc=2.5mm,boxrule=0.9pt,left=3mm]{\Large\bfseries\color{cnom}c)}\quad\begin{minipage}[c]{5.3cm}\raggedright\Large $3^{7}$\end{minipage}\end{tcolorbox}\hfill\begin{tcolorbox}[enhanced,nobeforeafter,width=7.1cm,height=2.1cm,valign=center,colback=cnom!6,colframe=cnom,arc=2.5mm,boxrule=0.9pt,left=3mm]{\Large\bfseries\color{cnom}d)}\quad\begin{minipage}[c]{5.3cm}\raggedright\Large $3^{5}$\end{minipage}\end{tcolorbox}
\end{tcolorbox}
\newpage
\begin{tcolorbox}[enhanced,width=16cm,height=14cm,colback=white,colframe=cnom,boxrule=2pt,arc=0mm,
left=6mm,right=6mm,top=0mm,bottom=4mm,sharp corners,
overlay={\fill[cnom] (frame.north west) rectangle ([yshift=-1.25cm]frame.north east);
\node[anchor=west,text=white,font=\sffamily\bfseries\LARGE] at ([xshift=5mm,yshift=-0.62cm]frame.north west) {Question 16 D};
\node[anchor=east,text=white,font=\sffamily\itshape\large] at ([xshift=-5mm,yshift=-0.62cm]frame.north east) {Nombres et calculs};}]
\vspace*{1.75cm}
\begin{minipage}[t][6.9cm][c]{\linewidth}\LARGE\centering
\begin{minipage}{0.95\linewidth}\raggedright\linespread{1.25}\selectfont Écrire $\dfrac{10^{-3}\times10^{8}}{10^{2}}$ sous la forme d'une puissance de $10$.\end{minipage}
\end{minipage}

\vfill
\begin{tcolorbox}[enhanced,nobeforeafter,width=7.1cm,height=2.1cm,valign=center,colback=cnom!6,colframe=cnom,arc=2.5mm,boxrule=0.9pt,left=3mm]{\Large\bfseries\color{cnom}a)}\quad\begin{minipage}[c]{5.3cm}\raggedright\Large $10^{7}$\end{minipage}\end{tcolorbox}\hfill\begin{tcolorbox}[enhanced,nobeforeafter,width=7.1cm,height=2.1cm,valign=center,colback=cnom!6,colframe=cnom,arc=2.5mm,boxrule=0.9pt,left=3mm]{\Large\bfseries\color{cnom}b)}\quad\begin{minipage}[c]{5.3cm}\raggedright\Large $10^{-26}$\end{minipage}\end{tcolorbox}\par\vspace{2.5mm}
\begin{tcolorbox}[enhanced,nobeforeafter,width=7.1cm,height=2.1cm,valign=center,colback=cnom!6,colframe=cnom,arc=2.5mm,boxrule=0.9pt,left=3mm]{\Large\bfseries\color{cnom}c)}\quad\begin{minipage}[c]{5.3cm}\raggedright\Large $10^{3}$\end{minipage}\end{tcolorbox}\hfill\begin{tcolorbox}[enhanced,nobeforeafter,width=7.1cm,height=2.1cm,valign=center,colback=cnom!6,colframe=cnom,arc=2.5mm,boxrule=0.9pt,left=3mm]{\Large\bfseries\color{cnom}d)}\quad\begin{minipage}[c]{5.3cm}\raggedright\Large $10^{-13}$\end{minipage}\end{tcolorbox}
\end{tcolorbox}
\newpage
\begin{tcolorbox}[enhanced,width=16cm,height=14cm,colback=white,colframe=csta,boxrule=2pt,arc=0mm,
left=6mm,right=6mm,top=0mm,bottom=4mm,sharp corners,
overlay={\fill[csta] (frame.north west) rectangle ([yshift=-1.25cm]frame.north east);
\node[anchor=west,text=white,font=\sffamily\bfseries\LARGE] at ([xshift=5mm,yshift=-0.62cm]frame.north west) {Question 17 D};
\node[anchor=east,text=white,font=\sffamily\itshape\large] at ([xshift=-5mm,yshift=-0.62cm]frame.north east) {Évolutions et statistiques};}]
\vspace*{1.75cm}
\begin{minipage}[t][6.9cm][c]{\linewidth}\LARGE\centering
\begin{minipage}{0.95\linewidth}\raggedright\linespread{1.25}\selectfont Quel est le coefficient multiplicateur associé à une hausse de $0{,}5\,\%$~?\end{minipage}
\end{minipage}

\vfill
\begin{tcolorbox}[enhanced,nobeforeafter,width=7.1cm,height=2.1cm,valign=center,colback=csta!6,colframe=csta,arc=2.5mm,boxrule=0.9pt,left=3mm]{\Large\bfseries\color{csta}a)}\quad\begin{minipage}[c]{5.3cm}\raggedright\Large $0{,}005$\end{minipage}\end{tcolorbox}\hfill\begin{tcolorbox}[enhanced,nobeforeafter,width=7.1cm,height=2.1cm,valign=center,colback=csta!6,colframe=csta,arc=2.5mm,boxrule=0.9pt,left=3mm]{\Large\bfseries\color{csta}b)}\quad\begin{minipage}[c]{5.3cm}\raggedright\Large $1{,}005$\end{minipage}\end{tcolorbox}\par\vspace{2.5mm}
\begin{tcolorbox}[enhanced,nobeforeafter,width=7.1cm,height=2.1cm,valign=center,colback=csta!6,colframe=csta,arc=2.5mm,boxrule=0.9pt,left=3mm]{\Large\bfseries\color{csta}c)}\quad\begin{minipage}[c]{5.3cm}\raggedright\Large $1{,}05$\end{minipage}\end{tcolorbox}\hfill\begin{tcolorbox}[enhanced,nobeforeafter,width=7.1cm,height=2.1cm,valign=center,colback=csta!6,colframe=csta,arc=2.5mm,boxrule=0.9pt,left=3mm]{\Large\bfseries\color{csta}d)}\quad\begin{minipage}[c]{5.3cm}\raggedright\Large $0{,}995$\end{minipage}\end{tcolorbox}
\end{tcolorbox}
\newpage
\begin{tcolorbox}[enhanced,width=16cm,height=14cm,colback=white,colframe=csta,boxrule=2pt,arc=0mm,
left=6mm,right=6mm,top=0mm,bottom=4mm,sharp corners,
overlay={\fill[csta] (frame.north west) rectangle ([yshift=-1.25cm]frame.north east);
\node[anchor=west,text=white,font=\sffamily\bfseries\LARGE] at ([xshift=5mm,yshift=-0.62cm]frame.north west) {Question 18 D};
\node[anchor=east,text=white,font=\sffamily\itshape\large] at ([xshift=-5mm,yshift=-0.62cm]frame.north east) {Évolutions et statistiques};}]
\vspace*{1.75cm}
\begin{minipage}[t][6.9cm][c]{\linewidth}\LARGE\centering
\begin{minipage}{0.95\linewidth}\raggedright\linespread{1.25}\selectfont Quel est le coefficient multiplicateur associé à une baisse de $45\,\%$~?\end{minipage}
\end{minipage}

\vfill
\begin{tcolorbox}[enhanced,nobeforeafter,width=7.1cm,height=2.1cm,valign=center,colback=csta!6,colframe=csta,arc=2.5mm,boxrule=0.9pt,left=3mm]{\Large\bfseries\color{csta}a)}\quad\begin{minipage}[c]{5.3cm}\raggedright\Large $-3{,}5$\end{minipage}\end{tcolorbox}\hfill\begin{tcolorbox}[enhanced,nobeforeafter,width=7.1cm,height=2.1cm,valign=center,colback=csta!6,colframe=csta,arc=2.5mm,boxrule=0.9pt,left=3mm]{\Large\bfseries\color{csta}b)}\quad\begin{minipage}[c]{5.3cm}\raggedright\Large $1{,}45$\end{minipage}\end{tcolorbox}\par\vspace{2.5mm}
\begin{tcolorbox}[enhanced,nobeforeafter,width=7.1cm,height=2.1cm,valign=center,colback=csta!6,colframe=csta,arc=2.5mm,boxrule=0.9pt,left=3mm]{\Large\bfseries\color{csta}c)}\quad\begin{minipage}[c]{5.3cm}\raggedright\Large $-0{,}45$\end{minipage}\end{tcolorbox}\hfill\begin{tcolorbox}[enhanced,nobeforeafter,width=7.1cm,height=2.1cm,valign=center,colback=csta!6,colframe=csta,arc=2.5mm,boxrule=0.9pt,left=3mm]{\Large\bfseries\color{csta}d)}\quad\begin{minipage}[c]{5.3cm}\raggedright\Large $0{,}55$\end{minipage}\end{tcolorbox}
\end{tcolorbox}
\newpage
\begin{tcolorbox}[enhanced,width=16cm,height=14cm,colback=white,colframe=csta,boxrule=2pt,arc=0mm,
left=6mm,right=6mm,top=0mm,bottom=4mm,sharp corners,
overlay={\fill[csta] (frame.north west) rectangle ([yshift=-1.25cm]frame.north east);
\node[anchor=west,text=white,font=\sffamily\bfseries\LARGE] at ([xshift=5mm,yshift=-0.62cm]frame.north west) {Question 19 D};
\node[anchor=east,text=white,font=\sffamily\itshape\large] at ([xshift=-5mm,yshift=-0.62cm]frame.north east) {Évolutions et statistiques};}]
\vspace*{1.75cm}
\begin{minipage}[t][6.9cm][c]{\linewidth}\LARGE\centering
\begin{minipage}{0.95\linewidth}\raggedright\linespread{1.25}\selectfont Un prix de $90$ € augmente de $50\,\%$. Quel est le nouveau prix~?\end{minipage}
\end{minipage}

\vfill
\begin{tcolorbox}[enhanced,nobeforeafter,width=7.1cm,height=2.1cm,valign=center,colback=csta!6,colframe=csta,arc=2.5mm,boxrule=0.9pt,left=3mm]{\Large\bfseries\color{csta}a)}\quad\begin{minipage}[c]{5.3cm}\raggedright\Large $136~\text{€}$\end{minipage}\end{tcolorbox}\hfill\begin{tcolorbox}[enhanced,nobeforeafter,width=7.1cm,height=2.1cm,valign=center,colback=csta!6,colframe=csta,arc=2.5mm,boxrule=0.9pt,left=3mm]{\Large\bfseries\color{csta}b)}\quad\begin{minipage}[c]{5.3cm}\raggedright\Large $45~\text{€}$\end{minipage}\end{tcolorbox}\par\vspace{2.5mm}
\begin{tcolorbox}[enhanced,nobeforeafter,width=7.1cm,height=2.1cm,valign=center,colback=csta!6,colframe=csta,arc=2.5mm,boxrule=0.9pt,left=3mm]{\Large\bfseries\color{csta}c)}\quad\begin{minipage}[c]{5.3cm}\raggedright\Large $135~\text{€}$\end{minipage}\end{tcolorbox}\hfill\begin{tcolorbox}[enhanced,nobeforeafter,width=7.1cm,height=2.1cm,valign=center,colback=csta!6,colframe=csta,arc=2.5mm,boxrule=0.9pt,left=3mm]{\Large\bfseries\color{csta}d)}\quad\begin{minipage}[c]{5.3cm}\raggedright\Large $140~\text{€}$\end{minipage}\end{tcolorbox}
\end{tcolorbox}
\newpage
\begin{tcolorbox}[enhanced,width=16cm,height=14cm,colback=white,colframe=csta,boxrule=2pt,arc=0mm,
left=6mm,right=6mm,top=0mm,bottom=4mm,sharp corners,
overlay={\fill[csta] (frame.north west) rectangle ([yshift=-1.25cm]frame.north east);
\node[anchor=west,text=white,font=\sffamily\bfseries\LARGE] at ([xshift=5mm,yshift=-0.62cm]frame.north west) {Question 20 D};
\node[anchor=east,text=white,font=\sffamily\itshape\large] at ([xshift=-5mm,yshift=-0.62cm]frame.north east) {Évolutions et statistiques};}]
\vspace*{1.75cm}
\begin{minipage}[t][6.9cm][c]{\linewidth}\LARGE\centering
\begin{minipage}{0.95\linewidth}\raggedright\linespread{1.25}\selectfont Un prix de $12{,}5$ € baisse de $40\,\%$. Quel est le nouveau prix~?\end{minipage}
\end{minipage}

\vfill
\begin{tcolorbox}[enhanced,nobeforeafter,width=7.1cm,height=2.1cm,valign=center,colback=csta!6,colframe=csta,arc=2.5mm,boxrule=0.9pt,left=3mm]{\Large\bfseries\color{csta}a)}\quad\begin{minipage}[c]{5.3cm}\raggedright\Large $5~\text{€}$\end{minipage}\end{tcolorbox}\hfill\begin{tcolorbox}[enhanced,nobeforeafter,width=7.1cm,height=2.1cm,valign=center,colback=csta!6,colframe=csta,arc=2.5mm,boxrule=0.9pt,left=3mm]{\Large\bfseries\color{csta}b)}\quad\begin{minipage}[c]{5.3cm}\raggedright\Large $-27{,}5~\text{€}$\end{minipage}\end{tcolorbox}\par\vspace{2.5mm}
\begin{tcolorbox}[enhanced,nobeforeafter,width=7.1cm,height=2.1cm,valign=center,colback=csta!6,colframe=csta,arc=2.5mm,boxrule=0.9pt,left=3mm]{\Large\bfseries\color{csta}c)}\quad\begin{minipage}[c]{5.3cm}\raggedright\Large $17{,}5~\text{€}$\end{minipage}\end{tcolorbox}\hfill\begin{tcolorbox}[enhanced,nobeforeafter,width=7.1cm,height=2.1cm,valign=center,colback=csta!6,colframe=csta,arc=2.5mm,boxrule=0.9pt,left=3mm]{\Large\bfseries\color{csta}d)}\quad\begin{minipage}[c]{5.3cm}\raggedright\Large $7{,}5~\text{€}$\end{minipage}\end{tcolorbox}
\end{tcolorbox}
\newpage
\begin{tcolorbox}[enhanced,width=16cm,height=14cm,colback=white,colframe=csta,boxrule=2pt,arc=0mm,
left=6mm,right=6mm,top=0mm,bottom=4mm,sharp corners,
overlay={\fill[csta] (frame.north west) rectangle ([yshift=-1.25cm]frame.north east);
\node[anchor=west,text=white,font=\sffamily\bfseries\LARGE] at ([xshift=5mm,yshift=-0.62cm]frame.north west) {Question 21 D};
\node[anchor=east,text=white,font=\sffamily\itshape\large] at ([xshift=-5mm,yshift=-0.62cm]frame.north east) {Évolutions et statistiques};}]
\vspace*{1.75cm}
\begin{minipage}[t][6.9cm][c]{\linewidth}\LARGE\centering
\begin{minipage}{0.95\linewidth}\raggedright\linespread{1.25}\selectfont Une quantité passe de $90$ à $99$. Quel est son taux d'évolution~?\end{minipage}
\end{minipage}

\vfill
\begin{tcolorbox}[enhanced,nobeforeafter,width=7.1cm,height=2.1cm,valign=center,colback=csta!6,colframe=csta,arc=2.5mm,boxrule=0.9pt,left=3mm]{\Large\bfseries\color{csta}a)}\quad\begin{minipage}[c]{5.3cm}\raggedright\Large $+5\,\%$\end{minipage}\end{tcolorbox}\hfill\begin{tcolorbox}[enhanced,nobeforeafter,width=7.1cm,height=2.1cm,valign=center,colback=csta!6,colframe=csta,arc=2.5mm,boxrule=0.9pt,left=3mm]{\Large\bfseries\color{csta}b)}\quad\begin{minipage}[c]{5.3cm}\raggedright\Large $+10\,\%$\end{minipage}\end{tcolorbox}\par\vspace{2.5mm}
\begin{tcolorbox}[enhanced,nobeforeafter,width=7.1cm,height=2.1cm,valign=center,colback=csta!6,colframe=csta,arc=2.5mm,boxrule=0.9pt,left=3mm]{\Large\bfseries\color{csta}c)}\quad\begin{minipage}[c]{5.3cm}\raggedright\Large $+9\,\%$\end{minipage}\end{tcolorbox}\hfill\begin{tcolorbox}[enhanced,nobeforeafter,width=7.1cm,height=2.1cm,valign=center,colback=csta!6,colframe=csta,arc=2.5mm,boxrule=0.9pt,left=3mm]{\Large\bfseries\color{csta}d)}\quad\begin{minipage}[c]{5.3cm}\raggedright\Large $-10\,\%$\end{minipage}\end{tcolorbox}
\end{tcolorbox}
\newpage
\begin{tcolorbox}[enhanced,width=16cm,height=14cm,colback=white,colframe=csta,boxrule=2pt,arc=0mm,
left=6mm,right=6mm,top=0mm,bottom=4mm,sharp corners,
overlay={\fill[csta] (frame.north west) rectangle ([yshift=-1.25cm]frame.north east);
\node[anchor=west,text=white,font=\sffamily\bfseries\LARGE] at ([xshift=5mm,yshift=-0.62cm]frame.north west) {Question 22 D};
\node[anchor=east,text=white,font=\sffamily\itshape\large] at ([xshift=-5mm,yshift=-0.62cm]frame.north east) {Évolutions et statistiques};}]
\vspace*{1.75cm}
\begin{minipage}[t][6.9cm][c]{\linewidth}\LARGE\centering
\begin{minipage}{0.95\linewidth}\raggedright\linespread{1.25}\selectfont Une quantité passe de $1\,500$ à $1\,140$. Quel est son taux d'évolution~?\end{minipage}
\end{minipage}

\vfill
\begin{tcolorbox}[enhanced,nobeforeafter,width=7.1cm,height=2.1cm,valign=center,colback=csta!6,colframe=csta,arc=2.5mm,boxrule=0.9pt,left=3mm]{\Large\bfseries\color{csta}a)}\quad\begin{minipage}[c]{5.3cm}\raggedright\Large $-12\,\%$\end{minipage}\end{tcolorbox}\hfill\begin{tcolorbox}[enhanced,nobeforeafter,width=7.1cm,height=2.1cm,valign=center,colback=csta!6,colframe=csta,arc=2.5mm,boxrule=0.9pt,left=3mm]{\Large\bfseries\color{csta}b)}\quad\begin{minipage}[c]{5.3cm}\raggedright\Large $+24\,\%$\end{minipage}\end{tcolorbox}\par\vspace{2.5mm}
\begin{tcolorbox}[enhanced,nobeforeafter,width=7.1cm,height=2.1cm,valign=center,colback=csta!6,colframe=csta,arc=2.5mm,boxrule=0.9pt,left=3mm]{\Large\bfseries\color{csta}c)}\quad\begin{minipage}[c]{5.3cm}\raggedright\Large $-360\,\%$\end{minipage}\end{tcolorbox}\hfill\begin{tcolorbox}[enhanced,nobeforeafter,width=7.1cm,height=2.1cm,valign=center,colback=csta!6,colframe=csta,arc=2.5mm,boxrule=0.9pt,left=3mm]{\Large\bfseries\color{csta}d)}\quad\begin{minipage}[c]{5.3cm}\raggedright\Large $-24\,\%$\end{minipage}\end{tcolorbox}
\end{tcolorbox}
\newpage
\begin{tcolorbox}[enhanced,width=16cm,height=14cm,colback=white,colframe=csta,boxrule=2pt,arc=0mm,
left=6mm,right=6mm,top=0mm,bottom=4mm,sharp corners,
overlay={\fill[csta] (frame.north west) rectangle ([yshift=-1.25cm]frame.north east);
\node[anchor=west,text=white,font=\sffamily\bfseries\LARGE] at ([xshift=5mm,yshift=-0.62cm]frame.north west) {Question 23 D};
\node[anchor=east,text=white,font=\sffamily\itshape\large] at ([xshift=-5mm,yshift=-0.62cm]frame.north east) {Évolutions et statistiques};}]
\vspace*{1.75cm}
\begin{minipage}[t][6.9cm][c]{\linewidth}\LARGE\centering
\begin{minipage}{0.95\linewidth}\raggedright\linespread{1.25}\selectfont Un prix subit une hausse de $50\,\%$, puis une hausse de $20\,\%$. Quelle est son évolution globale~?\end{minipage}
\end{minipage}

\vfill
\begin{tcolorbox}[enhanced,nobeforeafter,width=7.1cm,height=2.1cm,valign=center,colback=csta!6,colframe=csta,arc=2.5mm,boxrule=0.9pt,left=3mm]{\Large\bfseries\color{csta}a)}\quad\begin{minipage}[c]{5.3cm}\raggedright\Large $+70\,\%$\end{minipage}\end{tcolorbox}\hfill\begin{tcolorbox}[enhanced,nobeforeafter,width=7.1cm,height=2.1cm,valign=center,colback=csta!6,colframe=csta,arc=2.5mm,boxrule=0.9pt,left=3mm]{\Large\bfseries\color{csta}b)}\quad\begin{minipage}[c]{5.3cm}\raggedright\Large $+80\,\%$\end{minipage}\end{tcolorbox}\par\vspace{2.5mm}
\begin{tcolorbox}[enhanced,nobeforeafter,width=7.1cm,height=2.1cm,valign=center,colback=csta!6,colframe=csta,arc=2.5mm,boxrule=0.9pt,left=3mm]{\Large\bfseries\color{csta}c)}\quad\begin{minipage}[c]{5.3cm}\raggedright\Large $+81\,\%$\end{minipage}\end{tcolorbox}\hfill\begin{tcolorbox}[enhanced,nobeforeafter,width=7.1cm,height=2.1cm,valign=center,colback=csta!6,colframe=csta,arc=2.5mm,boxrule=0.9pt,left=3mm]{\Large\bfseries\color{csta}d)}\quad\begin{minipage}[c]{5.3cm}\raggedright\Large $+10\,\%$\end{minipage}\end{tcolorbox}
\end{tcolorbox}
\newpage
\begin{tcolorbox}[enhanced,width=16cm,height=14cm,colback=white,colframe=csta,boxrule=2pt,arc=0mm,
left=6mm,right=6mm,top=0mm,bottom=4mm,sharp corners,
overlay={\fill[csta] (frame.north west) rectangle ([yshift=-1.25cm]frame.north east);
\node[anchor=west,text=white,font=\sffamily\bfseries\LARGE] at ([xshift=5mm,yshift=-0.62cm]frame.north west) {Question 24 D};
\node[anchor=east,text=white,font=\sffamily\itshape\large] at ([xshift=-5mm,yshift=-0.62cm]frame.north east) {Évolutions et statistiques};}]
\vspace*{1.75cm}
\begin{minipage}[t][6.9cm][c]{\linewidth}\LARGE\centering
\begin{minipage}{0.95\linewidth}\raggedright\linespread{1.25}\selectfont Un prix subit une baisse de $10\,\%$, puis une hausse de $20\,\%$. Quelle est son évolution globale~?\end{minipage}
\end{minipage}

\vfill
\begin{tcolorbox}[enhanced,nobeforeafter,width=7.1cm,height=2.1cm,valign=center,colback=csta!6,colframe=csta,arc=2.5mm,boxrule=0.9pt,left=3mm]{\Large\bfseries\color{csta}a)}\quad\begin{minipage}[c]{5.3cm}\raggedright\Large $+9\,\%$\end{minipage}\end{tcolorbox}\hfill\begin{tcolorbox}[enhanced,nobeforeafter,width=7.1cm,height=2.1cm,valign=center,colback=csta!6,colframe=csta,arc=2.5mm,boxrule=0.9pt,left=3mm]{\Large\bfseries\color{csta}b)}\quad\begin{minipage}[c]{5.3cm}\raggedright\Large $+10\,\%$\end{minipage}\end{tcolorbox}\par\vspace{2.5mm}
\begin{tcolorbox}[enhanced,nobeforeafter,width=7.1cm,height=2.1cm,valign=center,colback=csta!6,colframe=csta,arc=2.5mm,boxrule=0.9pt,left=3mm]{\Large\bfseries\color{csta}c)}\quad\begin{minipage}[c]{5.3cm}\raggedright\Large $-2\,\%$\end{minipage}\end{tcolorbox}\hfill\begin{tcolorbox}[enhanced,nobeforeafter,width=7.1cm,height=2.1cm,valign=center,colback=csta!6,colframe=csta,arc=2.5mm,boxrule=0.9pt,left=3mm]{\Large\bfseries\color{csta}d)}\quad\begin{minipage}[c]{5.3cm}\raggedright\Large $+8\,\%$\end{minipage}\end{tcolorbox}
\end{tcolorbox}
\newpage
\begin{tcolorbox}[enhanced,width=16cm,height=14cm,colback=white,colframe=csta,boxrule=2pt,arc=0mm,
left=6mm,right=6mm,top=0mm,bottom=4mm,sharp corners,
overlay={\fill[csta] (frame.north west) rectangle ([yshift=-1.25cm]frame.north east);
\node[anchor=west,text=white,font=\sffamily\bfseries\LARGE] at ([xshift=5mm,yshift=-0.62cm]frame.north west) {Question 25 D};
\node[anchor=east,text=white,font=\sffamily\itshape\large] at ([xshift=-5mm,yshift=-0.62cm]frame.north east) {Évolutions et statistiques};}]
\vspace*{1.75cm}
\begin{minipage}[t][6.9cm][c]{\linewidth}\LARGE\centering
\begin{minipage}{0.95\linewidth}\raggedright\linespread{1.25}\selectfont Après une hausse de $300\,\%$, quelle évolution faut-il appliquer pour revenir au prix de départ~?\end{minipage}
\end{minipage}

\vfill
\begin{tcolorbox}[enhanced,nobeforeafter,width=7.1cm,height=2.1cm,valign=center,colback=csta!6,colframe=csta,arc=2.5mm,boxrule=0.9pt,left=3mm]{\Large\bfseries\color{csta}a)}\quad\begin{minipage}[c]{5.3cm}\raggedright\Large une baisse de $300\,\%$\end{minipage}\end{tcolorbox}\hfill\begin{tcolorbox}[enhanced,nobeforeafter,width=7.1cm,height=2.1cm,valign=center,colback=csta!6,colframe=csta,arc=2.5mm,boxrule=0.9pt,left=3mm]{\Large\bfseries\color{csta}b)}\quad\begin{minipage}[c]{5.3cm}\raggedright\Large une baisse de $75\,\%$\end{minipage}\end{tcolorbox}\par\vspace{2.5mm}
\begin{tcolorbox}[enhanced,nobeforeafter,width=7.1cm,height=2.1cm,valign=center,colback=csta!6,colframe=csta,arc=2.5mm,boxrule=0.9pt,left=3mm]{\Large\bfseries\color{csta}c)}\quad\begin{minipage}[c]{5.3cm}\raggedright\Large une baisse de $85\,\%$\end{minipage}\end{tcolorbox}\hfill\begin{tcolorbox}[enhanced,nobeforeafter,width=7.1cm,height=2.1cm,valign=center,colback=csta!6,colframe=csta,arc=2.5mm,boxrule=0.9pt,left=3mm]{\Large\bfseries\color{csta}d)}\quad\begin{minipage}[c]{5.3cm}\raggedright\Large une hausse de $75\,\%$\end{minipage}\end{tcolorbox}
\end{tcolorbox}
\newpage
\begin{tcolorbox}[enhanced,width=16cm,height=14cm,colback=white,colframe=csta,boxrule=2pt,arc=0mm,
left=6mm,right=6mm,top=0mm,bottom=4mm,sharp corners,
overlay={\fill[csta] (frame.north west) rectangle ([yshift=-1.25cm]frame.north east);
\node[anchor=west,text=white,font=\sffamily\bfseries\LARGE] at ([xshift=5mm,yshift=-0.62cm]frame.north west) {Question 26 D};
\node[anchor=east,text=white,font=\sffamily\itshape\large] at ([xshift=-5mm,yshift=-0.62cm]frame.north east) {Évolutions et statistiques};}]
\vspace*{1.75cm}
\begin{minipage}[t][6.9cm][c]{\linewidth}\LARGE\centering
\begin{minipage}{0.95\linewidth}\raggedright\linespread{1.25}\selectfont Après une baisse de $80\,\%$, quelle évolution faut-il appliquer pour revenir au prix de départ~?\end{minipage}
\end{minipage}

\vfill
\begin{tcolorbox}[enhanced,nobeforeafter,width=7.1cm,height=2.1cm,valign=center,colback=csta!6,colframe=csta,arc=2.5mm,boxrule=0.9pt,left=3mm]{\Large\bfseries\color{csta}a)}\quad\begin{minipage}[c]{5.3cm}\raggedright\Large une hausse de $400\,\%$\end{minipage}\end{tcolorbox}\hfill\begin{tcolorbox}[enhanced,nobeforeafter,width=7.1cm,height=2.1cm,valign=center,colback=csta!6,colframe=csta,arc=2.5mm,boxrule=0.9pt,left=3mm]{\Large\bfseries\color{csta}b)}\quad\begin{minipage}[c]{5.3cm}\raggedright\Large une baisse de $400\,\%$\end{minipage}\end{tcolorbox}\par\vspace{2.5mm}
\begin{tcolorbox}[enhanced,nobeforeafter,width=7.1cm,height=2.1cm,valign=center,colback=csta!6,colframe=csta,arc=2.5mm,boxrule=0.9pt,left=3mm]{\Large\bfseries\color{csta}c)}\quad\begin{minipage}[c]{5.3cm}\raggedright\Large une hausse de $80\,\%$\end{minipage}\end{tcolorbox}\hfill\begin{tcolorbox}[enhanced,nobeforeafter,width=7.1cm,height=2.1cm,valign=center,colback=csta!6,colframe=csta,arc=2.5mm,boxrule=0.9pt,left=3mm]{\Large\bfseries\color{csta}d)}\quad\begin{minipage}[c]{5.3cm}\raggedright\Large une hausse de $410\,\%$\end{minipage}\end{tcolorbox}
\end{tcolorbox}
\newpage
\begin{tcolorbox}[enhanced,width=16cm,height=14cm,colback=white,colframe=csta,boxrule=2pt,arc=0mm,
left=6mm,right=6mm,top=0mm,bottom=4mm,sharp corners,
overlay={\fill[csta] (frame.north west) rectangle ([yshift=-1.25cm]frame.north east);
\node[anchor=west,text=white,font=\sffamily\bfseries\LARGE] at ([xshift=5mm,yshift=-0.62cm]frame.north west) {Question 27 D};
\node[anchor=east,text=white,font=\sffamily\itshape\large] at ([xshift=-5mm,yshift=-0.62cm]frame.north east) {Évolutions et statistiques};}]
\vspace*{1.75cm}
\begin{minipage}[t][6.9cm][c]{\linewidth}\Large\centering
\begin{minipage}{0.95\linewidth}\raggedright\linespread{1.25}\selectfont Dans un lycée, $70\,\%$ des élèves sont internes et $30\,\%$ des internes viennent d'une autre commune. Quelle proportion des élèves du lycée sont des internes venant d'une autre commune~?\end{minipage}
\end{minipage}

\vfill
\begin{tcolorbox}[enhanced,nobeforeafter,width=7.1cm,height=2.1cm,valign=center,colback=csta!6,colframe=csta,arc=2.5mm,boxrule=0.9pt,left=3mm]{\Large\bfseries\color{csta}a)}\quad\begin{minipage}[c]{5.3cm}\raggedright\Large $50\,\%$\end{minipage}\end{tcolorbox}\hfill\begin{tcolorbox}[enhanced,nobeforeafter,width=7.1cm,height=2.1cm,valign=center,colback=csta!6,colframe=csta,arc=2.5mm,boxrule=0.9pt,left=3mm]{\Large\bfseries\color{csta}b)}\quad\begin{minipage}[c]{5.3cm}\raggedright\Large $21\,\%$\end{minipage}\end{tcolorbox}\par\vspace{2.5mm}
\begin{tcolorbox}[enhanced,nobeforeafter,width=7.1cm,height=2.1cm,valign=center,colback=csta!6,colframe=csta,arc=2.5mm,boxrule=0.9pt,left=3mm]{\Large\bfseries\color{csta}c)}\quad\begin{minipage}[c]{5.3cm}\raggedright\Large $100\,\%$\end{minipage}\end{tcolorbox}\hfill\begin{tcolorbox}[enhanced,nobeforeafter,width=7.1cm,height=2.1cm,valign=center,colback=csta!6,colframe=csta,arc=2.5mm,boxrule=0.9pt,left=3mm]{\Large\bfseries\color{csta}d)}\quad\begin{minipage}[c]{5.3cm}\raggedright\Large $40\,\%$\end{minipage}\end{tcolorbox}
\end{tcolorbox}
\newpage
\begin{tcolorbox}[enhanced,width=16cm,height=14cm,colback=white,colframe=csta,boxrule=2pt,arc=0mm,
left=6mm,right=6mm,top=0mm,bottom=4mm,sharp corners,
overlay={\fill[csta] (frame.north west) rectangle ([yshift=-1.25cm]frame.north east);
\node[anchor=west,text=white,font=\sffamily\bfseries\LARGE] at ([xshift=5mm,yshift=-0.62cm]frame.north west) {Question 28 D};
\node[anchor=east,text=white,font=\sffamily\itshape\large] at ([xshift=-5mm,yshift=-0.62cm]frame.north east) {Évolutions et statistiques};}]
\vspace*{1.75cm}
\begin{minipage}[t][6.9cm][c]{\linewidth}\Large\centering
\begin{minipage}{0.95\linewidth}\raggedright\linespread{1.25}\selectfont Dans un lycée, $55\,\%$ des élèves sont internes et $40\,\%$ des internes viennent d'une autre commune. Quelle proportion des élèves du lycée sont des internes venant d'une autre commune~?\end{minipage}
\end{minipage}

\vfill
\begin{tcolorbox}[enhanced,nobeforeafter,width=7.1cm,height=2.1cm,valign=center,colback=csta!6,colframe=csta,arc=2.5mm,boxrule=0.9pt,left=3mm]{\Large\bfseries\color{csta}a)}\quad\begin{minipage}[c]{5.3cm}\raggedright\Large $15\,\%$\end{minipage}\end{tcolorbox}\hfill\begin{tcolorbox}[enhanced,nobeforeafter,width=7.1cm,height=2.1cm,valign=center,colback=csta!6,colframe=csta,arc=2.5mm,boxrule=0.9pt,left=3mm]{\Large\bfseries\color{csta}b)}\quad\begin{minipage}[c]{5.3cm}\raggedright\Large $95\,\%$\end{minipage}\end{tcolorbox}\par\vspace{2.5mm}
\begin{tcolorbox}[enhanced,nobeforeafter,width=7.1cm,height=2.1cm,valign=center,colback=csta!6,colframe=csta,arc=2.5mm,boxrule=0.9pt,left=3mm]{\Large\bfseries\color{csta}c)}\quad\begin{minipage}[c]{5.3cm}\raggedright\Large $47{,}5\,\%$\end{minipage}\end{tcolorbox}\hfill\begin{tcolorbox}[enhanced,nobeforeafter,width=7.1cm,height=2.1cm,valign=center,colback=csta!6,colframe=csta,arc=2.5mm,boxrule=0.9pt,left=3mm]{\Large\bfseries\color{csta}d)}\quad\begin{minipage}[c]{5.3cm}\raggedright\Large $22\,\%$\end{minipage}\end{tcolorbox}
\end{tcolorbox}
\newpage
\begin{tcolorbox}[enhanced,width=16cm,height=14cm,colback=white,colframe=csta,boxrule=2pt,arc=0mm,
left=6mm,right=6mm,top=0mm,bottom=4mm,sharp corners,
overlay={\fill[csta] (frame.north west) rectangle ([yshift=-1.25cm]frame.north east);
\node[anchor=west,text=white,font=\sffamily\bfseries\LARGE] at ([xshift=5mm,yshift=-0.62cm]frame.north west) {Question 29 D};
\node[anchor=east,text=white,font=\sffamily\itshape\large] at ([xshift=-5mm,yshift=-0.62cm]frame.north east) {Évolutions et statistiques};}]
\vspace*{1.75cm}
\begin{minipage}[t][6.9cm][c]{\linewidth}\Large\centering
\begin{minipage}{0.95\linewidth}\raggedright\linespread{1.25}\selectfont Quelle est la moyenne de cette série statistique~? $\begin{array}{|l|c|c|c|}\hline \text{Valeur} & 9 & 11 & 16\\\hline \text{Effectif} & 5 & 5 & 10\\\hline\end{array}$\end{minipage}
\end{minipage}

\vfill
\begin{tcolorbox}[enhanced,nobeforeafter,width=7.1cm,height=2.1cm,valign=center,colback=csta!6,colframe=csta,arc=2.5mm,boxrule=0.9pt,left=3mm]{\Large\bfseries\color{csta}a)}\quad\begin{minipage}[c]{5.3cm}\raggedright\Large $12$\end{minipage}\end{tcolorbox}\hfill\begin{tcolorbox}[enhanced,nobeforeafter,width=7.1cm,height=2.1cm,valign=center,colback=csta!6,colframe=csta,arc=2.5mm,boxrule=0.9pt,left=3mm]{\Large\bfseries\color{csta}b)}\quad\begin{minipage}[c]{5.3cm}\raggedright\Large $14$\end{minipage}\end{tcolorbox}\par\vspace{2.5mm}
\begin{tcolorbox}[enhanced,nobeforeafter,width=7.1cm,height=2.1cm,valign=center,colback=csta!6,colframe=csta,arc=2.5mm,boxrule=0.9pt,left=3mm]{\Large\bfseries\color{csta}c)}\quad\begin{minipage}[c]{5.3cm}\raggedright\Large $13$\end{minipage}\end{tcolorbox}\hfill\begin{tcolorbox}[enhanced,nobeforeafter,width=7.1cm,height=2.1cm,valign=center,colback=csta!6,colframe=csta,arc=2.5mm,boxrule=0.9pt,left=3mm]{\Large\bfseries\color{csta}d)}\quad\begin{minipage}[c]{5.3cm}\raggedright\Large $\dfrac{260}{3}$\end{minipage}\end{tcolorbox}
\end{tcolorbox}
\newpage
\begin{tcolorbox}[enhanced,width=16cm,height=14cm,colback=white,colframe=csta,boxrule=2pt,arc=0mm,
left=6mm,right=6mm,top=0mm,bottom=4mm,sharp corners,
overlay={\fill[csta] (frame.north west) rectangle ([yshift=-1.25cm]frame.north east);
\node[anchor=west,text=white,font=\sffamily\bfseries\LARGE] at ([xshift=5mm,yshift=-0.62cm]frame.north west) {Question 30 D};
\node[anchor=east,text=white,font=\sffamily\itshape\large] at ([xshift=-5mm,yshift=-0.62cm]frame.north east) {Évolutions et statistiques};}]
\vspace*{1.75cm}
\begin{minipage}[t][6.9cm][c]{\linewidth}\Large\centering
\begin{minipage}{0.95\linewidth}\raggedright\linespread{1.25}\selectfont Quelle est la moyenne de cette série statistique~? $\begin{array}{|l|c|c|}\hline \text{Valeur} & 6 & 13\\\hline \text{Effectif} & 2 & 3\\\hline\end{array}$\end{minipage}
\end{minipage}

\vfill
\begin{tcolorbox}[enhanced,nobeforeafter,width=7.1cm,height=2.1cm,valign=center,colback=csta!6,colframe=csta,arc=2.5mm,boxrule=0.9pt,left=3mm]{\Large\bfseries\color{csta}a)}\quad\begin{minipage}[c]{5.3cm}\raggedright\Large $25{,}5$\end{minipage}\end{tcolorbox}\hfill\begin{tcolorbox}[enhanced,nobeforeafter,width=7.1cm,height=2.1cm,valign=center,colback=csta!6,colframe=csta,arc=2.5mm,boxrule=0.9pt,left=3mm]{\Large\bfseries\color{csta}b)}\quad\begin{minipage}[c]{5.3cm}\raggedright\Large $11{,}2$\end{minipage}\end{tcolorbox}\par\vspace{2.5mm}
\begin{tcolorbox}[enhanced,nobeforeafter,width=7.1cm,height=2.1cm,valign=center,colback=csta!6,colframe=csta,arc=2.5mm,boxrule=0.9pt,left=3mm]{\Large\bfseries\color{csta}c)}\quad\begin{minipage}[c]{5.3cm}\raggedright\Large $10{,}2$\end{minipage}\end{tcolorbox}\hfill\begin{tcolorbox}[enhanced,nobeforeafter,width=7.1cm,height=2.1cm,valign=center,colback=csta!6,colframe=csta,arc=2.5mm,boxrule=0.9pt,left=3mm]{\Large\bfseries\color{csta}d)}\quad\begin{minipage}[c]{5.3cm}\raggedright\Large $9{,}5$\end{minipage}\end{tcolorbox}
\end{tcolorbox}
\newpage
\begin{tcolorbox}[enhanced,width=16cm,height=14cm,colback=white,colframe=csta,boxrule=2pt,arc=0mm,
left=6mm,right=6mm,top=0mm,bottom=4mm,sharp corners,
overlay={\fill[csta] (frame.north west) rectangle ([yshift=-1.25cm]frame.north east);
\node[anchor=west,text=white,font=\sffamily\bfseries\LARGE] at ([xshift=5mm,yshift=-0.62cm]frame.north west) {Question 31 D};
\node[anchor=east,text=white,font=\sffamily\itshape\large] at ([xshift=-5mm,yshift=-0.62cm]frame.north east) {Évolutions et statistiques};}]
\vspace*{1.75cm}
\begin{minipage}[t][6.9cm][c]{\linewidth}\LARGE\centering
\begin{minipage}{0.95\linewidth}\raggedright\linespread{1.25}\selectfont Voici une série rangée dans l'ordre croissant~: $20\ ;\ 22\ ;\ 25\ ;\ 26\ ;\ 28\ ;\ 30\ ;\ 31\ ;\ 35\ ;\ 40$. Quel est son premier quartile $Q_1$~?\end{minipage}
\end{minipage}

\vfill
\begin{tcolorbox}[enhanced,nobeforeafter,width=7.1cm,height=2.1cm,valign=center,colback=csta!6,colframe=csta,arc=2.5mm,boxrule=0.9pt,left=3mm]{\Large\bfseries\color{csta}a)}\quad\begin{minipage}[c]{5.3cm}\raggedright\Large $26$\end{minipage}\end{tcolorbox}\hfill\begin{tcolorbox}[enhanced,nobeforeafter,width=7.1cm,height=2.1cm,valign=center,colback=csta!6,colframe=csta,arc=2.5mm,boxrule=0.9pt,left=3mm]{\Large\bfseries\color{csta}b)}\quad\begin{minipage}[c]{5.3cm}\raggedright\Large $31$\end{minipage}\end{tcolorbox}\par\vspace{2.5mm}
\begin{tcolorbox}[enhanced,nobeforeafter,width=7.1cm,height=2.1cm,valign=center,colback=csta!6,colframe=csta,arc=2.5mm,boxrule=0.9pt,left=3mm]{\Large\bfseries\color{csta}c)}\quad\begin{minipage}[c]{5.3cm}\raggedright\Large $25$\end{minipage}\end{tcolorbox}\hfill\begin{tcolorbox}[enhanced,nobeforeafter,width=7.1cm,height=2.1cm,valign=center,colback=csta!6,colframe=csta,arc=2.5mm,boxrule=0.9pt,left=3mm]{\Large\bfseries\color{csta}d)}\quad\begin{minipage}[c]{5.3cm}\raggedright\Large $28$\end{minipage}\end{tcolorbox}
\end{tcolorbox}
\newpage
\begin{tcolorbox}[enhanced,width=16cm,height=14cm,colback=white,colframe=csta,boxrule=2pt,arc=0mm,
left=6mm,right=6mm,top=0mm,bottom=4mm,sharp corners,
overlay={\fill[csta] (frame.north west) rectangle ([yshift=-1.25cm]frame.north east);
\node[anchor=west,text=white,font=\sffamily\bfseries\LARGE] at ([xshift=5mm,yshift=-0.62cm]frame.north west) {Question 32 D};
\node[anchor=east,text=white,font=\sffamily\itshape\large] at ([xshift=-5mm,yshift=-0.62cm]frame.north east) {Évolutions et statistiques};}]
\vspace*{1.75cm}
\begin{minipage}[t][6.9cm][c]{\linewidth}\Large\centering
\begin{minipage}{0.95\linewidth}\raggedright\linespread{1.25}\selectfont Voici une série rangée dans l'ordre croissant~: $4\ ;\ 8\ ;\ 9\ ;\ 12\ ;\ 15\ ;\ 16\ ;\ 18\ ;\ 22\ ;\ 23\ ;\ 27\ ;\ 30$. Quel est son premier quartile $Q_1$~?\end{minipage}
\end{minipage}

\vfill
\begin{tcolorbox}[enhanced,nobeforeafter,width=7.1cm,height=2.1cm,valign=center,colback=csta!6,colframe=csta,arc=2.5mm,boxrule=0.9pt,left=3mm]{\Large\bfseries\color{csta}a)}\quad\begin{minipage}[c]{5.3cm}\raggedright\Large $16$\end{minipage}\end{tcolorbox}\hfill\begin{tcolorbox}[enhanced,nobeforeafter,width=7.1cm,height=2.1cm,valign=center,colback=csta!6,colframe=csta,arc=2.5mm,boxrule=0.9pt,left=3mm]{\Large\bfseries\color{csta}b)}\quad\begin{minipage}[c]{5.3cm}\raggedright\Large $12$\end{minipage}\end{tcolorbox}\par\vspace{2.5mm}
\begin{tcolorbox}[enhanced,nobeforeafter,width=7.1cm,height=2.1cm,valign=center,colback=csta!6,colframe=csta,arc=2.5mm,boxrule=0.9pt,left=3mm]{\Large\bfseries\color{csta}c)}\quad\begin{minipage}[c]{5.3cm}\raggedright\Large $23$\end{minipage}\end{tcolorbox}\hfill\begin{tcolorbox}[enhanced,nobeforeafter,width=7.1cm,height=2.1cm,valign=center,colback=csta!6,colframe=csta,arc=2.5mm,boxrule=0.9pt,left=3mm]{\Large\bfseries\color{csta}d)}\quad\begin{minipage}[c]{5.3cm}\raggedright\Large $9$\end{minipage}\end{tcolorbox}
\end{tcolorbox}
\newpage
\begin{tcolorbox}[enhanced,width=16cm,height=14cm,colback=white,colframe=cfon,boxrule=2pt,arc=0mm,
left=6mm,right=6mm,top=0mm,bottom=4mm,sharp corners,
overlay={\fill[cfon] (frame.north west) rectangle ([yshift=-1.25cm]frame.north east);
\node[anchor=west,text=white,font=\sffamily\bfseries\LARGE] at ([xshift=5mm,yshift=-0.62cm]frame.north west) {Question 33 D};
\node[anchor=east,text=white,font=\sffamily\itshape\large] at ([xshift=-5mm,yshift=-0.62cm]frame.north east) {Fonctions};}]
\vspace*{1.75cm}
\begin{minipage}[t][6.9cm][c]{\linewidth}\LARGE\centering
\begin{minipage}{0.95\linewidth}\raggedright\linespread{1.25}\selectfont Soit $f(x)=4x^2+x-2$. Que vaut $f(-1)$~?\end{minipage}
\end{minipage}

\vfill
\begin{tcolorbox}[enhanced,nobeforeafter,width=7.1cm,height=2.1cm,valign=center,colback=cfon!6,colframe=cfon,arc=2.5mm,boxrule=0.9pt,left=3mm]{\Large\bfseries\color{cfon}a)}\quad\begin{minipage}[c]{5.3cm}\raggedright\Large $-1$\end{minipage}\end{tcolorbox}\hfill\begin{tcolorbox}[enhanced,nobeforeafter,width=7.1cm,height=2.1cm,valign=center,colback=cfon!6,colframe=cfon,arc=2.5mm,boxrule=0.9pt,left=3mm]{\Large\bfseries\color{cfon}b)}\quad\begin{minipage}[c]{5.3cm}\raggedright\Large $-7$\end{minipage}\end{tcolorbox}\par\vspace{2.5mm}
\begin{tcolorbox}[enhanced,nobeforeafter,width=7.1cm,height=2.1cm,valign=center,colback=cfon!6,colframe=cfon,arc=2.5mm,boxrule=0.9pt,left=3mm]{\Large\bfseries\color{cfon}c)}\quad\begin{minipage}[c]{5.3cm}\raggedright\Large $-11$\end{minipage}\end{tcolorbox}\hfill\begin{tcolorbox}[enhanced,nobeforeafter,width=7.1cm,height=2.1cm,valign=center,colback=cfon!6,colframe=cfon,arc=2.5mm,boxrule=0.9pt,left=3mm]{\Large\bfseries\color{cfon}d)}\quad\begin{minipage}[c]{5.3cm}\raggedright\Large $1$\end{minipage}\end{tcolorbox}
\end{tcolorbox}
\newpage
\begin{tcolorbox}[enhanced,width=16cm,height=14cm,colback=white,colframe=cfon,boxrule=2pt,arc=0mm,
left=6mm,right=6mm,top=0mm,bottom=4mm,sharp corners,
overlay={\fill[cfon] (frame.north west) rectangle ([yshift=-1.25cm]frame.north east);
\node[anchor=west,text=white,font=\sffamily\bfseries\LARGE] at ([xshift=5mm,yshift=-0.62cm]frame.north west) {Question 34 D};
\node[anchor=east,text=white,font=\sffamily\itshape\large] at ([xshift=-5mm,yshift=-0.62cm]frame.north east) {Fonctions};}]
\vspace*{1.75cm}
\begin{minipage}[t][6.9cm][c]{\linewidth}\LARGE\centering
\begin{minipage}{0.95\linewidth}\raggedright\linespread{1.25}\selectfont Soit $f(x)=2x^2+5x-3$. Que vaut $f(-4)$~?\end{minipage}
\end{minipage}

\vfill
\begin{tcolorbox}[enhanced,nobeforeafter,width=7.1cm,height=2.1cm,valign=center,colback=cfon!6,colframe=cfon,arc=2.5mm,boxrule=0.9pt,left=3mm]{\Large\bfseries\color{cfon}a)}\quad\begin{minipage}[c]{5.3cm}\raggedright\Large $-9$\end{minipage}\end{tcolorbox}\hfill\begin{tcolorbox}[enhanced,nobeforeafter,width=7.1cm,height=2.1cm,valign=center,colback=cfon!6,colframe=cfon,arc=2.5mm,boxrule=0.9pt,left=3mm]{\Large\bfseries\color{cfon}b)}\quad\begin{minipage}[c]{5.3cm}\raggedright\Large $-39$\end{minipage}\end{tcolorbox}\par\vspace{2.5mm}
\begin{tcolorbox}[enhanced,nobeforeafter,width=7.1cm,height=2.1cm,valign=center,colback=cfon!6,colframe=cfon,arc=2.5mm,boxrule=0.9pt,left=3mm]{\Large\bfseries\color{cfon}c)}\quad\begin{minipage}[c]{5.3cm}\raggedright\Large $9$\end{minipage}\end{tcolorbox}\hfill\begin{tcolorbox}[enhanced,nobeforeafter,width=7.1cm,height=2.1cm,valign=center,colback=cfon!6,colframe=cfon,arc=2.5mm,boxrule=0.9pt,left=3mm]{\Large\bfseries\color{cfon}d)}\quad\begin{minipage}[c]{5.3cm}\raggedright\Large $-55$\end{minipage}\end{tcolorbox}
\end{tcolorbox}
\newpage
\begin{tcolorbox}[enhanced,width=16cm,height=14cm,colback=white,colframe=cfon,boxrule=2pt,arc=0mm,
left=6mm,right=6mm,top=0mm,bottom=4mm,sharp corners,
overlay={\fill[cfon] (frame.north west) rectangle ([yshift=-1.25cm]frame.north east);
\node[anchor=west,text=white,font=\sffamily\bfseries\LARGE] at ([xshift=5mm,yshift=-0.62cm]frame.north west) {Question 35 D};
\node[anchor=east,text=white,font=\sffamily\itshape\large] at ([xshift=-5mm,yshift=-0.62cm]frame.north east) {Fonctions};}]
\vspace*{1.75cm}
\begin{minipage}[t][6.9cm][c]{\linewidth}\LARGE\centering
\begin{minipage}{0.95\linewidth}\raggedright\linespread{1.25}\selectfont Soit $f(x)=6x+1$. Quel est l'antécédent de $4$ par $f$~?\end{minipage}
\end{minipage}

\vfill
\begin{tcolorbox}[enhanced,nobeforeafter,width=7.1cm,height=2.1cm,valign=center,colback=cfon!6,colframe=cfon,arc=2.5mm,boxrule=0.9pt,left=3mm]{\Large\bfseries\color{cfon}a)}\quad\begin{minipage}[c]{5.3cm}\raggedright\Large $25$\end{minipage}\end{tcolorbox}\hfill\begin{tcolorbox}[enhanced,nobeforeafter,width=7.1cm,height=2.1cm,valign=center,colback=cfon!6,colframe=cfon,arc=2.5mm,boxrule=0.9pt,left=3mm]{\Large\bfseries\color{cfon}b)}\quad\begin{minipage}[c]{5.3cm}\raggedright\Large $18$\end{minipage}\end{tcolorbox}\par\vspace{2.5mm}
\begin{tcolorbox}[enhanced,nobeforeafter,width=7.1cm,height=2.1cm,valign=center,colback=cfon!6,colframe=cfon,arc=2.5mm,boxrule=0.9pt,left=3mm]{\Large\bfseries\color{cfon}c)}\quad\begin{minipage}[c]{5.3cm}\raggedright\Large $\dfrac{5}{6}$\end{minipage}\end{tcolorbox}\hfill\begin{tcolorbox}[enhanced,nobeforeafter,width=7.1cm,height=2.1cm,valign=center,colback=cfon!6,colframe=cfon,arc=2.5mm,boxrule=0.9pt,left=3mm]{\Large\bfseries\color{cfon}d)}\quad\begin{minipage}[c]{5.3cm}\raggedright\Large $\dfrac{1}{2}$\end{minipage}\end{tcolorbox}
\end{tcolorbox}
\newpage
\begin{tcolorbox}[enhanced,width=16cm,height=14cm,colback=white,colframe=cfon,boxrule=2pt,arc=0mm,
left=6mm,right=6mm,top=0mm,bottom=4mm,sharp corners,
overlay={\fill[cfon] (frame.north west) rectangle ([yshift=-1.25cm]frame.north east);
\node[anchor=west,text=white,font=\sffamily\bfseries\LARGE] at ([xshift=5mm,yshift=-0.62cm]frame.north west) {Question 36 D};
\node[anchor=east,text=white,font=\sffamily\itshape\large] at ([xshift=-5mm,yshift=-0.62cm]frame.north east) {Fonctions};}]
\vspace*{1.75cm}
\begin{minipage}[t][6.9cm][c]{\linewidth}\LARGE\centering
\begin{minipage}{0.95\linewidth}\raggedright\linespread{1.25}\selectfont Soit $f(x)=-5x-2$. Quel est l'antécédent de $8$ par $f$~?\end{minipage}
\end{minipage}

\vfill
\begin{tcolorbox}[enhanced,nobeforeafter,width=7.1cm,height=2.1cm,valign=center,colback=cfon!6,colframe=cfon,arc=2.5mm,boxrule=0.9pt,left=3mm]{\Large\bfseries\color{cfon}a)}\quad\begin{minipage}[c]{5.3cm}\raggedright\Large $-2$\end{minipage}\end{tcolorbox}\hfill\begin{tcolorbox}[enhanced,nobeforeafter,width=7.1cm,height=2.1cm,valign=center,colback=cfon!6,colframe=cfon,arc=2.5mm,boxrule=0.9pt,left=3mm]{\Large\bfseries\color{cfon}b)}\quad\begin{minipage}[c]{5.3cm}\raggedright\Large $-42$\end{minipage}\end{tcolorbox}\par\vspace{2.5mm}
\begin{tcolorbox}[enhanced,nobeforeafter,width=7.1cm,height=2.1cm,valign=center,colback=cfon!6,colframe=cfon,arc=2.5mm,boxrule=0.9pt,left=3mm]{\Large\bfseries\color{cfon}c)}\quad\begin{minipage}[c]{5.3cm}\raggedright\Large $-\dfrac{6}{5}$\end{minipage}\end{tcolorbox}\hfill\begin{tcolorbox}[enhanced,nobeforeafter,width=7.1cm,height=2.1cm,valign=center,colback=cfon!6,colframe=cfon,arc=2.5mm,boxrule=0.9pt,left=3mm]{\Large\bfseries\color{cfon}d)}\quad\begin{minipage}[c]{5.3cm}\raggedright\Large $-50$\end{minipage}\end{tcolorbox}
\end{tcolorbox}
\newpage
\begin{tcolorbox}[enhanced,width=16cm,height=14cm,colback=white,colframe=cfon,boxrule=2pt,arc=0mm,
left=6mm,right=6mm,top=0mm,bottom=4mm,sharp corners,
overlay={\fill[cfon] (frame.north west) rectangle ([yshift=-1.25cm]frame.north east);
\node[anchor=west,text=white,font=\sffamily\bfseries\LARGE] at ([xshift=5mm,yshift=-0.62cm]frame.north west) {Question 37 D};
\node[anchor=east,text=white,font=\sffamily\itshape\large] at ([xshift=-5mm,yshift=-0.62cm]frame.north east) {Fonctions};}]
\vspace*{1.75cm}
\begin{minipage}[t][6.9cm][c]{\linewidth}\LARGE\centering
\begin{minipage}{0.95\linewidth}\raggedright\linespread{1.25}\selectfont Quel est le coefficient directeur de la droite $(AB)$, avec $A(-1\,;\,0)$ et $B(2\,;\,9)$~?\end{minipage}
\end{minipage}

\vfill
\begin{tcolorbox}[enhanced,nobeforeafter,width=7.1cm,height=2.1cm,valign=center,colback=cfon!6,colframe=cfon,arc=2.5mm,boxrule=0.9pt,left=3mm]{\Large\bfseries\color{cfon}a)}\quad\begin{minipage}[c]{5.3cm}\raggedright\Large $9$\end{minipage}\end{tcolorbox}\hfill\begin{tcolorbox}[enhanced,nobeforeafter,width=7.1cm,height=2.1cm,valign=center,colback=cfon!6,colframe=cfon,arc=2.5mm,boxrule=0.9pt,left=3mm]{\Large\bfseries\color{cfon}b)}\quad\begin{minipage}[c]{5.3cm}\raggedright\Large $-3$\end{minipage}\end{tcolorbox}\par\vspace{2.5mm}
\begin{tcolorbox}[enhanced,nobeforeafter,width=7.1cm,height=2.1cm,valign=center,colback=cfon!6,colframe=cfon,arc=2.5mm,boxrule=0.9pt,left=3mm]{\Large\bfseries\color{cfon}c)}\quad\begin{minipage}[c]{5.3cm}\raggedright\Large $3$\end{minipage}\end{tcolorbox}\hfill\begin{tcolorbox}[enhanced,nobeforeafter,width=7.1cm,height=2.1cm,valign=center,colback=cfon!6,colframe=cfon,arc=2.5mm,boxrule=0.9pt,left=3mm]{\Large\bfseries\color{cfon}d)}\quad\begin{minipage}[c]{5.3cm}\raggedright\Large $\dfrac{1}{3}$\end{minipage}\end{tcolorbox}
\end{tcolorbox}
\newpage
\begin{tcolorbox}[enhanced,width=16cm,height=14cm,colback=white,colframe=cfon,boxrule=2pt,arc=0mm,
left=6mm,right=6mm,top=0mm,bottom=4mm,sharp corners,
overlay={\fill[cfon] (frame.north west) rectangle ([yshift=-1.25cm]frame.north east);
\node[anchor=west,text=white,font=\sffamily\bfseries\LARGE] at ([xshift=5mm,yshift=-0.62cm]frame.north west) {Question 38 D};
\node[anchor=east,text=white,font=\sffamily\itshape\large] at ([xshift=-5mm,yshift=-0.62cm]frame.north east) {Fonctions};}]
\vspace*{1.75cm}
\begin{minipage}[t][6.9cm][c]{\linewidth}\LARGE\centering
\begin{minipage}{0.95\linewidth}\raggedright\linespread{1.25}\selectfont Quel est le coefficient directeur de la droite $(AB)$, avec $A(4\,;\,5)$ et $B(7\,;\,5)$~?\end{minipage}
\end{minipage}

\vfill
\begin{tcolorbox}[enhanced,nobeforeafter,width=7.1cm,height=2.1cm,valign=center,colback=cfon!6,colframe=cfon,arc=2.5mm,boxrule=0.9pt,left=3mm]{\Large\bfseries\color{cfon}a)}\quad\begin{minipage}[c]{5.3cm}\raggedright\Large $1$\end{minipage}\end{tcolorbox}\hfill\begin{tcolorbox}[enhanced,nobeforeafter,width=7.1cm,height=2.1cm,valign=center,colback=cfon!6,colframe=cfon,arc=2.5mm,boxrule=0.9pt,left=3mm]{\Large\bfseries\color{cfon}b)}\quad\begin{minipage}[c]{5.3cm}\raggedright\Large $0$\end{minipage}\end{tcolorbox}\par\vspace{2.5mm}
\begin{tcolorbox}[enhanced,nobeforeafter,width=7.1cm,height=2.1cm,valign=center,colback=cfon!6,colframe=cfon,arc=2.5mm,boxrule=0.9pt,left=3mm]{\Large\bfseries\color{cfon}c)}\quad\begin{minipage}[c]{5.3cm}\raggedright\Large $\dfrac{10}{11}$\end{minipage}\end{tcolorbox}\hfill\begin{tcolorbox}[enhanced,nobeforeafter,width=7.1cm,height=2.1cm,valign=center,colback=cfon!6,colframe=cfon,arc=2.5mm,boxrule=0.9pt,left=3mm]{\Large\bfseries\color{cfon}d)}\quad\begin{minipage}[c]{5.3cm}\raggedright\Large $-1$\end{minipage}\end{tcolorbox}
\end{tcolorbox}
\newpage
\begin{tcolorbox}[enhanced,width=16cm,height=14cm,colback=white,colframe=cfon,boxrule=2pt,arc=0mm,
left=6mm,right=6mm,top=0mm,bottom=4mm,sharp corners,
overlay={\fill[cfon] (frame.north west) rectangle ([yshift=-1.25cm]frame.north east);
\node[anchor=west,text=white,font=\sffamily\bfseries\LARGE] at ([xshift=5mm,yshift=-0.62cm]frame.north west) {Question 39 D};
\node[anchor=east,text=white,font=\sffamily\itshape\large] at ([xshift=-5mm,yshift=-0.62cm]frame.north east) {Fonctions};}]
\vspace*{1.75cm}
\begin{minipage}[t][6.9cm][c]{\linewidth}\LARGE\centering
\begin{minipage}{0.95\linewidth}\raggedright\linespread{1.25}\selectfont Quelle est l'équation réduite de la droite de coefficient directeur $-2$ passant par $A(-2\,;\,1)$~?\end{minipage}
\end{minipage}

\vfill
\begin{tcolorbox}[enhanced,nobeforeafter,width=7.1cm,height=2.1cm,valign=center,colback=cfon!6,colframe=cfon,arc=2.5mm,boxrule=0.9pt,left=3mm]{\Large\bfseries\color{cfon}a)}\quad\begin{minipage}[c]{5.3cm}\raggedright\Large $y=-2x+1$\end{minipage}\end{tcolorbox}\hfill\begin{tcolorbox}[enhanced,nobeforeafter,width=7.1cm,height=2.1cm,valign=center,colback=cfon!6,colframe=cfon,arc=2.5mm,boxrule=0.9pt,left=3mm]{\Large\bfseries\color{cfon}b)}\quad\begin{minipage}[c]{5.3cm}\raggedright\Large $y=-2x-2$\end{minipage}\end{tcolorbox}\par\vspace{2.5mm}
\begin{tcolorbox}[enhanced,nobeforeafter,width=7.1cm,height=2.1cm,valign=center,colback=cfon!6,colframe=cfon,arc=2.5mm,boxrule=0.9pt,left=3mm]{\Large\bfseries\color{cfon}c)}\quad\begin{minipage}[c]{5.3cm}\raggedright\Large $y=-2x+5$\end{minipage}\end{tcolorbox}\hfill\begin{tcolorbox}[enhanced,nobeforeafter,width=7.1cm,height=2.1cm,valign=center,colback=cfon!6,colframe=cfon,arc=2.5mm,boxrule=0.9pt,left=3mm]{\Large\bfseries\color{cfon}d)}\quad\begin{minipage}[c]{5.3cm}\raggedright\Large $y=-2x-3$\end{minipage}\end{tcolorbox}
\end{tcolorbox}
\newpage
\begin{tcolorbox}[enhanced,width=16cm,height=14cm,colback=white,colframe=cfon,boxrule=2pt,arc=0mm,
left=6mm,right=6mm,top=0mm,bottom=4mm,sharp corners,
overlay={\fill[cfon] (frame.north west) rectangle ([yshift=-1.25cm]frame.north east);
\node[anchor=west,text=white,font=\sffamily\bfseries\LARGE] at ([xshift=5mm,yshift=-0.62cm]frame.north west) {Question 40 D};
\node[anchor=east,text=white,font=\sffamily\itshape\large] at ([xshift=-5mm,yshift=-0.62cm]frame.north east) {Fonctions};}]
\vspace*{1.75cm}
\begin{minipage}[t][6.9cm][c]{\linewidth}\LARGE\centering
\begin{minipage}{0.95\linewidth}\raggedright\linespread{1.25}\selectfont Quelle est l'équation réduite de la droite de coefficient directeur $3$ passant par $A(2\,;\,-1)$~?\end{minipage}
\end{minipage}

\vfill
\begin{tcolorbox}[enhanced,nobeforeafter,width=7.1cm,height=2.1cm,valign=center,colback=cfon!6,colframe=cfon,arc=2.5mm,boxrule=0.9pt,left=3mm]{\Large\bfseries\color{cfon}a)}\quad\begin{minipage}[c]{5.3cm}\raggedright\Large $y=3x-7$\end{minipage}\end{tcolorbox}\hfill\begin{tcolorbox}[enhanced,nobeforeafter,width=7.1cm,height=2.1cm,valign=center,colback=cfon!6,colframe=cfon,arc=2.5mm,boxrule=0.9pt,left=3mm]{\Large\bfseries\color{cfon}b)}\quad\begin{minipage}[c]{5.3cm}\raggedright\Large $y=2x-7$\end{minipage}\end{tcolorbox}\par\vspace{2.5mm}
\begin{tcolorbox}[enhanced,nobeforeafter,width=7.1cm,height=2.1cm,valign=center,colback=cfon!6,colframe=cfon,arc=2.5mm,boxrule=0.9pt,left=3mm]{\Large\bfseries\color{cfon}c)}\quad\begin{minipage}[c]{5.3cm}\raggedright\Large $y=3x+5$\end{minipage}\end{tcolorbox}\hfill\begin{tcolorbox}[enhanced,nobeforeafter,width=7.1cm,height=2.1cm,valign=center,colback=cfon!6,colframe=cfon,arc=2.5mm,boxrule=0.9pt,left=3mm]{\Large\bfseries\color{cfon}d)}\quad\begin{minipage}[c]{5.3cm}\raggedright\Large $y=3x-1$\end{minipage}\end{tcolorbox}
\end{tcolorbox}
\newpage
\begin{tcolorbox}[enhanced,width=16cm,height=14cm,colback=white,colframe=cfon,boxrule=2pt,arc=0mm,
left=6mm,right=6mm,top=0mm,bottom=4mm,sharp corners,
overlay={\fill[cfon] (frame.north west) rectangle ([yshift=-1.25cm]frame.north east);
\node[anchor=west,text=white,font=\sffamily\bfseries\LARGE] at ([xshift=5mm,yshift=-0.62cm]frame.north west) {Question 41 D};
\node[anchor=east,text=white,font=\sffamily\itshape\large] at ([xshift=-5mm,yshift=-0.62cm]frame.north east) {Fonctions};}]
\vspace*{1.75cm}
\begin{minipage}[t][6.9cm][c]{\linewidth}\LARGE\centering
\begin{minipage}{0.95\linewidth}\raggedright\linespread{1.25}\selectfont Pour quels réels $x$ a-t-on $3x+1>0$~?\end{minipage}
\end{minipage}

\vfill
\begin{tcolorbox}[enhanced,nobeforeafter,width=7.1cm,height=2.1cm,valign=center,colback=cfon!6,colframe=cfon,arc=2.5mm,boxrule=0.9pt,left=3mm]{\Large\bfseries\color{cfon}a)}\quad\begin{minipage}[c]{5.3cm}\raggedright\Large $x>-\dfrac{1}{3}$\end{minipage}\end{tcolorbox}\hfill\begin{tcolorbox}[enhanced,nobeforeafter,width=7.1cm,height=2.1cm,valign=center,colback=cfon!6,colframe=cfon,arc=2.5mm,boxrule=0.9pt,left=3mm]{\Large\bfseries\color{cfon}b)}\quad\begin{minipage}[c]{5.3cm}\raggedright\Large $x<\dfrac{1}{3}$\end{minipage}\end{tcolorbox}\par\vspace{2.5mm}
\begin{tcolorbox}[enhanced,nobeforeafter,width=7.1cm,height=2.1cm,valign=center,colback=cfon!6,colframe=cfon,arc=2.5mm,boxrule=0.9pt,left=3mm]{\Large\bfseries\color{cfon}c)}\quad\begin{minipage}[c]{5.3cm}\raggedright\Large $x<-\dfrac{1}{3}$\end{minipage}\end{tcolorbox}\hfill\begin{tcolorbox}[enhanced,nobeforeafter,width=7.1cm,height=2.1cm,valign=center,colback=cfon!6,colframe=cfon,arc=2.5mm,boxrule=0.9pt,left=3mm]{\Large\bfseries\color{cfon}d)}\quad\begin{minipage}[c]{5.3cm}\raggedright\Large $x>\dfrac{1}{3}$\end{minipage}\end{tcolorbox}
\end{tcolorbox}
\newpage
\begin{tcolorbox}[enhanced,width=16cm,height=14cm,colback=white,colframe=cfon,boxrule=2pt,arc=0mm,
left=6mm,right=6mm,top=0mm,bottom=4mm,sharp corners,
overlay={\fill[cfon] (frame.north west) rectangle ([yshift=-1.25cm]frame.north east);
\node[anchor=west,text=white,font=\sffamily\bfseries\LARGE] at ([xshift=5mm,yshift=-0.62cm]frame.north west) {Question 42 D};
\node[anchor=east,text=white,font=\sffamily\itshape\large] at ([xshift=-5mm,yshift=-0.62cm]frame.north east) {Fonctions};}]
\vspace*{1.75cm}
\begin{minipage}[t][6.9cm][c]{\linewidth}\LARGE\centering
\begin{minipage}{0.95\linewidth}\raggedright\linespread{1.25}\selectfont Pour quels réels $x$ a-t-on $-x-3>0$~?\end{minipage}
\end{minipage}

\vfill
\begin{tcolorbox}[enhanced,nobeforeafter,width=7.1cm,height=2.1cm,valign=center,colback=cfon!6,colframe=cfon,arc=2.5mm,boxrule=0.9pt,left=3mm]{\Large\bfseries\color{cfon}a)}\quad\begin{minipage}[c]{5.3cm}\raggedright\Large $x>3$\end{minipage}\end{tcolorbox}\hfill\begin{tcolorbox}[enhanced,nobeforeafter,width=7.1cm,height=2.1cm,valign=center,colback=cfon!6,colframe=cfon,arc=2.5mm,boxrule=0.9pt,left=3mm]{\Large\bfseries\color{cfon}b)}\quad\begin{minipage}[c]{5.3cm}\raggedright\Large $x>-3$\end{minipage}\end{tcolorbox}\par\vspace{2.5mm}
\begin{tcolorbox}[enhanced,nobeforeafter,width=7.1cm,height=2.1cm,valign=center,colback=cfon!6,colframe=cfon,arc=2.5mm,boxrule=0.9pt,left=3mm]{\Large\bfseries\color{cfon}c)}\quad\begin{minipage}[c]{5.3cm}\raggedright\Large $x<3$\end{minipage}\end{tcolorbox}\hfill\begin{tcolorbox}[enhanced,nobeforeafter,width=7.1cm,height=2.1cm,valign=center,colback=cfon!6,colframe=cfon,arc=2.5mm,boxrule=0.9pt,left=3mm]{\Large\bfseries\color{cfon}d)}\quad\begin{minipage}[c]{5.3cm}\raggedright\Large $x<-3$\end{minipage}\end{tcolorbox}
\end{tcolorbox}
\newpage
\begin{tcolorbox}[enhanced,width=16cm,height=14cm,colback=white,colframe=cfon,boxrule=2pt,arc=0mm,
left=6mm,right=6mm,top=0mm,bottom=4mm,sharp corners,
overlay={\fill[cfon] (frame.north west) rectangle ([yshift=-1.25cm]frame.north east);
\node[anchor=west,text=white,font=\sffamily\bfseries\LARGE] at ([xshift=5mm,yshift=-0.62cm]frame.north west) {Question 43 D};
\node[anchor=east,text=white,font=\sffamily\itshape\large] at ([xshift=-5mm,yshift=-0.62cm]frame.north east) {Fonctions};}]
\vspace*{1.75cm}
\begin{minipage}[t][6.9cm][c]{\linewidth}\LARGE\centering
\begin{minipage}{0.95\linewidth}\raggedright\linespread{1.25}\selectfont Quelles sont les solutions de l'équation $x^2=-1$~?\end{minipage}
\end{minipage}

\vfill
\begin{tcolorbox}[enhanced,nobeforeafter,width=7.1cm,height=2.1cm,valign=center,colback=cfon!6,colframe=cfon,arc=2.5mm,boxrule=0.9pt,left=3mm]{\Large\bfseries\color{cfon}a)}\quad\begin{minipage}[c]{5.3cm}\raggedright\Large aucune solution\end{minipage}\end{tcolorbox}\hfill\begin{tcolorbox}[enhanced,nobeforeafter,width=7.1cm,height=2.1cm,valign=center,colback=cfon!6,colframe=cfon,arc=2.5mm,boxrule=0.9pt,left=3mm]{\Large\bfseries\color{cfon}b)}\quad\begin{minipage}[c]{5.3cm}\raggedright\Large $\{-1\}$\end{minipage}\end{tcolorbox}\par\vspace{2.5mm}
\begin{tcolorbox}[enhanced,nobeforeafter,width=7.1cm,height=2.1cm,valign=center,colback=cfon!6,colframe=cfon,arc=2.5mm,boxrule=0.9pt,left=3mm]{\Large\bfseries\color{cfon}c)}\quad\begin{minipage}[c]{5.3cm}\raggedright\Large $\{1\}$\end{minipage}\end{tcolorbox}\hfill\begin{tcolorbox}[enhanced,nobeforeafter,width=7.1cm,height=2.1cm,valign=center,colback=cfon!6,colframe=cfon,arc=2.5mm,boxrule=0.9pt,left=3mm]{\Large\bfseries\color{cfon}d)}\quad\begin{minipage}[c]{5.3cm}\raggedright\Large $\{-1\,;\,1\}$\end{minipage}\end{tcolorbox}
\end{tcolorbox}
\newpage
\begin{tcolorbox}[enhanced,width=16cm,height=14cm,colback=white,colframe=cfon,boxrule=2pt,arc=0mm,
left=6mm,right=6mm,top=0mm,bottom=4mm,sharp corners,
overlay={\fill[cfon] (frame.north west) rectangle ([yshift=-1.25cm]frame.north east);
\node[anchor=west,text=white,font=\sffamily\bfseries\LARGE] at ([xshift=5mm,yshift=-0.62cm]frame.north west) {Question 44 D};
\node[anchor=east,text=white,font=\sffamily\itshape\large] at ([xshift=-5mm,yshift=-0.62cm]frame.north east) {Fonctions};}]
\vspace*{1.75cm}
\begin{minipage}[t][6.9cm][c]{\linewidth}\LARGE\centering
\begin{minipage}{0.95\linewidth}\raggedright\linespread{1.25}\selectfont Quelles sont les solutions de l'équation $x^2=12$~?\end{minipage}
\end{minipage}

\vfill
\begin{tcolorbox}[enhanced,nobeforeafter,width=7.1cm,height=2.1cm,valign=center,colback=cfon!6,colframe=cfon,arc=2.5mm,boxrule=0.9pt,left=3mm]{\Large\bfseries\color{cfon}a)}\quad\begin{minipage}[c]{5.3cm}\raggedright\Large aucune solution\end{minipage}\end{tcolorbox}\hfill\begin{tcolorbox}[enhanced,nobeforeafter,width=7.1cm,height=2.1cm,valign=center,colback=cfon!6,colframe=cfon,arc=2.5mm,boxrule=0.9pt,left=3mm]{\Large\bfseries\color{cfon}b)}\quad\begin{minipage}[c]{5.3cm}\raggedright\Large $\{-12\,;\,12\}$\end{minipage}\end{tcolorbox}\par\vspace{2.5mm}
\begin{tcolorbox}[enhanced,nobeforeafter,width=7.1cm,height=2.1cm,valign=center,colback=cfon!6,colframe=cfon,arc=2.5mm,boxrule=0.9pt,left=3mm]{\Large\bfseries\color{cfon}c)}\quad\begin{minipage}[c]{5.3cm}\raggedright\large $\{-2\sqrt{3}\,;\,2\sqrt{3}\}$\end{minipage}\end{tcolorbox}\hfill\begin{tcolorbox}[enhanced,nobeforeafter,width=7.1cm,height=2.1cm,valign=center,colback=cfon!6,colframe=cfon,arc=2.5mm,boxrule=0.9pt,left=3mm]{\Large\bfseries\color{cfon}d)}\quad\begin{minipage}[c]{5.3cm}\raggedright\Large $\{2\sqrt{3}\}$\end{minipage}\end{tcolorbox}
\end{tcolorbox}
\newpage
\begin{tcolorbox}[enhanced,width=16cm,height=14cm,colback=white,colframe=cfon,boxrule=2pt,arc=0mm,
left=6mm,right=6mm,top=0mm,bottom=4mm,sharp corners,
overlay={\fill[cfon] (frame.north west) rectangle ([yshift=-1.25cm]frame.north east);
\node[anchor=west,text=white,font=\sffamily\bfseries\LARGE] at ([xshift=5mm,yshift=-0.62cm]frame.north west) {Question 45 D};
\node[anchor=east,text=white,font=\sffamily\itshape\large] at ([xshift=-5mm,yshift=-0.62cm]frame.north east) {Fonctions};}]
\vspace*{1.75cm}
\begin{minipage}[t][6.9cm][c]{\linewidth}\LARGE\centering
\begin{minipage}{0.95\linewidth}\raggedright\linespread{1.25}\selectfont Quel est le sens de variation de la fonction $f(x)=-4x-4$ sur $\mathbb{R}$~?\end{minipage}
\end{minipage}

\vfill
\begin{tcolorbox}[enhanced,nobeforeafter,width=7.1cm,height=2.1cm,valign=center,colback=cfon!6,colframe=cfon,arc=2.5mm,boxrule=0.9pt,left=3mm]{\Large\bfseries\color{cfon}a)}\quad\begin{minipage}[c]{5.3cm}\raggedright\large croissante sur $\mathbb{R}$\end{minipage}\end{tcolorbox}\hfill\begin{tcolorbox}[enhanced,nobeforeafter,width=7.1cm,height=2.1cm,valign=center,colback=cfon!6,colframe=cfon,arc=2.5mm,boxrule=0.9pt,left=3mm]{\Large\bfseries\color{cfon}b)}\quad\begin{minipage}[c]{5.3cm}\raggedright\large constante sur $\mathbb{R}$\end{minipage}\end{tcolorbox}\par\vspace{2.5mm}
\begin{tcolorbox}[enhanced,nobeforeafter,width=7.1cm,height=2.1cm,valign=center,colback=cfon!6,colframe=cfon,arc=2.5mm,boxrule=0.9pt,left=3mm]{\Large\bfseries\color{cfon}c)}\quad\begin{minipage}[c]{5.3cm}\raggedright\large croissante puis décroissante\end{minipage}\end{tcolorbox}\hfill\begin{tcolorbox}[enhanced,nobeforeafter,width=7.1cm,height=2.1cm,valign=center,colback=cfon!6,colframe=cfon,arc=2.5mm,boxrule=0.9pt,left=3mm]{\Large\bfseries\color{cfon}d)}\quad\begin{minipage}[c]{5.3cm}\raggedright\large décroissante sur $\mathbb{R}$\end{minipage}\end{tcolorbox}
\end{tcolorbox}
\newpage
\begin{tcolorbox}[enhanced,width=16cm,height=14cm,colback=white,colframe=cfon,boxrule=2pt,arc=0mm,
left=6mm,right=6mm,top=0mm,bottom=4mm,sharp corners,
overlay={\fill[cfon] (frame.north west) rectangle ([yshift=-1.25cm]frame.north east);
\node[anchor=west,text=white,font=\sffamily\bfseries\LARGE] at ([xshift=5mm,yshift=-0.62cm]frame.north west) {Question 46 D};
\node[anchor=east,text=white,font=\sffamily\itshape\large] at ([xshift=-5mm,yshift=-0.62cm]frame.north east) {Fonctions};}]
\vspace*{1.75cm}
\begin{minipage}[t][6.9cm][c]{\linewidth}\LARGE\centering
\begin{minipage}{0.95\linewidth}\raggedright\linespread{1.25}\selectfont Quel est le sens de variation de la fonction $f(x)=-2$ sur $\mathbb{R}$~?\end{minipage}
\end{minipage}

\vfill
\begin{tcolorbox}[enhanced,nobeforeafter,width=7.1cm,height=2.1cm,valign=center,colback=cfon!6,colframe=cfon,arc=2.5mm,boxrule=0.9pt,left=3mm]{\Large\bfseries\color{cfon}a)}\quad\begin{minipage}[c]{5.3cm}\raggedright\large constante sur $\mathbb{R}$\end{minipage}\end{tcolorbox}\hfill\begin{tcolorbox}[enhanced,nobeforeafter,width=7.1cm,height=2.1cm,valign=center,colback=cfon!6,colframe=cfon,arc=2.5mm,boxrule=0.9pt,left=3mm]{\Large\bfseries\color{cfon}b)}\quad\begin{minipage}[c]{5.3cm}\raggedright\large décroissante sur $\mathbb{R}$\end{minipage}\end{tcolorbox}\par\vspace{2.5mm}
\begin{tcolorbox}[enhanced,nobeforeafter,width=7.1cm,height=2.1cm,valign=center,colback=cfon!6,colframe=cfon,arc=2.5mm,boxrule=0.9pt,left=3mm]{\Large\bfseries\color{cfon}c)}\quad\begin{minipage}[c]{5.3cm}\raggedright\large croissante sur $\mathbb{R}$\end{minipage}\end{tcolorbox}\hfill\begin{tcolorbox}[enhanced,nobeforeafter,width=7.1cm,height=2.1cm,valign=center,colback=cfon!6,colframe=cfon,arc=2.5mm,boxrule=0.9pt,left=3mm]{\Large\bfseries\color{cfon}d)}\quad\begin{minipage}[c]{5.3cm}\raggedright\large croissante puis décroissante\end{minipage}\end{tcolorbox}
\end{tcolorbox}
\newpage
\begin{tcolorbox}[enhanced,width=16cm,height=14cm,colback=white,colframe=cfon,boxrule=2pt,arc=0mm,
left=6mm,right=6mm,top=0mm,bottom=4mm,sharp corners,
overlay={\fill[cfon] (frame.north west) rectangle ([yshift=-1.25cm]frame.north east);
\node[anchor=west,text=white,font=\sffamily\bfseries\LARGE] at ([xshift=5mm,yshift=-0.62cm]frame.north west) {Question 47 D};
\node[anchor=east,text=white,font=\sffamily\itshape\large] at ([xshift=-5mm,yshift=-0.62cm]frame.north east) {Fonctions};}]
\vspace*{1.75cm}
\begin{minipage}[t][6.9cm][c]{\linewidth}\LARGE\centering
\begin{minipage}{0.95\linewidth}\raggedright\linespread{1.25}\selectfont Quelles sont les solutions de l'inéquation $-2x+1\leqslant x-5$~?\end{minipage}
\end{minipage}

\vfill
\begin{tcolorbox}[enhanced,nobeforeafter,width=7.1cm,height=2.1cm,valign=center,colback=cfon!6,colframe=cfon,arc=2.5mm,boxrule=0.9pt,left=3mm]{\Large\bfseries\color{cfon}a)}\quad\begin{minipage}[c]{5.3cm}\raggedright\Large $x\leqslant \dfrac{4}{3}$\end{minipage}\end{tcolorbox}\hfill\begin{tcolorbox}[enhanced,nobeforeafter,width=7.1cm,height=2.1cm,valign=center,colback=cfon!6,colframe=cfon,arc=2.5mm,boxrule=0.9pt,left=3mm]{\Large\bfseries\color{cfon}b)}\quad\begin{minipage}[c]{5.3cm}\raggedright\Large $x\geqslant 2$\end{minipage}\end{tcolorbox}\par\vspace{2.5mm}
\begin{tcolorbox}[enhanced,nobeforeafter,width=7.1cm,height=2.1cm,valign=center,colback=cfon!6,colframe=cfon,arc=2.5mm,boxrule=0.9pt,left=3mm]{\Large\bfseries\color{cfon}c)}\quad\begin{minipage}[c]{5.3cm}\raggedright\Large $x\leqslant 2$\end{minipage}\end{tcolorbox}\hfill\begin{tcolorbox}[enhanced,nobeforeafter,width=7.1cm,height=2.1cm,valign=center,colback=cfon!6,colframe=cfon,arc=2.5mm,boxrule=0.9pt,left=3mm]{\Large\bfseries\color{cfon}d)}\quad\begin{minipage}[c]{5.3cm}\raggedright\Large $x\geqslant -2$\end{minipage}\end{tcolorbox}
\end{tcolorbox}
\newpage
\begin{tcolorbox}[enhanced,width=16cm,height=14cm,colback=white,colframe=cfon,boxrule=2pt,arc=0mm,
left=6mm,right=6mm,top=0mm,bottom=4mm,sharp corners,
overlay={\fill[cfon] (frame.north west) rectangle ([yshift=-1.25cm]frame.north east);
\node[anchor=west,text=white,font=\sffamily\bfseries\LARGE] at ([xshift=5mm,yshift=-0.62cm]frame.north west) {Question 48 D};
\node[anchor=east,text=white,font=\sffamily\itshape\large] at ([xshift=-5mm,yshift=-0.62cm]frame.north east) {Fonctions};}]
\vspace*{1.75cm}
\begin{minipage}[t][6.9cm][c]{\linewidth}\LARGE\centering
\begin{minipage}{0.95\linewidth}\raggedright\linespread{1.25}\selectfont Quelles sont les solutions de l'inéquation $6x-1\leqslant 2x+7$~?\end{minipage}
\end{minipage}

\vfill
\begin{tcolorbox}[enhanced,nobeforeafter,width=7.1cm,height=2.1cm,valign=center,colback=cfon!6,colframe=cfon,arc=2.5mm,boxrule=0.9pt,left=3mm]{\Large\bfseries\color{cfon}a)}\quad\begin{minipage}[c]{5.3cm}\raggedright\Large $x\geqslant -2$\end{minipage}\end{tcolorbox}\hfill\begin{tcolorbox}[enhanced,nobeforeafter,width=7.1cm,height=2.1cm,valign=center,colback=cfon!6,colframe=cfon,arc=2.5mm,boxrule=0.9pt,left=3mm]{\Large\bfseries\color{cfon}b)}\quad\begin{minipage}[c]{5.3cm}\raggedright\Large $x\geqslant 2$\end{minipage}\end{tcolorbox}\par\vspace{2.5mm}
\begin{tcolorbox}[enhanced,nobeforeafter,width=7.1cm,height=2.1cm,valign=center,colback=cfon!6,colframe=cfon,arc=2.5mm,boxrule=0.9pt,left=3mm]{\Large\bfseries\color{cfon}c)}\quad\begin{minipage}[c]{5.3cm}\raggedright\Large $x\leqslant 2$\end{minipage}\end{tcolorbox}\hfill\begin{tcolorbox}[enhanced,nobeforeafter,width=7.1cm,height=2.1cm,valign=center,colback=cfon!6,colframe=cfon,arc=2.5mm,boxrule=0.9pt,left=3mm]{\Large\bfseries\color{cfon}d)}\quad\begin{minipage}[c]{5.3cm}\raggedright\Large $x\leqslant \dfrac{3}{2}$\end{minipage}\end{tcolorbox}
\end{tcolorbox}
\newpage
\begin{tcolorbox}[enhanced,width=16cm,height=14cm,colback=white,colframe=cvec,boxrule=2pt,arc=0mm,
left=6mm,right=6mm,top=0mm,bottom=4mm,sharp corners,
overlay={\fill[cvec] (frame.north west) rectangle ([yshift=-1.25cm]frame.north east);
\node[anchor=west,text=white,font=\sffamily\bfseries\LARGE] at ([xshift=5mm,yshift=-0.62cm]frame.north west) {Question 49 D};
\node[anchor=east,text=white,font=\sffamily\itshape\large] at ([xshift=-5mm,yshift=-0.62cm]frame.north east) {Vecteurs et repérage};}]
\vspace*{1.75cm}
\begin{minipage}[t][6.9cm][c]{\linewidth}\LARGE\centering
\begin{minipage}{0.95\linewidth}\raggedright\linespread{1.25}\selectfont Dans un repère, $A(2\,;\,-7)$ et $B(7\,;\,-2)$. Quelles sont les coordonnées du vecteur $\overrightarrow{AB}$~?\end{minipage}
\end{minipage}

\vfill
\begin{tcolorbox}[enhanced,nobeforeafter,width=7.1cm,height=2.1cm,valign=center,colback=cvec!6,colframe=cvec,arc=2.5mm,boxrule=0.9pt,left=3mm]{\Large\bfseries\color{cvec}a)}\quad\begin{minipage}[c]{5.3cm}\raggedright\Large $(5\,;\,-5)$\end{minipage}\end{tcolorbox}\hfill\begin{tcolorbox}[enhanced,nobeforeafter,width=7.1cm,height=2.1cm,valign=center,colback=cvec!6,colframe=cvec,arc=2.5mm,boxrule=0.9pt,left=3mm]{\Large\bfseries\color{cvec}b)}\quad\begin{minipage}[c]{5.3cm}\raggedright\Large $(-5\,;\,-5)$\end{minipage}\end{tcolorbox}\par\vspace{2.5mm}
\begin{tcolorbox}[enhanced,nobeforeafter,width=7.1cm,height=2.1cm,valign=center,colback=cvec!6,colframe=cvec,arc=2.5mm,boxrule=0.9pt,left=3mm]{\Large\bfseries\color{cvec}c)}\quad\begin{minipage}[c]{5.3cm}\raggedright\Large $(5\,;\,5)$\end{minipage}\end{tcolorbox}\hfill\begin{tcolorbox}[enhanced,nobeforeafter,width=7.1cm,height=2.1cm,valign=center,colback=cvec!6,colframe=cvec,arc=2.5mm,boxrule=0.9pt,left=3mm]{\Large\bfseries\color{cvec}d)}\quad\begin{minipage}[c]{5.3cm}\raggedright\Large $(9\,;\,-9)$\end{minipage}\end{tcolorbox}
\end{tcolorbox}
\newpage
\begin{tcolorbox}[enhanced,width=16cm,height=14cm,colback=white,colframe=cvec,boxrule=2pt,arc=0mm,
left=6mm,right=6mm,top=0mm,bottom=4mm,sharp corners,
overlay={\fill[cvec] (frame.north west) rectangle ([yshift=-1.25cm]frame.north east);
\node[anchor=west,text=white,font=\sffamily\bfseries\LARGE] at ([xshift=5mm,yshift=-0.62cm]frame.north west) {Question 50 D};
\node[anchor=east,text=white,font=\sffamily\itshape\large] at ([xshift=-5mm,yshift=-0.62cm]frame.north east) {Vecteurs et repérage};}]
\vspace*{1.75cm}
\begin{minipage}[t][6.9cm][c]{\linewidth}\LARGE\centering
\begin{minipage}{0.95\linewidth}\raggedright\linespread{1.25}\selectfont Dans un repère, $A(-4\,;\,0)$ et $B(0\,;\,5)$. Quelles sont les coordonnées du vecteur $\overrightarrow{AB}$~?\end{minipage}
\end{minipage}

\vfill
\begin{tcolorbox}[enhanced,nobeforeafter,width=7.1cm,height=2.1cm,valign=center,colback=cvec!6,colframe=cvec,arc=2.5mm,boxrule=0.9pt,left=3mm]{\Large\bfseries\color{cvec}a)}\quad\begin{minipage}[c]{5.3cm}\raggedright\Large $(-4\,;\,5)$\end{minipage}\end{tcolorbox}\hfill\begin{tcolorbox}[enhanced,nobeforeafter,width=7.1cm,height=2.1cm,valign=center,colback=cvec!6,colframe=cvec,arc=2.5mm,boxrule=0.9pt,left=3mm]{\Large\bfseries\color{cvec}b)}\quad\begin{minipage}[c]{5.3cm}\raggedright\Large $(5\,;\,4)$\end{minipage}\end{tcolorbox}\par\vspace{2.5mm}
\begin{tcolorbox}[enhanced,nobeforeafter,width=7.1cm,height=2.1cm,valign=center,colback=cvec!6,colframe=cvec,arc=2.5mm,boxrule=0.9pt,left=3mm]{\Large\bfseries\color{cvec}c)}\quad\begin{minipage}[c]{5.3cm}\raggedright\Large $(-4\,;\,-5)$\end{minipage}\end{tcolorbox}\hfill\begin{tcolorbox}[enhanced,nobeforeafter,width=7.1cm,height=2.1cm,valign=center,colback=cvec!6,colframe=cvec,arc=2.5mm,boxrule=0.9pt,left=3mm]{\Large\bfseries\color{cvec}d)}\quad\begin{minipage}[c]{5.3cm}\raggedright\Large $(4\,;\,5)$\end{minipage}\end{tcolorbox}
\end{tcolorbox}
\newpage
\begin{tcolorbox}[enhanced,width=16cm,height=14cm,colback=white,colframe=cvec,boxrule=2pt,arc=0mm,
left=6mm,right=6mm,top=0mm,bottom=4mm,sharp corners,
overlay={\fill[cvec] (frame.north west) rectangle ([yshift=-1.25cm]frame.north east);
\node[anchor=west,text=white,font=\sffamily\bfseries\LARGE] at ([xshift=5mm,yshift=-0.62cm]frame.north west) {Question 51 D};
\node[anchor=east,text=white,font=\sffamily\itshape\large] at ([xshift=-5mm,yshift=-0.62cm]frame.north east) {Vecteurs et repérage};}]
\vspace*{1.75cm}
\begin{minipage}[t][6.9cm][c]{\linewidth}\LARGE\centering
\begin{minipage}{0.95\linewidth}\raggedright\linespread{1.25}\selectfont Dans un repère, $A(7\,;\,-2)$ et $B(3\,;\,4)$. Quelles sont les coordonnées du milieu de $[AB]$~?\end{minipage}
\end{minipage}

\vfill
\begin{tcolorbox}[enhanced,nobeforeafter,width=7.1cm,height=2.1cm,valign=center,colback=cvec!6,colframe=cvec,arc=2.5mm,boxrule=0.9pt,left=3mm]{\Large\bfseries\color{cvec}a)}\quad\begin{minipage}[c]{5.3cm}\raggedright\Large $(5\,;\,1)$\end{minipage}\end{tcolorbox}\hfill\begin{tcolorbox}[enhanced,nobeforeafter,width=7.1cm,height=2.1cm,valign=center,colback=cvec!6,colframe=cvec,arc=2.5mm,boxrule=0.9pt,left=3mm]{\Large\bfseries\color{cvec}b)}\quad\begin{minipage}[c]{5.3cm}\raggedright\Large $(-2\,;\,3)$\end{minipage}\end{tcolorbox}\par\vspace{2.5mm}
\begin{tcolorbox}[enhanced,nobeforeafter,width=7.1cm,height=2.1cm,valign=center,colback=cvec!6,colframe=cvec,arc=2.5mm,boxrule=0.9pt,left=3mm]{\Large\bfseries\color{cvec}c)}\quad\begin{minipage}[c]{5.3cm}\raggedright\Large $(5\,;\,-3)$\end{minipage}\end{tcolorbox}\hfill\begin{tcolorbox}[enhanced,nobeforeafter,width=7.1cm,height=2.1cm,valign=center,colback=cvec!6,colframe=cvec,arc=2.5mm,boxrule=0.9pt,left=3mm]{\Large\bfseries\color{cvec}d)}\quad\begin{minipage}[c]{5.3cm}\raggedright\Large $(10\,;\,2)$\end{minipage}\end{tcolorbox}
\end{tcolorbox}
\newpage
\begin{tcolorbox}[enhanced,width=16cm,height=14cm,colback=white,colframe=cvec,boxrule=2pt,arc=0mm,
left=6mm,right=6mm,top=0mm,bottom=4mm,sharp corners,
overlay={\fill[cvec] (frame.north west) rectangle ([yshift=-1.25cm]frame.north east);
\node[anchor=west,text=white,font=\sffamily\bfseries\LARGE] at ([xshift=5mm,yshift=-0.62cm]frame.north west) {Question 52 D};
\node[anchor=east,text=white,font=\sffamily\itshape\large] at ([xshift=-5mm,yshift=-0.62cm]frame.north east) {Vecteurs et repérage};}]
\vspace*{1.75cm}
\begin{minipage}[t][6.9cm][c]{\linewidth}\LARGE\centering
\begin{minipage}{0.95\linewidth}\raggedright\linespread{1.25}\selectfont Dans un repère, $A(-1\,;\,0)$ et $B(6\,;\,5)$. Quelles sont les coordonnées du milieu de $[AB]$~?\end{minipage}
\end{minipage}

\vfill
\begin{tcolorbox}[enhanced,nobeforeafter,width=7.1cm,height=2.1cm,valign=center,colback=cvec!6,colframe=cvec,arc=2.5mm,boxrule=0.9pt,left=3mm]{\Large\bfseries\color{cvec}a)}\quad\begin{minipage}[c]{5.3cm}\raggedright\Large $(2{,}5\,;\,-2{,}5)$\end{minipage}\end{tcolorbox}\hfill\begin{tcolorbox}[enhanced,nobeforeafter,width=7.1cm,height=2.1cm,valign=center,colback=cvec!6,colframe=cvec,arc=2.5mm,boxrule=0.9pt,left=3mm]{\Large\bfseries\color{cvec}b)}\quad\begin{minipage}[c]{5.3cm}\raggedright\Large $(5\,;\,5)$\end{minipage}\end{tcolorbox}\par\vspace{2.5mm}
\begin{tcolorbox}[enhanced,nobeforeafter,width=7.1cm,height=2.1cm,valign=center,colback=cvec!6,colframe=cvec,arc=2.5mm,boxrule=0.9pt,left=3mm]{\Large\bfseries\color{cvec}c)}\quad\begin{minipage}[c]{5.3cm}\raggedright\Large $(3{,}5\,;\,2{,}5)$\end{minipage}\end{tcolorbox}\hfill\begin{tcolorbox}[enhanced,nobeforeafter,width=7.1cm,height=2.1cm,valign=center,colback=cvec!6,colframe=cvec,arc=2.5mm,boxrule=0.9pt,left=3mm]{\Large\bfseries\color{cvec}d)}\quad\begin{minipage}[c]{5.3cm}\raggedright\Large $(2{,}5\,;\,2{,}5)$\end{minipage}\end{tcolorbox}
\end{tcolorbox}
\newpage
\begin{tcolorbox}[enhanced,width=16cm,height=14cm,colback=white,colframe=cvec,boxrule=2pt,arc=0mm,
left=6mm,right=6mm,top=0mm,bottom=4mm,sharp corners,
overlay={\fill[cvec] (frame.north west) rectangle ([yshift=-1.25cm]frame.north east);
\node[anchor=west,text=white,font=\sffamily\bfseries\LARGE] at ([xshift=5mm,yshift=-0.62cm]frame.north west) {Question 53 D};
\node[anchor=east,text=white,font=\sffamily\itshape\large] at ([xshift=-5mm,yshift=-0.62cm]frame.north east) {Vecteurs et repérage};}]
\vspace*{1.75cm}
\begin{minipage}[t][6.9cm][c]{\linewidth}\LARGE\centering
\begin{minipage}{0.95\linewidth}\raggedright\linespread{1.25}\selectfont Dans un repère orthonormé, $A(0\,;\,3)$ et $B(1\,;\,1)$. Quelle est la distance $AB$~?\end{minipage}
\end{minipage}

\vfill
\begin{tcolorbox}[enhanced,nobeforeafter,width=7.1cm,height=2.1cm,valign=center,colback=cvec!6,colframe=cvec,arc=2.5mm,boxrule=0.9pt,left=3mm]{\Large\bfseries\color{cvec}a)}\quad\begin{minipage}[c]{5.3cm}\raggedright\Large $\sqrt{5}$\end{minipage}\end{tcolorbox}\hfill\begin{tcolorbox}[enhanced,nobeforeafter,width=7.1cm,height=2.1cm,valign=center,colback=cvec!6,colframe=cvec,arc=2.5mm,boxrule=0.9pt,left=3mm]{\Large\bfseries\color{cvec}b)}\quad\begin{minipage}[c]{5.3cm}\raggedright\Large $5$\end{minipage}\end{tcolorbox}\par\vspace{2.5mm}
\begin{tcolorbox}[enhanced,nobeforeafter,width=7.1cm,height=2.1cm,valign=center,colback=cvec!6,colframe=cvec,arc=2.5mm,boxrule=0.9pt,left=3mm]{\Large\bfseries\color{cvec}c)}\quad\begin{minipage}[c]{5.3cm}\raggedright\Large $\sqrt{3}$\end{minipage}\end{tcolorbox}\hfill\begin{tcolorbox}[enhanced,nobeforeafter,width=7.1cm,height=2.1cm,valign=center,colback=cvec!6,colframe=cvec,arc=2.5mm,boxrule=0.9pt,left=3mm]{\Large\bfseries\color{cvec}d)}\quad\begin{minipage}[c]{5.3cm}\raggedright\Large $3$\end{minipage}\end{tcolorbox}
\end{tcolorbox}
\newpage
\begin{tcolorbox}[enhanced,width=16cm,height=14cm,colback=white,colframe=cvec,boxrule=2pt,arc=0mm,
left=6mm,right=6mm,top=0mm,bottom=4mm,sharp corners,
overlay={\fill[cvec] (frame.north west) rectangle ([yshift=-1.25cm]frame.north east);
\node[anchor=west,text=white,font=\sffamily\bfseries\LARGE] at ([xshift=5mm,yshift=-0.62cm]frame.north west) {Question 54 D};
\node[anchor=east,text=white,font=\sffamily\itshape\large] at ([xshift=-5mm,yshift=-0.62cm]frame.north east) {Vecteurs et repérage};}]
\vspace*{1.75cm}
\begin{minipage}[t][6.9cm][c]{\linewidth}\LARGE\centering
\begin{minipage}{0.95\linewidth}\raggedright\linespread{1.25}\selectfont Dans un repère orthonormé, $A(6\,;\,-2)$ et $B(1\,;\,10)$. Quelle est la distance $AB$~?\end{minipage}
\end{minipage}

\vfill
\begin{tcolorbox}[enhanced,nobeforeafter,width=7.1cm,height=2.1cm,valign=center,colback=cvec!6,colframe=cvec,arc=2.5mm,boxrule=0.9pt,left=3mm]{\Large\bfseries\color{cvec}a)}\quad\begin{minipage}[c]{5.3cm}\raggedright\Large $169$\end{minipage}\end{tcolorbox}\hfill\begin{tcolorbox}[enhanced,nobeforeafter,width=7.1cm,height=2.1cm,valign=center,colback=cvec!6,colframe=cvec,arc=2.5mm,boxrule=0.9pt,left=3mm]{\Large\bfseries\color{cvec}b)}\quad\begin{minipage}[c]{5.3cm}\raggedright\Large $13$\end{minipage}\end{tcolorbox}\par\vspace{2.5mm}
\begin{tcolorbox}[enhanced,nobeforeafter,width=7.1cm,height=2.1cm,valign=center,colback=cvec!6,colframe=cvec,arc=2.5mm,boxrule=0.9pt,left=3mm]{\Large\bfseries\color{cvec}c)}\quad\begin{minipage}[c]{5.3cm}\raggedright\Large $\sqrt{119}$\end{minipage}\end{tcolorbox}\hfill\begin{tcolorbox}[enhanced,nobeforeafter,width=7.1cm,height=2.1cm,valign=center,colback=cvec!6,colframe=cvec,arc=2.5mm,boxrule=0.9pt,left=3mm]{\Large\bfseries\color{cvec}d)}\quad\begin{minipage}[c]{5.3cm}\raggedright\Large $17$\end{minipage}\end{tcolorbox}
\end{tcolorbox}
\newpage
\begin{tcolorbox}[enhanced,width=16cm,height=14cm,colback=white,colframe=cvec,boxrule=2pt,arc=0mm,
left=6mm,right=6mm,top=0mm,bottom=4mm,sharp corners,
overlay={\fill[cvec] (frame.north west) rectangle ([yshift=-1.25cm]frame.north east);
\node[anchor=west,text=white,font=\sffamily\bfseries\LARGE] at ([xshift=5mm,yshift=-0.62cm]frame.north west) {Question 55 D};
\node[anchor=east,text=white,font=\sffamily\itshape\large] at ([xshift=-5mm,yshift=-0.62cm]frame.north east) {Vecteurs et repérage};}]
\vspace*{1.75cm}
\begin{minipage}[t][6.9cm][c]{\linewidth}\LARGE\centering
\begin{minipage}{0.95\linewidth}\raggedright\linespread{1.25}\selectfont Pour quelle valeur de $k$ les vecteurs $\vec u\begin{pmatrix}-2\\-8\end{pmatrix}$ et $\vec v\begin{pmatrix}k\\12\end{pmatrix}$ sont-ils colinéaires~?\end{minipage}
\end{minipage}

\vfill
\begin{tcolorbox}[enhanced,nobeforeafter,width=7.1cm,height=2.1cm,valign=center,colback=cvec!6,colframe=cvec,arc=2.5mm,boxrule=0.9pt,left=3mm]{\Large\bfseries\color{cvec}a)}\quad\begin{minipage}[c]{5.3cm}\raggedright\Large $k=3$\end{minipage}\end{tcolorbox}\hfill\begin{tcolorbox}[enhanced,nobeforeafter,width=7.1cm,height=2.1cm,valign=center,colback=cvec!6,colframe=cvec,arc=2.5mm,boxrule=0.9pt,left=3mm]{\Large\bfseries\color{cvec}b)}\quad\begin{minipage}[c]{5.3cm}\raggedright\Large $k=-3$\end{minipage}\end{tcolorbox}\par\vspace{2.5mm}
\begin{tcolorbox}[enhanced,nobeforeafter,width=7.1cm,height=2.1cm,valign=center,colback=cvec!6,colframe=cvec,arc=2.5mm,boxrule=0.9pt,left=3mm]{\Large\bfseries\color{cvec}c)}\quad\begin{minipage}[c]{5.3cm}\raggedright\Large $k=48$\end{minipage}\end{tcolorbox}\hfill\begin{tcolorbox}[enhanced,nobeforeafter,width=7.1cm,height=2.1cm,valign=center,colback=cvec!6,colframe=cvec,arc=2.5mm,boxrule=0.9pt,left=3mm]{\Large\bfseries\color{cvec}d)}\quad\begin{minipage}[c]{5.3cm}\raggedright\Large $k=18$\end{minipage}\end{tcolorbox}
\end{tcolorbox}
\newpage
\begin{tcolorbox}[enhanced,width=16cm,height=14cm,colback=white,colframe=cvec,boxrule=2pt,arc=0mm,
left=6mm,right=6mm,top=0mm,bottom=4mm,sharp corners,
overlay={\fill[cvec] (frame.north west) rectangle ([yshift=-1.25cm]frame.north east);
\node[anchor=west,text=white,font=\sffamily\bfseries\LARGE] at ([xshift=5mm,yshift=-0.62cm]frame.north west) {Question 56 D};
\node[anchor=east,text=white,font=\sffamily\itshape\large] at ([xshift=-5mm,yshift=-0.62cm]frame.north east) {Vecteurs et repérage};}]
\vspace*{1.75cm}
\begin{minipage}[t][6.9cm][c]{\linewidth}\LARGE\centering
\begin{minipage}{0.95\linewidth}\raggedright\linespread{1.25}\selectfont Pour quelle valeur de $k$ les vecteurs $\vec u\begin{pmatrix}3\\1\end{pmatrix}$ et $\vec v\begin{pmatrix}k\\7\end{pmatrix}$ sont-ils colinéaires~?\end{minipage}
\end{minipage}

\vfill
\begin{tcolorbox}[enhanced,nobeforeafter,width=7.1cm,height=2.1cm,valign=center,colback=cvec!6,colframe=cvec,arc=2.5mm,boxrule=0.9pt,left=3mm]{\Large\bfseries\color{cvec}a)}\quad\begin{minipage}[c]{5.3cm}\raggedright\Large $k=-21$\end{minipage}\end{tcolorbox}\hfill\begin{tcolorbox}[enhanced,nobeforeafter,width=7.1cm,height=2.1cm,valign=center,colback=cvec!6,colframe=cvec,arc=2.5mm,boxrule=0.9pt,left=3mm]{\Large\bfseries\color{cvec}b)}\quad\begin{minipage}[c]{5.3cm}\raggedright\Large $k=21$\end{minipage}\end{tcolorbox}\par\vspace{2.5mm}
\begin{tcolorbox}[enhanced,nobeforeafter,width=7.1cm,height=2.1cm,valign=center,colback=cvec!6,colframe=cvec,arc=2.5mm,boxrule=0.9pt,left=3mm]{\Large\bfseries\color{cvec}c)}\quad\begin{minipage}[c]{5.3cm}\raggedright\Large $k=9$\end{minipage}\end{tcolorbox}\hfill\begin{tcolorbox}[enhanced,nobeforeafter,width=7.1cm,height=2.1cm,valign=center,colback=cvec!6,colframe=cvec,arc=2.5mm,boxrule=0.9pt,left=3mm]{\Large\bfseries\color{cvec}d)}\quad\begin{minipage}[c]{5.3cm}\raggedright\Large $k=\dfrac{7}{3}$\end{minipage}\end{tcolorbox}
\end{tcolorbox}
\newpage
\begin{tcolorbox}[enhanced,width=16cm,height=14cm,colback=white,colframe=cvec,boxrule=2pt,arc=0mm,
left=6mm,right=6mm,top=0mm,bottom=4mm,sharp corners,
overlay={\fill[cvec] (frame.north west) rectangle ([yshift=-1.25cm]frame.north east);
\node[anchor=west,text=white,font=\sffamily\bfseries\LARGE] at ([xshift=5mm,yshift=-0.62cm]frame.north west) {Question 57 D};
\node[anchor=east,text=white,font=\sffamily\itshape\large] at ([xshift=-5mm,yshift=-0.62cm]frame.north east) {Vecteurs et repérage};}]
\vspace*{1.75cm}
\begin{minipage}[t][6.9cm][c]{\linewidth}\Large\centering
\begin{minipage}{0.95\linewidth}\raggedright\linespread{1.25}\selectfont On donne $\vec u\begin{pmatrix}-3\\1\end{pmatrix}$ et $\vec v\begin{pmatrix}2\\2\end{pmatrix}$. Quelles sont les coordonnées du vecteur $-2\vec u+\vec v$~?\end{minipage}
\end{minipage}

\vfill
\begin{tcolorbox}[enhanced,nobeforeafter,width=7.1cm,height=2.1cm,valign=center,colback=cvec!6,colframe=cvec,arc=2.5mm,boxrule=0.9pt,left=3mm]{\Large\bfseries\color{cvec}a)}\quad\begin{minipage}[c]{5.3cm}\raggedright\Large $(4\,;\,-4)$\end{minipage}\end{tcolorbox}\hfill\begin{tcolorbox}[enhanced,nobeforeafter,width=7.1cm,height=2.1cm,valign=center,colback=cvec!6,colframe=cvec,arc=2.5mm,boxrule=0.9pt,left=3mm]{\Large\bfseries\color{cvec}b)}\quad\begin{minipage}[c]{5.3cm}\raggedright\Large $(8\,;\,0)$\end{minipage}\end{tcolorbox}\par\vspace{2.5mm}
\begin{tcolorbox}[enhanced,nobeforeafter,width=7.1cm,height=2.1cm,valign=center,colback=cvec!6,colframe=cvec,arc=2.5mm,boxrule=0.9pt,left=3mm]{\Large\bfseries\color{cvec}c)}\quad\begin{minipage}[c]{5.3cm}\raggedright\Large $(0\,;\,8)$\end{minipage}\end{tcolorbox}\hfill\begin{tcolorbox}[enhanced,nobeforeafter,width=7.1cm,height=2.1cm,valign=center,colback=cvec!6,colframe=cvec,arc=2.5mm,boxrule=0.9pt,left=3mm]{\Large\bfseries\color{cvec}d)}\quad\begin{minipage}[c]{5.3cm}\raggedright\Large $(-1\,;\,3)$\end{minipage}\end{tcolorbox}
\end{tcolorbox}
\newpage
\begin{tcolorbox}[enhanced,width=16cm,height=14cm,colback=white,colframe=cvec,boxrule=2pt,arc=0mm,
left=6mm,right=6mm,top=0mm,bottom=4mm,sharp corners,
overlay={\fill[cvec] (frame.north west) rectangle ([yshift=-1.25cm]frame.north east);
\node[anchor=west,text=white,font=\sffamily\bfseries\LARGE] at ([xshift=5mm,yshift=-0.62cm]frame.north west) {Question 58 D};
\node[anchor=east,text=white,font=\sffamily\itshape\large] at ([xshift=-5mm,yshift=-0.62cm]frame.north east) {Vecteurs et repérage};}]
\vspace*{1.75cm}
\begin{minipage}[t][6.9cm][c]{\linewidth}\Large\centering
\begin{minipage}{0.95\linewidth}\raggedright\linespread{1.25}\selectfont On donne $\vec u\begin{pmatrix}5\\0\end{pmatrix}$ et $\vec v\begin{pmatrix}-1\\-3\end{pmatrix}$. Quelles sont les coordonnées du vecteur $\vec u+2\vec v$~?\end{minipage}
\end{minipage}

\vfill
\begin{tcolorbox}[enhanced,nobeforeafter,width=7.1cm,height=2.1cm,valign=center,colback=cvec!6,colframe=cvec,arc=2.5mm,boxrule=0.9pt,left=3mm]{\Large\bfseries\color{cvec}a)}\quad\begin{minipage}[c]{5.3cm}\raggedright\Large $(4\,;\,-3)$\end{minipage}\end{tcolorbox}\hfill\begin{tcolorbox}[enhanced,nobeforeafter,width=7.1cm,height=2.1cm,valign=center,colback=cvec!6,colframe=cvec,arc=2.5mm,boxrule=0.9pt,left=3mm]{\Large\bfseries\color{cvec}b)}\quad\begin{minipage}[c]{5.3cm}\raggedright\Large $(3\,;\,-6)$\end{minipage}\end{tcolorbox}\par\vspace{2.5mm}
\begin{tcolorbox}[enhanced,nobeforeafter,width=7.1cm,height=2.1cm,valign=center,colback=cvec!6,colframe=cvec,arc=2.5mm,boxrule=0.9pt,left=3mm]{\Large\bfseries\color{cvec}c)}\quad\begin{minipage}[c]{5.3cm}\raggedright\Large $(7\,;\,6)$\end{minipage}\end{tcolorbox}\hfill\begin{tcolorbox}[enhanced,nobeforeafter,width=7.1cm,height=2.1cm,valign=center,colback=cvec!6,colframe=cvec,arc=2.5mm,boxrule=0.9pt,left=3mm]{\Large\bfseries\color{cvec}d)}\quad\begin{minipage}[c]{5.3cm}\raggedright\Large $(-6\,;\,3)$\end{minipage}\end{tcolorbox}
\end{tcolorbox}
\newpage
\begin{tcolorbox}[enhanced,width=16cm,height=14cm,colback=white,colframe=cvec,boxrule=2pt,arc=0mm,
left=6mm,right=6mm,top=0mm,bottom=4mm,sharp corners,
overlay={\fill[cvec] (frame.north west) rectangle ([yshift=-1.25cm]frame.north east);
\node[anchor=west,text=white,font=\sffamily\bfseries\LARGE] at ([xshift=5mm,yshift=-0.62cm]frame.north west) {Question 59 D};
\node[anchor=east,text=white,font=\sffamily\itshape\large] at ([xshift=-5mm,yshift=-0.62cm]frame.north east) {Vecteurs et repérage};}]
\vspace*{1.75cm}
\begin{minipage}[t][6.9cm][c]{\linewidth}\LARGE\centering
\begin{minipage}{0.95\linewidth}\raggedright\linespread{1.25}\selectfont Que vaut le déterminant $\det(\vec u,\vec v)$ pour $\vec u\begin{pmatrix}-2\\-4\end{pmatrix}$ et $\vec v\begin{pmatrix}1\\2\end{pmatrix}$~?\end{minipage}
\end{minipage}

\vfill
\begin{tcolorbox}[enhanced,nobeforeafter,width=7.1cm,height=2.1cm,valign=center,colback=cvec!6,colframe=cvec,arc=2.5mm,boxrule=0.9pt,left=3mm]{\Large\bfseries\color{cvec}a)}\quad\begin{minipage}[c]{5.3cm}\raggedright\Large $-8$\end{minipage}\end{tcolorbox}\hfill\begin{tcolorbox}[enhanced,nobeforeafter,width=7.1cm,height=2.1cm,valign=center,colback=cvec!6,colframe=cvec,arc=2.5mm,boxrule=0.9pt,left=3mm]{\Large\bfseries\color{cvec}b)}\quad\begin{minipage}[c]{5.3cm}\raggedright\Large $-10$\end{minipage}\end{tcolorbox}\par\vspace{2.5mm}
\begin{tcolorbox}[enhanced,nobeforeafter,width=7.1cm,height=2.1cm,valign=center,colback=cvec!6,colframe=cvec,arc=2.5mm,boxrule=0.9pt,left=3mm]{\Large\bfseries\color{cvec}c)}\quad\begin{minipage}[c]{5.3cm}\raggedright\Large $1$\end{minipage}\end{tcolorbox}\hfill\begin{tcolorbox}[enhanced,nobeforeafter,width=7.1cm,height=2.1cm,valign=center,colback=cvec!6,colframe=cvec,arc=2.5mm,boxrule=0.9pt,left=3mm]{\Large\bfseries\color{cvec}d)}\quad\begin{minipage}[c]{5.3cm}\raggedright\Large $0$\end{minipage}\end{tcolorbox}
\end{tcolorbox}
\newpage
\begin{tcolorbox}[enhanced,width=16cm,height=14cm,colback=white,colframe=cvec,boxrule=2pt,arc=0mm,
left=6mm,right=6mm,top=0mm,bottom=4mm,sharp corners,
overlay={\fill[cvec] (frame.north west) rectangle ([yshift=-1.25cm]frame.north east);
\node[anchor=west,text=white,font=\sffamily\bfseries\LARGE] at ([xshift=5mm,yshift=-0.62cm]frame.north west) {Question 60 D};
\node[anchor=east,text=white,font=\sffamily\itshape\large] at ([xshift=-5mm,yshift=-0.62cm]frame.north east) {Vecteurs et repérage};}]
\vspace*{1.75cm}
\begin{minipage}[t][6.9cm][c]{\linewidth}\LARGE\centering
\begin{minipage}{0.95\linewidth}\raggedright\linespread{1.25}\selectfont Que vaut le déterminant $\det(\vec u,\vec v)$ pour $\vec u\begin{pmatrix}1\\6\end{pmatrix}$ et $\vec v\begin{pmatrix}3\\-2\end{pmatrix}$~?\end{minipage}
\end{minipage}

\vfill
\begin{tcolorbox}[enhanced,nobeforeafter,width=7.1cm,height=2.1cm,valign=center,colback=cvec!6,colframe=cvec,arc=2.5mm,boxrule=0.9pt,left=3mm]{\Large\bfseries\color{cvec}a)}\quad\begin{minipage}[c]{5.3cm}\raggedright\Large $-20$\end{minipage}\end{tcolorbox}\hfill\begin{tcolorbox}[enhanced,nobeforeafter,width=7.1cm,height=2.1cm,valign=center,colback=cvec!6,colframe=cvec,arc=2.5mm,boxrule=0.9pt,left=3mm]{\Large\bfseries\color{cvec}b)}\quad\begin{minipage}[c]{5.3cm}\raggedright\Large $20$\end{minipage}\end{tcolorbox}\par\vspace{2.5mm}
\begin{tcolorbox}[enhanced,nobeforeafter,width=7.1cm,height=2.1cm,valign=center,colback=cvec!6,colframe=cvec,arc=2.5mm,boxrule=0.9pt,left=3mm]{\Large\bfseries\color{cvec}c)}\quad\begin{minipage}[c]{5.3cm}\raggedright\Large $-9$\end{minipage}\end{tcolorbox}\hfill\begin{tcolorbox}[enhanced,nobeforeafter,width=7.1cm,height=2.1cm,valign=center,colback=cvec!6,colframe=cvec,arc=2.5mm,boxrule=0.9pt,left=3mm]{\Large\bfseries\color{cvec}d)}\quad\begin{minipage}[c]{5.3cm}\raggedright\Large $16$\end{minipage}\end{tcolorbox}
\end{tcolorbox}
\newpage
\begin{tcolorbox}[enhanced,width=16cm,height=14cm,colback=white,colframe=cvec,boxrule=2pt,arc=0mm,
left=6mm,right=6mm,top=0mm,bottom=4mm,sharp corners,
overlay={\fill[cvec] (frame.north west) rectangle ([yshift=-1.25cm]frame.north east);
\node[anchor=west,text=white,font=\sffamily\bfseries\LARGE] at ([xshift=5mm,yshift=-0.62cm]frame.north west) {Question 61 D};
\node[anchor=east,text=white,font=\sffamily\itshape\large] at ([xshift=-5mm,yshift=-0.62cm]frame.north east) {Vecteurs et repérage};}]
\vspace*{1.75cm}
\begin{minipage}[t][6.9cm][c]{\linewidth}\LARGE\centering
\begin{minipage}{0.95\linewidth}\raggedright\linespread{1.25}\selectfont Dans un repère, $A(-1\,;\,2)$, $B(3\,;\,5)$ et $C(4\,;\,1)$. Quelles sont les coordonnées du point $D$ tel que $ABCD$ soit un parallélogramme~?\end{minipage}
\end{minipage}

\vfill
\begin{tcolorbox}[enhanced,nobeforeafter,width=7.1cm,height=2.1cm,valign=center,colback=cvec!6,colframe=cvec,arc=2.5mm,boxrule=0.9pt,left=3mm]{\Large\bfseries\color{cvec}a)}\quad\begin{minipage}[c]{5.3cm}\raggedright\Large $(1{,}5\,;\,1{,}5)$\end{minipage}\end{tcolorbox}\hfill\begin{tcolorbox}[enhanced,nobeforeafter,width=7.1cm,height=2.1cm,valign=center,colback=cvec!6,colframe=cvec,arc=2.5mm,boxrule=0.9pt,left=3mm]{\Large\bfseries\color{cvec}b)}\quad\begin{minipage}[c]{5.3cm}\raggedright\Large $(-2\,;\,6)$\end{minipage}\end{tcolorbox}\par\vspace{2.5mm}
\begin{tcolorbox}[enhanced,nobeforeafter,width=7.1cm,height=2.1cm,valign=center,colback=cvec!6,colframe=cvec,arc=2.5mm,boxrule=0.9pt,left=3mm]{\Large\bfseries\color{cvec}c)}\quad\begin{minipage}[c]{5.3cm}\raggedright\Large $(8\,;\,4)$\end{minipage}\end{tcolorbox}\hfill\begin{tcolorbox}[enhanced,nobeforeafter,width=7.1cm,height=2.1cm,valign=center,colback=cvec!6,colframe=cvec,arc=2.5mm,boxrule=0.9pt,left=3mm]{\Large\bfseries\color{cvec}d)}\quad\begin{minipage}[c]{5.3cm}\raggedright\Large $(0\,;\,-2)$\end{minipage}\end{tcolorbox}
\end{tcolorbox}
\newpage
\begin{tcolorbox}[enhanced,width=16cm,height=14cm,colback=white,colframe=cvec,boxrule=2pt,arc=0mm,
left=6mm,right=6mm,top=0mm,bottom=4mm,sharp corners,
overlay={\fill[cvec] (frame.north west) rectangle ([yshift=-1.25cm]frame.north east);
\node[anchor=west,text=white,font=\sffamily\bfseries\LARGE] at ([xshift=5mm,yshift=-0.62cm]frame.north west) {Question 62 D};
\node[anchor=east,text=white,font=\sffamily\itshape\large] at ([xshift=-5mm,yshift=-0.62cm]frame.north east) {Vecteurs et repérage};}]
\vspace*{1.75cm}
\begin{minipage}[t][6.9cm][c]{\linewidth}\LARGE\centering
\begin{minipage}{0.95\linewidth}\raggedright\linespread{1.25}\selectfont Dans un repère, $A(5\,;\,0)$, $B(2\,;\,-3)$ et $C(-1\,;\,1)$. Quelles sont les coordonnées du point $D$ tel que $ABCD$ soit un parallélogramme~?\end{minipage}
\end{minipage}

\vfill
\begin{tcolorbox}[enhanced,nobeforeafter,width=7.1cm,height=2.1cm,valign=center,colback=cvec!6,colframe=cvec,arc=2.5mm,boxrule=0.9pt,left=3mm]{\Large\bfseries\color{cvec}a)}\quad\begin{minipage}[c]{5.3cm}\raggedright\Large $(-4\,;\,-2)$\end{minipage}\end{tcolorbox}\hfill\begin{tcolorbox}[enhanced,nobeforeafter,width=7.1cm,height=2.1cm,valign=center,colback=cvec!6,colframe=cvec,arc=2.5mm,boxrule=0.9pt,left=3mm]{\Large\bfseries\color{cvec}b)}\quad\begin{minipage}[c]{5.3cm}\raggedright\Large $(2\,;\,4)$\end{minipage}\end{tcolorbox}\par\vspace{2.5mm}
\begin{tcolorbox}[enhanced,nobeforeafter,width=7.1cm,height=2.1cm,valign=center,colback=cvec!6,colframe=cvec,arc=2.5mm,boxrule=0.9pt,left=3mm]{\Large\bfseries\color{cvec}c)}\quad\begin{minipage}[c]{5.3cm}\raggedright\Large $(2\,;\,0{,}5)$\end{minipage}\end{tcolorbox}\hfill\begin{tcolorbox}[enhanced,nobeforeafter,width=7.1cm,height=2.1cm,valign=center,colback=cvec!6,colframe=cvec,arc=2.5mm,boxrule=0.9pt,left=3mm]{\Large\bfseries\color{cvec}d)}\quad\begin{minipage}[c]{5.3cm}\raggedright\Large $(8\,;\,-4)$\end{minipage}\end{tcolorbox}
\end{tcolorbox}
\newpage
\begin{tcolorbox}[enhanced,width=16cm,height=14cm,colback=white,colframe=cvec,boxrule=2pt,arc=0mm,
left=6mm,right=6mm,top=0mm,bottom=4mm,sharp corners,
overlay={\fill[cvec] (frame.north west) rectangle ([yshift=-1.25cm]frame.north east);
\node[anchor=west,text=white,font=\sffamily\bfseries\LARGE] at ([xshift=5mm,yshift=-0.62cm]frame.north west) {Question 63 D};
\node[anchor=east,text=white,font=\sffamily\itshape\large] at ([xshift=-5mm,yshift=-0.62cm]frame.north east) {Vecteurs et repérage};}]
\vspace*{1.75cm}
\begin{minipage}[t][6.9cm][c]{\linewidth}\LARGE\centering
\begin{minipage}{0.95\linewidth}\raggedright\linespread{1.25}\selectfont Quel est le couple solution du système $\left\{\begin{array}{l}x+3y=1\\2x+y=-3\end{array}\right.$~?\end{minipage}
\end{minipage}

\vfill
\begin{tcolorbox}[enhanced,nobeforeafter,width=7.1cm,height=2.1cm,valign=center,colback=cvec!6,colframe=cvec,arc=2.5mm,boxrule=0.9pt,left=3mm]{\Large\bfseries\color{cvec}a)}\quad\begin{minipage}[c]{5.3cm}\raggedright\Large $(2\,;\,1)$\end{minipage}\end{tcolorbox}\hfill\begin{tcolorbox}[enhanced,nobeforeafter,width=7.1cm,height=2.1cm,valign=center,colback=cvec!6,colframe=cvec,arc=2.5mm,boxrule=0.9pt,left=3mm]{\Large\bfseries\color{cvec}b)}\quad\begin{minipage}[c]{5.3cm}\raggedright\Large $(1\,;\,-2)$\end{minipage}\end{tcolorbox}\par\vspace{2.5mm}
\begin{tcolorbox}[enhanced,nobeforeafter,width=7.1cm,height=2.1cm,valign=center,colback=cvec!6,colframe=cvec,arc=2.5mm,boxrule=0.9pt,left=3mm]{\Large\bfseries\color{cvec}c)}\quad\begin{minipage}[c]{5.3cm}\raggedright\Large $(-2\,;\,1)$\end{minipage}\end{tcolorbox}\hfill\begin{tcolorbox}[enhanced,nobeforeafter,width=7.1cm,height=2.1cm,valign=center,colback=cvec!6,colframe=cvec,arc=2.5mm,boxrule=0.9pt,left=3mm]{\Large\bfseries\color{cvec}d)}\quad\begin{minipage}[c]{5.3cm}\raggedright\Large $(-2\,;\,-1)$\end{minipage}\end{tcolorbox}
\end{tcolorbox}
\newpage
\begin{tcolorbox}[enhanced,width=16cm,height=14cm,colback=white,colframe=cvec,boxrule=2pt,arc=0mm,
left=6mm,right=6mm,top=0mm,bottom=4mm,sharp corners,
overlay={\fill[cvec] (frame.north west) rectangle ([yshift=-1.25cm]frame.north east);
\node[anchor=west,text=white,font=\sffamily\bfseries\LARGE] at ([xshift=5mm,yshift=-0.62cm]frame.north west) {Question 64 D};
\node[anchor=east,text=white,font=\sffamily\itshape\large] at ([xshift=-5mm,yshift=-0.62cm]frame.north east) {Vecteurs et repérage};}]
\vspace*{1.75cm}
\begin{minipage}[t][6.9cm][c]{\linewidth}\LARGE\centering
\begin{minipage}{0.95\linewidth}\raggedright\linespread{1.25}\selectfont Quel est le couple solution du système $\left\{\begin{array}{l}2x-3y=3\\x+y=4\end{array}\right.$~?\end{minipage}
\end{minipage}

\vfill
\begin{tcolorbox}[enhanced,nobeforeafter,width=7.1cm,height=2.1cm,valign=center,colback=cvec!6,colframe=cvec,arc=2.5mm,boxrule=0.9pt,left=3mm]{\Large\bfseries\color{cvec}a)}\quad\begin{minipage}[c]{5.3cm}\raggedright\Large $(-3\,;\,1)$\end{minipage}\end{tcolorbox}\hfill\begin{tcolorbox}[enhanced,nobeforeafter,width=7.1cm,height=2.1cm,valign=center,colback=cvec!6,colframe=cvec,arc=2.5mm,boxrule=0.9pt,left=3mm]{\Large\bfseries\color{cvec}b)}\quad\begin{minipage}[c]{5.3cm}\raggedright\Large $(3\,;\,-1)$\end{minipage}\end{tcolorbox}\par\vspace{2.5mm}
\begin{tcolorbox}[enhanced,nobeforeafter,width=7.1cm,height=2.1cm,valign=center,colback=cvec!6,colframe=cvec,arc=2.5mm,boxrule=0.9pt,left=3mm]{\Large\bfseries\color{cvec}c)}\quad\begin{minipage}[c]{5.3cm}\raggedright\Large $(3\,;\,1)$\end{minipage}\end{tcolorbox}\hfill\begin{tcolorbox}[enhanced,nobeforeafter,width=7.1cm,height=2.1cm,valign=center,colback=cvec!6,colframe=cvec,arc=2.5mm,boxrule=0.9pt,left=3mm]{\Large\bfseries\color{cvec}d)}\quad\begin{minipage}[c]{5.3cm}\raggedright\Large $(1\,;\,3)$\end{minipage}\end{tcolorbox}
\end{tcolorbox}
\newpage
\begin{tcolorbox}[enhanced,width=16cm,height=14cm,colback=white,colframe=cari,boxrule=2pt,arc=0mm,
left=6mm,right=6mm,top=0mm,bottom=4mm,sharp corners,
overlay={\fill[cari] (frame.north west) rectangle ([yshift=-1.25cm]frame.north east);
\node[anchor=west,text=white,font=\sffamily\bfseries\LARGE] at ([xshift=5mm,yshift=-0.62cm]frame.north west) {Question 65 D};
\node[anchor=east,text=white,font=\sffamily\itshape\large] at ([xshift=-5mm,yshift=-0.62cm]frame.north east) {Arithmétique et ensembles de nombres};}]
\vspace*{1.75cm}
\begin{minipage}[t][6.9cm][c]{\linewidth}\LARGE\centering
\begin{minipage}{0.95\linewidth}\raggedright\linespread{1.25}\selectfont Lequel de ces nombres est premier~?\end{minipage}
\end{minipage}

\vfill
\begin{tcolorbox}[enhanced,nobeforeafter,width=7.1cm,height=2.1cm,valign=center,colback=cari!6,colframe=cari,arc=2.5mm,boxrule=0.9pt,left=3mm]{\Large\bfseries\color{cari}a)}\quad\begin{minipage}[c]{5.3cm}\raggedright\Large $113$\end{minipage}\end{tcolorbox}\hfill\begin{tcolorbox}[enhanced,nobeforeafter,width=7.1cm,height=2.1cm,valign=center,colback=cari!6,colframe=cari,arc=2.5mm,boxrule=0.9pt,left=3mm]{\Large\bfseries\color{cari}b)}\quad\begin{minipage}[c]{5.3cm}\raggedright\Large $119$\end{minipage}\end{tcolorbox}\par\vspace{2.5mm}
\begin{tcolorbox}[enhanced,nobeforeafter,width=7.1cm,height=2.1cm,valign=center,colback=cari!6,colframe=cari,arc=2.5mm,boxrule=0.9pt,left=3mm]{\Large\bfseries\color{cari}c)}\quad\begin{minipage}[c]{5.3cm}\raggedright\Large $111$\end{minipage}\end{tcolorbox}\hfill\begin{tcolorbox}[enhanced,nobeforeafter,width=7.1cm,height=2.1cm,valign=center,colback=cari!6,colframe=cari,arc=2.5mm,boxrule=0.9pt,left=3mm]{\Large\bfseries\color{cari}d)}\quad\begin{minipage}[c]{5.3cm}\raggedright\Large $117$\end{minipage}\end{tcolorbox}
\end{tcolorbox}
\newpage
\begin{tcolorbox}[enhanced,width=16cm,height=14cm,colback=white,colframe=cari,boxrule=2pt,arc=0mm,
left=6mm,right=6mm,top=0mm,bottom=4mm,sharp corners,
overlay={\fill[cari] (frame.north west) rectangle ([yshift=-1.25cm]frame.north east);
\node[anchor=west,text=white,font=\sffamily\bfseries\LARGE] at ([xshift=5mm,yshift=-0.62cm]frame.north west) {Question 66 D};
\node[anchor=east,text=white,font=\sffamily\itshape\large] at ([xshift=-5mm,yshift=-0.62cm]frame.north east) {Arithmétique et ensembles de nombres};}]
\vspace*{1.75cm}
\begin{minipage}[t][6.9cm][c]{\linewidth}\LARGE\centering
\begin{minipage}{0.95\linewidth}\raggedright\linespread{1.25}\selectfont Lequel de ces nombres est premier~?\end{minipage}
\end{minipage}

\vfill
\begin{tcolorbox}[enhanced,nobeforeafter,width=7.1cm,height=2.1cm,valign=center,colback=cari!6,colframe=cari,arc=2.5mm,boxrule=0.9pt,left=3mm]{\Large\bfseries\color{cari}a)}\quad\begin{minipage}[c]{5.3cm}\raggedright\Large $77$\end{minipage}\end{tcolorbox}\hfill\begin{tcolorbox}[enhanced,nobeforeafter,width=7.1cm,height=2.1cm,valign=center,colback=cari!6,colframe=cari,arc=2.5mm,boxrule=0.9pt,left=3mm]{\Large\bfseries\color{cari}b)}\quad\begin{minipage}[c]{5.3cm}\raggedright\Large $67$\end{minipage}\end{tcolorbox}\par\vspace{2.5mm}
\begin{tcolorbox}[enhanced,nobeforeafter,width=7.1cm,height=2.1cm,valign=center,colback=cari!6,colframe=cari,arc=2.5mm,boxrule=0.9pt,left=3mm]{\Large\bfseries\color{cari}c)}\quad\begin{minipage}[c]{5.3cm}\raggedright\Large $69$\end{minipage}\end{tcolorbox}\hfill\begin{tcolorbox}[enhanced,nobeforeafter,width=7.1cm,height=2.1cm,valign=center,colback=cari!6,colframe=cari,arc=2.5mm,boxrule=0.9pt,left=3mm]{\Large\bfseries\color{cari}d)}\quad\begin{minipage}[c]{5.3cm}\raggedright\Large $87$\end{minipage}\end{tcolorbox}
\end{tcolorbox}
\newpage
\begin{tcolorbox}[enhanced,width=16cm,height=14cm,colback=white,colframe=cari,boxrule=2pt,arc=0mm,
left=6mm,right=6mm,top=0mm,bottom=4mm,sharp corners,
overlay={\fill[cari] (frame.north west) rectangle ([yshift=-1.25cm]frame.north east);
\node[anchor=west,text=white,font=\sffamily\bfseries\LARGE] at ([xshift=5mm,yshift=-0.62cm]frame.north west) {Question 67 D};
\node[anchor=east,text=white,font=\sffamily\itshape\large] at ([xshift=-5mm,yshift=-0.62cm]frame.north east) {Arithmétique et ensembles de nombres};}]
\vspace*{1.75cm}
\begin{minipage}[t][6.9cm][c]{\linewidth}\LARGE\centering
\begin{minipage}{0.95\linewidth}\raggedright\linespread{1.25}\selectfont Quelle est la décomposition de $196$ en produit de facteurs premiers~?\end{minipage}
\end{minipage}

\vfill
\begin{tcolorbox}[enhanced,nobeforeafter,width=7.1cm,height=2.1cm,valign=center,colback=cari!6,colframe=cari,arc=2.5mm,boxrule=0.9pt,left=3mm]{\Large\bfseries\color{cari}a)}\quad\begin{minipage}[c]{5.3cm}\raggedright\Large $28\times7$\end{minipage}\end{tcolorbox}\hfill\begin{tcolorbox}[enhanced,nobeforeafter,width=7.1cm,height=2.1cm,valign=center,colback=cari!6,colframe=cari,arc=2.5mm,boxrule=0.9pt,left=3mm]{\Large\bfseries\color{cari}b)}\quad\begin{minipage}[c]{5.3cm}\raggedright\Large $2^{3}\times7^{2}$\end{minipage}\end{tcolorbox}\par\vspace{2.5mm}
\begin{tcolorbox}[enhanced,nobeforeafter,width=7.1cm,height=2.1cm,valign=center,colback=cari!6,colframe=cari,arc=2.5mm,boxrule=0.9pt,left=3mm]{\Large\bfseries\color{cari}c)}\quad\begin{minipage}[c]{5.3cm}\raggedright\Large $2^{2}\times7^{2}$\end{minipage}\end{tcolorbox}\hfill\begin{tcolorbox}[enhanced,nobeforeafter,width=7.1cm,height=2.1cm,valign=center,colback=cari!6,colframe=cari,arc=2.5mm,boxrule=0.9pt,left=3mm]{\Large\bfseries\color{cari}d)}\quad\begin{minipage}[c]{5.3cm}\raggedright\Large $2^{2}\times7^{3}$\end{minipage}\end{tcolorbox}
\end{tcolorbox}
\newpage
\begin{tcolorbox}[enhanced,width=16cm,height=14cm,colback=white,colframe=cari,boxrule=2pt,arc=0mm,
left=6mm,right=6mm,top=0mm,bottom=4mm,sharp corners,
overlay={\fill[cari] (frame.north west) rectangle ([yshift=-1.25cm]frame.north east);
\node[anchor=west,text=white,font=\sffamily\bfseries\LARGE] at ([xshift=5mm,yshift=-0.62cm]frame.north west) {Question 68 D};
\node[anchor=east,text=white,font=\sffamily\itshape\large] at ([xshift=-5mm,yshift=-0.62cm]frame.north east) {Arithmétique et ensembles de nombres};}]
\vspace*{1.75cm}
\begin{minipage}[t][6.9cm][c]{\linewidth}\LARGE\centering
\begin{minipage}{0.95\linewidth}\raggedright\linespread{1.25}\selectfont Quelle est la décomposition de $330$ en produit de facteurs premiers~?\end{minipage}
\end{minipage}

\vfill
\begin{tcolorbox}[enhanced,nobeforeafter,width=7.1cm,height=2.1cm,valign=center,colback=cari!6,colframe=cari,arc=2.5mm,boxrule=0.9pt,left=3mm]{\Large\bfseries\color{cari}a)}\quad\begin{minipage}[c]{5.3cm}\raggedright\large $2\times3\times5\times11^{2}$\end{minipage}\end{tcolorbox}\hfill\begin{tcolorbox}[enhanced,nobeforeafter,width=7.1cm,height=2.1cm,valign=center,colback=cari!6,colframe=cari,arc=2.5mm,boxrule=0.9pt,left=3mm]{\Large\bfseries\color{cari}b)}\quad\begin{minipage}[c]{5.3cm}\raggedright\Large $2\times3\times5\times11$\end{minipage}\end{tcolorbox}\par\vspace{2.5mm}
\begin{tcolorbox}[enhanced,nobeforeafter,width=7.1cm,height=2.1cm,valign=center,colback=cari!6,colframe=cari,arc=2.5mm,boxrule=0.9pt,left=3mm]{\Large\bfseries\color{cari}c)}\quad\begin{minipage}[c]{5.3cm}\raggedright\large $2^{2}\times3\times5\times11$\end{minipage}\end{tcolorbox}\hfill\begin{tcolorbox}[enhanced,nobeforeafter,width=7.1cm,height=2.1cm,valign=center,colback=cari!6,colframe=cari,arc=2.5mm,boxrule=0.9pt,left=3mm]{\Large\bfseries\color{cari}d)}\quad\begin{minipage}[c]{5.3cm}\raggedright\Large $6\times55$\end{minipage}\end{tcolorbox}
\end{tcolorbox}
\newpage
\begin{tcolorbox}[enhanced,width=16cm,height=14cm,colback=white,colframe=cari,boxrule=2pt,arc=0mm,
left=6mm,right=6mm,top=0mm,bottom=4mm,sharp corners,
overlay={\fill[cari] (frame.north west) rectangle ([yshift=-1.25cm]frame.north east);
\node[anchor=west,text=white,font=\sffamily\bfseries\LARGE] at ([xshift=5mm,yshift=-0.62cm]frame.north west) {Question 69 D};
\node[anchor=east,text=white,font=\sffamily\itshape\large] at ([xshift=-5mm,yshift=-0.62cm]frame.north east) {Arithmétique et ensembles de nombres};}]
\vspace*{1.75cm}
\begin{minipage}[t][6.9cm][c]{\linewidth}\LARGE\centering
\begin{minipage}{0.95\linewidth}\raggedright\linespread{1.25}\selectfont Quel est le plus petit ensemble de nombres auquel appartient $\dfrac{22}{7}$~?\end{minipage}
\end{minipage}

\vfill
\begin{tcolorbox}[enhanced,nobeforeafter,width=7.1cm,height=2.1cm,valign=center,colback=cari!6,colframe=cari,arc=2.5mm,boxrule=0.9pt,left=3mm]{\Large\bfseries\color{cari}a)}\quad\begin{minipage}[c]{5.3cm}\raggedright\Large $\mathbb{D}$\end{minipage}\end{tcolorbox}\hfill\begin{tcolorbox}[enhanced,nobeforeafter,width=7.1cm,height=2.1cm,valign=center,colback=cari!6,colframe=cari,arc=2.5mm,boxrule=0.9pt,left=3mm]{\Large\bfseries\color{cari}b)}\quad\begin{minipage}[c]{5.3cm}\raggedright\Large $\mathbb{Q}$\end{minipage}\end{tcolorbox}\par\vspace{2.5mm}
\begin{tcolorbox}[enhanced,nobeforeafter,width=7.1cm,height=2.1cm,valign=center,colback=cari!6,colframe=cari,arc=2.5mm,boxrule=0.9pt,left=3mm]{\Large\bfseries\color{cari}c)}\quad\begin{minipage}[c]{5.3cm}\raggedright\Large $\mathbb{Z}$\end{minipage}\end{tcolorbox}\hfill\begin{tcolorbox}[enhanced,nobeforeafter,width=7.1cm,height=2.1cm,valign=center,colback=cari!6,colframe=cari,arc=2.5mm,boxrule=0.9pt,left=3mm]{\Large\bfseries\color{cari}d)}\quad\begin{minipage}[c]{5.3cm}\raggedright\Large $\mathbb{N}$\end{minipage}\end{tcolorbox}
\end{tcolorbox}
\newpage
\begin{tcolorbox}[enhanced,width=16cm,height=14cm,colback=white,colframe=cari,boxrule=2pt,arc=0mm,
left=6mm,right=6mm,top=0mm,bottom=4mm,sharp corners,
overlay={\fill[cari] (frame.north west) rectangle ([yshift=-1.25cm]frame.north east);
\node[anchor=west,text=white,font=\sffamily\bfseries\LARGE] at ([xshift=5mm,yshift=-0.62cm]frame.north west) {Question 70 D};
\node[anchor=east,text=white,font=\sffamily\itshape\large] at ([xshift=-5mm,yshift=-0.62cm]frame.north east) {Arithmétique et ensembles de nombres};}]
\vspace*{1.75cm}
\begin{minipage}[t][6.9cm][c]{\linewidth}\LARGE\centering
\begin{minipage}{0.95\linewidth}\raggedright\linespread{1.25}\selectfont Quel est le plus petit ensemble de nombres auquel appartient $\pi$~?\end{minipage}
\end{minipage}

\vfill
\begin{tcolorbox}[enhanced,nobeforeafter,width=7.1cm,height=2.1cm,valign=center,colback=cari!6,colframe=cari,arc=2.5mm,boxrule=0.9pt,left=3mm]{\Large\bfseries\color{cari}a)}\quad\begin{minipage}[c]{5.3cm}\raggedright\Large $\mathbb{D}$\end{minipage}\end{tcolorbox}\hfill\begin{tcolorbox}[enhanced,nobeforeafter,width=7.1cm,height=2.1cm,valign=center,colback=cari!6,colframe=cari,arc=2.5mm,boxrule=0.9pt,left=3mm]{\Large\bfseries\color{cari}b)}\quad\begin{minipage}[c]{5.3cm}\raggedright\Large $\mathbb{R}$\end{minipage}\end{tcolorbox}\par\vspace{2.5mm}
\begin{tcolorbox}[enhanced,nobeforeafter,width=7.1cm,height=2.1cm,valign=center,colback=cari!6,colframe=cari,arc=2.5mm,boxrule=0.9pt,left=3mm]{\Large\bfseries\color{cari}c)}\quad\begin{minipage}[c]{5.3cm}\raggedright\Large $\mathbb{N}$\end{minipage}\end{tcolorbox}\hfill\begin{tcolorbox}[enhanced,nobeforeafter,width=7.1cm,height=2.1cm,valign=center,colback=cari!6,colframe=cari,arc=2.5mm,boxrule=0.9pt,left=3mm]{\Large\bfseries\color{cari}d)}\quad\begin{minipage}[c]{5.3cm}\raggedright\Large $\mathbb{Z}$\end{minipage}\end{tcolorbox}
\end{tcolorbox}
\newpage
\begin{tcolorbox}[enhanced,width=16cm,height=14cm,colback=white,colframe=cari,boxrule=2pt,arc=0mm,
left=6mm,right=6mm,top=0mm,bottom=4mm,sharp corners,
overlay={\fill[cari] (frame.north west) rectangle ([yshift=-1.25cm]frame.north east);
\node[anchor=west,text=white,font=\sffamily\bfseries\LARGE] at ([xshift=5mm,yshift=-0.62cm]frame.north west) {Question 71 D};
\node[anchor=east,text=white,font=\sffamily\itshape\large] at ([xshift=-5mm,yshift=-0.62cm]frame.north east) {Arithmétique et ensembles de nombres};}]
\vspace*{1.75cm}
\begin{minipage}[t][6.9cm][c]{\linewidth}\LARGE\centering
\begin{minipage}{0.95\linewidth}\raggedright\linespread{1.25}\selectfont Quelle est la forme irréductible de $\dfrac{24}{108}$~?\end{minipage}
\end{minipage}

\vfill
\begin{tcolorbox}[enhanced,nobeforeafter,width=7.1cm,height=2.1cm,valign=center,colback=cari!6,colframe=cari,arc=2.5mm,boxrule=0.9pt,left=3mm]{\Large\bfseries\color{cari}a)}\quad\begin{minipage}[c]{5.3cm}\raggedright\Large $\dfrac{12}{54}$\end{minipage}\end{tcolorbox}\hfill\begin{tcolorbox}[enhanced,nobeforeafter,width=7.1cm,height=2.1cm,valign=center,colback=cari!6,colframe=cari,arc=2.5mm,boxrule=0.9pt,left=3mm]{\Large\bfseries\color{cari}b)}\quad\begin{minipage}[c]{5.3cm}\raggedright\Large $\dfrac{9}{2}$\end{minipage}\end{tcolorbox}\par\vspace{2.5mm}
\begin{tcolorbox}[enhanced,nobeforeafter,width=7.1cm,height=2.1cm,valign=center,colback=cari!6,colframe=cari,arc=2.5mm,boxrule=0.9pt,left=3mm]{\Large\bfseries\color{cari}c)}\quad\begin{minipage}[c]{5.3cm}\raggedright\Large $\dfrac{2}{9}$\end{minipage}\end{tcolorbox}\hfill\begin{tcolorbox}[enhanced,nobeforeafter,width=7.1cm,height=2.1cm,valign=center,colback=cari!6,colframe=cari,arc=2.5mm,boxrule=0.9pt,left=3mm]{\Large\bfseries\color{cari}d)}\quad\begin{minipage}[c]{5.3cm}\raggedright\Large $\dfrac{8}{36}$\end{minipage}\end{tcolorbox}
\end{tcolorbox}
\newpage
\begin{tcolorbox}[enhanced,width=16cm,height=14cm,colback=white,colframe=cari,boxrule=2pt,arc=0mm,
left=6mm,right=6mm,top=0mm,bottom=4mm,sharp corners,
overlay={\fill[cari] (frame.north west) rectangle ([yshift=-1.25cm]frame.north east);
\node[anchor=west,text=white,font=\sffamily\bfseries\LARGE] at ([xshift=5mm,yshift=-0.62cm]frame.north west) {Question 72 D};
\node[anchor=east,text=white,font=\sffamily\itshape\large] at ([xshift=-5mm,yshift=-0.62cm]frame.north east) {Arithmétique et ensembles de nombres};}]
\vspace*{1.75cm}
\begin{minipage}[t][6.9cm][c]{\linewidth}\LARGE\centering
\begin{minipage}{0.95\linewidth}\raggedright\linespread{1.25}\selectfont Quelle est la forme irréductible de $\dfrac{132}{99}$~?\end{minipage}
\end{minipage}

\vfill
\begin{tcolorbox}[enhanced,nobeforeafter,width=7.1cm,height=2.1cm,valign=center,colback=cari!6,colframe=cari,arc=2.5mm,boxrule=0.9pt,left=3mm]{\Large\bfseries\color{cari}a)}\quad\begin{minipage}[c]{5.3cm}\raggedright\Large $\dfrac{44}{33}$\end{minipage}\end{tcolorbox}\hfill\begin{tcolorbox}[enhanced,nobeforeafter,width=7.1cm,height=2.1cm,valign=center,colback=cari!6,colframe=cari,arc=2.5mm,boxrule=0.9pt,left=3mm]{\Large\bfseries\color{cari}b)}\quad\begin{minipage}[c]{5.3cm}\raggedright\Large $\dfrac{4}{3}$\end{minipage}\end{tcolorbox}\par\vspace{2.5mm}
\begin{tcolorbox}[enhanced,nobeforeafter,width=7.1cm,height=2.1cm,valign=center,colback=cari!6,colframe=cari,arc=2.5mm,boxrule=0.9pt,left=3mm]{\Large\bfseries\color{cari}c)}\quad\begin{minipage}[c]{5.3cm}\raggedright\Large $\dfrac{3}{4}$\end{minipage}\end{tcolorbox}\hfill\begin{tcolorbox}[enhanced,nobeforeafter,width=7.1cm,height=2.1cm,valign=center,colback=cari!6,colframe=cari,arc=2.5mm,boxrule=0.9pt,left=3mm]{\Large\bfseries\color{cari}d)}\quad\begin{minipage}[c]{5.3cm}\raggedright\Large $\dfrac{5}{4}$\end{minipage}\end{tcolorbox}
\end{tcolorbox}
\newpage
\begin{tcolorbox}[enhanced,width=16cm,height=14cm,colback=white,colframe=cari,boxrule=2pt,arc=0mm,
left=6mm,right=6mm,top=0mm,bottom=4mm,sharp corners,
overlay={\fill[cari] (frame.north west) rectangle ([yshift=-1.25cm]frame.north east);
\node[anchor=west,text=white,font=\sffamily\bfseries\LARGE] at ([xshift=5mm,yshift=-0.62cm]frame.north west) {Question 73 D};
\node[anchor=east,text=white,font=\sffamily\itshape\large] at ([xshift=-5mm,yshift=-0.62cm]frame.north east) {Arithmétique et ensembles de nombres};}]
\vspace*{1.75cm}
\begin{minipage}[t][6.9cm][c]{\linewidth}\LARGE\centering
\begin{minipage}{0.95\linewidth}\raggedright\linespread{1.25}\selectfont Pour tout entier naturel $n$, que peut-on dire de l'entier $6n$~?\end{minipage}
\end{minipage}

\vfill
\begin{tcolorbox}[enhanced,nobeforeafter,width=7.1cm,height=2.1cm,valign=center,colback=cari!6,colframe=cari,arc=2.5mm,boxrule=0.9pt,left=3mm]{\Large\bfseries\color{cari}a)}\quad\begin{minipage}[c]{5.3cm}\raggedright\Large il est toujours premier\end{minipage}\end{tcolorbox}\hfill\begin{tcolorbox}[enhanced,nobeforeafter,width=7.1cm,height=2.1cm,valign=center,colback=cari!6,colframe=cari,arc=2.5mm,boxrule=0.9pt,left=3mm]{\Large\bfseries\color{cari}b)}\quad\begin{minipage}[c]{5.3cm}\raggedright\large il est pair ou impair selon $n$\end{minipage}\end{tcolorbox}\par\vspace{2.5mm}
\begin{tcolorbox}[enhanced,nobeforeafter,width=7.1cm,height=2.1cm,valign=center,colback=cari!6,colframe=cari,arc=2.5mm,boxrule=0.9pt,left=3mm]{\Large\bfseries\color{cari}c)}\quad\begin{minipage}[c]{5.3cm}\raggedright\Large il est toujours pair\end{minipage}\end{tcolorbox}\hfill\begin{tcolorbox}[enhanced,nobeforeafter,width=7.1cm,height=2.1cm,valign=center,colback=cari!6,colframe=cari,arc=2.5mm,boxrule=0.9pt,left=3mm]{\Large\bfseries\color{cari}d)}\quad\begin{minipage}[c]{5.3cm}\raggedright\Large il est toujours impair\end{minipage}\end{tcolorbox}
\end{tcolorbox}
\newpage
\begin{tcolorbox}[enhanced,width=16cm,height=14cm,colback=white,colframe=cari,boxrule=2pt,arc=0mm,
left=6mm,right=6mm,top=0mm,bottom=4mm,sharp corners,
overlay={\fill[cari] (frame.north west) rectangle ([yshift=-1.25cm]frame.north east);
\node[anchor=west,text=white,font=\sffamily\bfseries\LARGE] at ([xshift=5mm,yshift=-0.62cm]frame.north west) {Question 74 D};
\node[anchor=east,text=white,font=\sffamily\itshape\large] at ([xshift=-5mm,yshift=-0.62cm]frame.north east) {Arithmétique et ensembles de nombres};}]
\vspace*{1.75cm}
\begin{minipage}[t][6.9cm][c]{\linewidth}\LARGE\centering
\begin{minipage}{0.95\linewidth}\raggedright\linespread{1.25}\selectfont Pour tout entier naturel $n$, que peut-on dire de l'entier $2n-1$~?\end{minipage}
\end{minipage}

\vfill
\begin{tcolorbox}[enhanced,nobeforeafter,width=7.1cm,height=2.1cm,valign=center,colback=cari!6,colframe=cari,arc=2.5mm,boxrule=0.9pt,left=3mm]{\Large\bfseries\color{cari}a)}\quad\begin{minipage}[c]{5.3cm}\raggedright\Large il est toujours impair\end{minipage}\end{tcolorbox}\hfill\begin{tcolorbox}[enhanced,nobeforeafter,width=7.1cm,height=2.1cm,valign=center,colback=cari!6,colframe=cari,arc=2.5mm,boxrule=0.9pt,left=3mm]{\Large\bfseries\color{cari}b)}\quad\begin{minipage}[c]{5.3cm}\raggedright\large il est pair ou impair selon $n$\end{minipage}\end{tcolorbox}\par\vspace{2.5mm}
\begin{tcolorbox}[enhanced,nobeforeafter,width=7.1cm,height=2.1cm,valign=center,colback=cari!6,colframe=cari,arc=2.5mm,boxrule=0.9pt,left=3mm]{\Large\bfseries\color{cari}c)}\quad\begin{minipage}[c]{5.3cm}\raggedright\Large il est toujours premier\end{minipage}\end{tcolorbox}\hfill\begin{tcolorbox}[enhanced,nobeforeafter,width=7.1cm,height=2.1cm,valign=center,colback=cari!6,colframe=cari,arc=2.5mm,boxrule=0.9pt,left=3mm]{\Large\bfseries\color{cari}d)}\quad\begin{minipage}[c]{5.3cm}\raggedright\Large il est toujours pair\end{minipage}\end{tcolorbox}
\end{tcolorbox}
\newpage
\begin{tcolorbox}[enhanced,width=16cm,height=14cm,colback=white,colframe=cari,boxrule=2pt,arc=0mm,
left=6mm,right=6mm,top=0mm,bottom=4mm,sharp corners,
overlay={\fill[cari] (frame.north west) rectangle ([yshift=-1.25cm]frame.north east);
\node[anchor=west,text=white,font=\sffamily\bfseries\LARGE] at ([xshift=5mm,yshift=-0.62cm]frame.north west) {Question 75 D};
\node[anchor=east,text=white,font=\sffamily\itshape\large] at ([xshift=-5mm,yshift=-0.62cm]frame.north east) {Arithmétique et ensembles de nombres};}]
\vspace*{1.75cm}
\begin{minipage}[t][6.9cm][c]{\linewidth}\LARGE\centering
\begin{minipage}{0.95\linewidth}\raggedright\linespread{1.25}\selectfont Combien y a-t-il de multiples de $13$ compris entre $1$ et $300$~?\end{minipage}
\end{minipage}

\vfill
\begin{tcolorbox}[enhanced,nobeforeafter,width=7.1cm,height=2.1cm,valign=center,colback=cari!6,colframe=cari,arc=2.5mm,boxrule=0.9pt,left=3mm]{\Large\bfseries\color{cari}a)}\quad\begin{minipage}[c]{5.3cm}\raggedright\Large $23$\end{minipage}\end{tcolorbox}\hfill\begin{tcolorbox}[enhanced,nobeforeafter,width=7.1cm,height=2.1cm,valign=center,colback=cari!6,colframe=cari,arc=2.5mm,boxrule=0.9pt,left=3mm]{\Large\bfseries\color{cari}b)}\quad\begin{minipage}[c]{5.3cm}\raggedright\Large $22$\end{minipage}\end{tcolorbox}\par\vspace{2.5mm}
\begin{tcolorbox}[enhanced,nobeforeafter,width=7.1cm,height=2.1cm,valign=center,colback=cari!6,colframe=cari,arc=2.5mm,boxrule=0.9pt,left=3mm]{\Large\bfseries\color{cari}c)}\quad\begin{minipage}[c]{5.3cm}\raggedright\Large $277$\end{minipage}\end{tcolorbox}\hfill\begin{tcolorbox}[enhanced,nobeforeafter,width=7.1cm,height=2.1cm,valign=center,colback=cari!6,colframe=cari,arc=2.5mm,boxrule=0.9pt,left=3mm]{\Large\bfseries\color{cari}d)}\quad\begin{minipage}[c]{5.3cm}\raggedright\Large $24$\end{minipage}\end{tcolorbox}
\end{tcolorbox}
\newpage
\begin{tcolorbox}[enhanced,width=16cm,height=14cm,colback=white,colframe=cari,boxrule=2pt,arc=0mm,
left=6mm,right=6mm,top=0mm,bottom=4mm,sharp corners,
overlay={\fill[cari] (frame.north west) rectangle ([yshift=-1.25cm]frame.north east);
\node[anchor=west,text=white,font=\sffamily\bfseries\LARGE] at ([xshift=5mm,yshift=-0.62cm]frame.north west) {Question 76 D};
\node[anchor=east,text=white,font=\sffamily\itshape\large] at ([xshift=-5mm,yshift=-0.62cm]frame.north east) {Arithmétique et ensembles de nombres};}]
\vspace*{1.75cm}
\begin{minipage}[t][6.9cm][c]{\linewidth}\LARGE\centering
\begin{minipage}{0.95\linewidth}\raggedright\linespread{1.25}\selectfont Combien y a-t-il de multiples de $6$ compris entre $1$ et $1\,000$~?\end{minipage}
\end{minipage}

\vfill
\begin{tcolorbox}[enhanced,nobeforeafter,width=7.1cm,height=2.1cm,valign=center,colback=cari!6,colframe=cari,arc=2.5mm,boxrule=0.9pt,left=3mm]{\Large\bfseries\color{cari}a)}\quad\begin{minipage}[c]{5.3cm}\raggedright\Large $165$\end{minipage}\end{tcolorbox}\hfill\begin{tcolorbox}[enhanced,nobeforeafter,width=7.1cm,height=2.1cm,valign=center,colback=cari!6,colframe=cari,arc=2.5mm,boxrule=0.9pt,left=3mm]{\Large\bfseries\color{cari}b)}\quad\begin{minipage}[c]{5.3cm}\raggedright\Large $834$\end{minipage}\end{tcolorbox}\par\vspace{2.5mm}
\begin{tcolorbox}[enhanced,nobeforeafter,width=7.1cm,height=2.1cm,valign=center,colback=cari!6,colframe=cari,arc=2.5mm,boxrule=0.9pt,left=3mm]{\Large\bfseries\color{cari}c)}\quad\begin{minipage}[c]{5.3cm}\raggedright\Large $167$\end{minipage}\end{tcolorbox}\hfill\begin{tcolorbox}[enhanced,nobeforeafter,width=7.1cm,height=2.1cm,valign=center,colback=cari!6,colframe=cari,arc=2.5mm,boxrule=0.9pt,left=3mm]{\Large\bfseries\color{cari}d)}\quad\begin{minipage}[c]{5.3cm}\raggedright\Large $166$\end{minipage}\end{tcolorbox}
\end{tcolorbox}
\newpage
\begin{tcolorbox}[enhanced,width=16cm,height=14cm,colback=white,colframe=cari,boxrule=2pt,arc=0mm,
left=6mm,right=6mm,top=0mm,bottom=4mm,sharp corners,
overlay={\fill[cari] (frame.north west) rectangle ([yshift=-1.25cm]frame.north east);
\node[anchor=west,text=white,font=\sffamily\bfseries\LARGE] at ([xshift=5mm,yshift=-0.62cm]frame.north west) {Question 77 D};
\node[anchor=east,text=white,font=\sffamily\itshape\large] at ([xshift=-5mm,yshift=-0.62cm]frame.north east) {Arithmétique et ensembles de nombres};}]
\vspace*{1.75cm}
\begin{minipage}[t][6.9cm][c]{\linewidth}\LARGE\centering
\begin{minipage}{0.95\linewidth}\raggedright\linespread{1.25}\selectfont Combien l'entier $18$ a-t-il de diviseurs positifs~?\end{minipage}
\end{minipage}

\vfill
\begin{tcolorbox}[enhanced,nobeforeafter,width=7.1cm,height=2.1cm,valign=center,colback=cari!6,colframe=cari,arc=2.5mm,boxrule=0.9pt,left=3mm]{\Large\bfseries\color{cari}a)}\quad\begin{minipage}[c]{5.3cm}\raggedright\Large $5$\end{minipage}\end{tcolorbox}\hfill\begin{tcolorbox}[enhanced,nobeforeafter,width=7.1cm,height=2.1cm,valign=center,colback=cari!6,colframe=cari,arc=2.5mm,boxrule=0.9pt,left=3mm]{\Large\bfseries\color{cari}b)}\quad\begin{minipage}[c]{5.3cm}\raggedright\Large $7$\end{minipage}\end{tcolorbox}\par\vspace{2.5mm}
\begin{tcolorbox}[enhanced,nobeforeafter,width=7.1cm,height=2.1cm,valign=center,colback=cari!6,colframe=cari,arc=2.5mm,boxrule=0.9pt,left=3mm]{\Large\bfseries\color{cari}c)}\quad\begin{minipage}[c]{5.3cm}\raggedright\Large $8$\end{minipage}\end{tcolorbox}\hfill\begin{tcolorbox}[enhanced,nobeforeafter,width=7.1cm,height=2.1cm,valign=center,colback=cari!6,colframe=cari,arc=2.5mm,boxrule=0.9pt,left=3mm]{\Large\bfseries\color{cari}d)}\quad\begin{minipage}[c]{5.3cm}\raggedright\Large $6$\end{minipage}\end{tcolorbox}
\end{tcolorbox}
\newpage
\begin{tcolorbox}[enhanced,width=16cm,height=14cm,colback=white,colframe=cari,boxrule=2pt,arc=0mm,
left=6mm,right=6mm,top=0mm,bottom=4mm,sharp corners,
overlay={\fill[cari] (frame.north west) rectangle ([yshift=-1.25cm]frame.north east);
\node[anchor=west,text=white,font=\sffamily\bfseries\LARGE] at ([xshift=5mm,yshift=-0.62cm]frame.north west) {Question 78 D};
\node[anchor=east,text=white,font=\sffamily\itshape\large] at ([xshift=-5mm,yshift=-0.62cm]frame.north east) {Arithmétique et ensembles de nombres};}]
\vspace*{1.75cm}
\begin{minipage}[t][6.9cm][c]{\linewidth}\LARGE\centering
\begin{minipage}{0.95\linewidth}\raggedright\linespread{1.25}\selectfont Combien l'entier $100$ a-t-il de diviseurs positifs~?\end{minipage}
\end{minipage}

\vfill
\begin{tcolorbox}[enhanced,nobeforeafter,width=7.1cm,height=2.1cm,valign=center,colback=cari!6,colframe=cari,arc=2.5mm,boxrule=0.9pt,left=3mm]{\Large\bfseries\color{cari}a)}\quad\begin{minipage}[c]{5.3cm}\raggedright\Large $10$\end{minipage}\end{tcolorbox}\hfill\begin{tcolorbox}[enhanced,nobeforeafter,width=7.1cm,height=2.1cm,valign=center,colback=cari!6,colframe=cari,arc=2.5mm,boxrule=0.9pt,left=3mm]{\Large\bfseries\color{cari}b)}\quad\begin{minipage}[c]{5.3cm}\raggedright\Large $8$\end{minipage}\end{tcolorbox}\par\vspace{2.5mm}
\begin{tcolorbox}[enhanced,nobeforeafter,width=7.1cm,height=2.1cm,valign=center,colback=cari!6,colframe=cari,arc=2.5mm,boxrule=0.9pt,left=3mm]{\Large\bfseries\color{cari}c)}\quad\begin{minipage}[c]{5.3cm}\raggedright\Large $9$\end{minipage}\end{tcolorbox}\hfill\begin{tcolorbox}[enhanced,nobeforeafter,width=7.1cm,height=2.1cm,valign=center,colback=cari!6,colframe=cari,arc=2.5mm,boxrule=0.9pt,left=3mm]{\Large\bfseries\color{cari}d)}\quad\begin{minipage}[c]{5.3cm}\raggedright\Large $11$\end{minipage}\end{tcolorbox}
\end{tcolorbox}
\newpage
\begin{tcolorbox}[enhanced,width=16cm,height=14cm,colback=white,colframe=cari,boxrule=2pt,arc=0mm,
left=6mm,right=6mm,top=0mm,bottom=4mm,sharp corners,
overlay={\fill[cari] (frame.north west) rectangle ([yshift=-1.25cm]frame.north east);
\node[anchor=west,text=white,font=\sffamily\bfseries\LARGE] at ([xshift=5mm,yshift=-0.62cm]frame.north west) {Question 79 D};
\node[anchor=east,text=white,font=\sffamily\itshape\large] at ([xshift=-5mm,yshift=-0.62cm]frame.north east) {Arithmétique et ensembles de nombres};}]
\vspace*{1.75cm}
\begin{minipage}[t][6.9cm][c]{\linewidth}\LARGE\centering
\begin{minipage}{0.95\linewidth}\raggedright\linespread{1.25}\selectfont Que vaut l'intersection $[-1\,;\,10]\cap]5\,;\,+\infty[$~?\end{minipage}
\end{minipage}

\vfill
\begin{tcolorbox}[enhanced,nobeforeafter,width=7.1cm,height=2.1cm,valign=center,colback=cari!6,colframe=cari,arc=2.5mm,boxrule=0.9pt,left=3mm]{\Large\bfseries\color{cari}a)}\quad\begin{minipage}[c]{5.3cm}\raggedright\Large $[5\,;\,10]$\end{minipage}\end{tcolorbox}\hfill\begin{tcolorbox}[enhanced,nobeforeafter,width=7.1cm,height=2.1cm,valign=center,colback=cari!6,colframe=cari,arc=2.5mm,boxrule=0.9pt,left=3mm]{\Large\bfseries\color{cari}b)}\quad\begin{minipage}[c]{5.3cm}\raggedright\Large $[-1\,;\,+\infty[$\end{minipage}\end{tcolorbox}\par\vspace{2.5mm}
\begin{tcolorbox}[enhanced,nobeforeafter,width=7.1cm,height=2.1cm,valign=center,colback=cari!6,colframe=cari,arc=2.5mm,boxrule=0.9pt,left=3mm]{\Large\bfseries\color{cari}c)}\quad\begin{minipage}[c]{5.3cm}\raggedright\Large $]5\,;\,10[$\end{minipage}\end{tcolorbox}\hfill\begin{tcolorbox}[enhanced,nobeforeafter,width=7.1cm,height=2.1cm,valign=center,colback=cari!6,colframe=cari,arc=2.5mm,boxrule=0.9pt,left=3mm]{\Large\bfseries\color{cari}d)}\quad\begin{minipage}[c]{5.3cm}\raggedright\Large $]5\,;\,10]$\end{minipage}\end{tcolorbox}
\end{tcolorbox}
\newpage
\begin{tcolorbox}[enhanced,width=16cm,height=14cm,colback=white,colframe=cari,boxrule=2pt,arc=0mm,
left=6mm,right=6mm,top=0mm,bottom=4mm,sharp corners,
overlay={\fill[cari] (frame.north west) rectangle ([yshift=-1.25cm]frame.north east);
\node[anchor=west,text=white,font=\sffamily\bfseries\LARGE] at ([xshift=5mm,yshift=-0.62cm]frame.north west) {Question 80 D};
\node[anchor=east,text=white,font=\sffamily\itshape\large] at ([xshift=-5mm,yshift=-0.62cm]frame.north east) {Arithmétique et ensembles de nombres};}]
\vspace*{1.75cm}
\begin{minipage}[t][6.9cm][c]{\linewidth}\LARGE\centering
\begin{minipage}{0.95\linewidth}\raggedright\linespread{1.25}\selectfont Que vaut l'intersection $[2\,;\,9[\cap]-\infty\,;\,5[$~?\end{minipage}
\end{minipage}

\vfill
\begin{tcolorbox}[enhanced,nobeforeafter,width=7.1cm,height=2.1cm,valign=center,colback=cari!6,colframe=cari,arc=2.5mm,boxrule=0.9pt,left=3mm]{\Large\bfseries\color{cari}a)}\quad\begin{minipage}[c]{5.3cm}\raggedright\Large $]-\infty\,;\,9[$\end{minipage}\end{tcolorbox}\hfill\begin{tcolorbox}[enhanced,nobeforeafter,width=7.1cm,height=2.1cm,valign=center,colback=cari!6,colframe=cari,arc=2.5mm,boxrule=0.9pt,left=3mm]{\Large\bfseries\color{cari}b)}\quad\begin{minipage}[c]{5.3cm}\raggedright\Large $]2\,;\,5[$\end{minipage}\end{tcolorbox}\par\vspace{2.5mm}
\begin{tcolorbox}[enhanced,nobeforeafter,width=7.1cm,height=2.1cm,valign=center,colback=cari!6,colframe=cari,arc=2.5mm,boxrule=0.9pt,left=3mm]{\Large\bfseries\color{cari}c)}\quad\begin{minipage}[c]{5.3cm}\raggedright\Large $[2\,;\,5[$\end{minipage}\end{tcolorbox}\hfill\begin{tcolorbox}[enhanced,nobeforeafter,width=7.1cm,height=2.1cm,valign=center,colback=cari!6,colframe=cari,arc=2.5mm,boxrule=0.9pt,left=3mm]{\Large\bfseries\color{cari}d)}\quad\begin{minipage}[c]{5.3cm}\raggedright\Large $[2\,;\,5]$\end{minipage}\end{tcolorbox}
\end{tcolorbox}
\newpage
\begin{tcolorbox}[enhanced,width=16cm,height=14cm,colback=white,colframe=cpro,boxrule=2pt,arc=0mm,
left=6mm,right=6mm,top=0mm,bottom=4mm,sharp corners,
overlay={\fill[cpro] (frame.north west) rectangle ([yshift=-1.25cm]frame.north east);
\node[anchor=west,text=white,font=\sffamily\bfseries\LARGE] at ([xshift=5mm,yshift=-0.62cm]frame.north west) {Question 81 D};
\node[anchor=east,text=white,font=\sffamily\itshape\large] at ([xshift=-5mm,yshift=-0.62cm]frame.north east) {Probabilités};}]
\vspace*{1.75cm}
\begin{minipage}[t][6.9cm][c]{\linewidth}\LARGE\centering
\begin{minipage}{0.95\linewidth}\raggedright\linespread{1.25}\selectfont On a $P(A)=0{,}35$, $P(B)=0{,}25$ et $P(A\cap B)=0{,}05$. Que vaut $P(A\cup B)$~?\end{minipage}
\end{minipage}

\vfill
\begin{tcolorbox}[enhanced,nobeforeafter,width=7.1cm,height=2.1cm,valign=center,colback=cpro!6,colframe=cpro,arc=2.5mm,boxrule=0.9pt,left=3mm]{\Large\bfseries\color{cpro}a)}\quad\begin{minipage}[c]{5.3cm}\raggedright\Large $0{,}05$\end{minipage}\end{tcolorbox}\hfill\begin{tcolorbox}[enhanced,nobeforeafter,width=7.1cm,height=2.1cm,valign=center,colback=cpro!6,colframe=cpro,arc=2.5mm,boxrule=0.9pt,left=3mm]{\Large\bfseries\color{cpro}b)}\quad\begin{minipage}[c]{5.3cm}\raggedright\Large $0{,}65$\end{minipage}\end{tcolorbox}\par\vspace{2.5mm}
\begin{tcolorbox}[enhanced,nobeforeafter,width=7.1cm,height=2.1cm,valign=center,colback=cpro!6,colframe=cpro,arc=2.5mm,boxrule=0.9pt,left=3mm]{\Large\bfseries\color{cpro}c)}\quad\begin{minipage}[c]{5.3cm}\raggedright\Large $0{,}6$\end{minipage}\end{tcolorbox}\hfill\begin{tcolorbox}[enhanced,nobeforeafter,width=7.1cm,height=2.1cm,valign=center,colback=cpro!6,colframe=cpro,arc=2.5mm,boxrule=0.9pt,left=3mm]{\Large\bfseries\color{cpro}d)}\quad\begin{minipage}[c]{5.3cm}\raggedright\Large $0{,}55$\end{minipage}\end{tcolorbox}
\end{tcolorbox}
\newpage
\begin{tcolorbox}[enhanced,width=16cm,height=14cm,colback=white,colframe=cpro,boxrule=2pt,arc=0mm,
left=6mm,right=6mm,top=0mm,bottom=4mm,sharp corners,
overlay={\fill[cpro] (frame.north west) rectangle ([yshift=-1.25cm]frame.north east);
\node[anchor=west,text=white,font=\sffamily\bfseries\LARGE] at ([xshift=5mm,yshift=-0.62cm]frame.north west) {Question 82 D};
\node[anchor=east,text=white,font=\sffamily\itshape\large] at ([xshift=-5mm,yshift=-0.62cm]frame.north east) {Probabilités};}]
\vspace*{1.75cm}
\begin{minipage}[t][6.9cm][c]{\linewidth}\LARGE\centering
\begin{minipage}{0.95\linewidth}\raggedright\linespread{1.25}\selectfont On a $P(A)=0{,}5$, $P(B)=0{,}45$ et $P(A\cap B)=0{,}3$. Que vaut $P(A\cup B)$~?\end{minipage}
\end{minipage}

\vfill
\begin{tcolorbox}[enhanced,nobeforeafter,width=7.1cm,height=2.1cm,valign=center,colback=cpro!6,colframe=cpro,arc=2.5mm,boxrule=0.9pt,left=3mm]{\Large\bfseries\color{cpro}a)}\quad\begin{minipage}[c]{5.3cm}\raggedright\Large $0{,}3$\end{minipage}\end{tcolorbox}\hfill\begin{tcolorbox}[enhanced,nobeforeafter,width=7.1cm,height=2.1cm,valign=center,colback=cpro!6,colframe=cpro,arc=2.5mm,boxrule=0.9pt,left=3mm]{\Large\bfseries\color{cpro}b)}\quad\begin{minipage}[c]{5.3cm}\raggedright\Large $0{,}95$\end{minipage}\end{tcolorbox}\par\vspace{2.5mm}
\begin{tcolorbox}[enhanced,nobeforeafter,width=7.1cm,height=2.1cm,valign=center,colback=cpro!6,colframe=cpro,arc=2.5mm,boxrule=0.9pt,left=3mm]{\Large\bfseries\color{cpro}c)}\quad\begin{minipage}[c]{5.3cm}\raggedright\Large $0{,}225$\end{minipage}\end{tcolorbox}\hfill\begin{tcolorbox}[enhanced,nobeforeafter,width=7.1cm,height=2.1cm,valign=center,colback=cpro!6,colframe=cpro,arc=2.5mm,boxrule=0.9pt,left=3mm]{\Large\bfseries\color{cpro}d)}\quad\begin{minipage}[c]{5.3cm}\raggedright\Large $0{,}65$\end{minipage}\end{tcolorbox}
\end{tcolorbox}
\newpage
\begin{tcolorbox}[enhanced,width=16cm,height=14cm,colback=white,colframe=cpro,boxrule=2pt,arc=0mm,
left=6mm,right=6mm,top=0mm,bottom=4mm,sharp corners,
overlay={\fill[cpro] (frame.north west) rectangle ([yshift=-1.25cm]frame.north east);
\node[anchor=west,text=white,font=\sffamily\bfseries\LARGE] at ([xshift=5mm,yshift=-0.62cm]frame.north west) {Question 83 D};
\node[anchor=east,text=white,font=\sffamily\itshape\large] at ([xshift=-5mm,yshift=-0.62cm]frame.north east) {Probabilités};}]
\vspace*{1.75cm}
\begin{minipage}[t][6.9cm][c]{\linewidth}\LARGE\centering
\begin{minipage}{0.95\linewidth}\raggedright\linespread{1.25}\selectfont On a $P(A)=0{,}4$, $P(B)=0{,}4$ et $P(A\cup B)=0{,}65$. Que vaut $P(A\cap B)$~?\end{minipage}
\end{minipage}

\vfill
\begin{tcolorbox}[enhanced,nobeforeafter,width=7.1cm,height=2.1cm,valign=center,colback=cpro!6,colframe=cpro,arc=2.5mm,boxrule=0.9pt,left=3mm]{\Large\bfseries\color{cpro}a)}\quad\begin{minipage}[c]{5.3cm}\raggedright\Large $0{,}16$\end{minipage}\end{tcolorbox}\hfill\begin{tcolorbox}[enhanced,nobeforeafter,width=7.1cm,height=2.1cm,valign=center,colback=cpro!6,colframe=cpro,arc=2.5mm,boxrule=0.9pt,left=3mm]{\Large\bfseries\color{cpro}b)}\quad\begin{minipage}[c]{5.3cm}\raggedright\Large $0{,}25$\end{minipage}\end{tcolorbox}\par\vspace{2.5mm}
\begin{tcolorbox}[enhanced,nobeforeafter,width=7.1cm,height=2.1cm,valign=center,colback=cpro!6,colframe=cpro,arc=2.5mm,boxrule=0.9pt,left=3mm]{\Large\bfseries\color{cpro}c)}\quad\begin{minipage}[c]{5.3cm}\raggedright\Large $0{,}15$\end{minipage}\end{tcolorbox}\hfill\begin{tcolorbox}[enhanced,nobeforeafter,width=7.1cm,height=2.1cm,valign=center,colback=cpro!6,colframe=cpro,arc=2.5mm,boxrule=0.9pt,left=3mm]{\Large\bfseries\color{cpro}d)}\quad\begin{minipage}[c]{5.3cm}\raggedright\Large $0{,}35$\end{minipage}\end{tcolorbox}
\end{tcolorbox}
\newpage
\begin{tcolorbox}[enhanced,width=16cm,height=14cm,colback=white,colframe=cpro,boxrule=2pt,arc=0mm,
left=6mm,right=6mm,top=0mm,bottom=4mm,sharp corners,
overlay={\fill[cpro] (frame.north west) rectangle ([yshift=-1.25cm]frame.north east);
\node[anchor=west,text=white,font=\sffamily\bfseries\LARGE] at ([xshift=5mm,yshift=-0.62cm]frame.north west) {Question 84 D};
\node[anchor=east,text=white,font=\sffamily\itshape\large] at ([xshift=-5mm,yshift=-0.62cm]frame.north east) {Probabilités};}]
\vspace*{1.75cm}
\begin{minipage}[t][6.9cm][c]{\linewidth}\LARGE\centering
\begin{minipage}{0.95\linewidth}\raggedright\linespread{1.25}\selectfont On a $P(A)=0{,}35$, $P(B)=0{,}45$ et $P(A\cup B)=0{,}7$. Que vaut $P(A\cap B)$~?\end{minipage}
\end{minipage}

\vfill
\begin{tcolorbox}[enhanced,nobeforeafter,width=7.1cm,height=2.1cm,valign=center,colback=cpro!6,colframe=cpro,arc=2.5mm,boxrule=0.9pt,left=3mm]{\Large\bfseries\color{cpro}a)}\quad\begin{minipage}[c]{5.3cm}\raggedright\Large $0{,}1$\end{minipage}\end{tcolorbox}\hfill\begin{tcolorbox}[enhanced,nobeforeafter,width=7.1cm,height=2.1cm,valign=center,colback=cpro!6,colframe=cpro,arc=2.5mm,boxrule=0.9pt,left=3mm]{\Large\bfseries\color{cpro}b)}\quad\begin{minipage}[c]{5.3cm}\raggedright\Large $0{,}1575$\end{minipage}\end{tcolorbox}\par\vspace{2.5mm}
\begin{tcolorbox}[enhanced,nobeforeafter,width=7.1cm,height=2.1cm,valign=center,colback=cpro!6,colframe=cpro,arc=2.5mm,boxrule=0.9pt,left=3mm]{\Large\bfseries\color{cpro}c)}\quad\begin{minipage}[c]{5.3cm}\raggedright\Large $0{,}25$\end{minipage}\end{tcolorbox}\hfill\begin{tcolorbox}[enhanced,nobeforeafter,width=7.1cm,height=2.1cm,valign=center,colback=cpro!6,colframe=cpro,arc=2.5mm,boxrule=0.9pt,left=3mm]{\Large\bfseries\color{cpro}d)}\quad\begin{minipage}[c]{5.3cm}\raggedright\Large $0{,}35$\end{minipage}\end{tcolorbox}
\end{tcolorbox}
\newpage
\begin{tcolorbox}[enhanced,width=16cm,height=14cm,colback=white,colframe=cpro,boxrule=2pt,arc=0mm,
left=6mm,right=6mm,top=0mm,bottom=4mm,sharp corners,
overlay={\fill[cpro] (frame.north west) rectangle ([yshift=-1.25cm]frame.north east);
\node[anchor=west,text=white,font=\sffamily\bfseries\LARGE] at ([xshift=5mm,yshift=-0.62cm]frame.north west) {Question 85 D};
\node[anchor=east,text=white,font=\sffamily\itshape\large] at ([xshift=-5mm,yshift=-0.62cm]frame.north east) {Probabilités};}]
\vspace*{1.75cm}
\begin{minipage}[t][6.9cm][c]{\linewidth}\LARGE\centering
\begin{minipage}{0.95\linewidth}\raggedright\linespread{1.25}\selectfont On a $P(A)=\dfrac{9}{20}$. Que vaut $P(\overline{A})$, la probabilité de l'événement contraire~?\end{minipage}
\end{minipage}

\vfill
\begin{tcolorbox}[enhanced,nobeforeafter,width=7.1cm,height=2.1cm,valign=center,colback=cpro!6,colframe=cpro,arc=2.5mm,boxrule=0.9pt,left=3mm]{\Large\bfseries\color{cpro}a)}\quad\begin{minipage}[c]{5.3cm}\raggedright\Large $\dfrac{11}{20}$\end{minipage}\end{tcolorbox}\hfill\begin{tcolorbox}[enhanced,nobeforeafter,width=7.1cm,height=2.1cm,valign=center,colback=cpro!6,colframe=cpro,arc=2.5mm,boxrule=0.9pt,left=3mm]{\Large\bfseries\color{cpro}b)}\quad\begin{minipage}[c]{5.3cm}\raggedright\Large $\dfrac{9}{20}$\end{minipage}\end{tcolorbox}\par\vspace{2.5mm}
\begin{tcolorbox}[enhanced,nobeforeafter,width=7.1cm,height=2.1cm,valign=center,colback=cpro!6,colframe=cpro,arc=2.5mm,boxrule=0.9pt,left=3mm]{\Large\bfseries\color{cpro}c)}\quad\begin{minipage}[c]{5.3cm}\raggedright\Large $\dfrac{20}{9}$\end{minipage}\end{tcolorbox}\hfill\begin{tcolorbox}[enhanced,nobeforeafter,width=7.1cm,height=2.1cm,valign=center,colback=cpro!6,colframe=cpro,arc=2.5mm,boxrule=0.9pt,left=3mm]{\Large\bfseries\color{cpro}d)}\quad\begin{minipage}[c]{5.3cm}\raggedright\Large $\dfrac{11}{40}$\end{minipage}\end{tcolorbox}
\end{tcolorbox}
\newpage
\begin{tcolorbox}[enhanced,width=16cm,height=14cm,colback=white,colframe=cpro,boxrule=2pt,arc=0mm,
left=6mm,right=6mm,top=0mm,bottom=4mm,sharp corners,
overlay={\fill[cpro] (frame.north west) rectangle ([yshift=-1.25cm]frame.north east);
\node[anchor=west,text=white,font=\sffamily\bfseries\LARGE] at ([xshift=5mm,yshift=-0.62cm]frame.north west) {Question 86 D};
\node[anchor=east,text=white,font=\sffamily\itshape\large] at ([xshift=-5mm,yshift=-0.62cm]frame.north east) {Probabilités};}]
\vspace*{1.75cm}
\begin{minipage}[t][6.9cm][c]{\linewidth}\LARGE\centering
\begin{minipage}{0.95\linewidth}\raggedright\linespread{1.25}\selectfont On a $P(A)=\dfrac{13}{16}$. Que vaut $P(\overline{A})$, la probabilité de l'événement contraire~?\end{minipage}
\end{minipage}

\vfill
\begin{tcolorbox}[enhanced,nobeforeafter,width=7.1cm,height=2.1cm,valign=center,colback=cpro!6,colframe=cpro,arc=2.5mm,boxrule=0.9pt,left=3mm]{\Large\bfseries\color{cpro}a)}\quad\begin{minipage}[c]{5.3cm}\raggedright\Large $\dfrac{3}{32}$\end{minipage}\end{tcolorbox}\hfill\begin{tcolorbox}[enhanced,nobeforeafter,width=7.1cm,height=2.1cm,valign=center,colback=cpro!6,colframe=cpro,arc=2.5mm,boxrule=0.9pt,left=3mm]{\Large\bfseries\color{cpro}b)}\quad\begin{minipage}[c]{5.3cm}\raggedright\Large $\dfrac{16}{13}$\end{minipage}\end{tcolorbox}\par\vspace{2.5mm}
\begin{tcolorbox}[enhanced,nobeforeafter,width=7.1cm,height=2.1cm,valign=center,colback=cpro!6,colframe=cpro,arc=2.5mm,boxrule=0.9pt,left=3mm]{\Large\bfseries\color{cpro}c)}\quad\begin{minipage}[c]{5.3cm}\raggedright\Large $\dfrac{3}{16}$\end{minipage}\end{tcolorbox}\hfill\begin{tcolorbox}[enhanced,nobeforeafter,width=7.1cm,height=2.1cm,valign=center,colback=cpro!6,colframe=cpro,arc=2.5mm,boxrule=0.9pt,left=3mm]{\Large\bfseries\color{cpro}d)}\quad\begin{minipage}[c]{5.3cm}\raggedright\Large $\dfrac{13}{16}$\end{minipage}\end{tcolorbox}
\end{tcolorbox}
\newpage
\begin{tcolorbox}[enhanced,width=16cm,height=14cm,colback=white,colframe=cpro,boxrule=2pt,arc=0mm,
left=6mm,right=6mm,top=0mm,bottom=4mm,sharp corners,
overlay={\fill[cpro] (frame.north west) rectangle ([yshift=-1.25cm]frame.north east);
\node[anchor=west,text=white,font=\sffamily\bfseries\LARGE] at ([xshift=5mm,yshift=-0.62cm]frame.north west) {Question 87 D};
\node[anchor=east,text=white,font=\sffamily\itshape\large] at ([xshift=-5mm,yshift=-0.62cm]frame.north east) {Probabilités};}]
\vspace*{1.75cm}
\begin{minipage}[t][6.9cm][c]{\linewidth}\LARGE\centering
\begin{minipage}{0.95\linewidth}\raggedright\linespread{1.25}\selectfont On lance un dé équilibré à six faces. Quelle est la probabilité d'obtenir un nombre ni pair ni multiple de 3~?\end{minipage}
\end{minipage}

\vfill
\begin{tcolorbox}[enhanced,nobeforeafter,width=7.1cm,height=2.1cm,valign=center,colback=cpro!6,colframe=cpro,arc=2.5mm,boxrule=0.9pt,left=3mm]{\Large\bfseries\color{cpro}a)}\quad\begin{minipage}[c]{5.3cm}\raggedright\Large $\dfrac{1}{6}$\end{minipage}\end{tcolorbox}\hfill\begin{tcolorbox}[enhanced,nobeforeafter,width=7.1cm,height=2.1cm,valign=center,colback=cpro!6,colframe=cpro,arc=2.5mm,boxrule=0.9pt,left=3mm]{\Large\bfseries\color{cpro}b)}\quad\begin{minipage}[c]{5.3cm}\raggedright\Large $\dfrac{1}{2}$\end{minipage}\end{tcolorbox}\par\vspace{2.5mm}
\begin{tcolorbox}[enhanced,nobeforeafter,width=7.1cm,height=2.1cm,valign=center,colback=cpro!6,colframe=cpro,arc=2.5mm,boxrule=0.9pt,left=3mm]{\Large\bfseries\color{cpro}c)}\quad\begin{minipage}[c]{5.3cm}\raggedright\Large $\dfrac{2}{3}$\end{minipage}\end{tcolorbox}\hfill\begin{tcolorbox}[enhanced,nobeforeafter,width=7.1cm,height=2.1cm,valign=center,colback=cpro!6,colframe=cpro,arc=2.5mm,boxrule=0.9pt,left=3mm]{\Large\bfseries\color{cpro}d)}\quad\begin{minipage}[c]{5.3cm}\raggedright\Large $\dfrac{1}{3}$\end{minipage}\end{tcolorbox}
\end{tcolorbox}
\newpage
\begin{tcolorbox}[enhanced,width=16cm,height=14cm,colback=white,colframe=cpro,boxrule=2pt,arc=0mm,
left=6mm,right=6mm,top=0mm,bottom=4mm,sharp corners,
overlay={\fill[cpro] (frame.north west) rectangle ([yshift=-1.25cm]frame.north east);
\node[anchor=west,text=white,font=\sffamily\bfseries\LARGE] at ([xshift=5mm,yshift=-0.62cm]frame.north west) {Question 88 D};
\node[anchor=east,text=white,font=\sffamily\itshape\large] at ([xshift=-5mm,yshift=-0.62cm]frame.north east) {Probabilités};}]
\vspace*{1.75cm}
\begin{minipage}[t][6.9cm][c]{\linewidth}\LARGE\centering
\begin{minipage}{0.95\linewidth}\raggedright\linespread{1.25}\selectfont On lance un dé équilibré à six faces. Quelle est la probabilité d'obtenir un nombre multiple de 3 ou égal à 1~?\end{minipage}
\end{minipage}

\vfill
\begin{tcolorbox}[enhanced,nobeforeafter,width=7.1cm,height=2.1cm,valign=center,colback=cpro!6,colframe=cpro,arc=2.5mm,boxrule=0.9pt,left=3mm]{\Large\bfseries\color{cpro}a)}\quad\begin{minipage}[c]{5.3cm}\raggedright\Large $\dfrac{3}{5}$\end{minipage}\end{tcolorbox}\hfill\begin{tcolorbox}[enhanced,nobeforeafter,width=7.1cm,height=2.1cm,valign=center,colback=cpro!6,colframe=cpro,arc=2.5mm,boxrule=0.9pt,left=3mm]{\Large\bfseries\color{cpro}b)}\quad\begin{minipage}[c]{5.3cm}\raggedright\Large $\dfrac{1}{3}$\end{minipage}\end{tcolorbox}\par\vspace{2.5mm}
\begin{tcolorbox}[enhanced,nobeforeafter,width=7.1cm,height=2.1cm,valign=center,colback=cpro!6,colframe=cpro,arc=2.5mm,boxrule=0.9pt,left=3mm]{\Large\bfseries\color{cpro}c)}\quad\begin{minipage}[c]{5.3cm}\raggedright\Large $\dfrac{1}{2}$\end{minipage}\end{tcolorbox}\hfill\begin{tcolorbox}[enhanced,nobeforeafter,width=7.1cm,height=2.1cm,valign=center,colback=cpro!6,colframe=cpro,arc=2.5mm,boxrule=0.9pt,left=3mm]{\Large\bfseries\color{cpro}d)}\quad\begin{minipage}[c]{5.3cm}\raggedright\Large $\dfrac{2}{3}$\end{minipage}\end{tcolorbox}
\end{tcolorbox}
\newpage
\begin{tcolorbox}[enhanced,width=16cm,height=14cm,colback=white,colframe=cpro,boxrule=2pt,arc=0mm,
left=6mm,right=6mm,top=0mm,bottom=4mm,sharp corners,
overlay={\fill[cpro] (frame.north west) rectangle ([yshift=-1.25cm]frame.north east);
\node[anchor=west,text=white,font=\sffamily\bfseries\LARGE] at ([xshift=5mm,yshift=-0.62cm]frame.north west) {Question 89 D};
\node[anchor=east,text=white,font=\sffamily\itshape\large] at ([xshift=-5mm,yshift=-0.62cm]frame.north east) {Probabilités};}]
\vspace*{1.75cm}
\begin{minipage}[t][6.9cm][c]{\linewidth}\LARGE\centering
\begin{minipage}{0.95\linewidth}\raggedright\linespread{1.25}\selectfont On lance $5$ pièces équilibrées. Quelle est la probabilité de n'obtenir que des PILE~?\end{minipage}
\end{minipage}

\vfill
\begin{tcolorbox}[enhanced,nobeforeafter,width=7.1cm,height=2.1cm,valign=center,colback=cpro!6,colframe=cpro,arc=2.5mm,boxrule=0.9pt,left=3mm]{\Large\bfseries\color{cpro}a)}\quad\begin{minipage}[c]{5.3cm}\raggedright\Large $\dfrac{1}{2}$\end{minipage}\end{tcolorbox}\hfill\begin{tcolorbox}[enhanced,nobeforeafter,width=7.1cm,height=2.1cm,valign=center,colback=cpro!6,colframe=cpro,arc=2.5mm,boxrule=0.9pt,left=3mm]{\Large\bfseries\color{cpro}b)}\quad\begin{minipage}[c]{5.3cm}\raggedright\Large $\dfrac{31}{32}$\end{minipage}\end{tcolorbox}\par\vspace{2.5mm}
\begin{tcolorbox}[enhanced,nobeforeafter,width=7.1cm,height=2.1cm,valign=center,colback=cpro!6,colframe=cpro,arc=2.5mm,boxrule=0.9pt,left=3mm]{\Large\bfseries\color{cpro}c)}\quad\begin{minipage}[c]{5.3cm}\raggedright\Large $\dfrac{1}{32}$\end{minipage}\end{tcolorbox}\hfill\begin{tcolorbox}[enhanced,nobeforeafter,width=7.1cm,height=2.1cm,valign=center,colback=cpro!6,colframe=cpro,arc=2.5mm,boxrule=0.9pt,left=3mm]{\Large\bfseries\color{cpro}d)}\quad\begin{minipage}[c]{5.3cm}\raggedright\Large $\dfrac{1}{10}$\end{minipage}\end{tcolorbox}
\end{tcolorbox}
\newpage
\begin{tcolorbox}[enhanced,width=16cm,height=14cm,colback=white,colframe=cpro,boxrule=2pt,arc=0mm,
left=6mm,right=6mm,top=0mm,bottom=4mm,sharp corners,
overlay={\fill[cpro] (frame.north west) rectangle ([yshift=-1.25cm]frame.north east);
\node[anchor=west,text=white,font=\sffamily\bfseries\LARGE] at ([xshift=5mm,yshift=-0.62cm]frame.north west) {Question 90 D};
\node[anchor=east,text=white,font=\sffamily\itshape\large] at ([xshift=-5mm,yshift=-0.62cm]frame.north east) {Probabilités};}]
\vspace*{1.75cm}
\begin{minipage}[t][6.9cm][c]{\linewidth}\LARGE\centering
\begin{minipage}{0.95\linewidth}\raggedright\linespread{1.25}\selectfont On lance $5$ pièces équilibrées. Quelle est la probabilité d'obtenir au moins une fois FACE~?\end{minipage}
\end{minipage}

\vfill
\begin{tcolorbox}[enhanced,nobeforeafter,width=7.1cm,height=2.1cm,valign=center,colback=cpro!6,colframe=cpro,arc=2.5mm,boxrule=0.9pt,left=3mm]{\Large\bfseries\color{cpro}a)}\quad\begin{minipage}[c]{5.3cm}\raggedright\Large $\dfrac{1}{32}$\end{minipage}\end{tcolorbox}\hfill\begin{tcolorbox}[enhanced,nobeforeafter,width=7.1cm,height=2.1cm,valign=center,colback=cpro!6,colframe=cpro,arc=2.5mm,boxrule=0.9pt,left=3mm]{\Large\bfseries\color{cpro}b)}\quad\begin{minipage}[c]{5.3cm}\raggedright\Large $\dfrac{31}{32}$\end{minipage}\end{tcolorbox}\par\vspace{2.5mm}
\begin{tcolorbox}[enhanced,nobeforeafter,width=7.1cm,height=2.1cm,valign=center,colback=cpro!6,colframe=cpro,arc=2.5mm,boxrule=0.9pt,left=3mm]{\Large\bfseries\color{cpro}c)}\quad\begin{minipage}[c]{5.3cm}\raggedright\Large $\dfrac{1}{2}$\end{minipage}\end{tcolorbox}\hfill\begin{tcolorbox}[enhanced,nobeforeafter,width=7.1cm,height=2.1cm,valign=center,colback=cpro!6,colframe=cpro,arc=2.5mm,boxrule=0.9pt,left=3mm]{\Large\bfseries\color{cpro}d)}\quad\begin{minipage}[c]{5.3cm}\raggedright\Large $\dfrac{1}{10}$\end{minipage}\end{tcolorbox}
\end{tcolorbox}
\newpage
\begin{tcolorbox}[enhanced,width=16cm,height=14cm,colback=white,colframe=cpro,boxrule=2pt,arc=0mm,
left=6mm,right=6mm,top=0mm,bottom=4mm,sharp corners,
overlay={\fill[cpro] (frame.north west) rectangle ([yshift=-1.25cm]frame.north east);
\node[anchor=west,text=white,font=\sffamily\bfseries\LARGE] at ([xshift=5mm,yshift=-0.62cm]frame.north west) {Question 91 D};
\node[anchor=east,text=white,font=\sffamily\itshape\large] at ([xshift=-5mm,yshift=-0.62cm]frame.north east) {Probabilités};}]
\vspace*{1.75cm}
\begin{minipage}[t][6.9cm][c]{\linewidth}\Large\centering
\begin{minipage}{0.95\linewidth}\raggedright\linespread{1.25}\selectfont Voici la loi de probabilité d'une expérience aléatoire. Quelle est la probabilité manquante~? $\begin{array}{|l|c|c|c|c|}\hline \text{Issue} & 1 & 2 & 3 & 4\\\hline \text{Probabilité} & 0{,}32 & 0{,}18 &~? & 0{,}4\\\hline\end{array}$\end{minipage}
\end{minipage}

\vfill
\begin{tcolorbox}[enhanced,nobeforeafter,width=7.1cm,height=2.1cm,valign=center,colback=cpro!6,colframe=cpro,arc=2.5mm,boxrule=0.9pt,left=3mm]{\Large\bfseries\color{cpro}a)}\quad\begin{minipage}[c]{5.3cm}\raggedright\Large $0{,}5$\end{minipage}\end{tcolorbox}\hfill\begin{tcolorbox}[enhanced,nobeforeafter,width=7.1cm,height=2.1cm,valign=center,colback=cpro!6,colframe=cpro,arc=2.5mm,boxrule=0.9pt,left=3mm]{\Large\bfseries\color{cpro}b)}\quad\begin{minipage}[c]{5.3cm}\raggedright\Large $0{,}1$\end{minipage}\end{tcolorbox}\par\vspace{2.5mm}
\begin{tcolorbox}[enhanced,nobeforeafter,width=7.1cm,height=2.1cm,valign=center,colback=cpro!6,colframe=cpro,arc=2.5mm,boxrule=0.9pt,left=3mm]{\Large\bfseries\color{cpro}c)}\quad\begin{minipage}[c]{5.3cm}\raggedright\Large $0{,}9$\end{minipage}\end{tcolorbox}\hfill\begin{tcolorbox}[enhanced,nobeforeafter,width=7.1cm,height=2.1cm,valign=center,colback=cpro!6,colframe=cpro,arc=2.5mm,boxrule=0.9pt,left=3mm]{\Large\bfseries\color{cpro}d)}\quad\begin{minipage}[c]{5.3cm}\raggedright\Large $0{,}25$\end{minipage}\end{tcolorbox}
\end{tcolorbox}
\newpage
\begin{tcolorbox}[enhanced,width=16cm,height=14cm,colback=white,colframe=cpro,boxrule=2pt,arc=0mm,
left=6mm,right=6mm,top=0mm,bottom=4mm,sharp corners,
overlay={\fill[cpro] (frame.north west) rectangle ([yshift=-1.25cm]frame.north east);
\node[anchor=west,text=white,font=\sffamily\bfseries\LARGE] at ([xshift=5mm,yshift=-0.62cm]frame.north west) {Question 92 D};
\node[anchor=east,text=white,font=\sffamily\itshape\large] at ([xshift=-5mm,yshift=-0.62cm]frame.north east) {Probabilités};}]
\vspace*{1.75cm}
\begin{minipage}[t][6.9cm][c]{\linewidth}\Large\centering
\begin{minipage}{0.95\linewidth}\raggedright\linespread{1.25}\selectfont Voici la loi de probabilité d'une expérience aléatoire. Quelle est la probabilité manquante~? $\begin{array}{|l|c|c|c|c|}\hline \text{Issue} & 1 & 2 & 3 & 4\\\hline \text{Probabilité} & 0{,}25 & 0{,}25 &~? & 0{,}3\\\hline\end{array}$\end{minipage}
\end{minipage}

\vfill
\begin{tcolorbox}[enhanced,nobeforeafter,width=7.1cm,height=2.1cm,valign=center,colback=cpro!6,colframe=cpro,arc=2.5mm,boxrule=0.9pt,left=3mm]{\Large\bfseries\color{cpro}a)}\quad\begin{minipage}[c]{5.3cm}\raggedright\Large $0{,}5$\end{minipage}\end{tcolorbox}\hfill\begin{tcolorbox}[enhanced,nobeforeafter,width=7.1cm,height=2.1cm,valign=center,colback=cpro!6,colframe=cpro,arc=2.5mm,boxrule=0.9pt,left=3mm]{\Large\bfseries\color{cpro}b)}\quad\begin{minipage}[c]{5.3cm}\raggedright\Large $0{,}2$\end{minipage}\end{tcolorbox}\par\vspace{2.5mm}
\begin{tcolorbox}[enhanced,nobeforeafter,width=7.1cm,height=2.1cm,valign=center,colback=cpro!6,colframe=cpro,arc=2.5mm,boxrule=0.9pt,left=3mm]{\Large\bfseries\color{cpro}c)}\quad\begin{minipage}[c]{5.3cm}\raggedright\Large $0{,}8$\end{minipage}\end{tcolorbox}\hfill\begin{tcolorbox}[enhanced,nobeforeafter,width=7.1cm,height=2.1cm,valign=center,colback=cpro!6,colframe=cpro,arc=2.5mm,boxrule=0.9pt,left=3mm]{\Large\bfseries\color{cpro}d)}\quad\begin{minipage}[c]{5.3cm}\raggedright\Large $0{,}25$\end{minipage}\end{tcolorbox}
\end{tcolorbox}
\newpage
\begin{tcolorbox}[enhanced,width=16cm,height=14cm,colback=white,colframe=cpro,boxrule=2pt,arc=0mm,
left=6mm,right=6mm,top=0mm,bottom=4mm,sharp corners,
overlay={\fill[cpro] (frame.north west) rectangle ([yshift=-1.25cm]frame.north east);
\node[anchor=west,text=white,font=\sffamily\bfseries\LARGE] at ([xshift=5mm,yshift=-0.62cm]frame.north west) {Question 93 D};
\node[anchor=east,text=white,font=\sffamily\itshape\large] at ([xshift=-5mm,yshift=-0.62cm]frame.north east) {Probabilités};}]
\vspace*{1.75cm}
\begin{minipage}[t][6.9cm][c]{\linewidth}\Large\centering
\begin{minipage}{0.95\linewidth}\raggedright\linespread{1.25}\selectfont Une urne contient $9$ boules dont $6$ rouges. On tire une boule, on la remet, puis on en tire une deuxième. Quelle est la probabilité de tirer deux boules rouges~?\end{minipage}
\end{minipage}

\vfill
\begin{tcolorbox}[enhanced,nobeforeafter,width=7.1cm,height=2.1cm,valign=center,colback=cpro!6,colframe=cpro,arc=2.5mm,boxrule=0.9pt,left=3mm]{\Large\bfseries\color{cpro}a)}\quad\begin{minipage}[c]{5.3cm}\raggedright\Large $\dfrac{4}{9}$\end{minipage}\end{tcolorbox}\hfill\begin{tcolorbox}[enhanced,nobeforeafter,width=7.1cm,height=2.1cm,valign=center,colback=cpro!6,colframe=cpro,arc=2.5mm,boxrule=0.9pt,left=3mm]{\Large\bfseries\color{cpro}b)}\quad\begin{minipage}[c]{5.3cm}\raggedright\Large $\dfrac{5}{12}$\end{minipage}\end{tcolorbox}\par\vspace{2.5mm}
\begin{tcolorbox}[enhanced,nobeforeafter,width=7.1cm,height=2.1cm,valign=center,colback=cpro!6,colframe=cpro,arc=2.5mm,boxrule=0.9pt,left=3mm]{\Large\bfseries\color{cpro}c)}\quad\begin{minipage}[c]{5.3cm}\raggedright\Large $\dfrac{2}{3}$\end{minipage}\end{tcolorbox}\hfill\begin{tcolorbox}[enhanced,nobeforeafter,width=7.1cm,height=2.1cm,valign=center,colback=cpro!6,colframe=cpro,arc=2.5mm,boxrule=0.9pt,left=3mm]{\Large\bfseries\color{cpro}d)}\quad\begin{minipage}[c]{5.3cm}\raggedright\Large $1$\end{minipage}\end{tcolorbox}
\end{tcolorbox}
\newpage
\begin{tcolorbox}[enhanced,width=16cm,height=14cm,colback=white,colframe=cpro,boxrule=2pt,arc=0mm,
left=6mm,right=6mm,top=0mm,bottom=4mm,sharp corners,
overlay={\fill[cpro] (frame.north west) rectangle ([yshift=-1.25cm]frame.north east);
\node[anchor=west,text=white,font=\sffamily\bfseries\LARGE] at ([xshift=5mm,yshift=-0.62cm]frame.north west) {Question 94 D};
\node[anchor=east,text=white,font=\sffamily\itshape\large] at ([xshift=-5mm,yshift=-0.62cm]frame.north east) {Probabilités};}]
\vspace*{1.75cm}
\begin{minipage}[t][6.9cm][c]{\linewidth}\Large\centering
\begin{minipage}{0.95\linewidth}\raggedright\linespread{1.25}\selectfont Une urne contient $10$ boules dont $7$ rouges. On tire une boule, on la remet, puis on en tire une deuxième. Quelle est la probabilité de tirer deux boules rouges~?\end{minipage}
\end{minipage}

\vfill
\begin{tcolorbox}[enhanced,nobeforeafter,width=7.1cm,height=2.1cm,valign=center,colback=cpro!6,colframe=cpro,arc=2.5mm,boxrule=0.9pt,left=3mm]{\Large\bfseries\color{cpro}a)}\quad\begin{minipage}[c]{5.3cm}\raggedright\Large $\dfrac{49}{100}$\end{minipage}\end{tcolorbox}\hfill\begin{tcolorbox}[enhanced,nobeforeafter,width=7.1cm,height=2.1cm,valign=center,colback=cpro!6,colframe=cpro,arc=2.5mm,boxrule=0.9pt,left=3mm]{\Large\bfseries\color{cpro}b)}\quad\begin{minipage}[c]{5.3cm}\raggedright\Large $\dfrac{7}{10}$\end{minipage}\end{tcolorbox}\par\vspace{2.5mm}
\begin{tcolorbox}[enhanced,nobeforeafter,width=7.1cm,height=2.1cm,valign=center,colback=cpro!6,colframe=cpro,arc=2.5mm,boxrule=0.9pt,left=3mm]{\Large\bfseries\color{cpro}c)}\quad\begin{minipage}[c]{5.3cm}\raggedright\Large $1$\end{minipage}\end{tcolorbox}\hfill\begin{tcolorbox}[enhanced,nobeforeafter,width=7.1cm,height=2.1cm,valign=center,colback=cpro!6,colframe=cpro,arc=2.5mm,boxrule=0.9pt,left=3mm]{\Large\bfseries\color{cpro}d)}\quad\begin{minipage}[c]{5.3cm}\raggedright\Large $\dfrac{7}{15}$\end{minipage}\end{tcolorbox}
\end{tcolorbox}
\newpage
\begin{tcolorbox}[enhanced,width=16cm,height=14cm,colback=white,colframe=cpro,boxrule=2pt,arc=0mm,
left=6mm,right=6mm,top=0mm,bottom=4mm,sharp corners,
overlay={\fill[cpro] (frame.north west) rectangle ([yshift=-1.25cm]frame.north east);
\node[anchor=west,text=white,font=\sffamily\bfseries\LARGE] at ([xshift=5mm,yshift=-0.62cm]frame.north west) {Question 95 D};
\node[anchor=east,text=white,font=\sffamily\itshape\large] at ([xshift=-5mm,yshift=-0.62cm]frame.north east) {Probabilités};}]
\vspace*{1.75cm}
\begin{minipage}[t][6.9cm][c]{\linewidth}\LARGE\centering
\begin{minipage}{0.95\linewidth}\raggedright\linespread{1.25}\selectfont $A$ et $B$ sont deux événements incompatibles, avec $P(A)=0{,}22$ et $P(B)=0{,}3$. Que vaut $P(A\cup B)$~?\end{minipage}
\end{minipage}

\vfill
\begin{tcolorbox}[enhanced,nobeforeafter,width=7.1cm,height=2.1cm,valign=center,colback=cpro!6,colframe=cpro,arc=2.5mm,boxrule=0.9pt,left=3mm]{\Large\bfseries\color{cpro}a)}\quad\begin{minipage}[c]{5.3cm}\raggedright\Large $0{,}52$\end{minipage}\end{tcolorbox}\hfill\begin{tcolorbox}[enhanced,nobeforeafter,width=7.1cm,height=2.1cm,valign=center,colback=cpro!6,colframe=cpro,arc=2.5mm,boxrule=0.9pt,left=3mm]{\Large\bfseries\color{cpro}b)}\quad\begin{minipage}[c]{5.3cm}\raggedright\Large $0{,}454$\end{minipage}\end{tcolorbox}\par\vspace{2.5mm}
\begin{tcolorbox}[enhanced,nobeforeafter,width=7.1cm,height=2.1cm,valign=center,colback=cpro!6,colframe=cpro,arc=2.5mm,boxrule=0.9pt,left=3mm]{\Large\bfseries\color{cpro}c)}\quad\begin{minipage}[c]{5.3cm}\raggedright\Large $0{,}08$\end{minipage}\end{tcolorbox}\hfill\begin{tcolorbox}[enhanced,nobeforeafter,width=7.1cm,height=2.1cm,valign=center,colback=cpro!6,colframe=cpro,arc=2.5mm,boxrule=0.9pt,left=3mm]{\Large\bfseries\color{cpro}d)}\quad\begin{minipage}[c]{5.3cm}\raggedright\Large $0{,}066$\end{minipage}\end{tcolorbox}
\end{tcolorbox}
\newpage
\begin{tcolorbox}[enhanced,width=16cm,height=14cm,colback=white,colframe=cpro,boxrule=2pt,arc=0mm,
left=6mm,right=6mm,top=0mm,bottom=4mm,sharp corners,
overlay={\fill[cpro] (frame.north west) rectangle ([yshift=-1.25cm]frame.north east);
\node[anchor=west,text=white,font=\sffamily\bfseries\LARGE] at ([xshift=5mm,yshift=-0.62cm]frame.north west) {Question 96 D};
\node[anchor=east,text=white,font=\sffamily\itshape\large] at ([xshift=-5mm,yshift=-0.62cm]frame.north east) {Probabilités};}]
\vspace*{1.75cm}
\begin{minipage}[t][6.9cm][c]{\linewidth}\LARGE\centering
\begin{minipage}{0.95\linewidth}\raggedright\linespread{1.25}\selectfont $A$ et $B$ sont deux événements incompatibles, avec $P(A)=\dfrac{2}{9}$ et $P(B)=\dfrac{1}{3}$. Que vaut $P(A\cup B)$~?\end{minipage}
\end{minipage}

\vfill
\begin{tcolorbox}[enhanced,nobeforeafter,width=7.1cm,height=2.1cm,valign=center,colback=cpro!6,colframe=cpro,arc=2.5mm,boxrule=0.9pt,left=3mm]{\Large\bfseries\color{cpro}a)}\quad\begin{minipage}[c]{5.3cm}\raggedright\Large $\dfrac{2}{27}$\end{minipage}\end{tcolorbox}\hfill\begin{tcolorbox}[enhanced,nobeforeafter,width=7.1cm,height=2.1cm,valign=center,colback=cpro!6,colframe=cpro,arc=2.5mm,boxrule=0.9pt,left=3mm]{\Large\bfseries\color{cpro}b)}\quad\begin{minipage}[c]{5.3cm}\raggedright\Large $\dfrac{5}{9}$\end{minipage}\end{tcolorbox}\par\vspace{2.5mm}
\begin{tcolorbox}[enhanced,nobeforeafter,width=7.1cm,height=2.1cm,valign=center,colback=cpro!6,colframe=cpro,arc=2.5mm,boxrule=0.9pt,left=3mm]{\Large\bfseries\color{cpro}c)}\quad\begin{minipage}[c]{5.3cm}\raggedright\Large $\dfrac{13}{27}$\end{minipage}\end{tcolorbox}\hfill\begin{tcolorbox}[enhanced,nobeforeafter,width=7.1cm,height=2.1cm,valign=center,colback=cpro!6,colframe=cpro,arc=2.5mm,boxrule=0.9pt,left=3mm]{\Large\bfseries\color{cpro}d)}\quad\begin{minipage}[c]{5.3cm}\raggedright\Large $\dfrac{1}{9}$\end{minipage}\end{tcolorbox}
\end{tcolorbox}
\newpage
\end{document}